\documentclass[11pt,a4paper,oneside]{report}
\usepackage[utf8]{inputenc}
\usepackage[T1]{fontenc}
\usepackage{lmodern}
\usepackage[polish]{babel}
\usepackage[margin=2.5cm]{geometry}
\usepackage{amsmath,amssymb,amsthm,mathtools}
\usepackage{array}
\usepackage{enumitem}
\usepackage{float}
\usepackage{tikz}
\usetikzlibrary{arrows.meta,calc,positioning,decorations.pathreplacing,decorations.markings,patterns}
\usepackage{pgfplots}
\pgfplotsset{compat=1.18}
\definecolor{manifoldblue}{RGB}{35,91,140}
\definecolor{manifoldred}{RGB}{170,65,55}
\definecolor{manifoldgreen}{RGB}{55,125,85}
\definecolor{manifoldgold}{RGB}{205,145,35}
\usepackage{microtype}
\usepackage{fancyhdr}
\usepackage[unicode]{hyperref}
\hypersetup{colorlinks=true, linkcolor=blue!45!black, citecolor=blue!45!black, urlcolor=blue!45!black}

\newcommand{\F}{\mathbb{F}}
\newcommand{\R}{\mathbb{R}}
\newcommand{\C}{\mathbb{C}}
\newcommand{\N}{\mathbb{N}}
\newcommand{\Z}{\mathbb{Z}}
\newcommand{\Q}{\mathbb{Q}}
\newcommand{\id}{\mathrm{id}}
\newcommand{\Lie}{\mathcal L}
\newcommand{\X}{\mathfrak X}
\newcommand{\supp}{\operatorname{supp}}
\newcommand{\pr}{\operatorname{pr}}
\newcommand{\D}{\mathrm D}
\newcommand{\dd}{\mathrm d}

\theoremstyle{plain}
\newtheorem{twierdzenie}{Twierdzenie}[section]
\newtheorem{lemat}[twierdzenie]{Lemat}
\newtheorem{stwierdzenie}[twierdzenie]{Stwierdzenie}
\newtheorem{wniosek}[twierdzenie]{Wniosek}
\newtheorem{fakt}[twierdzenie]{Fakt}
\theoremstyle{definition}
\newtheorem{definicja}[twierdzenie]{Definicja}
\newtheorem{przyklad}[twierdzenie]{Przykład}
\theoremstyle{remark}
\newtheorem{uwaga}[twierdzenie]{Uwaga}
\renewcommand{\proofname}{Dowód}

\pagestyle{fancy}
\fancyhf{}
\fancyhead[L]{\small Topologia algebraiczna}
\renewcommand{\chaptermark}[1]{\markboth{Rozdział \thechapter}{}}
\fancyhead[R]{\small\nouppercase{\leftmark}}
\fancyfoot[C]{\thepage}
\renewcommand{\headrulewidth}{0.4pt}

\setlength{\parindent}{0pt}
\setlength{\parskip}{4pt}
\setlength{\emergencystretch}{2em}

\begin{document}
\hypersetup{pageanchor=false}
\begin{titlepage}
\centering
\vspace*{4cm}
{\Huge\bfseries Topologia algebraiczna\par}
\vspace{0.6cm}
{\Large Homotopia, homologia, kohomologia i niezmienniki rozmaitości\par}
\vspace{1.3cm}
{\large Skrypt zebrany z rozdziałów książki o geometrii i topologii różniczkowej\par}
\vspace{1.0cm}
{\large Wersja XIX: po globalnym audycie matematycznym i redakcyjnym, 18 listopada 2026 r.\par}
\vfill
{\normalsize Wersja samodzielna; zadania w osobnym pliku.\par}
\end{titlepage}
\hypersetup{pageanchor=true}
\begingroup
\footnotesize\linespread{0.92}\selectfont\setlength{\parskip}{0pt}
\tableofcontents
\endgroup
\clearpage
\chapter*{Wprowadzenie i zakres}
\addcontentsline{toc}{chapter}{Wprowadzenie i zakres}
Niniejszy skrypt wyodrębnia z wcześniejszej książki część poświęconą
metodom topologii algebraicznej. Zachowano zasadniczą treść rozdziałów
źródłowych: twierdzenia, dowody, przykłady i rysunki.
Dodano osobne wprowadzenie topologiczne, samodzielny rozdział
z algebry abstrakcyjnej, kohomologię singularną,
fibracje, teorię \v{C}echa, ciągi spektralne, przestrzenie klasyfikujące, teorię przeszkód, wieże Postnikowa, operacje kohomologiczne, homotopię stabilną, konstrukcję Pontriagina--Thoma, homologię Morse'a, twierdzenie Lefschetza, Poincarégo--Hopfa, homologię Borel--Moore'a, kohomologię snopów, dualność Verdiera, formalizm sześciu operacji i cykle pobliskie oraz topologię odmian projektowych, twierdzenia Lefschetza i rozdmuchania. Rozdział 29 uzupełnia granice odwrotne, teorię przeszkód, ekscyzję homotopijną i analizę Morse'a.
Część poświęcona dualności i klasom charakterystycznym pochodziła z
późniejszej części skryptu geometrycznego.

Punktem wyjścia są podstawy topologii i algebry liniowej;
potrzebną teorię grup, pierścieni i modułów przedstawiamy od początku.
W rozdziale o formach i w rozdziale o klasach charakterystycznych
korzystamy również z elementarnej teorii gładkich rozmaitości,
przestrzeni stycznych, wiązek wektorowych i transwersalności;
przypomnienie potrzebnych pojęć znajduje się w dodatku.
Pełne dowody twierdzeń o otoczeniach tubularnych, transwersalności
oraz teorii spinowej pozostają w książce o geometrii różniczkowej.

\noindent\textbf{Status dowodów.} Wyrażenie ``dowód oznacza miejscami
dowód względny: po przyjęciu wyraźnie wskazanego twierdzenia
zewnętrznego, np. o transwersalności, rozkładzie lub dualności Verdiera.
Globalny audyt XIX identyfikuje te zależności w osobnym raporcie.
Nie należy utożsamiać szkiców i rachunków wynikających z głębokiego
twierdzenia z niezależnym dowodem całego tego twierdzenia.

\noindent\textbf{Uwaga redakcyjna.} Ta wersja scala istniejący materiał;
nie oznacza pełnej przebudowy dowodów wszystkich dużych twierdzeń.
Osobny dokument z zadaniami zawiera wskazówki do trudniejszych ćwiczeń.
\clearpage
\chapter{Język i podstawowe konstrukcje topologiczne}
\label{ch:wstep-topologia}

Topologia algebraiczna bada przestrzenie za pomocą niezmienników, które można
obliczać algebraicznie: grup podstawowych, grup homologii i pierścieni
kohomologii. Nie mierzymy odległości; pytamy, jak przestrzeń jest sklejona.
Celem tego rozdziału jest zebranie podstawowego języka potrzebnego w dalszej
części książki. Zakładamy elementarną znajomość zbiorów, odwzorowań i algebry
liniowej. Wszystkie konstrukcje poniżej dotyczą przestrzeni topologicznych.

\section{Przestrzenie topologiczne, ciągłość i przykłady}

\begin{definicja}
Topologią na zbiorze $X$ nazywamy rodzinę $\tau\subseteq\mathcal P(X)$
zawierającą $\varnothing$ i $X$, zamkniętą na dowolne sumy oraz skończone
przecięcia. Parę $(X,\tau)$ nazywamy przestrzenią topologiczną; elementy
$\tau$ nazywamy zbiorami otwartymi. Odwzorowanie $f:X\to Y$ jest ciągłe,
jeśli przeciwobraz każdego otwartego podzbioru $Y$ jest otwarty w $X$.
Homeomorfizmem nazywamy ciągłą bijekcję o ciągłej odwrotności.
\end{definicja}

\begin{przyklad}
Podstawowe modele to $\R^n$ z topologią euklidesową, sfera
$S^n=\{x\in\R^{n+1}:\|x\|=1\}$ i dysk $D^n=\{x\in\R^n:\|x\|\leq1\}$.
Nie utożsamiamy $S^0$ z okręgiem: $S^0$ składa się z dwóch punktów.
Punkty $S^1$ można zapisać jako $e^{2\pi it}$ dla $t\in\R$.
\end{przyklad}

\begin{definicja}
Podprzestrzeń $A\subseteq X$ wyposażamy w topologię $\{A\cap U:U\in\tau_X\}$.
Na iloczynie $X\times Y$ przyjmujemy topologię generowaną przez prostokąty
$U\times V$ z $U\subseteq X$, $V\subseteq Y$ otwartymi. Przestrzeń $X$
nazywamy spójną, gdy nie jest sumą dwóch rozłącznych, niepustych zbiorów
otwartych; nazywamy ją łukowo spójną, gdy każde dwa punkty można połączyć
ciągłą drogą $[0,1]\to X$.
\end{definicja}

\begin{stwierdzenie}
Niech $X$ będzie łukowo spójna. Wówczas $X$ jest spójna.
\end{stwierdzenie}
\begin{proof}
Przypuśćmy, że $X=U\sqcup V$, gdzie $U,V$ są niepuste i otwarte.
Wybierzmy $u\in U$ i $v\in V$ oraz drogę $\gamma$ od $u$ do $v$.
Zbiory $\gamma^{-1}(U)$ oraz $\gamma^{-1}(V)$ byłyby rozłącznymi,
niepustymi zbiorami otwartymi o sumie $[0,1]$. Przeczyłoby to spójności
odcinka $[0,1]$.
\end{proof}

\section{Topologia ilorazowa i przestrzenie sklejone}

\begin{definicja}[Topologia ilorazowa]
\label{def:topologia-ilorazowa}
Niech $\sim$ będzie relacją równoważności na $X$ i
$q:X\to X/{\sim}$ odwzorowaniem $x\mapsto[x]$.
Zbiór $U\subseteq X/{\sim}$ nazywamy otwartym wtedy i tylko wtedy, gdy
$q^{-1}(U)$ jest otwarty w $X$. Tak otrzymaną topologię nazywamy
\emph{topologią ilorazową}.
\end{definicja}

\begin{twierdzenie}[Własność uniwersalna ilorazu]
Niech $q:X\to X/{\sim}$ będzie odwzorowaniem ilorazowym i niech $Y$ będzie
przestrzenią topologiczną. Odwzorowanie $F:X/{\sim}\to Y$ jest ciągłe wtedy
i tylko wtedy, gdy $F\circ q:X\to Y$ jest ciągłe.
W szczególności każde ciągłe $f:X\to Y$, stałe na klasach relacji $\sim$,
wyznacza jedyne ciągłe $\bar f:X/{\sim}\to Y$ takie, że $f=\bar f\circ q$.
\end{twierdzenie}
\begin{proof}
Gdy $F$ jest ciągłe, ciągłe jest także $F\circ q$. W drugą stronę,
dla każdego otwartego $V\subseteq Y$ mamy
$q^{-1}(F^{-1}(V))=(F\circ q)^{-1}(V)$, co jest otwarte.
Z definicji topologii ilorazowej wynika otwartość $F^{-1}(V)$.
Jeśli $f$ jest stałe na klasach, wzór $\bar f([x])=f(x)$ nie zależy od wyboru
przedstawiciela; ciągłość wynika z pierwszej części dowodu.
\end{proof}

\begin{przyklad}[Okrąg i torus]
Iloraz $[0,1]/(0\sim1)$ jest homeomorficzny z $S^1$ przez
$[t]\mapsto e^{2\pi it}$; ciągła bijekcja jest homeomorfizmem, ponieważ
jej dziedzina jest zwarta, a przeciwdziedzina Hausdorffa.
Podobnie $\R^2/\Z^2$, gdzie $(x,y)\sim(x+m,y+n)$ dla $m,n\in\Z$,
jest torusem $T^2=S^1\times S^1$.
\end{przyklad}

\begin{definicja}
Dla pary $(X,A)$ z $A\subseteq X$ przestrzeń $X/A$ powstaje przez
utożsamienie wszystkich punktów $A$ w jeden punkt. Dla przestrzeni z
punktami bazowymi $(X,x_0)$ i $(Y,y_0)$ \emph{suma klinowa}
$X\vee Y$ jest ilorazem $X\sqcup Y$ przez relację $x_0\sim y_0$.
Stożek, zawieszenie i produkt smash dla dobrze bazowanych
przestrzeni zapisujemy odpowiednio jako
$CX=X\times[0,1]/(X\times\{1\})$,
$\Sigma X=(X\times[0,1])/(X\times\{0\},X\times\{1\})$,
oraz $X\wedge Y=(X\times Y)/(X\vee Y)$,
gdzie $X\vee Y$ utożsamiamy z
$(X\times\{y_0\})\cup(\{x_0\}\times Y)$.
W zawieszeniu obie końcowe kopie $X$ są sklejane do \emph{dwóch różnych}
punktów.
\end{definicja}

\begin{uwaga}
Powyższy stożek $CX$ jest stożkiem niezredukowanym (jedna sklejona
podstawa). Termin „zredukowany” jest zarezerwowany dla dodatkowego
utożsamienia $\{x_0\}\times[0,1]$; w dalszej części nie będziemy
mylić tych dwóch konstrukcji.
\end{uwaga}

\begin{przyklad}
Zawieszenie $S^n$ jest homeomorficzne z $S^{n+1}$, a zawieszenie $S^0$
z dwiema rozłącznymi kopiami punktów daje okrąg $S^1$.
\end{przyklad}

\section{Homotopie i deformacyjne retrakcje}

\begin{definicja}
Dwa ciągłe odwzorowania $f,g:X\to Y$ są homotopijne, piszemy
$f\simeq g$, jeżeli istnieje ciągłe $H:X\times[0,1]\to Y$ takie, że
$H(x,0)=f(x)$ i $H(x,1)=g(x)$ dla każdego $x\in X$.
Podzbiór $A\subseteq X$ jest \emph{deformacyjnym retraktem} $X$, jeśli
istnieje homotopia $H:\id_X\simeq i\circ r$, gdzie $r:X\to A$ jest
retrakcją, $i:A\hookrightarrow X$ jest inkluzją oraz
$H(a,t)=a$ dla wszystkich $a\in A$, $t\in[0,1]$.
Przestrzeń jest ściągalna, gdy jej identyczność jest homotopijna z
odwzorowaniem stałym.
\end{definicja}

\begin{przyklad}[Dziurawa płaszczyzna]
Płaszczyzna $\R^2\setminus\{0\}$ ma deformacyjny retrakt $S^1$.
Dla $x\ne0$ kładziemy
$H(x,t)=\big((1-t)+t/\|x\|\big)x$.
Współczynnik jest dodatni, więc homotopia nie przechodzi przez zero;
$H(x,0)=x$, $H(x,1)=x/\|x\|$ i $H(a,t)=a$ dla $a\in S^1$.
Ten rachunek tłumaczy, dlaczego oba zbiory mają takie same grupy homotopii
i homologii, mimo że nie są homeomorficzne.
\end{przyklad}

\begin{twierdzenie}
Jeżeli $A$ jest deformacyjnym retraktem $X$, to inkluzja $i:A\to X$
jest równoważnością homotopijną, czyli ma homotopijną odwrotność.
\end{twierdzenie}
\begin{proof}
Retrakcja $r:X\to A$ spełnia $r\circ i=\id_A$.
Z definicji deformacyjnej retrakcji mamy $i\circ r\simeq\id_X$.
To właśnie dwa warunki definiujące odwrotność homotopijną.
\end{proof}

\section{Przestrzenie rzutowe i modele do obliczeń}

\begin{definicja}
Rzeczywista przestrzeń rzutowa $\R P^n$ to iloraz $S^n$ przez
utożsamienie punktów antypodycznych $x\sim-x$.
Zespolona przestrzeń rzutowa $\C P^n$ to zbiór zespolonych prostych
w $\C^{n+1}$ z topologią ilorazową pochodzącą z
$\C^{n+1}\setminus\{0\}$ i relacji $z\sim\lambda z$ dla $\lambda\in\C^*$.
\end{definicja}

\begin{uwaga}
Przestrzenie rzutowe i torusy będą naszymi laboratoriami: najpierw
rozpoznamy ich rozkłady komórkowe, następnie policzymy homologię oraz
pierścienie kohomologii. Ważne będzie odróżnienie niezmienników
\emph{grupowych} od dodatkowych struktur, takich jak iloczyn kubkowy.
\end{uwaga}

% Rozdzial przygotowany do samodzielnego skryptu topologii algebraicznej.
\chapter{Podstawy algebry abstrakcyjnej: grupy, pierścienie i moduły}
\label{ch:algebra-abstrakcyjna-v5}

W topologii algebraicznej każdej przestrzeni przypisujemy grupy, a odwzorowaniom
ciągłym --- homomorfizmy. Grupa podstawowa może być nieprzemienna, grupy
homologii są abelowe, a iloczyn tensorowy i funktory $\operatorname{Tor}$,
$\operatorname{Ext}$ opisują oddziaływanie współczynników. Dlatego zbieramy
potrzebne konstrukcje \emph{przed} rozdziałem o homologii algebraicznej.
Żadne twierdzenie niniejszego rozdziału nie odsyła do osobnego rozdziału
książki geometrycznej o algebrze tensorowej. Pracujemy nad pierścieniami
z jedynką, homomorfizmy pierścieni zachowują jedynkę, a moduły są lewostronne,
jeżeli nie wskazano inaczej.

\section{Grupy, homomorfizmy, podgrupy normalne}
\begin{definicja}
\emph{Grupa} to zbiór $G$ z działaniem $(g,h)\mapsto gh$ łącznym,
elementem neutralnym $1_G$ i odwrotnościami $g^{-1}$ wszystkich elementów.
Jest \emph{abelowa}, gdy $gh=hg$ dla wszystkich $g,h\in G$.
\emph{Homomorfizm} $\varphi:G\to H$ spełnia
$\varphi(gh)=\varphi(g)\varphi(h)$. Podgrupa $N\leq G$ jest
\emph{normalna}, $N\trianglelefteq G$, gdy $gNg^{-1}=N$ dla wszystkich
$g\in G$. Jądro homomorfizmu to $\ker\varphi=\varphi^{-1}(\{1_H\})$.
\end{definicja}
\begin{przyklad}
$(\Z,+)$, $(\Z/m\Z,+)$, grupa permutacji $S_m$ i grupa macierzy
$\mathrm{GL}_m(\R)$ są grupami. Pierwsze dwie są abelowe, dwie następne
zazwyczaj nie. Homomorfizm $\Z\to \Z/m\Z$ wysyłający $k$ na jego
resztę ma jądro $m\Z$. Wyznacznik jest homomorfizmem
$\det:\mathrm{GL}_m(\R)\to\R^\times$ z jądrem
$\mathrm{SL}_m(\R)$.
\end{przyklad}
\begin{stwierdzenie}
Dla każdego homomorfizmu grup $\varphi:G\to H$ jego jądro jest
normalną podgrupą $G$, a obraz jest podgrupą $H$.
\end{stwierdzenie}
\begin{proof}
Jeżeli $n\in\ker\varphi$ i $g\in G$, to
$\varphi(gng^{-1})=\varphi(g)1_H\varphi(g)^{-1}=1_H$.
Pozostałe aksjomaty wynikają z zachowywania iloczynu, jedynki
i odwrotności przez $\varphi$.
\end{proof}

\begin{definicja}
Dla $N\trianglelefteq G$ na zbiorze klas lewostronnych $G/N$ definiujemy
$(gN)(hN)=ghN$. Tak otrzymana grupa nazywa się \emph{ilorazową}.
\end{definicja}
\begin{lemat}
Iloczyn na $G/N$ jest dobrze określony wtedy i tylko wtedy, gdy $N$
jest normalna. Naturalny rzut $q:G\to G/N$ jest homomorfizmem.
\end{lemat}
\begin{proof}
Przy normalności $g'=gn$, $h'=hm$ implikuje
$g'h'=gh(h^{-1}nh)m$, zatem $g'h'N=ghN$. W drugą stronę,
niezależność iloczynu od reprezentantów daje
$(gN)(nN)(g^{-1}N)=(gN)(N)(g^{-1}N)=N$,
więc $gng^{-1}\in N$. Włączenie dla $g^{-1}$ daje równość.
\end{proof}
\begin{twierdzenie}[Pierwsze twierdzenie o izomorfizmie]
\label{thm:izo-grupy-v5}
Jeżeli $\varphi:G\to H$ jest homomorfizmem grup, to
$G/\ker\varphi\cong\operatorname{im}\varphi$, przy czym izomorfizm
wysyła $g\ker\varphi$ na $\varphi(g)$.
\end{twierdzenie}
\begin{proof}
Dla $\varphi(g)=\varphi(h)$ mamy $h^{-1}g\in\ker\varphi$,
a więc $g\ker\varphi=h\ker\varphi$. Wzór nie zależy zatem
od reprezentanta i jest różnowartościowy. Surjektywność na obraz
jest oczywista, a zgodność z mnożeniem wynika z własności $\varphi$.
\end{proof}
\begin{uwaga}
W topologii z tego twierdzenia otrzymamy np. abelianizację
$\pi_1(X)\twoheadrightarrow H_1(X;\Z)$. Dla niełukowo spójnych
przestrzeni twierdzenie stosuje się osobno do wyróżnionej składowej
łukowej.
\end{uwaga}

\section{Działania grup, orbity, stabilizatory i iloczyny półproste}
\begin{definicja}
\emph{Lewym działaniem} grupy $G$ na zbiorze $X$ nazywamy funkcję
$(g,x)\mapsto gx$ spełniającą $1_Gx=x$ i $(gh)x=g(hx)$.
\emph{Orbita} $x$ to $Gx=\{gx:g\in G\}$, a \emph{stabilizator}
$G_x=\{g\in G:gx=x\}$.
\end{definicja}
\begin{twierdzenie}[Orbita i stabilizator]
Dla działania $G\curvearrowright X$ i każdego $x\in X$ istnieje
bijekcja zbiorów $G/G_x\to Gx$, $gG_x\mapsto gx$.
Jeżeli $G$ jest skończona, to $|Gx|=[G:G_x]$.
\end{twierdzenie}
\begin{proof}
Równość $gx=hx$ jest równoważna $h^{-1}g\in G_x$,
czyli $gG_x=hG_x$. To daje poprawność i injektywność.
Każdy punkt orbity ma postać $gx$, zatem mapa jest na.
\end{proof}
\begin{przyklad}
Działanie $\Z$ na $\R$ przez $n\cdot t=t+n$ ma iloraz
$\R/\Z\cong S^1$; stabilizatory są trywialne.
Działanie $S_m$ na $\{1,\ldots,m\}$ jest przechodnie,
a stabilizator $1$ jest izomorficzny z $S_{m-1}$.
\end{przyklad}

\begin{definicja}[Iloczyn prosty i półprosty]
Dla grup $N,H$ iloczyn prosty $N\times H$ ma mnożenie
$(n,h)(n',h')=(nn',hh')$. Jeżeli
$\alpha:H\to\operatorname{Aut}(N)$ jest homomorfizmem,
to \emph{iloczyn półprosty} $N\rtimes_\alpha H$ jest zbiorem
$N\times H$ z mnożeniem
$(n,h)(n',h')=(n\alpha_h(n'),hh')$.
\end{definicja}
\begin{stwierdzenie}
Powyższe mnożenie jest grupowe; elementem neutralnym jest
$(1_N,1_H)$, a odwrotność $(n,h)$ wynosi
$(\alpha_{h^{-1}}(n^{-1}),h^{-1})$. Jeśli $\alpha$ jest
trywialne, otrzymujemy iloczyn prosty.
\end{stwierdzenie}
\begin{proof}
Łączność wynika z $\alpha_{hh'}=\alpha_h\alpha_{h'}$:
obie strony iloczynu trzech elementów mają pierwszą współrzędną
$n\alpha_h(n')\alpha_{hh'}(n'')$. Wstawienie podanego wzoru
na odwrotność daje jedynkę z obu stron.
\end{proof}
\begin{przyklad}[Grupa butelki Kleina]
Niech $H=\Z$, $N=\Z$ i $\alpha_k(n)=(-1)^k n$. Wtedy
$N\rtimes_\alpha H$ jest grupą o prezentacji
$\langle a,b\mid aba^{-1}=b^{-1}\rangle$.
To pojawi się przy obliczaniu grupy podstawowej butelki Kleina
z monodromii fibracji okręgowej.
\end{przyklad}

\section{Grupy wolne i prezentacje przez generatory i relacje}
\begin{definicja}
Niech $S$ będzie zbiorem. \emph{Grupa wolna} $F(S)$ składa się ze
zredukowanych skończonych słów w symbolach $s,s^{-1}$, $s\in S$,
z mnożeniem przez konkatenację i redukcję par $ss^{-1}$ i $s^{-1}s$.
\emph{Prezentacja} $\langle S\mid R\rangle$ oznacza iloraz $F(S)$
przez najmniejszą normalną podgrupę zawierającą słowa relacji $R$.
\end{definicja}
\begin{twierdzenie}[Własność uniwersalna grupy wolnej]
Dla dowolnej grupy $H$ i funkcji $u:S\to H$ istnieje dokładnie
jeden homomorfizm $\overline u:F(S)\to H$ przedłużający $u$.
Homomorfizm ten przechodzi na $\langle S\mid R\rangle$ wtedy
i tylko wtedy, gdy wszystkie relacje z $R$ mają obraz $1_H$.
\end{twierdzenie}
\begin{proof}
Definiujemy obraz słowa, mnożąc obrazy kolejnych liter,
a dla $s^{-1}$ biorąc $u(s)^{-1}$. Usunięcie sąsiednich
liter odwrotnych nie zmienia wartości. Każdy homomorfizm
musi działać na słowach w ten sposób, co daje jednoznaczność.
Przejście na iloraz to zawieranie domknięcia normalnego $R$
w jądrze, równoważne zerowaniu relacji.
\end{proof}
\begin{przyklad}
Grupa podstawowa bukietu $r$ okręgów jest wolna na $r$
generatorach. Dołączenie $2$-komórki wzdłuż słowa
$aba^{-1}b^{-1}$ wymusza przemienność generatorów i daje
$\langle a,b\mid aba^{-1}b^{-1}\rangle\cong\Z^2$.
Ścisły związek z dołączaniem komórek dowiedziemy z twierdzenia
van Kampena.
\end{przyklad}

\section{Pierścienie, ideały i konstrukcja ilorazu}
\begin{definicja}
\emph{Pierścień z jedynką} $(R,+,\cdot,0,1)$ ma abelową grupę
addytywną, łączne mnożenie z jedynką oraz rozdzielność mnożenia
względem dodawania. Jest \emph{przemienny}, gdy $ab=ba$.
Podzbiór $I\subseteq R$ jest \emph{ideałem dwustronnym}, gdy
jest podgrupą addytywną i $rI,Ir\subseteq I$ dla wszystkich $r\in R$.
\end{definicja}
\begin{przyklad}
$\Z$, $\Q[x]$, pierścień macierzy $M_n(\R)$ i pierścień grupowy
$\Z[G]$ są pierścieniami z jedynką. Pierścień $M_n(\R)$
przy $n\ge2$ jest nieprzemienny. W $\Z$ każdy ideał ma postać
$m\Z$; w $\Q[x]$ ideały mają postać $(f)$.
\end{przyklad}
\begin{twierdzenie}[Pierścień ilorazowy]
Jeżeli $I\triangleleft R$ jest ideałem dwustronnym, to na $R/I$
wzory $(r+I)+(s+I)=(r+s)+I$ i $(r+I)(s+I)=rs+I$
definiują pierścień. Jeżeli $\varphi:R\to S$ jest homomorfizmem
pierścieni, to $R/\ker\varphi\cong\operatorname{im}\varphi$.
\end{twierdzenie}
\begin{proof}
Dla $r'=r+i$, $s'=s+j$ mamy
$r's'-rs=rj+is+ij\in I$, stąd mnożenie nie zależy
od reprezentantów. Aksjomaty dziedziczymy z $R$.
Izomorfizm ilorazu z obrazem jest zadany przez $r+\ker\varphi
\mapsto\varphi(r)$; sprawdzenie jego poprawności i bijektywności
przebiega tak jak dla grup.
\end{proof}
\begin{definicja}
\emph{Dziedziną całkowitości} nazywamy niezerowy pierścień
przemienny z jedynką bez dzielników zera. \emph{Ciałem} jest
pierścień przemienny, w którym każdy element niezerowy jest
odwracalny. \emph{Dziedziną ideałów głównych} (PID) jest
dziedzina, w której każdy ideał jest generowany przez jeden element.
\end{definicja}
\begin{przyklad}
$\Z$ jest PID na mocy dzielenia z resztą, podobnie $K[x]$
dla dowolnego ciała $K$. Pierścień $\Z[x]$ nie jest PID:
ideał $(2,x)$ nie jest główny, ponieważ jego generator musiałby
dzielić jednocześnie $2$ i $x$, a ich jedyne wspólne dzielniki to jednostki.
Jednak $1\notin(2,x)$, co widać po podstawieniu $x=0$
i redukcji modulo $2$. Pierścień $\Z/6\Z$ ma dzielniki zera,
ponieważ $[2][3]=[0]$, więc nie jest dziedziną.
\end{przyklad}

\section{Moduły, podmoduły, ilorazy i ranga}
\begin{definicja}
Niech $R$ będzie pierścieniem z jedynką. \emph{Lewy $R$-moduł}
$M$ jest grupą abelową z działaniem $R\times M\to M$ takim,
że $(rs)m=r(sm)$, $1m=m$ oraz zachodzą obie rozdzielności.
Podmoduł $N\subseteq M$ jest podgrupą addytywną zamkniętą
na mnożenie przez skalar. Homomorfizm $R$-modułów zachowuje
sumy i mnożenie przez elementy $R$.
\end{definicja}
\begin{przyklad}
Moduły nad $\Z$ to dokładnie grupy abelowe, a moduły nad
ciałem to przestrzenie wektorowe. $2\Z\subset\Z$ jest
podmodułem, którego ilorazem jest $\Z/2\Z$.
Przy pierścieniu $R=\Z/4\Z$ podmoduł $2R$ nie jest wolnym
$R$-modułem: ma dwa elementy, podczas gdy każdy wolny
moduł niezerowej skończonej rangi ma co najmniej cztery.
\end{przyklad}
\begin{twierdzenie}[Pierwsze twierdzenie o izomorfizmie modułów]
Jeżeli $f:M\to N$ jest homomorfizmem $R$-modułów, to
$M/\ker f\cong\operatorname{im}f$ jako $R$-moduły.
\end{twierdzenie}
\begin{proof}
Mapa $m+\ker f\mapsto f(m)$ jest poprawna i bijektywna
dzięki rachunkowi jak dla grup; dodatkowo
$f(rm)=rf(m)$ pokazuje zgodność działania $R$.
\end{proof}

\begin{definicja}
Rodzina $(e_i)_{i\in I}$ jest \emph{bazą modułu} $M$,
jeżeli każdy element ma jednoznaczne przedstawienie jako
skończona kombinacja $\sum_{i\in I}r_ie_i$.
Moduł z bazą nazywamy \emph{wolnym}. Zapis $R^{(I)}$
oznacza sumę prostą $|I|$ kopii $R$, a nie iloczyn wszystkich
kopii. \emph{Ciąg dokładny} $M'\xrightarrow{u}M\xrightarrow{v}M''$
spełnia $\operatorname{im}u=\ker v$; \emph{krótki ciąg dokładny}
ma postać $0\to M'\xrightarrow{u}M\xrightarrow{v}M''\to0$.
\end{definicja}
\begin{twierdzenie}[Rozszczepienie krótkiego ciągu]
Dla krótkiego ciągu dokładnego $R$-modułów
$0\to A\xrightarrow{i}B\xrightarrow{p}C\to0$
następujące warunki są równoważne:
(i) istnieje $R$-liniowe $s:C\to B$ z $ps=\id_C$;
(ii) istnieje $R$-liniowe $r:B\to A$ z $ri=\id_A$;
(iii) $B\cong A\oplus C$ w sposób zgodny z danym ciągiem.
\end{twierdzenie}
\begin{proof}
Z (i) wynika, że $b=(b-sp(b))+sp(b)$, gdzie pierwszy
składnik należy do $\ker p=\operatorname{im}i$.
Suma jest prosta: jeśli $i(a)=s(c)$, to zastosowanie $p$
daje $c=0$, następnie $a=0$. To daje (iii).
Z (iii) otrzymujemy naturalne włączenie $C\to A\oplus C$
i rzut na $A$, co daje (i) i (ii). Alternatywnie z (ii) mamy
$B=i(A)\oplus\ker r$, a $p|_{\ker r}$ jest izomorfizmem na $C$.
\end{proof}
\begin{przyklad}[Ciąg nierozszczepialny]
$0\to\Z\xrightarrow{\cdot 2}\Z\to\Z/2\Z\to0$
jest dokładny, ale nie rozszczepia się. Gdyby $s:\Z/2\Z\to\Z$
był przekrojem addytywnym, obraz klasy $[1]$ musiałby być
jednocześnie niezerowy (bo $ps([1])=[1]$) i mieć rząd $2$,
co jest niemożliwe w $\Z$.
\end{przyklad}


\section{Iloczyn tensorowy modułów: konstrukcja bez zewnętrznych odsyłaczy}
\label{sec:tensor-algebra-v5}
Dotychczas operowaliśmy na grupach $M$, $N$ osobno. Iloczyn
tensorowy reprezentuje odwzorowania \emph{dwuliniowe}:
zamiast rozpatrywać funkcję dwu zmiennych, można rozpatrywać
jedno odwzorowanie liniowe z nowego modułu.

\begin{definicja}[Odwzorowanie zbilansowane]
Niech $R$ będzie pierścieniem przemiennym z jedynką,
a $M,N,P$ modułami nad $R$. Odwzorowanie
$b:M\times N\to P$ jest $R$-\emph{dwuliniowe}, jeśli jest
addytywne w każdej zmiennej oraz
$b(rm,n)=r b(m,n)=b(m,rn)$ dla wszystkich $r\in R$.
\end{definicja}
\begin{twierdzenie}[Istnienie i uniwersalność tensoru]
\label{thm:tensor-uni-v5}
Dla dwóch modułów $M,N$ nad przemiennym pierścieniem $R$
istnieje $R$-moduł $M\otimes_RN$ oraz odwzorowanie
dwuliniowe $\otimes:M\times N\to M\otimes_RN$ takie, że
każde odwzorowanie dwuliniowe $b:M\times N\to P$
ma dokładnie jedno faktoryzowanie
$b=\overline b\circ\otimes$ przez $R$-liniowe
$\overline b:M\otimes_RN\to P$.
Para ta jest jednoznaczna z dokładnością do jednoznacznego
izomorfizmu zachowującego elementy $m\otimes n$.
\end{twierdzenie}
\begin{proof}
Niech $F$ będzie wolnym $R$-modułem o bazie symboli $[m,n]$
indeksowanych przez pary $(m,n)\in M\times N$.
Niech $L$ będzie podmodułem generowanym przez elementy
\[
[m+m',n]-[m,n]-[m',n],\quad
[m,n+n']-[m,n]-[m,n'],\quad
[rm,n]-r[m,n],\quad[m,rn]-r[m,n].
\]
Definiujemy $M\otimes_RN=F/L$ i $m\otimes n=[m,n]+L$.
Relacje gwarantują dwuliniowość. Dla $b$ jak w tezie
wolność $F$ daje liniową mapę $B:F\to P$, $B([m,n])=b(m,n)$.
Dwuliniowość oznacza $L\subseteq\ker B$, więc $B$ przechodzi
na jedyną mapę $\overline b:F/L\to P$. Jedyność bierze się
stąd, że elementy $m\otimes n$ generują $F/L$.
Jeżeli $(T,t)$ ma tę samą własność, to uniwersalność daje
mapy $T\to M\otimes_RN$ i odwrotną, a obie kompozycje
są identycznościami na generatorach, więc są identycznościami.
\end{proof}
\begin{przyklad}
Dla grup abelowych $(\Z/m)\otimes_\Z(\Z/n)\cong\Z/d$,
$d=\gcd(m,n)$, przy $m,n\ge1$. Istotnie tensor jest
generowany przez $[1]\otimes[1]$, a relacje $m[1]=n[1]=0$
narzucają $m t=nt=0$; odwzorowanie dwuliniowe
$([a],[b])\mapsto[ab]\in\Z/d$ dowodzi, że element
o rzędzie $d$ faktycznie istnieje. W szczególności
$\Z/2\otimes\Z/3=0$, choć oba czynniki są niezerowe.
\end{przyklad}
\begin{stwierdzenie}[Rachunek tensorów]
Dla przemiennego $R$ i $R$-modułów $M,N,P$ istnieją naturalne
izomorfizmy
\[
R\otimes_RM\cong M,\qquad
(M\oplus N)\otimes_RP\cong(M\otimes_RP)\oplus(N\otimes_RP),
\qquad (M\otimes_RN)\otimes_RP\cong M\otimes_R(N\otimes_RP).
\]
\end{stwierdzenie}
\begin{proof}
Dla pierwszego izomorfizmu wzór $r\otimes m\mapsto rm$
jest liniowy z odwrotnością $m\mapsto1\otimes m$.
W drugim wzór $(m,n)\otimes p\mapsto(m\otimes p,n\otimes p)$
i jego odwrotność wynikają z uniwersalności, bo oba wyrażenia
są dwuliniowe. Dla trzeciego definiujemy na prostych tensorach
$(m\otimes n)\otimes p\mapsto m\otimes(n\otimes p)$;
kolejne zastosowania uniwersalności gwarantują poprawność
względem wszystkich relacji. Odwrotne przestawienie nawiasów
daje odwrotność.
\end{proof}
\begin{twierdzenie}[Prawostronna dokładność tensorowania]
Jeżeli $M'\xrightarrow{u}M\xrightarrow{v}M''\to0$ jest
dokładnym ciągiem $R$-modułów nad przemiennym $R$,
to dla każdego $R$-modułu $N$ ciąg
$M'\otimes_RN\xrightarrow{u\otimes1}M\otimes_RN
\xrightarrow{v\otimes1}M''\otimes_RN\to0$ jest dokładny.
\end{twierdzenie}
\begin{proof}
Surjektywność ostatniej mapy wynika z faktu, że każdy
prosty tensor $m''\otimes n$ ma podniesienie $m\otimes n$.
Ponieważ $M''\cong M/\operatorname{im}u$, mapa $M\otimes N\to
M''\otimes N$ jest uniwersalnym ilorazem, w którym
$(u(m'))\otimes n$ staje się zerem: po przejściu do
ilorazu $(M\otimes N)/\operatorname{im}(u\otimes1)$
wartość $m\otimes n$ zależy tylko od klasy $m$ modulo
$\operatorname{im}u$. Odwrotnie, ta zależność daje mapę
z $M''\otimes N$ do ilorazu. Mapy są wzajemnie odwrotne
na tensorach prostych.
\end{proof}
\begin{przyklad}[Dlaczego brak lewostronnej dokładności?]
Wstrzyknięcie $\Z\xrightarrow{\cdot2}\Z$ po tensorowaniu
z $\Z/2\Z$ staje się mapą zerową
$\Z/2\Z\xrightarrow{0}\Z/2\Z$, która nie jest injekcją.
Właśnie tę utratę informacji wykryje funktor $\operatorname{Tor}$.
\end{przyklad}

\begin{uwaga}[Ranga modułu nie jest liczbą generatorów]
Jeśli $R$ jest dziedziną z ciałem ułamków $K$, rangą $R$-modułu
$M$ nazywamy $\dim_K(K\otimes_R M)$; iloczyn tensorowy
został właśnie skonstruowany. Dla $R=\Z$
moduł $\Z/6\Z$ ma rangę $0$, choć potrzebuje generatora,
a $\Z^3\oplus\Z/2\Z$ ma rangę $3$.
Nie wolno nazywać rangą wymiaru $M/N$, chyba że pracujemy
w kategorii przestrzeni wektorowych lub po rozszerzeniu skalarów.
\end{uwaga}

\section{Projektywność, \texorpdfstring{$\operatorname{Tor}$ i $\operatorname{Ext}$}{Tor i Ext}}
\begin{definicja}
Moduł $P$ nad $R$ nazywamy \emph{projektywnym}, jeśli dla
każdej surjekcji $e:M\twoheadrightarrow N$ i każdej mapy
$f:P\to N$ istnieje $\widetilde f:P\to M$ z
$e\widetilde f=f$. Moduł $F$ nazywamy \emph{płaskim}, jeśli
$-\otimes_RF$ przekształca wszystkie krótkie ciągi dokładne
w krótkie ciągi dokładne.
\end{definicja}
\begin{stwierdzenie}
Wolny moduł jest projektywny, a projektywny jest płaski.
Każdy moduł jest ilorazem wolnego modułu.
\end{stwierdzenie}
\begin{proof}
Jeżeli $P$ jest wolny o bazie $(e_i)$, to dla każdego $f$
wybieramy preobrazy $f(e_i)$ przez surjekcję $e$ i przedłużamy
liniowo. To daje podniesienie. Dla dowolnego $M$ surjekcja
z wolnego modułu na generatorach symboli $[m]$ jest zadana
$[m]\mapsto m$. Moduł projektywny jest składnikiem prostym
wolnego: podnosząc $\id_P$ wzdłuż wolnej surjekcji $F\to P$
otrzymujemy $F\cong P\oplus\ker(F\to P)$.
Tensorowanie przez wolny moduł jest dokładne, bo jest sumą
prostą kopii odwzorowań danego ciągu; dokładność sum prostych
sprawdzamy po współrzędnych. Każdy składnik prosty funktora
dokładnego jest dokładny, stąd $P$ jest płaski.
\end{proof}
\begin{definicja}[Wolna prezentacja i pierwsze funktory pochodne]
Dla grup abelowych $A,B$ wybieramy wolną prezentację
$0\to F_1\xrightarrow{d}F_0\to A\to0$; istnieje, ponieważ
podgrupa grupy wolnej nad $\Z$ jest wolna. Definiujemy
$\operatorname{Tor}_1^{\Z}(A,B)=\ker(d\otimes1_B)$ oraz
$\operatorname{Ext}^1_{\Z}(A,B)=\operatorname{coker}
\bigl(\operatorname{Hom}(F_0,B)\to\operatorname{Hom}(F_1,B)\bigr)$.
Niezależność od prezentacji wynika z porównania prezentacji
za pomocą podnoszeń homomorfizmów wolnych grup; dowód
zapiszemy w lemacie niżej.
\end{definicja}
\begin{lemat}[Niezależność od prezentacji]
Jeżeli $F_1\to F_0$ i $F'_1\to F'_0$ są wolnymi prezentacjami
tej samej grupy $A$, to istnieją mapy łańcuchowe między
odpowiednimi dwustopniowymi kompleksami indukujące
$\id_A$ w stopniu zerowym, a ich kompozycje są łańcuchowo
homotopijne z identycznościami. Dlatego $\operatorname{Tor}_1$
i $\operatorname{Ext}^1$ powyżej nie zależą od prezentacji.
\end{lemat}
\begin{proof}
Niech $p:F_0\twoheadrightarrow A$, $p':F'_0\twoheadrightarrow A$.
Wybieramy dla każdego elementu bazy $F_0$ podniesienie przez
$p'$ jego obrazu przez $p$, uzyskując $f_0:F_0\to F'_0$
z $p'f_0=p$. Mapa $f_0d$ ląduje w $\ker p'=\operatorname{im}d'$, 
a $d':F'_1\to\operatorname{im}d'$ jest izomorfizmem;
można zatem położyć $f_1=(d')^{-1}f_0d$.
Analogicznie konstruujemy $g_0,g_1$ w przeciwną stronę.
Ponieważ $gf-\id$ indukuje zero na $A$, obraz
$g_0f_0-\id_{F_0}$ leży w $\operatorname{im}d$.
Definiujemy $h_0=d^{-1}(g_0f_0-\id)$ na $F_0$.
Wtedy $dh_0=g_0f_0-\id$ oraz, ponieważ $d$ jest injekcją,
$h_0d=g_1f_1-\id$; jest to homotopia łańcuchowa.
Podobnie dla $fg-\id$. Tensorowanie zachowuje homotopie
łańcuchowe, a zastosowanie $\operatorname{Hom}(-,B)$
odwraca ich kierunek i również zachowuje homotopie.
Stąd izomorfizmy odpowiednich grup homologii i kohomologii.
\end{proof}
\begin{przyklad}
Dla $A=\Z/m\Z$ prezentacja
$0\to\Z\xrightarrow{\cdot m}\Z\to A\to0$ daje
\[
\operatorname{Tor}_1^\Z(A,B)=\{b\in B:mb=0\},\qquad
\operatorname{Ext}^1_\Z(A,B)=B/mB.
\]
Jeśli $A$ jest wolna, oba funktory znikają, gdyż można
wybrać prezentację $F_1=0$.
\end{przyklad}
\begin{uwaga}[Połączenie z następnym rozdziałem]
Dalszy rozdział o algebrze homologicznej wprowadzi
kompleksy i krótkie ciągi dokładne kompleksów. Definicje
$\operatorname{Tor}$ i $\operatorname{Ext}$ wykorzystamy potem
w twierdzeniu Künnetha oraz twierdzeniu o współczynnikach
uniwersalnych. Powyższy dowód pokazuje, dlaczego tensoru
nie wolno traktować jak zwykłego iloczynu kartezjańskiego.
\end{uwaga}
\clearpage


\chapter{Algebra homologiczna: moduły, kompleksy i (ko)homologie}
\label{ch:algebra-homologiczna}

W poprzednim rozdziale zdefiniowaliśmy moduły, ilorazy,
iloczyny tensorowe i funktory $\operatorname{Tor}$ oraz $\operatorname{Ext}$.
Teraz badamy ciągi odwzorowań między modułami. Najważniejsza idea
jest prosta: mamy dwa rodzaje elementów, które spełniają
równanie $\partial x=0$. Jedne pochodzą już z poprzedniego
etapu, czyli są postaci $x=\partial y$; pozostałe tworzą
\emph{homologię}. W kohomologii strzałki biegną w
przeciwną stronę. Te konstrukcje będą później wspólnym
językiem dla różnych sposobów badania przestrzeni,
ale w tym rozdziale pracujemy wyłącznie algebraicznie.

\section{Moduły, homomorfizmy i ilorazy: przypomnienie}
\label{sec:moduly}
Definicje modułu, podmodułu, modułu wolnego, ilorazu i homomorfizmu
wraz z dowodem pierwszego twierdzenia o izomorfizmie podaliśmy już
w Rozdziale~\ref{ch:algebra-abstrakcyjna-v5}. W tej części przechodzimy
od konstrukcji modułowych do kompleksów łańcuchowych i ich homologii.
Na użytek rachunków nad $\Z$ przypomnijmy tylko torsję.

\begin{definicja}[Torsja grupy abelowej]
\label{def:torsja}
Niech $A$ będzie grupą abelową. Element $a\in A$ jest
\emph{torsyjny}, jeżeli $ka=0$ dla pewnego $k\in\Z\setminus\{0\}$.
Podgrupa torsyjna $\operatorname{Tors}(A)$ składa się ze
wszystkich takich elementów. Dla $m\ge2$ oznaczamy
$A[m]=\{a\in A:ma=0\}$.
Przykładowo $\operatorname{Tors}(\Z\oplus\Z/6)=0\oplus\Z/6$,
ale $A[2]$ w tym przypadku ma tylko dwa elementy.
\end{definicja}

\begin{twierdzenie}[Przypomnienie twierdzenia o izomorfizmie]
\label{thm:modul-izomorfizm}
Dla $R$-liniowego $f:M\to N$ mamy naturalny izomorfizm
$M/\ker f\cong\operatorname{im}f$, $[m]\mapsto f(m)$.
\end{twierdzenie}
\begin{proof}
Jest to Twierdzenie~\ref{thm:izo-grupy-v5} zastosowane do
addytywnych grup modułów, uzupełnione o zgodność działań skalarów
$f(rm)=rf(m)$. Szczegółowy dowód modułowy znajduje się
w Rozdziale~\ref{ch:algebra-abstrakcyjna-v5}.
\end{proof}

\section{Ciągi modułów i dokładność}
\label{sec:dokladnosc-modulow}

Mając kolejne homomorfizmy
\[
 \cdots\longrightarrow M_{i-1}
 \xrightarrow{f_{i-1}}M_i
 \xrightarrow{f_i}M_{i+1}
 \longrightarrow\cdots,
\]
pytamy, które elementy $M_i$ przychodzą z lewej
i które znikają po wysłaniu w prawo.
Samo słowo \emph{dokładny} odnosi się do
ciągu \emph{w danym miejscu}, a nie do
pojedynczego modułu.

\begin{definicja}[Dokładność ciągu]
\label{def:ciag-dokladny}
Ciąg jest \emph{dokładny w $M_i$}, jeśli
\[
 \operatorname{im}f_{i-1}=\ker f_i.
\]
Jest \emph{dokładny}, jeśli tak jest
w każdym jego miejscu. Równość zawiera
dwa żądania: $f_i f_{i-1}=0$ oraz to,
że każdy element zabity przez $f_i$
już pochodzi z $M_{i-1}$.
\emph{Długi ciąg dokładny} ma wiele
następujących po sobie takich miejsc;
może rozciągać się w obie strony
bez wyznaczonego początku.
\end{definicja}

Ponieważ jedyne odwzorowanie $0\to M$ ma obraz $0$,
warunek dokładności $0\to A\xrightarrow{i}B$
w $A$ mówi $\ker i=0$, czyli $i$ jest iniekcją.
Analogicznie dokładność
$B\xrightarrow{p}C\to0$ w $C$ mówi
$\operatorname{im}p=C$, czyli $p$ jest surjekcją.

\begin{definicja}[Krótki ciąg dokładny]
\label{def:krotki-ciag}
Ciąg
\begin{equation}\label{eq:krotki-ciag}
 0\longrightarrow A\xrightarrow{i}B
   \xrightarrow{p}C\longrightarrow0
\end{equation}
jest \emph{krótki i dokładny}, gdy $i$ jest
iniekcją, $p$ jest surjekcją oraz
$\operatorname{im}i=\ker p$.
Wtedy $C\cong B/i(A)$ przez
Twierdzenie~\ref{thm:modul-izomorfizm}
zastosowane do $p$. Mówimy też,
że $B$ jest \emph{rozszerzeniem} $C$ przez $A$.
\end{definicja}

\begin{przyklad}[Reszty z dzielenia]
\label{prz:ses-mod-m}
Dla $m\geq2$ mamy krótki ciąg dokładny
\[
 0\longrightarrow\Z\xrightarrow{\ \times m\ }\Z
   \xrightarrow{\ \pi\ }\Z/m\Z\longrightarrow0,
 \qquad \pi(k)=[k].
\]
Pierwsze odwzorowanie jest iniektywne: $mk=0$
w $\Z$ pociąga $k=0$. Odwzorowanie $\pi$ jest
surjektywne, bo każda klasa ma reprezentanta.
Jej jądro to dokładnie wielokrotności $m$,
czyli obraz pierwszego odwzorowania.
\end{przyklad}

Krótki ciąg może, ale nie musi, rozkładać się
na dwie niezależne części. Nie należy mylić
wyboru dowolnego przedstawiciela każdej klasy
z \emph{liniowym} wyborem przedstawicieli.

\begin{twierdzenie}[Kiedy krótki ciąg się rozszczepia?]
\label{thm:ses-rozszczepienie}
Dla~\eqref{eq:krotki-ciag} równoważne są:
\begin{enumerate}[leftmargin=*]
\item istnieje homomorfizm $s:C\to B$
      z $p s=\id_C$ (przekrój $p$);
\item istnieje homomorfizm $r:B\to A$
      z $r i=\id_A$ (retrakcja $i$);
\item istnieje izomorfizm $B\cong A\oplus C$
      zgodny z odwzorowaniami $i(a)=(a,0)$
      i $p(a,c)=c$.
\end{enumerate}
Wtedy ciąg nazywamy \emph{rozszczepionym}.
\end{twierdzenie}
\begin{proof}
Załóżmy najpierw, że mamy $s$.
Każde $b\in B$ można rozłożyć jako
\[
 b=(b-s(p(b)))+s(p(b)).
\]
Pierwszy składnik leży w $\ker p=\operatorname{im}i$.
Ponieważ $i$ jest iniekcją, istnieje dokładnie
jeden $a\in A$ taki, że $i(a)=b-s(p(b))$.
Odwzorowanie
$\Phi:A\oplus C\to B$, $\Phi(a,c)=i(a)+s(c)$
jest więc surjektywne. Jeżeli $\Phi(a,c)=0$,
to po zastosowaniu $p$ otrzymujemy $c=0$,
a wtedy $i(a)=0$, czyli $a=0$.
To dowodzi (1)$\Rightarrow$(3).

Z izomorfizmu $\Phi:A\oplus C\to B$ w (3) i rzutu
$\operatorname{pr}_A:A\oplus C\to A$ bierzemy
$r=\operatorname{pr}_A\circ\Phi^{-1}$. Dostajemy
$r:B\to A$ z $ri=\id_A$. To daje (3)$\Rightarrow$(2).
Wreszcie, mając $r$, dla $c\in C$ wybierzmy
jakikolwiek $b$ z $p(b)=c$ i połóżmy
$s(c)=b-i(r(b))$. Jeśli zmienimy $b$ na
$b+i(a)$, nowa wartość wyniesie
\[
 b+i(a)-i(r(b+i(a)))
 =b+i(a)-i(r(b)+a)=b-i(r(b)),
\]
ponieważ $r(i(a))=a$ i $r$ jest liniowa.
Tak otrzymane odwzorowanie $s$ nie zależy od
wybranego $b$ i jest liniowe: dla $c,c'$
możemy jako podniesienie $c+c'$ wziąć sumę
podniesień $c,c'$, a dla $tc$ podniesienie $tb$.
Ponadto $p(s(c))=p(b)=c$.
Otrzymaliśmy (2)$\Rightarrow$(1).
\end{proof}

\begin{uwaga}[Dwa kontrastujące przypadki]
Jeśli $R=\mathbb K$ jest ciałem, każdy krótki
ciąg przestrzeni liniowych rozszczepia się:
wybierzmy bazę $C$, podnieśmy jej wektory
do $B$ przez surjekcję $p$ i rozszerzmy
liniowo do odwzorowania $s$. Nie zależy nam na
kanoniczności tego wyboru.

Ciąg z Przykładu~\ref{prz:ses-mod-m} dla
$m\geq2$ \emph{nie} rozszczepia się jako
ciąg $\Z$-modułów. Gdyby przekrój $s$
istniał, dla $z=s([1])\in\Z$ byłoby
$mz=s(m[1])=0$, więc $z=0$.
Ale wtedy $\pi(s([1]))=[0]\ne[1]$.
Brak dzielenia przez $m$ uniemożliwia tu
liniowy wybór reprezentantów.
\end{uwaga}
\section{Kompleksy łańcuchowe i homologia algebraiczna}
\label{sec:kompleks-homologia}

W dokładnym ciągu każdy element zabity przez
następne odwzorowanie pochodzi z poprzedniego miejsca.
Do zdefiniowania homologii wystarczy najpierw
słabszy warunek: obraz poprzedniego odwzorowania ma
\emph{mieścić się} w jądrze następnego.
Homologia mierzy brak równości.

\begin{definicja}[Kompleks łańcuchowy]
\label{def:kompleks-lancuchowy}
\emph{Kompleks łańcuchowy} $C_\bullet$ nad $R$
to moduły $C_n$, $n\in\Z$, i homomorfizmy
$\partial_n:C_n\to C_{n-1}$ takie, że
\begin{equation}\label{eq:d-kwadrat-zero}
 \partial_{n-1}\circ \partial_n=0\qquad\text{dla każdego }n.
\end{equation}
Odwzorowanie $\partial_n$ nazywamy \emph{różniczką} albo
\emph{brzegiem}. Dla kompleksów łańcuchowych
będziemy używać symbolu $\partial$, a symbol $d$
pozostawimy różniczkom kołańcuchowym.
Gdy $C_n=0$ dla $n<0$,
mówimy o kompleksie w stopniach nieujemnych.
\end{definicja}

Warunek~\eqref{eq:d-kwadrat-zero} należy czytać
dosłownie: brzeg tego, co samo jest brzegiem,
wynosi zero. Nie zakładamy, że wszystkie
elementy o zerowym brzegu są brzegami.

\begin{definicja}[Cykle, brzegi i grupy homologii]
\label{def:homologia-algebraiczna}
W stopniu $n$ określamy podmoduły
\[
 Z_n(C):=\ker(\partial_n:C_n\to C_{n-1}),\qquad
 B_n(C):=\operatorname{im}(\partial_{n+1}:C_{n+1}\to C_n).
\]
Elementy $Z_n(C)$ to \emph{cykle}, a elementy
$B_n(C)$ to \emph{brzegi}. Ponieważ
$\partial_n\partial_{n+1}=0$, mamy $B_n(C)\subseteq Z_n(C)$,
więc można zdefiniować \emph{homologię stopnia $n$}
\begin{equation}\label{eq:homologia-def}
 H_n(C):=Z_n(C)/B_n(C)
 =\frac{\ker \partial_n}{\operatorname{im}\partial_{n+1}}.
\end{equation}
Klasę cyklu $z$ oznaczamy $[z]$.
\end{definicja}

Dwa cykle $z,z'$ opisują tę samą klasę dokładnie
wtedy, gdy $z-z'=\partial_{n+1}w$ dla pewnego
$w\in C_{n+1}$. To jest algebraiczne
znaczenie słowa \emph{homologiczne}.
Klasa $[z]$ znika dokładnie wtedy, gdy
$z$ jest brzegiem. Z Definicji~\ref{def:ciag-dokladny}
otrzymujemy zatem
\[
 H_n(C)=0
 \quad\Longleftrightarrow\quad
 \operatorname{im}\partial_{n+1}=\ker \partial_n
 \quad\Longleftrightarrow\quad
 C_\bullet\text{ jest dokładny w }C_n.
\]
Gdy znika homologia we wszystkich stopniach,
kompleks nazywamy \emph{acyklicznym}.
Nie znaczy to, że same moduły $C_n$ są zerowe.

\begin{przyklad}[Cztery rachunki]
\label{prz:homologia-rachunki}
\begin{enumerate}[leftmargin=*]
\item W kompleksie
      $0\to\Z\xrightarrow{\ \times2\ }\Z\to0$,
      gdzie niezerowe moduły są w stopniach
      $1$ i $0$, mamy $Z_1=0$, $Z_0=\Z$,
      $B_0=2\Z$. Stąd
      $H_1=0$ i $H_0=\Z/2\Z$.
\item Jeżeli na tych samych dwóch modułach
      postawimy różniczkę zerową, każdy element
      jest cyklem, a jedynym brzegiem jest $0$.
      Otrzymujemy $H_1\cong H_0\cong\Z$.
\item Dla
      $0\to R\xrightarrow{\ \id\ }R\to0$
      odwzorowanie w stopniu $1$ ma zerowe jądro,
      a w stopniu $0$ jego obraz jest całym $R$.
      Wszystkie homologie są zerowe, choć
      dwa moduły kompleksu są niezerowe.
\item Torsja może powstać nawet z modułów
      wolnych. Weźmy
      \[
       0\longrightarrow\Z
       \xrightarrow{\ \partial_2\ }\Z^2
       \xrightarrow{\ \partial_1\ }\Z
       \longrightarrow0,\qquad
       \partial_2(t)=(0,2t),\quad \partial_1(x,y)=x.
      \]
      Mamy $\partial_1\partial_2=0$.
      Cykle w stopniu $1$ to $(0,y)$,
      a brzegi to $(0,2t)$, więc
      $H_1\cong\Z/2\Z$. Odwzorowanie $\partial_1$ jest
      surjektywne, więc $H_0=0$; odwzorowanie
      $\partial_2$ jest iniektywne, więc $H_2=0$.
\end{enumerate}
\end{przyklad}

\begin{definicja}[Odwzorowanie łańcuchowe]
\label{def:mapa-lancuchowa}
Odwzorowanie $f:C_\bullet\to D_\bullet$ to rodzina
homomorfizmów $f_n:C_n\to D_n$ zgodnych
z różniczkami:
\begin{equation}\label{eq:mapa-lancuchowa}
 \partial^D_n f_n=f_{n-1}\partial^C_n.
\end{equation}
Oznacza to, że najpierw zastosować $f$,
a potem wziąć brzeg, daje ten sam wynik,
co zrobić to w odwrotnej kolejności.
\end{definicja}

\begin{stwierdzenie}[Odwzorowanie indukowane na homologii]
\label{prop:homologia-funktor}
Odwzorowanie łańcuchowe $f$ wyznacza homomorfizmy
$H_n(f):H_n(C)\to H_n(D)$,
$H_n(f)[z]=[f_n(z)]$. Zachodzą
$H_n(\id_C)=\id_{H_n(C)}$ oraz
$H_n(gf)=H_n(g)H_n(f)$.
\end{stwierdzenie}
\begin{proof}
Jeżeli $\partial^C_n z=0$, to z~\eqref{eq:mapa-lancuchowa}
$\partial^D_n f_nz=f_{n-1}\partial^C_nz=0$, czyli
$f_nz$ jest cyklem. Jeżeli zmienimy
$z$ na $z+\partial^C_{n+1}w$, jego obraz zmieni się o
\[
 f_n(\partial^C_{n+1}w)
 =\partial^D_{n+1}(f_{n+1}w),
\]
czyli o brzeg. Odwzorowanie na klasach jest więc
dobrze określone i liniowe.
Równości dla identyczności i złożenia
wynikają przez zastosowanie wzorów
do dowolnej klasy $[z]$.
\end{proof}

\section{Skąd bierze się długi ciąg dokładny homologii?}
\label{sec:dlugi-ciag-homologii}

Niech $A_\bullet,B_\bullet,C_\bullet$ będą
kompleksami, a $i:A_\bullet\to B_\bullet$,
$p:B_\bullet\to C_\bullet$ odwzorowaniami łańcuchowymi.
Mówimy, że
\[
 0\longrightarrow A_\bullet
 \xrightarrow{i}B_\bullet
 \xrightarrow{p}C_\bullet
 \longrightarrow0
\]
jest \emph{krótkim ciągiem dokładnym kompleksów},
jeśli dla \emph{każdego} $n$ ciąg modułów
\begin{equation}\label{eq:ses-kompleks}
 0\longrightarrow A_n\xrightarrow{i_n}B_n
 \xrightarrow{p_n}C_n\longrightarrow0
\end{equation}
jest dokładny. Stopniowa dokładność nie
mówi jeszcze, że ciąg powstałych homologii
jest krótki i dokładny: cykl podniesiony
z $C_n$ może przestać być cyklem w $B_n$.
Jego brzeg mierzy właśnie tę przeszkodę.

\begin{figure}[!htbp]
\centering
\begin{tikzpicture}[>=Stealth,line cap=round,
  every node/.style={font=\small}]
 \node[draw=manifoldblue!75!black,rounded corners,
  fill=manifoldblue!7,minimum width=2.25cm,
  minimum height=.9cm] (c) at (-5.25,0) {$c\in Z_n(C)$};
 \node[draw=manifoldgreen!75!black,rounded corners,
  fill=manifoldgreen!8,minimum width=2.05cm,
  minimum height=.9cm] (b) at (-1.75,0) {$b\in B_n$};
 \node[draw=manifoldred!75!black,rounded corners,
  fill=manifoldred!8,minimum width=2.15cm,
  minimum height=.9cm] (db) at (1.75,0) {$\partial_Bb\in B_{n-1}$};
 \node[draw=manifoldgold!85!black,rounded corners,
  fill=manifoldgold!10,minimum width=2.4cm,
  minimum height=.9cm] (a) at (5.25,0) {$a\in Z_{n-1}(A)$};
 \draw[->,thick,manifoldblue] (c) --
  node[midway,above,font=\footnotesize] {wybór $b$} (b);
 \draw[->,thick,manifoldred] (b) --
  node[midway,above,font=\footnotesize] {$\partial_B$} (db);
 \draw[->,thick,manifoldgold!85!black] (db) --
  node[midway,above,font=\footnotesize] {$i_{n-1}^{-1}$} (a);
 \node[font=\footnotesize,manifoldgreen!75!black]
  at (-1.75,-.9) {$p_n(b)=c$};
 \node[font=\footnotesize,manifoldred!85!black]
  at (1.75,-.9) {$p_{n-1}(\partial_Bb)=\partial_Cc=0$};
 \node[font=\footnotesize,manifoldgold!85!black]
  at (5.25,-.9) {$[a]=\delta_n[c]$};
\end{tikzpicture}
\caption{Odwzorowanie łączące: podnosimy cykl $c$, bierzemy brzeg
podniesienia $b$, po czym odczytujemy go w $A_{n-1}$.
Odwrócenie $i_{n-1}$ dotyczy tylko jego obrazu
$\ker p_{n-1}$; całego odwzorowania $i_{n-1}$ nie odwracamy.}
\label{fig:homologia-delta}
\end{figure}

\begin{twierdzenie}[Długi ciąg dokładny]
\label{thm:les-homologia}
Z krótkiego ciągu dokładnego kompleksów
otrzymujemy naturalny długi ciąg dokładny
\begin{equation}\label{eq:les-homologia}
 \cdots\longrightarrow H_n(A)
 \xrightarrow{i_*}H_n(B)
 \xrightarrow{p_*}H_n(C)
 \xrightarrow{\delta_n}H_{n-1}(A)
 \xrightarrow{i_*}H_{n-1}(B)
 \longrightarrow\cdots.
\end{equation}
Odwzorowanie \emph{łączące} $\delta_n$ ma następujący
konkretny wzór. Dla cyklu $c\in C_n$ wybieramy
podniesienie $b\in B_n$ z $p_n(b)=c$.
Wtedy $\partial_B b=i_{n-1}(a)$ dla dokładnie
jednego $a\in A_{n-1}$, a
\begin{equation}\label{eq:delta-wzor}
 \delta_n[c]=[a]\in H_{n-1}(A).
\end{equation}
Jeżeli kompleksy są zerowe w ujemnych stopniach,
końcowy fragment ma postać
$H_0(A)\to H_0(B)\to H_0(C)\to0$.
\end{twierdzenie}
\begin{proof}
\emph{Konstrukcja i cykl.} Surjektywność $p_n$
daje podniesienie $b$. Ponieważ $p$ jest odwzorowaniem łańcuchowym,
\[
 p_{n-1}(\partial_Bb)=\partial_C(p_nb)=\partial_Cc=0.
\]
Z dokładności~\eqref{eq:ses-kompleks} w
stopniu $n-1$ mamy
$\partial_Bb\in\ker p_{n-1}=\operatorname{im}i_{n-1}$,
więc istnieje $a$ z $i_{n-1}a=\partial_Bb$.
Iniektywność $i_{n-1}$ daje jego jednoznaczność.
Ponadto
\[
 i_{n-2}(\partial_Aa)
 =\partial_B(i_{n-1}a)=\partial_B^2b=0.
\]
Iniektywność $i_{n-2}$ wymusza $\partial_Aa=0$.
Możemy więc wziąć klasę $[a]$.

\emph{Niezależność od podniesienia.} Jeśli $b'$ też
podnosi $c$, to $b'-b=i_n(u)$ dla pewnego
$u\in A_n$. Wtedy
\[
 \partial_Bb'=\partial_Bb+\partial_Bi_nu
 =i_{n-1}(a+\partial_Au).
\]
Nowy cykl w $A$ różni się od $a$ o brzeg,
więc wyznacza tę samą klasę.

\emph{Niezależność od reprezentanta.}
Niech $c'=c+\partial_Cv$, $v\in C_{n+1}$.
Wybierzmy podniesienie $w\in B_{n+1}$ elementu $v$.
Wtedy $b+\partial_Bw$ podnosi $c'$ i ma ten sam
brzeg co $b$, bo $\partial_B^2w=0$.
Każde inne podniesienie $c'$ różni się od tego
o element $\operatorname{im}i_n$; właśnie
sprawdzona niezależność od podniesienia
kończy dowód niezależności od wyboru reprezentanta klasy $[c]$.
Suma podniesień podnosi sumę cykli, a
podniesienie $rb$ podnosi $rc$, więc $\delta_n$
jest homomorfizmem $R$-modułów.

\emph{Dokładność przy $H_n(B)$.}
Z $pi=0$ mamy $p_*i_*=0$.
Jeśli $[b]\in H_n(B)$ spełnia
$p_*[b]=0$, to dla cyklu $b$
istnieje $v\in C_{n+1}$ z
$p_nb=\partial_Cv$. Wybierzmy podniesienie
$w\in B_{n+1}$ elementu $v$.
Wtedy $b'=b-\partial_Bw$ nadal jest cyklem,
$[b']=[b]$ i
$p_nb'=p_nb-\partial_Cv=0$. Z dokładności
w stopniu $n$ mamy $b'=i_na$.
Ponieważ $\partial_Bb'=i_{n-1}\partial_Aa=0$
i $i_{n-1}$ jest iniektywna, $a$ jest
cyklem. Stąd $[b]=i_*[a]$.

\emph{Dokładność przy $H_n(C)$.}
Jeżeli $c=p_nb$ dla cyklu $b$, możemy
w definicji $\delta_n$ wybrać to właśnie
podniesienie: $\partial_Bb=0$, więc $\delta_n[c]=0$.
Odwrotnie, załóżmy $\delta_n[c]=0$
i podnieśmy cykl $c$ do $b$.
Jeśli $\partial_Bb=i_{n-1}a$, to $[a]=0$,
czyli $a=\partial_Au$ dla pewnego $u\in A_n$.
Element $b-i_nu$ jest cyklem, bo
$\partial_B(b-i_nu)=i_{n-1}(a-\partial_Au)=0$,
a $p_n(b-i_nu)=c$. Zatem
$[c]=p_*[b-i_nu]$.

\emph{Dokładność przy $H_{n-1}(A)$.}
Dla klasy $\delta_n[c]=[a]$ mamy
$i_{n-1}a=\partial_Bb$, czyli $i_*[a]=0$.
Odwrotnie, jeżeli $a\in Z_{n-1}(A)$
i $i_*[a]=0$, istnieje $b\in B_n$
takie, że $i_{n-1}a=\partial_Bb$.
Wówczas $c=p_nb$ jest cyklem, bo
$\partial_Cc=p_{n-1}\partial_Bb=p_{n-1}i_{n-1}a=0$,
a wzór~\eqref{eq:delta-wzor} daje
$\delta_n[c]=[a]$.
Przesuwając indeks $n$ o jeden,
obejmujemy tymi trzema sprawdzeniami
każde miejsce całego ciągu~\eqref{eq:les-homologia}.
W końcu, przy założonej zerowości kompleksów w ujemnych
stopniach mamy $A_{-1}=0$, a więc
$H_{-1}(A)=0$, więc $H_0(B)\to H_0(C)$
jest surjekcją na mocy właśnie wykazanej
dokładności.
\end{proof}

\begin{uwaga}[Naturalność odwzorowania łączącego]
\label{uwaga:les-naturalnosc}
Jeżeli mamy odwzorowanie między dwoma krótkimi
ciągami kompleksów, zgodne z $i,p$ i
różniczkami, to odwzorowania na homologiach
komutują także z $\delta$.
Dokładniej, są to odwzorowania łańcuchowe
$f_A:A\to A'$, $f_B:B\to B'$, $f_C:C\to C'$ spełniające
$f_Bi=i'f_A$ oraz $p'f_B=f_Cp$.
Istotnie, podniesienie $b$ cyklu $c$ przechodzi
w podniesienie obrazu $c$, jego brzeg przechodzi
w obraz $\partial_Bb$, a wstrzyknięcie $i$
pozwala odczytać obraz tego samego $a$.
Równość na reprezentantach ma postać
$\delta'_n(H_n(f_C)[c])
=H_{n-1}(f_A)(\delta_n[c])$.
\end{uwaga}

\begin{przyklad}[Niezerowe odwzorowanie łączące]
\label{prz:delta-izomorfizm}
Niech $A_\bullet$ ma jedyny niezerowy moduł
$A_0=R$, a $C_\bullet$ jedyny niezerowy
moduł $C_1=R$. Niech $B_\bullet$ ma
$B_1=B_0=R$ i różniczkę
$\partial_B=\id_R:B_1\to B_0$.
W stopniu $0$ odwzorowanie $i_0:A_0\to B_0$
jest identycznością, a w stopniu $1$
odwzorowanie $p_1:B_1\to C_1$ jest identycznością;
pozostałe odwzorowania to jedyne możliwe odwzorowania zerowe.
W każdym stopniu otrzymujemy rozszczepiony
krótki ciąg dokładny modułów, ale
$H_n(B)=0$ we wszystkich stopniach.
Cykl $c=1\in C_1$ podnosimy do $b=1\in B_1$.
Jego brzeg to $1=i_0(1)\in B_0$,
więc $\delta_1:R\to R$ jest identycznością.
Przykład pokazuje, że stopniowe rozszczepienie
nie oznacza zgodnego z różniczkami
rozszczepienia kompleksów.
\end{przyklad}
\section{Kohomologia algebraiczna i jej długi ciąg}
\label{sec:kohomologia-algebraiczna}

W homologii różniczka obniża stopień.
Kohomologia powstaje z równie prostej
konstrukcji, w której różniczka stopień
\emph{podwyższa}. Gdy elementy kompleksu
łańcuchowego są wektorami, elementy
kompleksu dualnego są funkcjonałami:
zamiast brać brzeg wektora, oceniamy
funkcjonał na brzegu.

\begin{definicja}[Kompleks kołańcuchowy i kohomologia]
\label{def:kohomologia-algebraiczna}
\emph{Kompleks kołańcuchowy} $K^\bullet$ to
moduły $K^n$ i homomorfizmy
$d^n:K^n\to K^{n+1}$ z
$d^{n+1}d^n=0$ dla $n\in\Z$. Jego \emph{kohomologią stopnia $n$}
jest
\begin{equation}\label{eq:kohomologia-def}
 H^n(K):=\frac{\ker(d^n:K^n\to K^{n+1})}
 {\operatorname{im}(d^{n-1}:K^{n-1}\to K^n)}.
\end{equation}
Elementy jądra nazywamy \emph{kocyklami},
a elementy obrazu \emph{kobrzegami}.
Odwzorowanie kołańcuchowe $F:K^\bullet\to L^\bullet$
składa się z $F^n:K^n\to L^n$ spełniających
$d_L^nF^n=F^{n+1}d_K^n$.
\end{definicja}

Tak jak w dowodzie
Stwierdzenia~\ref{prop:homologia-funktor},
odwzorowanie kołańcuchowe przenosi kocykle w
kocykle i kobrzegi w kobrzegi:
$d_LFz=Fdz=0$ oraz
$F(d_Ku)=d_L(Fu)$. Wobec tego definiuje
$H^n(F)[z]=[F^nz]$. Równość
$H^n(K)=0$ mówi dokładnie, że
$K^{n-1}\to K^n\to K^{n+1}$
jest dokładny w $K^n$.

\begin{przyklad}[Kohomologia dwuwyrazowego kompleksu]
Dla $K^0=K^1=\Z$, $d^0=\times2$
i wszystkich innych modułów równych $0$
otrzymujemy $H^0(K)=0$ oraz
$H^1(K)=\Z/2\Z$. W stopniu $1$ nie ma
już różniczki wychodzącej, więc
wszystkie elementy $\Z$ są kocyklami;
kobrzegami są liczby parzyste.
\end{przyklad}

\begin{stwierdzenie}[Kompleks dualny]
\label{prop:kompleks-dualny}
Dla kompleksu łańcuchowego $C_\bullet$
i $R$-modułu $G$ połóżmy
\[
 C^n(C;G):=\operatorname{Hom}_R(C_n,G),
 \qquad d^n\varphi:=\varphi\circ \partial_{n+1}.
\]
Wtedy $C^\bullet(C;G)$ jest kompleksem
kołańcuchowym. Odwzorowanie łańcuchowe
$f:C_\bullet\to D_\bullet$ daje odwzorowanie
kołańcuchowe w odwrotną stronę,
$f^*:C^\bullet(D;G)\to C^\bullet(C;G)$,
$f^*(\varphi)=\varphi\circ f_n$.
\end{stwierdzenie}
\begin{proof}
Dla $\varphi:C_n\to G$ mamy
\[
 d^{n+1}(d^n\varphi)
 =(\varphi\circ \partial_{n+1})\circ \partial_{n+2}
 =\varphi\circ(\partial_{n+1}\partial_{n+2})=0.
\]
Z kolei dla $\varphi:D_n\to G$
oraz $x\in C_{n+1}$ równość odwzorowań łańcuchowych
z~\eqref{eq:mapa-lancuchowa} daje
\[
 (d f^*\varphi)(x)
 =\varphi(f_n(\partial_C x))
 =\varphi(\partial_D(f_{n+1}x))
 =(f^*d\varphi)(x).
\]
To dowodzi zgodności z różniczką.
\end{proof}

Słowo \emph{dualny} wymaga tu ostrożności:
zastosowanie $\operatorname{Hom}_R(-,G)$
do krótkiego ciągu dokładnego modułów
nie musi dać krótkiego ciągu dokładnego.
Na przykład przy $G=\Z$ w ciągu
$0\to2\Z\hookrightarrow\Z\to\Z/2\Z\to0$
funkcjonał $2\Z\to\Z$, $2k\mapsto k$,
nie przedłuża się do homomorfizmu
$\Z\to\Z$: musiałby wysłać $1$
na liczbę $t$ z $2t=1$.
Z tego powodu nie będziemy budować
długiego ciągu kohomologii przez
bezrefleksyjne dualizowanie
Twierdzenia~\ref{thm:les-homologia}.

\begin{twierdzenie}[Długi ciąg kohomologii]
\label{thm:les-kohomologia}
Jeżeli w każdym stopniu dokładny jest
krótki ciąg \emph{kompleksów kołańcuchowych}
\[
 0\longrightarrow A^\bullet\xrightarrow{i}B^\bullet
 \xrightarrow{p}C^\bullet\longrightarrow0,
\]
to istnieje długi ciąg dokładny
\begin{equation}\label{eq:les-kohomologia}
 \cdots\longrightarrow H^n(A)
 \xrightarrow{i_*}H^n(B)
 \xrightarrow{p_*}H^n(C)
 \xrightarrow{\delta^n}H^{n+1}(A)
 \xrightarrow{i_*}H^{n+1}(B)\longrightarrow\cdots.
\end{equation}
\end{twierdzenie}
\begin{proof}
Dla kocyklu $c\in C^n$ wybierzmy podniesienie
$b\in B^n$ z $p(b)=c$.
Wtedy $p(d_Bb)=d_Cc=0$, więc dokładność
w stopniu $n+1$ i iniektywność $i$
dają dokładnie jedno $a\in A^{n+1}$
z $i(a)=d_Bb$.
Ponieważ $i(d_Aa)=d_B^2b=0$,
otrzymujemy kocykl $a$.
Definiujemy $\delta^n[c]=[a]$.
Jeśli podniesienie $b$ zastąpimy $b+i(u)$,
to $a$ zmienia się na $a+d_Au$.
Jeśli $c$ zastąpimy $c+d_Cv$, to po
wyborze podniesienia $w$ dla $v$ możemy
zastąpić $b$ przez $b+d_Bw$,
co w ogóle nie zmienia $d_Bb$.
Oba wybory prowadzą zatem do tej
samej klasy. Podniesienie sumy można wziąć
jako sumę podniesień, a $rb$ podnosi $rc$ dla $r\in R$.
Zatem $\delta^n$ jest $R$-liniowa.

Sprawdźmy każde z trzech rodzajów miejsc
w~\eqref{eq:les-kohomologia}.
Przy $H^n(B)$ mamy $p_*i_*=0$.
Jeśli $b\in B^n$ jest kocyklem
i $p_*[b]=0$, to $p(b)=d_Cv$
dla pewnego $v\in C^{n-1}$.
Podnieśmy $v$ do $w\in B^{n-1}$.
Wówczas $b-d_Bw$ jest kocyklem
z tą samą klasą co $b$, a jego
obraz przez $p$ wynosi $0$.
Równa się więc $i(a)$ dla pewnego
$a\in A^n$. Z $0=d_Bi(a)=i(d_Aa)$
i iniektywności $i$ dostajemy
$d_Aa=0$, a zatem $[b]=i_*[a]$.

Przy $H^n(C)$ inkluzja
$\operatorname{im}p_*\subseteq\ker\delta^n$
wynika z wybrania jako podniesienia $c=p(b)$
samego kocyklu $b$: jego różniczka jest $0$.
Jeśli odwrotnie $\delta^n[c]=0$,
to dla podniesienia $b$ mamy $d_Bb=i(a)$,
gdzie $a=d_Au$ dla pewnego $u\in A^n$.
Element $b-i(u)$ jest kocyklem,
którego obraz to $c$, więc
$[c]\in\operatorname{im}p_*$.

Przy $H^{n+1}(A)$ element
$i_*\delta^n[c]$ reprezentuje
$i(a)=d_Bb$ i dlatego jest zerowy.
Jeśli z drugiej strony $a\in A^{n+1}$
jest kocyklem i $i_*[a]=0$,
to $i(a)=d_Bb$ dla pewnego
$b\in B^n$. Obraz $c=p(b)$ jest
kocyklem, gdyż $d_Cc=p(d_Bb)=p(i(a))=0$,
a jego odwzorowanie łączące wynosi $[a]$.
Zmiana $n$ daje dokładność także
w pozostałych stopniach. Konstrukcja jest naturalna:
dla zgodnych odwzorowań $f_A,f_B,f_C$ między krótkimi ciągami
podniesieniem $f_C(c)$ jest $f_B(b)$, a jego różniczka to
$f_B(d_Bb)=f_B(i(a))=i'(f_A(a))$.
Stąd $\delta'^n H^n(f_C)=H^{n+1}(f_A)\delta^n$.
\end{proof}

\begin{twierdzenie}[Nad ciałem kohomologia jest dualna do homologii]
\label{thm:field-hom-dual}
Niech $C_\bullet$ będzie kompleksem
przestrzeni liniowych nad ciałem
$\mathbb K$, a
$C^n=\operatorname{Hom}_{\mathbb K}(C_n,\mathbb K)$.
Wtedy naturalne przypisanie
\begin{equation}\label{eq:field-duality}
 H^n(C^\bullet)
 \longrightarrow\operatorname{Hom}_{\mathbb K}
  (H_n(C_\bullet),\mathbb K),
 \qquad [\varphi]\longmapsto
     ([z]\longmapsto\varphi(z))
\end{equation}
jest izomorfizmem. Zakładamy, że
przestrzenie pozwalają na wybór
dopełnień podprzestrzeni
(w szczególności wystarcza skończony wymiar).
\end{twierdzenie}
\begin{proof}
Jeżeli $\varphi\in C^n$ jest kocyklem,
to $d^n\varphi=\varphi \partial_{n+1}=0$.
Zatem $\varphi$ znika na $B_n$,
więc jej wartość na $z\in Z_n$
zależy tylko od klasy $[z]\in H_n$.
Jeżeli zmienimy $\varphi$ na
$\varphi+d^{n-1}\psi
=\varphi+\psi \partial_n$, to dla cyklu $z$
mamy $\partial_nz=0$, więc otrzymujemy
ten sam funkcjonał na $H_n$.
Odwzorowanie~\eqref{eq:field-duality} jest
dobrze określone i liniowe. Jest też naturalne: dla odwzorowania
łańcuchowego $f:C\to D$, kocyklu $\varphi$ na $D$ i cyklu $z$ na $C$
mamy $(f^*\varphi)(z)=\varphi(fz)$. Zatem obie drogi, przez
$H^n(f^*)$ lub przez prekompozycję z $H_n(f)$, dają ten sam funkcjonał.

Aby pokazać surjektywność, weźmy
$\ell:H_n=Z_n/B_n\to\mathbb K$.
Złożenie $Z_n\to H_n\xrightarrow{\ell}\mathbb K$
jest liniowe i znika na $B_n$.
Wybieramy dopełnienie $Z_n$ w $C_n$
i przedłużamy funkcjonał liniowo
na $C_n$, na przykład jako zero
na tym dopełnieniu. Otrzymany
$\varphi$ spełnia
$\varphi \partial_{n+1}=0$, bo
$\operatorname{im}\partial_{n+1}=B_n$;
jest kocyklem wysyłanym na $\ell$.

Jeżeli klasa $[\varphi]$ trafia na zero,
to $\varphi$ znika na całym $Z_n$.
Możemy więc zdefiniować funkcjonał
$\eta$ na $B_{n-1}=\operatorname{im}\partial_n$
przez $\eta(\partial_nx)=\varphi(x)$.
Jest dobrze określony: z
$\partial_nx=\partial_nx'$ wynika $x-x'\in Z_n$,
a tam $\varphi$ znika.
Wybieramy dopełnienie $B_{n-1}$
w $C_{n-1}$ i przedłużamy $\eta$
do $\psi:C_{n-1}\to\mathbb K$.
Wtedy $\varphi=\psi \partial_n=d^{n-1}\psi$
jest kobrzegiem. Jądro odwzorowania
w~\eqref{eq:field-duality} jest więc zerowe.
\end{proof}

\begin{uwaga}[Dlaczego ważne są współczynniki?]
Nad $\Z$ odpowiednik~\eqref{eq:field-duality}
może być fałszywy. Dla kompleksu
$C_1=C_0=\Z$ z $\partial_1=\times2$ mamy
$H_0(C)=\Z/2$ i $H_1(C)=0$,
podczas gdy dla dualnego kompleksu
$H^1(C^\bullet)=\Z/2$.
Grupa $\operatorname{Hom}_{\Z}
(H_1(C),\Z)$ jest zerowa.
Ta dodatkowa klasa kohomologii
wykrywa informację torsyjną z
poprzedniego stopnia. Nie utożsamiamy
więc homologii z kohomologią bez
sprawdzenia założeń o współczynnikach.
\end{uwaga}
\section{Homotopia łańcuchowa i uproszczenia kompleksów}
\label{sec:homotopia-lancuchowa}

Jedna homologia może mieć bardzo różne opisy
przez generatory i różniczki. Poniższe kryterium
mówi, kiedy dwa odwzorowania kompleksów dają
\emph{dokładnie to samo} odwzorowanie na
homologiach. Jest to algebraiczny mechanizm
usuwania par generatorów, które wzajemnie
się kasują.

\begin{definicja}[Homotopia łańcuchowa]
\label{def:homotopia-lancuchowa}
Dwa odwzorowania łańcuchowe $f,g:C_\bullet\to D_\bullet$
są \emph{łańcuchowo homotopijne}, jeśli istnieją
homomorfizmy $R$-modułów $h_n:C_n\to D_{n+1}$ takie, że
\begin{equation}\label{eq:chain-homotopy}
 f_n-g_n=\partial^D_{n+1}h_n+h_{n-1}\partial^C_n
 \qquad\text{dla każdego }n.
\end{equation}
Kompleks $C_\bullet$ jest \emph{ściągalny},
jeśli $\id_C$ jest łańcuchowo homotopijna
z odwzorowaniem zerowym.
\end{definicja}

Wzór~\eqref{eq:chain-homotopy} przypomina
zasadę całkowania przez części: różnica
odwzorowań składa się z dwóch operacji, w których
$h$ i $\partial$ występują w przeciwnych
kolejnościach. Na cyklach drugi składnik
znika, a pierwszy jest brzegiem.

\begin{twierdzenie}[Homotopia nie zmienia homologii]
\label{thm:homotopia-homologia}
Jeśli $f$ i $g$ są łańcuchowo homotopijne,
to $H_n(f)=H_n(g)$ dla każdego $n$.
W szczególności kompleks ściągalny jest
acykliczny. Jeśli istnieją odwzorowania łańcuchowe
$f:C_\bullet\to D_\bullet$ i
$g:D_\bullet\to C_\bullet$ z
$gf\simeq\id_C$ i $fg\simeq\id_D$,
to $H_n(f)$ i $H_n(g)$ są wzajemnymi
izomorfizmami.
\end{twierdzenie}
\begin{proof}
Dla $z\in Z_n(C)$ równanie~\eqref{eq:chain-homotopy}
redukuje się do
\[
 f_n(z)-g_n(z)=\partial^D_{n+1}h_n(z),
\]
bo $\partial^C_nz=0$. Różnica obrazów jest
brzegiem, więc $[f_n(z)]=[g_n(z)]$.
Gdy $f=\id_C$, $g=0$, otrzymujemy
$[z]=[0]$ dla każdej klasy, czyli
$H_n(C)=0$.
Na koniec, dzięki
Stwierdzeniu~\ref{prop:homologia-funktor},
\[
 H_n(g)H_n(f)=H_n(gf)
 =H_n(\id_C)=\id,\qquad
 H_n(f)H_n(g)=\id.
\]
Stąd oba odwzorowania na homologiach są
odwrotne do siebie.
\end{proof}

\begin{przyklad}[Najprostsze ściągnięcie]
Dla kompleksu
$0\to R\xrightarrow{\id}R\to0$
z modułami w stopniach $1$ i $0$
połóżmy $h_0:R=C_0\to C_1=R$
równe identyczności, a pozostałe
$h_n=0$. W stopniu $0$ mamy
$\partial_1h_0=\id_{C_0}$, a w stopniu
$1$ mamy $h_0\partial_1=\id_{C_1}$.
Wzór~\eqref{eq:chain-homotopy}
daje $\id_C=\partial h+h\partial$.
\end{przyklad}

\begin{twierdzenie}[Algebraiczne skasowanie pary generatorów]
\label{thm:kasowanie-pary}
Niech $C_\bullet$ będzie kompleksem
wolnych $R$-modułów z ustalonymi
bazami, a $n\geq1$. Załóżmy, że
przy generatorze $f\in C_{n-1}$
we współrzędnych $\partial_n(e)$ dla pewnego
generatora $e\in C_n$ stoi skalar
odwracalny $u\in R$.
Wtedy przez zmiany baz można rozłożyć
$C_\bullet$ na sumę prostą
podkompleksu i dwuwyrazowego kompleksu
\[
 0\longrightarrow R e
 \xrightarrow{\ \partial_n\ }R\,\partial_n(e)
 \longrightarrow0,
\]
w którym $\partial_n$ jest izomorfizmem.
Usunięcie tej pary nie zmienia homologii.
\end{twierdzenie}
\begin{proof}
Zapiszmy
$C_n=R e\oplus M$ oraz
$C_{n-1}=R f\oplus N$, gdzie
$M,N$ są rozpięte przez pozostałe
wybrane wektory bazowe.
Z założenia
$\partial_n(e)=u f+v$ z $v\in N$.
Ponieważ $u$ jest odwracalny,
$f':=\partial_n(e)$ może zastąpić $f$
w bazie $C_{n-1}$:
$C_{n-1}=R f'\oplus N$.
Rzeczywiście, $f=u^{-1}(f'-v)$
pozwala przejść w obie strony
między rozkładami.

Dla $x\in M$ zapiszmy
$\partial_n(x)=a(x)f+w(x)$,
gdzie $a(x)\in R$ i $w(x)\in N$.
Zmieńmy też bazę w stopniu $n$,
zastępując każdy bazowy $x\in M$
wektorem $x':=x-u^{-1}a(x)e$.
Wtedy
\[
 \partial_n(x')
 =a(x)f+w(x)-u^{-1}a(x)(uf+v)
 =w(x)-u^{-1}a(x)v\in N.
\]
Nowe wektory $x'$ razem z $e$
tworzą bazę $C_n$: możemy odzyskać
$x=x'+u^{-1}a(x)e$.
Oznaczmy ich rozpiętość przez $M'$.
Mamy zatem $\partial_n(Re)=Rf'$
i $\partial_n(M')\subseteq N$.

Równość $\partial_{n-1}(f')=\partial_{n-1}\partial_n(e)=0$
mówi, że z $Rf'$ nie wychodzi
różniczka. Z kolei dla $y\in C_{n+1}$
zapiszmy $\partial_{n+1}(y)=t e+x'$,
$x'\in M'$. Po zastosowaniu $\partial_n$
dostajemy
$0=\partial_n\partial_{n+1}(y)=t f'+\partial_n(x')$.
Pierwszy składnik należy do $Rf'$,
drugi do $N$; rozkład jest prosty,
więc $t=0$. Żadna różniczka
z wyższego stopnia nie trafia do $Re$.
Pozostałe moduły wraz z $M'$ i $N$
tworzą więc podkompleks, a
$Re\xrightarrow{\cong}Rf'$
jest jego prostym składnikiem
w całym kompleksie.

Na tym dwuwyrazowym składniku
homotopia $h_{n-1}(f')=e$,
zerowa w innych stopniach,
spełnia $\id=\partial h+h\partial$.
Jest on ściągalny. Przedłużmy tę homotopię przez zero na
pozostałym podkompleksie $D$. Dla rzutu $p:C\to D$ i włączenia
$j:D\to C$ otrzymujemy $pj=\id_D$ oraz
$\id_C-jp=\partial h+h\partial$. Rzut i włączenie
wyznaczają odwrotne izomorfizmy
na homologiach na mocy
Twierdzenia~\ref{thm:homotopia-homologia}.
\end{proof}

Nad $\Z$ jednostkami są tylko $1$ i $-1$.
Współczynnik $2$ w różniczce
$\Z\xrightarrow{\times2}\Z$ nie pozwala
wykonać takiego skasowania: pozostała
klasa $\Z/2$ z
Przykładu~\ref{prz:homologia-rachunki}
jest rzeczywistą przeszkodą.

\begin{stwierdzenie}[Acykliczny kompleks przestrzeni liniowych]
\label{prop:acykliczny-sciagalny}
Jeżeli $C_\bullet$ jest kompleksem
przestrzeni liniowych nad ciałem
$\mathbb K$ i $H_n(C)=0$ dla wszystkich
$n$, to $C_\bullet$ jest ściągalny.
Wystarczy, aby każda podprzestrzeń
$Z_n(C)\subseteq C_n$ miała dopełnienie;
warunek ten zachodzi zawsze przy
skończonych wymiarach.
\end{stwierdzenie}
\begin{proof}
Acykliczność daje $Z_n=B_n$.
Wybierzmy dopełnienie $S_n$ tak, aby
$C_n=Z_n\oplus S_n$. Ograniczenie
\[
 \partial_n|_{S_n}:S_n\longrightarrow Z_{n-1}
\]
jest iniektywne: element jego jądra
należy do $S_n\cap Z_n=\{0\}$.
Jest też surjektywne, bo
$\partial_n(C_n)=B_{n-1}=Z_{n-1}$
i $\partial_n$ znika na $Z_n$.
Ma więc odwrotność.
Każde $x\in C_n$ zapiszmy
jednoznacznie jako $x=z+s$
z $z\in Z_n$, $s\in S_n$, i połóżmy
\[
 h_n(z+s)
 :=(\partial_{n+1}|_{S_{n+1}})^{-1}(z)
 \in S_{n+1}.
\]
Wówczas $\partial_{n+1}h_n(z+s)=z$.
Z drugiej strony $\partial_n(z+s)=\partial_n(s)$,
więc
$h_{n-1}\partial_n(z+s)
=(\partial_n|_{S_n})^{-1}(\partial_ns)=s$.
Zatem $(\partial h+h\partial)(x)=z+s=x$
w każdym stopniu.
\end{proof}

\begin{przyklad}[Acykliczny, lecz nieściągalny nad $\Z$]
\label{prz:acykliczny-niesciagalny}
Potraktujmy dokładny ciąg z
Przykładu~\ref{prz:ses-mod-m} dla $m=2$
jako kompleks w stopniach $2,1,0$:
\[
 0\longrightarrow\Z\xrightarrow{\ \times2\ }\Z
 \xrightarrow{\ \pi\ }\Z/2\Z\longrightarrow0.
\]
Jego wszystkie homologie są zerowe,
bo $\times2$ jest iniekcją,
$\ker\pi=2\Z$, a $\pi$ jest
surjekcją. Gdyby kompleks był
ściągalny, to w stopniu $0$
wzór $\id=\partial h+h\partial$ dałby
$\id_{\Z/2}=\pi h_0$.
Odwzorowanie $h_0:\Z/2\to\Z$ byłoby
przekrojem $\pi$, którego
nie ma według uwagi po
Twierdzeniu~\ref{thm:ses-rozszczepienie}.
Zatem implikacja \emph{acykliczny
$\Rightarrow$ ściągalny} wymaga
dodatkowych założeń. Implikacja
odwrotna z
Twierdzenia~\ref{thm:homotopia-homologia}
działa nad każdym pierścieniem.
\end{przyklad}

\section{Ilorazy, ciągi względne i filtracje}
\label{sec:iloraz-filtracja}

Podkompleks $A_\bullet\subseteq C_\bullet$
oznacza podmoduły $A_n\subseteq C_n$
we wszystkich stopniach, takie że
$\partial_C(A_n)\subseteq A_{n-1}$.
Definiujemy \emph{kompleks ilorazowy}
$(C/A)_n=C_n/A_n$ i
$\partial_{C/A}[c]=[\partial_Cc]$.
Jeśli $[c]=[c+a]$ dla $a\in A_n$,
to $[\partial_Cc]=[\partial_Cc+\partial_Ca]$;
różniczka ilorazowa jest dobrze
określona i jej kwadrat wynosi zero.

\begin{stwierdzenie}[Długi ciąg ilorazu]
\label{prop:ses-iloraz}
Dla podkompleksu $A_\bullet\subseteq C_\bullet$
ciąg
\[
 0\longrightarrow A_\bullet
 \longrightarrow C_\bullet
 \longrightarrow C_\bullet/A_\bullet
 \longrightarrow0
\]
jest stopniowo dokładny. Dlatego istnieje
długi ciąg
\begin{equation}\label{eq:les-iloraz}
 \cdots\to H_n(A)\to H_n(C)\to H_n(C/A)
 \xrightarrow{\delta_n}H_{n-1}(A)\to\cdots.
\end{equation}
Jeśli $\xi\in H_n(C/A)$, to możemy
reprezentować ją elementem $c\in C_n$
z $\partial_Cc\in A_{n-1}$, a wtedy
$\delta_n\xi=[\partial_Cc]$. Tutaj $c$ najpierw wyznacza cykl
$c+A_n$ w kompleksie ilorazowym, a dopiero jego klasa modulo
brzegi jest elementem $\xi$; są to dwa różne przejścia do ilorazu.
\end{stwierdzenie}
\begin{proof}
W każdym stopniu włączenie $A_n\to C_n$
jest iniekcją, projekcja $C_n\to C_n/A_n$
jest surjekcją, a jej jądrem jest $A_n$.
Oba odwzorowania komutują z różniczkami
z definicji podkompleksu i ilorazu.
Twierdzenie~\ref{thm:les-homologia}
stosuje się więc bez dodatkowych
założeń. Warunek
$\partial_{C/A}[c]=[0]$ jest równoważny
$\partial_Cc\in A_{n-1}$; w tej sytuacji
konstrukcja~\eqref{eq:delta-wzor}
przy podniesieniu $c$ daje dokładnie
klasę $[\partial_Cc]$.
\end{proof}

Symbol $H_n(C,A):=H_n(C/A)$ będziemy
nazywać \emph{względną homologią
algebraiczną} kompleksu $C$ względem
podkompleksu $A$. Jest to definicja
czysto algebraiczna. W przyszłych
konstrukcjach topologicznych trzeba
osobno sprawdzić, że odpowiednie
łańcuchy tworzą podkompleks.

\begin{stwierdzenie}[Względna kohomologia nad ciałem]
\label{prop:wzgledna-kohomologia}
Jeśli $A_\bullet\subseteq C_\bullet$
jest podkompleksem przestrzeni liniowych
nad ciałem $\mathbb K$, to wzięcie
funkcjonałów daje krótki ciąg dokładny
kompleksów kołańcuchowych
\[
 0\to\operatorname{Hom}_{\mathbb K}
      (C/A,\mathbb K)
 \to\operatorname{Hom}_{\mathbb K}(C,\mathbb K)
 \to\operatorname{Hom}_{\mathbb K}(A,\mathbb K)
 \to0.
\]
Jeśli $H^n(C,A;\mathbb K)$ oznacza
kohomologię pierwszego z nich, to mamy
długi ciąg
\[
 \cdots\to H^n(C,A;\mathbb K)\to H^n(C;\mathbb K)
 \to H^n(A;\mathbb K)\to H^{n+1}(C,A;\mathbb K)
 \to\cdots .
\]
W tym zapisie $H^n(C;\mathbb K)$
oznacza kohomologię kompleksu
$\operatorname{Hom}_{\mathbb K}(C,\mathbb K)$,
a analogicznie dla $A$.
\end{stwierdzenie}
\begin{proof}
W stopniu $n$ pierwsze odwzorowanie składa
funkcjonał $C_n/A_n\to\mathbb K$
z projekcją $C_n\to C_n/A_n$,
więc jest iniektywne.
Drugie ogranicza funkcjonał na $C_n$
do podprzestrzeni $A_n$.
Jego jądrem są dokładnie funkcjonały
znikające na $A_n$, czyli faktoryzujące
się przez $C_n/A_n$; to obraz
pierwszego odwzorowania. Każdy funkcjonał
na $A_n$ przedłuża się do $C_n$:
wybieramy dopełnienie $A_n$
i stawiamy na nim wartość $0$.
Drugie odwzorowanie jest zatem surjektywne.
Odwzorowania komutują z różniczkami
kołańcuchowymi, bo są zdefiniowane
przez składanie z odwzorowaniami łańcuchowymi.
Stosujemy
Twierdzenie~\ref{thm:les-kohomologia}.
\end{proof}

\begin{stwierdzenie}[Trzy zagnieżdżone kompleksy]
\label{prop:les-potrojny}
Jeśli $A_\bullet\subseteq B_\bullet
\subseteq C_\bullet$, to istnieje
krótki ciąg dokładny kompleksów
\[
 0\longrightarrow B_\bullet/A_\bullet
 \longrightarrow C_\bullet/A_\bullet
 \longrightarrow C_\bullet/B_\bullet
 \longrightarrow0
\]
i wynikający z niego długi ciąg
\[
 \cdots\to H_n(B,A)\to H_n(C,A)\to H_n(C,B)
 \to H_{n-1}(B,A)\to\cdots.
\]
\end{stwierdzenie}
\begin{proof}
W stopniu $n$ pierwsze odwzorowanie
$[b]_A\mapsto[b]_A$ jest iniektywne:
jeśli $[b]_A=0$ w $C_n/A_n$, to
$b\in A_n$, więc była już zerowa w
$B_n/A_n$. Ostatnie odwzorowanie
$[c]_A\mapsto[c]_B$ jest dobrze
określone, bo $A_n\subseteq B_n$,
i surjektywne, bo każda klasa modulo
$B_n$ ma reprezentanta $c\in C_n$.
Jej jądro składa się z klas
$[c]_A$ z $c\in B_n$, czyli jest
obrazem pierwszego odwzorowania.
Różniczki zachowują wszystkie trzy
podkompleksy, więc są to odwzorowania
kompleksów. Stosujemy
Twierdzenie~\ref{thm:les-homologia}.
\end{proof}

\begin{definicja}[Filtracja kompleksu]
\label{def:filtracja-kompleksu}
Przez \emph{skończoną filtrację rosnącą} rozumiemy tutaj ciąg
podkompleksów
\[
 0=F_{-1}C_\bullet\subseteq
 F_0C_\bullet\subseteq\cdots\subseteq
 F_NC_\bullet=C_\bullet.
\]
Iloraz $F_pC/F_{p-1}C$ opisuje
algebraicznie \emph{nową warstwę}
w kroku $p$. W każdym kroku mamy
krótki ciąg dokładny
$0\to F_{p-1}C\to F_pC\to
F_pC/F_{p-1}C\to0$
i odpowiadający mu długi ciąg homologii.
\end{definicja}

\begin{wniosek}[Zmiana skoncentrowana w jednym stopniu]
\label{wn:filtracja-jeden-stopien}
Niech $A\subseteq B$ będzie podkompleksem
w nieujemnych stopniach i niech
$H_j(B/A)=0$ dla $j\ne q$, $q\geq1$.
Wówczas $H_j(A)\to H_j(B)$ jest
izomorfizmem dla $j\notin\{q,q-1\}$,
a pozostałe zmiany opisuje ciąg
\begin{equation}\label{eq:jedna-warstwa}
 0\to H_q(A)\to H_q(B)\to H_q(B/A)
 \xrightarrow{\delta_q}H_{q-1}(A)
 \to H_{q-1}(B)\to0.
\end{equation}
\end{wniosek}
\begin{proof}
Dla dowolnego $j$ wytnijmy z~\eqref{eq:les-iloraz}
fragment
\[
 H_{j+1}(B/A)\to H_j(A)\to H_j(B)\to H_j(B/A).
\]
Jeśli oba zewnętrzne moduły są zerowe,
środkowe odwzorowanie jest jednocześnie
iniektywne i surjektywne.
Tak jest, gdy $j\ne q,q-1$.
Dla $j=q$ i $j=q-1$ zachowujemy
pełen fragment długiego ciągu.
Jego lewa skrajna strzałka wychodzi
z $H_{q+1}(B/A)=0$, a prawa
wchodzi do $H_{q-1}(B/A)=0$,
co daje dokładnie~\eqref{eq:jedna-warstwa}.
\end{proof}

\begin{uwaga}[Co mówi odwzorowanie łączące?]
W~\eqref{eq:jedna-warstwa} element
nowej warstwy może zwiększyć
$H_q(B)$, jeśli ma zerowy obraz
przez $\delta_q$. Może też sprawić,
że wcześniejsza klasa w
$H_{q-1}(A)$ stanie się brzegiem:
tak dzieje się, jeśli leży w
$\operatorname{im}\delta_q$.
Ten jeden fragment ciągu odróżnia
\emph{powstanie klasy} od
\emph{zabicia klasy}, bez odwoływania
się do szczególnej realizacji
geometrycznej kompleksu.
\end{uwaga}
\section{Suma podkompleksów i ciąg Mayera--Vietorisa}
\label{sec:mv-algebraiczny}

Przy rozbiciu obiektu na dwie części trzeba
uwzględnić elementy wspólne: bez tego
zliczylibyśmy je dwa razy. Poniższy wynik
jest czysto algebraicznym źródłem
późniejszych ciągów Mayera--Vietorisa.
Nie wymaga jeszcze żadnej konstrukcji
łańcuchów z przestrzeni topologicznej.

\begin{twierdzenie}[Algebraiczny Mayer--Vietoris]
\label{thm:mv-algebraiczny}
Niech $A_\bullet$ i $B_\bullet$ będą
podkompleksami $C_\bullet$. Wtedy
$A\cap B$ i $A+B$ też są podkompleksami,
a ciąg
\begin{equation}\label{eq:mv-ses}
 0\longrightarrow A\cap B
 \xrightarrow{\ j\ }A\oplus B
 \xrightarrow{\ q\ }A+B
 \longrightarrow0,
 \quad j(x)=(x,-x),\quad q(a,b)=a+b
\end{equation}
jest krótkim ciągiem dokładnym kompleksów.
W szczególności mamy długi ciąg dokładny
\begin{equation}\label{eq:mv-les}
 \cdots\to H_n(A\cap B)\to H_n(A)\oplus H_n(B)
 \to H_n(A+B)\xrightarrow{\delta}
 H_{n-1}(A\cap B)\to\cdots.
\end{equation}
Jeśli cykl $c\in(A+B)_n$ zapiszemy
$c=a+b$ z $a\in A_n$, $b\in B_n$,
to $\delta[c]=[\partial a]=[-\partial b]$.
\end{twierdzenie}
\begin{proof}
Jeśli $x\in A_n\cap B_n$, to
$\partial x\in A_{n-1}$ i $\partial x\in B_{n-1}$,
więc $\partial$ zachowuje przecięcie.
Podobnie $\partial(a+b)=\partial a+\partial b$ należy do
$A_{n-1}+B_{n-1}$, zatem suma
też jest podkompleksem.
Różniczka w $A\oplus B$ działa
składowo, dlatego
$\partial\,j(x)=(\partial x,-\partial x)=j(\partial x)$ i
$\partial\,q(a,b)=\partial(a+b)=q(\partial a,\partial b)$.
Są to odwzorowania łańcuchowe.

W każdym stopniu odwzorowanie $j$ jest
iniektywne, bo $j(x)=(0,0)$
implikuje $x=0$. Odwzorowanie $q$ jest
surjektywne na $A_n+B_n$ z samej
definicji sumy. Jeśli $q(a,b)=0$,
to $a=-b$ należy jednocześnie do
$A_n$ i do $B_n$, a para ma postać
$(a,-a)=j(a)$. Stąd
$\ker q=\operatorname{im}j$.
Twierdzenie~\ref{thm:les-homologia}
daje~\eqref{eq:mv-les}. Używamy przy tym kanonicznego izomorfizmu
$H_n(A\oplus B)\cong H_n(A)\oplus H_n(B)$:
różniczka działa składowo, więc cykle i brzegi są odpowiednio
$Z_n(A)\oplus Z_n(B)$ i $B_n(A)\oplus B_n(B)$,
a $[(a,b)]\mapsto([a],[b])$ jest izomorfizmem ich ilorazów.

Na koniec, jeśli $\partial(a+b)=0$, to
$\partial a=-\partial b$. Lewa strona należy do
$A_{n-1}$, prawa do $B_{n-1}$,
więc obie są elementem $A\cap B$.
W konstrukcji odwzorowania łączącego
podniesienie cyklu $c$ jest parą $(a,b)$,
a jej brzeg $(\partial a,\partial b)=(\partial a,-\partial a)$
jest obrazem $j(\partial a)$. To daje
podany wzór i objaśnia znak.
\end{proof}

\begin{przyklad}[Niezerowy brzeg z przecięcia]
\label{prz:mv-niezerowa-delta}
Nad ciałem $\mathbb K$ weźmy kompleks
z $C_1=\mathbb K u\oplus\mathbb K v$,
$C_0=\mathbb K x\oplus\mathbb K y$,
$\partial(u)=x-y$ i $\partial(v)=y-x$,
zerowy w pozostałych stopniach.
Niech $A_1=\mathbb K u$, $B_1=\mathbb K v$,
ale $A_0=B_0=C_0$. Są to podkompleksy,
$A+B=C$, zaś $A\cap B$ jest
skoncentrowany w stopniu $0$
i ma tam moduł $\mathbb K x\oplus\mathbb K y$.
Element $u+v$ jest cyklem w $C_1$.
Jego podniesienie do $A_1\oplus B_1$ to $(u,v)$,
którego brzeg wynosi
$(x-y,y-x)=j(x-y)$. Stąd
\[
 \delta[u+v]=[x-y]
 \in H_0(A\cap B)=\mathbb K x\oplus\mathbb K y,
\]
klasa niezerowa. Widzimy, jak brzeg
jednej części zostaje zniesiony
przez brzeg drugiej i powstaje
cykl w ich sumie.
\end{przyklad}

\section{Współczynniki, torsja i charakterystyka Eulera}
\label{sec:wspolczynniki-euler}

Wybór pierścienia $R$ wpływa na
wynik homologiczny. Kompleks
wolnych grup abelowych
można sprowadzić modulo $m$,
lecz homologia po tej operacji
nie polega zawsze wyłącznie
na redukcji dawnej homologii:
dochodzi torsja z poprzedniego stopnia.

\begin{twierdzenie}[Redukcja współczynników modulo $m$]
\label{thm:mod-m-homologia}
Niech $C_\bullet$ będzie kompleksem
wolnych grup abelowych, a $m\geq2$.
Różniczki wyznaczają kompleks
$(C/mC)_n=C_n/mC_n$. Dla każdego
$n$ istnieje krótki ciąg dokładny
\begin{equation}\label{eq:mod-m-homologia}
 0\longrightarrow H_n(C)/mH_n(C)
 \longrightarrow H_n(C/mC)
 \longrightarrow H_{n-1}(C)[m]
 \longrightarrow0,
\end{equation}
gdzie $H_{n-1}(C)[m]:=
\{x\in H_{n-1}(C):mx=0\}$
to $m$-torsja z Definicji~\ref{def:torsja}.
Nie twierdzimy, że ten ciąg ma
kanoniczne rozszczepienie.
\end{twierdzenie}
\begin{proof}
Wolna grupa abelowa nie ma
niezerowych elementów zabijanych
przez $m$. W każdym stopniu mamy
zatem dokładny ciąg
\[
 0\to C_n\xrightarrow{\ \times m\ }C_n
 \xrightarrow{\ \pi\ }C_n/mC_n\to0.
\]
Ponieważ $\partial(mx)=m\,\partial(x)$,
odwzorowania $\times m$ i $\pi$
komutują z różniczkami.
To krótki ciąg dokładny
kompleksów, więc
Twierdzenie~\ref{thm:les-homologia}
daje fragment
\[
 H_n(C)\xrightarrow{\ \times m\ }H_n(C)
 \xrightarrow{\ \pi_*\ }H_n(C/mC)
 \xrightarrow{\ \delta\ }H_{n-1}(C)
 \xrightarrow{\ \times m\ }H_{n-1}(C).
\]
Dokładność mówi
$\ker\pi_*=mH_n(C)$, a więc
$\pi_*$ wyznacza iniekcję
$H_n(C)/mH_n(C)\hookrightarrow H_n(C/mC)$
o obrazie $\ker\delta$.
Mówi też
$\operatorname{im}\delta
=\ker(\times m:H_{n-1}(C)\to H_{n-1}(C))
=H_{n-1}(C)[m]$.
Odwzorowanie $\delta$ po ograniczeniu
przeciwdziedziny do tego jądra
jest więc surjektywne.
To dokładnie trzy warunki
ciągu~\eqref{eq:mod-m-homologia}.
Warto zapisać ostatnią strzałkę jawnie. Jeśli $\bar c$ jest cyklem
modulo $m$ i $c\in C_n$ jest jego przedstawicielem, to
$\partial c=ma$ dla pewnego $a\in C_{n-1}$.
Element $a$ jest jednoznaczny, bo mnożenie przez $m$ jest iniektywne.
Ponadto $m\partial a=0$ wymusza $\partial a=0$,
a $m[a]=[\partial c]=0$. Zatem strzałka wysyła $[\bar c]$ na $[a]$.
Jest naturalna względem odwzorowań łańcuchowych, gdyż
$\partial f(c)=mf(a)$; pierwsza strzałka wysyła klasę cyklu
$[z]$ modulo $mH_n(C)$ na klasę $z$ modulo $m$.
\end{proof}

\begin{przyklad}[Nowa homologia po redukcji]
Dla $C_1=C_0=\Z$ i $\partial_1=\times2$
otrzymaliśmy $H_1(C)=0$ oraz
$H_0(C)=\Z/2$. Po redukcji modulo
$2$ różniczka $\times2$ staje się
zerowa, a więc
$H_1(C/2C)\cong H_0(C/2C)\cong\Z/2$.
Nowa klasa stopnia $1$ pochodzi
w~\eqref{eq:mod-m-homologia} z
$H_0(C)[2]=\Z/2$, choć wcześniejsze
$H_1(C)$ było zerowe.
\end{przyklad}

\begin{uwaga}[Iloczyn tensorowy modułów i rozszerzenie skalarów]
\label{uwaga:tensor-modulow}
W konstrukcji algebry tensorowej tensorowaliśmy przestrzenie
nad ciałem. Potrzebny teraz $M\otimes_R N$ dla modułów nad
pierścieniem przemiennym $R$ powstaje podobnie: bierzemy wolny
$R$-moduł na symbolach $[x,y]$, $x\in M$, $y\in N$, i dzielimy
przez podmoduł generowany przez relacje addytywności w obu argumentach
oraz
\[
 [rx,y]-r[x,y],\qquad [x,ry]-r[x,y].
\]
Klasę $[x,y]$ oznaczamy $x\otimes y$. Każde odwzorowanie
$R$-dwuliniowe z $M\times N$ do modułu $P$ przedłuża się liniowo
na moduł wolny i zeruje te relacje, więc przechodzi jednoznacznie
na homomorfizm $M\otimes_R N\to P$. To dowodzi własności uniwersalnej;
daje też jednoznaczność konstrukcji przez homomorfizmy w obie strony,
ustalone na tensorach prostych.

Jeśli $S$ jest pierścieniem przemiennym z ustalonym homomorfizmem
$R\to S$ zachowującym jedynkę (czyli przemienną $R$-algebrą),
na $M\otimes_R S$ mamy
działanie $s(x\otimes a)=x\otimes sa$; zachowuje ono relacje,
więc określa $S$-moduł. Dla modułu wolnego
$M=\bigoplus_i Re_i$ dostajemy
$M\otimes_R S\cong\bigoplus_i Se_i$:
$ (\sum_i r_i e_i)\otimes a\mapsto\sum_i(r_i a)e_i$,
a odwrotność wysyła $\sum_i a_i e_i$ na $\sum_i e_i\otimes a_i$.
Wszystkie te sumy są skończone.
Ta sama reguła, zastosowana do rzutów i włączeń, pokazuje,
że tensorowanie zachowuje sumy proste. Homomorfizm $f$ indukuje
$f\otimes\id$ przez $x\otimes a\mapsto f(x)\otimes a$;
izomorfizm pozostaje izomorfizmem, z odwrotnością
$f^{-1}\otimes\id$. Nie wnioskujemy stąd, że tensorowanie
zachowuje dowolne iniekcje.
\end{uwaga}

\begin{definicja}[Homologia ze współczynnikami w $\Q$, $\R$ lub $\C$]
\label{def:homologia-ciala-char-zero}
Niech $(C_\bullet,\partial)$ będzie kompleksem
wolnych grup abelowych, na przykład
kompleksem łańcuchów symplicjalnych.
Dla $\mathbb F\in\{\Q,\R,\C\}$ tworzymy
\[
 C_n(C;\mathbb F):=C_n\otimes_{\Z}\mathbb F,
 \qquad
 \partial_{\mathbb F}(c\otimes a)
       :=(\partial c)\otimes a.
\]
Jeśli $C_n$ ma bazę sympleksów, nowy kompleks
ma tę samą bazę, lecz można przy nich stawiać
współczynniki z $\mathbb F$ zamiast liczb całkowitych.
Jego homologię oznaczamy $H_n(C;\mathbb F)$;
dla kompleksu przypisanego przestrzeni $X$
piszemy $H_n(X;\mathbb F)$.
Piszemy też $H_n(C;\Z)=H_n(C)$ dla pierwotnego kompleksu.
Różniczka jest dobrze określona przez własność uniwersalną,
a $\partial_{\mathbb F}^2=0$ sprawdzamy na tensorach prostych.
\end{definicja}

\begin{twierdzenie}[Co zmienia przejście do charakterystyki zero]
\label{thm:homologia-char-zero}
Dla każdego kompleksu wolnych grup abelowych
i $\mathbb F=\Q,\R$ lub $\C$ istnieje naturalny
izomorfizm
\begin{equation}\label{eq:homologia-char-zero}
 H_n(C;\mathbb F)
 \cong H_n(C;\Z)\otimes_{\Z}\mathbb F.
\end{equation}
Jeśli $H_n(C;\Z)\cong\Z^{b_n}\oplus T_n$,
gdzie $T_n$ jest skończoną grupą torsyjną, to
$H_n(C;\mathbb F)\cong\mathbb F^{b_n}$.
W szczególności w tych trzech ciałach
pozostaje ta sama liczba $b_n$ niezależnych
klas, a torsja znika.
\end{twierdzenie}
\begin{proof}
Najpierw rozważmy $\mathbb F=\Q$.
Element $C_n\otimes_{\Z}\Q$ jest skończoną
sumą prostych tensorów; sprowadzając
mianowniki do wspólnego, zapisujemy go jako
$c/d$ z $c\in C_n$ i niezerowym
$d\in\Z$. Ponieważ każda grupa $C_j$
jest wolna, odwzorowania
$C_j\to C_j\otimes_{\Z}\Q$ są iniektywne.
Jeśli $c/d$ jest cyklem, to
$\partial c/d=0$, a więc $\partial c=0$.
Każda klasa racjonalna pochodzi zatem
całkowitego cyklu po dopuszczeniu mianownika:
odwzorowanie
\[
 H_n(C;\Z)\otimes_{\Z}\Q\longrightarrow H_n(C;\Q),
 \qquad [c]\otimes q\longmapsto[c\otimes q]
\]
jest surjektywne. Jest ono dobrze określone,
bo cykl pozostaje cyklem, brzeg pozostaje
brzegiem, a reguła jest dwuliniowa nad $\Z$.
Każdy tensor w dziedzinie ma postać $[c]\otimes(1/d)$
po sprowadzeniu skończenie wielu ułamków do wspólnego mianownika.

Sprawdźmy, czy nie utożsamia zbyt wielu klas.
Załóżmy, że cykl $c/d$ jest brzegiem
w kompleksie racjonalnym. Istnieją
$b\in C_{n+1}$ i $e\ne0$ takie, że
$c/d=\partial b/e$. Po pomnożeniu przez
$de$ i użyciu iniektywności
$C_n\to C_n\otimes\Q$ otrzymujemy
$ec=d\,\partial b=\partial(db)$
już w $C_n$. To znaczy, że $e[c]=0$
w $H_n(C;\Z)$. Po przejściu do $\Q$
liczba $e$ jest odwracalna, więc
$[c]\otimes(1/d)=0$. Odwzorowanie jest
również iniektywne i wzór został
udowodniony dla $\Q$.

Teraz $\mathbb F=\R$ albo $\C$.
Są to przestrzenie liniowe nad $\Q$, a
rozszerzenie skalarów daje izomorfizm
kompleksów
\[
 C_\bullet\otimes_{\Z}\mathbb F
 \cong (C_\bullet\otimes_{\Z}\Q)
             \otimes_{\Q}\mathbb F.
\]
Jawnie, $(c\otimes q)\otimes a\mapsto c\otimes qa$,
a odwrotnie $c\otimes a\mapsto(c\otimes1)\otimes a$.
Relacje tensorowe pokazują, że oba wzory są dobrze określone
i wzajemnie odwrotne; komutują też z różniczkami.
Wyjaśnijmy, dlaczego tensorowanie kompleksu
przestrzeni liniowych nad ciałem zachowuje
homologię. Oznaczmy na chwilę
$D_\bullet=C_\bullet\otimes_{\Z}\Q$,
$Z_n^\Q=\ker(\partial:D_n\to D_{n-1})$
oraz $B_n^\Q=\operatorname{im}(\partial:D_{n+1}\to D_n)$.
W każdym stopniu uzupełniamy bazy, aby zapisać
$Z_n^\Q=B_n^\Q\oplus V_n$ i
$D_n=Z_n^\Q\oplus L_n$.
Wtedy $\partial:L_n\to B_{n-1}^\Q$ jest
izomorfizmem: surjektywność wynika z definicji
$B_{n-1}^\Q$, a jądro to
$L_n\cap Z_n^\Q=0$. Ponadto
$V_n\cong Z_n^\Q/B_n^\Q=H_n(D)$.
Po tensorowaniu nad $\Q$ rozkłady nadal
są sumami prostymi, a ograniczenie
$\partial\otimes1:L_n\otimes\mathbb F
\to B_{n-1}^\Q\otimes\mathbb F$ pozostaje
izomorfizmem. Nowa homologia to więc
$V_n\otimes\mathbb F
\cong H_n(D)\otimes_{\Q}\mathbb F$.
Łącząc to z przypadkiem $\Q$, otrzymujemy
\eqref{eq:homologia-char-zero} dla $\R,\C$.
Choć dopełnienia ułatwiły dowód bijektywności, sam izomorfizm
jest kanoniczny: we wszystkich trzech przypadkach jest zadany
przez $[z]\otimes a\mapsto[z\otimes a]$.
Równość $f(z)\otimes a=(f\otimes\id)(z\otimes a)$ dowodzi
naturalności względem odwzorowań łańcuchowych.

Wreszcie $\Z\otimes_{\Z}\mathbb F\cong\mathbb F$.
Jeśli $mx=0$ w $T_n$, to
$m(x\otimes a)=0$; ponieważ $m\ne0$
ma odwrotność w $\mathbb F$,
$x\otimes a=0$. Zatem $T_n\otimes\mathbb F=0$,
a $b_n$ generatorów części wolnej daje
$\mathbb F^{b_n}$.
\end{proof}

Warto porównać to z Twierdzeniem~\ref{thm:mod-m-homologia}.
W ciele $\Q$, $\R$ lub $\C$ każda niezerowa
liczba całkowita staje się odwracalna.
Przy współczynnikach modulo $p$ liczba $p$
staje się natomiast \emph{zerem}; torsja
może wtedy dać nowe klasy także w innym stopniu.
Przejście do $\R$ lub $\C$ po obliczeniu
homologii nad $\Q$ jest już tylko
rozszerzeniem skalarów:
\[
 H_n(C;\R)\cong H_n(C;\Q)\otimes_{\Q}\R,
 \qquad
 H_n(C;\C)\cong H_n(C;\R)\otimes_{\R}\C.
\]

\begin{twierdzenie}[Algebraiczna charakterystyka Eulera]
\label{thm:euler-homologia}
Niech $C_\bullet$ będzie kompleksem
skończenie wymiarowych przestrzeni
liniowych nad ciałem $\mathbb K$
i niech tylko skończenie wiele $C_n$
będzie niezerowych. Wtedy
\begin{equation}\label{eq:euler-homologia}
 \sum_n(-1)^n\dim_{\mathbb K}C_n
 =\sum_n(-1)^n\dim_{\mathbb K}H_n(C).
\end{equation}
Lewą stronę nazywamy \emph{charakterystyką
Eulera kompleksu}.
\end{twierdzenie}
\begin{proof}
Odwzorowanie $\partial_n:C_n\to B_{n-1}$ jest
surjektywne z definicji obrazu,
a jego jądro to $Z_n$.
Wybór bazy $Z_n$ i jej uzupełnienie
do bazy $C_n$ pokazują
\[
 \dim C_n=\dim Z_n+\dim B_{n-1}.
\]
Podobnie projekcja $Z_n\to H_n=Z_n/B_n$
ma jądro $B_n$, więc
$\dim Z_n=\dim B_n+\dim H_n$.
Łącząc te równości, otrzymujemy
\[
 \dim C_n=\dim B_n+\dim H_n+\dim B_{n-1}.
\]
Mnożymy przez $(-1)^n$ i sumujemy
po $n$. Suma
$\sum_n(-1)^n\dim B_n$
znosi się z
$\sum_n(-1)^n\dim B_{n-1}$:
po podstawieniu $j=n-1$
druga jest równa
$-\sum_j(-1)^j\dim B_j$.
Skończona liczba niezerowych
wyrazów pozwala to wykonać bez
kłopotów ze zbieżnością.
Pozostaje~\eqref{eq:euler-homologia}.
\end{proof}

\begin{przyklad}[Wykrywanie brakującego wymiaru]
Niech nad $\mathbb K$ mamy tylko
$C_2\cong\mathbb K^2$, $C_1\cong\mathbb K^5$
i $C_0\cong\mathbb K^4$.
Niezależnie od szczegółów różniczek
spełniających $\partial^2=0$, musi zachodzić
\[
 \dim H_0-\dim H_1+\dim H_2
 =4-5+2=1.
\]
Nie można zatem uzyskać
acyklicznego kompleksu z takimi
wymiarami. Charakterystyka Eulera
nie mówi, w którym stopniu leży
niezerowa klasa; daje jednak
szybki test każdego obliczenia.
\end{przyklad}

\section{Porównywanie długich ciągów}
\label{sec:lemat-pieciu}

Czasem umiemy opisać cztery składniki
dwóch długich ciągów, a piąty jest trudny.
Następny lemat zamienia dane o odwzorowaniach
sąsiednich w informację o odwzorowaniu środkowym.
Przyda się do porównywania kompleksów
i różnych modeli tej samej homologii.

\begin{lemat}[Lemat pięciu]
\label{lem:pieciu}
Załóżmy, że mamy przemienny diagram
homomorfizmów $R$-modułów o dokładnych
wierszach
\[
\begin{array}{ccccccccc}
A_1&\longrightarrow&A_2&\longrightarrow&A_3&
 \longrightarrow&A_4&\longrightarrow&A_5\\
\downarrow f_1&&\downarrow f_2&&\downarrow f_3&&
\downarrow f_4&&\downarrow f_5\\
B_1&\longrightarrow&B_2&\longrightarrow&B_3&
 \longrightarrow&B_4&\longrightarrow&B_5 .
\end{array}
\]
Jeśli $f_1,f_2,f_4,f_5$ są izomorfizmami,
to $f_3$ także jest izomorfizmem.
\end{lemat}
\begin{proof}
Oznaczmy kolejne poziome odwzorowania w górnym
wierszu przez $u_1,u_2,u_3,u_4$,
a w dolnym przez $v_1,v_2,v_3,v_4$.
Przemienność oznacza
$f_{j+1}u_j=v_jf_j$.

\emph{Iniektywność.} Niech $x\in A_3$
spełnia $f_3x=0$.
Wtedy
$f_4(u_3x)=v_3(f_3x)=0$;
iniektywność $f_4$ daje $u_3x=0$.
Z dokładności w $A_3$ mamy
$x=u_2y$ dla pewnego $y\in A_2$.
Ponieważ
$v_2f_2y=f_3u_2y=f_3x=0$,
dokładność w $B_2$ daje
$f_2y=v_1z$ dla pewnego $z\in B_1$.
Surjektywność $f_1$ pozwala napisać
$z=f_1w$ dla $w\in A_1$.
Stąd $f_2y=v_1f_1w=f_2u_1w$.
Iniektywność $f_2$ daje $y=u_1w$,
a więc $x=u_2u_1w=0$.

\emph{Surjektywność.} Weźmy $z\in B_3$.
Surjektywność $f_4$ daje $a_4\in A_4$
z $f_4a_4=v_3z$. Teraz
$f_5u_4a_4=v_4f_4a_4=v_4v_3z=0$.
Z iniektywności $f_5$ wynika
$u_4a_4=0$. Dokładność w $A_4$
daje $a_4=u_3x$ dla pewnego
$x\in A_3$. Różnica
$z-f_3x$ leży w $\ker v_3$,
bo $v_3(z-f_3x)=v_3z-f_4u_3x=0$.
Z dokładności w $B_3$ jest więc
równa $v_2y$ dla pewnego $y\in B_2$.
Surjektywność $f_2$ daje
$y=f_2w$ dla $w\in A_2$, a wtedy
$z=f_3x+v_2f_2w
=f_3(x+u_2w)$.
Otrzymaliśmy surjektywność $f_3$.
\end{proof}

\section*{Dalsza lektura}
\addcontentsline{toc}{section}{Dalsza lektura}
Algebraiczny długi ciąg homologii i kohomologii
oraz homotopie kompleksów są opisane w
\emph{The Stacks Project}, rozdział
\emph{Homological Algebra},
\url{https://stacks.math.columbia.edu/tag/010V}.
Zastosowania długich ciągów do geometrii
topologicznej rozwija Allen Hatcher,
\emph{Algebraic Topology}, rozdział 2,
\url{https://pi.math.cornell.edu/~hatcher/AT/AT.pdf}.
Definicje potrzebne w obecnym rozdziale
i dowody podstawowych twierdzeń podaliśmy
bez używania jeszcze modeli
symplicjalnych ani komórkowych.
\clearpage

\chapter{Homotopia, nakrycia i komórkowe modele przestrzeni}
\label{ch:topologia-algebraiczna-podstawy}

W Rozdziale~\ref{ch:algebra-homologiczna} kompleks łańcuchowy
był ciągiem modułów i homomorfizmów. Aby zbudować taki
kompleks z przestrzeni, nauczymy się rozkładać przestrzeń
na proste części. Równolegle zapytamy, kiedy dwa odwzorowania
można ciągle deformować jedno w drugie. Są to pytania
topologiczne: gładkość nie jest tu potrzebna.

Zaczniemy od dróg i nakryć, następnie od sympleksów
przejdziemy do CW-kompleksów. Homologia tych obiektów
pojawi się w następnym rozdziale. Homotopii odwzorowań
nie należy mylić z homotopią \emph{łańcuchową}
z Definicji~\ref{def:homotopia-lancuchowa};
ta druga jest algebraicznym odpowiednikiem pierwszej.

\section{Homotopie, drogi i grupa podstawowa}
\label{sec:homotopie-pi1}

W całym rozdziale $I=[0,1]$. Parametr $t\in I$
można interpretować jako czas ciągłej deformacji.
Obrazy pośrednie mogą być bardzo różne, lecz
ciągłość wyklucza skoki i rozdarcia.

\begin{definicja}[Homotopia i równoważność homotopijna]
\label{def:homotopia-topologiczna}
Dwa ciągłe odwzorowania $f,g:X\to Y$ są
\emph{homotopijne}, zapisujemy $f\simeq g$,
jeśli istnieje ciągłe $H:X\times I\to Y$ z
$H(x,0)=f(x)$ i $H(x,1)=g(x)$.
Jeśli $f|_A=g|_A$ oraz $H(a,t)=f(a)$
dla $a\in A\subseteq X$ i $t\in I$,
mówimy o homotopii \emph{względem $A$}.
Przestrzenie $X,Y$ są \emph{homotopijnie
równoważne}, gdy istnieją ciągłe $f:X\to Y$,
$g:Y\to X$ z $gf\simeq\id_X$ i
$fg\simeq\id_Y$.
\end{definicja}

Homeomorfizm ma prawdziwą odwrotność.
Równoważność homotopijna dopuszcza odwrotność
z dokładnością do deformacji. Dysk i punkt
nie są homeomorficzne, ale są homotopijnie
równoważne.

\begin{definicja}[Retrakt deformacyjny]
\label{def:retrakt-deformacyjny}
Podprzestrzeń $A\subseteq X$ jest
\emph{silnym retraktem deformacyjnym} $X$,
jeśli istnieje ciągłe $H:X\times I\to X$ takie, że
\[
 H(x,0)=x,\qquad H(x,1)\in A,\qquad
 H(a,t)=a\quad(a\in A).
\]
Odwzorowanie $r(x):=H(x,1)$ jest wówczas
retrakcją $X\to A$. Przestrzeń homotopijnie
równoważną punktowi nazywamy \emph{ściągalną}.
\end{definicja}

\begin{przyklad}[Promieniowa deformacja]
\label{prz:retrakcja-pierscien}
Dla $x\in\R^n\setminus\{0\}$ wzór
\[
 H(x,t)=\left((1-t)+\frac{t}{\|x\|}\right)x
\]
daje silny retrakt deformacyjny
$\R^n\setminus\{0\}$ na $S^{n-1}$.
Współczynnik jest dodatni, więc tor
nie przechodzi przez usunięty początek.
Dla $\|x\|=1$ punkt pozostaje nieruchomy,
a dla $t=1$ otrzymujemy $x/\|x\|$.
Każdy niepusty zbiór wypukły $C\subseteq\R^N$ jest ściągalny:
dla wybranego $c_0\in C$ stosujemy
$H(x,t)=(1-t)x+tc_0$.
\end{przyklad}

Do badania dziur potrzebujemy dróg,
których końce pozostają nieruchome w czasie deformacji.

\begin{definicja}[Drogi, pętle i homotopia dróg]
\label{def:droga-petla}
\emph{Droga} w $X$ to ciągłe
$\alpha:I\to X$; jej początek i koniec
to $\alpha(0)$ i $\alpha(1)$.
Gdy oba są równe $x_0$, nazywamy ją
\emph{pętlą opartą w $x_0$}.
Drogi $\alpha,\beta$ o tych samych końcach
są homotopijne \emph{względem końców},
jeśli łączy je homotopia $H:I\times I\to X$
z $H(0,t)=\alpha(0)$ i
$H(1,t)=\alpha(1)$ dla każdego $t$.
Gdy $\alpha(1)=\beta(0)$, określamy
\[
 (\alpha*\beta)(s)=
 \begin{cases}
  \alpha(2s),&0\leq s\leq\tfrac12,\\
  \beta(2s-1),&\tfrac12\leq s\leq1,
 \end{cases}
 \qquad \bar\alpha(s)=\alpha(1-s).
\]
Symbol $\alpha*\beta$ oznacza: najpierw $\alpha$,
potem $\beta$.
\end{definicja}

Homotopia względem końców jest relacją równoważności.
Stała deformacja daje zwrotność, zamiana $t$ na $1-t$ daje
symetrię, a sklejenie dwóch homotopii na połówkach przedziału
czasu daje przechodniość; ciągłość wynika z lematu o sklejaniu.

\begin{figure}[!htbp]
\centering
\begin{tikzpicture}[>=Stealth,line cap=round,line join=round,scale=.9]
 \fill[manifoldblue!5,rounded corners=18pt]
  (-.4,-1.15) rectangle (10.2,2.45);
 \draw[manifoldblue!35,thick,rounded corners=18pt]
  (-.4,-1.15) rectangle (10.2,2.45);
 \coordinate (x) at (.5,.25);
 \coordinate (y) at (9.25,.25);
 \draw[manifoldred,very thick,-{Stealth[length=2.8mm]}]
  (x) .. controls (2.6,2.4) and (6.7,2.0) .. (y)
  node[pos=.48,above=4pt] {$\alpha$};
 \draw[manifoldgreen!80!black,very thick,-{Stealth[length=2.8mm]}]
  (x) .. controls (2.7,-1.0) and (6.5,-.7) .. (y)
  node[pos=.51,below=4pt] {$\beta$};
 \draw[manifoldgold!90!black,thick,-{Stealth[length=2.5mm]}]
  (4.9,1.5) -- (4.9,.35)
  node[pos=.52,right=4pt,font=\small] {deformacja};
 \fill[manifoldblue] (x) circle(2.4pt) (y) circle(2.4pt);
 \node[below left=3pt] at (x) {$x_0$};
 \node[below right=3pt] at (y) {$x_1$};
\end{tikzpicture}
\caption{Homotopia dróg przesuwa ich wnętrza,
ale zachowuje oba punkty końcowe.}
\label{fig:homotopia-drog}
\end{figure}

\begin{lemat}[Zmiana prędkości nie zmienia klasy]
\label{lem:reparametr-drogi}
Jeśli $r:I\to I$ jest ciągła, niemalejąca
i spełnia $r(0)=0$, $r(1)=1$, to
$\alpha\circ r$ jest homotopijna z $\alpha$
względem końców.
\end{lemat}
\begin{proof}
Połóżmy $r_t(s)=(1-t)r(s)+ts$.
Każde $r_t$ jest ciągłe, ma te same końce
i wartości w $I$. Funkcja
$H(s,t)=\alpha(r_t(s))$ jest ciągła.
Dla $t=0$ daje $\alpha\circ r$, dla $t=1$
daje $\alpha$, a końce pozostają stałe.
\end{proof}

\begin{twierdzenie}[Grupa podstawowa]
\label{thm:grupa-podstawowa}
Klasy homotopii pętli opartych w $x_0$,
z działaniem $[\alpha][\beta]=[\alpha*\beta]$,
tworzą grupę $\pi_1(X,x_0)$.
Element neutralny to stała pętla $c_{x_0}$,
a odwrotnością $[\alpha]$ jest $[\bar\alpha]$.
\end{twierdzenie}
\begin{proof}
Jeśli $H$ deformuje $\alpha$ w $\alpha'$,
a $K$ deformuje $\beta$ w $\beta'$
względem końców, sklejenie
\[
 L(s,t)=
 \begin{cases}
   H(2s,t),&s\leq\tfrac12,\\
   K(2s-1,t),&s\geq\tfrac12
 \end{cases}
\]
jest ciągłe, bo na wspólnej krawędzi
obie części mają wartość $x_0$.
Działanie nie zależy więc od reprezentantów.

Przy trzech pętlach $(\alpha*\beta)*\gamma$
przeznacza na $\alpha,\beta,\gamma$
odcinki czasu długości $1/4,1/4,1/2$,
zaś $\alpha*(\beta*\gamma)$ długości
$1/2,1/4,1/4$. Obie drogi przebiegają
te same kawałki w tej samej kolejności
i różnią się zmianą parametru.
Lemat~\ref{lem:reparametr-drogi} daje łączność
na klasach. Pętle $c_{x_0}*\alpha$
i $\alpha*c_{x_0}$ różnią się od $\alpha$
tylko chwilą postoju, więc
$[c_{x_0}]$ jest neutralny.

Dla odwrotności określmy
\[
 H(s,t)=
 \begin{cases}
  \alpha(2(1-t)s),&0\leq s\leq\tfrac12,\\
  \alpha(2(1-t)(1-s)),&\tfrac12\leq s\leq1.
 \end{cases}
\]
Przy $s=1/2$ obie gałęzie są równe
$\alpha(1-t)$, a przy $s=0,1$ mają
wartość $x_0$. Jest to homotopia
$\alpha*\bar\alpha$ do pętli stałej
względem punktu bazowego. Zamiana
$\alpha$ z $\bar\alpha$ daje też
$[\bar\alpha][\alpha]=[c_{x_0}]$.
\end{proof}

\begin{stwierdzenie}[Odwzorowania indukowane i punkt bazowy]
\label{prop:pi1-funktor}
Odwzorowanie ciągłe $f:(X,x_0)\to(Y,y_0)$
z $f(x_0)=y_0$ wyznacza homomorfizm
\[
 f_*:\pi_1(X,x_0)\to\pi_1(Y,y_0),
 \qquad [\alpha]\mapsto[f\circ\alpha].
\]
Zachodzi $(gf)_*=g_*f_*$ i
$(\id_X)_*=\id$. Homotopia odwzorowań
utrzymująca $x_0$ w $y_0$ daje
równe homomorfizmy. Jeśli
$u$ jest drogą od $x_0$ do $x_1$,
to $b_u([\alpha])=[\bar u*\alpha*u]$
jest izomorfizmem
$\pi_1(X,x_0)\to\pi_1(X,x_1)$.
\end{stwierdzenie}
\begin{proof}
Złożenie $f\circ H$ przenosi homotopię
pętli w homotopię ich obrazów, a
$f\circ(\alpha*\beta)
=(f\circ\alpha)*(f\circ\beta)$.
Równości dla złożeń wynikają stąd wprost.
Jeśli $F$ jest homotopią odwzorowań utrzymującą
punkt bazowy, $F(\alpha(s),t)$ jest
homotopią pętli opartych w $y_0$.
Wreszcie $b_u$ zachowuje iloczyn:
w środku złożenia pojawia się
$u*\bar u$, homotopijna z pętlą stałą.
Odwrotność stanowi $b_{\bar u}$.
\end{proof}

\begin{wniosek}[Niezmienniczość grupy podstawowej]
\label{cor:pi1-homotopia}
Jeśli $f\simeq g:X\to Y$ przez homotopię
$H$, a $u(t)=H(x_0,t)$ prowadzi od
$f(x_0)$ do $g(x_0)$, to
$g_*=b_u\circ f_*$. Każda równoważność
homotopijna $f:X\to Y$ daje więc
izomorfizm grup podstawowych,
po uwzględnieniu zmiany punktu bazowego.
\end{wniosek}
\begin{proof}
Dla pętli $\alpha$ rozpatrzmy kwadrat
$(s,t)\mapsto H(\alpha(s),t)$.
Jego dolny i górny bok to $f\alpha$
i $g\alpha$, a oba pionowe boki to
$u$. Brzeg kwadratu jest homotopijny do pętli stałej względem
wybranego narożnika: w wypukłym kwadracie ściągamy jego parametryzację
liniowo do tego narożnika, a następnie składamy z $H(\alpha(s),t)$.
Zatem w grupie z punktem
bazowym $g(x_0)$ mamy
$[g\alpha]=[\bar u*(f\alpha)*u]$.
To daje pierwszą równość.

Niech teraz $g$ będzie homotopijną
odwrotnością $f$ i oznaczmy
$y_0=f(x_0)$. Homotopia
$gf\simeq\id_X$ daje drogę $u$
od $gf(x_0)$ do $x_0$ i
izomorfizm $b_u\circ(gf)_*$.
Homotopia $fg\simeq\id_Y$ daje pewną drogę $v$
od $fg(y_0)$ do $y_0$ z $b_v\circ(fg)_*=\id$.
Droga $fu$ ma te same końce, choć może nie być równa $v$.
Ponieważ $b_{fu}\circ b_v^{-1}$ jest izomorfizmem,
również $b_{fu}\circ(fg)_*$ jest izomorfizmem grupy
$\pi_1(Y,y_0)$ na siebie
(każda zmiana punktu bazowego
jest izomorfizmem).
Z naturalności konkatenacji
$f_*\circ b_u=b_{f u}\circ f_*$
na odpowiednich grupach.
Zatem dla $g'_*:=b_u\circ g_*$
oba złożenia $g'_*f_*$
oraz $f_*g'_*$ są izomorfizmami.
Pierwsze zapewnia iniektywność
$f_*$, drugie jego surjektywność.
\end{proof}

Silna retrakcja deformacyjna
z~\ref{def:retrakt-deformacyjny} daje
izomorfizm grup podstawowych w punktach
podprzestrzeni. Dla przestrzeni wypukłej
każdą pętlę $\alpha$ można deformować
wzorem $(1-t)\alpha(s)+tx_0$ do pętli
stałej, nie ruszając punktu bazowego.

\section{Nakrycia: lokalnie wiele kopii jednej przestrzeni}
\label{sec:nakrycia}

Nakrycie przypomina nałożenie kilku arkuszy
na tę samą małą okolicę. Gdy idziemy po bazie,
wybranie początkowego arkusza jednoznacznie
ustala przebieg drogi na nakryciu.
Po zamkniętej drodze możemy wrócić na inny arkusz.

\begin{definicja}[Nakrycie i arkusze]
\label{def:nakrycie}
Surjekcja ciągła $p:E\to X$ jest \emph{nakryciem},
jeśli każdy $x\in X$ ma otwarte otoczenie $U$
takie, że
\[
 p^{-1}(U)=\coprod_{\lambda\in\Lambda}V_\lambda,
 \qquad
 p|_{V_\lambda}:V_\lambda\xrightarrow{\cong}U,
\]
gdzie $V_\lambda$ są parami rozłącznymi
otwartymi podzbiorami $E$, a ograniczenia $p|_{V_\lambda}$
są homeomorfizmami.
Takie $U$ nazywamy \emph{równomiernie nakrytym},
a $V_\lambda$ \emph{arkuszami}.
Włókno $p^{-1}(x)$ jest dyskretne.
\end{definicja}

\begin{przyklad}[Prosta nad okręgiem]
\label{prz:exp-nakrycie}
Odwzorowanie
\[
 p:\R\to S^1,\qquad p(t)=e^{2\pi i t},
\]
jest nakryciem. Jeśli $U$ jest krótkim,
otwartym łukiem okręgu, wybierzmy przedział
$J$ długości mniejszej niż $1$, na którym
$p|_J$ jest homeomorfizmem na $U$.
Wtedy $p^{-1}(U)=\coprod_{k\in\Z}(J+k)$.
Włókno nad $1$ to $\Z$.
Podobnie $q_m:S^1\to S^1$, $q_m(z)=z^m$,
jest $m$-krotnym nakryciem dla $m\geq1$.
Dwukrotne nakrycie $\operatorname{Spin}(n)
\to\mathrm{SO}(n)$ pojawiło się już w teorii grup Spin.
\end{przyklad}

\begin{figure}[!htbp]
\centering
\begin{tikzpicture}[>=Stealth,line cap=round,line join=round,
 x={(1.55cm,0cm)},y={(.40cm,.37cm)},z={(0cm,1.25cm)}]
 % Okrąg u dołu, widoczna i zasłonięta połowa.
 \draw[manifoldblue!40,densely dashed,line width=1pt]
  plot[domain=0:180,samples=60,variable=\u]
   ({cos(\u)},{sin(\u)},0);
 \draw[manifoldblue,line width=1.8pt]
  plot[domain=180:360,samples=60,variable=\u]
   ({cos(\u)},{sin(\u)},0);
 % Włókno nad 1 oraz promień projekcji dla punktu t=1.25.
 \draw[manifoldblue!30,densely dotted]
  (1,0,0)--(1,0,2.9);
 \draw[manifoldgold!90!black,thick,dashed,
       -{Stealth[length=2.2mm]}]
  (0,1,1.65)--(0,1,0);
 \node[manifoldgold!70!black,right=4pt]
  at (0,1,.75) {$p$};
 % Fragment helisy: jeden obrót podnosi wysokość o jedną jednostkę.
 \draw[manifoldred,line width=2.1pt,
       -{Stealth[length=2.8mm]}]
  plot[domain=0:2.38,samples=180,variable=\t]
   ({cos(360*\t)},{sin(360*\t)},{.4+\t});
 \foreach \k in {0,1,2} {
  \fill[manifoldred] (1,0,{.4+\k}) circle(2.4pt);
  \node[manifoldred,right=5pt,font=\small]
   at (1,0,{.4+\k}) {$t=\k$};
 }
 \fill[manifoldblue] (1,0,0) circle(2.8pt);
 \node[manifoldblue,below right=3pt] at (1,0,0) {$1\in S^1$};
 \node[manifoldred!85!black,left=7pt]
  at (-1,0,2.7) {$\widetilde X\cong\R$};
 \node[manifoldblue!85!black,left=11pt]
  at (-1,0,-.22) {$X=S^1$};
\end{tikzpicture}
\caption{Nakrycie $p:\R\to S^1$, $p(t)=e^{2\pi i t}$,
przedstawione jako linia śrubowa nad okręgiem. Rzut pionowy
odpowiada $p$. Punkty linii śrubowej oznaczone $t=0,1,2$
leżą nad tym samym punktem $1\in S^1$.
Pokazano fragment linii śrubowej, która biegnie dalej
w obie strony.}
\label{fig:nakrycie-okragu}
\end{figure}

W poniższych dowodach użyjemy \emph{liczby Lebesgue'a}:
dla otwartego pokrycia zwartej przestrzeni metrycznej istnieje
$\delta>0$ takie, że każdy niepusty zbiór o średnicy mniejszej
niż $\delta$ leży w jednym elemencie pokrycia.
Oto uzasadnienie. Dla każdego $x$ wybierzmy $\varepsilon_x>0$
tak, aby $B(x,2\varepsilon_x)$ mieściła się w elemencie pokrycia.
Ze zwartości skończenie wiele kul $B(x_i,\varepsilon_{x_i})$
pokrywa przestrzeń. Wystarczy $\delta=\min_i\varepsilon_{x_i}$:
jeśli $y$ należy do jednej z tych mniejszych kul, wszystkie
punkty zbioru zawierającego $y$ i o średnicy $<\delta$
leżą w odpowiadającej jej większej kuli. Stosujemy ten fakt
do odcinka i kwadratu z ich zwykłą metryką.

\begin{twierdzenie}[Podnoszenie drogi i homotopii]
\label{thm:podnoszenie-nakrycie}
Niech $p:E\to X$ będzie nakryciem.
\begin{enumerate}[label=(\alph*),leftmargin=*]
\item Każda droga $\gamma:I\to X$ i wybór
      $e_0\in p^{-1}(\gamma(0))$ wyznaczają
      dokładnie jedną drogę
      $\widetilde\gamma:I\to E$ z
      $\widetilde\gamma(0)=e_0$ i
      $p\widetilde\gamma=\gamma$.
\item Jeśli $H:I\times I\to X$ jest
      homotopią dróg, a $\widetilde H(s,0)$
      jest zadanym podniesieniem $H(s,0)$,
      to istnieje dokładnie jedno ciągłe
      podniesienie $\widetilde H$ całego $H$.
      Gdy $H$ zachowuje końce dróg,
      podniesione końce również są stałe.
\end{enumerate}
\end{twierdzenie}
\begin{proof}
\emph{Droga.} Przeciwobrazy równomiernie
nakrytych otwartych $U$ tworzą otwarte
pokrycie odcinka $I$. Z liczby Lebesgue'a
tego pokrycia wybieramy podział
$0=t_0<\cdots<t_N=1$ tak drobny, aby
każdy obraz $\gamma([t_{j-1},t_j])$
mieścił się w jednym $U_j$.
Arkusz nad $U_1$ zawierający $e_0$
określa na pierwszym przedziale wzór
\[
 \widetilde\gamma(t)
 =(p|_{V_1})^{-1}(\gamma(t)).
\]
W punkcie $t_1$ wybieramy arkusz nad $U_2$
zawierający uzyskany punkt włókna i powtarzamy.
Wzory zgadzają się na końcach przedziałów,
więc sklejają się w drogę ciągłą.
Jeśli dwa podniesienia zgadzają się w pewnej
chwili, to lokalnie leżą w tym samym arkuszu
i są identyczne. Gdy w pewnej chwili
się różnią, leżą nad tym samym punktem
w różnych arkuszach, więc również w małej
okolicy chwili pozostają różne. Zbiór
chwil zgodności jest zatem otwarty
i domknięty w spójnym $I$ oraz
zawiera $0$, co daje jednoznaczność.

\emph{Homotopia.} Pokrycie $I^2$ zbiorami
$H^{-1}(U)$ ma liczbę Lebesgue'a.
Dzielimy kwadrat na skończoną siatkę
małych prostokątów, z których każdy ma
obraz w jednym równomiernie nakrytym $U$.
Na dolnej krawędzi mamy zadaną drogę.
Podnosimy prostokąty kolejno od dołu i
od lewej: na każdym wybieramy arkusz
zawierający już ustalony obraz jego narożnika
i stosujemy $(p|_V)^{-1}\circ H$.
Na wspólnej krawędzi oba wzory są
podniesieniami tej samej drogi i zgadzają
się w narożniku, więc na mocy (a) są równe.
Skończone sklejanie daje ciągłą
$\widetilde H$; lokalna jednoznaczność
na spójnym $I^2$ daje jej jedyność.
Jeśli $H(0,t)$ albo $H(1,t)$ jest stała,
jej podniesienie leży w dyskretnym włóknie,
więc też jest stałe.
\end{proof}

Stąd $p_*:\pi_1(E,e_0)\to\pi_1(X,p(e_0))$
jest iniektywny. Jeśli obraz pętli jest
homotopijny do stałej pętli względem bazy,
podnosimy homotopię od danej pętli w $E$.
Końcowa podniesiona pętla leży w dyskretnym
włóknie i jest stała.

\begin{twierdzenie}[Kryterium podnoszenia odwzorowań]
\label{thm:kryterium-podnoszenia}
Niech $Y$ będzie łukowo spójna i lokalnie
łukowo spójna, $f:Y\to X$ ciągła,
$p:E\to X$ nakryciem, a $p(e_0)=f(y_0)$.
Istnieje ciągła $\widetilde f:Y\to E$
z $p\widetilde f=f$ i
$\widetilde f(y_0)=e_0$ wtedy i tylko wtedy, gdy
\[
 f_*\pi_1(Y,y_0)\subseteq
 p_*\pi_1(E,e_0).
\]
Takie podniesienie jest jedyne.
\end{twierdzenie}
\begin{proof}
Gdy $\widetilde f$ istnieje,
$f=p\widetilde f$, zatem
$f_*=p_*\widetilde f_*$, co daje konieczność.
Załóżmy inkluzję. Dla $y\in Y$ wybierzmy
drogę $\gamma$ od $y_0$ do $y$ i
zdefiniujmy $\widetilde f(y)$ jako koniec
podniesienia $f\circ\gamma$ z początkiem $e_0$.
Gdy $\gamma'$ jest innym wyborem,
pętla $f\circ(\gamma'*\bar\gamma)$
ma klasę w obrazie $p_*$. Jest więc
homotopijna względem punktu bazowego
do obrazu pewnej pętli w $E$ o początku
$e_0$. Podnosząc tę homotopię za pomocą
Twierdzenia~\ref{thm:podnoszenie-nakrycie},
otrzymujemy, że podniesienie
$f\circ(\gamma'*\bar\gamma)$ wraca do $e_0$.
Druga część tego zamkniętego podniesienia jest podniesieniem
$f\circ\bar\gamma$ kończącym się w $e_0$.
Po odwróceniu jest, z jednoznaczności, podniesieniem
$f\circ\gamma$ zaczynającym się w $e_0$.
Podniesienia $f\circ\gamma$ i $f\circ\gamma'$ mają więc ten sam koniec.

Sprawdźmy ciągłość, której nie daje sama
definicja punkt po punkcie. Wybierzmy
równomiernie nakryte $U$ zawierające $f(y)$
oraz łukowo spójne otwarte $W\ni y$ z
$f(W)\subseteq U$. Niech $V$ będzie
arkuszem nad $U$ zawierającym $\widetilde f(y)$.
Dla $z\in W$ połączmy $y$ z $z$ drogą
w $W$ i dopiszmy ją do drogi od $y_0$
do $y$. Definicja i podnoszenie drogi
pokazują wtedy, że
$\widetilde f(z)=(p|_V)^{-1}(f(z))$.
Na $W$ odwzorowanie jest więc ciągłe.
Jedyność wynika z jedyności podniesienia
każdej drogi od $y_0$.
\end{proof}

\begin{twierdzenie}[Grupa podstawowa okręgu]
\label{thm:pi1-okragu}
Podnoszenie względem $p:\R\to S^1$
z Przykładu~\ref{prz:exp-nakrycie} daje
izomorfizm
\[
 \operatorname{wind}:\pi_1(S^1,1)\xrightarrow{\ \cong\ }\Z,
 \qquad [\gamma]\longmapsto
 \widetilde\gamma(1),
 \quad \widetilde\gamma(0)=0.
\]
Liczba po prawej to \emph{liczba owinięć}.
\end{twierdzenie}
\begin{proof}
Ponieważ $\gamma(1)=1$, koniec
podniesienia leży w $p^{-1}(1)=\Z$.
Homotopijne pętle dają homotopię
podniesień z nieruchomymi końcami,
więc liczba zależy tylko od klasy.
Jeśli podniesienie $\alpha$ kończy się w $k$,
to podniesienie następnej pętli $\beta$
zaczynające się w $k$ jest przesunięciem
o $k$ podniesienia zaczynającego się w $0$.
Stąd $\operatorname{wind}([\alpha*\beta])
=\operatorname{wind}([\alpha])
+\operatorname{wind}([\beta])$.
Pętla $\gamma_k(s)=e^{2\pi i ks}$
ma podniesienie $ks$, więc homomorfizm
jest surjektywny. Gdy podniesienie
$\widetilde\gamma$ wraca do $0$,
wzór $(s,t)\mapsto(1-t)\widetilde\gamma(s)$
deformuje ją w $\R$ do stałej drogi
z zachowaniem końców. Projekcja tej
homotopii do $S^1$ deformuje $\gamma$
do pętli stałej. Jądro jest zerowe.
\end{proof}

Z Przykładu~\ref{prz:retrakcja-pierscien}
otrzymujemy $\pi_1(\R^2\setminus\{0\})\cong\Z$.
W szczególności okrąg i punkt nie są
homotopijnie równoważne. Obliczenie
$\pi_1(S^{n})=0$ dla $n\geq2$
wyniknie niżej z twierdzenia van Kampena.

\begin{definicja}[Lokalna spójność i warunek półlokalny]
\label{def:semilokalnie-prosto-spojna}
Przestrzeń $X$ jest \emph{lokalnie łukowo spójna},
gdy każdy punkt ma bazę otoczeń otwartych,
które są łukowo spójne. Jest
\emph{półlokalnie jednospójna}, gdy każdy
$x\in X$ ma otoczenie $U$, w którym
każda pętla oparta w $x$ jest homotopijna do stałej
względem punktu bazowego \emph{w całym $X$}. Nie wymagamy, aby
ta homotopia mieściła się w $U$.
Przestrzeń łukowo spójna z
$\pi_1(X,x_0)=0$ jest \emph{jednospójna}.
\end{definicja}

Każda rozmaitość topologiczna spełnia
lokalne warunki z tej definicji, gdyż
w otwartej mapie współrzędnych można
wybrać małą kulę. Warunek półlokalny
jest ważny: bez niego nakrycie uniwersalne
może w ogóle nie istnieć.

\begin{twierdzenie}[Konstrukcja nakrycia uniwersalnego]
\label{thm:uniwersalne-nakrycie}
Jeśli $X$ jest łukowo spójna, lokalnie
łukowo spójna i półlokalnie jednospójna,
to ma nakrycie $q:\widetilde X\to X$
z jednospójną przestrzenią $\widetilde X$.
\end{twierdzenie}
\begin{proof}
Ustalmy $x_0\in X$. Punktem $\widetilde X$
niech będzie klasa $[\gamma]$ drogi w $X$
zaczynającej się w $x_0$, przy czym
homotopie zachowują oba końce.
Definiujemy $q([\gamma])=\gamma(1)$.
Istnieje baza otwartych łukowo spójnych
zbiorów $U\subseteq X$, w których
każda pętla jest zerowa po włączeniu do $X$:
najpierw wybieramy otoczenie z warunku
półlokalnego, następnie w nim małe
łukowo spójne otoczenie z lokalnej
łukowej spójności. Każdą pętlę opartą w innym punkcie
takiego $U$ przenosimy do jego wyróżnionego punktu drogą w $U$;
izomorfizm zmiany punktu bazowego pokazuje, że ona też jest zerowa w $X$.

Dla drogi $\gamma$ kończącej się w $x\in U$
niech $W([\gamma],U)$ składa się z klas
$[\gamma*\eta]$, gdzie $\eta$ jest drogą
w $U$ od $x$ do dowolnego $y\in U$.
Jeśli dwie takie drogi $\eta,\eta'$
mają ten sam koniec, to
$\eta*\bar\eta'$ jest pętlą w $U$,
zerową w $X$. Zatem $[\gamma*\eta]
=[\gamma*\eta']$. Otrzymujemy bijekcję
$q:W([\gamma],U)\to U$. Bierzemy wszystkie
te zbiory jako bazę topologii. Sprawdźmy warunek przecięć:
jeśli $[\beta]\in W([\gamma],U)\cap W([\gamma'],U')$,
wybieramy otwarte łukowo spójne $O$ z naszej bazy,
zawierające $\beta(1)$ i zawarte w $U\cap U'$.
Wtedy $W([\beta],O)$ zawiera $[\beta]$ i mieści się w obu zbiorach:
do drogi kończącej się w $\beta(1)$ dopisujemy tylko drogę w $O$.
To dowodzi, że rzeczywiście otrzymujemy bazę.
Ograniczenie $q$ do każdego $W$ jest ciągłe, ponieważ przeciwobrazy
otwartych podzbiorów $U$ są sumami takich mniejszych zbiorów bazowych;
jest też otwarte, bo mniejsze $W([\beta],O)$ przechodzą na otwarte $O$.
Wcześniej wykazana bijekcja jest więc homeomorfizmem.
Jeśli dwa $W$ nad tym samym
$U$ mają wspólny punkt, ich drogi początkowe
różnią się drogą w $U$ z dokładnością
do homotopii w $X$, więc zbiory $W$
są równe. Dla mniejszych otoczeń
używamy ograniczenia tych samych odwzorowań.
Stąd $q^{-1}(U)$ jest rozłączną sumą
arkuszy $W([\gamma],U)$: $q$ jest nakryciem.

Jest ono łukowo spójne. Drogę od klasy
stałej $[c_{x_0}]$ do $[\gamma]$
otrzymujemy, biorąc po kolei klasy
odcinków $\gamma|_{[0,t]}$ po
przeparametryzowaniu na $I$; lokalnie
leżą one w jednym $W$, więc tworzą
ciągłą drogę. Jeśli pętla
$\lambda$ w $\widetilde X$ jest oparta
w $[c_{x_0}]$, jej rzut $\gamma=q\lambda$
ma podniesienie $\lambda$ i kończy
się w tej samej klasie drogi.
Konstrukcja podniesienia pokazuje,
że ta klasa jest $[\gamma]$, więc
$\gamma$ jest homotopijna do stałej
pętli względem bazowego końca.
Podnosimy tę homotopię, zaczynając
od $\lambda$, za pomocą
Twierdzenia~\ref{thm:podnoszenie-nakrycie}.
Na końcu dostajemy stałą pętlę
w dyskretnym włóknie. Zatem
$\pi_1(\widetilde X)=0$.
\end{proof}

\begin{uwaga}[Skąd podgrupy grupy podstawowej?]
\label{uwaga:nakrycia-podgrupy}
Dla spójnego nakrycia $p:(E,e_0)\to(X,x_0)$
iniekcja $p_*$ wyróżnia podgrupę
$H=p_*\pi_1(E,e_0)\leq\pi_1(X,x_0)$.
W warunkach Twierdzenia~\ref{thm:uniwersalne-nakrycie}
każda taka podgrupa daje nakrycie:
na klasach dróg w $\widetilde X$ grupa
$\pi_1(X,x_0)$ działa przez
$[\omega]\cdot[\gamma]=[\omega*\gamma]$,
a iloraz $H\backslash\widetilde X$
jest nakryciem $X$. Istotnie, działanie przenosi
$W([\gamma],U)$ na $W([\omega*\gamma],U)$,
więc jest działaniem przez homeomorfizmy. Jest wolne:
z $[\omega*\gamma]=[\gamma]$ po dopisaniu $\bar\gamma$
wynika $[\omega]=1$.
Co więcej, jeśli $hW\cap W\ne\varnothing$, to $hW=W$,
bo są to arkusze nad tym samym $U$. Ponieważ $q$ jest
iniektywne na $W$ i $q(hz)=q(z)$, element $h$ ustala punkt $z\in W$;
stąd $h=1$. Różne przesunięcia jednego arkusza są więc rozłączne.
Projekcja na orbity jest otwarta: przeciwobraz obrazu zbioru
otwartego $O$ jest sumą otwartych zbiorów $hO$.
Obrazy arkuszy w ilorazie są zatem otwarte i każdy odwzorowuje się
homeomorficznie na $U$. To uzasadnia własność nakrycia.
Przestrzeń ilorazowa jest łukowo spójna jako obraz $\widetilde X$.
Podniesienie pętli $\gamma$ z $x_0$ kończy się w orbicie $[\gamma]$;
wraca do orbity drogi stałej dokładnie wtedy, gdy $[\gamma]\in H$.
Ponieważ obraz homomorfizmu indukowanego przez nakrycie składa się
właśnie z klas mających zamknięte podniesienie, skonstruowane nakrycie
realizuje dokładnie podgrupę $H$.
Z drugiej strony kryterium
\ref{thm:kryterium-podnoszenia} podnosi
$q:\widetilde X\to X$ do $E$, z ustaloną wartością $e_0$.
Przy obecnych założeniach $E$ jest lokalnie łukowo spójna,
więc jej spójność daje łukową spójność: składowe łukowe są otwarte.
Każdy punkt $E$ można połączyć z $e_0$ drogą; rzut tej drogi
pokazuje, że podniesione odwzorowanie $\widetilde X\to E$ jest surjektywne.
Dwa punkty $[\gamma],[\gamma']$ nad tym samym punktem $X$
trafiają w ten sam punkt $E$
dokładnie wtedy, gdy $[\gamma'*\bar\gamma]\in H$.
W jedną stronę składamy podniesione drogi o wspólnym końcu,
otrzymując pętlę w $E$; w drugą używamy podnoszenia homotopii
jak w dowodzie kryterium podnoszenia. Ten warunek jest też
równoważny przynależności do jednej orbity, gdyż
$[\gamma'*\bar\gamma]\cdot[\gamma]=[\gamma']$.
Lokalnie obie projekcje są
homeomorfizmami na $U$, co daje
$E\cong H\backslash\widetilde X$
nad $X$. Zmiana wybranego punktu
$e_0$ sprzęga podgrupę $H$: droga w $E$ między dwoma punktami
nad $x_0$ rzutuje się na pętlę $u$, a zmiana bazy zastępuje
$H$ podgrupą $[\bar u]H[u]$.
Odwrotnie, dowolną pętlę $u$ w $X$ można podnieść z $e_0$;
jej koniec wybiera punkt włókna realizujący właśnie to sprzężenie.
Tak powstaje klasyfikacja spójnych
nakryć przez klasy sprzężoności
podgrup $\pi_1(X,x_0)$.
\end{uwaga}

\section{Twierdzenie van Kampena i pierwsze obliczenia}
\label{sec:van-kampen}

Nakrycie okręgu mierzy jedną liczbę owinięć.
Gdy przestrzeń składa się z kilku części,
potrzebujemy sposobu łączenia informacji
o pętlach w tych częściach.

\begin{twierdzenie}[van Kampen: dwa otwarte zbiory]
\label{thm:van-kampen}
Niech $X=U\cup V$, gdzie $U,V$ są otwarte,
łukowo spójne, a $U\cap V$ jest łukowo
spójne i zawiera $x_0$. Jeśli
$G_U=\pi_1(U,x_0)$, $G_V=\pi_1(V,x_0)$
oraz $G_0=\pi_1(U\cap V,x_0)$, to
\[
 \pi_1(X,x_0)\cong
 (G_U*G_V)/
 \left\langle\!\left\langle
 (i_U)_*(h)(i_V)_*(h)^{-1}:h\in G_0
 \right\rangle\!\right\rangle.
\]
W iloczynie wolnym $G_U*G_V$ każde słowo ma postać zredukowaną:
niepuste słowa mają niejednostkowe litery z obu grup
występujące naprzemiennie. Przy mnożeniu sklejamy słowa,
mnożymy sąsiednie litery należące do tego samego czynnika i usuwamy
jedynki; słowo puste jest jedynką. Każda para homomorfizmów
z $G_U,G_V$ do jednej grupy przedłuża się jednoznacznie przez
mnożenie obrazów liter. Nawias
$\langle\!\langle\cdot\rangle\!\rangle$
oznacza najmniejszą podgrupę normalną
zawierającą wskazane relacje.
\end{twierdzenie}
\begin{proof}
Włączenia $U,V\hookrightarrow X$
określają homomorfizm
$\Phi:G_U*G_V\to\pi_1(X,x_0)$.
Pętla w $U\cap V$ daje ten sam obraz
niezależnie od wyboru strony.
Zatem wszystkie relacje z tezy
leżą w $\ker\Phi$.

\emph{Surjektywność.} Dla pętli
$\alpha:I\to X$ przeciwobrazy $U,V$
pokrywają zwarty odcinek. Dzielimy
go na tyle małych przedziałów,
aby obraz każdego leżał w $U$ albo $V$.
Sąsiednie kawałki przypisane temu
samemu zbiorowi łączymy. Każdy
punkt przejścia między kawałkiem
w $U$ i kawałkiem w $V$ leży w
$U\cap V$. Łączymy go z $x_0$
pomocniczą drogą w $U\cap V$.
Jeśli $\lambda_j$ prowadzi od $x_0$
do $j$-tego punktu przejścia,
to z kawałka $\alpha_j$ budujemy
pętlę $\lambda_{j-1}*\alpha_j*
\bar\lambda_j$ całkowicie w $U$
albo w $V$. Iloczyn tych pętli
odtwarza $\alpha$: sąsiednie
$\bar\lambda_j*\lambda_j$
kasują się na mocy
Twierdzenia~\ref{thm:grupa-podstawowa}.

Zmiana drogi $\lambda_j$ na $\lambda'_j$
wyznacza pętlę $h=\lambda'_j*\bar\lambda_j$ w $U\cap V$
\emph{opartą w $x_0$}. Poprzedni czynnik słowa mnożymy z prawej
przez $h^{-1}$, a następny z lewej przez $h$.
Oba zapisy $h$ są równe w ilorazie z tezy, więc się skracają.
Również dopisanie podziału wewnątrz
jednego otwartego zbioru nie zmienia
otrzymanej klasy słowa.

\emph{Jądro.} Przypuśćmy, że słowo
z $G_U*G_V$ daje pętlę homotopijną
do stałej w $X$. Wypełnijmy ją
homotopią $F:I^2\to X$, stałą na
pozostałych trzech bokach.
Zwartość kwadratu pozwala wybrać
siatkę tak drobną, by obraz
każdego małego prostokąta leżał
w całości w $U$ lub w $V$.
Podział dolnego boku dobieramy także tak, aby zawierał końce
pętli reprezentujących litery wyjściowego słowa.
Dla każdego wierzchołka $z$ siatki wybierzmy drogę $\mu_z$
od $x_0$ do $F(z)$. Jeśli $F(z)\in U\cap V$, wybieramy ją
w przecięciu; jeśli $F(z)$ należy tylko do jednego zbioru,
wybieramy drogę w tym zbiorze. Jest to możliwe dzięki założeniom
o łukowej spójności. Przy $F(z)=x_0$ wybieramy drogę stałą.

Zorientowanej krawędzi $e:z\to w$ przypisujemy klasę
$[\mu_z*(F|_e)*\bar\mu_w]$ w grupie zbioru zawierającego
obraz sąsiedniego prostokąta. Obie drogi pomocnicze leżą w tym
zbiorze. Jeżeli dwa sąsiednie prostokąty mają różne oznaczenia,
ich wspólna krawędź i obie drogi pomocnicze leżą w $U\cap V$,
więc przypisania dają ten sam element ilorazu z tezy.
Odwrócenie krawędzi daje element odwrotny.

Iloczyn czterech elementów wokół jednego prostokąta jest jedynką:
drogi pomocnicze skracają się parami, a jego brzeg jest homotopijny
do stałej we właściwym $U$ lub $V$, przez mapę $F$ na tym prostokącie.
Możemy zatem zastępować dwa kolejne boki prostokąta pozostałymi
dwoma, z tymi samymi końcami, bez zmiany elementu ilorazu.
Przesuwając tak łamaną po kolejnych prostokątach, wiersz po wierszu,
zamieniamy dolny bok kwadratu na drogę wzdłuż pozostałych trzech boków.
Te trzy boki są stałe. Na dolnym boku drogi pomocnicze wewnątrz
każdej pierwotnej litery skracają się w jej grupie, a przy końcach
liter są stałe. Odtwarzamy więc dokładnie wyjściowe słowo,
które okazało się jedynką w ilorazie. Każdy krok używał tylko
relacji wewnątrz obu grup i relacji pochodzących z przecięcia.
Stąd $\ker\Phi$ zawiera się w
wskazanym domknięciu normalnym relacji, co kończy dowód.
\end{proof}

\begin{przyklad}[Dwa okręgi złączone w punkcie]
\label{prz:pi1-dwa-okregi}
Niech $X=S^1_a\vee S^1_b$.
Wybierzmy na każdym okręgu po jednym
punkcie różnym od punktu sklejenia
$x_0$. Usuwając punkt z drugiego
okręgu, otrzymujemy otwarty $U$
retrakcyjny na $S^1_a$; analogicznie
$V$ retrakcyjny na $S^1_b$.
Przecięcie $U\cap V$ jest ściągalne
(to dwa otwarte łuki spotykające się
w $x_0$). Wobec tego
\[
 \pi_1(S^1\vee S^1,x_0)
 \cong\Z*\Z=:F(a,b).
\]
Słowa $ab$ i $ba$ są różne:
grupa podstawowa nie musi być przemienna.
\end{przyklad}

\begin{przyklad}[Sfera wymiaru co najmniej dwa]
\label{prz:pi1-sfera}
Niech $n\geq2$. Zbiory
$U=S^n\setminus\{N\}$ i
$V=S^n\setminus\{S\}$, gdzie $N,S$
są biegunami, są homeomorficzne
z $\R^n$, więc mają trywialne
grupy podstawowe. Ich przecięcie
jest łukowo spójne: deformuje się
na równik $S^{n-1}$, który jest
łukowo spójny dla $n-1\geq1$.
Twierdzenie~\ref{thm:van-kampen}
daje $\pi_1(S^n)=0$.
Stąd $\pi_1(\R^{n+1}\setminus\{0\})=0$
dla $n\geq2$, co odróżnia ten
przypadek od płaszczyzny.
\end{przyklad}

\section{Sympleksy i triangulacje}
\label{sec:sympleksy}

Sympleks jest najprostszym wielowymiarowym
klockiem: odcinek w wymiarze $1$, trójkąt
w wymiarze $2$, czworościan w wymiarze $3$.
Współrzędne barycentryczne opisują
zarazem punkt i to, na której ścianie leży.

\begin{definicja}[Sympleks standardowy i ściany]
\label{def:sympleks-standardowy}
Dla $n\geq0$ \emph{standardowy $n$-sympleks}
to
\[
 \Delta^n=
 \left\{(t_0,\ldots,t_n)\in\R^{n+1}:
 t_i\geq0,\ \sum_{i=0}^{n}t_i=1\right\}.
\]
Jego wierzchołki to wektory $e_0,\ldots,e_n$.
\emph{Wnętrze} oznacza $t_i>0$ dla wszystkich $i$
(względem afinicznej otoczki sympleksu).
Dla $n\geq1$ \emph{ściana kowymiaru jeden} naprzeciw $e_i$
jest zadana równaniem $t_i=0$ i jest kopią $\Delta^{n-1}$.
Ogólniej ścianę wyznacza dowolny niepusty podzbiór wierzchołków;
jest właściwa, gdy pomija choć jeden z nich.
Suma wszystkich właściwych ścian to $\partial\Delta^n$.
Przy $n=0$ nie ma niepustych właściwych ścian i $\partial\Delta^0=\varnothing$.
Punkty $v_0,\ldots,v_n\in\R^N$
są \emph{afinicznie niezależne}, jeśli
$v_1-v_0,\ldots,v_n-v_0$ są liniowo
niezależne. Ich otoczka wypukła
$[v_0,\ldots,v_n]$ jest
\emph{geometrycznym $n$-sympleksem}.
\end{definicja}

Każdy punkt sympleksu ma postać
$x=\sum_i t_iv_i$. Współczynniki
$t_i$ są jednoznaczne: jeśli
$\sum_i(t_i-s_i)v_i=0$ i obie sumy
współczynników są równe $1$, to
$\sum_{i=1}^n(t_i-s_i)(v_i-v_0)=0$;
afiniczna niezależność wymusza
$t_i=s_i$ dla wszystkich $i$.
W szczególności usunięcie jednego
wierzchołka daje ścianę
o wymiarze $n-1$.

\begin{figure}[!htbp]
\centering
\begin{tikzpicture}[line cap=round,line join=round]
 \node[font=\bfseries] at (0,3.45) {$\Delta^0$};
 \node[font=\bfseries] at (4.05,3.45) {$\Delta^1$};
 \node[font=\bfseries] at (8.63,3.45) {$\Delta^2$};
 % Wymiar zero.
 \fill[manifoldblue!10] (0,1.15) circle(.26);
 \fill[manifoldred] (0,1.15) circle(3.2pt);
 \node[below=9pt] at (0,1.15) {$v_0$};
 % Wymiar jeden.
 \coordinate (b0) at (2.75,1.15);
 \coordinate (b1) at (5.35,1.15);
 \draw[manifoldblue,line width=2.6pt] (b0)--(b1);
 \fill[manifoldred] (b0) circle(3.2pt)
  (b1) circle(3.2pt);
 \node[below=8pt] at (b0) {$v_0$};
 \node[below=8pt] at (b1) {$v_1$};
 % Wymiar dwa.
 \coordinate (c0) at (7.00,.35);
 \coordinate (c1) at (10.28,.35);
 \coordinate (c2) at (8.21,2.45);
 \fill[manifoldblue!13] (c0)--(c1)--(c2)--cycle;
 \draw[manifoldblue,line width=2pt] (c0)--(c1)--(c2)--cycle;
 \fill[manifoldred] (c0) circle(3.2pt)
  (c1) circle(3.2pt) (c2) circle(3.2pt);
 \node[below left=3pt] at (c0) {$v_0$};
 \node[below right=3pt] at (c1) {$v_1$};
 \node[above=4pt] at (c2) {$v_2$};
\end{tikzpicture}
\caption{Sympleksy wymiarów $0,1,2$: punkt,
odcinek i wypełniony trójkąt. Wymiar
$n$ odpowiada $n+1$ afinicznie niezależnym
wierzchołkom; ściany powstają przez ich usunięcie.}
\label{fig:sympleks-trojkat}
\end{figure}

\begin{stwierdzenie}[Brzeg sympleksu jest sferą]
\label{prop:brzeg-sympleksu}
Dla $n\geq1$ przestrzeń $\Delta^n$
jest homeomorficzna z dyskiem $D^n$,
a $\partial\Delta^n$ z $S^{n-1}$.
\end{stwierdzenie}
\begin{proof}
Pracujemy w $n$-wymiarowej płaszczyźnie
afinicznej $P=\{\sum t_i=1\}$.
Środek $c=(1/(n+1),\ldots,1/(n+1))$
leży we wnętrzu. Każdy promień z $c$
w $P$ przecina brzeg dokładnie raz:
nierówności $t_i\geq0$ wyznaczają
na nim odcinek od $c$ do pierwszej
ściany, na której pewna współrzędna
staje się zerem.
Przypisanie punktowi brzegowemu
$x$ kierunku $(x-c)/\|x-c\|$
jest ciągłą bijekcją
$\partial\Delta^n\to S^{n-1}$
po utożsamieniu $P-c\cong\R^n$.
Dziedzina jest zwarta, a sfera
Hausdorffa, więc odwrotność jest
ciągła. Jeśli $b(u)$ jest tą odwrotnością dla kierunku
$u\in S^{n-1}$, wzór $ru\mapsto c+r(b(u)-c)$, $0<r\leq1$,
z wartością $c$ w zerze daje ciągłą bijekcję z $D^n$ na $\Delta^n$.
Ciągłość w zerze wynika z ograniczoności $b(u)-c$.
Ponownie zwartość i własność Hausdorffa dają homeomorfizm,
który na brzegu jest poprzednią identyfikacją. Dla $n=1$ brzeg
składa się z dwóch punktów,
czyli jest $S^0$.
\end{proof}

\begin{definicja}[Skończony kompleks symplicjalny]
\label{def:kompleks-symplicjalny}
Niech $V$ będzie skończonym zbiorem
wierzchołków. \emph{Kompleks symplicjalny}
$K$ to rodzina niepustych skończonych
podzbiorów $V$ taka, że jeśli
$\sigma\in K$ i $\varnothing\ne\tau\subseteq\sigma$,
to $\tau\in K$. Zbiór o $n+1$
wierzchołkach nazywamy
\emph{$n$-sympleksem} $K$.
Jego \emph{realizacja geometryczna} to
\[
 |K|=\left\{\sum_{v\in V}t_v e_v:
 t_v\geq0,\ \sum_v t_v=1,\
 \{v:t_v>0\}\in K\right\}\subseteq\R^{V}.
\]
Kompleks $L\subseteq K$ zamknięty
na przechodzenie do ścian jest
\emph{podkompleksem}.
\emph{Skończona triangulacja} przestrzeni $X$
to homeomorfizm $|K|\to X$. W tej sekcji ograniczamy się do
skończonych kompleksów; ich realizacje są zwarte. Triangulacje
przestrzeni niezwartych wymagają dopuszczenia nieskończenie wielu sympleksów.
\end{definicja}

W tej definicji dwie różne ściany
mają dokładnie to wspólne, co
wynika ze wspólnych wierzchołków.
Na przykład trzy krawędzie
$[v_0,v_1]$, $[v_1,v_2]$,
$[v_2,v_0]$, bez wnętrza trójkąta,
realizują okrąg. Po dodaniu
$[v_0,v_1,v_2]$ otrzymujemy dysk.
Słowo \emph{podkompleks} jest
ważniejsze niż dowolny podzbiór:
zawiera wraz z sympleksem
wszystkie jego ściany.

\begin{uwaga}[Nie każdy podział na trójkąty
jest kompleksem symplicjalnym]
\label{uwaga:delta-motywacja}
Kwadrat z utożsamionymi przeciwległymi
bokami przedstawia torus. Gdy podzielimy
go jedną przekątną, po sklejeniu wszystkie
cztery narożniki stają się jednym
wierzchołkiem, a każda z dwóch
trójkątnych części nadal ma sens.
Nie jest to \emph{kompleks symplicjalny}
z jednym wierzchołkiem: trójkąt
potrzebuje trzech różnych wierzchołków.
Wygodniejszym opisem są
$\Delta$-kompleksy.
\end{uwaga}

\begin{definicja}[$\Delta$-kompleks]
\label{def:delta-kompleks}
\emph{$\Delta$-kompleks} $X$ jest
przestrzenią Hausdorffa z rodziną ciągłych odwzorowań charakterystycznych
$\sigma_\alpha:\Delta^{n_\alpha}\to X$:
\begin{enumerate}[label=(\roman*),leftmargin=*]
\item ograniczenie $\sigma_\alpha$ do wnętrza
      jest iniektywne, a obrazy wnętrz
      różnych sympleksów są rozłączne
      i pokrywają $X$;
\item ograniczenie $\sigma_\alpha$ do każdej
      ściany jest jednym z odwzorowań
      charakterystycznych niższego wymiaru
      po standardowym afinicznym
      utożsamieniu tej ściany z sympleksem, zachowującym
      kolejność pozostałych wierzchołków;
\item $A\subseteq X$ jest domknięty dokładnie
      wtedy, gdy $\sigma_\alpha^{-1}(A)$
      jest domknięty w każdym
      $\Delta^{n_\alpha}$.
\end{enumerate}
Różne ściany tego samego sympleksu
mogą być tu sklejone, o ile zgodne
są odwzorowania na ich dalszych ścianach.
\end{definicja}

Ostatni warunek to \emph{topologia słaba}
względem odwzorowań charakterystycznych:
odwzorowanie z $X$ do innej przestrzeni
jest ciągłe dokładnie wtedy, gdy jego
złożenia z każdą $\sigma_\alpha$ są ciągłe.
Jest to topologia ilorazowa z
Definicji~\ref{def:topologia-ilorazowa}
na rozłącznej sumie wszystkich sympleksów.

\begin{przyklad}[Okrąg i torus jako $\Delta$-kompleksy]
\label{prz:delta-torus}
Okrąg ma jeden $0$-sympleks $v$
i jeden $1$-sympleks, którego oba końce
przechodzą w $v$. Dla torusa
podzielmy kwadrat $[0,1]^2$
przekątną i sklejmy
$(0,t)\sim(1,t)$ oraz
$(s,0)\sim(s,1)$. Otrzymujemy
jeden wierzchołek, trzy krawędzie
(poziomą $a$, pionową $b$ i przekątną
$c$) oraz dwa $2$-sympleksy. Porządkujemy ich wierzchołki jako
$((0,0),(0,1),(1,1))$ i $((0,0),(1,0),(1,1))$.
Oba porządki nadają krawędziom poziomym kierunek w prawo,
pionowym w górę, a przekątnej kierunek od $(0,0)$ do $(1,1)$.
Mapy ścian są zatem zgodne także z porządkiem wierzchołków.
Iloraz jest homeomorficzny z
$(\R/\Z)\times(\R/\Z)$:
odwzorowanie $(s,t)\mapsto([s],[t])$
identyfikuje dokładnie wskazane
punkty brzegowe, jest ciągłe
i przechodzi na bijekcję
między przestrzenią zwartą
a przestrzenią Hausdorffa.
\end{przyklad}

\begin{figure}[!htbp]
\centering
\begin{tikzpicture}[>=Stealth,line cap=round,line join=round,
 x=1cm,y=1cm]
 \coordinate (a) at (0,0);
 \coordinate (b) at (4.4,0);
 \coordinate (c) at (4.4,4.4);
 \coordinate (d) at (0,4.4);
 \fill[manifoldblue!8] (a)--(b)--(c)--cycle;
 \fill[manifoldgreen!10] (a)--(d)--(c)--cycle;
 \draw[manifoldblue,very thick] (a)--(b)--(c)--(d)--cycle;
 \draw[manifoldgold!90!black,thick] (a)--(c);
 \draw[manifoldgold!90!black,thick,-{Stealth[length=2.7mm]}]
  (1.8,1.8)--(2.7,2.7);
 \draw[manifoldred,very thick,-{Stealth[length=2.7mm]}]
  (1.3,0)--(2.6,0);
 \draw[manifoldred,very thick,-{Stealth[length=2.7mm]}]
  (1.3,4.4)--(2.6,4.4);
 \draw[manifoldgreen!75!black,very thick,-{Stealth[length=2.7mm]}]
  (0,1.3)--(0,2.6);
 \draw[manifoldgreen!75!black,very thick,-{Stealth[length=2.7mm]}]
  (4.4,1.3)--(4.4,2.6);
 \node[manifoldred,below=2pt] at (3.2,0) {$a$};
 \node[manifoldgreen!65!black,right=2pt] at (4.4,3.2) {$b$};
 \node[manifoldgold!65!black,above left=2pt]
  at (2.6,2.6) {$c$};
 \foreach \p in {a,b,c,d}
  \fill[manifoldblue] (\p) circle(2.3pt);
 \node[font=\small] at (1.42,2.48) {$\Delta^2_1$};
 \node[font=\small] at (3.17,1.74) {$\Delta^2_2$};
\end{tikzpicture}
\caption{Torus z dwóch sympleksów: boki o tych samych
kolorach i kierunkach sklejamy ze sobą. Przekątna
pozostaje trzecią krawędzią. Wszystkie cztery narożniki
reprezentują ten sam wierzchołek torusa.}
\label{fig:delta-torus}
\end{figure}

\section{CW-kompleksy i dołączanie komórek}
\label{sec:cw-kompleksy}

Trójkąty są wygodne do obliczeń, lecz
potrafią wymagać wielu wierzchołków.
Komórka pozwala dołączyć naraz cały
otwarty dysk, przy czym jego brzeg
może biec po starszej części przestrzeni
w sposób nietrywialny.

\begin{definicja}[Dołączenie komórki]
\label{def:dolaczenie-komorki}
Niech $Y$ będzie przestrzenią i
$\varphi:S^{n-1}=\partial D^n\to Y$
odwzorowaniem ciągłym, $n\geq1$.
\emph{Dołączeniem $n$-komórki} jest
przestrzeń ilorazowa
\[
 Y\cup_\varphi D^n
 :=(Y\sqcup D^n)/(z\sim\varphi(z)
 \text{ dla }z\in S^{n-1}).
\]
Obraz wnętrza dysku oznaczamy $e^n$.
Odwzorowanie z dysku do ilorazu jest
\emph{odwzorowaniem charakterystycznym}
tej komórki; $\varphi$ nazywamy
\emph{odwzorowaniem dołączającym}.
W wymiarze $0$ po prostu dołączamy
izolowany punkt, gdyż
$S^{-1}=\varnothing$.
\end{definicja}

Naturalne włączenie $Y$ w przestrzeń dołączoną jest osadzeniem
na zbiór domknięty. Dla każdego domkniętego $F\subseteq Y$
przeciwobraz jego obrazu w $Y\sqcup D^n$ jest równy
$F\sqcup\varphi^{-1}(F)$ i jest domknięty; dotyczy to również $F=Y$.
Ponieważ na punktach $Y$ nie wprowadzamy nowych utożsamień,
dostajemy zamknięte osadzenie. Wnętrze $D^n$ jest natomiast
otwartym zbiorem nasyconym, bez żadnych utożsamień, więc przechodzi
homeomorficznie na otwartą komórkę $e^n$.

Topologia ilorazowa z
Definicji~\ref{def:topologia-ilorazowa}
ma tu konkretne znaczenie. Jeśli
$F:Y\cup_\varphi D^n\to Z$
ma być ciągła, wystarczy podać
ciągłe $F_Y:Y\to Z$ i
$F_D:D^n\to Z$ zgodne na brzegu:
$F_D(z)=F_Y(\varphi(z))$.
Istotnie, takie odwzorowania dają ciągłe
odwzorowanie na rozłącznej sumie, stałe
na każdej klasie sklejenia;
uniwersalna własność ilorazu
przenosi je na iloraz.
To jest praktyczny sposób
konstruowania odwzorowań na przestrzeniach
sklejonych.

\begin{definicja}[CW-kompleks]
\label{def:cw-kompleks}
\emph{CW-kompleks} $X$ jest przestrzenią Hausdorffa budowaną
ze szkieletów $X^0\subseteq X^1
\subseteq X^2\subseteq\cdots$:
$X^0$ jest dyskretnym zbiorem punktów,
a $X^n$ powstaje z $X^{n-1}$
przez dołączenie rodziny $n$-dysków
odwzorowaniami $S^{n-1}\to X^{n-1}$.
Przestrzeń $X=\bigcup_n X^n$
ma \emph{topologię słabą}:
$A\subseteq X$ jest domknięty,
dokładnie wtedy, gdy przeciwobraz $A$ pod każdym
odwzorowaniem charakterystycznym $D^n\to X$
jest domknięty. Wymagamy ponadto
\emph{skończoności domknięć komórek}:
domknięcie każdej komórki
spotyka tylko skończenie wiele
komórek. Litery CW przypominają
te dwa warunki
(\emph{closure-finite}, \emph{weak topology}).
\emph{Podkompleks CW} to suma komórek
domknięta w $X$; jeśli zawiera komórkę,
zawiera także całe jej domknięcie.
\end{definicja}

\begin{figure}[!htbp]
\centering
\begin{tikzpicture}[>=Stealth,line cap=round,line join=round]
 \node[font=\bfseries] at (0,3.02) {$X^0$};
 \node[font=\bfseries] at (4.33,3.02) {$X^1$};
 \node[font=\bfseries] at (8.68,3.02) {$X^2$};
 % 0-komórka.
 \fill[manifoldred] (0,1.39) circle(3.7pt);
 \node[below=7pt] at (0,1.39) {$e^0=\{v\}$};
 % Dołączenie 1-komórki: oba końce trafiają do v.
 \draw[manifoldblue,line width=2.5pt]
  (4.33,1.45) circle(1.03);
 \fill[manifoldred] (4.33,.42) circle(3.7pt);
 \node[manifoldblue!85!black] at (4.33,1.45) {$e^1$};
 \draw[->,manifoldblue!75!black,thin]
  (4.62,1.74)--({4.33+1.03*cos(45)},{1.45+1.03*sin(45)});
 \node[below=7pt] at (4.33,.42) {$v$};
 \draw[->,manifoldgold!80!black,thick]
  (1.22,1.54)--(3.08,1.54)
  node[midway,above=5pt,font=\small] {dołącz $e^1$};
 % Dołączenie 2-komórki: brzeg dysku biegnie po okręgu.
 \fill[manifoldgreen!18] (8.68,1.45) circle(1.03);
 \draw[manifoldblue,line width=2.5pt]
  (8.68,1.45) circle(1.03);
 \fill[manifoldred] (8.68,.42) circle(3.7pt);
 \node[manifoldgreen!65!black] at (8.68,1.45) {$e^2$};
 \node[below=7pt] at (8.68,.42) {$v$};
 \draw[->,manifoldgold!80!black,thick]
  (5.58,1.54)--(7.35,1.54)
  node[midway,above=5pt,font=\small] {dołącz $e^2$};
\end{tikzpicture}
\caption{Przykłady kolejnych wymiarów CW: punkt
$X^0=\{v\}$, okrąg $X^1\cong S^1$ otrzymany
przez sklejenie obu końców $1$-komórki z $v$,
oraz dysk $X^2\cong D^2$ po dołączeniu
$2$-komórki wzdłuż całego okręgu.
Czerwony punkt, niebieski łuk i zielone wnętrze
oznaczają odpowiednio otwarte komórki
wymiarów $0,1,2$.}
\label{fig:cw-wymiary}
\end{figure}

W kompleksie skończonym warunki
topologiczne łatwo sprawdzić
bez osobnej teorii granic:
topologia słaba jest ilorazową
topologią sklejania skończenie
wielu dysków, a domknięcie
każdej komórki spotyka
tylko skończenie wiele innych.
Przy nieskończenie wielu komórkach
topologia słaba jest konieczna,
aby ciągłość dało się sprawdzać
na wszystkich domkniętych komórkach.

\begin{stwierdzenie}[Sympleksy dają komórki]
\label{prop:sympleksy-cw}
Skończony kompleks symplicjalny i
skończony $\Delta$-kompleks są
CW-kompleksami: ich otwarte komórki
to obrazy wnętrz sympleksów.
\end{stwierdzenie}
\begin{proof}
Indukujemy po wymiarze. Szkielet
zerowy składa się z wierzchołków.
Wnętrze każdego $n$-sympleksu
jest rozłączne z obrazami innych
wnętrz, a brzeg jest sumą ścian
zawartych w poprzednim szkielecie.
Stwierdzenie~\ref{prop:brzeg-sympleksu}
pozwala traktować sympleks
jako dysk $D^n$, którego
odwzorowanie na brzegu jest odwzorowaniem dołączającym.
Warunek o topologii jest w
Definicji~\ref{def:delta-kompleks}
podany wprost, a dla skończonej
realizacji symplicjalnej wynika
z topologii skończonej sumy zwartych sympleksów:
ciągła surjekcja z tej sumy na $|K|\subseteq\R^V$ jest domknięta,
gdyż obrazy zbiorów domkniętych są zwarte w przestrzeni Hausdorffa.
Jest więc odwzorowaniem ilorazowym, co daje właśnie wymaganą topologię.
Skończona liczba sympleksów
zapewnia skończoność domknięć komórek.
\end{proof}

\begin{przyklad}[Minimalne struktury komórkowe]
\label{prz:cw-sfera-torus}
\begin{enumerate}[label=(\arabic*),leftmargin=*]
\item Dla $n\geq1$ sfera $S^n$
      ma jedną $0$-komórkę i jedną
      $n$-komórkę:
      $D^n/S^{n-1}\cong S^n$.
      Jawny homeomorfizm opisujemy
      poniżej.
      Dla $S^0$ używamy dwóch
      $0$-komórek.
\item Torus ma jedną $0$-komórkę,
      dwie $1$-komórki $a,b$ oraz
      jedną $2$-komórkę. Obchodząc
      brzeg kwadratu z
      Rysunku~\ref{fig:delta-torus},
      otrzymujemy słowo
      $aba^{-1}b^{-1}$ w
      $S^1\vee S^1$.
      Dwa trójkąty realizacji
      $\Delta$-kompleksowej zostały
      zastąpione jednym dyskiem,
      którego odwzorowanie dołączające
      zapamiętuje sklejanie.
\end{enumerate}
\end{przyklad}

\begin{przyklad}[Sfera $S^2$ bez $1$-komórek]
\label{prz:cw-s2}
Weźmy $X^0=\{v\}$ i nie dołączajmy
żadnej $1$-komórki, więc $X^1=X^0$.
Niech odwzorowanie dołączające
$\varphi:\partial D^2=S^1\to\{v\}$
będzie stałe. Wtedy
\[
 X^2=X^1\cup_\varphi D^2
 \cong D^2/\partial D^2\cong S^2.
\]
Cały brzeg dysku staje się \emph{jednym}
punktem $v$; obraz jego wnętrza
jest jedną otwartą $2$-komórką
$e^2=S^2\setminus\{v\}\cong\R^2$.
Rysunek~\ref{fig:cw-s2} pokazuje to
sklejenie. W szczególności
$X^1$ nie jest okręgiem:
okrąg po lewej stronie rysunku
to brzeg \emph{dołączanego dysku},
a nie komórka przestrzeni wynikowej.
\end{przyklad}

\begin{figure}[!htbp]
\centering
\begin{tikzpicture}[>=Stealth,line cap=round,line join=round]
 % Dysk przed sklejeniem: czerwony brzeg zapada się do v.
 \fill[manifoldblue!10] (0,1.85) circle(1.33);
 \draw[manifoldred,line width=2.4pt] (0,1.85) circle(1.33);
 \node[font=\bfseries] at (0,3.58) {$D^2$};
 \node[manifoldblue!80!black] at (0,1.85) {$\operatorname{int}D^2$};
 \node[manifoldred!85!black,below=7pt]
  at (0,.52) {$\partial D^2=S^1$};
 % Sklejenie brzegu do jednego punktu.
 \draw[->,manifoldgold!85!black,line width=1.5pt]
  (1.73,1.85)--(4.8,1.85)
  node[midway,above=7pt] {$q$};
 \node[font=\small,manifoldgold!70!black]
  at (3.26,1.25) {sklejamy cały brzeg};
 % Sfera w stylistyce rysunku 1.4: jasny środek, tył i przód łuków.
 \shade[inner color=white,outer color=manifoldblue!43]
  (6.8,1.85) circle(1.43);
 \draw[manifoldblue!75!black,very thick]
  (6.8,1.85) circle(1.43);
 \draw[manifoldblue!50,densely dashed,thin]
  (8.23,1.85) arc[start angle=0,end angle=180,x radius=1.43,y radius=.34];
 \draw[manifoldblue!70!black,thin]
  (5.37,1.85) arc[start angle=180,end angle=360,x radius=1.43,y radius=.34];
 \draw[manifoldblue!38,densely dashed,thin]
  (6.8,3.28) .. controls (6.35,2.76) and (6.35,.94) .. (6.8,.42);
 \draw[manifoldblue!65!black,thin]
  (6.8,3.28) .. controls (7.25,2.76) and (7.25,.94) .. (6.8,.42);
 \fill[manifoldred] (6.8,.42) circle(3.8pt);
 \node[font=\bfseries] at (6.8,3.58) {$X^2\cong S^2$};
 \node[manifoldblue!85!black] at (6.8,2.44) {$e^2$};
 \node[manifoldred!85!black,below=7pt] at (6.8,.42) {$v$};
\end{tikzpicture}
\caption{Struktura CW sfery: $D^2$ dołączamy do
punktu $v$ odwzorowaniem stałym na całym brzegu;
$q:D^2\to D^2/\partial D^2$ jest mapą ilorazową.
W powstałej sferze czerwony punkt to
jedyne $e^0$, a cała pozostała powierzchnia
to jedna otwarta komórka $e^2$.
Nie występują komórki wymiaru $1$.}
\label{fig:cw-s2}
\end{figure}

\begin{uwaga}[Uwaga o ilorazie dysku]
W punkcie (1) nie potrzebujemy
identyfikować półkuli z ilorazem
za pomocą niejawnego rysunku.
Istnieje jawne ciągłe odwzorowanie
$D^n\to S^n\subseteq\R^{n+1}$,
\[
 x\longmapsto
 \bigl(\sin(\pi\|x\|)\,x/\|x\|,
       \cos(\pi\|x\|)\bigr)
 \quad(x\ne0),
\]
z wartością $(0,\ldots,0,1)$
w $x=0$. Cały brzeg $\|x\|=1$
przechodzi w biegun południowy,
a na wnętrzu odwzorowanie jest iniektywne
i pokrywa sferę bez tego bieguna.
Przechodzi więc na ciągłą bijekcję
$D^n/S^{n-1}\to S^n$, która jest
homeomorfizmem przez zwartość
dziedziny i własność Hausdorffa
sfery.
\end{uwaga}

\begin{twierdzenie}[Co robi dołączenie komórki z $\pi_1$?]
\label{thm:pi1-dolaczenie-komorki}
Niech $Y$ będzie łukowo spójna,
a $X=Y\cup_\varphi D^n$.
Jeśli $n=2$, po wyborze drogi od
punktu bazowego do $\varphi(S^1)$
grupa $\pi_1(X)$ jest ilorazem
$\pi_1(Y)$ przez domknięcie normalne
klasy pętli $\varphi$.
Jeśli $n\geq3$, włączenie
$Y\hookrightarrow X$ wyznacza
izomorfizm na $\pi_1$.
\end{twierdzenie}
\begin{proof}
W dysku wybierzmy promień $r=\|x\|$.
Otwarty w $X$ zbiór $U$ złożony
z $Y$ i pasa $r>1/3$ deformuje
się na $Y$: przesuwamy punkty
pasa promieniowo do $r=1$,
a punkty $Y$ zostawiamy. Na części dyskowej homotopia ma wzór
$ru\mapsto((1-t)r+t)u$, gdzie $u\in S^{n-1}$.
Wyjaśnijmy ciągłość przy sklejanym $Y$. Dla otwartego otoczenia
$O$ punktu $y\in Y$ zmniejszamy jego część dyskową do punktów $ru$,
dla których $\varphi(u)\in O\cap Y$ i cały odcinek
$\{\rho u:r\leq\rho\leq1\}$ leży w przeciwobrazie $O$.
Ta część jest otwarta: zawieranie zwartego odcinka w zbiorze
otwartym zachowuje się przy małej zmianie $r,u$.
Wraz z $O\cap Y$ tworzy otwarte otoczenie $y$ w ilorazie,
którego obraz przez homotopię pozostaje w $O$ dla wszystkich $t$.
Dowodzi to ciągłości przy $Y$; na wnętrzu dysku wynika ona z podanego wzoru.
Drugi otwarty zbiór $V$ jest
obrazem $r<2/3$; jest to
ściągalne wnętrze dysku.
Przecięcie $U\cap V$ to pierścień
$1/3<r<2/3$, homotopijnie
równoważny $S^{n-1}$.
Jest łukowo spójne dla $n\geq2$.
Otwartość obu zbiorów sprawdzamy przed sklejeniem:
przeciwobrazy $U$ to całe $Y$ i otwarty w $D^n$ pas $r>1/3$,
a przeciwobraz $V$ to $r<2/3$, bez punktów $Y$.

Twierdzenie~\ref{thm:van-kampen}
z $G_V=1$ mówi, że $\pi_1(X)$
jest ilorazem $\pi_1(U)\cong\pi_1(Y)$
przez domknięcie normalne obrazu
$\pi_1(U\cap V)$. Przy $n=2$
generator $\pi_1(S^1)\cong\Z$
po przejściu przez retrakcję
$U\to Y$ staje się pętlą
$\varphi$ (z uwzględnieniem
drogi do punktu bazowego).
Przy $n\geq3$ mamy
$\pi_1(S^{n-1})=0$ z
Przykładu~\ref{prz:pi1-sfera},
więc nie pojawia się żadna
relacja ani nowy generator.
\end{proof}

\begin{przyklad}[Torus i powierzchnia rzutowa]
\label{prz:pi1-torus-rp2}
Dla torusa część jednowymiarowa
jest $S^1\vee S^1$, z grupą
$F(a,b)$ z Przykładu~\ref{prz:pi1-dwa-okregi}.
Dołączona $2$-komórka wymusza
$aba^{-1}b^{-1}=1$, czyli
\[
 \pi_1(T^2)\cong
 \langle a,b\mid ab=ba\rangle
 \cong\Z^2.
\]
Drugim przykładem jest
$\R P^2$: dysk będący górną
półkulą $S^2$ ma na brzegu
antypody utożsamione, więc
$\R P^2$ ma $0$-komórkę,
      $1$-komórkę i $2$-komórkę
      dołączoną do $S^1\cong\R P^1$
      odwzorowaniem odpowiadającym
      $z\mapsto z^2$ (liczba owinięć wynosi $2$).
      Istotnie, $S^1/(z\sim-z)\to S^1$, $[z]\mapsto z^2$,
      jest homeomorfizmem jako ciągła bijekcja ze zwartej
      przestrzeni do przestrzeni Hausdorffa; odwzorowanie dołączające
      przechodzi przez ten iloraz. Pełny obieg brzegu daje zatem dwa obiegi
      $1$-szkieletu, co wyjaśnia relację $a^2=1$.
      Zatem
\[
 \pi_1(\R P^2)
 \cong\langle a\mid a^2=1\rangle
 \cong\Z/2\Z.
\]
W obu przykładach relacja grupowa pochodzi bezpośrednio
z przebiegu brzegu dołączanej $2$-komórki.
\end{przyklad}

\begin{uwaga}[Granica obecnego rozdziału]
Sympleksy i komórki dają
konkretne przestrzenie, a ich
odwzorowania dołączające kodują sposób
sklejenia. W następnym rozdziale
zorientowanym sympleksom przypiszemy
generatory łańcuchów, a ich ścianom
operator brzegu $\partial$.
Sprawdzimy $\partial^2=0$,
porównamy homologię symplicjalną,
singularną i komórkową oraz
wykorzystamy długie ciągi
z Rozdziału~\ref{ch:algebra-homologiczna}.
Wyższe grupy homotopii i twierdzenie
Hurewicza wymagają już tak
zbudowanej homologii przestrzeni.
\end{uwaga}

\section*{Dalsza lektura}
\addcontentsline{toc}{section}{Dalsza lektura}
Pełne rozwinięcie teorii homotopii,
nakryć, kompleksów CW i
$\Delta$-kompleksów podaje
Allen Hatcher, \emph{Algebraic Topology},
rozdziały 0 i 1 oraz początek rozdziału 2,
\url{https://pi.math.cornell.edu/~hatcher/AT/AT.pdf}.
W tym rozdziale wyprowadziliśmy
podstawowe własności, które
wykorzystamy do budowania homologii.

\clearpage

\chapter{Homologia przestrzeni}
\label{ch:homologia-przestrzeni}

W Rozdziale~\ref{ch:algebra-homologiczna} homologia była
niezmiennikiem \emph{kompleksu łańcuchowego}.
Teraz przypiszemy taki kompleks przestrzeni. Najpierw
użyjemy sympleksów z Rozdziału~\ref{ch:topologia-algebraiczna-podstawy};
potem zastąpimy je wszystkimi ciągłymi obrazami sympleksów,
aby móc badać przestrzenie bez wybranej triangulacji. Na końcu
zobaczymy, dlaczego do obliczeń zwykle wystarczy
kilka komórek CW.

Intuicja jest jedna: cykl to suma zorientowanych kawałków
bez brzegu, a brzeg to cykl, który ogranicza kawałek
o wymiar wyższy. Iloraz
\[
 H_n(X;\Z)=
 \frac{\{\text{$n$-cykle w $X$}\}}
 {\{\text{brzegi $(n+1)$-łańcuchów w $X$}\}}
\]
odróżnia na przykład pętlę będącą brzegiem na tarczy od pętli
obiegającej dziurę. Pracujemy nad $\Z$; zmianę
współczynników omawia
Twierdzenie~\ref{thm:homologia-char-zero},
a konkretny rachunek znajduje się poniżej.

\section{Łańcuchy symplicjalne: orientacja i brzeg}
\label{sec:homologia-symplicjalna}

\begin{definicja}[Kompleks łańcuchów symplicjalnych]
\label{def:homologia-symplicjalna}
Niech $K$ będzie skończonym kompleksem symplicjalnym
z Definicji~\ref{def:kompleks-symplicjalny}.
Ustalmy porządek jego wierzchołków. Grupa
$C_n^\Delta(K)$ jest wolną grupą abelową z bazą
wszystkich $n$-sympleksów zapisanych w porządku
$[v_0,\ldots,v_n]$. Jeśli wierzchołki wypiszemy
w innej kolejności, przyjmujemy zasadę
\[
 [v_{\pi(0)},\ldots,v_{\pi(n)}]
 =\operatorname{sgn}(\pi)[v_0,\ldots,v_n].
\]
Zatem odwrócenie orientacji zmienia znak, a dwa
różne sympleksy nadal dają niezależne generatory.
Dla $n\geq1$ określamy
\begin{equation}\label{eq:brzeg-symplicjalny}
 \partial_n[v_0,\ldots,v_n]
 =\sum_{i=0}^n(-1)^i
 [v_0,\ldots,\widehat v_i,\ldots,v_n],
 \qquad \partial_0=0.
\end{equation}
Rozszerzamy $\partial$ liniowo na sumy
o całkowitych współczynnikach. Dla $n<0$ przyjmujemy $C_n^\Delta(K)=0$.
\end{definicja}

Znak ściany nie jest ozdobą: kierunek odziedziczony
przez brzeg wypełnionego trójkąta musi być spójny.
Dla $e=[v_0,v_1]$ i $t=[v_0,v_1,v_2]$ dostajemy
\[
 \partial e=[v_1]-[v_0],\qquad
 \partial t=[v_1,v_2]-[v_0,v_2]+[v_0,v_1].
\]
Drugą równość można zapisać
$\partial t=[v_0,v_1]+[v_1,v_2]+[v_2,v_0]$.
Jest to obieg brzegu trójkąta.

\begin{figure}[!htbp]
\centering
\begin{tikzpicture}[>=Stealth,line cap=round,line join=round]
 \coordinate (a) at (0,0); \coordinate (b) at (3.3,0);
 \coordinate (c) at (1.35,2.55);
 \fill[manifoldblue!13] (a)--(b)--(c)--cycle;
 \draw[manifoldblue!65!black,line width=1.6pt]
  (a)--(b)--(c)--cycle;
 \draw[->,manifoldred,line width=1.5pt]
  (.7,-.13)--(2.5,-.13);
 \draw[->,manifoldred,line width=1.5pt]
  (3.1,.38)--(1.8,2.07);
 \draw[->,manifoldred,line width=1.5pt]
  (1.04,2.1)--(.13,.39);
 \foreach \p in {a,b,c}{\fill[manifoldgold!85!black] (\p) circle(2.8pt);}
 \node[below left] at (a) {$v_0$};
 \node[below right] at (b) {$v_1$};
 \node[above] at (c) {$v_2$};
 \node[manifoldblue!75!black] at (1.4,.84) {$t$};
 \node[align=left,anchor=west,text width=5.5cm] at (4.1,1.38)
  {$\partial t=e_{01}+e_{12}+e_{20}$\\[3pt]
   zewnętrzna pętla jest brzegiem, gdy wnętrze $t$
   należy do przestrzeni};
\end{tikzpicture}
\caption{Zorientowany brzeg trójkąta. Usunięcie jego
wnętrza pozostawia cykl, który w samym okręgu
nie jest już brzegiem żadnego $2$-łańcucha.}
\label{fig:homologia-trojkat}
\end{figure}

\begin{twierdzenie}[Brzeg brzegu znika]
\label{thm:brzeg-brzegu}
Odwzorowania~\eqref{eq:brzeg-symplicjalny}
spełniają $\partial_{n-1}\partial_n=0$.
Wobec tego
\[
 H_n^\Delta(K):=
 \ker(\partial_n:C_n^\Delta(K)\to C_{n-1}^\Delta(K))
 \big/
 \operatorname{im}(\partial_{n+1})
\]
jest dobrze określoną grupą. Licznik składa
się z \emph{cykli}, mianownik z \emph{brzegów}.
\end{twierdzenie}
\begin{proof}
Dla $n=1$ wynika to z $\partial_0=0$.
Niech $n\geq2$ i weźmy generator
$[v_0,\ldots,v_n]$. Każda jego ściana
o wymiarze $n-2$ powstaje po usunięciu
pewnych dwóch wierzchołków $v_i,v_j$,
$i<j$, dokładnie na dwa sposoby.
Usunięcie najpierw $v_i$, a następnie
$v_j$ (ma on już pozycję $j-1$) daje
współczynnik $(-1)^i(-1)^{j-1}$.
W przeciwnej kolejności współczynnik
wynosi $(-1)^j(-1)^i$.
Ich suma to
$(-1)^{i+j-1}+(-1)^{i+j}=0$.
Dotyczy to każdej ściany z osobna,
więc $\partial^2$ znika na bazie,
a przez liniowość na całej grupie.
Stąd $\operatorname{im}\partial_{n+1}
\subseteq\ker\partial_n$, co pozwala
utworzyć podany iloraz.
\end{proof}

\begin{przyklad}[Odcinek, okrąg, tarcza]
\label{prz:homologia-odcinek-okrag}
Odcinek ma dwa wierzchołki $v_0,v_1$
i jedną krawędź $e$. Odwzorowanie
$\partial_1:\Z e\to\Z v_0\oplus\Z v_1$
przypisuje $me\mapsto m(v_1-v_0)$.
Ma zerowe jądro, więc $H_1^\Delta=0$,
a w ilorazie $H_0^\Delta$ obie klasy
wierzchołków są równe. Otrzymujemy $\Z$.

Teraz weźmy obwód trójkąta, bez wnętrza,
i ustawmy kierunki $e_{01},e_{12},e_{20}$.
Każdy $1$-łańcuch
$ae_{01}+be_{12}+ce_{20}$ jest cyklem
dokładnie wtedy, gdy $a=b=c$
(współczynnik przy każdym wierzchołku
w jej brzegu musi być zerem).
Nie ma tu $2$-sympleksów, a więc
$H_1^\Delta\cong\Z$ z generatorem
$e_{01}+e_{12}+e_{20}$.
Jeśli wstawimy wnętrze $t$, ten sam
cykl staje się $\partial t$ i
$H_1^\Delta$ znika. Rysunek~\ref{fig:homologia-trojkat}
pokazuje dokładnie tę zmianę.
\end{przyklad}

W $\Delta$-kompleksie z
Definicji~\ref{def:delta-kompleks} postępujemy
podobnie: każdy \emph{sympleks charakterystyczny}
jest osobnym generatorem, nawet jeżeli jego
różne ściany mają ten sam obraz. Brzeg
jest sumą obrazów ścian z ich znakami;
po identyfikacji mogą się one wzajemnie
skasować. Tak właśnie sklejony odcinek
z jednym wierzchołkiem ma $\partial e=v-v=0$.

\begin{stwierdzenie}[Składowe i $H_0$]
\label{prop:h0-symplicjalne}
Jeżeli $K$ ma $c$ składowych łukowych,
to $H_0^\Delta(K)\cong\Z^c$.
\end{stwierdzenie}
\begin{proof}
Ponieważ $\partial_0=0$, grupa $H_0$
to wolna grupa na wierzchołkach
modulo relacje
$v_j-v_i=\partial[v_i,v_j]$
dla krawędzi. Idąc po drodze krawędziowej,
utożsamiamy klasy wszystkich wierzchołków
jej składowej. Dwa wierzchołki tej samej
składowej $|K|$ można połączyć taką drogą:
wszystkie wierzchołki jednego sympleksu należą do tej samej
składowej grafu krawędzi. Sympleksy przypisane jednej takiej
składowej tworzą podkompleks. Jego realizacja i jej dopełnienie
są domknięte jako skończone sumy domkniętych sympleksów,
więc są też otwarte. Droga nie może przejść między tymi zbiorami.
Każdy punkt sympleksu można natomiast połączyć odcinkiem z wierzchołkiem.
Składowe grafu odpowiadają zatem dokładnie składowym łukowym $|K|$. Wreszcie żadna
relacja nie łączy różnych składowych,
bo obie końcówki każdej krawędzi leżą
w jednej. W każdej składowej pozostaje
więc jeden wolny generator.
\end{proof}

\section{Łańcuchy singularne i niezmienniczość homotopijna}
\label{sec:homologia-singularna}

Triangulacja bywa trudna do znalezienia. W konstrukcji
singularnej sam wybór \emph{ciągłego} odwzorowania
standardowego sympleksu do przestrzeni tworzy
generator; nie żądamy różnowartościowości ani
ładnego kształtu jego obrazu.

\begin{definicja}[Homologia singularna]
\label{def:homologia-singularna}
\emph{Singularny $n$-sympleks} w $X$ to dowolne
ciągłe odwzorowanie $\sigma:\Delta^n\to X$.
Niech $C_n(X)$ będzie wolną grupą abelową
na wszystkich takich odwzorowaniach; dla $n<0$ przyjmujemy $C_n(X)=0$.
W szczególności
dwa różne odwzorowania o tym samym obrazie są
różnymi generatorami. Niech
$\iota_i:\Delta^{n-1}\to\Delta^n$
będzie afinicznym włączeniem ściany
naprzeciw $i$-tego wierzchołka. Kładziemy
\begin{equation}\label{eq:singular-boundary}
 \partial\sigma=\sum_{i=0}^n(-1)^i
 \sigma\circ\iota_i\quad(n\geq1),
 \qquad \partial:C_0(X)\to0,
 \qquad H_n(X):=H_n(C_\bullet(X),\partial).
\end{equation}
\end{definicja}

Obliczenie $\partial^2=0$ jest słowo w słowo
rachunkiem dwóch usuwanych wierzchołków
z dowodu Twierdzenia~\ref{thm:brzeg-brzegu}.
Teraz pojawia się nawet singularny sympleks
stały: nie odrzucamy go na poziomie
łańcuchów, gdyż później jego wkład może
być brzegiem.

\begin{figure}[!htbp]
\centering
\begin{tikzpicture}[>=Stealth,line cap=round,line join=round]
 \coordinate (a) at (0,0); \coordinate (b) at (2.45,0);
 \coordinate (c) at (.9,1.9);
 \fill[manifoldblue!12] (a)--(b)--(c)--cycle;
 \draw[manifoldblue!70!black,thick] (a)--(b)--(c)--cycle;
 \node[below left] at (a) {$v_0$};
 \node[below right] at (b) {$v_1$};
 \node[above] at (c) {$v_2$};
 \node at (1.05,.55) {$\Delta^2$};
 \draw[->,manifoldred,very thick] (2.95,.95)--(4.28,.95)
  node[midway,above] {$\sigma$};
 \fill[manifoldgreen!12]
  (4.95,.12) .. controls (5.7,-.65) and (6.6,.24) ..
  (7.75,-.08) .. controls (8.8,.3) and (8.3,1.4) ..
  (7.5,1.55) .. controls (6.5,2.2) and (5.2,1.66) ..
  (4.95,.12);
 \draw[manifoldgreen!70!black,thick]
  (4.95,.12) .. controls (5.7,-.65) and (6.6,.24) ..
  (7.75,-.08) .. controls (8.8,.3) and (8.3,1.4) ..
  (7.5,1.55) .. controls (6.5,2.2) and (5.2,1.66) ..
  (4.95,.12);
 \fill[manifoldblue!25,opacity=.7]
  (5.6,.45) .. controls (6.3,.1) and (7.3,.2) .. (7.9,.65)
  .. controls (8.05,1.1) and (7.3,1.6) .. (6.8,1.5)
  .. controls (6.1,1.7) and (5.8,1.05) .. (5.6,.45);
 \draw[manifoldred,thick,->]
  (5.6,.45) .. controls (6.3,.1) and (7.3,.2) .. (7.9,.65);
 \draw[manifoldred,thick,->]
  (7.9,.65) .. controls (8.05,1.1) and (7.3,1.6) .. (6.8,1.5);
 \draw[manifoldred,thick,->]
  (6.8,1.5) .. controls (6.1,1.7) and (5.8,1.05) .. (5.6,.45);
 \foreach \x/\y in {5.6/.45,7.9/.65,6.8/1.5}
  {\fill[manifoldred] (\x,\y) circle(1.8pt);}
 \node[font=\scriptsize,below left] at (5.6,.45) {$\sigma(v_0)$};
 \node[font=\scriptsize,right] at (7.9,.65) {$\sigma(v_1)$};
 \node[font=\scriptsize,above] at (6.8,1.5) {$\sigma(v_2)$};
 \node[manifoldgreen!70!black] at (8.05,1.87) {$X$};
 \node[manifoldblue!70!black] at (6.7,.8) {$\sigma(\Delta^2)$};
\end{tikzpicture}
\caption{Singularny sympleks jest odwzorowaniem do $X$, niekoniecznie
płaskim trójkątem. Jego brzeg liczymy
z ograniczeń odwzorowania $\sigma$ do trzech ścian.}
\label{fig:singular-simplex}
\end{figure}

\begin{stwierdzenie}[Ciągłe odwzorowania działają na homologii]
\label{prop:funktor-singularny}
Jeśli $f:X\to Y$ jest ciągłe, to
$f_\#:C_n(X)\to C_n(Y)$,
$f_\#(\sigma)=f\circ\sigma$,
jest odwzorowaniem łańcuchowym. Wyznacza homomorfizm
$f_*:H_n(X)\to H_n(Y)$, przy czym
$(g\circ f)_*=g_*\circ f_*$ oraz
$(\id_X)_*=\id_{H_n(X)}$.
\end{stwierdzenie}
\begin{proof}
Na pojedynczym generatorze
\[
 \partial(f\circ\sigma)
 =\sum_i(-1)^i f\circ\sigma\circ\iota_i
 =f_\#(\partial\sigma).
\]
Równość na sumach wynika z liniowości.
Zatem obraz cyklu jest cyklem, a obraz
brzegu brzegiem, więc $f_*$ jest
dobrze określony na ilorazie.
Równości dotyczące składania sprawdzamy
już na $\sigma$:
$g_\#f_\#(\sigma)=g\circ f\circ\sigma$.
\end{proof}

\begin{przyklad}[Punkt]
\label{prz:homologia-punktu}
Dla $X=\{*\}$ w każdym stopniu jest
dokładnie jeden singularny sympleks
stały $c_n$, więc $C_n(X)\cong\Z$.
Jego brzeg to
$\partial c_n=(\sum_{i=0}^n(-1)^i)c_{n-1}$:
dla $n$ nieparzystego współczynnik
jest $0$, dla dodatniego parzystego $1$.
Kompleks ma więc postać
\[
 \cdots\xrightarrow{\;0\;}\Z
 \xrightarrow{\;1\;}\Z
 \xrightarrow{\;0\;}\Z
 \xrightarrow{\;1\;}\Z
 \xrightarrow{\;0\;}\Z\longrightarrow0,
\]
gdzie prawy skraj to $C_0$; precyzyjnie
$\partial_1=0$, $\partial_2=1$.
Stąd $H_0(*)=\Z$ i $H_n(*)=0$ dla $n>0$.
\end{przyklad}

\begin{stwierdzenie}[Składowe łukowe]
\label{prop:h0-singularne}
Grupa $H_0(X)$ jest wolna na składowych
łukowych $X$. W szczególności
dla przestrzeni niepustej łukowo spójnej
$H_0(X)\cong\Z$.
\end{stwierdzenie}
\begin{proof}
Każdy punkt $x$ jest singularnym
$0$-sympleksem, więc $C_0(X)$
jest wolną grupą na punktach.
Dla drogi $\gamma:[0,1]=\Delta^1\to X$
mamy $\partial\gamma=[\gamma(1)]-
[\gamma(0)]$. Punkty w jednej składowej
łukowej mają zatem tę samą klasę.
Z drugiej strony każdy składnik
brzegu dowolnego $1$-łańcucha
jest różnicą punktów połączonych drogą:
nie może utożsamić dwóch różnych
składowych. Sumowanie współczynników
punktów osobno w każdej składowej
daje odwrotny homomorfizm do
wolnej grupy na składowych.
\end{proof}

\begin{twierdzenie}[Niezmienniczość homotopijna]
\label{thm:homotopia-singularna}
Jeśli $f,g:X\to Y$ są homotopijne, to
$f_*=g_*:H_n(X)\to H_n(Y)$ w każdym stopniu.
Jeśli $X$ i $Y$ są homotopijnie równoważne,
ich grupy homologii są izomorficzne.
\end{twierdzenie}
\begin{proof}
Niech $F:X\times[0,1]\to Y$ będzie
homotopią, $F(x,0)=f(x)$,
$F(x,1)=g(x)$.
Graniastosłup $\Delta^n\times[0,1]$
dzielimy na $n+1$ zorientowanych
$(n+1)$-sympleksów. Dla $i=0,\ldots,n$
ich uporządkowane wierzchołki to
\[
 (e_0,0),\ldots,(e_i,0),
 (e_i,1),\ldots,(e_n,1).
\]
Niech $T_i:\Delta^{n+1}\to
\Delta^n\times[0,1]$ będzie odwzorowaniem
afinicznym posyłającym kolejne wierzchołki
właśnie na tę listę. Określmy
\[
 P_n(\sigma)=\sum_{i=0}^n(-1)^i\,
 F\circ(\sigma\times\id_{[0,1]})\circ T_i
 \in C_{n+1}(Y).
\]
Rozpiszmy brzeg sumy. Wewnętrzne
$n$-ściany, na których spotykają się
sąsiednie części graniastosłupa,
pojawiają się z przeciwnymi znakami.
Ściany boczne nad każdą ścianą
$\Delta^n$ tworzą dokładnie
$-P_{n-1}(\partial\sigma)$:
jeśli $j<i$, usuwamy dolny wierzchołek $(e_j,0)$
z $T_i$, z łącznym znakiem $(-1)^{i+j}$.
Na pryzmacie ściany naprzeciw $e_j$ punkt podziału ma indeks $i-1$,
więc znak w $-P_{n-1}\partial$ wynosi
$-(-1)^j(-1)^{i-1}=(-1)^{i+j}$.
Jeśli $j>i$, usuwamy górny wierzchołek $(e_j,1)$,
na pozycji $j+1$; oba rachunki dają $(-1)^{i+j+1}$.
Wewnętrzna ściana między $T_{i-1}$ i $T_i$ powstaje
przez usunięcie odpowiednio $(e_{i-1},1)$ i $(e_i,0)$,
ze znakami $-1$ i $+1$.
Pozostałe dwie ściany to górna
$g\circ\sigma$ ze znakiem $+1$
oraz dolna $f\circ\sigma$ ze
znakiem $-1$. Zatem
\begin{equation}\label{eq:prism-identity}
 \partial P_n+P_{n-1}\partial
 =g_\#-f_\#.
\end{equation}
Przyjmujemy $P_{-1}=0$. Dla $n=0$ jest to znana równość
,,koniec drogi minus początek'';
powyższe parowanie ścian dowodzi
jej w każdym wymiarze. W świetle
Twierdzenia~\ref{thm:homotopia-homologia}
oba odwzorowania indukują to samo na homologii.
Jeśli $u:X\to Y$, $v:Y\to X$
są odwrotne z dokładnością do homotopii,
to $(vu)_*=\id$ i $(uv)_*=\id$,
więc $u_*$ i $v_*$ są odwrotne.
\end{proof}

W szczególności przestrzeń ściągalna
ma homologię punktu. Jeśli
$A\subseteq X$ jest retraktem
deformacyjnym, to $H_n(A)\cong H_n(X)$.

\section{Pary przestrzeni i homologia zredukowana}
\label{sec:homologia-wzgledna}

\begin{definicja}[Grupy względne i zredukowane]
\label{def:homologia-wzgledna}
Dla $A\subseteq X$ włączenie daje
podkompleks $C_\bullet(A)\subseteq C_\bullet(X)$.
Kładziemy
\[
 C_n(X,A)=C_n(X)/C_n(A),\qquad
 H_n(X,A)=H_n(C_\bullet(X,A)).
\]
\emph{Augmentacja} $\varepsilon:C_0(X)\to\Z$ sumuje
współczynniki punktów. Dołączamy stopień $C_{-1}=\Z$,
z różniczką $\partial_0=\varepsilon$; niższe grupy są zerowe.
Ponieważ $\varepsilon\partial_1=0$, otrzymujemy kompleks,
którego homologie oznaczamy $\widetilde H_n(X)$.
Dla niepustego $X$ mamy
$\widetilde H_0(X)=\ker(\varepsilon:H_0(X)\to\Z)$,
$\widetilde H_n(X)=H_n(X)$ dla $n>0$ i
$\widetilde H_n(X)=0$ dla $n<0$.
Dla $X=\varnothing$ wszystkie grupy łańcuchów w stopniach
nieujemnych są zerowe, więc $\widetilde H_{-1}(\varnothing)=\Z$
i pozostałe grupy zredukowane są zerowe.
\end{definicja}

Względny cykl może mieć brzeg w $A$:
to, co w $X$ jest otwartą powierzchnią,
po ,,pominięciu $A$'' staje się cyklem.
Na przykład zorientowany dysk reprezentuje klasę względną
względem swojego okręgu brzegowego.
Jeżeli $X$ ma skończoną liczbę $c>0$ składowych łukowych,
to $\widetilde H_0(X)\cong\Z^{c-1}$.
Ogólnie, po wyborze jednej składowej $P_0$, bazę jądra
augmentacji tworzą różnice $[P]-[P_0]$ dla $P\ne P_0$:
jest to suma prosta kopii $\Z$, także dla nieskończenie wielu składowych.

\begin{twierdzenie}[Długi ciąg pary]
\label{thm:les-pary}
Każda para $(X,A)$ daje naturalny
długi ciąg dokładny
\begin{equation}\label{eq:les-pary}
 \cdots\to H_n(A)\xrightarrow{i_*}H_n(X)
 \xrightarrow{j_*}H_n(X,A)
 \xrightarrow{\delta}H_{n-1}(A)
 \to H_{n-1}(X)\to\cdots.
\end{equation}
Jeśli $\partial c\in C_{n-1}(A)$,
to $\delta[c]=[\partial c]$.
\end{twierdzenie}
\begin{proof}
W każdej grupie łańcuchów mamy
krótki ciąg dokładny
$0\to C_n(A)\to C_n(X)\to C_n(X,A)\to0$.
Jest zgodny z brzegami, więc
Stwierdzenie~\ref{prop:ses-iloraz}
daje długi ciąg i wzór na $\delta$.
Naturalność pochodzi z przemienności
odwzorowań łańcuchów dla odwzorowań
par $(X,A)\to(Y,B)$.
\end{proof}

Jeśli $A\ne\varnothing$, możemy w tym ciągu zastąpić homologie
$A$ i $X$ zredukowanymi, pozostawiając zwykłą homologię względną.
Istotnie, w stopniu $-1$ włączenie kompleksów augmentowanych
jest identycznością $\Z\to\Z$, więc ich iloraz ma tam zero
i jest właśnie kompleksem $C_\bullet(X,A)$.

Homotopia par (w każdym czasie obraz $A$
leży w $B$) daje te same odwzorowania na $H_n(X,A)$:
pryzmat z dowodu Twierdzenia~\ref{thm:homotopia-singularna}
przenosi $C_\bullet(A)$ do $C_{\bullet+1}(B)$,
więc wzór~\eqref{eq:prism-identity}
schodzi na kompleksy ilorazowe.

\section{Podział barycentryczny, wycinanie i sklejanie}
\label{sec:wycinanie-homologia}

Homologia singularna ma nieskończenie wiele generatorów nawet
dla prostych przestrzeni. Podział dużego sympleksu na małe
części pozwala jednak sprowadzić rachunek do lokalnych
fragmentów. Z tego mechanizmu wynika wycinanie.

\begin{lemat}[Małe łańcuchy]
\label{lem:male-lancuchy}
Niech $\mathcal U=\{U_\alpha\}$ będzie otwartym
pokryciem $X$. Podkompleks
$C_n^{\mathcal U}(X)\subseteq C_n(X)$
generowany przez sympleksy, których obrazy
leżą w jednym $U_\alpha$, ma te same
grupy homologii co $C_\bullet(X)$.
Włączenie indukuje izomorfizmy w każdym stopniu.
\end{lemat}
\begin{proof}
Podamy konstrukcję potrzebną w obu kierunkach.
Podzielmy liniowo $\Delta^n$ na sympleksy
barycentryczne: wierzchołkami małego sympleksu
są barycentra łańcucha niepustych ścian
$F_0\subsetneq F_1\subsetneq\cdots\subsetneq F_n$,
zapisane w tej właśnie kolejności. Każdy taki łańcuch odpowiada
permutacji $\pi$ wierzchołków: $F_j=[e_{\pi(0)},\ldots,e_{\pi(j)}]$.
Niech $a_\pi:\Delta^n\to\Delta^n$ będzie afinicznym odwzorowaniem
na uporządkowany mały sympleks. Określamy
\[
 S_n(\sigma)=\sum_{\pi\in S_{n+1}}\operatorname{sgn}(\pi)
                 \,\sigma\circ a_\pi.
\]
Znak pochodzi z orientacji: dla permutacji identycznościowej
macierz kolejnych barycentrów ma dodatni wyznacznik w otoczce afinicznej,
a permutacja wierzchołków zmienia go o $\operatorname{sgn}(\pi)$.
Kolejność barycentrów jest ważna: w kompleksie singularnym
odwrócenie parametryzacji nie jest z definicji ujemnym generatorem.
Ściany wewnętrzne dwóch części
występują z przeciwnymi znakami,
a ich ściany na zewnątrz dają
podział odpowiedniej ściany $\sigma$.
Dokładniej, wewnętrzną ścianę uzyskujemy, pomijając barycentrum $F_j$
dla $j<n$. Dwie odpowiadające jej permutacje zamieniają miejscami
pozycje $j,j+1$, więc mają przeciwne znaki, a pozostałe
uporządkowane barycentra są identyczne. Dla $j=n$ otrzymujemy
sympleks podziału ściany naprzeciw $e_{\pi(n)}$.
Jeśli $\rho$ jest permutacją pozostałych wierzchołków w ich
odziedziczonej kolejności, to
$(-1)^n\operatorname{sgn}(\pi)=(-1)^{\pi(n)}\operatorname{sgn}(\rho)$.
To właśnie znak tej ściany w $S\partial$; stąd $\partial S=S\partial$.

Wyjaśnijmy też, dlaczego $S$ nie zmienia
klasy homologii. Na standardowym sympleksie
oznaczmy jego identycznościowe odwzorowanie
przez $u_n$ i indukcyjnie zbudujmy
łańcuchy $t_n$ tak, aby dla $n\geq1$
\[
 \partial t_n+\sum_{i=0}^n(-1)^i(\iota_i)_\#t_{n-1}
       =S(u_n)-u_n. \tag{*}
\]
W stopniu $0$ obie strony są zerowe,
więc $t_0=0$. Jeśli mamy już konstrukcję
na wszystkich ścianach, to oznaczmy ich sumę
$b_n=\sum_i(-1)^i(\iota_i)_\#t_{n-1}$.
Łańcuch $z=S(u_n)-u_n-b_n$
jest cyklem: stosując $\partial$,
używamy $\partial S=S\partial$,
równości (*) na ścianach
i $\partial^2=0$.
W wypukłym $\Delta^n$ taki cykl jest
brzegiem: z barycentrum $b$ łączymy
każdy jego singularny sympleks prostym
stożkiem. Dla $\tau:\Delta^k\to\Delta^n$ stożek $b*\tau$
ma wzór
\[
 (b*\tau)(s_0,\ldots,s_{k+1})
 =s_0b+(1-s_0)\tau\left(\frac{s_1}{1-s_0},\ldots,
                         \frac{s_{k+1}}{1-s_0}\right)
 \quad(s_0<1),
\]
z wartością $b$ dla $s_0=1$. Jest ciągły także w tym wierzchołku,
bo obraz $\tau$ jest ograniczony. Brzeg tego stożka to
$\tau-b*(\partial\tau)$ (dla $0$-sympleksu
dodatkowo odejmujemy punkt $b$);
po zsumowaniu nad cyklem składniki
stożkowe się kasują. Daje to $t_n$.
Przenosząc $t_n$ przez dowolne
$\sigma:\Delta^n\to X$, otrzymujemy
operator $T_n(\sigma)=\sigma_\#t_n$ o nośniku zawartym
w $\sigma(\Delta^n)$, który spełnia
\[
 \partial T+T\partial=S-\id.
\]
Zatem $S^r-\id=\partial T^{(r)}+
T^{(r)}\partial$, gdzie
$T^{(r)}=T+ST+\cdots+S^{r-1}T$.

Średnice części po kolejnych podziałach
barycentrycznych dążą do zera: dla
$n$-sympleksu każdy nowy wierzchołek
jest średnią wierzchołków pewnej ściany,
a nowa krawędź łączy barycentra zagnieżdżonych ścian.
Jeśli mają one $p<q\leq n+1$ wierzchołków, to
$b_G=(p/q)b_F+((q-p)/q)b_{G\setminus F}$.
Zatem $\|b_G-b_F\|\leq(q-p)\operatorname{diam}(\Delta)/q
\leq n\operatorname{diam}(\Delta)/(n+1)$.
Średnica wypukłego sympleksu jest największą odległością jego
wierzchołków, więc dla $n\geq1$ otrzymujemy geometryczny spadek
średnic (także na ścianach). Dla $n=0$ sympleks jest już punktem. Dla danego $\sigma$
przeciwobrazy $U_\alpha$ tworzą
otwarte pokrycie zwartego $\Delta^n$;
jego liczba Lebesgue'a jest dodatnia.
Po dostatecznie wielu podziałach
każda część trafia przez $\sigma$
do jednego $U_\alpha$.
Każdy łańcuch jest skończoną sumą,
więc można wybrać wspólne $r$
dla wszystkich jego sympleksów.

Jeśli $z$ jest cyklem, $S^rz$
jest mały i różni się od $z$
o $\partial T^{(r)}z$; to dowodzi
surjektywności na homologii.
Jeśli mały cykl $z$ jest brzegiem
$c$ w pełnym kompleksie, podzielmy
$c$ tak wiele razy, by $S^rc$
był mały. Łańcuch $T^{(r)}z$
też jest mały, ponieważ każda
jego część leży w obrazie
jednego pierwotnego, małego sympleksu.
Równości $\partial c=z$ i
$\partial T^{(r)}z=S^rz-z$ dają
\[
 \partial\bigl(S^rc-T^{(r)}z\bigr)=z.
\]
Zatem $z$ jest brzegiem już
w podkompleksie małych łańcuchów:
otrzymujemy iniektywność.
\end{proof}

Stożkowanie wypełnia każdy cykl dodatniego stopnia w wypukłym
sympleksie, a w stopniu zero tylko cykl o sumie współczynników zero.
W konstrukcji $T$ stopień zero załatwiamy przez $t_0=0$,
bo $S_0=\id$. Lokalność $T$ jest istotna dla dowodu iniektywności.

\begin{twierdzenie}[Wycinanie]
\label{thm:wycinanie}
Jeżeli $Z\subseteq A\subseteq X$ oraz
$\overline Z\subseteq\operatorname{int}_X A$,
to włączenie par
\[
 (X\setminus Z,A\setminus Z)\longrightarrow(X,A)
\]
indukuje izomorfizm na wszystkich
grupach homologii względnej.
\end{twierdzenie}
\begin{proof}
Oznaczmy $U=X\setminus\overline Z$
oraz $V=\operatorname{int}_X A$.
To otwarte pokrycie $X$;
po ograniczeniu do $X\setminus Z$
otrzymujemy pokrycie
$\{U,V\setminus Z\}$.
Lemat~\ref{lem:male-lancuchy} działa
także po wzięciu ilorazu przez
łańcuchy w $A$: izomorfizmy
dla $A$ i $X$, zestawione
z długimi ciągami par, dają
izomorfizm dla odpowiadających
kompleksów względnych
(Lemat~\ref{lem:pieciu}).

Policzmy teraz te małe kompleksy.
Ponieważ $V\subseteq A$,
\[
 C_\bullet^{\{U,V\}}(X)
   =C_\bullet(U)+C_\bullet(V),\qquad
 C_\bullet^{\{U,V\}}(A)
   =C_\bullet(U\cap A)+C_\bullet(V).
\]
Ich iloraz jest kanonicznie
$C_\bullet(U)/C_\bullet(U\cap A)$:
sympleks należący jednocześnie do
$U$ i $V$ należy do $U\cap A$,
więc jądro tej identyfikacji jest
dokładnie wskazanym podkompleksem.
Analogiczny iloraz dla
$(X\setminus Z,A\setminus Z)$
i jego pokrycia jest \emph{tym samym}
kompleksem
$C_\bullet(U)/C_\bullet(U\cap A)$.
Włączenie par staje się na tych
modelach identycznością, co
kończy dowód.
\end{proof}

Wycinanie mówi, że usunięcie części
leżącej głęboko we wnętrzu $A$
nie zmienia tego, co para $(X,A)$
mierzy \emph{poza} $A$. Warunek
o domknięciu ma znaczenie:
zapewnia otwarte pokrycia użyte
w dowodzie.

\begin{twierdzenie}[Ciąg Mayera--Vietorisa]
\label{thm:mv-przestrzenie}
Jeśli $X=U\cup V$, gdzie $U,V$
są otwarte w $X$, to istnieje
długi ciąg dokładny
\begin{equation}\label{eq:mv-przestrzenie}
 \cdots\to H_n(U\cap V)
 \xrightarrow{(i_*,-j_*)}
 H_n(U)\oplus H_n(V)
 \xrightarrow{k_*+\ell_*}H_n(X)
 \xrightarrow{\delta}H_{n-1}(U\cap V)\to\cdots.
\end{equation}
Tak samo działają grupy zredukowane.
\end{twierdzenie}
\begin{proof}
Wewnątrz $C_\bullet(X)$ rozważmy
$B_\bullet=C_\bullet(U)+C_\bullet(V)$.
Mamy krótki ciąg kompleksów
\[
 0\to C_\bullet(U\cap V)
 \xrightarrow{c\mapsto(c,-c)}
 C_\bullet(U)\oplus C_\bullet(V)
 \xrightarrow{(a,b)\mapsto a+b}
 B_\bullet\to0.
\]
Iniektywność pierwszego odwzorowania jest jasna;
jądro drugiej składa się z par
$(c,-c)$, w których każdy singularny
sympleks występujący w $c$ należy
zarazem do $U$ i do $V$.
Surjektywność wynika z definicji sumy.
Ponadto $B_\bullet=C_\bullet^{\{U,V\}}(X)$,
więc Lemat~\ref{lem:male-lancuchy}
identyfikuje $H_*(B)$ z $H_*(X)$.
Zastosowanie Twierdzenia~\ref{thm:les-homologia}
daje~\eqref{eq:mv-przestrzenie}.
W wersji zredukowanej dołączamy
do każdego kompleksu stopień $-1$
równy $\Z$ i różniczkę
$\varepsilon:C_0\to\Z$.
Na stopniu $-1$ ten sam krótki
ciąg ma postać
$0\to\Z\xrightarrow{c\mapsto(c,-c)}
\Z\oplus\Z\xrightarrow{(a,b)\mapsto a+b}
\Z\to0$, nadal dokładną.
W stopniach nieujemnych
obowiązuje poprzedni rachunek.
Homologie tych powiększonych
kompleksów są właśnie
homologiami zredukowanymi;
dotyczy to również pustego
przecięcia, dla którego
$\widetilde H_{-1}(\varnothing)=\Z$.
\end{proof}

Jeśli $c=a+b\in B_n$ jest cyklem,
gdzie $a\in C_n(U)$ i $b\in C_n(V)$,
to $\partial a=-\partial b$ jest
łańcuchem w $U\cap V$, a
$\delta[c]=[\partial a]$ z wyborem
znaku zgodnym z pierwszym odwzorowaniem.
To praktyczna recepta na odwzorowanie łączące.

\begin{przyklad}[Okrąg z dwóch łuków]
\label{prz:mv-okrag}
Pokryjmy $S^1$ dwoma otwartymi,
ściągalnymi łukami $U,V$, których
przecięcie składa się z dwóch
łuków. Mamy
$\widetilde H_0(U)=
\widetilde H_0(V)=0$,
a $\widetilde H_0(U\cap V)\cong\Z$
generuje różnica klas punktów
z dwóch składowych.
Zredukowany ciąg
Mayera--Vietorisa daje
\[
 0=\widetilde H_1(U)\oplus
 \widetilde H_1(V)\to
 \widetilde H_1(S^1)
 \xrightarrow[\cong]{\delta}
 \widetilde H_0(U\cap V)\to0.
\]
Zatem $H_1(S^1)\cong\Z$.
Geometria odwzorowania $\delta$ jest prosta:
rozcinając cykl okręgu na łuki
leżące w $U$ i $V$, otrzymujemy
dwa końce w różnych składowych
ich przecięcia.
\end{przyklad}

\begin{twierdzenie}[Homologia sfer]
\label{thm:homologia-sfer}
Dla $n\geq0$:
\[
 \widetilde H_k(S^n;\Z)=
 \begin{cases}\Z,&k=n,\\0,&k\ne n.\end{cases}
\]
Ponadto dla $n\geq1$
\[
 H_k(D^n,S^{n-1};\Z)=
 \begin{cases}\Z,&k=n,\\0,&k\ne n.\end{cases}
\]
\end{twierdzenie}
\begin{proof}
$S^0$ ma dwie składowe punktowe:
$H_0(S^0)=\Z^2$ oraz
$\widetilde H_0(S^0)=\Z$.
Wyższa homologia znika, bo jest
to suma dwóch kopii homologii
punktu z Przykładu~\ref{prz:homologia-punktu}.

Dla $n\geq1$ weźmy lekko
powiększone otwarte półkule
$U,V\subset S^n$. Każda jest
ściągalna, a $U\cap V$
deformacyjnie zwija się na
sferę równikową $S^{n-1}$.
Zredukowany ciąg
Mayera--Vietorisa, wraz z
zerowaniem homologii $U,V$,
daje dla każdego $k$ izomorfizm
\[
 \widetilde H_k(S^n)
 \xrightarrow{\ \cong\ }
 \widetilde H_{k-1}(S^{n-1}).
\]
W stopniu $k=1$ używa się
właśnie zredukowanego $H_0$;
gdy $n=1$, obie składowe
przecięcia są niezbędne.
Indukcja od $S^0$ daje pierwszy wzór.

Dysk $D^n$ jest ściągalny. W ciągu
zredukowanym pary $(D^n,S^{n-1})$
grupy dysku są zerowe, a odwzorowanie
łączące daje
$H_k(D^n,S^{n-1})\cong
\widetilde H_{k-1}(S^{n-1})$
dla $k\geq1$. Dla $k=0$
odwzorowanie $H_0(S^{n-1})\to H_0(D^n)$
jest surjekcją, więc grupa
względna również znika.
To daje drugi wzór.
\end{proof}

\begin{uwaga}[Klasa fundamentalna i orientacja]
Generator $H_n(D^n,S^{n-1})\cong\Z$
nazywamy \emph{klasą orientacji dysku}.
Wybór jednego z jego dwóch znaków
określa orientację; w długim ciągu
pary $\delta$ posyła ją na
generator $\widetilde H_{n-1}(S^{n-1})$.
Różniczkowe ujęcie orientacji
i całkowania pozostawiamy
osobnemu rozdziałowi.
\end{uwaga}

\section{Od triangulacji do homologii przestrzeni}
\label{sec:porownanie-symplicjalne}

Dla kompleksu symplicjalnego zachowujemy ustalony porządek
wierzchołków, a dla $\Delta$-kompleksu porządki jego map charakterystycznych.
Każdy bazowy sympleks $\sigma:\Delta^n\to|K|$ jest wtedy
jednocześnie singularnym sympleksem. Daje to odwzorowanie łańcuchowe
\[
 I:C_\bullet^\Delta(K)\longrightarrow C_\bullet(|K|),
 \qquad [\sigma]\longmapsto \sigma,
\]
bo oba operatory $\partial$ sumują te same
ściany z tymi samymi znakami.
Mapę $I$ rozszerzamy liniowo: na przykład $I(-[\sigma])=-\sigma$.
Nie zastępujemy prawej strony osobnym generatorem otrzymanym
przez odwrócenie parametryzacji. Nie jest to
izomorfizm grup łańcuchów: po prawej są
wszystkie odwzorowania ciągłe, po lewej tylko
wybrane sympleksy. Niezwykłe jest to,
że po przejściu do homologii różnica znika.

Potrzebny będzie prosty sposób traktowania
komórki z całym brzegiem zwiniętym do punktu.

\begin{lemat}[Iloraz pary komórkowej]
\label{lem:iloraz-cw-homologia}
Niech $A\ne\varnothing$ będzie podkompleksem skończonego
kompleksu CW $X$. Jeśli istnieje otwarte
otoczenie $N$ zbioru $A$, które silnie
deformacyjnie retrahuje się na $A$, to
odwzorowanie ilorazowe indukuje izomorfizm
\[
 H_k(X,A)\xrightarrow{\ \cong\ }
 H_k(X/A,\{*\})=\widetilde H_k(X/A)
 \quad(k\geq0).
\]
Takie otoczenie istnieje, gdy
$A=X^{n-1}\ne\varnothing$ i $X=X^n$ jest skończonym
szkieletem CW dla $n\geq1$, a także gdy do skończonego $A$
dołączamy pojedynczą komórkę dodatniego wymiaru.
\end{lemat}
\begin{proof}
Włączenie $A\hookrightarrow N$ jest
równoważnością homotopijną. Porównując
długie ciągi par $(X,A)$ i $(X,N)$
za pomocą Lematu~\ref{lem:pieciu},
dostajemy $H_k(X,A)\cong H_k(X,N)$.
Po zwinięciu $A$ do punktu homotopia schodzi na $N/A$.
Uzasadnijmy jej ciągłość również względem czasu.
Skończony kompleks CW jest zwarty, bo jest obrazem skończonej
sumy zwartych dysków; domknięty $A$ też jest zwarty.
Projekcja $q:N\to N/A$ jest domknięta: dla domkniętego $F\subseteq N$
zbiór $q^{-1}q(F)$ jest równy $F$ lub $F\cup A$.
Ma zwarte włókna, więc $q\times\id_I$ jest również domknięta.
Istotnie, jeśli domknięty zbiór omija włókno nad $(y,t)$,
skończone pokrycie tego włókna prostokątami otwartymi daje
otoczenia $O\supseteq q^{-1}(y)$ i $J\ni t$ z $O\times J$
rozłącznym z tym zbiorem; domkniętość $q$ daje otoczenie
$W\ni y$ takie, że $q^{-1}(W)\subseteq O$.
Stąd $(q\times\id_I)$ jest odwzorowaniem ilorazowym,
co pozwala przenieść homotopię. Zatem $N/A$ silnie zwija się do $\{*\}$,
więc tak samo
$H_k(X/A,\{*\})\cong H_k(X/A,N/A)$.

Porównajmy środkowe grupy.
Zbiór $A$ jest domknięty
(dla skończonego CW wynika to
z topologii słabej), zatem
$\{X\setminus A,N\}$ jest
otwartym pokryciem $X$.
W ilorazie mamy pokrycie
$\{(X/A)\setminus\{*\},N/A\}$.
Lemat o małych łańcuchach,
zastosowany również do $N$,
pozwala obliczać homologie
tych dwóch par odpowiednio za pomocą kompleksów
\[
 \frac{C_\bullet(X\setminus A)}
      {C_\bullet(N\setminus A)}
 \quad\text{oraz}\quad
 \frac{C_\bullet((X/A)\setminus\{*\})}
      {C_\bullet((N/A)\setminus\{*\})}.
\]
Poza $A$ odwzorowanie ilorazowe jest
homeomorfizmem, więc te dwa
kompleksy są izomorficzne.
Składając trzy izomorfizmy,
otrzymujemy tezę.

Dla $A=X^{n-1}\subset X^n$
weźmy $N$ złożony z $A$ oraz,
w każdej dołączonej $n$-komórce,
obrazu pierścienia
$\{x\in D^n:\|x\|>1/2\}$.
Jest otwarty w topologii
nadanej przez odwzorowania dołączające.
W każdym pierścieniu promieniowo
przesuwamy $r$ od wartości
$r>1/2$ do $1$, pozostawiając
brzeg i cały $A$ nieruchome.
Na częściach dyskowych używamy wzoru
$ru\mapsto((1-t)r+t)u$. Ciągłość przy $A$ sprawdzamy jak
w dowodzie Twierdzenia~\ref{thm:pi1-dolaczenie-komorki}:
w każdym otoczeniu zachowujemy tylko te punkty pasa, których
cały odcinek radialny do brzegu mieści się w tym otoczeniu.
Te mniejsze zbiory są otwarte i pozostają w pierwotnym otoczeniu
przez cały czas homotopii. Daje to silną retrakcję $N\to A$.
Ten sam argument działa przy dołączeniu pojedynczej komórki.
\end{proof}

\begin{lemat}[Generator pojedynczego sympleksu]
\label{lem:orientacja-simpleksu}
Identycznościowe odwzorowanie singularne
$u_n:\Delta^n\to\Delta^n$ określa
generator
$H_n(\Delta^n,\partial\Delta^n)\cong\Z$
dla $n\geq1$.
\end{lemat}
\begin{proof}
Jest cyklem względnym, bo jego
brzeg leży w $\partial\Delta^n$.
Brzeg sympleksu jest sferą
(Stwierdzenie~\ref{prop:brzeg-sympleksu}),
a jej homologię policzyliśmy
w Twierdzeniu~\ref{thm:homologia-sfer}.
Wynika stąd, że odwzorowanie łączące
\[
 \delta:H_n(\Delta^n,\partial\Delta^n)
 \longrightarrow
 \widetilde H_{n-1}(\partial\Delta^n)
 \]
jest izomorfizmem i posyła
$[u_n]$ na klasę cyklu
$\partial u_n$.
Pokażemy indukcyjnie, że ten
cykl jest generatorem.
Dla $n=1$ jest to różnica
końców, generator
$\widetilde H_0(S^0)$.
W wyższym wymiarze rozdzielmy
$\partial\Delta^n$ na jedną
ścianę $F\cong\Delta^{n-1}$
i sumę $B$ pozostałych ścian.
$F$ i $B$ są ściągalne:
$B$ jest stożkiem z przeciwległego
wierzchołka nad $\partial F$;
ich przecięciem jest $\partial F$.
W ciągu Mayera--Vietorisa
odwzorowanie łączące
$\widetilde H_{n-1}(\partial\Delta^n)
\to\widetilde H_{n-2}(\partial F)$
jest izomorfizmem. Aby zastosować wersję dla zbiorów otwartych, przedstawmy
oba dyski $F,B$ jako stożki nad ich wspólnym brzegiem $\partial F$.
Niech $r\in[0,1]$ oznacza w każdym z nich parametr radialny,
z $r=1$ na $\partial F$. Zbiory
$U=F\cup\{x\in B:r(x)>1/2\}$ oraz
$V=B\cup\{x\in F:r(x)>1/2\}$ są otwarte w $F\cup B$:
na każdej z dwóch części ich przeciwobrazy są otwarte.
Przesuwając dodatkowe pasy do $r=1$, retrahujemy je na $F,B$,
a ich przecięcie na $\partial F$.
Cykl $\partial u_n$ jest sumą
$\pm u_F$ na $F$ i pozostałych
ścian na $B$; jego obrazem przez
odwzorowanie łączące jest
$\pm[\partial u_F]$.
To generator na mocy
założenia indukcyjnego.
Zatem $\partial u_n$, a więc
również $u_n$ względnie,
jest generatorem.
\end{proof}

\begin{twierdzenie}[Porównanie homologii]
\label{thm:symplicjalna-singularna}
Dla skończonego kompleksu symplicjalnego
$K$ naturalne odwzorowanie indukuje
\[
 I_*:H_n^\Delta(K)\xrightarrow{\ \cong\ }
 H_n(|K|;\Z)
 \qquad(n\geq0).
\]
To samo zachodzi dla skończonego
$\Delta$-kompleksu, jeśli generatory
po lewej to jego sympleksy
charakterystyczne z ich ścianami.
\end{twierdzenie}
\begin{proof}
Uporządkujmy sympleksy tak, aby
każdy pojawiał się po wszystkich
swoich ścianach. Budujemy kompleks
po jednym sympleksie:
$L\subset K=L\cup_\varphi\Delta^n$.
Załóżmy indukcyjnie, że porównanie
jest izomorfizmem dla $L$.
Iloraz kompleksów symplicjalnych
$C^\Delta(K)/C^\Delta(L)$
ma jeden generator w stopniu $n$
i zero w pozostałych stopniach,
więc jego homologia jest $\Z$
tylko w stopniu $n$.

Jeśli $n\geq1$, po zwinięciu $L$
do punktu przestrzeń $|K|/|L|$
jest $\Delta^n/\partial\Delta^n
\cong S^n$; analogiczny opis
działa dla $\Delta$-kompleksu,
nawet gdy jego ściany były
między sobą utożsamione.
Lemat~\ref{lem:iloraz-cw-homologia}
i Twierdzenie~\ref{thm:homologia-sfer}
dają $H_q(|K|,|L|)=\Z$ dla
$q=n$ i zero w innych stopniach.
Odwzorowanie względne przenosi jedyny
generator na klasę charakterystycznego sympleksu.
Mapa par $(\Delta^n,\partial\Delta^n)\to(|K|,|L|)$
indukuje homeomorfizm ich ilorazów: jest ciągłą bijekcją
ze zwartej przestrzeni do przestrzeni Hausdorffa.
Naturalność map do ilorazów w Lemacie~\ref{lem:iloraz-cw-homologia}
pokazuje, że także na homologii względnej jest izomorfizmem.
Na mocy Lematu~\ref{lem:orientacja-simpleksu} przenosi więc
generator na generator, a nie na jego dowolną wielokrotność.
Dla $n=0$ nowy wierzchołek jest
osobną składową i odpowiednie
odwzorowanie względne jest izomorfizmem
$\Z\to\Z$ w stopniu $0$.

Mamy teraz przemienny diagram
długich ciągów par dla kompleksów
symplicjalnych i singularnych.
Na wyrazach dotyczących $L$
odwzorowania pionowe są izomorfizmami
z założenia indukcyjnego,
a na względnych właśnie
to pokazaliśmy. Lemat~\ref{lem:pieciu}
zastosowany do kolejnych pięciu
wyrazów daje izomorfizm dla $K$.
Zaczynamy od kompleksu pustego,
dla którego wszystkie zwykłe
grupy homologii są zerowe.
\end{proof}

W praktyce wolno zatem obliczać
homologię przestrzeni z jednej
wygodnej triangulacji: otrzymany
wynik nie zależy od wyboru
triangulacji. Podziały sympleksów
zmieniają grupy łańcuchów,
ale nie grupy homologii.

\section{Homologia komórkowa}
\label{sec:homologia-komorkowa}

W triangulacji torusa potrzeba
trójkątów i ich ścian. Struktura
CW z Sekcji~\ref{sec:cw-kompleksy}
może opisywać go jednym punktem,
dwoma pętlami i jedną
$2$-komórką. Homologia komórkowa
zachowuje tę oszczędność.

\begin{twierdzenie}[Względna homologia warstwy komórek]
\label{thm:warstwa-komorek}
Niech $X$ będzie skończonym
kompleksem CW, $X^n$ jego
$n$-szkieletem i $X^{-1}=\varnothing$.
Dla $n\geq1$:
\[
 H_k(X^n,X^{n-1})
 \cong
 \begin{cases}
 \displaystyle\bigoplus_{\alpha\in I_n}\Z[e_\alpha^n],
 &k=n,\\[3pt]0,&k\ne n,
 \end{cases}
\]
gdzie $I_n$ oznacza zbiór $n$-komórek.
Dla $n=0$ mamy
$H_0(X^0)=\bigoplus_{\alpha\in I_0}\Z[e_\alpha^0]$
i $H_k(X^0)=0$ dla $k\ne0$.
\end{twierdzenie}
\begin{proof}
Jeśli $X^{n-1}=\varnothing$ i $n\geq1$, to również $X^n=\varnothing$:
nie można dołączyć dysku o niepustym brzegu do pustej przestrzeni.
Teza jest wtedy oczywista. W pozostałych przypadkach
w ilorazie $X^n/X^{n-1}$
brzeg każdej dołączonej kopii
$D^n$ staje się jednym punktem,
a wnętrza komórek pozostają
rozłączne. Z definicji topologii
CW otrzymujemy bukiet
$\bigvee_{\alpha\in I_n}
(D^n/\partial D^n)\cong
\bigvee_{\alpha\in I_n}S^n$.
Lemat~\ref{lem:iloraz-cw-homologia}
zamienia homologię względną
na zredukowaną homologię bukietu.
Dla skończonego bukietu
zredukowana homologia jest sumą
zredukowanych homologii składników:
dołączamy kolejne sfery. Dla bukietu $Y\vee S^n$
bierzemy $U=Y\cup B$ i $V=S^n\cup N$, gdzie $B$ jest małą
otwartą kulą wokół punktu wspólnego w nowej sferze, a $N$
jest sumą takich kul w dotychczasowych sferach.
Te zbiory są otwarte w topologii ilorazowej bukietu,
retrahują się na $Y,S^n$, a $U\cap V=N\vee B$ jest ściągalne.
Ciąg Mayera--Vietorisa daje izomorfizm indukowany włączeniami
składników, a więc właśnie ich sumę prostą.
Twierdzenie~\ref{thm:homologia-sfer}
daje wskazane grupy.
Znaki generatorów zależą
od wyboru orientacji komórek.
Dla $n=0$ mamy skończony
zbiór punktów; rachunek wynika
z homologii punktu.
\end{proof}

\begin{definicja}[Kompleks komórkowy]
\label{def:kompleks-komorkowy}
Kładziemy
$C_n^{\mathrm{CW}}(X)=H_n(X^n,X^{n-1})$.
Dla $n\geq1$ jego operator brzegu
jest złożeniem
\begin{equation}\label{eq:cellular-boundary}
 H_n(X^n,X^{n-1})
 \xrightarrow{\ \delta_n\ }
 H_{n-1}(X^{n-1})
 \xrightarrow{\ j_{n-1}\ }
 H_{n-1}(X^{n-1},X^{n-2}),
\end{equation}
gdzie $\delta_n$ pochodzi z
długiego ciągu pary, a
$j_{n-1}$ z ilorazu łańcuchów.
Przyjmujemy $\partial_0^{\mathrm{CW}}=0$.
\end{definicja}

\begin{twierdzenie}[Komórki obliczają homologię]
\label{thm:komorkowa-singularna}
Operator~\eqref{eq:cellular-boundary}
spełnia $(\partial^{\mathrm{CW}})^2=0$.
Dla skończonego CW istnieją
izomorfizmy
\[
 H_n(C_\bullet^{\mathrm{CW}}(X),
       \partial^{\mathrm{CW}})
 \cong H_n(X;\Z)
 \qquad\text{dla każdego }n\geq0.
\]
\end{twierdzenie}
\begin{proof}
Dla $n=1$ równość brzegu brzegu wynika z $\partial_0^{\mathrm{CW}}=0$.
Dla $n\geq2$ z dokładności długiego ciągu
pary $(X^{n-1},X^{n-2})$
wynika $\delta_{n-1}j_{n-1}=0$.
Stąd
\[
 \partial^{\mathrm{CW}}_{n-1}
 \partial^{\mathrm{CW}}_n
 =j_{n-2}\delta_{n-1}j_{n-1}\delta_n=0.
\]

Rozpiszmy porównanie homologii.
Indukcja po szkieletach,
z użyciem długiego ciągu pary
i Twierdzenia~\ref{thm:warstwa-komorek},
daje $H_k(X^m)=0$ dla $k>m$.
Włączenie $X^{m-1}\hookrightarrow X^m$
indukuje ponadto izomorfizm na
$H_k$ dla $k<m-1$: dwa sąsiednie
wyrazy względne o stopniach
$k+1$ i $k$ są wtedy zerowe.
W szczególności komórki wymiaru
$\geq n+2$ nie zmieniają $H_n$.

Dla $n\geq1$ w ciągu pary
$(X^n,X^{n-1})$ mamy odcinek
\[
 0=H_n(X^{n-1})\longrightarrow
 H_n(X^n)\xrightarrow{\ j_n\ }
 C_n^{\mathrm{CW}}
 \xrightarrow{\ \delta_n\ }
 H_{n-1}(X^{n-1}).
\]
Odwzorowanie $j_{n-1}:H_{n-1}(X^{n-1})
\to C_{n-1}^{\mathrm{CW}}$
jest iniektywne: jego jądro
jest obrazem
$H_{n-1}(X^{n-2})=0$.
Dla $n=1$ używamy
$X^{-1}=\varnothing$, co daje
ten sam wniosek.
Zatem $\ker\partial_n^{\mathrm{CW}}
=\ker\delta_n=\operatorname{im}j_n$
i $j_n$ identyfikuje cykle
komórkowe z $H_n(X^n)$.

W następnym długim ciągu,
dla $(X^{n+1},X^n)$,
\[
 C_{n+1}^{\mathrm{CW}}
 \xrightarrow{\ \delta_{n+1}\ }
 H_n(X^n)\longrightarrow
 H_n(X^{n+1})\longrightarrow
 H_n(X^{n+1},X^n)=0.
\]
Jądro przejścia do $H_n(X^{n+1})$
jest więc $\operatorname{im}\delta_{n+1}$.
Po włożeniu w $C_n^{\mathrm{CW}}$
przez $j_n$ jest to dokładnie
$\operatorname{im}\partial_{n+1}^{\mathrm{CW}}$.
Otrzymujemy
$H_n^{\mathrm{CW}}(X)\cong
H_n(X^{n+1})\cong H_n(X)$.
Dla $n=0$ ten sam argument działa
z $H_0(X^0)=C_0^{\mathrm{CW}}$:
$H_0(X^1)$ jest ilorazem przez
obrazy brzegów $1$-komórek,
a wyższe komórki już nie
zmieniają $H_0$.
\end{proof}

\section{Odwzorowania dołączające i konkretne obliczenia}
\label{sec:obliczenia-homologii}

\begin{definicja}[Stopień odwzorowania sfer]
\label{def:stopien-sfer}
Dla $m\geq1$ każde ciągłe odwzorowanie
$h:S^m\to S^m$ indukuje
endomorfizm $H_m(S^m)\cong\Z$.
Istnieje więc dokładnie jedna
liczba całkowita $\deg h$ taka,
że $h_*([S^m])=(\deg h)[S^m]$.
Nazywamy ją \emph{stopniem} $h$.
\end{definicja}

Z definicji i funktorialności
$\deg(\id)=1$ oraz
$\deg(h\circ g)=\deg h\deg g$.
Homotopijne odwzorowania mają ten sam
stopień. Na $S^1=\{z\in\C:|z|=1\}$
odwzorowanie $z\mapsto z^r$ ma stopień $r\in\Z$.
Uzasadnijmy przejście od liczby obiegów do homologii.
Pętla $\gamma(t)=e^{2\pi it}$ reprezentuje generator $H_1(S^1)$:
po podziale na łuki jej obraz przez izomorfizm Mayera--Vietorisa
z Przykładu~\ref{prz:mv-okrag} jest różnicą punktów dwóch składowych przecięcia.
Dla dowolnych dwóch pętli $\alpha,\beta$ opartych w tym samym punkcie
mamy $[\alpha*\beta]=[\alpha]+[\beta]$ w $H_1$.
Istotnie, złożenie $\alpha*\beta$ z afinicznym odwzorowaniem
$\Delta^2\to[0,1]$ o wartościach $0,1/2,1$ w kolejnych wierzchołkach
jest singularnym $2$-sympleksem o brzegu
$\beta-(\alpha*\beta)+\alpha$.
Ponadto homotopia pętli względem końców daje homologiczne cykle
przez wzór pryzmatyczny; pętla stała jest brzegiem stałego $2$-sympleksu.
Stąd $[\bar\alpha]=-[\alpha]$, bo $\alpha*\bar\alpha$ jest
homotopijna do stałej. Pętla $t\mapsto e^{2\pi irt}$ jest,
z dokładnością do zmiany parametryzacji, złożeniem $r$ obiegów dla $r>0$,
$|r|$ obiegów przeciwnych dla $r<0$
i jest stała dla $r=0$. Jej klasa to zatem $r[\gamma]$.

\begin{twierdzenie}[Współczynniki brzegu komórkowego]
\label{thm:formula-komorkowa}
Niech $\varphi_\alpha:S^{n-1}\to X^{n-1}$
będzie odwzorowaniem dołączającym
zorientowanej $n$-komórki
$e_\alpha^n$, gdzie $n\geq2$.
Dla każdej zorientowanej
$(n-1)$-komórki $e_\beta^{n-1}$
zwińmy w $X^{n-1}$ do punktu
dopełnienie jej otwartego wnętrza.
Otrzymamy odwzorowanie
\[
 q_\beta:X^{n-1}\longrightarrow
 X^{n-1}/(X^{n-1}\setminus e_\beta^{n-1})
 \cong S^{n-1}.
\]
Sferę docelową orientujemy zgodnie z komórką $e_\beta^{n-1}$,
a dziedzinę $\varphi_\alpha$ jako brzeg zorientowanej komórki $e_\alpha^n$.
Wówczas
\begin{equation}\label{eq:formula-komorkowa}
 \partial_n^{\mathrm{CW}}(e_\alpha^n)
 =\sum_{\beta}
 \deg(q_\beta\circ\varphi_\alpha)
 e_\beta^{n-1}.
\end{equation}
W stopniu $1$ brzeg zorientowanej
komórki od początku $v_-$ do końca
$v_+$ wynosi $\partial_1e=v_+-v_-$.
\end{twierdzenie}
\begin{proof}
Generator $e_\alpha^n$ jest
obrazem klasy orientacji
$(D^n,S^{n-1})$ w parze
$(X^n,X^{n-1})$. Pierwsza
strzałka~\eqref{eq:cellular-boundary}
zwraca klasę obrazu
zorientowanego brzegu $S^{n-1}$
przez $\varphi_\alpha$:
tak działa odwzorowanie łączące,
bo na łańcuchach bierze
dosłowny brzeg podniesionego
cyklu względnego.
Druga strzałka przenosi tę
klasę do
$H_{n-1}(X^{n-1},X^{n-2})$.
Według
Twierdzenia~\ref{thm:warstwa-komorek}
jest to suma kopii $\Z$,
po jednej dla każdego $\beta$.
Rzutowanie na składnik $\beta$
jest, poprzez iloraz z
Lematu~\ref{lem:iloraz-cw-homologia},
tym samym co $q_{\beta*}$.
Zatem współczynnik przy
$e_\beta^{n-1}$ to liczba,
przez którą
$(q_\beta\circ\varphi_\alpha)_*$
mnoży generator homologii
$S^{n-1}$: z definicji
jest to jego stopień.
Dla $n=1$ używamy
bezpośrednio
$\partial[0,1]=[1]-[0]$;
oba końce mogą w $X^0$
oznaczać ten sam wierzchołek.
\end{proof}

Warto przeczytać wzór geometrycznie:
liczymy, ile razy, z jakimi znakami,
brzeg wyższej komórki przechodzi
przez daną komórkę niższą. Komórka
z niezerowym współczynnikiem brzegu
może zabić klasę homologii.

\begin{figure}[!htbp]
\centering
\begin{tikzpicture}[>=Stealth,line cap=round,line join=round]
 \node[draw=manifoldblue!65!black,fill=manifoldblue!9,
  rounded corners=4pt,minimum width=2.15cm,
  minimum height=.8cm] (c2) at (1.2,0) {$C_2=\Z$};
 \node[draw=manifoldgreen!65!black,fill=manifoldgreen!9,
  rounded corners=4pt,minimum width=2.35cm,
  minimum height=.8cm] (c1) at (4.5,0) {$C_1=\Z^2$};
 \node[draw=manifoldgold!75!black,fill=manifoldgold!10,
  rounded corners=4pt,minimum width=2.15cm,
  minimum height=.8cm] (c0) at (7.8,0) {$C_0=\Z$};
 \draw[->,thick,manifoldred] (c2)--(c1)
  node[midway,above=5pt] {$\partial_2=0$};
 \draw[->,thick,manifoldred] (c1)--(c0)
  node[midway,above=5pt] {$\partial_1=0$};
 \node[align=center,font=\small,manifoldblue!75!black]
  at (1.2,-.9) {jedna\\$2$-komórka};
 \node[align=center,font=\small,manifoldgreen!75!black]
  at (4.5,-.9) {dwie\\pętle};
 \node[font=\small,manifoldgold!75!black]
  at (7.8,-.9) {punkt};
\end{tikzpicture}
\caption{Kompleks komórkowy torusa. Ma trzy
niezerowe grupy łańcuchów, a oba odwzorowania
brzegu znikają: pętle kończą się w tym
samym punkcie, a słowo dołączające
komórkę $2$-wymiarową jest komutatorem.}
\label{fig:kompleks-torusa}
\end{figure}

\begin{przyklad}[Sfera z dwóch komórek]
\label{prz:homologia-cw-sfera}
W modelu z
Przykładu~\ref{prz:cw-s2}
jest jedna komórka $e^0=v$
i jedna $e^2$, bez $1$-komórek.
Kompleks komórkowy to
\[
 0\longrightarrow\Z e^2
 \xrightarrow{\;0\;}0
 \xrightarrow{\;0\;}\Z v
 \longrightarrow0.
\]
Stąd $H_2(S^2)=\Z$, $H_1(S^2)=0$,
$H_0(S^2)=\Z$. Ogólniej
$S^n=D^n/\partial D^n$ ma
jedną komórkę w stopniu $0$
i jedną w stopniu $n$
($n\geq2$). Dla $S^1$
jedna $1$-komórka jest pętlą,
więc $\partial_1=0$.
W ten sposób odzyskujemy
Twierdzenie~\ref{thm:homologia-sfer}.
\end{przyklad}

\begin{przyklad}[Torus]
\label{prz:homologia-cw-torus}
Model CW z
Przykładu~\ref{prz:pi1-torus-rp2}
ma komórki $v$,
$a,b$ oraz $e^2$.
Krawędzie są pętlami, więc
$\partial_1a=\partial_1b=0$.
Brzeg $e^2$ jest dołączony
według słowa $aba^{-1}b^{-1}$.
Współczynnik przy $a$
to $1-1=0$, a przy $b$
to $1-1=0$; jest to
szczególny przypadek
Wzoru~\eqref{eq:formula-komorkowa}.
Zatem $\partial_2=0$ i
\[
 H_0(T^2)=\Z,\qquad
 H_1(T^2)=\Z^2,\qquad
 H_2(T^2)=\Z,\qquad
 H_k(T^2)=0\quad(k>2).
\]
Rysunek~\ref{fig:kompleks-torusa}
pokazuje cały rachunek.
Grupa $\Z^2$ rozpoznaje
dwa niezależne kierunki pętli.
\end{przyklad}

\begin{przyklad}[Płaszczyzna rzutowa i torsja]
\label{prz:homologia-rp2}
Przestrzeń $\mathbb{RP}^2$
ma po jednej komórce
$e^0,e^1,e^2$.
Jedyna $1$-komórka jest pętlą,
więc $\partial_1=0$.
Odwzorowanie dołączające $e^2$
okrąża $\mathbb{RP}^1\cong S^1$
dwa razy, co było widoczne
w Przykładzie~\ref{prz:pi1-torus-rp2}.
Ma stopień $2$ po ustaleniu
orientacji, zatem
\[
 0\longrightarrow\Z
 \xrightarrow{\;\times2\;}\Z
 \xrightarrow{\;0\;}\Z
 \longrightarrow0.
\]
Jądro mnożenia przez $2$
jest zerowe, a iloraz
$\Z/2\Z$ ma element
niezerowy, który po
dwukrotnym dodaniu znika.
Otrzymujemy
\[
 H_0(\mathbb{RP}^2)=\Z,\quad
 H_1(\mathbb{RP}^2)=\Z/2\Z,\quad
 H_2(\mathbb{RP}^2)=0.
\]
To konkretny przykład \emph{torsji}
zdefiniowanej w
Sekcji~\ref{sec:wspolczynniki-euler}.
\end{przyklad}

\begin{uwaga}[Homologia ze współczynnikami]
\label{rem:homologia-wspolczynniki-przestrzeni}
Uzasadnijmy zmianę współczynników także w dodatniej charakterystyce.
Dla przemiennego pierścienia $R$ z jedynką definiujemy $C_n(X;R)$
jako wolny $R$-moduł na singularnych sympleksach, z tym samym wzorem brzegu.
Równoważnie jest to $C_n(X;\Z)\otimes_{\Z}R$.
Homotopie pryzmatyczne i podział barycentryczny mają całkowite współczynniki,
więc wszystkie ich tożsamości pozostają prawdziwe nad $R$.
Ciągi łańcuchów par oraz małych łańcuchów Mayera--Vietorisa są
rozszczepione w każdym stopniu na wybranych bazach, więc są dokładne także nad $R$.
Rachunek punktu, sfer i kolejnych szkieletów przebiega wtedy tak samo:
$C_n^{\mathrm{CW}}(X;R)$ ma po jednej kopii $R$ na komórkę.
Naturalność odwzorowań łączących względem $\Z\to R$ pokazuje,
że całkowity współczynnik brzegu $d$ staje się $d\cdot1_R$.
To uzasadnia obliczenia komórkowe nad $R$; nie zakładamy,
że tensorowanie dowolnego ciągu dokładnego zachowuje dokładność.
\end{uwaga}

\begin{przyklad}[Te same przestrzenie nad różnymi ciałami]
\label{prz:zmiana-wspolczynnikow-rp2-torus}
Wstawmy teraz do powyższych kompleksów komórkowych
w miejsce $\Z$ dowolne
$\mathbb F\in\{\Q,\R,\C\}$.
Dla płaszczyzny rzutowej dostajemy
\[
 0\longrightarrow\mathbb F
 \xrightarrow{\;\times2\;}\mathbb F
 \xrightarrow{\;0\;}\mathbb F
 \longrightarrow0.
\]
Ponieważ $2$ ma w $\mathbb F$ odwrotność
$1/2$, środkowe odwzorowanie jest
izomorfizmem. Zatem
\[
 H_0(\mathbb{RP}^2;\mathbb F)=\mathbb F,
 \qquad H_k(\mathbb{RP}^2;\mathbb F)=0\quad(k>0).
\]
Klasa $\Z/2$ zniknęła. Dla torusa oba
operatory brzegu nadal są zerowe, więc
\[
 H_0(T^2;\mathbb F)=\mathbb F,\qquad
 H_1(T^2;\mathbb F)=\mathbb F^2,\qquad
 H_2(T^2;\mathbb F)=\mathbb F.
\]
Wymiary tych przestrzeni, odpowiednio
$1,2,1$, to liczby Betti'ego torusa.
Tak działa ogólny
Wzór~\eqref{eq:homologia-char-zero}.
Dla porównania nad $\Z/2$ strzałka
$\times2$ w kompleksie $\mathbb{RP}^2$
staje się zerowa. Wtedy oprócz $H_0$
otrzymujemy również
$H_1(\mathbb{RP}^2;\Z/2)\cong\Z/2$
oraz $H_2(\mathbb{RP}^2;\Z/2)\cong\Z/2$.
\end{przyklad}

\begin{przyklad}[Bukiet okręgów i charakterystyka Eulera]
\label{prz:homologia-bukiet}
Bukiet $r$ okręgów ma jeden
wierzchołek i $r$ pętli,
bez komórek wyższych:
$C_1^{\mathrm{CW}}=\Z^r$,
$C_0^{\mathrm{CW}}=\Z$ i
$\partial_1=0$.
Stąd $H_1=\Z^r$, $H_0=\Z$,
a wyższe homologie znikają.
Dla skończonego CW z
$c_n$ komórkami w wymiarze $n$
równanie Eulera z
Twierdzenia~\ref{thm:euler-homologia}
przyjmuje postać
\[
 \sum_n(-1)^n c_n
 =\sum_n(-1)^n
   \operatorname{rank}_{\Z} H_n(X;\Z).
\]
W szczególności dla torusa
$1-2+1=0$, a dla
$\mathbb{RP}^2$ mamy
$1-1+1=1$; torsja
$\Z/2\Z$ ma rangę zero.
\end{przyklad}

\section{Lokalna homologia i dalsza droga}
\label{sec:lokalna-homologia}

\begin{stwierdzenie}[Lokalna homologia rozmaitości]
\label{prop:lokalna-homologia}
Jeśli $M$ jest topologiczną
rozmaitością wymiaru $m$
bez brzegu i $p\in M$, to
\[
 H_k(M,M\setminus\{p\};\Z)
 \cong
 \begin{cases}\Z,&k=m,\\0,&k\ne m.\end{cases}
\]
Wybór generatora grupy w stopniu $m$
jest \emph{lokalną orientacją}
w punkcie $p$.
\end{stwierdzenie}
\begin{proof}
Wybierzmy otwarte otoczenie $U\subseteq M$ i mapę
$\psi:U\to B^m$ na otwartą kulę jednostkową,
z $\psi(p)=0$. Zbiór
$Z=M\setminus U$
jest domknięty i leży
we wnętrzu $M\setminus\{p\}$,
więc wycinanie
(Twierdzenie~\ref{thm:wycinanie})
daje
\[
 H_k(M,M\setminus\{p\})
 \cong H_k(U,U\setminus\{p\})
\cong H_k(B^m,B^m\setminus\{0\}).
\]
Dla $m\geq1$ kula jest
ściągalna, a kula pozbawiona
środka silnie deformacyjnie retrahuje się na sferę
o promieniu $1/2$, homeomorficzną z $S^{m-1}$:
używamy $ru\mapsto((1-t)r+t/2)u$, co pozostaje w $0<r<1$.
Długi ciąg pary i
Twierdzenie~\ref{thm:homologia-sfer}
dają wynik. Dla $m=0$
punkt jest otwartą składową
mapy lokalnej, a grupa
$H_0(\{p\},\varnothing)=\Z$.
\end{proof}

Istnienie lokalnego generatora
nie oznacza jeszcze, że można
wybrać generatory \emph{spójnie}
na całej rozmaitości. To będzie
treść pojęcia orientowalności.
Na obecnym etapie mamy
już narzędzia do dalszej nauki:
kompleksy łańcuchowe i długie ciągi,
homotopijną niezmienniczość,
wycinanie, Mayera--Vietorisa
oraz model komórkowy, w którym
odwzorowania dołączające kontrolują
operator brzegu. Pozwalają one
później opisać, jak dołączenie
uchwytu zmienia homologię,
bez wprowadzania tutaj
teorii Morse'a i chirurgii.

\section*{Dalsza lektura}
\addcontentsline{toc}{section}{Dalsza lektura}
Pełne rozwinięcie homologii symplicjalnej,
singularnej i komórkowej, wraz ze
szczegółami technicznymi dotyczącymi
podziału barycentrycznego, znajduje się
w Allen Hatcher,
\emph{Algebraic Topology}, rozdział 2,
\url{https://pi.math.cornell.edu/~hatcher/AT/AT.pdf}.
\clearpage

\chapter{Wyższe grupy homotopii i związek z homologią}
\label{ch:wyzsza-homotopia}

Grupa podstawowa z Rozdziału~\ref{ch:topologia-algebraiczna-podstawy}
rejestruje pętle. Wyższe grupy homotopii mierzą, czy można
ściągnąć do punktu odwzorowania całej sfery $S^n$. Homologia z
Rozdziału~\ref{ch:homologia-przestrzeni} bierze z kolei
zorientowane cykle i zapomina część informacji o sposobie ich
odwzorowania. Twierdzenie Hurewicza wskazuje dokładnie sytuację,
w której te dwa opisy mają ten sam pierwszy niezerowy wynik.
Twierdzenie Whiteheada mówi, kiedy wszystkie grupy homotopii
rozstrzygają typ homotopijny kompleksu CW.

W tym rozdziale pracujemy z przestrzeniami z wybranym punktem bazowym,
przede wszystkim z kompleksami CW. W argumentach komórkowych
wybieramy punkt bazowy będący wierzchołkiem. Gdy piszemy
$\pi_n(X)$ bez punktu bazowego, zakładamy spójność łukową i
wybieramy punkt bazowy; jego zmiana wzdłuż drogi daje izomorfizm, ale
sam wybór drogi może mieć znaczenie.

\section{Od pętli do odwzorowań sfer}
\label{sec:grupy-pi-n}

\begin{definicja}[Grupa $\pi_n$ i punkt bazowy]
\label{def:pi-n}
Niech $I=[0,1]$, $I^n=I\times\cdots\times I$, a
$\partial I^n$ niech oznacza sumę wszystkich jego ścian.
Dla przestrzeni z wybranym punktem $x_0$ kładziemy
\[
 \pi_n(X,x_0)
 =\big[(I^n,\partial I^n),(X,x_0)\big]_*
 =\big[(S^n,*),(X,x_0)\big]_*
 \qquad(n\geq1).
\]
Nawias oznacza klasy homotopii, w których
\emph{przez cały czas} brzeg kostki lub wyróżniony
punkt sfery jest wysyłany do $x_0$.
Utożsamienie pochodzi z $I^n/\partial I^n\cong S^n$.
Dla $n=1$ otrzymujemy $\pi_1(X,x_0)$ z
Definicji~\ref{def:droga-petla}. Symbol $\pi_0(X)$
oznacza jedynie zbiór składowych łukowych,
zwykle bez działania grupowego.
\end{definicja}

Obraz odwzorowania $S^2\to X$ można sobie wyobrażać jako
elastyczną powłokę rozpiętą w przestrzeni; jej
ściągalność do punktu wymaga ciągłej deformacji
\emph{całego odwzorowania}, nie tylko zniknięcia klasy
homologii reprezentowanej przez to odwzorowanie. Sfera $S^n$ jest brzegiem dysku $D^{n+1}$,
a także jest homeomorficzna z brzegiem kostki $I^{n+1}$.

\begin{twierdzenie}[Działanie i przemienność]
\label{thm:pi-n-abelowe}
Dla $n\geq1$ klasy w $\pi_n(X,x_0)$ tworzą grupę.
Dla $n\geq2$ grupa ta jest abelowa.
\end{twierdzenie}
\begin{proof}
Jeśli $f,g:(I^n,\partial I^n)\to(X,x_0)$,
przedstawmy ich sumę sklejeniem w pierwszej
współrzędnej:
\[
 (f+_1g)(s_1,\ldots,s_n)=
 \begin{cases}
 f(2s_1,s_2,\ldots,s_n),&0\leq s_1\leq\tfrac12,\\
 g(2s_1-1,s_2,\ldots,s_n),&\tfrac12\leq s_1\leq1.
 \end{cases}
\]
Na płaszczyźnie sklejenia oba wzory dają
$x_0$; tak samo jest na zewnętrznym brzegu.
Homotopie $f$ lub $g$ zachowujące brzeg
sklejają się tym samym wzorem, więc działanie
nie zależy od reprezentantów. Łączność
wynika z liniowej zmiany podziału pierwszej
współrzędnej na trzy odcinki, dokładnie jak
w dowodzie Twierdzenia~\ref{thm:grupa-podstawowa}.
Stałe odwzorowanie jest elementem neutralnym:
dodatkową połowę, na której odwzorowanie jest stałe,
można stopniowo zwęzić. Odbicie
$s_1\mapsto1-s_1$ reprezentuje odwrotność,
bo dwie przeciwne połówki ściągamy
równocześnie do punktu.

Gdy $n\geq2$, możemy niezależnie zdefiniować
drugie działanie $+_2$, sklejając w kierunku
$s_2$. Ma ono ten sam element neutralny
i tę samą konstrukcję homotopii. Podział
kwadratu $(s_1,s_2)$ na cztery ćwiartki
daje \emph{równość odwzorowań}, a nie tylko klas:
\[
 (a+_1 b)+_2(c+_1d)
   =(a+_2c)+_1(b+_2d).
\]
Teraz przechodzimy do klas homotopii i oznaczamy klasę stałą przez $0$.
Korzystając z neutralności $0$, otrzymujemy
\[
 a+_1b=(a+_2 0)+_1(0+_2 b)
      =(a+_1 0)+_2(0+_1 b)=a+_2b.
\]
Z kolei
\[
 a+_1b=(0+_2a)+_1(b+_2 0)
      =(0+_1b)+_2(a+_1 0)=b+_2a=b+_1a.
\]
Te ostatnie równości dotyczą klas; samo usuwanie stałych połówek
zmienia parametryzację odwzorowania.
To argument Eckmanna--Hiltona: dwa
zgodne sposoby sklejania wymuszają
przemienność. Użycie drugiej
współrzędnej jest niemożliwe dla
$n=1$, dlatego $\pi_1$ może być
nieabelowa.
\end{proof}

\begin{figure}[!htbp]
\centering
\begin{tikzpicture}[>=Stealth,line cap=round,line join=round]
 \fill[manifoldblue!8] (0,0) rectangle (4,4);
 \draw[manifoldblue!75!black,line width=1.4pt] (0,0) rectangle (4,4);
 \draw[manifoldblue!55] (2,0)--(2,4) (0,2)--(4,2);
 \fill[manifoldred!20] (.28,2.28) rectangle (1.72,3.72);
 \fill[manifoldgreen!22] (2.28,.28) rectangle (3.72,1.72);
 \node[font=\large,manifoldred!80!black] at (1,3) {$f$};
 \node[font=\large,manifoldgreen!70!black] at (3,1) {$g$};
 \node at (3,3) {$*$}; \node at (1,1) {$*$};
 \draw[->,manifoldgold!85!black,thick] (0,-.35)--(4,-.35)
  node[midway,below=4pt] {sklejanie w kierunku $s_1$};
 \draw[->,manifoldgold!85!black,thick] (-.35,0)--(-.35,4)
  ;
 \node[rotate=90,manifoldgold!85!black] at (-.8,2)
  {sklejanie w kierunku $s_2$};
 \node[font=\small,align=left,anchor=west,text width=4.2cm]
  at (4.6,2)
  {Puste pola są odwzorowaniem stałym. Oba kierunki
   sklejania opisują to samo odwzorowanie czterech
   ćwiartek; zamiana kolejności daje
   $[f]+[g]=[g]+[f]$.};
\end{tikzpicture}
\caption{Dwa niezależne kierunki sklejania w $\pi_2$.
W $\pi_n$ dla $n\geq2$ pozostałe współrzędne
traktujemy jako parametry.}
\label{fig:pi2-sklejanie}
\end{figure}

\begin{stwierdzenie}[Odwzorowania, iloczyny i homotopie]
\label{prop:pi-n-funktor}
Ciągłe odwzorowanie oparte $F:(X,x_0)\to(Y,y_0)$ wyznacza
homomorfizm $F_*:\pi_n(X,x_0)\to\pi_n(Y,y_0)$
przez $[f]\mapsto[F\circ f]$.
Odwzorowania homotopijne \emph{przez odwzorowania oparte}
indukują ten sam homomorfizm. Dla
$n\geq1$ i przestrzeni opartych $X,Y$:
\[
 \pi_n(X\times Y,(x_0,y_0))
 \cong \pi_n(X,x_0)\times\pi_n(Y,y_0).
\]
\end{stwierdzenie}
\begin{proof}
Złożenie $F\circ f$ nadal wysyła
$\partial I^n$ do $y_0$, a homotopia
$f_t$ daje homotopię $F\circ f_t$.
Ponieważ $F$ działa po sklejaniu
obu połówek dokładnie tak samo
jak przed nim, $F_*([f]+[g])
=F_*[f]+F_*[g]$. Homotopia
oparta $F_t$ daje homotopię
$F_t\circ f$ z nieruchomym
obrazem brzegu. Wreszcie odwzorowanie
$I^n\to X\times Y$ to dokładnie
para odwzorowań do $X$ i $Y$; to
samo dotyczy homotopii.
Projekcje i składanie pary
stanowią wzajemne odwrotności
podanego izomorfizmu.
\end{proof}

\begin{uwaga}[Zmiana punktu bazowego]
Droga $\gamma$ od $x_0$ do $x_1$
pozwala umieścić odwzorowanie oparte w $x_1$
w mniejszej kostce i w otaczającym
ją pasie doprowadzić wartości
wzdłuż $\gamma$ do $x_0$.
Otrzymujemy izomorfizm grup
$\pi_n(X,x_1)\to\pi_n(X,x_0)$
z odwrotnością od drogi
przebytej w przeciwną stronę.
Dwie różne drogi mogą dać
różne izomorfizmy; dla zamkniętej
drogi dostajemy \emph{działanie}
$\pi_1(X,x_0)$ na $\pi_n(X,x_0)$.
Konstrukcja zależy tylko od klasy homotopii drogi względem końców:
homotopię drogi stosujemy na tym samym pasie. Złożenie dwóch pasów
odpowiada złożeniu dróg, a droga i jej odwrotność dają pas ściągalny.
To uzasadnia zarówno izomorfizm, jak i działanie grupy.
W przypadku przestrzeni jednospójnej wybór drogi
nie powoduje tej niejednoznaczności.
Jeśli $F_t:X\to Y$ jest homotopią bez warunku oparcia i
$\gamma(t)=F_t(x_0)$, to po zmianie punktu bazowego wzdłuż $\gamma$
homomorfizm $(F_1)_*$ staje się $(F_0)_*$: pas odwzorowania
$F_t\circ f$ jest właśnie pasem tej zmiany bazy.
W szczególności zwykła równoważność homotopijna indukuje
izomorfizmy wszystkich grup homotopii, z odpowiednio dobranymi bazami.
\end{uwaga}

\section{Grupy względne i długi ciąg homotopii}
\label{sec:pi-wzgledne}

Tak jak $H_n(X,A)$ mierzy cykle,
których brzeg może trafić do $A$,
tak względna homotopia mierzy
dyski w $X$ z brzegiem w $A$.
Brzeg dysku może się przy tym
przesuwać w $A$ podczas homotopii.

\begin{definicja}[Względne grupy homotopii]
\label{def:pi-wzgledne}
Niech $x_0\in A\subseteq X$.
W kostce $I^n$ wyróżnijmy
dolną ścianę $B=I^{n-1}\times\{0\}$,
a przez $J$ oznaczmy sumę
pozostałych ścian kostki $I^n$.
Klasa względna jest klasą odwzorowań
\[
 F:(I^n,\partial I^n,J)
       \longrightarrow(X,A,x_0);
\]
zatem $F(\partial I^n)\subseteq A$
i $F(J)=x_0$. Homotopie spełniają
te same warunki. Ich zbiór oznaczamy
$\pi_n(X,A,x_0)$.
W równoważnym modelu są to
oparte odwzorowania trójek
$(D^n,S^{n-1},*)\to(X,A,x_0)$.
Dla $n\geq2$ składanie w pierwszej
współrzędnej daje grupę; dla
$n\geq3$ jest ona abelowa.
Dla $n=1$ pozostaje \emph{zbiór
z wyróżnionym elementem},
ponieważ nie ma wolnej współrzędnej
do składania.
\end{definicja}

W przeciwieństwie do $\pi_n(X)$,
względna $\pi_2(X,A)$ może być
nieprzemienna. Drugi kierunek
sklejania zajmuje wówczas rolę
wyróżnionej ściany, a argument
z Rysunku~\ref{fig:pi2-sklejanie}
nie ma zastosowania.

\begin{figure}[!htbp]
\centering
\begin{tikzpicture}[>=Stealth,line cap=round,line join=round]
 \filldraw[rounded corners=8pt,fill=manifoldblue!7,draw=manifoldblue!60!black]
  (0,0) rectangle (8,3.4);
 \fill[manifoldblue!18] (2.5,1.65) ellipse (1.65 and 1.05);
 \draw[manifoldgreen!28,line width=12pt] (2.5,1.65) ellipse (1.65 and 1.05);
 \draw[manifoldred,line width=1.4pt] (2.5,1.65) ellipse (1.65 and 1.05);
 \fill[manifoldgold!80!black] (.85,1.65) circle(2.6pt);
 \node[left,font=\small] at (.78,1.65) {$x_0$};
 \node[manifoldblue!75!black] at (2.5,1.65) {$F(\operatorname{int}D^2)$};
 \node[anchor=west,manifoldgreen!60!black] at (5.1,2.5) {$A$ (zielony pas)};
 \draw[->,manifoldgreen!60!black] (5,2.5)--(3.85,2.32);
 \node[anchor=west,manifoldred!80!black] at (5.1,1.05) {$F(S^1)\subset A$};
 \draw[->,manifoldred!80!black] (5,1.05)--(4.05,1.3);
 \node[manifoldblue!75!black] at (7.65,3.05) {$X$};
\end{tikzpicture}
\caption{Przykład odwzorowania względnego $(D^2,S^1,*)\to(X,A,x_0)$.
Czerwony obraz brzegu leży w zielonym pasie $A$, podczas gdy
obraz wnętrza dysku może wyjść poza $A$. W ogólnej definicji
nie wymaga się różnowartościowości $F$.}
\label{fig:pi-wzgledny-dysk}
\end{figure}

\begin{twierdzenie}[Długi ciąg homotopii pary]
\label{thm:les-pi}
Dla $x_0\in A\subseteq X$ otrzymujemy
naturalny ciąg
\begin{equation}\label{eq:les-pi}
 \cdots\to\pi_n(A,x_0)\xrightarrow{i_*}
 \pi_n(X,x_0)\xrightarrow{j_*}
 \pi_n(X,A,x_0)\xrightarrow{\partial}
 \pi_{n-1}(A,x_0)\xrightarrow{i_*}
 \pi_{n-1}(X,x_0)\to\cdots .
\end{equation}
Jest on dokładny: obraz każdej
strzałki to jądro następnej.
Odwzorowanie $\partial$ ogranicza
reprezentujący dysk do jego
wyróżnionej ściany $B$.
Na końcu ciągu występują
zbiory $\pi_1(X,A,x_0)$ i
$\pi_0(A),\pi_0(X)$; dokładność
oznacza tam przeciwobraz
wyróżnionego elementu.
\end{twierdzenie}
\begin{proof}
W modelu kostkowym $i_*$ jest
włączeniem $A\to X$, a $j_*$
traktuje odwzorowanie z całym brzegiem
w $x_0$ jako odwzorowanie względne.
Wzór $\partial[F]=[F|_B]$
jest poprawny, ponieważ
$\partial B\subset J$ i
$F(J)=x_0$. Homotopia
względna ogranicza się
do homotopii opartych
odwzorowań na $B$, więc
$\partial$ nie zależy
od reprezentanta.

Sprawdźmy trzy kolejne miejsca.
\emph{W $\pi_n(X)$.} Jeśli $f$ ma cały obraz w $A$, homotopia
$f_t(s',r)=f(s',(1-t)r+t)$ jest względną homotopią do stałej;
tutaj $s'\in I^{n-1}$, a $r$ jest ostatnią współrzędną.
Odwrotnie, niech $H(s',r,t)$ będzie względną homotopią
od $f$ do stałej. Ściana $g(s',t)=H(s',0,t)$ leży w $A$
i ma cały brzeg w $x_0$, więc reprezentuje element $\pi_n(A)$.
W kwadracie $(r,t)$ porównujemy drogi od $(0,0)$ do $(1,1)$:
dolny, a potem prawy bok oraz lewy, a potem górny bok.
Są homotopijne względem końców przez liniową interpolację w kwadracie.
Po złożeniu z $H$ dają odpowiednio $f$ i $g$ ze stałym dodatkowym
odcinkiem; na brzegu współrzędnych $s'$ wszystko pozostaje w $x_0$.
Stąd $[f]=i_*[g]$.

\emph{W $\pi_n(X,A)$.} Dysk
pochodzący z $\pi_n(X)$
ma stały brzeg, więc jego
$\partial$ jest zerem.
Jeśli natomiast brzeg
$F|_B:S^{n-1}\to A$ jest
zerowy w $\pi_{n-1}(A)$,
wybieramy jego wypełnienie
przez $D^n\to A$. Sklejenie
tego wypełnienia z $F$
wzdłuż wspólnego brzegu
daje oparte odwzorowanie
$S^n=D^n\cup_{S^{n-1}}D^n\to X$.
Po zapomnieniu drugiego
dysku jego obraz względny
jest klasą $[F]$.

\emph{W $\pi_{n-1}(A)$.} Brzeg
dysku względnego wypełnia
się w $X$, więc ginie
po włączeniu do $X$.
Jeśli oparte odwzorowanie
$S^{n-1}\to A$ jest
zerowe w $\pi_{n-1}(X)$,
jego homotopia do odwzorowania
stałego wyznacza wypełnienie
$D^n\to X$ i tym samym
klasę względną, której
brzeg jest daną klasą.
Te trzy rachunki powtarzają się w sąsiednich stopniach.
Na końcu ciągu element $\pi_1(X,A,x_0)$ jest drogą
od pewnego $a\in A$ do $x_0$, a $\partial$ przypisuje mu
składową $a$ w $A$. Trafia ona do wyróżnionej składowej $X$.
Odwrotnie, punkt tej składowej można połączyć drogą z $x_0$.
Jeśli $a$ leży w składowej $x_0$ już w $A$, dopinamy drogę
z $x_0$ do $a$ w $A$ i otrzymujemy pętlę w $X$ reprezentującą
daną klasę względną. To dowodzi dokładności także dla zbiorów wskazanych.
Wszystkie konstrukcje zachowują się przy składaniu z mapą par,
co daje naturalność ciągu.
\end{proof}

\begin{definicja}[Spójność homotopijna]
\label{def:spojnosc-homotopijna}
Przestrzeń $X$ jest
\emph{$r$-spójna} ($r\geq0$),
jeśli jest łukowo spójna
i $\pi_i(X)=0$ dla $1\leq i\leq r$.
W szczególności $1$-spójna
oznacza jednospójną.
Para $(X,A)$ jest
\emph{$r$-spójna}, jeśli
każda składowa łukowa $X$
spotyka $A$ oraz
$\pi_i(X,A)=0$ dla
$1\leq i\leq r$ przy każdym wyborze punktu bazowego w $A$
(w stopniu $1$ jest to
trywialność zbioru
z wyróżnioną klasą).
\end{definicja}

Dokładność~\eqref{eq:les-pi}
pokazuje, że dla $r$-spójnej
pary włączenie $A\hookrightarrow X$
jest izomorfizmem na
$\pi_i$ dla $i<r$ oraz
surjekcją na $\pi_r$.
Wymiar komórek CW pozwoli
za chwilę uzyskać takie
zanikanie bez liczenia odwzorowań
pojedynczo.

\section{Homotopie na komórkach CW}
\label{sec:pi-n-cw}

W rozdziale o homologii komórki dawały
generatory grup łańcuchów. W homotopii
ich wymiary pozwalają \emph{usuwać
przeszkody do deformacji}. Żeby
deformować jedną komórkę bez psucia
tego, co zrobiliśmy na jej brzegu,
potrzebujemy przedłużania homotopii.

\begin{lemat}[Zwarte obrazy i homotopie w CW]
\label{lem:zwartosc-cw-pi}
Każdy zwarty podzbiór kompleksu CW leży w skończonym podkompleksie.
Ponadto odwzorowanie $H:X\times I\to Y$ jest ciągłe, jeśli jego złożenia
z każdym $\chi_\alpha\times\id_I:D^{d_\alpha}\times I\to X\times I$
są ciągłe, gdzie $\chi_\alpha$ są mapami charakterystycznymi komórek.
\end{lemat}
\begin{proof}
Gdyby zwarty zbiór $C$ spotykał nieskończenie wiele otwartych komórek,
wybralibyśmy z każdej z nich jeden punkt $c_\alpha\in C$.
Każda domknięta komórka spotyka tylko skończenie wiele spośród tych punktów.
Każdy podzbiór $E$ zbioru wybranych punktów ma więc domknięty
przeciwobraz przez każdą mapę charakterystyczną: jest to skończona
suma przeciwobrazów punktów w przestrzeni Hausdorffa.
Topologia słaba mówi, że $E$ jest domknięty w $X$.
Zatem zbiór wszystkich $c_\alpha$ jest domkniętym, dyskretnym
podzbiorem zwartego $C$, co jest niemożliwe, jeśli jest nieskończony.
Do skończenie wielu spotkanych komórek dołączamy ich ściany;
warunek skończoności domknięć i malejące wymiary dają skończony podkompleks.

Topologia słaba oznacza, że surjekcja
$q:\bigsqcup_\alpha D^{d_\alpha}\to X$ jest ilorazowa.
Iloczyn dowolnej mapy ilorazowej $q:E\to X$ z $\id_I$ też jest ilorazowy.
Aby to sprawdzić, niech przeciwobraz $W\subseteq X\times I$ będzie otwarty.
Dla $(x,t)\in W$ wybieramy zwarte otoczenie $J$ punktu $t$ w $I$
takie, że $\{x\}\times J\subset W$. Zbiór
$V=\{y:\{y\}\times J\subset W\}$ ma otwarty przeciwobraz przez $q$:
dla każdego $e$ nad $V$ skończone pokrycie zwartego $J$
prostokątami zawartymi w przeciwobrazie $W$ daje otoczenie $e$
o tej samej własności. Stąd $V$ jest otwarty i
$V\times\operatorname{int}_I J\subseteq W$ jest otoczeniem $(x,t)$.
To dowodzi tezy o iloczynie, a więc także kryterium ciągłości $H$.
\end{proof}

Z lematu wynika również rozszerzenie homologii komórkowej
z Rozdziału~\ref{ch:homologia-przestrzeni} na dowolne CW i pary CW.
Dla pary skończonych CW definiujemy najpierw
$C_\bullet^{\mathrm{CW}}(X,A)=C_\bullet^{\mathrm{CW}}(X)/C_\bullet^{\mathrm{CW}}(A)$:
jest to kompleks wolny na komórkach spoza $A$, z brzegiem modulo komórki $A$.
Krótki ciąg tych kompleksów daje długi ciąg homologii.
Naturalne porównanie ze zwykłym ciągiem pary i lemat pięciu
przenoszą izomorfizmy dla $A,X$ na grupy względne.
Każdy singularny łańcuch, jego brzeg oraz świadectwo homologiczności
są skończonymi sumami map zwartych sympleksów, więc mieszczą się
w skończonym podkompleksie. Cykl pochodzi z takiego podkompleksu,
a relacja między cyklami jest już prawdziwa w pewnym większym,
również skończonym. Tak samo jest dla skończonych sum komórek
i dla par: przecięcie skończonego podkompleksu z podkompleksem $A$
jest skończonym podkompleksem. Izomorfizmy dla modeli skończonych
są zgodne z włączeniami, bo pochodzą z naturalnych ciągów par.
Przechodzą więc na izomorfizmy dla całego CW.
Grupy łańcuchów komórkowych są przy tym \emph{sumami prostymi},
a nie iloczynami: każda kombinacja nadal ma skończony nośnik.

\begin{lemat}[Przedłużanie homotopii dla par CW]
\label{lem:hep-cw}
Niech $A\subseteq X$ będzie podkompleksem
CW. Jeśli $f:X\to Y$ i homotopia
$h:A\times I\to Y$ spełniają
$h(a,0)=f(a)$, to istnieje
$H:X\times I\to Y$ takie, że
$H(x,0)=f(x)$ oraz $H|_{A\times I}=h$.
W szczególności homotopię
zadaną na brzegu dołączonego
dysku można przedłużyć do
całego dysku.
\end{lemat}
\begin{proof}
Najpierw rozważmy jedną komórkę
$D^k$, której brzeg $S^{k-1}$
leży w już zbudowanej części.
Na
\[
 (D^k\times\{0\})\cup
 (S^{k-1}\times I)
\subset D^k\times I
\]
mamy dane: $f$ na dole
i homotopię $h$ na boku.
Ta podprzestrzeń jest
retraktem całego cylindra.
Można to zobaczyć jawnie:
pisząc $r=\|x\|$ i
\[
 \lambda(x,t)=
 \min\left\{\frac1r,\frac2{2-t}\right\},
 \qquad
 R(x,t)=\bigl(\lambda x,\,
                  2-\lambda(2-t)\bigr),
\]
przy $r=0$ przyjmujemy
$1/r=+\infty$. Dla $r\leq1/2$ minimum wynosi $2/(2-t)$,
więc wzór jest ciągły także przy $r=0$. Wartość
$R(x,t)$ leży albo na boku
$\|\lambda x\|=1$, albo
na dole z drugą współrzędną
równą $0$. Jeśli $t=0$
lub $r=1$, mamy $\lambda=1$,
więc $R$ pozostawia dane
punkty nieruchome.
Złożenie danych z $R$
przedłuża homotopię
na $D^k\times I$.

Dołączajmy teraz komórki
$X\setminus A$ według
wymiarów. Na każdym kroku
stosujemy poprzednią
konstrukcję na wszystkich
dyskach danego wymiaru;
na ich brzegach odwzorowania
już się zgadzają. Ponieważ
topologia CW jest zadana
przez odwzorowania charakterystyczne,
sklejone homotopie są
ciągłe. Dla nieskończenie wielu komórek uzasadnieniem ciągłości
z parametrem czasu jest Lemat~\ref{lem:zwartosc-cw-pi},
zastosowany do wszystkich map $D^k\times I\to Y$.
Dla $0$-komórek wybieramy homotopię stałą równą wartości $f$;
wzór na $R$ jest potrzebny dla $k\geq1$.
\end{proof}

\begin{twierdzenie}[Przybliżenie komórkowe]
\label{thm:przyblizenie-komorkowe}
Każde odwzorowanie $f:X\to Y$ kompleksów CW
jest homotopijne z odwzorowaniem
\emph{komórkowym}, czyli takim,
które spełnia
$f(X^k)\subseteq Y^k$ dla
każdego $k$. Jeśli na
podkompleksie $A\subseteq X$
odwzorowanie już jest komórkowe,
homotopię można wybrać
stałą na $A$.
\end{twierdzenie}
\begin{proof}
Wyjaśnijmy lokalny krok geometryczny. Niech
$u:D^k\to Y^m$ wysyła brzeg do $Y^{k-1}$, gdzie $m>k$.
Obraz jest zawarty w skończonym podkompleksie na mocy
Lematu~\ref{lem:zwartosc-cw-pi}. W jednej otwartej $m$-komórce
używamy współrzędnych $\operatorname{int}D^m$ i oznaczamy przez $v$
współrzędnościowy zapis $u$ na otwartym przeciwobrazie tej komórki.
Wybierzmy $0<a<b<c<1$. Zbiór $K=v^{-1}(\overline B_b)$
jest zwarty i leży we wnętrzu $D^k$: obraz $\overline B_b$
jest zwarty i domknięty w przestrzeni Hausdorffa $Y$.
Wybieramy skończone otoczenie wielościenne $P$ zbioru $K$
takie, że $K\subset\operatorname{int}P\subset P\subset v^{-1}(B_c)$.
Można je uzyskać jako sumę domkniętych kostek dostatecznie drobnej siatki
spotykających $K$. Triangulujemy $P$ i przez interpolację wartości $v$
w wierzchołkach dostajemy odcinkowo afiniczne $w$ z
$\|w-v\|<\varepsilon$, gdzie $\varepsilon<\min(b-a,1-c)$.
Istnienie takiego przybliżenia wynika z jednostajnej ciągłości $v$ na $P$:
na każdym małym sympleksie interpolujemy wartości różniące się o mniej niż $\varepsilon$.

Niech $\eta:D^k\to[0,1]$ będzie ciągła, równa $1$ na $K$
i $0$ poza $\operatorname{int}P$. Na $P$ stosujemy homotopię
$v_t=v+t\eta(w-v)$, a poza $P$ nic nie zmieniamy.
Wszystkie wartości pozostają we wnętrzu komórki i homotopia nie rusza brzegu dysku.
Poza $K$ mamy $\|v\|>b$, więc $\|v_t\|>a$.
W $K$ końcowe odwzorowanie jest równe $w$.
Jego obraz jest zawarty w skończonej sumie podprzestrzeni afinicznych
wymiaru co najwyżej $k<m$, która nie zawiera otwartej kuli $B_a$.
Istnieje więc $z\in B_a$ pominięty przez całe nowe odwzorowanie.

Promieniowa retrakcja $D^m\setminus\{z\}$ na $\partial D^m$
zostawia brzeg nieruchomy, więc po sklejeniu komórki daje homotopię
usuwającą jej wnętrze z obrazu $u$. Ciągłość z parametrem czasu
wynika z ilorazowego kryterium z poprzedniego lematu.
Usuwamy kolejno komórki najwyższego wymiaru i schodzimy z wymiarem do $k$.
Pracujemy w skończonym podkompleksie zawierającym obraz i domknięcia
spotkanych komórek, zatem proces kończy się po skończenie wielu krokach.

Zaczynamy od $X^0$:
obraz każdego wierzchołka
przesuwamy do odpowiedniego
wierzchołka w jego składowej
$Y$, a na $A$ niczego
nie zmieniamy. Załóżmy
indukcyjnie, że
$f(X^{k-1})\subseteq Y^{k-1}$.
Dla każdej $k$-komórki
stosujemy lokalny krok
do jej odwzorowania
charakterystycznego:
usuwamy po kolei obrazy
komórek celu o wymiarach
$m>k$, zachowując
jej brzeg. Dostajemy
odwzorowanie tej komórki do
$Y^k$. Dzięki
Lematowi~\ref{lem:hep-cw}
homotopię na dotychczas
zbudowanej części
przedłużamy na całe $X$
i przechodzimy do
następnego wymiaru.
Podane kroki wykonuje
się względnie $A$.
Etapowi wymiaru $k$ przydzielamy przedział
$[1-2^{-k},1-2^{-k-1}]$, a w chwili $1$ bierzemy odwzorowanie końcowe.
Każdy ustalony szkielet pozostaje nieruchomy po zakończeniu swojego etapu.
Na każdym $D^k\times I$ mamy więc skończenie wiele sklejonych homotopii
i stały odcinek końcowy. Lemat~\ref{lem:zwartosc-cw-pi}
zapewnia ciągłość całej homotopii, także w chwili $1$.
\end{proof}

\begin{uwaga}[Przybliżenie w równych wymiarach]
\label{rem:przyblizenie-rowne-wymiary}
Ta sama lokalna konstrukcja dla $k=m$ pozwala wybrać $z\in B_a$
poza obrazami ścian triangulacji $P$ i tych $m$-sympleksów,
na których $w$ ma rząd mniejszy od $m$. Jest to skończona suma
zbiorów o wymiarze mniejszym od $m$. Każdy przeciwobraz $z$
leży wtedy we wnętrzu jednego sympleksu, gdzie $w$ jest afinicznym homeomorfizmem.
Przeciwobrazów jest skończenie wiele, a każdy ma znak wyznacznika $+1$ lub $-1$.
Możemy wokół $z$ wybrać małą kulę, której przeciwobrazy są rozłącznymi
dyskami, odwzorowanymi homeomorficznie na tę kulę.
To wynika również ze zwartości: poza małymi otoczeniami znalezionych
przeciwobrazów obraz omija całe dostatecznie małe otoczenie $z$.
Ten wariant będzie użyty w dowodzie o pierwszych komórkach.
\end{uwaga}

\begin{wniosek}[Szkielet widziany przez $\pi_i$]
\label{cor:pi-szkielety}
Dla kompleksu CW $X$ para
$(X,X^n)$ jest $n$-spójna.
Włączenie $X^n\hookrightarrow X$
daje więc izomorfizm na
$\pi_i$ dla $i<n$
i surjekcję dla $i=n$.
Jeśli $1\leq i<k$, to
$\pi_i(S^k)=0$.
\end{wniosek}
\begin{proof}
Reprezentant klasy
$\pi_i(X,X^n)$ jest
odwzorowaniem $(D^i,\partial D^i)\to
(X,X^n)$, gdzie $i\leq n$.
Najpierw przybliżamy komórkowo mapę brzegu do $X^n$,
zachowując punkt bazowy, i przedłużamy tę homotopię na dysk
Lematem~\ref{lem:hep-cw}. Dopiero teraz brzeg jest komórkowy,
więc względne przybliżenie na dysku przesuwa cały obraz do
$X^i\subseteq X^n$. Klasa względna znika.
Każda składowa $X$ zawiera wierzchołek: punkt komórki można
połączyć z jej brzegiem i powtarzać to w malejących wymiarach.
To daje także stwierdzenie o składowych i przypadek $n=0$.
Długi ciąg~\eqref{eq:les-pi}
daje twierdzenie o
włączeniu szkieletu.
Sfera $S^k$ ma strukturę CW
z jednym wierzchołkiem
i jedną $k$-komórką,
więc jej $i$-szkielet
jest punktem dla $0\leq i<k$.
Każde oparte odwzorowanie
$S^i\to S^k$ jest
homotopijne do odwzorowania
komórkowego i wtedy
jego obraz jest punktem.
\end{proof}

\begin{stwierdzenie}[Wyższe grupy i nakrycia]
\label{prop:pi-n-nakrycie}
Jeśli $p:\widetilde X\to X$
jest nakryciem, $\widetilde x_0$
leży nad $x_0$, to
\[
 p_*:\pi_n(\widetilde X,\widetilde x_0)
 \xrightarrow{\ \cong\ }
 \pi_n(X,x_0)\qquad(n\geq2).
\]
\end{stwierdzenie}
\begin{proof}
Sfera $S^n$ dla $n\geq2$
jest jednospójna na mocy
Wniosku~\ref{cor:pi-szkielety}
(dla $n=2$ znika $\pi_1$,
a dla większych wymiarów
również). Każde oparte
odwzorowanie $S^n\to X$ podnosi
się więc jednoznacznie
po wskazaniu obrazu
punktu bazowego,
zgodnie z kryterium
z Twierdzenia~\ref{thm:kryterium-podnoszenia}.
Daje to surjektywność $p_*$.
Jeśli dwa podniesienia
po zrzutowaniu są
homotopijne przez odwzorowania oparte,
ich homotopię
$S^n\times I\to X$ podnosimy z kryterium podnoszenia:
$S^n\times I$ jest łukowo i lokalnie łukowo spójna oraz jednospójna,
bo retrahuje się na $S^n$. Ustalamy obraz $(\ast,0)$.
Ścieżka punktu bazowego leży w dyskretnym włóknie, więc jest stała,
a ograniczenie do $S^n\times\{0\}$ jest zadanym podniesieniem.
Jednoznaczność podnoszenia
sprawia, że na końcu
otrzymujemy drugie
podniesienie, ponieważ
oba zgadzają się w
punkcie bazowym.
To daje iniektywność.
\end{proof}

\begin{przyklad}[Okrąg i torus nie mają wyższych $\pi_n$]
\label{prz:pi-torus}
Nakrycia $\R\to S^1$ i
$\R^2\to T^2=S^1\times S^1$
mają ściągalne przestrzenie
całkowite. Zatem dla $n\geq2$
\[
 \pi_n(S^1)=0,\qquad
 \pi_n(T^2)=0.
\]
Jednocześnie
$H_2(T^2)=\Z$ z
Przykładu~\ref{prz:homologia-cw-torus}.
Już tu widać, że w
przestrzeni z nietrywialną
$\pi_1$ wyższa homologia
nie musi być obrazem
wyższej grupy homotopii.
\end{przyklad}

\section{Pierwsze odwzorowanie Hurewicza: od pętli do cykli}
\label{sec:hurewicz-jeden}

\begin{twierdzenie}[Hurewicz w stopniu pierwszym]
\label{thm:hurewicz-h1}
Dla przestrzeni łukowo spójnej $X$
odwzorowanie przypisujące opartej pętli
jej klasę singularnego
$1$-cyklu indukuje naturalny
izomorfizm
\[
 h_1:\pi_1(X,x_0)_{\mathrm{ab}}
 =\pi_1(X,x_0)/
 [\pi_1(X,x_0),\pi_1(X,x_0)]
 \xrightarrow{\ \cong\ }H_1(X;\Z).
\]
\end{twierdzenie}
\begin{proof}
Pętla $\ell:I\to X$ ma
$\partial\ell=[x_0]-[x_0]=0$.
Złożenie dwóch pętli
ma tę samą klasę homologii
co suma ich łańcuchów:
singularny $2$-sympleks, którego boki
przebiegają $\ell_1$, $\ell_2$ i ich złożenie, ma
brzeg $\ell_1+\ell_2-
(\ell_1*\ell_2)$
(z odpowiednią
parametryzacją boków).
Homotopia oparta pętli daje homologiczne cykle przez operator pryzmatyczny
z Twierdzenia~\ref{thm:homotopia-singularna}.
Zatem $[\ell]\mapsto
[\ell]_{H_1}$ jest
homomorfizmem do grupy
abelowej i zabija komutatory.

Skonstruujemy odwrotność.
Dla każdego $x\in X$
wybierzmy drogę $u_x$
od $x_0$ do $x$,
przy czym $u_{x_0}$
jest drogą stałą.
Singularnej krawędzi
$\sigma:x\to y$
przypiszmy w
$\pi_1(X,x_0)_{\mathrm{ab}}$
klasę pętli
\[
 L(\sigma)
 =[u_x*\sigma*\bar u_y].
\]
Niech $\tau:\Delta^2\to X$,
a $x_0',x_1',x_2'$
będą obrazami jej
wierzchołków.
Pętla złożona z boków
$\tau_{01}*\tau_{12}*
\bar\tau_{02}$ jest
zerowa w $\pi_1(X,x_0')$,
gdyż wypełnia ją
sam trójkąt $\tau$.
Po doprowadzeniu jej
do bazy drogami $u_{x_i'}$
otrzymujemy
\[
 L(\partial\tau)
 =L(\tau_{12})-L(\tau_{02})
  +L(\tau_{01})=0
 \quad\text{w }\pi_1(X,x_0)_{\mathrm{ab}}.
\]
Przez liniowość $L$
zeruje więc wszystkie
singularne brzegi
$2$-wymiarowe.
Na cyklach $z\in C_1(X)$
daje homomorfizm
$\bar L:H_1(X)\to
\pi_1(X,x_0)_{\mathrm{ab}}$.

Dla opartej pętli
$\ell$ mamy $L(\ell)=[\ell]$,
więc $\bar Lh_1=\id$.
W drugą stronę,
w homologii
obraz cyklu
$z=\sum m_\sigma\sigma$
rozwija się na
sumę łańcuchów $u_x$, samych krawędzi i $-u_y$,
z dokładnością do brzegów $2$-łańcuchów.
Pojedyncze $u_x$ zwykle nie są cyklami, więc nie przypisujemy
im osobnych klas w $H_1$. Rachunek przeprowadzamy w $C_1(X)/\operatorname{im}\partial_2$.
Tożsamość dla złożenia dróg wynika z trójkąta użytego na początku,
także gdy ich końce są różne; droga odwrotna daje ujemny łańcuch
modulo brzegi, bo droga złożona z własną odwrotnością jest homotopijna do stałej.
Współczynnik przy
każdej pomocniczej
drodze $u_x$ jest
równy minus współczynnikowi
przy $[x]$ w $\partial z$,
a zatem zerowy.
Pozostaje klasa
samego $z$ w $H_1(X)$;
formalnie używamy
homologii złożenia
dróg udowodnionej
na początku. Stąd
$h_1\bar L=\id$
i oba odwzorowania są
odwrotne.
\end{proof}

\begin{przyklad}[Od wolnej grupy do liczby pętli]
Bukiet $r$ okręgów ma
grupę podstawową
$F_r$ z
Przykładu~\ref{prz:pi1-dwa-okregi}
(dla $r=2$ rachunek
podano tam wprost).
Po narzuceniu
przemienności generatorów
$F_r^{\mathrm{ab}}\cong\Z^r$,
zgodnie z
Przykładem~\ref{prz:homologia-bukiet}.
Homologia pamięta tu
liczbę niezależnych
kierunków, ale nie
kolejność, w jakiej
przebywamy pętle.
\end{przyklad}

\section{Twierdzenie Whiteheada}
\label{sec:whitehead}

Równość wszystkich grup $\pi_n$
nie oznacza automatycznie, że
dowolne odwzorowanie przestrzeni
topologicznych ma homotopijną
odwrotność. Dla kompleksów CW
komórki pozwalają jednak
przekształcić informację
o sferach w deformację
całej przestrzeni.

\begin{definicja}[Słaba równoważność homotopijna]
\label{def:slaba-rownowaznosc}
Odwzorowanie $f:X\to Y$ jest
\emph{słabą równoważnością
homotopijną}, jeśli daje
bijekcję $\pi_0(X)\to\pi_0(Y)$
i dla wszystkich punktów
bazowych $x\in X$ oraz
wszystkich $n\geq1$
izomorfizmy
$f_*:\pi_n(X,x)\to
\pi_n(Y,f(x))$.
\end{definicja}

\begin{lemat}[Ściskanie komórek]
\label{lem:kompresja-cw}
Niech $(K,L)$ będzie parą CW,
a $B\subseteq Y$ niepustą
podprzestrzenią. Załóżmy, że
każda składowa łukowa $Y$ spotyka
$B$ oraz wszystkie grupy
$\pi_k(Y,B,b)$ dla
$k\geq1$ i $b\in B$
są trywialne (dla $k=1$ chodzi o zbiór z wyróżnioną klasą). Wówczas
każde odwzorowanie par
$F:(K,L)\to(Y,B)$
jest homotopijne
\emph{względem $L$}
z odwzorowaniem o obrazie w $B$.
\end{lemat}
\begin{proof}
Dla $0$-komórek $K\setminus L$
wybieramy drogi z ich obrazów
do $B$. Homotopia przesuwa
te obrazy do $B$ i jest
stała na $L$.

Załóżmy, że tak zrobiliśmy
na wszystkich komórkach
o wymiarze mniejszym od $k$.
Dla dołączonej $k$-komórki
odwzorowanie charakterystyczne
daje parę
$(D^k,S^{k-1})\to(Y,B)$.
Przy $k\geq1$ jej klasa
w $\pi_k(Y,B)$ jest zerowa.
Znaczy to, że można
deformować dysk do $B$,
pozwalając początkowo
brzegowi poruszać się w $B$.
Wyjaśnijmy, jak unieruchomić brzeg. Niech
$G:D^k\times I\to Y$ będzie względną homotopią do stałej,
z $G_0=F$ i $G_t(S^{k-1})\subset B$.
Połóżmy $a_t=1-t/2$ i dla $x=ru$ określmy
\[
 K_t(ru)=
 \begin{cases}
 G_t(ru/a_t),&r\leq a_t,\\
 G_{2(1-r)}(u),&a_t\leq r\leq1.
 \end{cases}
\]
Na $r=a_t$ oba wzory dają $G_t(u)$, na $r=1$ mamy stale $F(u)$,
a dla $t=0$ otrzymujemy $F$. Ciągłość przy zanikającym kołnierzu
wynika z ciągłości $G$; dla $t=1$ cały obraz leży w $B$.
Tak dostajemy deformację względem całego brzegu, a nie jedynie
homotopię, która pozwala mu przesuwać się w $B$.
Wykonujemy je na
wszystkich $k$-komórkach
równocześnie.
Lemat~\ref{lem:hep-cw}
przedłuża otrzymaną
homotopię do reszty
$K$, bez ruszania $L$.
Indukcja usuwa obrazy
wszystkich komórek
spoza $L$ z $Y\setminus B$.
Dla nieskończonego CW
przeznaczamy na etap $k$
coraz krótszy przedział
czasu kończący się w $1$.
Każdy ustalony szkielet staje się w końcu nieruchomy.
Ograniczenie do każdej domkniętej komórki, wraz z całym przedziałem czasu,
jest skończonym sklejeniem ciągłych homotopii i stałego końca.
Lemat~\ref{lem:zwartosc-cw-pi} daje ciągłość także w chwili $1$.
\end{proof}

\begin{twierdzenie}[Whitehead]
\label{thm:whitehead}
Słaba równoważność homotopijna
$f:X\to Y$ między kompleksami
CW jest równoważnością
homotopijną. Jeśli
$i:X\hookrightarrow Y$
jest włączeniem podkompleksu
i indukuje izomorfizmy
wszystkich $\pi_n$ oraz
bijekcję na $\pi_0$, to
$X$ jest silnym retraktem
deformacyjnym $Y$.
\end{twierdzenie}
\begin{proof}
Jeśli $X=\varnothing$, bijekcja na składowych wymusza $Y=\varnothing$
i teza jest natychmiastowa. Dalej zakładamy $X\ne\varnothing$.
Zacznijmy od włączenia.
W długim ciągu~\eqref{eq:les-pi}
pary $(Y,X)$ znajdują się
fragmenty
\[
 \pi_n(X)\xrightarrow[\cong]{i_*}
 \pi_n(Y)\longrightarrow
 \pi_n(Y,X)\longrightarrow
 \pi_{n-1}(X)
 \xrightarrow[\cong]{i_*}
 \pi_{n-1}(Y).
\]
Dokładność wymusza
$\pi_n(Y,X)=0$ dla
wszystkich $n\geq2$.
Dla stopnia $1$ warto użyć dróg, ponieważ nie jest to już ciąg grup.
Droga względna od $a\in X$ do $x_0$ ma oba końce w jednej składowej $X$,
gdyż $\pi_0(X)\to\pi_0(Y)$ jest injekcją. Dopinamy drogę
od $x_0$ do $a$ w $X$, otrzymując pętlę w $Y$.
Surjektywność $\pi_1(X)\to\pi_1(Y)$ pozwala zastąpić ją pętlą w $X$.
Każda taka droga jest trywialna względnie $X$, więc $\pi_1(Y,X)$
ma tylko wyróżniony element. Każda składowa $Y$ spotyka $X$,
bo odwzorowanie na $\pi_0$ jest także surjekcją.
Lemat~\ref{lem:kompresja-cw}
zastosowany do
$F=\id_Y:(Y,X)\to(Y,X)$
daje homotopię nieruchomą
na $X$, prowadzącą od
$\id_Y$ do odwzorowania
$r:Y\to X$.
Ponieważ $r|_X=\id_X$,
jest to silna retrakcja
deformacyjna.

Niech teraz $f$ będzie
dowolną słabą równoważnością
CW. Twierdzenie
\ref{thm:przyblizenie-komorkowe}
pozwala zastąpić ją
odwzorowaniem komórkowym $f'\simeq f$.
Zmiana punktu bazowego wzdłuż tej homotopii pokazuje,
że $f'$ nadal jest słabą równoważnością.
Utwórzmy jego
\emph{cylinder odwzorowania}
\[
 M_{f'}=(X\times I)\sqcup Y/
       \bigl((x,1)\sim f'(x)\bigr).
\]
Cylinder ma strukturę CW: poza komórkami $Y$ bierzemy dolne
komórki $e^k\times\{0\}$ i komórki $e^k\times(0,1)$ wymiaru $k+1$.
Brzeg każdej z tych ostatnich składa się z dołu, boków nad
brzegiem $e^k$ i góry odwzorowanej przez $f'$ do $Y^k$.
Komórkowość $f'$ zapewnia zgodność wymiarów, a topologia
ilorazowa i Lemat~\ref{lem:zwartosc-cw-pi} dają topologię słabą.
Dolna kopia $X=X\times\{0\}$ jest więc podkompleksem $M_{f'}$.
Zwijając każdy odcinek
$\{x\}\times I$ do
$f'(x)\in Y$, otrzymujemy
retrakcyjną deformację
$M_{f'}\to Y$.
Złożenie włączenia
$X\hookrightarrow M_{f'}$
z tą retrakcją to $f'$.
Ponieważ $f'_*$ i
retrakcyjne odwzorowanie na $Y$
indukują izomorfizmy
na $\pi_n$, również
włączenie $X\hookrightarrow
M_{f'}$ je indukuje.
Pierwsza część dowodu
ściska więc $M_{f'}$
do $X$. Oznaczmy dolne włączenie przez $j$, włączenie $Y$ do cylindra przez $i$,
a retrakcje na $X,Y$ odpowiednio przez $s,r$.
Wtedy $f'=rj$, a odwrotnością homotopijną jest $g=si$:
$gf'=sirj\simeq sj=\id_X$ oraz $f'g=rjsi\simeq ri=\id_Y$.
Ponieważ $f\simeq f'$, to samo $g$ jest odwrotnością homotopijną $f$.
\end{proof}

\begin{figure}[!htbp]
\centering
\begin{tikzpicture}[>=Stealth,line cap=round,line join=round]
 \fill[manifoldblue!10]
  (0,0)--(4.1,0)--(4.35,2.0)--(.25,2.0)--cycle;
 \draw[manifoldblue!70!black,thick]
  (0,0)--(4.1,0)--(4.35,2.0)--(.25,2.0)--cycle;
 \foreach \x in {.7,1.8,2.9,3.7}{
  \draw[manifoldblue!35] (\x,.07)--({\x+.23},1.94);
 }
 \fill[manifoldgreen!18]
  (-.38,2.0) rectangle (4.86,2.5);
 \draw[manifoldgreen!70!black,thick]
  (-.38,2.0) rectangle (4.86,2.5);
 \node at (2.2,2.25) {$Y$};
 \node[below] at (2.08,0) {$X\times\{0\}$};
 \node[manifoldblue!70!black] at (2.2,1.02)
  {$X\times I$};
 \node[manifoldred!85!black,above]
  at (2.2,2.62) {$(x,1)\sim f'(x)$};
 \draw[->,manifoldred,thick]
  (3.35,1.48)--(3.35,2.2)
  node[midway,right=3pt] {$r$};
 \node[font=\small,align=left,anchor=west,text width=4.8cm]
  at (5.15,1.15)
  {Cylinder zwija się na $Y$.
   Gdy $f'$ jest słabą równoważnością,
   poprzednia część dowodu zwija go
   również na dolną kopię $X$.};
\end{tikzpicture}
\caption{Cylinder odwzorowania zamienia
odwzorowanie $f':X\to Y$ na włączenie
podkompleksu $X\hookrightarrow M_{f'}$.
Rysunek pokazuje mechanizm dowodu
twierdzenia Whiteheada.}
\label{fig:cylinder-whiteheada}
\end{figure}

\begin{uwaga}[Zakres twierdzenia]
Sprawdzenie tylko $\pi_1$
i kilku pierwszych grup
$\pi_n$ nie wystarcza
do zastosowania Whiteheada:
założenie dotyczy
\emph{każdego} $n$.
Również samo podobieństwo
grup homologii dwóch
przestrzeni nie wskazuje
jeszcze odwzorowania, które
mogłoby być równoważnością.
Niżej wyprowadzimy
użyteczne kryterium
homologiczne przy
dodatkowej jednospójności.
\end{uwaga}

\section{Odwzorowanie Hurewicza i pierwszy niezerowy wymiar}
\label{sec:hurewicz-wyzszy}

\begin{definicja}[Odwzorowanie Hurewicza]
\label{def:mapa-hurewicza}
Homologia $H_n(S^n;\Z)$
jest izomorficzna z $\Z$
na mocy
Twierdzenia~\ref{thm:homologia-sfer}.
Wybierzmy generatory zgodnie z orientacjami dysków i ich brzegów,
tak aby $\delta[D^n,S^{n-1}]=[S^{n-1}]$.
Dla $n\geq2$ określamy
\[
 h_n:\pi_n(X,x_0)\longrightarrow H_n(X;\Z),
 \qquad h_n([f])=f_*[S^n].
\]
Dla pary $(X,A)$ bierzemy
analogicznie generator
$[D^n,S^{n-1}]
\in H_n(D^n,S^{n-1})$ i
kładziemy
\[
 h_n^{\mathrm{rel}}([F])
 =F_*[D^n,S^{n-1}]
 \in H_n(X,A).
\]
W stopniu $1$ dla przestrzeni używamy homomorfizmu
$\pi_1(X)\to\pi_1(X)_{\mathrm{ab}}\xrightarrow{h_1}H_1(X)$.
Względny homomorfizm określamy tutaj dla $n\geq2$;
zbiór $\pi_1(X,A)$ nie jest w ogólności grupą.
\end{definicja}

Homotopie oparte są w
szczególności zwykłymi
homotopiami odwzorowań, więc
$h_n$ jest dobrze
określone dzięki
Twierdzeniu~\ref{thm:homotopia-singularna}.
Odwzorowanie jest homomorfizmem:
sklejanie dwóch połówek
sfery odpowiada sumie
ich klas fundamentalnych.
Dokładniej, sklejanie jest złożeniem mapy ściskającej równik
$S^n\to S^n\vee S^n$ z $f\vee g$. Mapa ściskająca posyła
generator na sumę dwóch generatorów: rzut na każdy składnik
ma stopień $+1$, co widać z klasy orientacji odpowiedniej półkuli
i naturalności odwzorowania łączącego dla dysku i brzegu.
Względnie używamy kostki podzielonej na dwie połowy w pierwszej współrzędnej.
Po zorientowaniu obu połówek zgodnie z całą kostką ich wewnętrzne
ściany kasują się; obraz łańcucha fundamentalnego jest sumą obu obrazów.
To daje również homomorfizm względny. Naturalność wynika od razu ze składania map.

Do dowodu potrzebujemy
przeliczyć grupy homotopii
prostych modeli komórkowych.
Najważniejszy jest
następujący krok: w
pierwszym nowym wymiarze
komórki są niezależne,
a komórki o wymiar
wyższym wprowadzają
relacje.

\begin{lemat}[Pierwsze komórki i ich relacje]
\label{lem:pierwsze-komorki-pi}
Niech $A$ będzie jednospójnym
podkompleksem CW kompleksu $Z$.
Załóżmy, że wszystkie komórki
$Z\setminus A$ mają wymiar
co najmniej $n\geq2$.
Oznaczmy przez $I_n$
zbiór nowych $n$-komórek.
Wtedy
\[
 \pi_n(Z^n\cup A,A)
 \cong\bigoplus_{\alpha\in I_n}\Z[e_\alpha^n]
 \xrightarrow[\cong]{\ h_n^{\mathrm{rel}}\ }
 H_n(Z^n\cup A,A).
\]
Po dołączeniu $(n+1)$-komórek
obie grupy otrzymujemy
z tej wolnej grupy przez
te \emph{same} relacje:
obraz brzegu odwzorowania
dołączającego $S^n\to
Z^n\cup A$. Komórki
o wymiarach co najmniej
$n+2$ nie zmieniają
tych grup.
\end{lemat}
\begin{proof}
Niech $W=Z^n\cup A$ oraz $W'=Z^{n+1}\cup A$.
Przestrzeń $W$ jest jednospójna: do jednospójnego $A$
dołączamy tylko komórki wymiaru co najmniej $2$, co nie tworzy
nowych pętli (lokalny krok przybliżenia komórkowego przesuwa pętlę do $A$).
W szczególności także $\pi_2(W,A)$ jest abelowa, bo z ciągu pary
jest ilorazem abelowej $\pi_2(W)$ przy $\pi_1(A)=0$.
Możemy więc mówić o homomorfizmie z wolnej grupy abelowej
również w przypadku $n=2$.
Po zwinięciu $A$
otrzymujemy bukiet
$\bigvee_{\alpha\in I_n}S^n$.
Względna homologia
$H_n(W,A)$ jest
wolna na komórkach
na mocy
Twierdzenia~\ref{thm:warstwa-komorek}
i rozszerzenia na pary CW po Lemacie~\ref{lem:zwartosc-cw-pi};
to jeszcze nie jest
argument o homotopii,
więc rozpiszmy go
osobno.

Każdą mapę charakterystyczną nowej komórki opieramy w $x_0$
przez dołączenie drogi w $A$ do obrazu wybranego punktu jej brzegu.
Wybór drogi nie zmienia klasy, bo $A$ jest jednospójna.
Tak rozumiemy symbole $[e_\alpha^n]$ po stronie homotopii.
Weźmy teraz odwzorowanie względne
$F:(D^n,S^{n-1},*)\to(W,A,x_0)$.
Zwarty obraz spotyka
tylko skończenie wiele
otwartych komórek.
W każdej z nowych $n$-komórek stosujemy przybliżenie
z Uwagi~\ref{rem:przyblizenie-rowne-wymiary}, nieruchome przy $S^{n-1}$.
Wybieramy regularny punkt $z_\alpha$ poza obrazami ścian triangulacji.
Ma on skończenie wiele przeciwobrazów; na małych otoczeniach
$F$ jest afinicznym homeomorfizmem o znaku $+1$ lub $-1$.
Wybierzmy wokół $z_\alpha$ małą kulę. Jej przeciwobrazy są rozłącznymi
dyskami $D_j$, a poza nimi $F$ omija tę kulę.
Parametryzując $D_j$ przez odwrotność lokalnej mapy afinicznej,
otrzymujemy dokładnie lokalny fragment mapy charakterystycznej;
ujemny znak odpowiada odbiciu pierwszej współrzędnej, czyli odwrotności w grupie. Po usunięciu
środków każda komórka
promieniowo zwija się
na swój brzeg w $A$;
zatem odwzorowanie na dopełnieniu
małych dysków możemy
przesunąć do $A$.
Przedłużamy tę deformację na $D_j$ przez rozszerzanie ich obrazów
od małej kuli do całej komórki, zostawiając zgodność na wspólnych brzegach.
Połączmy małe dyski i punkt bazowy drzewem łuków w ich dopełnieniu.
Rozcinając wzdłuż cienkich pasów tego drzewa, przedstawiamy dysk
jako sumę dysków z doprowadzeniem do bazy; pozostałe części mają obraz w $A$
i reprezentują zero względne. Takie rozcięcie realizuje kolejno
kostkowe dodawanie z Definicji~\ref{def:pi-wzgledne}:
wewnętrzny pas ma obraz w $A$ i można go przesuwać podczas homotopii względnej.
Zmiana drogi doprowadzającej modyfikuje składnik przez pętlę w $A$.
Ponieważ $\pi_1(A)=0$, te zmiany nie wpływają na klasę.
Ostatecznie $[F]$ jest sumą klas charakterystycznych dysków
$e_\alpha^n$, z ich znakami. Każda suma jest skończona.
To dowodzi surjektywności
odwzorowania od wolnej grupy
na nowych komórkach
do $\pi_n(W,A)$.
Obrazy tych generatorów
przez $h_n^{\mathrm{rel}}$
tworzą bazę
$H_n(W,A)$. Gdyby ich
kombinacja dawała
zerową klasę względnej
homotopii, jej obraz
homologiczny również
byłby zerowy; niezależność
bazy wymusza wszystkie
współczynniki równe $0$.
Mamy więc izomorfizm.

Niech teraz dołączamy dyski $D^{n+1}_\beta$, otrzymując $W'$.
Każdą mapę $D^n\to W'$ względnie $A$ przesuwamy do $W$
przez lokalne usuwanie punktów w komórkach wymiaru $n+1$.
Włączenie indukuje więc surjekcję na $\pi_n$.
Relację od $\varphi_\beta:S^n\to W$ rozumiemy jako jej obraz
przez $\pi_n(W)\to\pi_n(W,A)$.
Ich odwzorowania charakterystyczne
dają oczywiste relacje:
brzeg każdego takiego
dysku staje się zerem
w $\pi_n(W',A)$.
Dlaczego nie pojawiają
się inne relacje?
Homotopię względną
między dwoma sumami
$n$-dysków przedstawiamy
na $D^n\times I$.
Uwaga~\ref{rem:przyblizenie-rowne-wymiary}, zastosowana do cylindra
triangulowanego jako $(n+1)$-dysk i nieruchomo na całym jego brzegu,
zapewnia, że przeciwobrazy wybranych punktów $(n+1)$-komórek
są skończonym zbiorem
we wnętrzu $(n+1)$-wymiarowego
cylindra. Usuwamy małe
otoczenia tych punktów.
Pozostała część odwzorowania
omija środki i promieniowo
przesuwa się do $W$.
Na brzegu każdego
usuniętego otoczenia
otrzymujemy, ze znakiem,
odwzorowanie dołączające
odpowiedniej komórki.
Łączymy te otoczenia z brzegiem cylindra drzewem łuków.
Po rozcięciu wzdłuż drzewa zorientowany brzeg pozostałego obszaru
daje różnicę końcowych $n$-dysków oraz sumę nowych sfer brzegowych.
Jego wypełnienie leży w $W$, a pozostałe ściany w $A$,
więc w $\pi_n(W,A)$ otrzymujemy równość różnicy klas z sumą
obrazów $[\varphi_\beta]$, ze znakami lokalnych homeomorfizmów.
Drogi doprowadzające leżą teraz w $W$; nie zmieniają klas,
ponieważ $W$ jest jednospójna. To wyjaśnia, dlaczego dostajemy
zwykłe całkowite kombinacje, bez dodatkowych sprzężeń.
Każda relacja w
$\pi_n(W',A)$
jest więc kombinacją
brzegów dołączonych
komórek.
W homologii to samo
stwierdza wzór
na brzeg komórkowy
\eqref{eq:formula-komorkowa}:
stopnie odwzorowań dołączających
są dokładnie
współczynnikami
tych relacji.

Wreszcie dla $k\geq n+2$
homotopię dysku $D^n$
można przybliżyć
komórkowo, zachowując
brzeg w $A$: jej obraz
nie musi wchodzić
do nowych $k$-komórek.
To samo dotyczy
$(n+1)$-wymiarowej
homotopii, bo $n+1<k$.
Zatem grupy $\pi_n$
nie zmieniają się.
Również względny
kompleks komórkowy
nie zmienia swoich
grup i brzegów
w stopniach $n$ i $n+1$.
\end{proof}

\begin{wniosek}[Zabijanie wybranych klas]
\label{cor:zabijanie-klas-pi}
Niech $Y$ będzie jednospójnym CW i $k\geq2$.
Dołączenie $(k+1)$-komórek wzdłuż klas $a_\beta\in\pi_k(Y)$
zmienia $\pi_k(Y)$ na iloraz przez podgrupę generowaną przez $a_\beta$,
a nie zmienia $\pi_i$ dla $i<k$.
\end{wniosek}
\begin{proof}
Oznaczmy nową przestrzeń przez $V$. Poprzedni lemat z bazą $A=Y$
i pierwszym nowym wymiarem $k+1$ mówi, że
$\pi_{k+1}(V,Y)$ jest wolna na dołączonych komórkach.
Mapa brzegowa do $\pi_k(Y)$ posyła ich generatory na $a_\beta$.
Niższe grupy względne znikają przez lokalny krok przybliżenia
komórkowego, ponieważ dyski tych wymiarów można przesunąć do $Y$.
Długi ciąg pary daje dokładnie podany iloraz i izomorfizmy w niższych stopniach.
\end{proof}

\begin{uwaga}[O przybliżeniu w lemacie]
Warunek $\pi_1(A)=0$ ma
konkretną rolę: pozwala
jednoznacznie przenosić
lokalne dyski do bazy.
Gdy $A$ nie jest
jednospójna, powstają
przesunięte przez pętle
kopie generatorów.
Wersja względnego
twierdzenia Hurewicza
bez tego warunku
wymaga ilorazu przez
działanie $\pi_1(A)$;
nie utożsamiamy wtedy
bezpośrednio całej
$\pi_n(X,A)$ z
$H_n(X,A)$.
\end{uwaga}

\begin{lemat}[Model bez niższych komórek]
\label{lem:model-wysoko-spojny}
Jeśli kompleks CW $X$
jest $(n-1)$-spójny,
$n\geq2$, to jest
homotopijnie równoważny
kompleksowi $Z$ z
jedną $0$-komórką
i bez komórek
wymiarów $1,\ldots,n-1$.
Odwzorowanie $Z\to X$
można dobrać tak,
by przenosiło
wyróżniony punkt
na $x_0$.
\end{lemat}
\begin{proof}
Zaczynamy od jednej
$0$-komórki $Z^0=\{*\}$
wysłanej do $x_0$.
Ponieważ $\pi_i(X)=0$
dla $i<n$, odwzorowanie to
jest już izomorfizmem
na tych grupach.
Dołączamy do punktu
po jednej $n$-komórce
dla wybranego zbioru
generatorów $\pi_n(X)$.
Każda taka komórka
z brzegiem sklejonym
do punktu jest $S^n$;
odwzorowanie do $X$ definiujemy
na niej reprezentantem
wybranej klasy.
Uzyskujemy surjekcję
na $\pi_n(X)$.

Teraz dla każdego
generatora jądra
$\pi_n(Z^n)\to\pi_n(X)$
dołączamy $(n+1)$-komórkę
wzdłuż reprezentującego
go odwzorowania $S^n\to Z^n$.
Jej obraz w $X$
jest zerowy w $\pi_n(X)$,
więc wybieramy
wypełnienie $D^{n+1}\to X$
i przedłużamy odwzorowanie
$Z\to X$.
Wniosek~\ref{cor:zabijanie-klas-pi}, zastosowany do $Y=Z^n$,
mówi, że dokładnie te relacje zabijają wybrane jądro; w
stopniu $n$ odwzorowanie
staje się izomorfizmem.
Dołączone komórki
nie zmieniają niższych
grup dzięki
Wnioskowi~\ref{cor:pi-szkielety}.

Powtarzamy procedurę
w stopniu $k=n+1,n+2,\ldots$:
komórki $k$-wymiarowe
dołączone do punktu
realizują brakujące
elementy $\pi_k(X)$,
a komórki $(k+1)$-wymiarowe
dołączone wzdłuż klas
z jądra zabijają
dokładnie te klasy.
Dołączenie sfer $S^k$ do punktu daje surjekcję na brakujące klasy
i nie narusza już ustalonych niższych grup: nowy bukiet ma retrakcję
na poprzedni etap, a względne dyski wymiaru $<k$ można przesunąć do niego.
Do zabijania jądra w każdym stopniu stosujemy
Wniosek~\ref{cor:zabijanie-klas-pi} z bazą równą \emph{całemu bieżącemu etapowi}.
Nie wolno tu ponownie traktować punktu jako jedynej bazy pierwszych komórek,
bo wcześniej dołączono już komórki niższych wymiarów.
Mapy dołączające przybliżamy komórkowo, dzięki czemu otrzymujemy strukturę CW;
homotopia takiej mapy nie zmienia jej klasy ani możliwości wypełnienia w $X$.
Wszystkie odwzorowania do $X$ przedłużają się,
bo obrazy odwzorowań
dołączających z jądra
są zerowe w $\pi_k(X)$.
Komórki późniejszego
wymiaru nie naruszają
ustalonych już
izomorfizmów w
niższych stopniach.
Po wykonaniu
wszystkich etapów otrzymujemy CW $Z$.
Mapy sfer i homotopie mają zwarte dziedziny, więc z Lematu~\ref{lem:zwartosc-cw-pi}
każda klasa i każda relacja pojawiają się już w skończonym etapie.
Stabilizacja grup oznacza zatem, że odwzorowanie $Z\to X$
indukuje izomorfizmy
na wszystkich
$\pi_k$ oraz na
$\pi_0$. Twierdzenie
\ref{thm:whitehead}
zamienia tę słabą
równoważność CW
w równoważność
homotopijną.
\end{proof}

\begin{twierdzenie}[Hurewicz: pierwszy niezerowy wymiar]
\label{thm:hurewicz-wyzszy}
Jeżeli $X$ jest
$(n-1)$-spójnym
kompleksem CW,
$n\geq2$, to
\[
 \widetilde H_i(X;\Z)=0
 \quad\text{dla }i<n,
 \qquad
 h_n:\pi_n(X,x_0)
 \xrightarrow{\ \cong\ }H_n(X;\Z).
\]
Izomorfizm jest naturalny:
dla odwzorowania opartego
$f:X\to Y$ odpowiedni
kwadrat z $f_*$ na
homotopii i homologii
jest przemienny.
\end{twierdzenie}
\begin{proof}
Lemat~\ref{lem:model-wysoko-spojny}
pozwala zastąpić
$X$ homotopijnie
równoważnym $Z$
z $Z^{n-1}=\{*\}$.
Homologia i homotopia
nie zmieniają się
przy tej operacji
(Twierdzenia
\ref{thm:homotopia-singularna}
i~\ref{thm:whitehead}).
W kompleksie komórkowym
$Z$ nie ma grup
łańcuchów w
stopniach $1,\ldots,n-1$,
więc
$\widetilde H_i(Z)=0$
dla $i<n$.

Niech
$E_n=\bigoplus_{\alpha\in I_n}
\Z[e_\alpha^n]$.
Z Lematu~\ref{lem:pierwsze-komorki-pi}
dla $A=\{*\}$ mamy
$\pi_n(Z^n)\cong E_n$.
Włączenie
$Z^n\hookrightarrow Z$
jest surjekcją
na $\pi_n$,
a jądrem są relacje
narzucone przez
$(n+1)$-komórki.
Po lewej otrzymujemy
zatem
\[
 \pi_n(Z)
 \cong
 E_n/\langle
   [\varphi_\beta]:
   \varphi_\beta:S^n\to Z^n
   \text{ dołącza }e_\beta^{n+1}
 \rangle.
\]
Po prawej
$\partial_n^{\mathrm{CW}}=0$,
bo $C_{n-1}^{\mathrm{CW}}=0$,
więc
\[
 H_n(Z)
 \cong E_n/
     \operatorname{im}
       \partial_{n+1}^{\mathrm{CW}}.
\]
Obraz $[\varphi_\beta]$
przez $h_n$ jest
klasą brzegową
$(n+1)$-komórki:
jej współczynniki
w bazie $E_n$
to te same stopnie
odwzorowań dołączających
z Twierdzenia
\ref{thm:formula-komorkowa}.
Oba ilorazy są
więc ilorazami
\emph{tej samej}
wolnej grupy przez
\emph{ten sam}
podmoduł.
Ponieważ $h_n$
na generatorach
$[e_\alpha^n]$
przypisuje im
klasy komórkowe
$[e_\alpha^n]$,
jest izomorfizmem.

Na koniec
$h_n(f_*[u])
=(f\circ u)_*[S^n]
=f_*(u_*[S^n])$,
co dowodzi
naturalności.
\end{proof}

\begin{twierdzenie}[Hurewicz względny dla par jednospójnych]
\label{thm:hurewicz-wzgledny}
Niech $(X,A)$ będzie parą CW,
$A\ne\varnothing$ i $A$
jednospójna. Jeśli para
jest $(n-1)$-spójna,
$n\geq2$, to
\[
 H_i(X,A;\Z)=0\quad(i<n),
 \qquad
 h_n^{\mathrm{rel}}:
 \pi_n(X,A,x_0)
 \xrightarrow{\ \cong\ }
 H_n(X,A;\Z).
\]
\end{twierdzenie}
\begin{proof}
Zbudujemy najpierw
model pary bez nowych
komórek poniżej
wymiaru $n$.
Długi ciąg~\eqref{eq:les-pi}
oraz $\pi_i(X,A)=0$
dają izomorfizmy
$\pi_i(A)\to\pi_i(X)$
dla $i<n-1$ i
surjekcję dla
$i=n-1$.
Zaczynamy od
identyczności $A\to A$
oraz włączenia $A\to X$.
Dołączamy $n$-komórki
wzdłuż reprezentantów
jądra
$\pi_{n-1}(A)\to
\pi_{n-1}(X)$.
Każdy taki reprezentant
wypełnia się w $X$,
więc odwzorowanie do $X$
przedłuża się na
nową komórkę.
W długim ciągu
pary nowych komórek
odwzorowanie brzegowe
$\pi_n(Z^n,A)\to
\pi_{n-1}(A)$
posyła każdą
charakterystyczną
$n$-komórkę na
klasę jej odwzorowania
dołączającego.
Lemat~\ref{lem:pierwsze-komorki-pi},
zastosowany do
jednospójnego $A$,
pokazuje, że
te dołączenia
zabijają dokładnie
podgrupę przez
nie generowaną,
czyli wybrane jądro.
W stopniach
niższych niczego
nie zmieniają.
Gdy $n=2$,
wyjściowe $\pi_1(A)$
jest już zerowe
i ten krok nie
wymaga komórek.

Teraz dla $k=n,n+1,\ldots$
postępujemy jak w
dowodzie Lematu
\ref{lem:model-wysoko-spojny}:
dołączamy $k$-komórki
z brzegiem wysłanym
do $x_0$, aby
otrzymać surjekcję
na $\pi_k(X)$,
a potem
$(k+1)$-komórki
wzdłuż generatorów
jądra, aby
otrzymać izomorfizm.
Nigdy nie dołączamy
komórek wymiaru
mniejszego od $n$.
W każdym kroku wniosek o zabijaniu klas stosujemy z bieżącym,
jednospójnym etapem jako bazą. Przybliżamy komórkowo mapy dołączające.
Powstaje odwzorowanie par $F:(Z,A)\to(X,A)$
równe identyczności na $A$, które
indukuje izomorfizmy
wszystkich grup
homotopii przestrzeni
oraz składowych.
Twierdzenie~\ref{thm:whitehead}
zapewnia równoważność
homotopijną $Z\simeq X$.
Porównując długie ciągi
homotopii i homologii
obu par, widzimy
ponadto, że $F$
indukuje izomorfizmy
grup \emph{względnych}:
odwzorowania na wyrazach dotyczących $A$ są identycznością,
a na wyrazach dotyczących $Z,X$ izomorfizmami.
Dla stopni co najmniej $3$ stosujemy lemat pięciu do grup abelowych.
W stopniu $2$ względne grupy są ilorazami $\pi_2(Z)$ i $\pi_2(X)$
przez obrazy $\pi_2(A)$, ponieważ $\pi_1(A)=0$;
to daje izomorfizm bez stosowania lematu pięciu do zbiorów.
Stopień $1$ jest trywialny po obu stronach, bo wszystkie trzy przestrzenie
są jednospójne. Dla homologii stosujemy zwykły lemat pięciu.

W modelu $(Z,A)$
nie ma względnych
łańcuchów komórkowych
w wymiarach $<n$,
więc $H_i(Z,A)=0$
dla $i<n$.
Lemat~\ref{lem:pierwsze-komorki-pi}
oblicza zarówno
$\pi_n(Z,A)$,
jak i $H_n(Z,A)$
jako tę samą
wolną grupę na
$n$-komórkach
modulo te same
relacje od
$(n+1)$-komórek.
Na generatorach
odwzorowanie $h_n^{\mathrm{rel}}$
przypisuje dyskowi
jego klasę orientacji,
więc jest izomorfizmem.
Naturalność
$h_n^{\mathrm{rel}}$
przenosi ten wynik
przez $F$
na parę $(X,A)$.
\end{proof}

\begin{przyklad}[Dysk względem brzegu]
\label{prz:pi-dysk-brzeg}
Dla pary $(D^{n+1},S^n)$
długi ciąg homotopii
i ściągalność dysku
dają w szczególności
\[
 \partial:
 \pi_{n+1}(D^{n+1},S^n)
 \xrightarrow{\ \cong\ }
 \pi_n(S^n)
 \qquad(n\geq1).
\]
Generator po lewej jest reprezentowany przez dysk tożsamościowy,
a jego brzegiem jest
tożsamość sfery.
W homologii dzieje się
to samo na mocy
Twierdzenia~\ref{thm:homologia-sfer};
oba ciągi łączy odwzorowanie
Hurewicza.
\end{przyklad}

\begin{wniosek}[Homologiczne kryterium Whiteheada]
\label{cor:whitehead-homologiczny}
Niech $f:X\to Y$ będzie odwzorowaniem
między jednospójnymi
kompleksami CW.
Jeśli
$f_*:H_i(X;\Z)\to H_i(Y;\Z)$
jest izomorfizmem
dla wszystkich $i$,
to $f$ jest
równoważnością
homotopijną.
\end{wniosek}
\begin{proof}
Zastępujemy $f$
odwzorowaniem komórkowym
homotopijnym z nim
i tworzymy cylinder
$M_f$ z dowodu
Twierdzenia~\ref{thm:whitehead}.
Włączenie
$i:X\hookrightarrow M_f$
ma te same odwzorowania
na homologii co $f$
po złożeniu z
retrakcyjną deformacją
$M_f\to Y$.
Zatem w długim
ciągu homologii
pary $(M_f,X)$
wszystkie grupy
$H_i(M_f,X)$
są zerowe.
Ponieważ $X$
i $M_f\simeq Y$
są jednospójne,
długi ciąg homotopii
daje
$\pi_1(M_f,X)=0$.

Załóżmy, że pewna
wyższa względna
grupa homotopii
nie znika, i weźmy
najmniejsze $n\geq2$
z $\pi_n(M_f,X)\ne0$.
Para jest wtedy
$(n-1)$-spójna,
a $X$ jednospójna.
Twierdzenie
\ref{thm:hurewicz-wzgledny}
dałoby
$H_n(M_f,X)\cong
\pi_n(M_f,X)\ne0$,
sprzeczność.
Wszystkie względne
$\pi_n$ są więc
zerowe, a długi
ciąg homotopii
pokazuje, że
$i$, a zatem $f$,
indukuje izomorfizmy
wszystkich $\pi_n$.
Twierdzenie
\ref{thm:whitehead}
kończy dowód.
\end{proof}

\section{Obliczenia i ograniczenia twierdzeń}
\label{sec:pi-przyklady}

\begin{przyklad}[Pierwsza grupa homotopii sfery]
\label{prz:pi-sfera-hurewicz}
Z Wniosku~\ref{cor:pi-szkielety}
wiemy, że $S^n$
jest $(n-1)$-spójna.
Twierdzenia
\ref{thm:hurewicz-wyzszy}
i~\ref{thm:homologia-sfer}
dają dla $n\geq2$
\[
 \pi_i(S^n)=0\quad(1\leq i<n),
 \qquad
 \pi_n(S^n)\cong\Z.
\]
Tożsamość $S^n\to S^n$
odpowiada generatorowi,
a odwzorowanie $g:S^n\to S^n$
odpowiada liczbie
$\deg g$ z
Definicji~\ref{def:stopien-sfer}.
Zatem oparte $g$ jest homotopijne do tożsamości lub do ustalonego
odbicia sfery zachowującego punkt bazowy odpowiednio wtedy,
gdy stopień wynosi $+1$ lub $-1$.
Odbicie ma stopień $-1$, bo odwraca orientację dysku i jego brzegu.
Dla nieopartego $g$ najpierw przesuwamy obraz wyróżnionego punktu do bazy
i przedłużamy tę homotopię na sferę przez własność HEP.
Jednospójność sfery sprawia, że wybór drogi nie wpływa na wynik. W szczególności
odwzorowanie sfery stopnia
$\pm1$ jest
równoważnością
homotopijną:
wywołuje izomorfizm
na $H_n$ i $H_0$,
a inne grupy homologii
sfery są zerowe,
więc można użyć
Wniosku~\ref{cor:whitehead-homologiczny}.
\end{przyklad}

\begin{przyklad}[Kompleks Moore'a]
\label{prz:moore-hurewicz}
Niech $n\geq2$ i $m\geq2$.
Do $S^n$ dołączmy
jedną $(n+1)$-komórkę
odwzorowaniem stopnia $m$.
Powstała przestrzeń
$M(\Z/m\Z,n)$ ma
kompleks komórkowy
\[
 0\longrightarrow\Z
 \xrightarrow{\;\times m\;}\Z
 \longrightarrow0
\]
w dodatnich stopniach $n+1$ i $n$; ponadto $C_0=\Z$,
a pozostałe grupy łańcuchów są zerowe. Zatem
$H_n(M)=\Z/m\Z$,
$H_{n+1}(M)=0$.
Ponieważ przestrzeń
jest $(n-1)$-spójna,
Hurewicz daje
$\pi_n(M)=\Z/m\Z$.
Intuicja jest
dosłowna: dołączony
dysk czyni
$m$-krotne
owinięcie sfery
zerem w pierwszej
niezerowej grupie.
\end{przyklad}

\begin{przyklad}[Dlaczego potrzeba jednospójności]
\label{prz:rp2-pi2}
Nakrycie antypodalne
$p:S^2\to\mathbb{RP}^2$
indukuje izomorfizm
$\pi_2(S^2)\cong
\pi_2(\mathbb{RP}^2)$
na mocy
Stwierdzenia~\ref{prop:pi-n-nakrycie}.
Zatem
$\pi_2(\mathbb{RP}^2)\cong\Z$.
Tymczasem
$H_2(\mathbb{RP}^2)=0$
z Przykładu~\ref{prz:homologia-rp2},
więc tutaj
$h_2:\Z\to0$
nie jest izomorfizmem.
Przestrzeń rzutowa
nie jest jednospójna:
$\pi_1(\mathbb{RP}^2)=\Z/2\Z$.
To konkretna granica
Twierdzenia~\ref{thm:hurewicz-wyzszy}.
\end{przyklad}

\begin{przyklad}[Bukiet $S^1\vee S^2$]
\label{prz:pi-bukiet-pokrycie}
Przestrzeń
$X=S^1\vee S^2$
ma $\pi_1(X)\cong\Z$.
Jej nakrycie uniwersalne
składa się z prostej
z kopią $S^2$ w każdym
punkcie całkowitym.
Pokrycie otrzymujemy przez przesunięcia prostej o liczby całkowite,
przenoszące jednocześnie $S_k^2$ na $S_{k+1}^2$.
Jego przestrzeń jest jednospójna: każda pętla ma obraz
w skończonym odcinku z dołączonymi skończenie wieloma sferami,
a taka przestrzeń jest jednospójna z van Kampena.
Można obliczyć homologię bez zwijania całej prostej:
mamy wierzchołki w punktach całkowitych, krawędzie między nimi
i po jednej $2$-komórce każdej sfery, dołączonej stale do jej wierzchołka.
Z homologii komórkowej
\[
 H_2(\widetilde X)
 \cong\bigoplus_{k\in\Z}\Z.
\]
Nakrycie nie zmienia
$\pi_2$, a
$\widetilde X$
jest jednospójna,
więc Hurewicz daje
\[
 \pi_2(X)\cong
 \pi_2(\widetilde X)
 \cong\bigoplus_{k\in\Z}\Z,
 \qquad
 H_2(X)\cong\Z.
\]
Każda podniesiona
sfera $S_k^2$
rzutuje się
na tę samą sferę
wewnątrz $X$.
W tych bazach
odwzorowanie Hurewicza
sumuje współczynniki:
$(a_k)_{k\in\Z}
\mapsto\sum_k a_k$.
Homologia traci
informację, którą
z podniesionych
sfer reprezentuje
dana klasa.
\end{przyklad}

\begin{figure}[!htbp]
\centering
\begin{tikzpicture}[>=Stealth,line cap=round,line join=round]
 \draw[manifoldgold!80!black,line width=2pt]
  (-4.5,0)--(4.5,0);
 \foreach \x/\lab in {-3/{-1},0/{0},3/{1}}{
  \shade[inner color=white,outer color=manifoldblue!37]
   (\x,.95) circle(.95);
  \draw[manifoldblue!75!black,thick]
   (\x,.95) circle(.95);
  \draw[manifoldblue!45,densely dashed,thin]
   ({\x+.95},.95)
   arc[start angle=0,end angle=180,x radius=.95,y radius=.23];
  \draw[manifoldblue!70!black,thin]
   ({\x-.95},.95)
   arc[start angle=180,end angle=360,x radius=.95,y radius=.23];
  \draw[manifoldblue!32,densely dashed,thin]
   (\x,1.9)..controls({\x-.28},1.6)
     and({\x-.28},.5)..(\x,0);
  \draw[manifoldblue!62!black,thin]
   (\x,1.9)..controls({\x+.28},1.6)
     and({\x+.28},.5)..(\x,0);
  \fill[manifoldred] (\x,0) circle(2.8pt);
  \node[below=5pt,font=\small] at (\x,0)
   {$k=\lab$};
 }
 \node[manifoldgold!75!black,font=\small]
  at (4.12,-.35) {$\R$};
 \node[font=\large] at (-4.28,1.02) {$\cdots$};
 \node[font=\large] at (4.28,1.02) {$\cdots$};
\end{tikzpicture}
\caption{Fragment nakrycia uniwersalnego
$S^1\vee S^2$: kopia sfery w każdym
punkcie całkowitym nad pętlą okręgu.
Każda sfera daje inny generator
$\pi_2(X)$, ale wszystkie mają
ten sam obraz w $H_2(X)$.}
\label{fig:nakrycie-bukiet-sfer}
\end{figure}

\begin{uwaga}[Który Whitehead?]
Twierdzenie~\ref{thm:whitehead}
dotyczy słabych równoważności
homotopijnych CW.
\emph{Torsja Whiteheada}
jest innym, późniejszym
niezmiennikiem używanym
przy prostej homotopii
i ogólniejszych
twierdzeniach o
kobordyzmach; sama
nazwa nie czyni jej
częścią tego twierdzenia.
\end{uwaga}

\section*{Dalsza droga}
\addcontentsline{toc}{section}{Dalsza droga}
Następny naturalny blok to
formy różniczkowe,
orientacja, całkowanie,
twierdzenie Stokesa
i kohomologia de Rhama.
Po nich rozwiniemy
geometrię riemannowską:
tensor metryczny, koneksję,
symbole Christoffela,
transport równoległy,
geodezyjne, odwzorowanie
wykładnicze, twierdzenie
Hopfa--Rinowa, krzywiznę
i tensory krzywizny,
pola Jacobiego,
twierdzenia porównawcze
oraz twierdzenie
Gaussa--Bonneta.
\emph{Dopiero po tym bloku}
przejdziemy do teorii
Morse'a, uchwytów,
twierdzenia Smale'a
i chirurgii.
Nie wprowadzamy
tych późniejszych
teorii w obecnym rozdziale.

\section*{Dalsza lektura}
\addcontentsline{toc}{section}{Dalsza lektura}
Allen Hatcher,
\emph{Algebraic Topology},
rozdział 4, §§4.1--4.2,
\url{https://pi.math.cornell.edu/~hatcher/AT/ATch4.pdf},
zawiera pełne rozwinięcie
grup względnych,
przybliżenia komórkowego,
twierdzeń Whiteheada
i Hurewicza.
\clearpage

\chapter{Fibracje i długi ciąg homotopii}
\label{ch:fibracje-v4}

\section{Dlaczego samo włókno nie wystarcza?}
Wiązka nakryciowa pozwala jednoznacznie podnosić drogi. W wiązce
okręgowej podniesień jest zazwyczaj wiele, ale można je wybierać
\emph{ciągle w rodzinach}. Właśnie tej cechy potrzeba, aby porównywać
homotopię włókna, przestrzeni całkowitej i podstawy. Wszystkie
przestrzenie w tym rozdziale rozpatrujemy z ich topologiami; dla
przestrzeni odwzorowań stosujemy topologię zwarto-otwartą.

\begin{definicja}[Własność podnoszenia homotopii]
Niech $p:E\to B$ będzie ciągłe, $Z$ przestrzenią topologiczną,
$f:Z\to E$ ciągłe oraz $H:Z\times I\to B$ homotopią spełniającą
$H(z,0)=p(f(z))$. Mówimy, że $p$ ma \emph{własność podnoszenia
homotopii względem $Z$}, jeżeli istnieje ciągłe
$\widetilde H:Z\times I\to E$ o
$\widetilde H(z,0)=f(z)$ oraz $p\widetilde H=H$.
Odwzorowanie $p$ nazywamy \emph{fibracją Hurewicza}, jeśli posiada tę
własność dla każdego $Z$, oraz \emph{fibracją Serre'a}, jeśli posiada ją
dla wszystkich dysków $Z=D^k$, $k\ge0$. Włóknem nad $b\in B$ nazywamy
$F_b=p^{-1}(b)$.
\end{definicja}

\begin{przyklad}
Rzut $B\times F\to B$ jest fibracją Hurewicza: przy
$f(z)=(f_B(z),f_F(z))$ można wziąć
$\widetilde H(z,t)=(H(z,t),f_F(z))$. Nakrycie także jest fibracją
Hurewicza: lokalna jednoznaczność podnoszenia dróg umożliwia
podnoszenie całych rodzin. Włókna tej samej fibracji nad różnymi punktami podstawy nie
muszą być homeomorficzne, lecz nad łukowo spójną podstawą
fibracji Hurewicza są równoważne homotopijnie.
\end{przyklad}

\begin{lemat}[Względne podnoszenie na komórce]
\label{lem:fib-wzgledne-v4}
Niech $p:E\to B$ będzie fibracją Serre'a, $k\ge0$,
$H:D^k\times I\to B$ i
$L=(D^k\times\{0\})\cup(\partial D^k\times I)$.
Każda mapa $h:L\to E$ spełniająca $ph=H|_L$ rozszerza się do mapy
$\widetilde H:D^k\times I\to E$ spełniającej
$p\widetilde H=H$. Dla pary CW $(X,A)$ analogicznie można podnieść
homotopię na $X\times I$, jeżeli podniesienie zadano na
$(X\times\{0\})\cup(A\times I)$.
\end{lemat}
\begin{proof}
Dla $k=0$ jest to własność podnoszenia dla punktu. Gdy $k\ge1$,
$D^k\times I$ jest kulą topologiczną, a $L$ jest jej brzegiem
z usuniętym wnętrzem górnej $k$-ściany. Istnieje homeomorfizm par
$(D^k\times I,L)\cong(D^k\times I,D^k\times\{0\})$:
obie podprzestrzenie są domkniętymi $k$-dyskami w brzegu kuli;
identyfikujemy brzeg ze sferą i przesuwamy odpowiedni dysk na dolną
półsferę, po czym przedłużamy homeomorfizm radialnie do kuli.
Po przeniesieniu diagramu mamy standardową własność podnoszenia.

Dla CW-pary przeprowadzamy indukcję po szkieletach $X$.
Wchodząc na nową komórkę $e^k$ używamy jej mapy
charakterystycznej $D^k\to X$; podniesienie jest już określone
na $\partial D^k\times I$ i $D^k\times\{0\}$, więc działa pierwszy
krok. Na wspólnych brzegach otrzymane mapy zgadzają się ze
starszym podniesieniem. Topologia słaba CW, a następnie fakt,
że $I$ jest lokalnie zwarty, zapewniają ciągłość mapy sklejonej
na $X\times I$. Nie zakładamy jednoznaczności podniesienia.
\end{proof}

\begin{twierdzenie}[Długi ciąg dokładny fibracji]
\label{thm:fib-les-v4}
Niech $p:E\to B$ będzie fibracją Serre'a, $e_0\in E$,
$b_0=p(e_0)$ oraz $F=p^{-1}(b_0)$. Wówczas istnieje naturalny
ciąg dokładny grup homotopii
\[
\cdots\to\pi_n(F,e_0)\xrightarrow{i_*}\pi_n(E,e_0)
\xrightarrow{p_*}\pi_n(B,b_0)\xrightarrow{\partial}\pi_{n-1}(F,e_0)
\to\cdots.
\]
Na końcu otrzymujemy $\pi_1(B,b_0)\to\pi_0(F)\to\pi_0(E)\to\pi_0(B)$;
tu dokładność rozumie się dla zbiorów z wyróżnionym punktem.
Dla $n=2$ grupa po prawej może być nieprzemienna, choć obraz
homomorfizmu $\partial$ jest wtedy przemienny.
\end{twierdzenie}
\begin{proof}
Z długiego ciągu pary $(E,F)$ wystarczy otrzymać naturalne
izomorfizmy $\pi_n(E,F,e_0)\cong\pi_n(B,b_0)$ dla $n\ge2$
oraz odpowiednią bijekcję w stopniu $1$.
Użyjmy kostkowego opisu grup względnych: wyróżniona ściana
$I^{n-1}\times\{0\}$ jest odwzorowywana do $F$, a wszystkie
pozostałe ściany do $e_0$. Aby podnieść mapę
$(I^n,\partial I^n)\to(B,b_0)$, traktujemy jej ostatnią
współrzędną jako czas; na jednej skrajnej ścianie i na ścianach
bocznych ustalamy $e_0$. Lemat~\ref{lem:fib-wzgledne-v4}
rozszerza podniesienie. Otrzymujemy surjektywność.

Jeśli dwie mapy względne $u,v$ mają homotopijne obrazy w $B$,
podnosimy najpierw homotopię ich obrazów, przytrzymując ściany
z wartością $e_0$. Otrzymujemy $u\simeq u'$ z $pu'=pv$.
Następnie podnosimy stałą w czasie homotopię $pv$, przy
z góry ustalonych wartościach $u'$ i $v$ na obu podstawach
kostki $I^n\times I$ i na pozostałych ścianach bocznych.
Ten ustalony fragment brzegu jest dyskiem $n$-wymiarowym:
komplementem w brzegu całej kostki jest wnętrze
$\bigl(I^{n-1}\times\{0\}\bigr)\times I$. Ten sam argument
z homeomorfizmem par, co w lemacie, daje homotopię
$u'\simeq v$ przez mapy względne. Jest to iniektywność.
Zgodność z działaniem grupowym wynika z zachowania
sklejania kostek przez $p$. Wstawienie uzyskanych identyfikacji
do długiego ciągu pary daje tezę.
\end{proof}

\begin{uwaga}[Jak obliczać homomorfizm łączący?]
Dla $a:(D^n,S^{n-1})\to(B,b_0)$ podnosimy $a$ po
ustaleniu $e_0$ na odpowiedniej części brzegu. Restrykcja
podniesienia do pozostałej ściany leży w $F$ i reprezentuje
$\partial[a]\in\pi_{n-1}(F,e_0)$. Wybór orientacji kostki
może zmienić znak $\partial$, dlatego należy zachować
jedną konwencję w całym rachunku.
\end{uwaga}

\section{Transport wzdłuż drogi i monodromia}
\begin{stwierdzenie}
Niech $p:E\to B$ będzie fibracją Hurewicza, a $\gamma:I\to B$
drogą od $b_0$ do $b_1$. Istnieje transport ciągły
$T_\gamma:F_{b_0}\to F_{b_1}$ będący równoważnością
homotopijną. Jego klasa homotopii zależy tylko od klasy
$\gamma$ względem końców.
\end{stwierdzenie}
\begin{proof}
Stosujemy własność podnoszenia do $Z=F_{b_0}$,
$f:F_{b_0}\hookrightarrow E$ i $H(x,t)=\gamma(t)$, a następnie
kładziemy $T_\gamma(x)=\widetilde H(x,1)$.
Podobnie wybieramy transport wzdłuż drogi odwróconej.
Aby wykazać, że są to odwrotności z dokładnością do homotopii,
zauważmy, że $\gamma*\overline\gamma$ jest homotopijna względem
końców drodze stałej. Złączone podniesienia dają transport wzdłuż
złożenia; własność podnoszenia homotopii rodzin dróg zapewnia
homotopię z transportem wzdłuż drogi stałej. Ten ostatni można
wybrać jako identyczność. Analogicznie dowodzimy
$T_\gamma T_{\overline\gamma}\simeq\id$.
\end{proof}

\begin{przyklad}[Butelka Kleina]
Dla $K=(S^1\times I)/((z,1)\sim(\overline z,0))$
rzut na $I/(0\sim1)\cong S^1$ jest wiązką o włóknie $S^1$.
Monodromia działa na $\pi_1(S^1)=\Z$ przez $m\mapsto-m$.
Z tego wynika prezentacja
$\pi_1(K)\cong\langle a,b\mid aba^{-1}=b^{-1}\rangle$.
Grupa nie jest iloczynem prostym dwóch kopii $\Z$.
\end{przyklad}

\section{Fibracja przestrzeni dróg}
Dla punktowanej przestrzeni $(B,b_0)$ określmy
$PB=\{\gamma:I\to B:\gamma(0)=b_0\}$ oraz
$\Omega B=\{\gamma\in PB:\gamma(1)=b_0\}$.
\begin{twierdzenie}
Ewaluacja $\mathrm{ev}_1:PB\to B$, $\gamma\mapsto\gamma(1)$,
jest fibracją Hurewicza z włóknem $\Omega B$, a $PB$ jest
ściągalna. Dla $n\ge2$ zachodzi
$\pi_n(B,b_0)\cong\pi_{n-1}(\Omega B,c_{b_0})$;
ponadto $\pi_1(B,b_0)$ jest w naturalnej bijekcji z
$\pi_0(\Omega B)$.
\end{twierdzenie}
\begin{proof}
Niech $f:Z\to PB$ oraz $H:Z\times I\to B$ spełniają
$H(z,0)=f(z)(1)$. Wybieramy podniesienie $\widetilde H(z,t)$
łącząc najpierw $f(z)$, a następnie przebiegając odcinek
$H(z,s)$ dla $0\le s\le t$. Żeby zachować warunek
$\widetilde H(z,0)=f(z)$, stosujemy zmienny podział parametru:
\[
\widetilde H(z,t)(u)=
\begin{cases}
f(z)\!\left(\dfrac{u}{1-t/2}\right),&0\le u\le1-t/2,\\[4pt]
H\!\left(z,2u-2+t\right),&1-t/2\le u\le1.
\end{cases}
\]
Przy $t=0$ drugi przedział jest jednopunktowy; obie gałęzie
mają zgodną wartość na styku, a koniec drogi to $H(z,t)$.
Lemat o sklejaniu i własność wykładnicza topologii
zwarto-otwartej dla zwartego, lokalnie zwartego $I$ dają
ciągłość całej rodziny. Homotopia
$R_s(\gamma)(u)=\gamma((1-s)u)$ ściąga $PB$ do drogi stałej.
Teraz stosujemy Twierdzenie~\ref{thm:fib-les-v4}.
\end{proof}

\section{Fibracja Hopfa: poprawny wzór i obliczenie}
\begin{twierdzenie}
\label{thm:hopf-v4}
Odwzorowanie
$h:S^3\to S^2\subset\C\times\R$ dane wzorem
\[
h(z,w)=(2z\overline w,\ |z|^2-|w|^2),\qquad
|z|^2+|w|^2=1,
\]
jest lokalnie trywialną wiązką o włóknie $S^1$.
W konsekwencji $\pi_3(S^2)\cong\Z$, a $[h]$ jest generatorem.
\end{twierdzenie}
\begin{proof}
Równość $4|z\overline w|^2+(|z|^2-|w|^2)^2=1$
pokazuje, że obraz należy do $S^2$.
Dla $\lambda\in S^1$ mamy
$h(\lambda z,\lambda w)=h(z,w)$, ponieważ
$\lambda\overline\lambda=1$.
Dlatego $h$ przechodzi na $\C P^1$.
Dla $z\ne0$ współrzędna $u=w/z$ jednoznacznie opisuje
prostą zespoloną, której jednostkowy przedstawiciel to
$s(u)=(1,u)/\sqrt{1+|u|^2}$. Mapa
$(u,\lambda)\mapsto\lambda s(u)$ jest homeomorfizmem
lokalnej trywializacji. Analogicznie postępujemy dla $w\ne0$.
W tych współrzędnych $h$ daje standardową identyfikację
$\C P^1\cong S^2$, więc jest wiązką $S^1\to S^3\to S^2$.
Z ciągu dokładnego oraz $\pi_k(S^1)=0$ dla $k\ge2$ dostajemy
$\pi_3(S^3)\xrightarrow{\cong}\pi_3(S^2)$.
Ponieważ $S^3$ jest $2$-spójna, twierdzenie Hurewicza daje
$\pi_3(S^3)\cong H_3(S^3;\Z)=\Z$; obraz klasy identyczności
jest klasą $h$.
\end{proof}
\begin{uwaga}
Zapis $2zw$ bez sprzężenia jest \emph{błędny}:
po pomnożeniu obu współrzędnych przez $\lambda\in S^1$
pierwsza współrzędna zmieniłaby się przez czynnik $\lambda^2$.
Włókna nie byłyby zatem orbitami działania $S^1$.
\end{uwaga}

\section{Zastąpienie odwzorowania fibracją}
\begin{twierdzenie}
Niech $f:X\to B$ będzie ciągłe. Zdefiniujmy
$E_f=\{(x,\gamma)\in X\times B^I:\gamma(0)=f(x)\}$,
$q(x,\gamma)=\gamma(1)$ i $j(x)=(x,c_{f(x)})$.
Wtedy $q:E_f\to B$ jest fibracją Hurewicza,
$j:X\to E_f$ jest równoważnością homotopijną oraz $qj=f$.
\end{twierdzenie}
\begin{proof}
Podnosząc homotopię w $B$, dopisujemy do drugiej współrzędnej
$(x,\gamma)$ odpowiedni odcinek drogi; zmienną parametryzację
wybieramy jak w dowodzie dla $PB$. Pierwsza współrzędna jest
niezmienna, więc $\gamma(0)=f(x)$ pozostaje prawdziwe.
Rzut $r:E_f\to X$ jest lewą odwrotnością $j$,
a $K_s(x,\gamma)=(x,\gamma((1-s)\,\cdot))$ jest
homotopią od $\id_{E_f}$ do $jr$.
\end{proof}

Włóknem homotopijnym $f$ nad $b_0$ nazywamy
$\operatorname{hofib}_{b_0}(f)=q^{-1}(b_0)$.
Dla każdej mapy można więc zastosować ciąg dokładny
\emph{włókno homotopijne -- dziedzina -- przeciwdziedzina}.

\section{Granice metody i zapowiedź ciągu Serre'a}
Długi ciąg dokładny mierzy grupy homotopii, ale nie wyznacza
sam z siebie pierścienia kohomologii $E$. Dla fibracji Serre'a
$F\to E\to B$ o podstawie łukowo spójnej i przestrzeniach
odpowiedniego typu CW ciąg spektralny Serre'a w homologii
ma drugą stronę
$E^2_{p,q}=H_p(B;\mathcal H_q(F;\Z))$,
gdzie $\mathcal H_q(F;\Z)$ jest systemem lokalnym wyznaczonym
przez monodromię. Jego zbieżność oznacza filtrację $H_{p+q}(E)$,
a nie bezpośrednie utożsamienie z sumą wyrazów strony drugiej.
Gdy monodromia jest trywialna, lokalne współczynniki można
zastąpić stałymi. Nie stosujemy tu tego ciągu bez osobnego
dowodu jego konstrukcji i warunków zbieżności.
\clearpage

\chapter{Kohomologia singularna i iloczyn kubkowy}
\label{ch:kohomologia-singularna}

Homologia przypisuje przestrzeni kompleks łańcuchowy i bierze jądra
modulo obrazy różniczek. W kohomologii stosujemy do łańcuchów
kontrawariantny funktor homomorfizmów do grupy współczynników.
W rezultacie uzyskujemy nie tylko grupy, lecz także naturalne mnożenie.
Ten rozdział łączy konstrukcje homologiczne z dualnością Poincarégo
i porównaniem de Rhama omawianym dalej.

\section{Kołańcuchy, koróżniczka i funktorialność}

\begin{definicja}
Niech $X$ będzie przestrzenią topologiczną, $G$ grupą abelową i
$(C_n(X),\partial_n)$ singularnym kompleksem łańcuchowym z całkowitymi
współczynnikami. Dla $n\geq0$ określamy grupę \emph{kołańcuchów}
$C^n(X;G)=\operatorname{Hom}_{\Z}(C_n(X),G)$, a dla $n<0$ przyjmujemy
$C^n(X;G)=0$. Homomorfizm
$\delta^n:C^n(X;G)\to C^{n+1}(X;G)$ jest dany wzorem
$(\delta^n\varphi)(c)=\varphi(\partial_{n+1}c)$.
Grupę kohomologii stopnia $n$ definiujemy przez
\[
H^n(X;G)=\ker\delta^n/\operatorname{im}\delta^{n-1}.
\]
Elementy $\ker\delta^n$ są kocyklami, a elementy
$\operatorname{im}\delta^{n-1}$ kobrzegami.
\end{definicja}

\begin{stwierdzenie}
Dla każdej przestrzeni topologicznej $X$ i grupy abelowej $G$
zachodzi $\delta^{n+1}\circ\delta^n=0$ dla wszystkich $n\geq0$.
\end{stwierdzenie}
\begin{proof}
Dla $\varphi\in C^n(X;G)$ i $c\in C_{n+2}(X)$ mamy
$(\delta^{n+1}\delta^n\varphi)(c)
=\varphi(\partial_{n+1}\partial_{n+2}c)=0$,
ponieważ $\partial^2=0$ w kompleksie łańcuchów singularnych.
Zatem każdy kobrzeg jest kocyklem i iloraz w definicji ma sens.
\end{proof}

\begin{definicja}
Niech $f:X\to Y$ będzie ciągłe. Dla $\varphi\in C^n(Y;G)$
określamy $(f^*\varphi)(c)=\varphi(f_\#c)$ dla $c\in C_n(X)$,
gdzie $f_\#$ jest odwzorowaniem łańcuchowym indukowanym przez $f$.
\end{definicja}

\begin{twierdzenie}[Kontrawariantność i niezmienniczość homotopijna]
Każde ciągłe $f:X\to Y$ indukuje homomorfizmy
$f^*:H^n(Y;G)\to H^n(X;G)$, dla których
$(g\circ f)^*=f^*\circ g^*$, a $(\id_X)^*=\id$.
Jeżeli $f\simeq g$, to $f^*=g^*$.
\end{twierdzenie}
\begin{proof}
Tożsamość $\partial f_\#=f_\#\partial$ po przejściu do
$\operatorname{Hom}(-,G)$ daje $\delta f^*=f^*\delta$.
Stąd $f^*$ odwzorowuje kocykle w kocykle i kobrzegi w kobrzegi.
Tożsamości związane ze składaniem wynikają już na poziomie łańcuchów.
Jeśli $f\simeq g$, w dowodzie niezmienniczości homotopijnej
homologii skonstruowano operator pryzmowy $P$ spełniający
$g_\#-f_\#=\partial P+P\partial$.
Po złożeniu z kołańcuchem $\varphi$ otrzymujemy
$g^*\varphi-f^*\varphi=\delta(P^*\varphi)+P^*(\delta\varphi)$,
gdzie $P^*$ obniża stopień o jeden. Dla kocyklu drugi składnik znika,
a pierwszy jest kobrzegiem. Indukowane odwzorowania są równe.
\end{proof}

\begin{przyklad}
Dla niepustej łukowo spójnej przestrzeni $X$ mamy $H^0(X;G)\cong G$:
kocykl stopnia zero jest funkcją na punktach, która przyjmuje tę samą
wartość na końcach każdej drogi. Nie ma nietrywialnych kobrzegów
stopnia zero. Dla przestrzeni o $k$ składowych łukowych otrzymujemy
$H^0(X;G)\cong G^k$; przy nieskończenie wielu składowych
łukowych pojawia się iloczyn, nie suma prosta.
\end{przyklad}

\section{Kohomologia względna i ciąg dokładny}

\begin{definicja}
Dla pary przestrzeni $A\subseteq X$ przyjmujemy
$C^n(X,A;G)=\{\varphi\in C^n(X;G):\varphi|_{C_n(A)}=0\}$.
Ponieważ brzeg łańcucha z $A$ pozostaje w $A$, jest to podkompleks
kołańcuchowy. Definiujemy $H^n(X,A;G)=H^n(C^*(X,A;G))$.
\end{definicja}

\begin{twierdzenie}[Długi ciąg dokładny pary]
Dla każdej pary przestrzeni topologicznych $A\subseteq X$ oraz
każdej grupy abelowej $G$ istnieje naturalny długi ciąg dokładny
\[
\cdots\to H^n(X,A;G)\to H^n(X;G)\to H^n(A;G)
\xrightarrow{\,\partial^*\,}H^{n+1}(X,A;G)\to\cdots.
\]
\end{twierdzenie}
\begin{proof}
Łańcuchy singularne są wolnymi grupami abelowymi na singularnych
sympleksach. Każdy sympleks w $A$ jest też sympleksem w $X$, więc
$C_n(A)$ jest wolnym składnikiem prostym $C_n(X)$.
Wobec tego ograniczenie kołańcuchów z $X$ do $A$ jest surjekcją
$C^n(X;G)\to C^n(A;G)$, a jej jądrem jest $C^n(X,A;G)$.
Otrzymujemy krótki ciąg dokładny kompleksów kołańcuchowych.
Konstrukcja homomorfizmu łączącego, przedstawiona w rozdziale
o algebrze homologicznej, daje wskazany długi ciąg. Naturalność
wynika z naturalności restrykcji i konstrukcji homomorfizmu łączącego.
\end{proof}

\section{Iloczyn kubkowy: skąd bierze się mnożenie klas}

\begin{definicja}[Iloczyn kubkowy Alexandra--Whitneya]
Niech $R$ będzie przemiennym pierścieniem z jedynką,
$\alpha\in C^p(X;R)$ i $\beta\in C^q(X;R)$, gdzie $p,q\geq0$.
Dla singularnego sympleksu $\sigma:\Delta^{p+q}\to X$ definiujemy
\[
(\alpha\smile\beta)(\sigma)=
\alpha\bigl(\sigma|_{[v_0,\ldots,v_p]}\bigr)\,
\beta\bigl(\sigma|_{[v_p,\ldots,v_{p+q}]}\bigr).
\]
Restrykcje do ścian są rozumiane przez standardowe liniowe
parametryzacje sympleksów o wskazanych wierzchołkach.
Definicję przedłużamy liniowo na wszystkie łańcuchy.
\end{definicja}

\begin{stwierdzenie}[Reguła Leibniza]
Jeżeli $\alpha\in C^p(X;R)$ i $\beta\in C^q(X;R)$, to
$\delta(\alpha\smile\beta)
=(\delta\alpha)\smile\beta+(-1)^p\alpha\smile(\delta\beta)$.
W konsekwencji iloczyn kocykli jest kocyklem, a gdy co najmniej
jeden czynnik jest kobrzegiem, iloczyn jest kobrzegiem.
\end{stwierdzenie}
\begin{proof}
Obliczamy lewą stronę na $\sigma:\Delta^{p+q+1}\to X$,
rozwijając $\partial\sigma$ jako naprzemienną sumę ścian.
Dla usunięcia wierzchołka o indeksie $i<p$ pierwsza restrykcja
ma ścianę z pominiętym wierzchołkiem, a druga jest niezmieniona.
Dla $i>p+1$ sytuacja jest odwrotna i znak $(-1)^p$ pojawia się
z przesunięcia numeracji ścian drugiej restrykcji.
Dwa wyrazy ze środkowych ścian $i=p$ i $i=p+1$ są identyczne
po ograniczeniu do wspólnego wierzchołka i mają przeciwne znaki;
redukują się. Pozostałe wyrazy to dokładnie rozwinięcia
$(\delta\alpha)\smile\beta$ oraz $(-1)^p\alpha\smile(\delta\beta)$.
Jeżeli $\delta\alpha=0$ i $\beta=\delta\gamma$, to
$\alpha\smile\beta=(-1)^p\delta(\alpha\smile\gamma)$;
przypadek kobrzegu w pierwszym czynniku jest analogiczny.
\end{proof}

\begin{twierdzenie}
Dla każdej przestrzeni $X$ i przemiennego pierścienia $R$ iloczyn
kubkowy indukuje dwuliniowe odwzorowania
$H^p(X;R)\times H^q(X;R)\to H^{p+q}(X;R)$.
Z tym mnożeniem $H^*(X;R)=\bigoplus_{n\geq0}H^n(X;R)$ jest
łącznym, unitalnym, \emph{stopniowo przemiennym} pierścieniem,
tzn. dla jednorodnych $a,b$ zachodzi
$a\smile b=(-1)^{\deg(a)\deg(b)}b\smile a$.
Ponadto $f^*(a\smile b)=f^*a\smile f^*b$ dla ciągłego $f$.
\end{twierdzenie}
\begin{proof}[Argument z użyciem standardowej homotopii diagonalnych]
Poprawność definicji na kohomologii wynika z reguły Leibniza.
Łączność wynika bezpośrednio z trzykrotnego rozcięcia
uporządkowanego zbioru wierzchołków sympleksu.
Elementem jednostkowym jest klasa stałego kocyklu
$1\in C^0(X;R)$, przyjmującego wartość $1_R$ na każdym
wierzchołku. Naturalność wynika z faktu, że złożenie z $f$
komutuje z ograniczeniem sympleksu do ściany.
Stopniowa przemienność nie zachodzi na ogół na poziomie
kołańcuchów: jej dowód wymaga homotopii między
odwzorowaniem diagonalnym Alexandra--Whitneya a jego
transpozycją. Istnienie takiej homotopii wynika z kołańcuchowego lematu
o porównaniu diagonali (jego niezależny dowód wymaga jawnej
konstrukcji homotopii pryzmowej). Daje ona
$\alpha\smile\beta-(-1)^{pq}\beta\smile\alpha$ jako kobrzeg,
gdy $\alpha$ i $\beta$ są kocyklami.
\end{proof}

\begin{przyklad}[Torus]
Niech $a,b\in H^1(T^2;\Z)$ będą cofnięciami generatora
$H^1(S^1;\Z)\cong\Z$ przez obie projekcje $T^2\to S^1$.
Wówczas $a\smile a=b\smile b=0$, natomiast $a\smile b$
jest generatorem $H^2(T^2;\Z)\cong\Z$ przy odpowiednim
wyborze orientacji. Pierścień kohomologii torusa jest więc
algebrą zewnętrzną na generatorach $a,b$ stopnia jeden.
To więcej informacji niż same grupy $H^0\cong\Z$,
$H^1\cong\Z^2$, $H^2\cong\Z$.
\end{przyklad}

\section{Współczynniki i twierdzenie o współczynnikach uniwersalnych}

\begin{uwaga}[Co oznacza $\operatorname{Ext}^1$?]
Dla grupy abelowej $A$ wybieramy dokładny ciąg
$0\to F_1\xrightarrow{u} F_0\to A\to0$ z wolnymi grupami
abelowymi $F_0,F_1$ i definiujemy
$\operatorname{Ext}^1_{\Z}(A,G)=
\operatorname{coker}\bigl(\operatorname{Hom}(F_0,G)
\xrightarrow{u^*}\operatorname{Hom}(F_1,G)\bigr)$.
Wynik nie zależy, z dokładnością do kanonicznego izomorfizmu,
od wyboru wolnej prezentacji. Na przykład
$\operatorname{Ext}^1_{\Z}(\Z/m\Z,G)\cong G/mG$,
co wynika z prezentacji $0\to\Z\xrightarrow{\cdot m}\Z
\to\Z/m\Z\to0$.
\end{uwaga}

\begin{twierdzenie}[Współczynniki uniwersalne dla kohomologii]
Niech $X$ będzie przestrzenią topologiczną, $G$ grupą abelową,
a $n\geq0$ liczbą całkowitą. Istnieje naturalny krótki ciąg
dokładny grup abelowych
\[
0\longrightarrow\operatorname{Ext}^1_{\Z}(H_{n-1}(X;\Z),G)
\longrightarrow H^n(X;G)
\longrightarrow\operatorname{Hom}_{\Z}(H_n(X;\Z),G)
\longrightarrow0,
\]
przy czym $H_{-1}(X;\Z)=0$. Ciąg rozszczepia się jako ciąg
grup abelowych, lecz nie ma na ogół naturalnego wyboru rozszczepienia.
\end{twierdzenie}



\begin{proof}
Piszemy $C_n=C_n(X;\Z)$,
$Z_n=\ker\partial_n$, $B_n=\operatorname{im}\partial_{n+1}$ oraz
$H_n=Z_n/B_n$. Wszystkie $C_n,Z_n,B_n$ są wolnymi grupami
abelowymi: $C_n$ z definicji, a pozostałe jako podgrupy
grup wolnych. Ponieważ $B_{n-1}$ jest wolna, krótki ciąg
$0\to Z_n\to C_n\xrightarrow{\partial_n}B_{n-1}\to0$
rozszczepia się. Można więc wybrać
$C_n\cong Z_n\oplus B_{n-1}$, a w tych współrzędnych
różniczka jest projekcją na drugi składnik, po której
następuje włączenie do $C_{n-1}$.

Dla kocyklu $\varphi\in\operatorname{Hom}(C_n,G)$ restrykcja
$\varphi|_{Z_n}$ znika na $B_n$, a więc definiuje
homomorfizm $H_n\to G$. Powstaje naturalna mapa
$\rho:H^n(X;G)\to\operatorname{Hom}(H_n,G)$,
$\rho[\varphi]([z])=\varphi(z)$. Jest dobrze określona:
kobrzeg znika na cyklach, a kocykl na brzegach.
Mapa $\rho$ jest surjektywna: dowolny homomorfizm
$H_n\to G$ cofamy do $Z_n$ i przedłużamy do
$C_n=Z_n\oplus B_{n-1}$ przez zero na drugim składniku.

Przedstawiciel elementu $\ker\rho$ może być wybrany
jako kocykl znikający na $Z_n$. Rzeczywiście, warunek
$\rho[\varphi]=0$ oznacza dokładnie $\varphi|_{Z_n}=0$.
Takie kocykle są równoważne homomorfizmom
$B_{n-1}=C_n/Z_n\to G$. Kobrzegi odpowiadają tym
homomorfizmom $B_{n-1}\to G$, które są restrykcjami
homomorfizmów $C_{n-1}\to G$. Ponieważ
$C_{n-1}\cong Z_{n-1}\oplus B_{n-2}$, otrzymujemy
\[
\ker\rho\cong
\operatorname{coker}\!\left(
\operatorname{Hom}(Z_{n-1},G)\longrightarrow
\operatorname{Hom}(B_{n-1},G)\right).
\]
Dokładny ciąg
$0\to B_{n-1}\to Z_{n-1}\to H_{n-1}\to0$
jest wolną prezentacją $H_{n-1}$, więc ostatni
kokernel to $\operatorname{Ext}^1_{\Z}(H_{n-1},G)$.
Dostajemy żądany krótki ciąg dokładny.

Wybrane rozszczepienia $C_n=Z_n\oplus B_{n-1}$
rozkładają grupę kocykli na część odpowiadającą
$\operatorname{Hom}(H_n,G)$ i część odpowiadającą
$\operatorname{Hom}(B_{n-1},G)$; kobrzegi leżą tylko
w drugiej. To zapewnia rozszczepienie otrzymanego
ciągu, ale wybór dopełnień $Z_n$ w $C_n$ nie jest
naturalny względem odwzorowań łańcuchowych.
\end{proof}

\begin{przyklad}[Torsja widoczna dopiero w kohomologii]
Dla $\R P^2$ mamy $H_1(\R P^2;\Z)\cong\Z/2\Z$
i $H_2(\R P^2;\Z)=0$. Z twierdzenia wynika
$H^1(\R P^2;\Z)=0$ oraz
$H^2(\R P^2;\Z)\cong\Z/2\Z$.
Nie wolno więc utożsamiać kohomologii integralnej
z dualem homologii integralnej.
\end{przyklad}

\section{Iloczyny przestrzeni: Eilenberg--Zilber i Künneth}
\label{sec:ez-kunneth-v4}

Samo znajdowanie grup $H_p(X)$ i $H_q(Y)$ nie wystarcza do
obliczenia $H_n(X\times Y)$: przy współczynnikach całkowitych
pojawia się poprawka torsyjna. Zaczniemy więc od \emph{kompleksów},
a nie od grup homologii.

\begin{definicja}[Iloczyn tensorowy kompleksów i funktor $\operatorname{Tor}$]
Dla nieujemnie stopniowanych kompleksów $C,D$ nad $\Z$ niech
$(C\otimes D)_n=\bigoplus_{p+q=n}C_p\otimes_\Z D_q$ i
$\partial(c\otimes d)=\partial c\otimes d+(-1)^p c\otimes\partial d$
przy $c\in C_p$. Bezpośredni rachunek daje $\partial^2=0$.
Jeśli $A,B$ są grupami abelowymi i
$0\to F_1\xrightarrow{u}F_0\to A\to0$ jest wolną prezentacją,
to $\operatorname{Tor}^{\Z}_1(A,B)$ oznacza
$\ker(u\otimes\id_B)$. Niezależność od prezentacji wynika
z porównywania jej z innymi prezentacjami poprzez podniesienia
homomorfizmów wolnych grup.
\end{definicja}

\begin{twierdzenie}[Eilenberg--Zilber, łącznie z homotopiami]
\label{thm:ez-v4}
Dla dowolnych przestrzeni topologicznych $X,Y$ istnieją naturalne
odwzorowania łańcuchowe
$\mathrm{AW}:C_*(X\times Y;\Z)\to C_*(X;\Z)\otimes C_*(Y;\Z)$
oraz $\mathrm{sh}$ w przeciwną stronę, takie że
$\mathrm{sh}\,\mathrm{AW}\simeq\id$ i
$\mathrm{AW}\,\mathrm{sh}\simeq\id$ jako mapy łańcuchowe.
\end{twierdzenie}
\begin{proof}
\emph{Krok 1: odwzorowanie Alexandra--Whitneya.}
Dla $\sigma=(\sigma_X,\sigma_Y):\Delta^n\to X\times Y$
przyjmujemy
\[
\mathrm{AW}(\sigma)=\sum_{i=0}^n
\sigma_X|[0,\ldots,i]\otimes\sigma_Y|[i,\ldots,n].
\]
Ograniczenia mają standardową rosnącą kolejność wierzchołków.
Sprawdźmy różniczkę. W $\partial\mathrm{AW}(\sigma)$
usunięcie wierzchołka przed indeksem rozcięcia odpowiada
usunięciu tej samej ściany z $\mathrm{AW}(\partial\sigma)$;
podobnie po rozcięciu, gdzie znak tensoru daje $(-1)^i$.
Wyrazy związane z usunięciem samego punktu rozcięcia
występują dwa razy z przeciwnymi znakami. Pozostają zatem
dokładnie wyrazy $\mathrm{AW}(\partial\sigma)$.

\emph{Krok 2: przeplot.}
Niech $\sigma:\Delta^p\to X$, $\tau:\Delta^q\to Y$.
$(p,q)$-przeplot to słowo długości $p+q$ o $p$ krokach
poziomych i $q$ pionowych, dające wierzchołki drogi kratowej
$(a_j,b_j)$ od $(0,0)$ do $(p,q)$. Niech
$\lambda_s:\Delta^{p+q}\to\Delta^p\times\Delta^q$
będzie mapą afiniczną na te wierzchołki. Oznaczmy przez
$\varepsilon_s$ znak permutacji przenoszącej uporządkowane
$p$ kroki poziome przed $q$ pionowymi do słowa $s$.
Definiujemy
$\mathrm{sh}(\sigma\otimes\tau)
=\sum_s\varepsilon_s(\sigma\times\tau)\lambda_s$.
Wszystkie wewnętrzne ściany triangulacji produktu
$\Delta^p\times\Delta^q$ pojawiają się dwukrotnie ze znakami
przeciwnymi: odpowiadają przeplotom różniącym się pojedynczym
zamienieniem sąsiednich kroków. Zewnętrzne ściany dają
$\mathrm{sh}(\partial\sigma\otimes\tau)$ i
$(-1)^p\mathrm{sh}(\sigma\otimes\partial\tau)$.
To wyjaśnia zarówno wystąpienie znaku, jak i tożsamość
$\partial\mathrm{sh}=\mathrm{sh}\partial$.

\emph{Krok 3: lemat o rozszerzaniu homotopii na modelach.}
Załóżmy, że $F,G:C\to D$ są mapami łańcuchowymi i $C_n$
jest wolna na wskazanych generatorach-modelach. Mamy skonstruować
$K_n:C_n\to D_{n+1}$ tak, aby $\partial K+K\partial=F-G$.
Przy konstrukcji indukcyjnej na generatorze $u$ stopnia $n$
wyraz $z=(F-G)u-K(\partial u)$ jest cyklem, ponieważ
$\partial z=(F-G)(\partial u)-(F-G)(\partial u)+K\partial^2u=0$.
Jeśli modelowa przestrzeń wartości ma zerową homologię
zredukowaną w stopniu $n$, można wybrać $K(u)$ z
$\partial K(u)=z$. Przy $n=0$ wymagamy dodatkowo zerowej
augmentacji $z$. Budując $K$ w modelu, a następnie
przenosząc go przez mapy $\sigma_\#$, otrzymujemy naturalność.

Uzasadnijmy wspomniane znikanie: dla wypukłego $V$ wybieramy
$v\in V$ i stożkujemy sympleksy ku $v$. Operator $S$ spełnia
$\partial S+S\partial=\id-\eta\epsilon$, gdzie
$\epsilon:C_0(V)\to\Z$ jest augmentacją i
$\eta:\Z\to C_0(V)$ wysyła $1$ na stały sympleks w $v$.
W dodatnich stopniach każdy cykl jest więc brzegiem.
Analogicznie tensor dwóch takich kompleksów ma kontrakcję:
na elementach jednorodnych działa
$S_X\otimes\id+(\eta\epsilon)_X\otimes S_Y$,
ze standardową konwencją znaków dla map stopnia $1$.
Po podstawieniu otrzymujemy
$\partial S_{\otimes}+S_{\otimes}\partial
=\id-(\eta\epsilon)_X\otimes(\eta\epsilon)_Y$.

\emph{Krok 4: $\mathrm{sh}\,\mathrm{AW}\simeq\id$.}
Dla uniwersalnego sympleksu w iloczynie wybieramy
$\Delta^n\xrightarrow{\Delta_n}\Delta^n\times\Delta^n$,
$t\mapsto(t,t)$. Jego obrazy w dowolnym $X\times Y$
powstają przez mapy $\sigma_X\times\sigma_Y$.
W stopniu zero obie mapy są równe. Jeśli $K$ znamy na
ścianach $\Delta_n$, to
$z=(\mathrm{sh}\,\mathrm{AW}-\id)(\Delta_n)-K(\partial\Delta_n)$
jest cyklem w $C_n(\Delta^n\times\Delta^n)$. Model jest
wypukły, więc cykl ma wypełnienie $K(\Delta_n)$.
Dla $n=0$ augmentacja zniknęła już na starcie.
Przenosząc wybrane wypełnienia przez
$(\sigma_X\times\sigma_Y)_\#$, otrzymujemy żądaną
homotopię naturalną. Istotne jest, że pracujemy z
kompleksem \emph{produktu} $\Delta^n\times\Delta^n$.

\emph{Krok 5: $\mathrm{AW}\,\mathrm{sh}\simeq\id$.}
Tym razem modele dla generatorów $\sigma\otimes\tau$ to
$\id_{\Delta^p}\otimes\id_{\Delta^q}$, a grupy docelowe
modelowej indukcji to
$C_*(\Delta^p)\otimes C_*(\Delta^q)$.
Dla $p+q=0$ mapy są identyczne. W wyższym stopniu
$z=(\mathrm{AW}\,\mathrm{sh}-\id)
(\id_{\Delta^p}\otimes\id_{\Delta^q})-K\partial
(\id_{\Delta^p}\otimes\id_{\Delta^q})$ jest cyklem.
Kontrakcja tensorowa z kroku 3 wypełnia ten cykl.
Przeniesienie przez $\sigma_\#\otimes\tau_\#$ daje $K$
na wszystkich generatorach. Indukcja dowodzi tożsamości
$\partial K+K\partial=\mathrm{AW}\,\mathrm{sh}-\id$.
\end{proof}

\begin{twierdzenie}[Künneth dla homologii singularnej o współczynnikach całkowitych]
\label{thm:kunneth-v4}
Niech $X,Y$ będą dowolnymi przestrzeniami topologicznymi, a $n\geq0$.
Istnieje naturalny względem obu przestrzeni krótki ciąg dokładny
\[
\begin{gathered}
0\longrightarrow
\bigoplus_{p+q=n}H_p(X;\Z)\otimes_\Z H_q(Y;\Z)
\xrightarrow{\ \times\ }H_n(X\times Y;\Z)\\
\xrightarrow{\ \kappa\ }
\bigoplus_{p+q=n-1}\operatorname{Tor}^{\Z}_1(H_p(X;\Z),H_q(Y;\Z))
\longrightarrow0.
\end{gathered}
\]
Ciąg rozszczepia się jako ciąg grup abelowych, ale na ogół nie istnieje
rozszczepienie naturalne. Nie zakładamy, że którakolwiek grupa homologii
jest skończenie generowana.
\end{twierdzenie}
\begin{proof}[Część algebraiczna i standardowa naturalność ciągu]
Dowód podzielimy na cztery kroki, oddzielając istnienie ciągu
\emph{naturalnego} od konstrukcji jego \emph{nienaturalnego} rozszczepienia.

\emph{1. Rozkład kompleksu wolnego.}
Niech $C$ będzie nieujemnie stopniowanym kompleksem wolnych grup
abelowych. Piszemy $Z_k=\ker(\partial_k:C_k\to C_{k-1})$ oraz
$B_k=\operatorname{im}(\partial_{k+1}:C_{k+1}\to C_k)$.
Każda podgrupa grupy wolnej abelowej jest wolna, również
przy nieskończonej randze. Wobec tego $B_{k-1}$ jest grupą
projektywną i ciąg
$0\to Z_k\to C_k\xrightarrow{\partial_k}B_{k-1}\to0$
rozszczepia się. Wybierzmy rozszczepienia
$s_{k-1}:B_{k-1}\to C_k$. Otrzymujemy izomorfizm
$C_k\cong Z_k\oplus B_{k-1}$, przy którym
$\partial_k(z,b)=(b,0)\in Z_{k-1}\oplus B_{k-2}$.
Stąd $C$ jest (niekanonicznie) izomorficzny z sumą prostą
kompleksów $P_k=(B_k\hookrightarrow Z_k)$ w stopniach
$k+1,k$. Każdy $P_k$ jest dwustopniową wolną rezolwentą
$H_k(C)=Z_k/B_k$. Ten argument nie korzysta z twierdzenia
strukturalnego o skończenie generowanych grupach abelowych.

\emph{2. Tensor z dowolnym kompleksem wolnym.}
Niech $D$ będzie drugim kompleksem grup wolnych.
Dla ustalonego $k$ użyjmy dokładnego ciągu kompleksów
$0\to Z_k\otimes D\to P_k\otimes D\to (B_k\otimes D)[1]\to0$,
gdzie $[1]$ oznacza przesunięcie stopnia homologicznego o jeden.
Dokładność wynika z tego, że w obu stopniach $P_k$
rozkład rozpatruje bezpośrednio kompleks stożka
$B_k\otimes D\to Z_k\otimes D$.
Ponieważ $Z_k,B_k$ są wolne, mamy
$H_j(Z_k\otimes D)\cong Z_k\otimes H_j(D)$ i analogicznie
dla $B_k$. Długi ciąg homologii stożka daje zatem
\[
\begin{split}
0\to&H_k(C)\otimes H_j(D)\longrightarrow
H_j(P_k\otimes D)\\
\longrightarrow&\operatorname{Tor}^{\Z}_1(H_k(C),H_{j-1}(D))\to0.
\end{split}
\]
W powyższym zapisie $P_k$ na chwilę traktujemy bez przesunięcia,
jako kompleks w stopniach $1,0$; przy wstawieniu go w stopnie
$k+1,k$ oba indeksy homologiczne należy zwiększyć o $k$.
Identyfikacja jądra z $\operatorname{Tor}_1$ pochodzi z wolnej
rezolwenty $0\to B_k\to Z_k\to H_k(C)\to0$.
W tym kroku nie wybraliśmy żadnego rozkładu samych grup homologii.

\emph{3. Naturalność ciągu.}
Aby uniknąć pozornej naturalności wybranych rozszczepień
$C_k\cong Z_k\oplus B_{k-1}$, definiujemy lewą strzałkę
bezpośrednio jako iloczyn zewnętrzny cykli
$[c]\otimes[d]\mapsto[c\otimes d]$.
Opis powyższy pokazuje, że jest ona monomorfizmem,
a jej kokernel to wskazana suma grup $\operatorname{Tor}_1$.
Naturalna strzałka $\kappa$ wymaga osobnej konstrukcji:
można ją otrzymać z homomorfizmów łączących w ciągach
dokładnych filtracji tensora wolnych kompleksów, równoważnie
z naturalnej sekwencji Künnetha. Same niekanoniczne rozkłady
$C_k=Z_k\oplus B_{k-1}$ nie dowodzą jej naturalności.
W tym miejscu przyjmujemy standardowy lemat o naturalności
konstrukcji $\kappa$; pozostałe kroki wykazują dokładność,
identyfikację członów i nienaturalne rozszczepienie.

\emph{4. Rozszczepienie bez skończonej generacji.}
Przeprowadźmy rozkład z kroku 1 \emph{dla obu} kompleksów.
Tensor rozdziela się nad dowolnymi sumami prostymi:
$C\otimes D\cong\bigoplus_{k,l}P_k(C)\otimes P_l(D)$
z uwzględnieniem przesunięć.
Dla dwóch wolnych rezolwent
$P(A):F_1\hookrightarrow F_0$ i
$P(B):G_1\hookrightarrow G_0$ kompleks tensorowy ma
zaledwie trzy stopnie $0,1,2$. Jego homologia w stopniu $2$
zanika: różniczka
$F_1\otimes G_1\to(F_0\otimes G_1)\oplus(F_1\otimes G_0)$
jest injektywna, bo $F_1\hookrightarrow F_0$ i $G_1$
jest płaska. Homologia w stopniu $0$ to $A\otimes B$,
a w stopniu $1$ to $\operatorname{Tor}_1(A,B)$:
aby to sprawdzić, rozpatrzmy dwustopniową filtrację
bikompleksu $P(A)\otimes P(B)$ według stopnia $P(A)$.
Dla ustalonego wolnego $F_i$ kompleks pionowy
$F_i\otimes(G_1\hookrightarrow G_0)$ ma homologię
$F_i\otimes B$ w stopniu $0$ i zero w stopniu $1$,
bo $F_i$ jest płaska. Zatem homologia całego tensoru
jest homologią dwustopniowego kompleksu
$F_1\otimes B\to F_0\otimes B$, a jego grupy
w stopniach $1,0$ są z definicji odpowiednio
$\operatorname{Tor}_1(A,B)$ i $A\otimes B$. Zatem homologia \emph{każdego}
$P_k(C)\otimes P_l(D)$ ma odrębne składowe
$H_k(C)\otimes H_l(D)$ w stopniu $k+l$ i
$\operatorname{Tor}_1(H_k(C),H_l(D))$ w stopniu $k+l+1$.
Ponieważ homologia komutuje z sumami prostymi kompleksów,
wybrane rozkłady obu kompleksów dają poszukiwane
rozszczepienie, choć zależne od wyborów.

Na koniec stosujemy twierdzenie Eilenberga--Zilbera:
$C_*(X\times Y)$ jest naturalnie homotopijnie równoważny
$C_*(X)\otimes C_*(Y)$; grupy łańcuchów singularnych są wolne
bez żadnych założeń o skończoności. Wobec tego powyższy
ciąg algebraiczny daje tezę, a jego lewa strzałka jest
klasycznym iloczynem zewnętrznym.
\end{proof}
\begin{przyklad}[Nieskończenie generowana grupa i brak problemu ze skończonością]
Dla $A=\bigoplus_{j\ge1}\Z/2^j$ wybierzmy dwustopniową wolną
rezolwentę
$\bigoplus_{j\ge1}\Z\xrightarrow{\oplus_j\,2^j}
 \bigoplus_{j\ge1}\Z$.
Wówczas dla $B=\Z/2$ mamy
$A\otimes B\cong\bigoplus_{j\ge1}\Z/2$ oraz
$\operatorname{Tor}^{\Z}_1(A,B)\cong\bigoplus_{j\ge1}\Z/2$.
Wyprowadzenie wynika bezpośrednio z tensora wolnych rezolwent;
nie wolno użyć tu rozkładu dla \emph{skończenie} generowanych
grup w zastępstwie dowodu ogólnego.
\end{przyklad}


\begin{twierdzenie}[Künneth dla kohomologii skończonych CW-kompleksów]
Niech $X,Y$ będą skończonymi CW-kompleksami. Dla $n\ge0$
istnieje naturalny krótki ciąg dokładny
\[
0\to\bigoplus_{p+q=n}H^p(X;\Z)\otimes H^q(Y;\Z)
\xrightarrow{\times}H^n(X\times Y;\Z)
\to\bigoplus_{p+q=n+1}\operatorname{Tor}^{\Z}_1
(H^p(X;\Z),H^q(Y;\Z))\to0.
\]
Rozszczepienie jest na ogół nienaturalne.
\end{twierdzenie}
\begin{proof}
Skończony rozkład CW daje kompleksy komórkowe z wolnych
grup skończonej rangi. Komórki produktu mają postać
$e^p\times e^q$, a wzór na ich granice to
$\partial(e^p\times e^q)=\partial e^p\times e^q
+(-1)^p e^p\times\partial e^q$.
Dla skończonych wolnych grup kanoniczna mapa
$\operatorname{Hom}(A,\Z)\otimes\operatorname{Hom}(B,\Z)
\to\operatorname{Hom}(A\otimes B,\Z)$ jest izomorfizmem.
Otrzymujemy więc iloczyn tensorowy komórkowych kompleksów
\emph{kołańcuchowych}, a do niego stosujemy powyższe
twierdzenie algebraiczne, zmieniając kierunek stopniowania.
Stąd po prawej występuje warunek $p+q=n+1$.
Lewa strzałka jest iloczynem zewnętrznym
$\alpha\times\beta=\mathrm{pr}_X^*\alpha\smile\mathrm{pr}_Y^*\beta$;
identyfikację potwierdza sprawdzenie na bazach dualnych komórek.
\end{proof}

\begin{przyklad}
Dla $\R P^2$ mamy $H^0(\R P^2;\Z)=\Z$,
$H^2(\R P^2;\Z)=\Z/2$ i brak innych dodatnich grup
kohomologii. Wzór kohomologiczny w stopniu $3$ daje
$H^3(\R P^2\times\R P^2;\Z)\cong\Z/2$ wyłącznie dzięki
$\operatorname{Tor}_1(\Z/2,\Z/2)$.
\end{przyklad}

\begin{uwaga}[Stopniowa przemienność iloczynu kubkowego]
Wcześniejszy dowód stopniowej przemienności można domknąć
konstrukcją z kroku 3. Niech $D=\mathrm{AW}\circ\Delta_\#$,
gdzie $\Delta:X\to X\times X$ jest diagonalą, a
$\tau(a\otimes b)=(-1)^{pq}b\otimes a$ dla stopni $p,q$.
Mapy $D$ i $\tau D$ są naturalnymi mapami łańcuchowymi
$C_*(X)\to C_*(X)\otimes C_*(X)$, zgodnymi w stopniu zero.
Ich różnicę w stopniu $n$ na modelu $\Delta^n$ wypełniamy
indukcyjnie kontrakcją tensoru
$C_*(\Delta^n)\otimes C_*(\Delta^n)$, identycznie jak w
dowodzie Eilenberga--Zilbera. Powstaje homotopia
$D-\tau D=\partial K+K\partial$. Po zastosowaniu
$\alpha\otimes\beta$ do tego wzoru, dla kocykli stopni
$p,q$, dostajemy
$\alpha\smile\beta-(-1)^{pq}\beta\smile\alpha=\delta\gamma$.
To podaje brakujący argument bez fałszywego twierdzenia o
przemienności na poziomie samych kołańcuchów.
\end{uwaga}


\section{Związek z kohomologią de Rhama}

\begin{definicja}
Niech $M$ będzie gładką rozmaitością, a $\Omega^k(M)$ przestrzenią
gładkich form różniczkowych stopnia $k$. Różniczka zewnętrzna
$\mathrm d:\Omega^k(M)\to\Omega^{k+1}(M)$ spełnia $\mathrm d^2=0$.
Kohomologią de Rhama nazywamy
$H^k_{\mathrm{dR}}(M)=\ker(\mathrm d:\Omega^k(M)\to\Omega^{k+1}(M))
/\operatorname{im}(\mathrm d:\Omega^{k-1}(M)\to\Omega^k(M))$,
przyjmując $\Omega^{-1}(M)=0$. Pełne konstrukcje i przykłady
pojawią się w następnym rozdziale.
\end{definicja}

\begin{twierdzenie}[Twierdzenie de Rhama, postać porównawcza]
Niech $M$ będzie gładką rozmaitością, a $n\geq0$. Całkowanie
gładkich form $n$-stopnia po gładkich sympleksach singularnych
indukuje naturalny izomorfizm
$H^n_{\mathrm{dR}}(M)\cong H^n(M;\R)$.
Izomorfizmy te zachowują iloczyn zewnętrzny po lewej i iloczyn
kubkowy po prawej na poziomie pierścieni kohomologii.
\end{twierdzenie}

\begin{uwaga}
Dla formy zamkniętej $\omega$ jej klasie odpowiada kocykl
$\sigma\mapsto\int_{\Delta^n}\sigma^*\omega$.
Twierdzenie Stokesa gwarantuje, że jest to kocykl,
a forma dokładna daje kobrzeg. Pełne uzasadnienie, że
konstrukcja jest izomorfizmem, korzysta z lematu Poincarégo
i ciągu Mayera--Vietorisa; jest rozwinięte w następnym
rozdziale. Zgodność iloczynów na kohomologii wynika z
porównania przekątnej singularnej z iloczynem zewnętrznym,
nie z punktowej równości wartości na pojedynczym sympleksie.
\end{uwaga}

\chapter{Formy różniczkowe, orientacja i kohomologia de Rhama}
\label{ch:formy-de-rham}

W konstrukcji algebry tensorowej zbudowaliśmy algebrę
zewnętrzną jednej przestrzeni wektorowej. Teraz wykonamy tę
konstrukcję w każdym punkcie rozmaitości i pozwolimy, by jej
współczynniki zmieniały się gładko. Otrzymamy formy różniczkowe.
Różniczkowanie form, ich całkowanie i twierdzenie Stokesa
połączą rachunek różniczkowy z homologią przestrzeni z
Rozdziału~\ref{ch:homologia-przestrzeni}. Na końcu porównamy
obie teorie kohomologii. Nadal nie potrzebujemy metryki,
koneksji ani krzywizny; tym pojęciom poświęcimy późniejszy blok.

\section{Czym jest forma różniczkowa?}
\label{sec:formy-definicje}

\begin{definicja}[Formy i ich wartości]
\label{def:forma-rozniczkowa}
Niech $M$ będzie gładką rozmaitością wymiaru $n$.
\emph{Formą różniczkową stopnia $k$} na $M$ nazywamy
gładkie pole naprzemiennych form $k$-liniowych:
\[
 \omega_p\in\Lambda^k T_p^*M,
 \qquad \omega\in\Omega^k(M)
 :=\Gamma(\Lambda^k T^*M).
\]
Przy $k=0$ jest to funkcja gładka,
$\Omega^0(M)=C^\infty(M)$; dla $k>n$ otrzymujemy
$\Omega^k(M)=0$. Przyjmujemy też $\Omega^k(M)=0$ dla $k<0$.
Na rozmaitości z brzegiem gładkość
oznacza lokalne przedłużenie współczynników do
otwartego otoczenia w $\R^n$.
\end{definicja}

Forma $1$-stopnia mierzy wektor styczny i zwraca liczbę.
Forma $2$-stopnia mierzy \emph{zorientowane pole} zadane
parą wektorów; zamiana wektorów zmienia znak.
Jeśli $(U,x)$ jest mapą $x=(x^1,\ldots,x^n)$, to
\begin{equation}\label{eq:forma-lokalna}
 \omega|_U=\sum_{i_1<\cdots<i_k}
 a_{i_1\ldots i_k}\,\dd x^{i_1}\wedge\cdots\wedge\dd x^{i_k},
 \qquad a_I\in C^\infty(U).
\end{equation}
Tutaj $\dd x^i|_p(v)=v(x^i)$ jest kowektorem z definicji przestrzeni kostycznej. Jedyność współczynników
wynika z bazy algebry zewnętrznej
(wynika to z niezależności standardowych iloczynów zewnętrznych bazowych kowektorów). Zmiana mapy zmienia
współczynniki, lecz nie samą formę.

\begin{przyklad}[Dwie formy na płaszczyźnie]
Na $\R^2$ forma $\alpha=x\,\dd y-y\,\dd x$
przypisuje wektorowi $(u,v)$ liczbę $xv-yu$.
Forma $\omega=\dd x\wedge\dd y$ spełnia
\[
 \omega((u_1,v_1),(u_2,v_2))
 =u_1v_2-v_1u_2.
\]
Jest to wyznacznik: wartość bezwzględna daje pole
równoległoboku, znak zaś jego orientację.
\end{przyklad}

\begin{figure}[!htbp]
\centering
\begin{tikzpicture}[>=Stealth,line cap=round,line join=round]
 \fill[manifoldblue!12] (0,0)--(3.2,.45)--(4.25,2.65)--(1.05,2.2)--cycle;
 \draw[manifoldblue!70!black,thick]
  (0,0)--(3.2,.45)--(4.25,2.65)--(1.05,2.2)--cycle;
 \draw[->,manifoldred,very thick] (0,0)--(3.2,.45)
  node[pos=.78,below=5pt] {$u$};
 \draw[->,manifoldgreen!70!black,very thick] (0,0)--(1.05,2.2)
  node[pos=.8,left=5pt] {$v$};
 \draw[->,manifoldgold!90!black,thick]
  (8:1.45) arc[start angle=8,end angle=61,radius=1.45];
 \node[font=\small,align=left,anchor=west,text width=5.3cm]
  at (5.05,1.42)
  {$(\dd x\wedge\dd y)(u,v)=\det(u,v)>0$.
   Po zamianie $u$ i $v$ ten sam obszar
   dostaje przeciwny znak.};
\end{tikzpicture}
\caption{Forma $\dd x\wedge\dd y$ daje pole ze znakiem
w standardowych współrzędnych płaszczyzny. Ogólna $2$-forma
przypisuje liczbę parze wektorów bez wyboru metryki;
nie musi być euklidesowym polem ich równoległoboku.}
\label{fig:forma-zorientowane-pole}
\end{figure}

\section{Iloczyn zewnętrzny, cofnięcie i różniczka}
\label{sec:formy-dzialania}

\begin{stwierdzenie}[Działania punkt po punkcie]
\label{prop:formy-wedge-pullback}
Dla $\alpha\in\Omega^k(M)$ i $\beta\in\Omega^\ell(M)$
iloczyn $\alpha\wedge\beta\in\Omega^{k+\ell}(M)$
otrzymujemy, biorąc iloczyn zewnętrzny w każdym
$T_p^*M$. Zachodzi $\alpha\wedge\beta
=(-1)^{k\ell}\beta\wedge\alpha$.
Jeżeli $F:M\to N$ jest gładkie, to \emph{cofnięcie}
$F^*:\Omega^k(N)\to\Omega^k(M)$ ma wzór
\begin{equation}\label{eq:pullback-formy}
 (F^*\omega)_p(v_1,\ldots,v_k)
 =\omega_{F(p)}(\dd F_pv_1,\ldots,\dd F_pv_k).
\end{equation}
Spełnia $(G\circ F)^*=F^*G^*$ i
$F^*(\alpha\wedge\beta)=F^*\alpha\wedge F^*\beta$.
\end{stwierdzenie}
\begin{proof}
Gładkość iloczynu w mapie wynika z faktu, że jego
współczynniki są sumami iloczynów gładkich $a_I,b_J$;
przemienność z odpowiednim znakiem jest własnością
algebry zewnętrznej z algebry tensorowej.
Wzór~\eqref{eq:pullback-formy} jest wieloliniowy i
naprzemienny, a w mapach używa pochodnych współrzędnych
$F$, więc wyznacza formę gładką. Reguła łańcucha
$\dd(G\circ F)_p=\dd G_{F(p)}\circ\dd F_p$
daje własność złożenia. Iloczyn zewnętrzny definiuje
się sumą wartości na przestawionych argumentach; podstawienie
do niej tych samych $\dd F_pv_i$ dowodzi ostatniej równości.
\end{proof}

W szczególności $F^*(\dd y^j)=\dd(y^j\circ F)$.
Jeżeli $F:\R^2\to\R^2$, $F(u,v)=(u^2,v+u)$,
to $F^*(\dd x\wedge\dd y)
=\dd(u^2)\wedge\dd(v+u)=2u\,\dd u\wedge\dd v$.
Współczynnik $2u$ jest wyznacznikiem macierzy Jacobiego
i rejestruje także odwrócenie orientacji przy $u<0$.

\begin{twierdzenie}[Różniczka zewnętrzna]
\label{thm:formy-d}
Istnieje dokładnie jedna rodzina $\R$-liniowych
operatorów $\dd:\Omega^k(M)\to\Omega^{k+1}(M)$
o następujących własnościach:
\begin{enumerate}[leftmargin=*]
\item dla funkcji $f$ forma $\dd f$ jest jej różniczką;
\item $\dd(\alpha\wedge\beta)=\dd\alpha\wedge\beta
 +(-1)^k\alpha\wedge\dd\beta$, gdy $\deg\alpha=k$;
\item $\dd^2=0$.
\end{enumerate}
W mapie jej wzór jest szczególnie prosty:
\begin{equation}\label{eq:zewnetrzna-lokalnie}
 \dd\!\left(\sum_I a_I\,\dd x^I\right)
  =\sum_I\sum_j\frac{\partial a_I}{\partial x^j}
     \dd x^j\wedge\dd x^I,
 \qquad \dd x^I=\dd x^{i_1}\wedge\cdots\wedge\dd x^{i_k}.
\end{equation}
Dla każdego gładkiego $F:M\to N$ zachodzi
$F^*\dd\omega=\dd(F^*\omega)$.
\end{twierdzenie}
\begin{proof}
Najpierw sprawdźmy lokalność dowolnego operatora spełniającego
podane warunki. Jeśli $\omega$ znika w otoczeniu $p$, wybierzmy
gładką funkcję $\chi$, równą $1$ blisko $p$, z nośnikiem
w tym otoczeniu. Taką pojedynczą funkcję buduje się w mapie
z $e^{-1/t}$, jak w dowodzie Twierdzenia
\ref{thm:podzial-jednosci-formy}; nie wymaga to różniczki form.
Ponieważ $\chi\omega=0$, reguła Leibniza daje w $p$
$0=\dd(\chi\omega)_p=(\dd\omega)_p$.
Operator zależy więc tylko od formy w otoczeniu punktu.
Lokalne funkcje i formy możemy pomnożyć przez taką funkcję
odcinającą i przedłużyć przez zero; wynik różniczkowania
blisko $p$ nie zależy od przedłużenia.

Każda forma lokalnie jest sumą wyrazów $a_I\dd x^I$.
Możemy już stosować warunki do lokalnych współrzędnych.
Ponieważ $\dd(\dd x^i)=\dd^2x^i=0$,
reguła Leibniza \emph{wymusza} wzór
\eqref{eq:zewnetrzna-lokalnie}; mamy więc jednoznaczność.
Zdefiniujmy operator tym wzorem. Po powtórnym użyciu
otrzymujemy współczynniki
$(\partial_{x^j}\partial_{x^h}a_I)
\dd x^j\wedge\dd x^h\wedge\dd x^I$.
Wyrazy z $(j,h)$ i $(h,j)$ znoszą się, gdyż pochodne
mieszane są równe, a iloczyn zewnętrzny zmienia znak.
Stąd $\dd^2=0$.

Dla dwóch wyrazów $a\dd x^I$ i $b\dd x^J$
zwykła reguła iloczynu $\dd(ab)=a\dd b+b\dd a$
oraz przestawienie $\dd b$ przez $k$ czynników
$\dd x^I$ dają dokładnie znak $(-1)^k$.
Liniowość rozszerza to na sumy.
Sprawdźmy zgodność na przecięciu map. Cofnięcie przez
zmianę współrzędnych zachowuje $\dd f$ dzięki regule
łańcucha i zachowuje $\wedge$ z poprzedniego
stwierdzenia. Zatem dwa lokalne wzory, obliczone po obu
stronach zmiany współrzędnych, są równe na generatorach
$f,\dd x^i$, a przez regułę Leibniza na wszystkich formach.
Wzory sklejają się w globalny operator. Ten sam argument
dla dowolnego $F$ daje $F^*\dd=\dd F^*$.
\end{proof}

\begin{przyklad}[Rachunek na $\R^3$]
Jeżeli $\alpha=P\,\dd x+Q\,\dd y+R\,\dd z$, to
\[
 \dd\alpha=(Q_x-P_y)\,\dd x\wedge\dd y
 +(R_x-P_z)\,\dd x\wedge\dd z
 +(R_y-Q_z)\,\dd y\wedge\dd z.
\]
Dla $\beta=P\,\dd y\wedge\dd z
+Q\,\dd z\wedge\dd x+R\,\dd x\wedge\dd y$
otrzymujemy $\dd\beta=(P_x+Q_y+R_z)
\dd x\wedge\dd y\wedge\dd z$.
Są to wyrażenia występujące w klasycznych wzorach na rotację
i dywergencję. Utożsamianie tu pól wektorowych z formami
korzysta ze standardowej struktury euklidesowej $\R^3$;
sama różniczka zewnętrzna na $M$ nie wymaga metryki.
\end{przyklad}

\begin{definicja}[Kontrakcja i pochodna Liego]
\label{def:formy-kontrakcja-lie}
Dla pola wektorowego $X$ operator
$\iota_X:\Omega^k(M)\to\Omega^{k-1}(M)$ określamy przez
\[
 (\iota_X\omega)_p(v_1,\ldots,v_{k-1})
 =\omega_p(X_p,v_1,\ldots,v_{k-1}),
 \qquad \iota_Xf=0.
\]
Jeśli $M$ ma brzeg, w definicji przez dwustronny przepływ
zakładamy, że $X$ jest styczne do $\partial M$.
Jeśli $\Phi_t$ jest lokalnym przepływem $X$ z teorii przepływów pól wektorowych, kładziemy
$\Lie_X\omega=\left.\frac{\dd}{\dd t}\right|_{t=0}\Phi_t^*\omega$.
\end{definicja}

\begin{twierdzenie}[Wzór Cartana]
\label{thm:cartan-formy}
Dla każdej formy $\omega$ mamy
\begin{equation}\label{eq:cartan-formula}
 {\ \Lie_X\omega=\dd(\iota_X\omega)
       +\iota_X(\dd\omega)\ }.
\end{equation}
Ponadto $\Lie_X\dd=\dd\Lie_X$ oraz
$\iota_X(\alpha\wedge\beta)=
(\iota_X\alpha)\wedge\beta
+(-1)^{\deg\alpha}\alpha\wedge\iota_X\beta$.
\end{twierdzenie}
Dla pola niestycznego do brzegu prawa strona wzoru Cartana
nadal określa gładką formę. Możemy nią zdefiniować
$\Lie_X\omega$ na brzegu: jest to przedłużenie pochodnej
Liego z wnętrza, niezależne od przedłużeń w mapach.
Nie zakładamy wówczas istnienia dwustronnego przepływu w $M$.
\begin{proof}
Ostatnia równość wynika wprost z definicji iloczynu
zewnętrznego: po wstawieniu $X$ w pierwsze miejsce
może on wejść do pierwszej grupy argumentów lub,
po przejściu przez $\deg\alpha$ argumentów, do drugiej.
Ponieważ $\Phi_t^*$ zachowuje iloczyn zewnętrzny,
$\Lie_X$ jest derywacją stopnia $0$ algebry form.
Prawa strona~\eqref{eq:cartan-formula} również nią jest:
rozwijamy oba składniki na $\alpha\wedge\beta$
regułami dla $\dd$ i $\iota_X$; wyrazy mieszane mają
przeciwne znaki i znoszą się.

Wystarczy zatem porównać obie strony na lokalnych
generatorach $f,\dd x^i$. Dla funkcji lewa strona to
$Xf$, a prawa to $\iota_X\dd f=\dd f(X)=Xf$.
Dla $\dd x^i$ lewa strona wynosi
$\dd(\Lie_Xx^i)=\dd(Xx^i)$, gdyż cofnięcie
komutuje z $\dd$; prawa wynosi
$\dd(\iota_X\dd x^i)+\iota_X\dd^2x^i
=\dd(Xx^i)$. Wzór obowiązuje więc dla
każdej sumy wyrazów $f\,\dd x^I$ i globalnie.
Wreszcie, z $\dd^2=0$ dostajemy
$\dd\Lie_X=\dd\iota_X\dd
=\Lie_X\dd$.
\end{proof}

\begin{przyklad}[Pochodna Liego formy]
Na $\R^2$ weźmy $X=\partial_x$ i $\omega=x\,\dd y$.
Przepływ przesuwa punkty poziomo, więc
$\Phi_t^*\omega=(x+t)\dd y$ oraz
$\Lie_X\omega=\dd y$. Ze wzoru Cartana:
$\iota_X\omega=0$, $\dd\omega=\dd x\wedge\dd y$
i $\iota_X\dd\omega=\dd y$. Oba obliczenia są zgodne.
\end{przyklad}
\section{Lokalne istnienie funkcji pierwotnej}
\label{sec:formy-poincare}

Równość $\dd^2=0$ mówi, że każda forma postaci
$\dd\eta$ jest \emph{zamknięta}. Odwrotność nie zachodzi
globalnie: przeszkodę może stanowić topologia. Lokalnie
przeszkoda znika, co wyjaśnia następujący rachunek.
Dla zamkniętej $1$-formy $\omega$ szukamy \emph{funkcji
pierwotnej} $f$ spełniającej $\omega=\dd f$.
Dla formy stopnia $k>1$ jej odpowiednikiem jest
\emph{forma pierwotna} $\eta$ stopnia $k-1$ z
$\omega=\dd\eta$. Lemat Poincarégo obejmuje
oba przypadki.

\begin{lemat}[Homotopia form]
\label{lem:homotopia-form}
Niech $H:[0,1]\times M\to N$ będzie gładką homotopią
odwzorowań $H_0,H_1$. Przy końcach przedziału, a także
przy narożach iloczynu, gdy $M$ ma brzeg, gładkość
rozumiemy jako gładkie przedłużanie składowych w mapach. Istnieje operator liniowy
$K:\Omega^k(N)\to\Omega^{k-1}(M)$ taki, że
\begin{equation}\label{eq:homotopia-form}
 H_1^*\omega-H_0^*\omega
 =\dd(K\omega)+K(\dd\omega).
\end{equation}
\end{lemat}
\begin{proof}
Dla $k=0$ przyjmujemy $K\omega=0$; poniższy rachunek
obejmuje ten przypadek z $a=0$. Każdą formę na
$[0,1]\times M$ zapisujemy jednoznacznie
w postaci $H^*\omega=\dd t\wedge a(t)+b(t)$,
gdzie $a(t)$ jest rodziną form stopnia $k-1$ na $M$,
a $b(t)$ rodziną form stopnia $k$. Położenie
$K\omega:=\int_0^1 a(t)\,\dd t$ oznacza całkowanie
\emph{współczynników} formy po parametrze; można je
wykonać w każdej mapie, a zgodność na przecięciach
wynika z liniowości zmiany współrzędnych.
Obliczenie różniczki daje
\[
 \dd(H^*\omega)
 =\dd t\wedge(\partial_tb-\dd_Ma)+\dd_Mb.
\]
Ponieważ $H^*\dd\omega=\dd H^*\omega$, współczynniki
przy $\dd t$ mówią, że
$a_{\dd\omega}(t)=\partial_tb(t)-\dd_Ma(t)$.
Po scałkowaniu od $0$ do $1$ otrzymujemy
$K\dd\omega=b(1)-b(0)-\dd_MK\omega$.
Ale $b(i)=H_i^*\omega$ dla $i=0,1$,
więc jest to dokładnie~\eqref{eq:homotopia-form}.
\end{proof}

\begin{twierdzenie}[Lemat Poincarégo]
\label{thm:lemat-poincare-formy}
Niech $U\subset\R^n$ będzie otwarty i gwiaździsty
względem $0$. Jeżeli $k\geq1$ i
$\omega\in\Omega^k(U)$ spełnia $\dd\omega=0$,
to istnieje $\eta\in\Omega^{k-1}(U)$ z
$\omega=\dd\eta$. Każda zamknięta forma stopnia
$0$ jest stała na $U$.
\end{twierdzenie}
\begin{proof}
W Lemacie~\ref{lem:homotopia-form} bierzemy
$H(t,x)=tx$. Dla $k\geq1$ odwzorowanie $H_0$
jest stałe, więc $H_0^*\omega=0$; dla zamkniętej
$\omega$ drugi składnik po prawej stronie
\eqref{eq:homotopia-form} także znika.
Dostajemy $\omega=\dd(K\omega)$.
Jeśli chcemy \emph{obliczyć} formę pierwotną
(dla $k=1$: funkcję pierwotną), to
\begin{equation}\label{eq:poincare-pierwotna}
 (K\omega)_x(v_1,\ldots,v_{k-1})
 =\int_0^1 t^{k-1}
   \omega_{tx}(x,v_1,\ldots,v_{k-1})\,\dd t.
\end{equation}
Istotnie, $\dd H(\partial_t)=x$, podczas gdy
$\dd H(v_i)=tv_i$, skąd $k-1$ czynników $t$.
Gwiaździstość zapewnia, że cały odcinek $tx$ leży w $U$.
Dla $k=0$ warunek $\dd f=0$ mówi, że wszystkie
pochodne kierunkowe znikają. Zatem pochodna
$t\mapsto f(tx)$ jest zerowa i $f(x)=f(0)$.
\end{proof}

\begin{przyklad}[Obliczenie funkcji pierwotnej]
Na $\R^2$ forma $\omega=2xy\,\dd x+x^2\,\dd y$
jest zamknięta, bo
$\dd\omega=(2x-2x)\dd x\wedge\dd y=0$.
Wzór~\eqref{eq:poincare-pierwotna} dla $k=1$ daje
\[
 (K\omega)(x,y)=\int_0^1
 \omega_{(tx,ty)}(x,y)\,\dd t
 =\int_0^1 3t^2x^2y\,\dd t=x^2y.
\]
Zatem $f(x,y)=x^2y$ jest funkcją pierwotną
$1$-formy $\omega$, ponieważ $\omega=\dd f$.
Obliczenie pokazuje mechanizm lematu.
\end{przyklad}

\begin{wniosek}[Niezmienniczość przez homotopię]
\label{cor:de-rham-homotopia}
Odwzorowania gładko homotopijne indukują tę samą
operację cofnięcia na kohomologii de Rhama
z Definicji~\ref{def:de-rham-cohomology}.
\end{wniosek}
\begin{proof}
Jeśli $\dd\omega=0$, wzór~\eqref{eq:homotopia-form}
mówi, że $H_1^*\omega-H_0^*\omega=\dd(K\omega)$.
Różnica jest więc dokładna i obie klasy są równe.
\end{proof}

\section{Podział jedności}
\label{sec:formy-podzial-jednosci}

Przy całkowaniu trzeba scalać dane z wielu map.
Służy do tego wynik, który wyjaśnia także wcześniejsze
użycia funkcji odcinających w skrypcie.

\begin{twierdzenie}[Gładki podział jedności]
\label{thm:podzial-jednosci-formy}
Niech $(U_\alpha)$ będzie otwartym pokryciem rozmaitości
$M$ Hausdorffa i o przeliczalnej bazie. Istnieją
nieujemne funkcje gładkie $(\rho_j)$ o lokalnie skończonej
rodzinie nośników (każdy punkt ma otoczenie spotykające
tylko skończenie wiele nośników) oraz indeksy $\alpha(j)$ takie, że
\[
 \supp\rho_j\subset U_{\alpha(j)},
 \qquad \sum_j\rho_j(p)=1\quad(p\in M).
\]
Każdą $\rho_j$ można wybrać o nośniku zwartym
w jednej dziedzinie mapy zawartej w $U_{\alpha(j)}$.
\end{twierdzenie}
\begin{proof}
Najpierw budujemy jedną funkcję odcinającą.
Jeśli w mapie kula domknięta $\overline B_r$
leży w większej kuli $B_R$, przy czym $0<r<R$
i $\overline B_R$ zawiera się w obrazie mapy,
wybierzmy gładką funkcję $\chi$ zależną tylko
od $\|x\|^2$, równą $1$ na $\overline B_r$,
dodatnią między $r$ i $R$, a zerową poza $B_R$.
Można ją złożyć z funkcji $h(t)=e^{-1/t}$ dla
$t>0$ i $h(t)=0$ dla $t\leq0$; jej gładkość w
$t=0$ wynika z zanikania wszystkich pochodnych
$t^{-a}e^{-1/t}$. Przedłużenie $\chi$ przez $0$
na resztę $M$ jest gładkie, bo nośnik ma zwartą
domkniętą otoczkę wewnątrz mapy. Przy brzegu bierzemy
przecięcia tych kul z półprzestrzenią i ograniczamy do nich
te same funkcje. Wszystkie wnętrza i domknięcia poniżej
rozumiemy w topologii $M$, także przy brzegu.

Z lokalnej zwartości i przeliczalności bazy
konstruujemy wyczerpanie zwartymi zbiorami
$K_1\subset\operatorname{int}K_2\subset\cdots$,
$\bigcup_jK_j=M$ (weźmy kolejno skończone
sumy domkniętych, stosunkowo zwartych kul
obejmujące pierwsze elementy przeliczalnego
pokrycia i powiększmy je przy każdym kroku).
Ustalmy $K_j=\varnothing$ dla $j\leq0$.
Zwarty pas $K_j\setminus\operatorname{int}K_{j-1}$
można pokryć skończenie wieloma małymi kulami
w mapach o domknięciach zawartych jednocześnie
w którymś $U_\alpha$ i w
$\operatorname{int}K_{j+1}\setminus K_{j-2}$.
Wybieramy funkcje odcinające dodatnie na
mniejszych kulach pokrywających pas.
Rodzina tych kul dla wszystkich $j$ jest
lokalnie skończona: otoczenie punktu wewnątrz
pewnego $K_N$ nie spotyka kul pochodzących
z pasów o dostatecznie dużym indeksie.

Niech $\chi_j$ oznaczają wszystkie otrzymane
funkcje. Suma $S=\sum_j\chi_j$ jest lokalnie
skończona, gładka i wszędzie dodatnia, bo
każdy punkt należy do któregoś pasa. Kładziemy
$\rho_j=\chi_j/S$. Wtedy nośniki pozostają
w wybranych kulach i
$\sum_j\rho_j=S/S=1$.
Gdy potrzebujemy funkcji indeksowanych samymi zbiorami
pokrycia, sumujemy $\rho_j$ o tym samym $\alpha(j)$.
Nośnik takiej sumy nadal leży w $U_\alpha$:
lokalnie skończona suma zamkniętych nośników jest
zbiorem domkniętym zawartym w $U_\alpha$.
\end{proof}
\section{Orientowalność i orientacje indukowane}
\label{sec:formy-orientacja}

\begin{definicja}[Orientacja]
\label{def:orientacja-formy}
Dwie uporządkowane bazy $m$-wymiarowej przestrzeni
rzeczywistej $V$, dla $m\geq1$, mają tę samą orientację, jeśli
macierz przejścia ma dodatni wyznacznik.
Powstają dwie klasy baz. Na rozmaitości $M^m$
\emph{orientacja} to gładki wybór takiej klasy w
każdym $T_pM$. Równoważnie jest to klasa
nigdzie niezerowych form $m$-stopnia: formy
$\mu,\nu$ wyznaczają tę samą orientację, jeśli
$\nu=f\mu$ dla gładkiej funkcji $f>0$.
Rozmaitość dopuszczająca taki wybór jest
\emph{orientowalna}; wraz z wyborem orientacji
jest \emph{zorientowana}.
\end{definicja}

W wymiarze $0$ stosujemy opis przez formy najwyższego
stopnia: $\Lambda^0 V^*=\R$, a orientacja jest wyborem
jednej z dwóch półprostych, czyli znaku $+1$ lub $-1$.
Każdemu punktowi rozmaitości dyskretnej przypisujemy taki
znak. Opis przez dwie klasy baz dotyczy tylko $m\geq1$,
gdyż przestrzeń zerowa ma tylko jedną bazę pustą.

Dla $m\geq1$ forma $\mu_p$ określa bazę dodatnią $(v_1,\ldots,v_m)$
przez nierówność $\mu_p(v_1,\ldots,v_m)>0$.
To definicja \emph{znaku}, a nie długości ani pola:
zamiana $\mu$ na $2\mu$ nie zmienia orientacji.
Na każdej spójnej orientowalnej rozmaitości istnieją
dwa wybory globalne, wzajemnie przeciwne.

\begin{twierdzenie}[Trzy opisy orientowalności]
\label{thm:orientacja-atlas-forma}
Dla rozmaitości gładkiej $M^m$ równoważne są
poniższe warunki. W mapach przy brzegu dopuszczamy
obie półprzestrzenie $x^m\geq0$ i $x^m\leq0$;
jest to ta sama struktura gładka, a pozwala stosować
odbicie współrzędnej także w wymiarze $1$.
\begin{enumerate}[leftmargin=*]
\item orientowalność $M$;
\item istnienie atlasu, w którym wszystkie wyznaczniki
      odwzorowań przejścia są dodatnie;
\item istnienie globalnej nigdzie niezerowej formy
      $\mu\in\Omega^m(M)$.
\end{enumerate}
\end{twierdzenie}
\begin{proof}
Dla $m=0$ wszystkie trzy warunki są spełnione:
rozmaitość jest dyskretna, dowolny wybór znaków jest
gładki, a funkcja $1$ jest nigdzie niezerową $0$-formą.
Atlas zerowymiarowy sam nie zapisuje wyboru znaków.
Załóżmy więc $m\geq1$.
Forma $\mu$ wybiera dodatnie bazy w każdym punkcie,
więc (3) daje (1). Jeśli wybór (1) jest ustalony,
zmniejszamy dziedziny map do spójnych otoczeń.
Każdą mapę można w razie potrzeby złożyć z odbiciem
jednej współrzędnej, aby jej standardowa baza była
dodatnia. W nowym atlasie dwie bazy lokalne opisują
tę samą orientację, zatem ich wyznacznik przejścia
jest dodatni: (1) daje (2).

Załóżmy (2). W każdej mapie z tego atlasu lokalna
forma $\mu_j=\dd x_j^1\wedge\cdots\wedge\dd x_j^m$
jest dodatnia. Wybierzmy podział jedności $(\rho_j)$
z Twierdzenia~\ref{thm:podzial-jednosci-formy},
o nośnikach w dziedzinach tych map, i przedłużmy
$\rho_j\mu_j$ przez zero. Suma
$\mu=\sum_j\rho_j\mu_j$ jest lokalnie skończona.
W dowolnym punkcie co najmniej jedno $\rho_j$ jest
dodatnie; wszystkie występujące $\mu_j$ są dodatnimi
wielokrotnościami jednej z nich, bo przejścia
mają dodatni wyznacznik. Suma nie może się
wyzerować. Tak dostajemy (3).
\end{proof}

\begin{stwierdzenie}[Dwa podstawowe sposoby przenoszenia orientacji]
\label{prop:orientacja-iloczyn-pullback}
Jeżeli $F:M\to N$ jest dyfeomorfizmem i $N$ ma
orientację $[\nu]$, to $F^*\nu$ orientuje $M$;
$F$ zachowuje orientację wtedy, gdy zgadza się ona
z uprzednio wybraną orientacją $M$. Jeśli $M^m,N^n$
są zorientowane, orientację $M\times N$ wyznacza
$\pr_M^*\mu\wedge\pr_N^*\nu$.
Po zamianie kolejności czynników znak zmienia się
o $(-1)^{mn}$.
\end{stwierdzenie}
\begin{proof}
Różniczka dyfeomorfizmu jest izomorfizmem, więc
cofnięcie nigdzie niezerowej formy najwyższego stopnia
pozostaje nigdzie niezerowe. Na iloczynie wartości
$\mu\wedge\nu$ na bazie najpierw $TM$, potem $TN$
są dodatnie. Przestawienie $m$ wektorów z pierwszego
czynnika z $n$ wektorami drugiego wymaga $mn$
zamian; każda zmienia znak formy.
\end{proof}

Orientacja brzegu i orientacja regularnego zbioru
poziomicowego są dwiema wersjami tej samej zasady:
\emph{najpierw kierunki normalne, potem styczne}.

\begin{definicja}[Orientacja brzegu: normalna na zewnątrz]
\label{def:orientacja-brzegu}
Niech $M^m$ będzie zorientowaną rozmaitością z brzegiem.
Dla $m\geq2$ baza $(v_1,\ldots,v_{m-1})$ przestrzeni $T_p\partial M$
jest dodatnia wtedy i tylko wtedy, gdy
$(n,v_1,\ldots,v_{m-1})$ jest dodatnią bazą $T_pM$
dla dowolnego wektora $n$ skierowanego na zewnątrz.
Wymóg nie zależy od konkretnego $n$: dwa takie wektory
mają dodatnio proporcjonalne składowe poprzeczne,
a dodanie wektora stycznego nie zmienia wyznacznika.
Równoważnie, dla formy $\mu$ reprezentującej orientację,
brzeg orientuje $\iota_n\mu$ ograniczona do $T\partial M$.
Ten opis działa również dla $m=1$: znak w punkcie brzegu
jest znakiem liczby $\mu_p(n)$.
\end{definicja}

\begin{stwierdzenie}[Brzeg jest orientowalny]
\label{prop:brzeg-orientowalny}
Powyższa reguła nadaje $\partial M$ gładką orientację.
Dla standardowej orientacji $[a,b]$ końcowy punkt $b$
ma znak $+1$, a początkowy $a$ ma znak $-1$.
Dla dysku $D^2$ ze standardową orientacją płaszczyzny
brzeg biegnie przeciwnie do ruchu wskazówek zegara.
\end{stwierdzenie}
\begin{proof}
W mapie półprzestrzeni $\{x^m\geq0\}$ kierunek
na zewnątrz to $-\partial_{x^m}$; reguła daje
stały wybór znaku baz współrzędnych stycznych.
Zmiany map zachowują orientację $M$ i wysyłają
kierunek na zewnątrz na kierunek na zewnątrz,
więc uzyskane wybory na przecięciach są zgodne.
W wymiarze $1$ przy $b$ wektor zewnętrzny
$+\partial_t$ jest dodatni, a przy $a$ wektor
$-\partial_t$ ujemny. W dysku przy $(1,0)$
wektor zewnętrzny jest $(1,0)$; aby para była
dodatnia, wektor styczny musi wskazywać $(0,1)$.
Ciągłość przenosi ten kierunek na cały okrąg.
\end{proof}

\begin{figure}[!htbp]
\centering
\begin{tikzpicture}[>=Stealth,line cap=round,line join=round]
 \fill[manifoldblue!9] (0,0) circle (1.65);
 \draw[manifoldblue!75!black,line width=1.2pt] (0,0) circle (1.65);
 \draw[->,manifoldred,very thick] (1.65,0)--(2.73,0)
  node[right,font=\small] {$n$};
 \draw[->,manifoldgreen!70!black,very thick] (1.65,0)--(1.65,1.15)
  node[above,font=\small] {$v$ styczny};
 \draw[->,manifoldgold!85!black,thick]
  (210:1.65) arc[start angle=210,end angle=255,radius=1.65];
 \draw[->,manifoldgold!85!black,thick]
  (120:1.65) arc[start angle=120,end angle=165,radius=1.65];
 \fill[manifoldred] (1.65,0) circle(2.5pt)
  node[below=5pt] {$p$};
 \node[font=\small,align=left,anchor=west,text width=4.6cm]
  at (4.75,.55)
  {Para $(n,v)$ ma dodatni wyznacznik.
   Dodatni kierunek na $\partial D^2$
   biegnie przeciwnie do wskazówek zegara.};
\end{tikzpicture}
\caption{Orientację brzegu wyznacza reguła
„normalna na zewnątrz jako pierwszy wektor”.}
\label{fig:orientacja-brzegu-dysk}
\end{figure}

\begin{twierdzenie}[Orientacja regularnego przeciwobrazu]
\label{thm:orientacja-poziomicy}
Niech $F:M^m\to N^q$ będzie gładkie,
$M,N$ zorientowane, $N$ bez brzegu, a $y\in N$ wartością
regularną. Jeśli $M$ ma brzeg, zakładamy ponadto, że $y$
jest wartością regularną $F|_{\partial M}$.
Wtedy $S=F^{-1}(y)$ jest orientowalną podrozmaitością
wymiaru $m-q$ (o ile nie jest pusta). Orientację
można określić dla wszystkich wymiarów następująco.
Dla dodatnich form objętości $\mu_p$ na $T_pM$
i $\nu_y$ na $T_yN$, dowolnej bazy $(e_1,\ldots,e_q)$
oraz podniesień $\dd F_pw_i=e_i$ kładziemy
\[
 \eta_p(v_1,\ldots,v_{m-q})
 =\frac{\mu_p(w_1,\ldots,w_q,v_1,\ldots,v_{m-q})}
        {\nu_y(e_1,\ldots,e_q)}.
\]
Puste listy argumentów oznaczają wartości form stopnia $0$.
Forma $\eta$ określa orientację $S$. Gdy $q\geq1$
i $m-q\geq1$, jest to reguła: $(v_1,\ldots,v_{m-q})$
w $T_pS=\ker\dd F_p$ jest dodatnia, jeśli dla
wektorów $w_1,\ldots,w_q\in T_pM$, których obrazy
$(\dd F_pw_1,\ldots,\dd F_pw_q)$ tworzą dodatnią
bazę $T_yN$, ciąg
$(w_1,\ldots,w_q,v_1,\ldots,v_{m-q})$
jest dodatni w $T_pM$. Gdy $m=q$, punktowi $p\in S$
przypisujemy znak izomorfizmu $\dd F_p$ między
zorientowanymi przestrzeniami. Przy brzegu
$\partial S=S\cap\partial M$.
\end{twierdzenie}
\begin{proof}
Regularność znaczy, że $\dd F_p$ jest surjekcją,
więc takie $w_i$ istnieją. Przy $\partial M$ dodatkowa
surjekcja na przestrzeni stycznej do brzegu pozwala wybrać
$q$ współrzędnych stycznych, w których pochodne $F$ są
niezależne. Twierdzenie o funkcji odwrotnej zastosowane
do $F$ wraz z pozostałymi współrzędnymi, w tym współrzędną
normalną $s\geq0$, daje lokalną postać poziomicy
$\R^{m-q-1}\times[0,\infty)$. Stąd podana formuła brzegu.
Bez założenia o $F|_{\partial M}$ wniosek mógłby zawieść:
dla $F(x,s)=x^2+s$ na $s\geq0$ wartość $0$ jest regularna
dla $F$, lecz przeciwobraz jest pojedynczym punktem,
a nie oczekiwaną jednowymiarową podrozmaitością z brzegiem. Zastąpienie $w_i$
innymi podniesieniami zmienia je o elementy
$\ker\dd F_p$; macierz zmiany bazy ma wówczas
blok trójkątny z jednostką na przekątnej w części
normalnej, zatem nie zmienia znaku. Zmiana
dodatniej bazy w $T_yN$ ma dodatni wyznacznik
i także zachowuje znak. Reguła jest więc
jednoznaczna. W ilorazowym wzorze na $\eta_p$ zmiana
dowolnej bazy $(e_i)$ mnoży licznik i mianownik przez
ten sam niezerowy wyznacznik; zmiana dodatnich form
$\mu_p,\nu_y$ mnoży wynik przez liczbę dodatnią.
W lokalnej postaci submersji wynikającej z twierdzenia o lokalnej postaci submersji wybory normalnych
i stycznych są gładkie; znaki są zgodne na przecięciach.
To dowodzi orientowalności $S$ i wzoru.
\end{proof}

Ta konstrukcja działa też dla zorientowanej
\emph{wiązki normalnej} podrozmaitości $S\subset M$:
jej baza dodatnia, a następnie baza $TS$ mają razem
dawać bazę dodatnią $TM$. Wybór orientacji normalnej
nazywamy \emph{koorientacją}. Bez takiego wyboru sama
orientacja $M$ nie zawsze orientuje $S$.

\begin{przyklad}[Okrąg jako poziomica]
W $\R^2$ i $\R$ wybierzmy orientacje standardowe
i weźmy $F(x,y)=x^2+y^2$. Na $S^1=F^{-1}(1)$
różniczka ma wartość dodatnią na wektorze
skierowanym promieniowo na zewnątrz.
Twierdzenie~\ref{thm:orientacja-poziomicy} daje więc
orientację przeciwną do ruchu wskazówek zegara,
taką samą jak orientacja brzegu $D^2$.
\end{przyklad}

\begin{przyklad}[Co zawodzi na wstędze i przestrzeni rzutowej]
Wstęga Möbiusa (klasyczny przykład nieorientowalnej powierzchni)
po jednym okrążeniu odwraca
kierunek poprzeczny, zachowując kierunek
wzdłuż okręgu. Orientacja płaszczyzny stycznej
musiałaby więc zmienić znak i nie może być globalna.
Jej brzeg jest jednak okręgiem, a zatem orientowalnym.

Dla $n\geq1$ przestrzeń $\mathbb{RP}^n$ jest
orientowalna dokładnie wtedy, gdy $n$ jest nieparzyste.
Istotnie, nakrycie $S^n\to\mathbb{RP}^n$ identyfikuje
$x$ z $-x$. Antypodalne $A(x)=-x$ zmienia orientację
$S^n$ o znak $(-1)^{n+1}$: dodatnią bazę przestrzeni otaczającej
otrzymujemy z $(x,v_1,\ldots,v_n)$ w $\R^{n+1}$,
a $A$ mnoży wszystkie $n+1$ wektorów przez $-1$.
Orientacja na ilorazie istnieje wtedy i tylko wtedy,
gdy to utożsamienie zachowuje orientację sfery:
jeśli zachowuje, orientacja schodzi przez lokalne
dyfeomorfizmy; jeśli odwraca, jej podniesienie z
ilorazu byłoby zarazem zachowane i odwrócone, co
jest niemożliwe na spójnej sferze.
\end{przyklad}
\section{Całkowanie form na rozmaitości}
\label{sec:formy-calkowanie}

Na zorientowanej rozmaitości $M^m$ całkujemy formy
stopnia $m$: mają dokładnie tyle argumentów, ile
niezależnych kierunków ma $M$. Bez orientacji można
całkować \emph{gęstości}, ale wymagają osobnej definicji;
całka z formy najwyższego stopnia potrzebuje znaku.

\begin{definicja}[Całka z formy o zwartym nośniku]
\label{def:calka-z-formy}
Niech $M^m$ będzie zorientowana i
$\omega\in\Omega_c^m(M)$, gdzie dolny indeks $c$
oznacza zwarty nośnik. W mapie $(U,x)$ zgodnej z
orientacją, jeśli $\supp\omega\subset U$ i
$\omega=f\,\dd x^1\wedge\cdots\wedge\dd x^m$,
definiujemy
\begin{equation}\label{eq:calka-lokalna}
 \int_M\omega:=\int_{x(U)} f(x^{-1}(u))\,\dd u^1\cdots\dd u^m.
\end{equation}
Dla dowolnego zwartego nośnika bierzemy podział
jedności $(\rho_j)$ z Twierdzenia
\ref{thm:podzial-jednosci-formy}, w mapach zgodnych
z orientacją, i kładziemy
\begin{equation}\label{eq:calka-globalna}
 \int_M\omega:=\sum_j\int_M\rho_j\omega.
\end{equation}
Tylko skończenie wiele składników może być niezerowych:
lokalnie skończona rodzina nośników spotyka zwarty
$\supp\omega$ tylko dla skończenie wielu indeksów.
W wymiarze $0$ definicję zastępuje suma
$\int_M f=\sum_{p\in M}\varepsilon_p f(p)$,
gdzie $\varepsilon_p\in\{1,-1\}$ jest orientacją punktu.
Zwarty nośnik w przestrzeni dyskretnej jest skończony.
\end{definicja}

\begin{twierdzenie}[Niezależność od map i podziału]
\label{thm:calka-dobrze-okreslona}
Definicja~\ref{def:calka-z-formy} nie zależy od
wybranego dodatniego atlasu ani podziału jedności.
Jeżeli $F:M\to N$ jest dyfeomorfizmem zachowującym
orientację, to dla $\omega\in\Omega_c^m(N)$ zachodzi
$\int_M F^*\omega=\int_N\omega$; gdy $F$ odwraca
orientację i $M,N$ są spójne, po prawej pojawia się
znak minus.
\end{twierdzenie}
\begin{proof}
W wymiarze $0$ teza wynika z sumy ze znakami i z tego,
że dyfeomorfizm jest bijekcją punktów. Dla $m\geq1$
najpierw niech nośnik leży w dwóch dodatnich mapach
$x$ i $y$. Z równości
\[
 \dd y^1\wedge\cdots\wedge\dd y^m
 =\det D(y\circ x^{-1})\,
   \dd x^1\wedge\cdots\wedge\dd x^m
\]
wynika, że współczynniki opisujące tę samą formę
w obu mapach różnią się o wyznacznik Jacobiego.
Jest on dodatni. Twierdzenie o zmianie zmiennych
w całce na otwartych podzbiorach $\R^m$ pokazuje,
że oba wzory~\eqref{eq:calka-lokalna} dają tę samą
liczbę. Dla mapy odwracającej orientację wyznacznik
jest ujemny, a zwykły wzór na zmianę zmiennych
zawiera jego wartość bezwzględną; stąd minus.

Jeśli $(\rho_j)$ i $(\sigma_k)$ są dwoma podziałami,
to na nośniku $\omega$ możemy wstawić
$1=\sum_k\sigma_k$ do każdej całki $\rho_j\omega$.
Otrzymujemy sumę $\sum_{j,k}\int_M\rho_j\sigma_k\omega$.
Każdy składnik ma nośnik w przecięciu dwóch map.
Rozcinamy go w razie potrzeby na mniejsze mapy,
stosujemy poprzedni rachunek i zamieniamy kolejność
obu skończonych sum. Wynik jest taki sam, gdy
zaczniemy od $(\sigma_k)$. Wzór dla $F$ sprawdzamy
na składnikach lokalnych przez $y\circ F\circ x^{-1}$
i sumujemy.
\end{proof}

Na zwartej rozmaitości każda gładka forma ma zwarty
nośnik. Dla zorientowanej krzywej parametryzowanej
$\gamma:[a,b]\to M$ całka z $1$-formy jest
$\int_\gamma\alpha:=\int_a^b\gamma^*\alpha$;
przeparametryzowanie zachowujące kierunek nie zmienia
jej wartości. Dla parametryzowanej powierzchni
analogicznie cofamy formę $2$-stopnia.

\begin{przyklad}[Całka po okręgu]
Dla $\gamma(t)=(\cos t,\sin t)$, $0\leq t\leq2\pi$,
zorientowanej przeciwnie do ruchu wskazówek zegara,
forma $\alpha=x\,\dd y-y\,\dd x$ daje
$\gamma^*\alpha=(\cos^2t+\sin^2t)\dd t=\dd t$.
Stąd $\int_{S^1}\alpha=2\pi$.
Nie jest to długość łuku wyprowadzona z metryki:
całkujemy $1$-formę wybraną w otaczającej płaszczyźnie.
\end{przyklad}

\section{Twierdzenie Stokesa}
\label{sec:formy-stokes}

Intuicja jest jedna: gdy zsumujemy brzegi
małych, zgodnie zorientowanych fragmentów,
ich wewnętrzne krawędzie wystąpią po dwa razy
z przeciwnymi znakami. Pozostanie brzeg całości.
Rachunek na formach pozwala zapisać tę zasadę
bez ograniczania się do konkretnych wymiarów.

\begin{figure}[!htbp]
\centering
\begin{tikzpicture}[>=Stealth,line cap=round,line join=round]
 \fill[manifoldblue!10] (0,0) rectangle (3,2.35);
 \fill[manifoldgreen!12] (3,0) rectangle (6,2.35);
 \draw[manifoldblue!70!black,thick] (0,0) rectangle (3,2.35);
 \draw[manifoldgreen!70!black,thick] (3,0) rectangle (6,2.35);
 \draw[->,manifoldred,very thick] (2.92,.6)--(2.92,1.68);
 \draw[->,manifoldgold!85!black,very thick] (3.08,1.68)--(3.08,.6);
 \draw[->,manifoldblue!70!black,thick] (.65,0)--(1.75,0);
 \draw[->,manifoldgreen!70!black,thick] (4.25,0)--(5.35,0);
 \node[manifoldblue!80!black] at (1.5,1.25) {$P_1$};
 \node[manifoldgreen!70!black] at (4.5,1.25) {$P_2$};
 \node[font=\small,align=left,anchor=west,text width=5.0cm]
  at (6.55,1.2)
  {Wspólna krawędź jest brzegiem $P_1$
   i $P_2$ z przeciwnymi kierunkami.
   W sumie znika; zostaje tylko brzeg zewnętrzny.};
\end{tikzpicture}
\caption{Znak orientacji usuwa wewnętrzną krawędź
przy składaniu dwóch fragmentów.}
\label{fig:stokes-anulowanie}
\end{figure}

\begin{twierdzenie}[Stokes dla form]
\label{thm:stokes-formy}
Niech $M^m$ będzie zorientowaną gładką rozmaitością
z brzegiem, $m\geq1$, a $\omega\in\Omega_c^{m-1}(M)$.
Brzeg ma orientację z Definicji
\ref{def:orientacja-brzegu}, a
$i:\partial M\hookrightarrow M$ jest \emph{inkluzją}.
Wtedy
\begin{equation}\label{eq:stokes-formy}
 {\displaystyle \int_M\dd\omega
        =\int_{\partial M}i^*\omega.}
\end{equation}
Gdy $\partial M=\varnothing$, prawa strona wynosi $0$.
\end{twierdzenie}

Zapis $i^*\omega$ oznacza \emph{cofnięcie formy}
przez inkluzję $i$, a nie nową pochodną czy
ograniczenie samego współczynnika. Z definicji
cofnięcia~\eqref{eq:pullback-formy}, dla
$p\in\partial M$ i $v_j\in T_p\partial M$,
\begin{equation}\label{eq:inkluzja-formy}
 (i^*\omega)_p(v_1,\ldots,v_{m-1})
 =\omega_p(\dd i_pv_1,\ldots,\dd i_pv_{m-1}).
\end{equation}
Różniczka $\dd i_p$ po prostu traktuje wektor
styczny do brzegu jako wektor styczny do $M$.
W mapie przybrzegowej $x^m\geq0$ mamy
$i(x^1,\ldots,x^{m-1})=(x^1,\ldots,x^{m-1},0)$,
więc $i^*(\dd x^j)=\dd x^j$ dla $j<m$, natomiast
$i^*(\dd x^m)=0$. Współczynniki formy podstawiamy
przy $x^m=0$; składniki zawierające $\dd x^m$
znikają. Znak orientacji brzegu rozpatrujemy
\emph{osobno}: nie jest częścią operacji $i^*$.

\begin{proof}
\emph{Przypadek mapy wewnętrznej.}
Załóżmy najpierw, że nośnik $\omega$ jest zwarty
i zawarty w dziedzinie mapy, która nie spotyka
$\partial M$. We współrzędnych jej obraz jest
otwartym podzbiorem $\R^m$. Ponieważ nośnik
leży wewnątrz obrazu, przedłużenie współczynników
przez zero na całą $\R^m$ jest gładkie i ma
zwarty nośnik. Dla standardowej orientacji piszemy
\[
 \omega=\sum_{i=1}^m f_i\,
 \dd x^1\wedge\cdots\wedge\widehat{\dd x^i}
      \wedge\cdots\wedge\dd x^m.
\]
Symbol daszka oznacza opuszczenie jednego czynnika.
Ze wzoru na różniczkę zewnętrzną wyraz
$\partial_j f_i\,\dd x^j$ mnożymy przez wszystkie
$\dd x^r$ z $r\ne i$. Jeśli $j\ne i$, występuje
dwukrotnie $\dd x^j$ i wyraz znika. Jeśli $j=i$,
przestawiamy $\dd x^i$ przez $i-1$ czynników,
co daje znak $(-1)^{i-1}$. Zatem
\[
 \dd\omega=\sum_{i=1}^m(-1)^{i-1}
 (\partial_i f_i)\,\dd x^1\wedge\cdots\wedge\dd x^m.
\]
Ustalmy wszystkie współrzędne poza $x^i$.
Funkcja $x^i\mapsto f_i(x)$ znika przy
$x^i\to\pm\infty$. Podstawowe twierdzenie
rachunku całkowego daje
\[
 \int_{-\infty}^{\infty}\partial_i f_i(x)\,\dd x^i
 =\lim_{b\to\infty}f_i(x^1,\ldots,b,\ldots,x^m)
  -\lim_{a\to-\infty}f_i(x^1,\ldots,a,\ldots,x^m)=0.
\]
Zamieniamy kolejność całkowania dzięki zwartemu
nośnikowi i sumujemy po $i$. Wychodzi
$\int_M\dd\omega=0$. Także
$i^*\omega=0$, ponieważ nośnik $\omega$ nie
spotyka brzegu, więc równość Stokesa zachodzi.

\emph{Przypadek mapy przybrzegowej.}
Teraz niech nośnik leży w mapie na otwarty
podzbiór półprzestrzeni
$\mathbb H^m=\{(x',s):x'\in\R^{m-1},\ s\geq0\}$,
gdzie $s=x^m$. Na razie używamy standardowej
orientacji $\dd x^1\wedge\cdots\wedge\dd x^m$.
Przedłużamy $f_i$ przez zero w obrębie
$\mathbb H^m$ poza obraz mapy; zwarty nośnik
sprawia, że to przedłużenie jest gładkie.
Nie przedłużamy funkcji przez $s=0$ na $s<0$:
wartość na brzegu może przecież być niezerowa.
Wzór na $\dd\omega$ jest taki sam jak wyżej.
Jeśli $i<m$, to po ustaleniu $s$ oraz pozostałych
współrzędnych stycznych całka z $\partial_i f_i$
po całej prostej $x^i$ wynosi zero.
Pozostaje wyraz $i=m$. Dla każdego $x'$ funkcja
$s\mapsto f_m(x',s)$ znika dla dostatecznie
dużego $s$. Stąd
\[
 \int_0^\infty\partial_s f_m(x',s)\,\dd s
 =\lim_{s\to\infty}f_m(x',s)-f_m(x',0)
 =-f_m(x',0).
\]
Po uwzględnieniu znaku $(-1)^{m-1}$ we wzorze
na $\dd\omega$ otrzymujemy
\[
 \int_{\mathbb H^m}\dd\omega
 =(-1)^{m-1}\int_{\R^{m-1}}
  \left(\int_0^\infty\partial_s f_m(x',s)\,\dd s\right)
          \dd x'
 =(-1)^m\int_{\R^{m-1}} f_m(x',0)\,\dd x'.
\]
Dlaczego jest to całka po \emph{zorientowanym}
brzegu? Na $s=0$ normalna na zewnątrz to
$n=-\partial_s$. Dla standardowej formy
objętościowej $\mu=\dd x^1\wedge\cdots\wedge\dd x^m$
liczymy bez skrótu
\[
 \mu(-\partial_s,\partial_1,\ldots,\partial_{m-1})
 =-(-1)^{m-1}=(-1)^m.
\]
Zasada „normalna na zewnątrz jako pierwsza”
z Definicji~\ref{def:orientacja-brzegu} mówi więc,
że orientacja brzegu względem standardowego
$\dd x^1\wedge\cdots\wedge\dd x^{m-1}$
ma znak $(-1)^m$. Jednocześnie wzór
\eqref{eq:inkluzja-formy} daje
\[
 i^*\omega=f_m(x',0)\,
 \dd x^1\wedge\cdots\wedge\dd x^{m-1},
\]
bo tylko wyraz z opuszczonym $\dd x^m$
nie zawiera różniczki normalnej. Całka tej
formy z orientacją brzegu jest zatem równa
obliczonej całce z $\dd\omega$.
Dla $m=1$ rozumiemy $x'$ jako pusty układ
współrzędnych: brzeg w tym modelu składa się
z punktu $s=0$ o znaku $-1$, a wzór nadal działa.
Gdy wybrana mapa odwraca orientację $M$,
całka po mapie i indukowana orientacja jej brzegu
zmieniają znaki jednocześnie, więc równość
pozostaje prawdziwa.

\emph{Przejście do całej rozmaitości.}
Wybierzmy podział jedności $(\rho_j)$
podporządkowany mapom wewnętrznym i przybrzegowym,
tak że każdy $\rho_j$ ma zwarty nośnik w swojej
mapie (Twierdzenie~\ref{thm:podzial-jednosci-formy}).
Ponieważ rodzina jest lokalnie skończona, a
$\supp\omega$ zwarty, tylko skończenie wiele
iloczynów $\rho_j\omega$ jest niezerowych.
Każdy ma nośnik zwarty w odpowiedniej mapie,
więc należy do jednego z dwóch powyższych
przypadków. Mamy $\sum_j\rho_j\omega=\omega$;
stosując liniowość $\dd$ do tej \emph{skończonej}
sumy, dostajemy $\sum_j\dd(\rho_j\omega)=\dd\omega$.
Z drugiej strony cofnięcie przez inkluzję jest
liniowe, a $i^*(\rho_j\omega)
=(\rho_j|_{\partial M})\,i^*\omega$.
Sumując lokalne równości Stokesa, otrzymujemy
\[
 \int_M\dd\omega
 =\sum_j\int_M\dd(\rho_j\omega)
 =\sum_j\int_{\partial M}i^*(\rho_j\omega)
 =\int_{\partial M}i^*\omega,
\]
ponieważ $\sum_j\rho_j|_{\partial M}=1$.
\end{proof}

\begin{przyklad}[Cofnięcie $1$-formy na brzeg dysku]
Niech $M=D^2$, $\omega=P(x,y)\,\dd x+Q(x,y)\,\dd y$
i niech $\gamma(t)=(\cos t,\sin t)$ parametryzuje
dodatnio zorientowany brzeg. Ponieważ
$\dd x(\gamma'(t))=-\sin t$ oraz
$\dd y(\gamma'(t))=\cos t$, dostajemy
\[
 \gamma^*(i^*\omega)
 =(i\circ\gamma)^*\omega
 =\bigl[-P(\cos t,\sin t)\sin t
       +Q(\cos t,\sin t)\cos t\bigr]\,\dd t.
\]
Tak obliczamy prawą stronę twierdzenia Stokesa
jako zwykłą całkę po $t\in[0,2\pi]$.
\end{przyklad}

\begin{wniosek}[Podstawowe twierdzenie rachunku całkowego]
\label{cor:stokes-ftc}
Jeśli $f$ jest gładka na $[a,b]$, to
\[
 \int_a^b f'(t)\,\dd t=f(b)-f(a).
\]
\end{wniosek}
\begin{proof}
Na zorientowanym przedziale $[a,b]$ forma $\dd f$
ma postać $f'(t)\dd t$. Jego brzeg ma orientację
$\partial[a,b]=\{b\}-\{a\}$, co sprawdziliśmy
w Stwierdzeniu~\ref{prop:brzeg-orientowalny}.
Całka $0$-formy $f$ po tym zorientowanym
zbiorze punktów jest z definicji $f(b)-f(a)$.
Wzór~\eqref{eq:stokes-formy} daje tezę.
\end{proof}
To identyfikacja jednowymiarowego przypadku Stokesa.
Podstawowego twierdzenia rachunku całkowego użyliśmy
już w jego dowodzie lokalnym, więc powyższy rachunek
nie stanowi niezależnego dowodu tego faktu analizy.

\begin{wniosek}[Twierdzenie Greena]
\label{cor:stokes-green}
Niech $D\subset\R^2$ będzie zwartym obszarem
o gładkim brzegu, z orientacją $\dd x\wedge\dd y$,
a $P,Q$ niech będą gładkie w otoczeniu $D$. Wtedy
\[
 \oint_{\partial D}P\,\dd x+Q\,\dd y
 =\iint_D\left(\frac{\partial Q}{\partial x}
             -\frac{\partial P}{\partial y}\right)
       \dd x\,\dd y.
\]
Orientacja $\partial D$ jest dodatnia, czyli
podczas obiegania każdej składowej brzegu
wnętrze $D$ pozostaje po lewej stronie.
\end{wniosek}
\begin{proof}
Dla $\alpha=P\,\dd x+Q\,\dd y$ reguła
Leibniza i $\dd x\wedge\dd x
=\dd y\wedge\dd y=0$ dają
\[
 \dd\alpha
 =P_y\,\dd y\wedge\dd x
  +Q_x\,\dd x\wedge\dd y
 =(Q_x-P_y)\,\dd x\wedge\dd y.
\]
Lewa strona wzoru Greena jest całką z
$i^*\alpha$ po brzegu; cofnięcie ogranicza
$\alpha$ do kierunków stycznych do krzywej,
jak w powyższym przykładzie z dyskiem. Podstawiamy $\alpha$ do
Twierdzenia~\ref{thm:stokes-formy}.
Znak orientacji sprawdzamy w punkcie
$(1,0)$ brzegu dodatnio zorientowanego dysku:
normalna zewnętrzna wskazuje w prawo, więc
dodatni wektor styczny wskazuje w górę.
Ta sama reguła działa na każdej składowej
brzegu, także na brzegu otworu.
\end{proof}

\begin{wniosek}[Klasyczne twierdzenie Stokesa o rotacji]
\label{cor:stokes-curl}
Niech $S\subset\R^3$ będzie zwartą gładką
zorientowaną powierzchnią z brzegiem,
$n$ jej dodatnią normalną jednostkową,
a $\mathbf F=(P,Q,R)$ gładkim polem w otoczeniu $S$.
Dla dodatniego kierunku $\partial S$ mamy
\[
 \oint_{\partial S}\mathbf F\cdot\dd\mathbf r
 =\iint_S(\nabla\times\mathbf F)\cdot n\,\dd A.
\]
\end{wniosek}
\begin{proof}
Bierzemy $1$-formę
$\alpha=P\,\dd x+Q\,\dd y+R\,\dd z$.
Bezpośredni rachunek pochodnych daje
\[
 \begin{aligned}
 \dd\alpha={}&(R_y-Q_z)\,\dd y\wedge\dd z\\
 &+(P_z-R_x)\,\dd z\wedge\dd x
 +(Q_x-P_y)\,\dd x\wedge\dd y.
 \end{aligned}
\]
Wektor współczynników przy trzech formach
bazowych to dokładnie
$\nabla\times\mathbf F
=(R_y-Q_z,P_z-R_x,Q_x-P_y)$.
Dla dodatniej parametryzacji
$r(u,v)$ fragmentu $S$ podstawmy
$a=r_u$, $b=r_v$. Z definicji iloczynu
zewnętrznego
\[
 (\dd\alpha)(a,b)
 =(\nabla\times\mathbf F)\cdot(a\times b)
 =(\nabla\times\mathbf F)\cdot n\,
    \lVert a\times b\rVert.
\]
Po scałkowaniu po $(u,v)$ otrzymujemy prawą
stronę tezy. Na krzywej brzegowej
$\alpha(\gamma')=\mathbf F(\gamma)\cdot\gamma'$,
więc całka z $i^*\alpha$ jest lewą stroną.
Twierdzenie~\ref{thm:stokes-formy} utożsamia
te dwie całki. Kierunek brzegu jest
indukowany przez orientację $S$; na płaskim
dysku z normalną w górę jest przeciwny
do ruchu wskazówek zegara.
\end{proof}

\begin{wniosek}[Twierdzenie Gaussa--Ostrogradskiego]
\label{cor:stokes-gauss}
Niech $D\subset\R^3$ będzie zwartym obszarem
o gładkim brzegu, a $\mathbf F=(P,Q,R)$
gładkim polem w otoczeniu $D$.
Jeśli $n$ oznacza normalną jednostkową
na zewnątrz, to
\[
 \iint_{\partial D}\mathbf F\cdot n\,\dd A
 =\iiint_D\operatorname{div}\mathbf F\,
                  \dd x\,\dd y\,\dd z.
\]
\end{wniosek}
\begin{proof}
Strumieniowi $\mathbf F$ odpowiada $2$-forma
\[
 \beta=P\,\dd y\wedge\dd z
      +Q\,\dd z\wedge\dd x
      +R\,\dd x\wedge\dd y.
\]
Różniczkując trzy składniki, dostajemy
\[
 \dd\beta=(P_x+Q_y+R_z)\,
           \dd x\wedge\dd y\wedge\dd z
          =(\operatorname{div}\mathbf F)\,
           \dd x\wedge\dd y\wedge\dd z.
\]
Dla dodatnio zorientowanej pary stycznych
wektorów $a,b\in T_p\partial D$ mamy
\[
 (i^*\beta)_p(a,b)
 =\mathbf F(p)\cdot(a\times b).
\]
Z reguły orientowania brzegu wynika, że
$a\times b$ wskazuje na zewnątrz; dla
ortonormalnej pary jest równy $n$.
Ogólnie jego długość dostarcza czynnika
elementu pola $\dd A$. Całka z $i^*\beta$
jest więc strumieniem przez brzeg, a całka
z $\dd\beta$ całką z dywergencji po $D$.
Równość daje Twierdzenie~\ref{thm:stokes-formy}.
\end{proof}

\begin{wniosek}[Zerowanie całki form zamkniętych i dokładnych]
\label{cor:stokes-top-form}
Jeśli $M^m$ jest zwartą zorientowaną rozmaitością
z brzegiem, a $\eta\in\Omega^{m-1}(M)$ jest
zamknięta, to $\int_{\partial M}i^*\eta=0$.
Jeżeli $M$ jest zwartą zorientowaną rozmaitością
bez brzegu i $\omega=\dd\eta\in\Omega^m(M)$,
to $\int_M\omega=0$. W szczególności forma
najwyższego stopnia o całce niezerowej nie jest
dokładna.
\end{wniosek}
\begin{proof}
W pierwszej sytuacji
$\int_{\partial M}i^*\eta=\int_M\dd\eta=0$.
W drugiej
$\int_M\omega=\int_{\partial M}i^*\eta=0$,
bo brzeg jest pusty. Oba kroki są
bezpośrednimi zastosowaniami~\eqref{eq:stokes-formy}.
\end{proof}
\section{Kohomologia de Rhama i ciąg Mayera--Vietorisa}
\label{sec:de-rham-mv}

\begin{definicja}[Kohomologia de Rhama]
\label{def:de-rham-cohomology}
\emph{Kompleks de Rhama} to ciąg
\[
 0\longrightarrow\Omega^0(M)\xrightarrow{\dd}
 \Omega^1(M)\xrightarrow{\dd}\cdots
 \xrightarrow{\dd}\Omega^m(M)\longrightarrow0.
\]
Jest kompleksem kołańcuchowym, bo $\dd^2=0$.
Formę z $\ker\dd$ nazywamy \emph{zamkniętą},
a formę z $\operatorname{im}\dd$ \emph{dokładną}. Definiujemy
\[
 H^k_{\mathrm{dR}}(M)
 :=\ker(\dd:\Omega^k\to\Omega^{k+1})/
   \operatorname{im}(\dd:\Omega^{k-1}\to\Omega^k).
\]
W stopniu $0$ jest to przestrzeń funkcji lokalnie stałych.
Z uwagi na brak form ujemnego stopnia zerowa funkcja
jest jedyną dokładną $0$-formą.
\end{definicja}

Iloczyn zewnętrzny przechodzi na klasy: jeśli
$\dd\alpha=\dd\beta=0$ i $\alpha$ zastąpimy przez
$\alpha+\dd\eta$, to różnica iloczynów wynosi
$(\dd\eta)\wedge\beta=\dd(\eta\wedge\beta)$.
Otrzymujemy \emph{pierścień kohomologii de Rhama}
$H^*_{\mathrm{dR}}(M)$ z przemiennością
$[\alpha][\beta]=(-1)^{k\ell}[\beta][\alpha]$.
Cofnięcie $F^*$ daje homomorfizm pierścieni
$H^*_{\mathrm{dR}}(N)\to H^*_{\mathrm{dR}}(M)$;
Wniosek~\ref{cor:de-rham-homotopia} pokazuje, że
gładkie homotopie dają tę samą operację.

\begin{twierdzenie}[Mayer--Vietoris dla form]
\label{thm:mv-de-rham}
Jeśli $M=U\cup V$ dla zbiorów otwartych, to
w każdym stopniu krótki ciąg
\begin{equation}\label{eq:mv-formy-kompleks}
\begin{gathered}
0\to\Omega^k(M)\xrightarrow{\ r\ }
\Omega^k(U)\oplus\Omega^k(V)
\xrightarrow{\ s\ }\Omega^k(U\cap V)\to0,\\
r(\omega)=(\omega|_U,\omega|_V),\qquad
s(\alpha,\beta)=\beta|_{U\cap V}-\alpha|_{U\cap V}.
\end{gathered}
\end{equation}
jest dokładny i zgodny z $\dd$. W konsekwencji
istnieje długi ciąg dokładny
\[
 \cdots\to H^{k-1}_{\mathrm{dR}}(U\cap V)
 \xrightarrow{\delta}H^k_{\mathrm{dR}}(M)
 \to H^k_{\mathrm{dR}}(U)\oplus H^k_{\mathrm{dR}}(V)
 \to H^k_{\mathrm{dR}}(U\cap V)
 \xrightarrow{\delta}H^{k+1}_{\mathrm{dR}}(M)\to\cdots.
\]
\end{twierdzenie}
\begin{proof}
Jeżeli $r(\omega)=0$, forma znika na $U$ i $V$,
więc znika na $M$: $r$ jest iniektywne.
Para $(\alpha,\beta)$ należy do $\ker s$ dokładnie
wtedy, gdy jej składniki zgadzają się na przecięciu;
wówczas sklejają się w jedną gładką formę na $M$.
To dowodzi $\operatorname{im}r=\ker s$.

Aby dowieść surjektywności $s$, weźmy
$\gamma\in\Omega^k(U\cap V)$ i podział jedności
$\rho_U+\rho_V=1$ podporządkowany $(U,V)$.
Na $U$ kładziemy $\alpha=-\rho_V\gamma$ na
$U\cap V$ i przedłużamy przez $0$ na $U\setminus V$;
na $V$ kładziemy $\beta=\rho_U\gamma$
i podobnie przedłużamy przez $0$.
Przedłużenia są gładkie: nośnik $\rho_V$ nie
dochodzi lokalnie do punktów poza $V$, a nośnik
$\rho_U$ do punktów poza $U$. Na przecięciu
$s(\alpha,\beta)=(\rho_U+\rho_V)\gamma=\gamma$.
Ponieważ obcięcie komutuje z $\dd$, otrzymujemy
krótki ciąg dokładny \emph{kompleksów}.
Twierdzenie~\ref{thm:les-kohomologia} z
Rozdziału~\ref{ch:algebra-homologiczna} tworzy
z niego długi ciąg. Konkretnie, dla zamkniętej
$\gamma$ wybieramy powyższe $(\alpha,\beta)$;
formy $\dd\alpha$ i $\dd\beta$ zgadzają się
na przecięciu, bo $\dd\gamma=0$. Sklejona
forma $\eta$ na $M$ jest zamknięta i
$\delta[\gamma]=[\eta]$. Zmiana podniesienia
lub reprezentanta $[\gamma]$ zmienia $\eta$
o formę dokładną, dokładnie jak w dowodzie
długiego ciągu kohomologii.
\end{proof}

\begin{przyklad}[Okrąg i forma kąta]
\label{prz:de-rham-s1}
Na $\R^2\setminus\{0\}$ forma
\[
 \vartheta=\frac{-y\,\dd x+x\,\dd y}{x^2+y^2}
\]
jest zamknięta: po zastosowaniu wzoru na $\dd$
otrzymujemy $\dd\vartheta=0$. Na okręgu
$\gamma(t)=(\cos t,\sin t)$ mamy
$\gamma^*\vartheta=\dd t$ i
$\int_{S^1}\vartheta=2\pi$. Gdyby $\vartheta$
była dokładna, Stokes na zamkniętej krzywej
dałby $0$, sprzeczność. Lokalnie możemy
pisać $\vartheta=\dd\theta$, ale funkcja kąta
po pełnym obrocie zmienia się o $2\pi$ i nie
jest globalną funkcją rzeczywistą.

W istocie $H^1_{\mathrm{dR}}(S^1)\cong\R$,
z izomorfizmem $[\alpha]\mapsto\int_{S^1}\alpha$.
Dowód surjektywności daje $c\vartheta$.
Jeśli całka z $\alpha$ znika, zapiszmy
$\gamma^*\alpha=g(t)\dd t$ dla funkcji
$2\pi$-okresowej $g$ i połóżmy
$h(t)=\int_0^t g(s)\,\dd s$. Wtedy
$h(t+2\pi)-h(t)=\int_t^{t+2\pi}g(s)\,\dd s
=\int_0^{2\pi}g(s)\,\dd s=0$ dla każdego $t$.
Funkcja $h$ jest zatem gładka i okresowa, więc schodzi
do gładkiej funkcji na $S^1$, której różniczka jest $\alpha$.
To dowodzi iniektywności. Ponadto
$H^0_{\mathrm{dR}}(S^1)=\R$ i wyższe grupy znikają.
\end{przyklad}

\begin{figure}[!htbp]
\centering
\begin{tikzpicture}[>=Stealth,line cap=round,line join=round]
 \fill[manifoldblue!9] (0,0) circle(2.0);
 \draw[manifoldblue!65!black,thick] (0,0) circle(2.0);
 \draw[manifoldgold!85!black,line width=1.4pt] (0,0) circle(1.35);
 \foreach \ang in {5,50,95,140,185,230,275,320}{
  \draw[->,manifoldred,thick]
   (\ang:1.35) arc[start angle=\ang,end angle={\ang+24},radius=1.35];
 }
 \fill[white] (0,0) circle(2.4pt);
 \draw[manifoldred,thick] (-.09,-.09)--(.09,.09)
  (-.09,.09)--(.09,-.09);
 \node[font=\small,align=left,anchor=west,text width=5.7cm]
  at (2.65,.5)
  {$\dd\vartheta=0$ poza środkiem,
   ale $\int_{S^1}\vartheta=2\pi$.
   Funkcja kąta jest pierwotną tylko lokalnie:
   po obejściu dziury zmienia się o $2\pi$.};
\end{tikzpicture}
\caption{Zamknięta $1$-forma na płaszczyźnie bez początku
może nie mieć globalnej funkcji pierwotnej.
Krzyżykiem zaznaczono jedyny usunięty punkt; niebieski
dysk przedstawia fragment dziedziny.}
\label{fig:de-rham-kat}
\end{figure}

\begin{przyklad}[Sfery i torus]
\label{prz:de-rham-sfery-torus}
Dwie dziedziny map na $S^n$ otrzymane przez usunięcie
przeciwnych biegunów są ściągalne, a ich przecięcie
zwija się na $S^{n-1}$. Długi ciąg z
Twierdzenia~\ref{thm:mv-de-rham}, zaczynając od
obliczenia dla $S^1$, daje indukcyjnie
\[
 H^k_{\mathrm{dR}}(S^n)=
 \begin{cases}\R,&k=0,n,\\0,&\text{w pozostałych stopniach,}\end{cases}
 \qquad n\geq1.
\]
Dla $n\geq2$ dokładność w stopniu $0$ mówi,
że $H^0(S^n)=\R$ i $H^1(S^n)=0$:
przecięcie jest spójne, a mapa
$H^0(U)\oplus H^0(V)=\R^2\to H^0(U\cap V)=\R$
ma postać $(a,b)\mapsto b-a$, więc jest surjektywna;
w stopniach $k\geq2$ odwzorowanie łączące
przenosi $H^{k-1}(S^{n-1})$ na $H^k(S^n)$.
Jedyny niezerowy stopień dodatni to $k=n$.

Na $T^2=S^1\times S^1$ niech
$\alpha=\pr_1^*\vartheta$ i
$\beta=\pr_2^*\vartheta$.
Ich całki po dwóch podstawowych okręgach
wynoszą odpowiednio $(2\pi,0)$ i $(0,2\pi)$;
$\int_{T^2}\alpha\wedge\beta=(2\pi)^2$.
Klasy $[\alpha],[\beta]$ są liniowo niezależne w stopniu $1$,
a $[\alpha\wedge\beta]$ jest niezerowa w stopniu $2$.
W obu przypadkach wynika to ze Stokesa: forma dokładna
ma zerową całkę po każdym z tych cykli.
Twierdzenie de Rhama poniżej i kompleks komórkowy
torusa z Przykładu~\ref{prz:homologia-cw-torus}
pokazują, że dają całe
$H^1_{\mathrm{dR}}(T^2)\cong\R^2$ oraz
$H^2_{\mathrm{dR}}(T^2)\cong\R$.
\end{przyklad}

\section{Dlaczego formy obliczają kohomologię przestrzeni?}
\label{sec:twierdzenie-de-rhama}

Forma zamknięta może nie mieć pierwotnej właśnie wtedy,
gdy wykrywa globalny „otwór”. Aby porównać ten język
z homologią singularną z Rozdziału~\ref{ch:homologia-przestrzeni},
całkujemy formę po zorientowanych sympleksach.
Używamy tu współczynników rzeczywistych: całka przyjmuje
wartości w $\R$, a torsja całkowitoliczbowa znika po
przejściu do tych współczynników.

\begin{definicja}[Kohomologia singularna z wartościami w $\R$]
\label{def:singular-cohom-r}
Jeżeli $C_k(M;\R)$ jest przestrzenią skończonych
$\R$-liniowych kombinacji ciągłych sympleksów
$\sigma:\Delta^k\to M$, to
$C^k(M;\R)=\operatorname{Hom}_{\R}(C_k(M;\R),\R)$.
Kobrzeg $\delta:C^k\to C^{k+1}$ określa wzór
$(\delta c)(\sigma)=c(\partial\sigma)$.
Z $\partial^2=0$ wynika $\delta^2=0$; iloraz
$\ker\delta/\operatorname{im}\delta$ oznaczamy przez
$H^k_{\mathrm{sing}}(M;\R)$.
\end{definicja}

Sympleks $\sigma:\Delta^k\to M$ nazywamy \emph{gładkim},
gdy w mapach jego składowe mają gładkie przedłużenia
do otoczeń punktów $\Delta^k$ w jego otoczce afinicznej.
Przy brzegu $M$ przedłużamy składowe do $\R^m$;
nie wymagamy, aby przedłużenie poza sympleks przyjmowało
wartości w półprzestrzeni modelującej $M$.
Na przykład $t\mapsto t$ z $[0,1]$ do $[0,\infty)$
jest w tym sensie gładkie. Dla takiego $\sigma$ definiujemy
$I(\omega)(\sigma)=\int_{\Delta^k}\sigma^*\omega$.
Znak orientacji $\Delta^k$ jest określony przez
kolejność jego wierzchołków. Na sympleksie
$\Delta^{k+1}$ twierdzenie Stokesa daje
\begin{equation}\label{eq:integracja-lancuchowa}
 \begin{aligned}
 (\delta I\omega)(\sigma)
 &= I\omega(\partial\sigma)
  =\sum_{j=0}^{k+1}(-1)^j
      \int_{[v_0,\ldots,\widehat v_j,\ldots,v_{k+1}]}
          \sigma^*\omega \\
 &=\int_{\Delta^{k+1}}\dd(\sigma^*\omega)
  =\int_{\Delta^{k+1}}\sigma^*(\dd\omega)
  =I(\dd\omega)(\sigma).
 \end{aligned}
\end{equation}
Uzasadnijmy użycie Stokesa na sympleksie, który ma naroża.
W otoczeniu punktu należącego do $r$ ścian jego dziedzinę
opisują współrzędne
$[0,\infty)^r\times\R^{k+1-r}$: równania przecinających
się ścian są niezależne. Rozkładamy formę podziałem
jedności i stosujemy lokalny dowód Stokesa do każdego
składnika o zwartym nośniku. Dla składnika z pominiętym
$\dd x^i$ całkowanie po $x^i$ daje zero, gdy $x^i\in\R$,
a gdy $x^i\geq0$, daje ujemną wartość współczynnika
na $x^i=0$. Są to dokładnie wkłady ścian z orientacją
„normalna na zewnątrz najpierw”. Przecięcia dwóch ścian
mają miarę zero w każdej ścianie i nie tworzą dodatkowych
całek. Dla orientacji sympleksu zadanej kolejnością
wierzchołków orientacja ściany naprzeciw $v_j$ jest
$(-1)^j$ razy orientacja pozostałych wierzchołków:
przesunięcie pominiętego wierzchołka na pierwsze miejsce
wymaga $j$ zamian, a przypadek $j=0$ daje orientację
brzegową dodatnią. Otrzymujemy więc dokładnie znaki
w \eqref{eq:integracja-lancuchowa}, bez zastępowania
naroży brzegiem gładkim.
Równanie mówi, że „całka po brzegu” i „różniczka
formy” opisują ten sam rachunek.

\begin{twierdzenie}[de Rham]
\label{thm:de-rham}
Dla każdej gładkiej rozmaitości $M$ (również z brzegiem)
całkowanie po gładkich sympleksach indukuje naturalny
izomorfizm przestrzeni wektorowych
\begin{equation}\label{eq:de-rham-iso}
 H^k_{\mathrm{dR}}(M)\xrightarrow{\ \cong\ }
 H^k_{\mathrm{sing}}(M;\R),
 \qquad [\omega]\longmapsto
 \bigl[\sigma\longmapsto\textstyle\int_\sigma\omega\bigr].
\end{equation}
Zapis po prawej najpierw określa kołańcuch na
\emph{gładkich} sympleksach. Przenosimy jego klasę do
zwykłej kohomologii singularnej przez odwrotność
kanonicznego izomorfizmu ograniczenia kołańcuchów,
którego istnienie udowodnimy poniżej. Nie wybieramy
niezależnie wygładzenia każdego ciągłego sympleksu:
takie wybory nie musiałyby być zgodne z jego ścianami.
\end{twierdzenie}
\begin{proof}
Dowód rozdzielimy na rachunek lokalny i precyzyjną
zasadę sklejania kompleksów. Równocześnie uzasadnimy,
dlaczego wystarczają gładkie sympleksy. Nie będziemy
potrzebować ani triangulacji $M$, ani metryki.

\emph{1. Trzy kompleksy i prostokąty wypukłe.}
Niech $C_*^\infty(U;\R)$ oznacza podkompleks
generowany przez gładkie sympleksy, a
$C_\infty^*(U;\R)=\operatorname{Hom}_{\R}
(C_*^\infty(U;\R),\R)$ jego kompleks dualny.
Ograniczenie kołańcuchów i całkowanie dają naturalne
odwzorowania kompleksów
\[
 C^*(U;\R)\xrightarrow{\ r\ }C_\infty^*(U;\R)
 \xleftarrow{\ I\ }\Omega^*(U).
\]
Pierwsze komutuje z kobrzegiem, ponieważ ściany
gładkiego sympleksu są gładkie; drugie dzięki
\eqref{eq:integracja-lancuchowa}.

Gdy $U$ jest niepustym wypukłym zbiorem otwartym
w $\R^m$ lub względnie otwartym w
$\mathbb H^m=\{x^m\geq0\}$, wybierzmy $a\in U$.
Homotopia $H(t,x)=(1-t)a+tx$ pozostaje w $U$.
Lemat~\ref{lem:homotopia-form} oblicza kohomologię
form, a operator pryzmatowy z Twierdzenia
\ref{thm:homotopia-singularna} oblicza obie kohomologie
singularne. Pryzmat i jego triangulacja są określone
afinicznie w dziedzinie, więc zachowują gładkość
sympleksów w przyjętym sensie. Każdy z trzech
kompleksów ma zatem kohomologię $\R$ w stopniu $0$
i zero w stopniach dodatnich. W stopniu $0$ oba
odwzorowania są identycznością na funkcjach stałych.
W szczególności $r$ i $I$ indukują izomorfizmy.

\emph{2. Algebraiczna zasada sklejania.}
Dla otwartego pokrycia $(U_i)$ rozważmy trzy tabele
\[
 \begin{aligned}
 A^{p,q}&=\prod_{i_0<\cdots<i_p}\Omega^q(U_{i_0\cdots i_p}),\\
 B^{p,q}&=\prod_{i_0<\cdots<i_p}C_\infty^q(U_{i_0\cdots i_p};\R),\\
 E^{p,q}&=\prod_{i_0<\cdots<i_p}C^q(U_{i_0\cdots i_p};\R),
 \qquad U_{i_0\cdots i_p}=U_{i_0}\cap\cdots\cap U_{i_p}.
 \end{aligned}
\]
Puste przecięcia wnoszą zera. W poziomie działa
naprzemienna suma ograniczeń $\delta_{\mathcal U}$,
np. $(\delta_{\mathcal U}a)_{ij}=a_j|_{U_{ij}}-a_i|_{U_{ij}}$.
W pionie działa odpowiednio $\dd$ lub kobrzeg singularny
$d_v$. Kompleks całkowity ma w stopniu $n$ sumę
składników $p+q=n$ i różniczkę
$D=\delta_{\mathcal U}+(-1)^p d_v$ na kolumnie $p$.
Ograniczenia komutują z $d_v$, więc $D^2=0$.

Użyjemy dwóch elementarnych własności takich tabel.
Jeśli wszystkie kolumny są dokładne, kompleks całkowity
jest dokładny. Istotnie, dla kocyklu $z$ jego składnik
w kolumnie o najmniejszym indeksie $p$ jest pionowo
zamknięty. Z dokładności wybieramy jego pionową pierwotną,
uwzględniając znak $(-1)^p$, i odejmujemy jej $D$.
Usuwamy tę kolumnę, ewentualnie zmieniając następną.
Powtarzamy to, aż usuniemy wszystkie składniki.
Procedura kończy się: w danym stopniu całkowitym jest
skończenie wiele par $(p,q)$. W najniższym pionowym
stopniu dokładność oznacza zerowe jądro, jeśli brak
już stopnia na pierwotną. Ten sam argument, z zamianą
ról indeksów, działa dla dokładnych wierszy.

Wynika stąd również zasada porównania: odwzorowanie
tabel, które indukuje izomorfizm kohomologii w każdej
kolumnie, indukuje izomorfizm kohomologii całkowitej.
Aby to sprawdzić, tworzymy stożek odwzorowania $f$.
Dla kompleksów $P,Q$ ma on składniki
$Q^n\oplus P^{n+1}$ i różniczkę
$(b,a)\mapsto(d_Qb+f(a),-d_Pa)$.
Z krótkiego ciągu dokładnego
$0\to Q\to\operatorname{Cone}(f)\to P[1]\to0$
i jego ciągu kohomologii wynika, że stożek jest dokładny
wtedy i tylko wtedy, gdy $f$ indukuje izomorfizmy
kohomologii. Dla odwzorowania tabel $f:P\to Q$ bierzemy
w pozycji $(p,q)$ przestrzeń $Q^{p,q}\oplus P^{p,q+1}$,
z różniczką pionową
$(b,a)\mapsto(d_vb+(-1)^p f(a),-d_va)$
i poziomą $(b,a)\mapsto(\delta_{\mathcal U}b,-\delta_{\mathcal U}a)$.
Te różniczki komutują, a po zastosowaniu znaku $(-1)^p$
w totalizacji otrzymujemy dokładnie stożek odwzorowania
kompleksów całkowitych. Kolumna $p$ jest stożkiem
$(-1)^p f$, więc jest dokładna na mocy założenia.
Stosujemy teraz poprzedni argument usuwania kolumn.
Dopuszczony dodatkowy stopień pionowy $-1$ nie przeszkadza:
przekątne nadal mają skończoną długość. Analogicznie,
jeśli do dokładnych wierszy dopiszemy w kolumnie $-1$
kompleks globalny $P$, to jego odwzorowanie do kompleksu
całkowitego $Q$ jest izomorfizmem kohomologii.
Rozszerzona tabela w stopniu całkowitym $n$ ma składniki
$Q^n\oplus P^{n+1}$. W kolumnie $-1$ pionowa różniczka
ma znak minus, a pozioma jest augmentacją $f$;
różniczka całkowita to zatem
$(b,a)\mapsto(D_Qb+f(a),-d_Pa)$.
Jest to dokładnie stożek tego odwzorowania, a jego
dokładne wiersze pozwalają zastosować ten sam argument.

\emph{3. Dlaczego wiersze obliczają kompleksy globalne.}
W przypadku $A$ dopisujemy $\Omega^q(M)$ na początku
wiersza. Wybierzmy gładki podział jedności
$\sum_j\rho_j=1$ podporządkowany pokryciu. Dla rodziny
naprzemiennej określamy
\[
 (ha)_{i_0\cdots i_{p-1}}
   =\sum_j\rho_j a_{j i_0\cdots i_{p-1}}.
\]
Składniki przedłużamy przez zero; lokalna skończoność
nośników zapewnia gładkość sumy. Dla $p=0$ otrzymujemy
globalną formę $\sum_j\rho_j a_j$.
Rozwinięcie sumy daje
$h\delta_{\mathcal U}+\delta_{\mathcal U}h=\id$:
wyrazy z pominiętym $i_r$ znoszą się parami, a pozostałe
dają $(\sum_j\rho_j)a$. Na formach globalnych złożenie
ograniczenia i $h$ także jest identycznością.
Rozszerzone wiersze są więc dokładne.

Dla $B$ oraz $E$ dopisujemy na początku kołańcuchy
na \emph{małych} sympleksach, czyli na sympleksach
o obrazie zawartym w jednym z $U_i$. Dla konkretnego
sympleksu $\sigma$ zbiór
$J_\sigma=\{i:\sigma(\Delta^q)\subset U_i\}$ jest
niepusty. Poziomy ciąg jego wartości to rozszerzony
kompleks kołańcuchów pełnego sympleksu o wierzchołkach
$J_\sigma$. Wybranie $i(\sigma)\in J_\sigma$ i wstawienie
go na początek indeksów daje homotopię $h$ spełniającą
tę samą tożsamość. Wiersz jest iloczynem takich ciągów
po wszystkich małych sympleksach; pozostaje dokładny,
ponieważ pierwotne wybieramy współrzędna po współrzędnej.

Lemat~\ref{lem:male-lancuchy} mówi, że włączenie małych
łańcuchów indukuje izomorfizm homologii. Jego dowód
działa też dla gładkich łańcuchów: podziały i homotopie
podziału są złożeniami z odwzorowaniami afinicznymi.
Nad $\R$ ograniczenie dualnych kołańcuchów indukuje
zatem izomorfizm kohomologii. Uzasadnienie algebraiczne
nie wymaga skończonego wymiaru: wybieramy dopełnienia
brzegów w cyklach i cykli w łańcuchach; kompleks jest
sumą części o zerowej różniczce (homologii) i par,
na których różniczka jest izomorfizmem. Dualizacja
zachowuje ten rozkład i dokładność części z parami.
W konsekwencji tabele $A,B,E$ obliczają odpowiednio
$H^*_{\mathrm{dR}}(M)$, $H^*(C_\infty^*(M;\R))$
i $H^*_{\mathrm{sing}}(M;\R)$.
Wszystkie ograniczenia i augmentacje komutują z $I$ i $r$.

\emph{4. Dwa zastosowania sklejania.}
Najpierw niech $W$ będzie otwartym podzbiorem $\R^m$
lub względnie otwartym podzbiorem $\mathbb H^m$.
Pokrywamy go prostokątami otwartymi, odpowiednio
przeciętymi z $\mathbb H^m$, i zawartymi w $W$.
Możemy wybrać pokrycie lokalnie skończone: w konstrukcji
wyczerpania zwartymi zbiorami z dowodu Twierdzenia
\ref{thm:podzial-jednosci-formy} zastępujemy małe kule
małymi prostokątami. Każde niepuste skończone przecięcie
jest wypukłe. Krok 1 pokazuje, że $I:A\to B$ oraz
$r:E\to B$ indukują izomorfizmy w każdej kolumnie.
Iloczyny przestrzeni wektorowych zachowują tu jądra,
obrazy i dokładność, więc także iloczyn lokalnych
izomorfizmów jest izomorfizmem. Kroki 2 i 3 dowodzą
obu porównań dla \emph{każdego} takiego $W$, nie tylko
dla zbioru wypukłego.

Teraz pokryjmy dowolną rozmaitość $M$ lokalnie skończoną
rodziną dziedzin map. Każde ich skończone przecięcie
jest otwartym podzbiorem jednej z tych dziedzin,
a więc jest dyfeomorficzne z pewnym $W$ rozpatrzonym
w poprzednim akapicie. Nie musi być ściągalne ani spójne.
Porównania $I$ i $r$ są na nim już udowodnione.
Ponowne zastosowanie kroków 2 i 3 daje izomorfizmy
\[
 H^*_{\mathrm{sing}}(M;\R)
 \xrightarrow{\ H(r)\ }H^*(C_\infty^*(M;\R))
 \xleftarrow{\ H(I)\ }H^*_{\mathrm{dR}}(M).
\]
Szukany izomorfizm to $H(r)^{-1}H(I)$.
Oba odwzorowania komutują z przenoszeniem sympleksów
i cofnięciem form przez gładkie odwzorowania, więc
izomorfizm jest naturalny. Cały argument obejmuje
rozmaitości z brzegiem dzięki modelom w $\mathbb H^m$.

Warto też wyjaśnić interpretację przez cykle. Nad ciałem
dualność kompleksów opisana w kroku 3 identyfikuje
kohomologię z dualną przestrzenią homologii.
Z izomorfizmu $H(r)$ wynika zatem, że włączenie
$C_*^\infty(M;\R)\to C_*(M;\R)$ indukuje izomorfizm
homologii: odwzorowanie liniowe, którego dualne jest
izomorfizmem, samo jest izomorfizmem (niezerowe jądro
lub niezerowy iloraz wykryłby funkcjonał liniowy).
Każda klasa cyklu ma więc gładkiego reprezentanta,
a dwaj tacy reprezentanci tej samej klasy różnią się
brzegiem gładkiego łańcucha. Dla formy zamkniętej ich
całki są równe przez \eqref{eq:integracja-lancuchowa}.
To uzasadnia niezależność okresów od reprezentantów
bez dokonywania niezgodnych wyborów na ścianach.
\end{proof}

\begin{przyklad}[Co widzą okresy, a czego nie widzą]
Jeżeli $\omega$ jest zamkniętą $k$-formą, jej
\emph{okres} na cyklu $z$ to $\int_z\omega$.
Jest niezależny od wyboru reprezentanta klasy
homologii: jeżeli $z$ zmienimy o $\partial c$,
różnica całek to $\int_c\dd\omega=0$.
Zmiana $\omega$ o formę dokładną także nie zmienia
okresu, bo $\int_z\dd\eta=\int_{\partial z}\eta=0$.
Na okręgu forma $\vartheta/(2\pi)$ z
Przykładu~\ref{prz:de-rham-s1} ma okres $1$ na
podstawowym cyklu. W $\mathbb{RP}^2$ grupa
$H_1(-;\Z)\cong\Z/2$ ma torsję, lecz
$H^1_{\mathrm{dR}}(\mathbb{RP}^2)=0$:
homomorfizm z $\Z/2$ do $\R$ musi wysłać
generator na liczbę $a$ z $2a=0$, zatem $a=0$.
Twierdzenie de Rhama wykrywa kohomologię
rzeczywistą; nie zastępuje informacji o torsji.
\end{przyklad}

\section{Porównanie homologii, kohomologii de Rhama i homotopii}
\label{sec:porownanie-trzech-teorii}

Pierwsza różnica dotyczy \emph{przestrzeni, na których
można zbudować daną teorię}. Homologię symplicjalną
definiujemy dla kompleksu symplicjalnego $K$:
sklejamy wierzchołki, odcinki, trójkąty i sympleksy
wyższych wymiarów wzdłuż ich ścian. Taki kompleks
\emph{nie musi być rozmaitością}. Na przykład dwa
wypełnione trójkąty stykające się tylko w jednym
wierzchołku tworzą kompleks, ale otoczenie
wspólnego wierzchołka nie przypomina dysku
ani półdysku: po usunięciu tego wierzchołka ma dwie
składowe, podczas gdy dysk oraz półdysk po usunięciu
punktu środkowego, odpowiednio punktu na prostym brzegu,
pozostają spójne. Mimo to łańcuchy, operator brzegu
$\partial$ i grupy $H_k^{\mathrm{simp}}(K;\Z)$
są określone.

\begin{figure}[!htbp]
\centering
\begin{tikzpicture}[line cap=round,line join=round]
 \fill[manifoldblue!14] (-3.1,-.95)--(-3.1,.95)--(0,0)--cycle;
 \fill[manifoldgold!19] (0,0)--(3.1,.95)--(3.1,-.95)--cycle;
 \draw[manifoldblue!75!black,thick]
  (-3.1,-.95)--(-3.1,.95)--(0,0)--cycle;
 \draw[manifoldgold!80!black,thick]
  (0,0)--(3.1,.95)--(3.1,-.95)--cycle;
 \fill[manifoldred] (0,0) circle (2.5pt);
 \node[below=5pt,font=\small] at (0,0) {$p$};
 \node[font=\small,manifoldblue!80!black] at (-2.05,0) {$\Delta^2$};
 \node[font=\small,manifoldgold!80!black] at (2.05,0) {$\Delta^2$};
\end{tikzpicture}
\caption{Dwa trójkąty ze wspólnym wierzchołkiem $p$.
To kompleks symplicjalny, ale w $p$ nie jest
powierzchnią gładką, nawet z brzegiem.}
\label{fig:kompleks-nie-rozmaitosc}
\end{figure}

Homologia singularna idzie jeszcze dalej: jej
sympleksy są \emph{ciągłymi odwzorowaniami}
$\Delta^k\to X$, dlatego definiuje się ją
dla dowolnej przestrzeni topologicznej $X$,
również bez wybranej triangulacji. Kohomologia
de Rhama wymaga natomiast \emph{gładkiej
rozmaitości} $M$: żeby pisać formy i ich
różniczkę $\dd$, potrzebujemy przestrzeni
stycznych i gładkich odwzorowań. Dla kompleksu
z Rysunku~\ref{fig:kompleks-nie-rozmaitosc}
zwykła kohomologia de Rhama w sensie
Definicji~\ref{def:de-rham-cohomology}
nie jest zdefiniowana. Grupy homotopii
można z kolei definiować dla przestrzeni
topologicznych z wybranym punktem bazowym.

Do czego służą te niezmienniki, jeśli słowo
„dziura” nie daje jeszcze ścisłej odpowiedzi?
Pozwalają rozstrzygać, czy cykl jest brzegiem,
czy zamknięta forma jest różniczką innej formy,
czy pętlę albo odwzorowanie sfery można
ściągnąć do punktu. Pomagają też odróżnić
przestrzenie oraz badać odwzorowania między nimi.
Istotny jest zarówno obiekt, od którego
zaczynamy, jak i relacja uznająca dwa takie
obiekty za równoważne:

\begin{center}
\small
\renewcommand{\arraystretch}{1.18}
\begin{tabular}{>{\raggedright\arraybackslash}p{3.05cm}
                >{\raggedright\arraybackslash}p{4.4cm}
                >{\raggedright\arraybackslash}p{6.75cm}}
\hline
\textbf{Typ}&\textbf{Na czym konstruowana}&
\textbf{Co wykrywa i o czym mówi}\\
\hline
$H_k^{\mathrm{simp}}(K;\Z)$,
homologia symplicjalna&
Zorientowane sympleksy kompleksu $K$,
całkowite łańcuchy, operator brzegu
$\partial$; cykle modulo brzegi.&
Rozpoznaje składowe ($H_0$), klasy cykli,
które nie są brzegami, ich liczbę niezależnych
generatorów i torsję. Na zamkniętej spójnej
rozmaitości jej najwyższa homologia z $\Z$
pomaga rozpoznać orientowalność.\\
\hline
$H_k^{\mathrm{sing}}(X;\Z)$,
homologia singularna&
Wszystkie ciągłe odwzorowania
$\Delta^k\to X$ jako generatory łańcuchów;
ten sam operator $\partial$ i iloraz
cykli przez brzegi.&
Daje homologiczne informacje bez wybierania
triangulacji, dla dowolnej przestrzeni
topologicznej. Ciągłe $f:X\to Y$ wyznacza
$f_*:H_k(X)\to H_k(Y)$; można nim badać,
jak odwzorowanie przenosi klasy cykli.\\
\hline
$H^k_{\mathrm{dR}}(M)$,
kohomologia de Rhama&
Gładkie formy różniczkowe na rozmaitości $M$,
różniczka $\dd$; formy zamknięte modulo dokładne.&
Mówi, czy zamknięta forma ma \emph{globalnie}
formę pierwotną. Całki po cyklach mierzą jej
okresy; iloczyn zewnętrzny form pozwala badać
iloczyny klas. Zwykła wersja ma wartości
rzeczywiste, więc nie wykrywa torsji.\\
\hline
$\pi_1(X,p)$,
grupa podstawowa&
Pętle o początku i końcu w $p$;
utożsamiamy pętle homotopijne, a działaniem
jest ich składanie.&
Mówi, których pętli nie można ściągnąć do
punktu. Zawiera informację o nakryciach
i może być nieprzemienna; przez to odróżnia
przestrzenie o tym samym $H_1$.\\
\hline
$\pi_k(X,p)$, $k\geq2$,
wyższe grupy homotopii&
Odwzorowania sfer $S^k\to X$ z punktem
bazowym, modulo homotopia zachowująca
ten punkt.&
Mówią, czy odwzorowanie sfery można
ściągnąć do punktu (równoważnie: przedłużyć
na dysk $D^{k+1}$). Mogą być niezerowe,
nawet gdy $H_k(X;\Z)=0$.\\
\hline
\end{tabular}
\end{center}

Na przykład $H_1$ pozwala wykazać, że
nie istnieje retrakcja dysku $D^2$ na jego
brzeg $S^1$: gdyby $r:D^2\to S^1$
i inkluzja $i:S^1\to D^2$ spełniały
$r\circ i=\id$, to na homologii byłoby
$r_*\circ i_*=\id$ na $H_1(S^1;\Z)\cong\Z$.
Jest to niemożliwe, bo $H_1(D^2;\Z)=0$.
Kohomologia de Rhama rozstrzyga, czy da się
zbudować globalny potencjał dla zamkniętej
formy, a $\pi_1$ pomaga opisać nakrycia.
„Dziury” są więc tylko jednym z obrazowych
skrótów dla tych konkretnych pytań.

Różni się też koszt rachunku. Bez triangulacji
kompleks łańcuchów singularnych ma ogromnie
wiele generatorów. Skończona triangulacja
sprowadza homologię symplicjalną do macierzy
operatorów brzegu, lecz znalezienie odpowiedniej
triangulacji i obliczenie jąder oraz obrazów
bywa trudne. Formy pozwalają czasem wykorzystać
współrzędne i symetrie rozmaitości, choć
ustalenie, czy forma ma \emph{globalną}
formę pierwotną, także wymaga pracy.
Nie ma jednej metody zawsze najłatwiejszej.

\subsection*{Od sympleksów do całek}

W tej części $H_k^{\mathrm{simp}}$ oznacza tę samą homologię,
którą wcześniej zapisywaliśmy $H_k^\Delta$, a
$H_k^{\mathrm{sing}}$ to wcześniejsze $H_k$ przestrzeni.
Indeksy służą tylko odróżnieniu konstrukcji w porównaniu.
Dla triangulacji $K$ Twierdzenie
\ref{thm:symplicjalna-singularna} utożsamia
$H_k^{\mathrm{simp}}(K;\Z)$ z
$H_k^{\mathrm{sing}}(M;\Z)$: wynik nie zależy
od wyboru triangulacji, chociaż liczba i kształt
sympleksów od niej zależą. Dla nieskończonej, lokalnie
skończonej triangulacji stosujemy ten sam wynik do skończonych
podkompleksów: skończony łańcuch symplicjalny ma skończony
nośnik, a obraz każdego łańcucha singularnego jest zwarty,
więc spotyka tylko skończenie wiele sympleksów.
To dotyczy także łańcuchów wypełniających; izomorfizmy
na skończonych podkompleksach dają więc zarówno
surjektywność, jak i iniektywność porównania globalnego.
Cyklem jest łańcuch
$z$ z $\partial z=0$. Jeżeli $z=\partial c$,
to reprezentuje zero w homologii.
W najprostszym przykładzie trzy krawędzie
obwodu trójkąta tworzą $1$-cykl. Jeśli
pozostawimy sam obwód, nie ma $2$-sympleksu,
którego ten cykl byłby brzegiem, i
$H_1\cong\Z$. Jeśli dodamy wypełniony trójkąt,
ten sam cykl staje się jego brzegiem:
\[
 \partial[v_0v_1v_2]
 =[v_1v_2]-[v_0v_2]+[v_0v_1].
\]
Wtedy $H_1=0$. To dosłowny sens
„wykrycia dziury” przez homologię.
Na okręgu
jedno pełne obejście daje generator $H_1(S^1;\Z)
\cong\Z$; na $S^2$ cała zorientowana powierzchnia
daje generator $H_2(S^2;\Z)\cong\Z$.

Klasa zamkniętej formy przypisuje liczby
klasom cykli: całkujemy formę po cyklu.
Wybieramy przy tym gładkiego reprezentanta całej klasy
cyklu, co uzasadniono w dowodzie Twierdzenia~\ref{thm:de-rham}.
Taką liniową regułę przypisującą wektorom
liczby nazywamy \emph{funkcjonałem liniowym}.
Właśnie w tym sensie kohomologia de Rhama
wiąże się z homologią. Jeśli
$\dd\omega=0$ i $\partial z=0$, definiujemy
\begin{equation}\label{eq:parowanie-de-rham-homologia}
 \langle[\omega],[z]\rangle:=\int_z\omega.
\end{equation}
Zmiana $\omega$ o $\dd\eta$ zmienia całkę o
$\int_{\partial z}\eta=0$; zmiana $z$ o
$\partial c$ zmienia ją o
$\int_c\dd\omega=0$. Jest więc dobrze określona
na klasach po obu stronach. To dokładnie
rachunek Stokesa~\eqref{eq:integracja-lancuchowa}.

Czy \emph{każdą} funkcję liniową przypisującą
klasom homologii liczby rzeczywiste można
otrzymać przez całkowanie zamkniętej formy?
Twierdzenie de Rhama daje odpowiedź twierdzącą.
Aby zobaczyć, skąd bierze się ta odpowiedź,
pokażmy najpierw krok algebraiczny.
Dla kompleksu
łańcuchowego przestrzeni wektorowych nad $\R$
kołańcuch $c:C_k\to\R$ jest kocyklem
($\delta c=0$) dokładnie wtedy, gdy znika
na $B_k=\partial C_{k+1}$. Jego ograniczenie
do $Z_k=\ker(\partial:C_k\to C_{k-1})$ zależy
więc tylko od klasy w $H_k=Z_k/B_k$.
Jeżeli ten funkcjonał na $H_k$ jest zerowy,
to $c$ znika na całym $Z_k$. Definiujemy
na $B_{k-1}=\partial C_k$ funkcjonał
$b(\partial a):=c(a)$. Jest dobrze określony:
z $\partial a=\partial a'$ wynika
$a-a'\in Z_k$, a zatem $c(a)=c(a')$.
Rozszerzamy $b$ liniowo na $C_{k-1}$.
Wtedy $(\delta b)(a)=b(\partial a)=c(a)$,
czyli $c$ był kobrzegiem. Odwrotnie, każdy
funkcjonał $\lambda:H_k\to\R$ definiuje
funkcjonał na $Z_k$, który można liniowo
przedłużyć na $C_k$; ponieważ znika na
$B_k$, przedłużenie jest kocyklem.
Dowiedliśmy zatem
\[
 H^k_{\mathrm{sing}}(M;\R)
 \cong \operatorname{Hom}_{\R}
       (H_k^{\mathrm{sing}}(M;\R),\R).
\]
Symbol $\operatorname{Hom}_{\R}(V,\R)$ oznacza
przestrzeń \emph{wszystkich} funkcjonałów liniowych
$V\to\R$: każdy przypisuje wektorowi z $V$
liczbę rzeczywistą i zachowuje dodawanie
oraz mnożenie przez skalar.
Łącząc to z Twierdzeniem~\ref{thm:de-rham}
i porównaniem symplicjalnym, otrzymujemy
dla triangulowanej rozmaitości
\begin{equation}\label{eq:porownanie-simp-de-rham}
 H^k_{\mathrm{dR}}(M)
 \cong \operatorname{Hom}_{\R}
       (H_k^{\mathrm{simp}}(K;\R),\R),
\end{equation}
przy czym izomorfizm jest dany przez całkę
z~\eqref{eq:parowanie-de-rham-homologia}.
Jeśli grupy homologii są skończenie generowane,
wymiar $H^k_{\mathrm{dR}}$ to $k$-ta liczba
Betti'ego, czyli liczba swobodnych generatorów
$H_k(M;\Z)$. Forma nie zobaczy torsji: gdy
$nz=\partial c$ dla $n\ne0$, to dla zamkniętej
$\omega$ mamy
$n\int_z\omega=\int_c\dd\omega=0$,
więc $\int_z\omega=0$.

Wybór współczynników wyjaśnia tu ważną różnicę.
Według Twierdzenia~\ref{thm:homologia-char-zero}
przejście z $\Z$ do $\Q$, $\R$ lub $\C$
zachowuje liczby Betti'ego, lecz usuwa torsję;
Przykład~\ref{prz:zmiana-wspolczynnikow-rp2-torus}
pokazuje to dla $\mathbb{RP}^2$ i torusa.
Zwykły kompleks form de Rhama jest rzeczywisty
i oblicza $H^k_{\mathrm{sing}}(M;\R)$.
Jeśli dopuścimy formy zespolone,
$\Omega^\bullet(M;\C)
=\Omega^\bullet(M;\R)\otimes_{\R}\C$,
otrzymamy analogicznie
$H^k_{\mathrm{dR}}(M;\C)
\cong H^k_{\mathrm{dR}}(M;\R)\otimes_{\R}\C$.
Współczynniki $\Q$ stosujemy bezpośrednio
w łańcuchach i kołańcuchach singularnych;
nie wymagają one osobnej teorii gładkich
form o wymiernych współczynnikach.

\subsection*{Od sfer do cykli}

W tym porównaniu zakładamy, że $M$ jest spójna drogowo;
w przeciwnym razie pracujemy na składowej zawierającej $p$.
Element $\pi_k(M,p)$ jest klasą odwzorowania
$(S^k,*)\to(M,p)$. Po przeniesieniu generatora
$H_k(S^k;\Z)$ otrzymujemy klasę w $H_k(M;\Z)$:
to odwzorowanie Hurewicza $h_k$
z Rozdziału~\ref{ch:wyzsza-homotopia}.
W stopniu $1$ dokładna odpowiedź z
Twierdzenia~\ref{thm:hurewicz-h1} brzmi
\[
 H_1(M;\Z)\cong\pi_1(M,p)_{\mathrm{ab}}.
\]
Komutatory pętli znikają po przejściu
do homologii, ponieważ dodawanie cykli
jest przemienne; $\pi_1$ może być nieprzemienna.
W wyższych stopniach $h_k$ nie musi być
izomorfizmem. Dla $k\geq2$, gdy $M$ jest
$(k-1)$-spójna, Twierdzenie
\ref{thm:hurewicz-wyzszy} utożsamia
pierwszą możliwą niezerową grupę
$\pi_k(M)$ z $H_k(M;\Z)$.
Aby zastosować tu wersję twierdzenia dla kompleksów CW,
używamy dodatkowego twierdzenia o triangulacji gładkich
rozmaitości: każda taka rozmaitość dopuszcza lokalnie
skończoną triangulację, a przy brzegu można wybrać
$\partial M$ jako podkompleks. Sympleksy tej triangulacji
są komórkami CW. Istnienie triangulacji jest osobnym
wynikiem topologii różniczkowej, nie konsekwencją de Rhama;
jego dowód wykracza poza to porównanie. Źródło wskazano
w dalszej lekturze. W rozważanych przykładach okręgu,
sfer, torusa i przestrzeni rzutowej strukturę CW znamy
już jawnie z wcześniejszych rozdziałów.
Teza dotyczy właśnie pierwszego
niezerowego stopnia przy podanej spójności;
nie utożsamia wszystkich późniejszych grup.

\begin{przyklad}[Ta sama informacja w stopniu pierwszym]
Na $S^1$ zachodzi
$\pi_1(S^1)\cong H_1(S^1;\Z)\cong\Z$,
a $H^1_{\mathrm{dR}}(S^1)\cong\R$.
Klasa formy $\vartheta/(2\pi)$ z
Przykładu~\ref{prz:de-rham-s1} daje wartość
$1$ na dodatnim generatorze homologii.
Liczba całkowita zlicza okrążenia, a całka
z formy przypisuje im liczby rzeczywiste
w sposób liniowy.
\end{przyklad}

\begin{przyklad}[Ta sama pierwsza homologia,
różne grupy podstawowe]
Torus $T^2$ ma $\pi_1(T^2)\cong\Z^2$.
Jeśli usuniemy otwarty dysk o domknięciu zawartym
we wnętrzu jego dwuwymiarowej komórki,
otrzymamy zwartą powierzchnię z brzegiem
$P=T^2\setminus\operatorname{int}D^2$.
W modelu komórkowym z Przykładu
\ref{prz:pi1-torus-rp2} część
dwuwymiarowa po usunięciu dysku ściąga się
na część jednowymiarową $S^1\vee S^1$.
Dlatego $\pi_1(P)\cong F_2$, wolna
nieprzemienna grupa na dwóch generatorach,
podczas gdy $H_1(P;\Z)\cong\Z^2$.
Zarówno $P$, jak i $T^2$ mają więc
$H^1_{\mathrm{dR}}\cong\R^2$:
w stopniu pierwszym formy i homologia
nie wykrywają nieprzemienności $\pi_1(P)$.
Pełne grupy homologii tych przestrzeni
różnią się jednak w stopniu $2$:
$H_2(P;\Z)=0$, a $H_2(T^2;\Z)=\Z$.
\end{przyklad}

\begin{przyklad}[Informacja sferyczna bez klasy homologii]
Dla $\mathbb{RP}^2$ wcześniejsze obliczenia dają
\[
 \pi_1(\mathbb{RP}^2)=\Z/2,\qquad
 \pi_2(\mathbb{RP}^2)=\Z,\qquad
 H_1(\mathbb{RP}^2;\Z)=\Z/2,\qquad
 H_2(\mathbb{RP}^2;\Z)=0.
\]
Porównaj Przykłady~\ref{prz:rp2-pi2}
i~\ref{prz:homologia-rp2}. Nakrycie
$S^2\to\mathbb{RP}^2$ pozostawia nietrywialną
klasę w $\pi_2$, ale jej obraz przez $h_2$
jest zerowy. Dodatkowo
$H^k_{\mathrm{dR}}(\mathbb{RP}^2)=0$
dla każdego $k>0$: całki form nie wykryją
ani torsji w $H_1$, ani tej klasy sferycznej.
Jest to wyraźny przykład, dlaczego trzech teorii
nie wolno po prostu utożsamiać.
\end{przyklad}

\section*{Dalsza lektura}
\addcontentsline{toc}{section}{Dalsza lektura}
\href{https://ocw.mit.edu/courses/18-965-geometry-of-manifolds-fall-2004/pages/lecture-notes/}
{Notatki kursu \emph{Geometry of Manifolds} (MIT OpenCourseWare)},
wykłady o formach, twierdzeniu Stokesa i twierdzeniu de Rhama.
Zob. także \href{https://www.math.toronto.edu/mgualt/courses/17-1300/docs/17-1300-notes-13.pdf}
{notatki o formach różniczkowych i twierdzeniu Stokesa}
Uniwersytetu w Toronto.
Porównanie z homologią i homotopią rozwija
\href{https://pi.math.cornell.edu/~hatcher/AT/AT.pdf}
{Allen Hatcher, \emph{Algebraic Topology}},
rozdziały 2--4. Twierdzenie o triangulacji rozmaitości
gładkich omawia Jacob Lurie w notatkach
\href{https://www.math.ias.edu/~lurie/937notes/937Lecture3.pdf}
{\emph{Whitehead Triangulations}, wykład 3} i kolejnych
wykładach tego kursu.


\chapter{Homologia i kohomologia \v{C}echa: od pokryć do granic przestrzeni}
\label{ch:cech-v4}

\section{Motywacja: co dzieje się przy nieskończonym przybliżaniu?}
Homologia singularna buduje łańcuchy z pojedynczych ciągłych map
standardowych sympleksów do $X$. Dobrze współpracuje z homotopią,
ale w przestrzeni o coraz drobniejszej, nieskończonej strukturze
lokalnej nie zawsze zachowuje się jak granica homologii
skończonych przybliżeń. Można zamiast tego obserwować $X$
przez \emph{otwarte pokrycia}: każdy układ przecinających się
zbiorów opisujemy skończonym kompleksem symplicjalnym,
a następnie pokrycia sukcesywnie zagęszczamy.

Ta zmiana punktu widzenia rozwiązuje szczególnie ważny problem:
jak obliczyć niezmienniki zwartego zbioru $A$ za pomocą malejących
otoczeń $U\supset A$ i jak związać je z homologią dopełnienia
$S^n\setminus A$ (dualność Aleksandra). Trzeba jednak od razu
rozróżnić \emph{homologię} i \emph{kohomologię} \v{C}echa:
pierwsza bierze granicę odwrotną, druga granicę prostą.

\section{Nerw pokrycia i odwzorowania przy rozdrobnieniu}
\begin{definicja}[Nerw]
Niech $\mathcal U=(U_i)_{i\in J}$ będzie pokryciem otwartym
przestrzeni $X$. \emph{Nerwem} $N(\mathcal U)$ jest abstrakcyjny
kompleks symplicjalny o wierzchołkach $i\in J$; skończona rodzina
różnych wierzchołków $\{i_0,\ldots,i_k\}$ wyznacza sympleks
wtedy i tylko wtedy, gdy $U_{i_0}\cap\cdots\cap U_{i_k}\ne\varnothing$.
Do skończonego pokrycia przyporządkowujemy zwykły skończony
kompleks symplicjalny, do którego stosujemy homologię
symplicjalną oraz kohomologię.
\end{definicja}

\begin{przyklad}[Okrąg]
Pokryjmy $S^1$ trzema dostatecznie krótkimi otwartymi łukami
$U_0,U_1,U_2$ tak, aby każde dwa się przecinały, lecz
$U_0\cap U_1\cap U_2=\varnothing$.
Nerw ma trzy wierzchołki i trzy krawędzie, ale nie ma
wypełnionego trójkąta; jego realizacja geometryczna jest więc
homeomorficzna z $S^1$.
Jeżeli dopuścimy niepuste przecięcie potrójne, pojawi się
$2$-sympleks, a nerw stanie się ściągalny. Samo pojedyncze
pokrycie nie może zatem definiować niezmiennika topologicznego:
trzeba rozważać \emph{wszystkie} coraz drobniejsze pokrycia.
\end{przyklad}

\begin{definicja}[Rozdrobnienie]
Pokrycie $\mathcal V$ jest \emph{drobniejsze} od $\mathcal U$,
co zapisujemy $\mathcal V\succ\mathcal U$, jeśli dla każdego
$V\in\mathcal V$ istnieje $U\in\mathcal U$ z $V\subseteq U$.
Wybór takich przyporządkowań daje mapę symplicjalną
$r_{\mathcal V\mathcal U}:N(\mathcal V)\to N(\mathcal U)$:
każdy wierzchołek $V$ wysyłamy na wybrany $U$ zawierający $V$.
\end{definicja}

\begin{lemat}[Niezależność wyborów]
\label{lem:cech-contiguity-v4}
Dwie mapy nerwów wybrane dla jednego rozdrobnienia
$\mathcal V\succ\mathcal U$ są przyległe symplicjalnie,
a zatem indukują te same homomorfizmy homologii i kohomologii.
Przy trzykrotnym rozdrobnieniu $\mathcal W\succ\mathcal V\succ
\mathcal U$ mapy indukowane są zgodne ze składaniem.
\end{lemat}
\begin{proof}
Niech $\{V_0,\ldots,V_k\}$ będzie sympleksem nerwu
$N(\mathcal V)$ i $x\in\bigcap_j V_j$.
Jeśli $U_j$ i $U'_j$ to dwa możliwe wybory zbiorów z
$\mathcal U$ zawierających $V_j$, to wszystkie $U_j,U'_j$
zawierają $x$. Ich wierzchołki tworzą więc wspólnie sympleks
w $N(\mathcal U)$, co jest dokładnie warunkiem przyległości.
Homotopia liniowa między realizacjami map przyległych
pozostaje w tym samym sympleksie, jest więc ciągła.
Mapy homotopijne działają identycznie na homologię i kohomologię.
Przy składaniu oba wybory z $\mathcal W$ do $\mathcal U$
również są przyległe, stąd ostatnie twierdzenie.
\end{proof}

\section{Definicje i kierunki granic}
W dalszym ciągu przyjmujemy, że $X$ jest \emph{zwartą przestrzenią
Hausdorffa}, a $G$ grupą abelową. Wszystkie otwarte pokrycia
mają skończone podpokrycia, a każde dwa skończone pokrycia
mają wspólne drobniejsze pokrycie (ich niepuste przecięcia).
Zatem skończone pokrycia tworzą zbiór skierowany przez
relację rozdrobnienia.

\begin{definicja}[Homologia i kohomologia \v{C}echa]
Dla każdego $n\ge0$ definiujemy
\[
\check H_n(X;G)=\varprojlim_{\mathcal U} H_n(N(\mathcal U);G),
\qquad
\check H^n(X;G)=\varinjlim_{\mathcal U}H^n(N(\mathcal U);G).
\]
W pierwszym układzie mapy $H_n(N(\mathcal V);G)\to
H_n(N(\mathcal U);G)$ biegną od drobniejszego do grubszego
pokrycia. W drugim układzie przez kontrawariantność
kohomologii otrzymujemy mapy
$H^n(N(\mathcal U);G)\to H^n(N(\mathcal V);G)$.
\emph{Zredukowane} grupy $\widetilde{\check H}_n$ i
$\widetilde{\check H}^{\,n}$ definiuje się za pomocą
zredukowanych (ko)homologii nerwów.
\end{definicja}

\begin{uwaga}[Co dosłownie znaczy granica odwrotna?]
Elementem $\check H_n(X;G)$ jest rodzina
$(a_{\mathcal U})_{\mathcal U}$, gdzie
$a_{\mathcal U}\in H_n(N(\mathcal U);G)$, spełniająca
$(r_{\mathcal V\mathcal U})_*a_{\mathcal V}=a_{\mathcal U}$
przy $\mathcal V\succ\mathcal U$.
Nie wystarcza wybrać dowolny cykl dla każdego pokrycia:
klasy muszą być zgodne przy rozdrobnieniu.
Natomiast element kohomologii \v{C}echa reprezentuje pewna
klasa w nerwie \emph{jednego} pokrycia; klasy reprezentują
ten sam element, jeśli po przejściu do wspólnego drobniejszego
pokrycia stają się równe. Jest to istota granicy prostej.
\end{uwaga}

\begin{stwierdzenie}[Funktorialność]
Jeśli $f:X\to Y$ jest ciągła między zwartymi przestrzeniami
Hausdorffa, to indukuje
$f_*:\check H_n(X;G)\to\check H_n(Y;G)$
oraz $f^*:\check H^n(Y;G)\to\check H^n(X;G)$.
Odwzorowania te są zgodne ze składaniem i identycznością.
\end{stwierdzenie}
\begin{proof}
Dla skończonego pokrycia $\mathcal U$ przestrzeni $Y$
przeciwobrazy $f^{-1}(U)$ stanowią skończone pokrycie $X$.
Zachodzi naturalna mapa nerwów
$N(f^{-1}\mathcal U)\to N(\mathcal U)$ (puste przeciwobrazy
pomijamy). Zgodność z rozdrobnieniami wynika z
Lematu~\ref{lem:cech-contiguity-v4}. Przechodząc do granicy
odwrotnej, otrzymujemy $f_*$; po odwróceniu strzałek
kohomologii dostajemy $f^*$.
\end{proof}

\section{Porównanie z homologią singularną}
\begin{stwierdzenie}[Porównanie dla wielościanów]
\label{prop:cech-poly-v4}
Jeżeli $X$ jest zwartym wielościanem (w szczególności
skończonym CW-kompleksem, który jest triangulowalny), to
naturalne homomorfizmy porównawcze są izomorfizmami
$H_n(X;G)\cong\check H_n(X;G)$ i
$H^n(X;G)\cong\check H^n(X;G)$ dla $n\ge0$.
\end{stwierdzenie}
\begin{proof}
Ustalmy skończoną triangulację $K$. Pokrycia otwartymi
gwiazdami w kolejnych podziałach barycentrycznych są
kofinalne wśród skończonych pokryć otwartych: z lematu
Lebesgue'a o liczbie pokrycia i malejących średnic sympleksów
wynika, że po dostatecznej liczbie podziałów każda taka
gwiazda mieści się w jednym zbiorze zadanego pokrycia.
Nerw pokrycia otwartymi gwiazdami wierzchołków odpowiedniego
podziału jest izomorficzny z jego kompleksem symplicjalnym.
Mapy nerwów dla kolejnych rozdrobnienień odpowiadają
mapom aproksymacji symplicjalnej identyczności $|K|\to|K|$,
a zatem indukują izomorfizmy na (ko)homologii.
Układy odwrotny i prosty stają się więc układami
stałymi z dokładnością do kanonicznych izomorfizmów.
\end{proof}

\begin{uwaga}
Nie można bez dodatkowych hipotez ogłosić powyższych
izomorfizmów dla wszystkich przestrzeni zwartych.
W szczególności spójność i spójność łukowa nie są tym samym:
$H_0^{\mathrm{sing}}(X;\Z)$ wykrywa składowe łukowe,
natomiast $\check H_0$ w zwartych przykładach z poniższej
sekcji może rejestrować jedynie spójność widoczną w
kolejnych skończonych przybliżeniach.
\end{uwaga}

\section{Kolczyk hawajski: coraz więcej okręgów}
\begin{definicja}
Kolczykiem hawajskim nazywamy zwarty podzbiór płaszczyzny
$\mathbb H=\bigcup_{n\ge1}C_n$, gdzie
$C_n=\{(x,y)\in\R^2:(x-1/n)^2+y^2=1/n^2\}$.
Wszystkie okręgi są styczne w $(0,0)$, a ich średnice
maleją do zera.
\end{definicja}

\begin{stwierdzenie}
Dla każdego $k\ge1$ zachodzi
$\check H_1(\mathbb H;\Z)\cong\prod_{n=1}^{\infty}\Z$,
zaś $\check H_k(\mathbb H;\Z)=0$ dla $k\ge2$.
\end{stwierdzenie}
\begin{proof}
Niech $X_m=C_1\vee\cdots\vee C_m$ będzie bukietem
pierwszych $m$ okręgów. Zbiory coraz mniejszych ogonów
$\bigcup_{n>m}C_n$ kurczą się w dowolnie małym otoczeniu
punktu wspólnego. Pozwala to dobrać kofinalny układ
skończonych pokryć $\mathbb H$, których nerwy mają typ
homotopijny $X_m$, a mapy od etapu $m+1$ do $m$ odpowiadają
zgnieceniu ostatniego okręgu. W szczególności układ
pierwszych grup homologii to
\[
\Z\longleftarrow\Z^2\longleftarrow\Z^3\longleftarrow\cdots,
\qquad (a_1,\ldots,a_{m+1})\longmapsto(a_1,\ldots,a_m).
\]
Rodzina zgodna jest wyznaczona przez nieskończony ciąg
liczb całkowitych $(a_1,a_2,\ldots)$, zatem granica odwrotna
to iloczyn $\prod_{n\ge1}\Z$. Dla $k\ge2$ wszystkie grupy
w tym układzie są zerowe. Uzasadnienie kofinalności można
przeprowadzić, dobierając dla danego pokrycia otoczenie
punktu $(0,0)$ mieszczące wszystkie dostatecznie małe okręgi
i dzieląc skończoną liczbę pozostałych łuków.
\end{proof}

\begin{uwaga}[Nie mylić z grupą fundamentalną]
Nie stwierdzamy, że $\pi_1(\mathbb H)$ jest wolną grupą o
przeliczalnej liczbie generatorów. Występują pętle
przebiegające nieskończenie wiele okręgów w coraz krótszym
czasie. To właśnie typ zjawiska, przy którym model przez
układ skończonych nerwów daje inną perspektywę niż
pojedynczy skończony kompleks CW.
\end{uwaga}

\section{Solenoid: granica okręgów i zanik pierwszej homologii}
Niech $p\ge2$ będzie liczbą całkowitą. Definiujemy
\emph{solenoid $p$-adyczny}
$\Sigma_p=\varprojlim(S^1\xleftarrow{z\mapsto z^p}S^1
\xleftarrow{z\mapsto z^p}\cdots)$.
Punktami solenoidu są ciągi $(z_0,z_1,\ldots)$ spełniające
$z_j=z_{j+1}^p$. Jest to zwarta przestrzeń Hausdorffa,
a każdy rzut na skończony poziom jest ciągły.

\begin{twierdzenie}[Dwie granice o przeciwnych wynikach]
\label{thm:solenoid-v4}
Dla $\Sigma_p$ zachodzi
\[
\check H_0(\Sigma_p;\Z)\cong\Z,\qquad
\check H_1(\Sigma_p;\Z)=0,\qquad
\check H^{1}(\Sigma_p;\Z)\cong\Z[1/p].
\]
Solenoid jest spójny, ale nie jest łukowo spójny.
\end{twierdzenie}
\begin{proof}
W przypadku granic odwrotnych zwartych wielościanów
$X=\varprojlim X_j$ z surjektywnymi mapami przejścia
pokrycia wstecznymi obrazami pokryć poziomów tworzą
rodzinę kofinalną. Wynika to z faktu, że ich obrazy zwrotne
tworzą bazę topologii granicy, a zwartość pozwala przejść
od pokryć do skończonych podpokryć. Stąd
$\check H_n(X;G)\cong\varprojlim H_n(X_j;G)$ oraz
$\check H^n(X;G)\cong\varinjlim H^n(X_j;G)$
dla tego układu wielościanów.

Na każdym poziomie mamy $H_0(S^1;\Z)=\Z$ i
$H_1(S^1;\Z)=\Z$. Odwzorowanie $z\mapsto z^p$ działa
na $H_0$ identycznie, a na $H_1$ mnożeniem przez $p$.
Zatem element $\check H_1$ to rodzina
$(a_0,a_1,\ldots)$ z $a_j=pa_{j+1}$.
Wymaga to, aby $a_0$ było podzielne przez dowolnie
wysoką potęgę $p$, zatem wszystkie $a_j=0$.
W kohomologii strzałki biegną odwrotnie, więc
$\check H^1$ jest granicą prostą
$\Z\xrightarrow{\cdot p}\Z\xrightarrow{\cdot p}\cdots$.
Klasa liczby $m$ z poziomu $j$ odpowiada $m/p^j$,
co daje $\Z[1/p]$.

Spójność wynika z tego, że granica odwrotna niepustych
zwartych przestrzeni spójnych przy surjektywnych mapach
przejścia jest spójna: gdyby granica rozpadła się na dwa
niepuste domknięto-otwarte zbiory, zwartość pozwoliłaby
wykryć to rozdzielenie na dostatecznie wysokim poziomie,
wbrew spójności okręgu.

Dla braku łukowej spójności rozważmy włókno rzutu
$\Sigma_p\to S^1$ nad $1$. Jest ono homeomorficzne z
$\Z_p$, zbiorem ciągów reszt zgodnych modulo $p^j$.
Ścieżka od punktu $(1,1,\ldots)$ do punktu tego włókna
po złożeniu z pierwszym rzutem daje pętlę o całkowitej
liczbie owinięć $m$. Kolejne współrzędne ścieżki są
jednoznacznymi podniesieniami tej pętli dla pokryć stopnia
$p^j$, zatem jej końcowe współrzędne odpowiadają resztom
$m\bmod p^j$. Możliwych takich końców jest przeliczalnie
wiele (po jednym dla $m\in\Z$), podczas gdy $\Z_p$ jest
nieprzeliczalny. Stąd nie każde dwa punkty można połączyć
ścieżką.
\end{proof}

\begin{wniosek}
Dla solenoidu $H_0^{\mathrm{sing}}(\Sigma_p;\Z)$ nie jest
izomorficzna z $\check H_0(\Sigma_p;\Z)=\Z$:
pierwsza grupa jest wolną grupą na składowych łukowych,
których jest więcej niż jedna. Nie należy zatem bez
hipotez lokalnej regularności zamieniać teorii singularnej
na teorię \v{C}echa.
\end{wniosek}

\section{Kohomologia \v{C}echa przez otoczenia i rola dualności}
\begin{twierdzenie}[Opis przez otoczenia w sferze]
\label{thm:cech-neighbourhood-v4}
Niech $A\subset S^n$ będzie niepustym zbiorem zwartym,
$G$ grupą abelową. Wówczas
$\check H^k(A;G)\cong\varinjlim_{U\supset A}H^k(U;G)$,
gdzie granicę bierzemy po otwartych otoczeniach $U$ zbioru
$A$, przy przejściach danych przez restrykcje.
\end{twierdzenie}
\begin{proof}
Wybieramy triangulacje sfery o coraz mniejszych średnicach
komórek. Dla ustalonego otwartego $U\supset A$ zwartość $A$
pozwala dobrać skończone pokrycie otwartymi gwiazdami
w dostatecznie drobnym podziale, którego suma leży w $U$.
Po niewielkim pogrubieniu otrzymujemy regularne otoczenie
$N$ zbioru $A$ o typie homotopijnym nerwu tego pokrycia
(lemat o nerwie, ponieważ wszystkie niepuste przecięcia
otwartych gwiazd są ściągalne). Dla odwrotnego przejścia,
każde skończone pokrycie otwarte $A$ można przedłużyć do
pokrycia pewnego otoczenia $U$ poprzez wybór otwartych
zbiorów sfery mających zadane ślady na $A$; po zwężeniu
otoczenia nie pojawiają się nowe niepuste przecięcia
skończonych rodzin, które były puste na $A$.
Te dwa kroki dają kofinalne układy nerwów i otoczeń,
a więc identyczną granicę prostą kohomologii.
\end{proof}

\begin{uwaga}[Dlaczego nie pojawia się tu homologia singularna $A$?]
W dualności Aleksandra dla dowolnego zwartego $A\subset S^n$
po prawej stronie musi stać kohomologia \v{C}echa.
Dla skończonych wielościanów oba rodzaje kohomologii
się zgadzają (Stwierdzenie~\ref{prop:cech-poly-v4}),
lecz w przypadku zbiorów dzikich nie wolno wykonywać
takiego podstawienia. Granica \emph{prosta} w powyższym
wzorze jest również zgodna z kierunkiem map wywołanym
przez malejące otoczenia.
\end{uwaga}

\section{Granice odwrotne a dokładność}
Granica prosta grup abelowych jest funktorem dokładnym,
ale granica odwrotna w ogólności nie jest prawostronnie
dokładna. Wynika stąd, że naiwnie zdefiniowana homologia
\v{C}echa nie musi spełniać wszystkich aksjomatów dokładności
homologii Eilenberga--Steenroda dla dowolnych par zwartych
przestrzeni. Przy bardziej zaawansowanym badaniu granic
stosuje się pochodny funktor $\varprojlim{}^1$ i tzw.
silne homologie. W tej książce rozróżniamy te konstrukcje,
nie utożsamiając ich bez dowodu.

\begin{uwaga}
Dla zadanego układu odwrotnego grup abelowych z
surjektywnymi mapami przejścia (przy przeliczalnym
indeksowaniu) zachodzą przydatne warunki zaniku
$\varprojlim{}^1$. Nie należy jednak przypisywać tej
własności każdemu układowi pochodzącemu od pokryć;
przykład solenoidu ma na $H_1$ mapy mnożenia przez $p$,
które nie są surjekcjami.
\end{uwaga}
\clearpage

\chapter{Przecięcia, dualność i klasy charakterystyczne}
\label{ch:przeciecia-dualnosc-klasy}

Otoczenie tubularne z twierdzenia o otoczeniu tubularnym
pozwala odróżnić ruch \emph{wzdłuż} podrozmaitości
od ruchu \emph{poprzecznego}. Teraz przypiszemy
przecięciom liczby, a wiązkom -- klasy kohomologii.
W ten sposób geometryczne dane zanurzenia stają się
niezmiennikami, które można obliczać i porównywać.
To przygotowanie do teorii Morse'a, Smale'a i chirurgii;
w końcowej sekcji wykonamy też prostą operację
wycinania i sklejania rozmaitości.

Przez \emph{zamkniętą rozmaitość} rozumiemy w tym
rozdziale rozmaitość zwartą bez brzegu. W dowodzie
komórkowym będziemy używać triangulacji gładkiej
rozmaitości. Jej istnienie jest osobnym, głębszym
twierdzeniem o triangulacji; opisujemy dokładnie
argument dualności \emph{po} wybraniu takiej
triangulacji. Rozróżnienie tych dwóch kroków
zapobiega ukryciu twierdzenia o triangulacji
w rysunku.

\section{Kołańcuchy symplicjalne i iloczyn kubkowy}
\label{sec:kubek-16}

Kompleksy łańcuchowe i ich duale poznaliśmy
w Rozdziałach~\ref{ch:algebra-homologiczna}
i~\ref{ch:homologia-przestrzeni}. Dla
skończonej triangulacji $K$ i ciała $\mathbb F$
zapisujemy precyzyjnie
\[
 C^k(K;\mathbb F)
   =\operatorname{Hom}_{\mathbb F}
       (C_k(K;\mathbb F),\mathbb F),
 \qquad
 (\dd\varphi)(c)=\varphi(\partial c).
\]
\emph{Kołańcuch} $\varphi\in C^k$ jest więc
funkcją liniową \emph{na łańcuchach} stopnia $k$.
Nie jest kolejnym rodzajem sympleksu:
przypisuje każdemu łańcuchowi liczbę z
$\mathbb F$. To polska nazwa angielskiego
\emph{cochain}; łańcuch należy do $C_k$,
a kołańcuch do przestrzeni dualnej $C^k$.
Znak $\partial$ pozostaje zarezerwowany
dla łańcuchów, a $\dd$ dla kołańcuchów.
Równość $\dd^2=0$ wynika od razu z
$\partial^2=0$. Klasa w $H^k$ to
funkcja na cyklach, która nie zmienia
się po dodaniu brzegu. Nad ciałem
$H^k(K;\mathbb F)\cong
\operatorname{Hom}_{\mathbb F}(H_k(K;\mathbb F),\mathbb F)$:
można wybrać uzupełnienia baz jąder
i obrazów i zdualizować powstały
rozkład. Nad $\Z$ analogiczne zdanie
wymaga poprawki $\operatorname{Ext}$,
co było omawiane przy zmianie
współczynników w Rozdziale~\ref{ch:algebra-homologiczna}.
Pracę nad ciałem rozpoczniemy tu dlatego,
że łatwiej zobaczyć parowanie;
dualność Poincarégo ma również
wersję całkowitą.

\begin{definicja}[Iloczyn kubkowy]
\label{def:kubek-16}
Ustalmy porządek wierzchołków $K$ i przemienny
pierścień współczynników $R$ (w szczególności $\Z$ lub ciało).
Dla $\alpha\in C^p(K;R)$ i $\beta\in C^q(K;R)$ definiujemy
na sympleksach o wierzchołkach zapisanych rosnąco
\[
(\alpha\smile\beta)
 [v_0,\ldots,v_{p+q}]
=\alpha[v_0,\ldots,v_p]\,
 \beta[v_p,\ldots,v_{p+q}].
\]
Na przeciwnie zorientowanym sympleksie przedłużamy ten
wzór liniowo ze znakiem. Obie ściany dzielą wierzchołek $v_p$.
Po przejściu do kohomologii otrzymamy
iloczyn $H^p\times H^q\to H^{p+q}$.
\end{definicja}

\begin{stwierdzenie}[Reguła Leibniza]
\label{prop:kubek-leibniz-16}
Dla kołańcuchów zachodzi
\[
 \dd(\alpha\smile\beta)
 =(\dd\alpha)\smile\beta
       +(-1)^p\alpha\smile(\dd\beta).
\]
Zatem iloczyn klas kohomologii jest
dobrze określony; jest łączny.
\end{stwierdzenie}
\begin{proof}
Obliczamy prawą stronę na $[v_0,\ldots,v_{p+q+1}]$.
W $(\dd\alpha)\smile\beta$ wyrazy z usunięciem
$v_i$ dla $0\leq i\leq p$ odpowiadają dokładnie
tym samym wyrazom w $\dd(\alpha\smile\beta)$.
W $(-1)^p\alpha\smile\dd\beta$ odpowiadają im
pozostałe wyrazy, o indeksach $p+1\leq i\leq p+q+1$.
Na prawej stronie są ponadto dwa dodatkowe wyrazy:
\[
 \bigl((-1)^{p+1}+(-1)^p\bigr)
 \alpha[v_0,\ldots,v_p]\,
 \beta[v_{p+1},\ldots,v_{p+q+1}]=0.
\]
Pochodzą one odpowiednio z ostatniej ściany
$\dd\alpha$ i pierwszej ściany $\dd\beta$.
To one się znoszą; żaden wyraz lewego kobrzegu
nie zostaje usunięty. Otrzymujemy regułę Leibniza.
Jeśli $\dd\alpha=\dd\beta=0$,
iloczyn jest kocyklem.
Jeśli $\alpha$ zastąpimy przez
$\alpha+\dd\gamma$, to zmiana
iloczynu wynosi
$\dd(\gamma\smile\beta)$
przy zamkniętej $\beta$;
przy zmianie $\beta$ na $\beta+\dd\eta$ różnica
wynosi $(-1)^p\dd(\alpha\smile\eta)$ przy zamkniętej $\alpha$.
Łączność widać bezpośrednio:
obie iteracje na sympleksie
stopnia $p+q+r$ dają iloczyn
wartości na trzech kolejnych
ścianach
$[v_0,\ldots,v_p]$,
$[v_p,\ldots,v_{p+q}]$ i
$[v_{p+q},\ldots,v_{p+q+r}]$.
\end{proof}

\begin{stwierdzenie}[Przemienność z właściwym znakiem]
\label{prop:kubek-graded-16}
W kohomologii
$\alpha\smile\beta
=(-1)^{pq}\beta\smile\alpha$
dla $\alpha\in H^p$,
$\beta\in H^q$.
\end{stwierdzenie}
\begin{proof}
Pracujemy najpierw nad $\Z$ i oznaczamy $C=C_*(K;\Z)$.
Na kompleksie tensorowym przyjmujemy
$\partial(a\otimes b)=\partial a\otimes b+(-1)^{|a|}a\otimes\partial b$.
Odwzorowanie Alexandera--Whitneya ma postać
\[
 A[v_0,\ldots,v_m]
 =\sum_{j=0}^m[v_0,\ldots,v_j]\otimes[v_j,\ldots,v_m].
\]
Jest odwzorowaniem łańcuchowym: po rozpisaniu brzegów
wyrazy wewnętrzne przy miejscu podziału znoszą się
tak jak w poprzednim dowodzie. Również
$T(a\otimes b)=(-1)^{|a||b|}b\otimes a$
jest odwzorowaniem łańcuchowym.

Skonstruujemy homotopię $H$ między $A$ i $TA$.
Dla wierzchołków oba odwzorowania są równe, więc
kładziemy $H=0$ w stopniu zero. Jeżeli $H$ jest
już określone na ścianach sympleksu $\sigma^m$,
łańcuch
\[
 z=TA(\sigma)-A(\sigma)-H(\partial\sigma)
\]
jest cyklem: stosujemy $\partial$, korzystamy
z łańcuchowości obu map i założenia indukcyjnego.
Cały $z$ leży w $C_*(\sigma)\otimes C_*(\sigma)$.
Ten kompleks ma zerową homologię w dodatnich stopniach:
stożkowanie sympleksu do jego wierzchołka jest
ściągnięciem łańcuchowym do $\Z$ w stopniu zero,
a tensorowanie tych ściągnięć daje to samo dla iloczynu.
Możemy więc wybrać $H(\sigma)$ z $\partial H(\sigma)=z$.
Indukcja daje $\partial H+H\partial=TA-A$.

Dla kocykli $\alpha,\beta$ kołańcuch
$\alpha\otimes\beta$ znika na brzegach kompleksu
tensorowego. Po ewaluacji ostatniej tożsamości
otrzymujemy, że
$(-1)^{pq}\beta\smile\alpha-\alpha\smile\beta$
jest kobrzegiem kołańcucha $(\alpha\otimes\beta)H$.
Dowód działa po zmianie współczynników na $R$.
Ta sama indukcja w ściągalnych nośnikach pokazuje
niezależność od porządku wierzchołków i wyboru
przybliżenia przekątnej. W połączeniu z aproksymacją
symplicjalną daje naturalność iloczynu w kohomologii.
\end{proof}

W przeciwieństwie do samego wymiaru
grup $H^k$, iloczyn kubkowy pamięta,
\emph{jak} różne cykle są względem siebie
położone. Ta dodatkowa informacja
wyjaśni później liczby przecięcia.
Oznaczenie $\smile$ odróżnia
iloczyn kołańcuchów od iloczynu
zewnętrznego form $\wedge$ z
Rozdziału~\ref{ch:formy-de-rham}.

\begin{definicja}[Iloczyn kapowy i konwencja znaków]
\label{def:kapowy-16}
Dla $m$-sympleksu $c=[v_0,\ldots,v_m]$ i kołańcucha
$\alpha$ stopnia $k\leq m$ kładziemy
\[
 c\frown\alpha=\alpha[v_0,\ldots,v_k]\,[v_k,\ldots,v_m].
\]
Wzór przedłużamy liniowo. Możemy wziąć całkowity łańcuch
i kołańcuch o wartościach w dowolnej grupie abelowej $G$;
wtedy wynik jest łańcuchem z $G$. Bezpośrednie rozpisanie
ścian (z tym samym skróceniem dwóch wyrazów przy podziale)
daje
\begin{equation}\label{eq:kapowy-brzeg-16}
 \partial(c\frown\alpha)=(-1)^k
  \bigl((\partial c)\frown\alpha-c\frown\dd\alpha\bigr).
\end{equation}
Dlatego cykl i kocykl wyznaczają klasę homologii,
niezależną od reprezentantów. Dla współczynników w $R$
mamy ponadto, prosto ze wzorów na ściany,
\[
 \langle\beta,c\frown\alpha\rangle
 =\langle\alpha\smile\beta,c\rangle.
\]
Ta kolejność czynników będzie obowiązywała w całym rozdziale.
\end{definicja}

\section{Orientacja jako klasa homologii}
\label{sec:klasa-fundamentalna-16}

W Stwierdzeniu~\ref{prop:lokalna-homologia}
ustaliliśmy, że
$H_n(M,M\setminus\{x\};\Z)\cong\Z$.
Orientacja oznacza spójny wybór
generatorów tych lokalnych grup.
Teraz skleimy je w jedną klasę
globalną.

\begin{twierdzenie}[Klasa fundamentalna]
\label{thm:klasa-fundamentalna-16}
Niech $M^n$ będzie spójną, zamkniętą,
zorientowaną rozmaitością. Istnieje
jedyna klasa
$[M]\in H_n(M;\Z)$, która w
$H_n(M,M\setminus\{x\};\Z)$
przechodzi w wybrany lokalny
generator dla każdego $x$.
W szczególności $H_n(M;\Z)\cong\Z$.
Dla zwartej zorientowanej rozmaitości z brzegiem
otrzymujemy zamiast tego
$[M,\partial M]\in H_n(M,\partial M;\Z)$.
Bez orientacji zawsze istnieje
klasa fundamentalna ze
współczynnikami $\mathbb F_2=\Z/2$.
\end{twierdzenie}
\begin{proof}
Wybierzmy triangulację $K$ zgodną
ze strukturą gładką. Orientacja $M$
wyznacza orientację każdego
$n$-sympleksu. Niech
$c=\sum_{\sigma^n\in K}\sigma^n$
oznacza sumę tak zorientowanych
sympleksów (w zapisie łańcuchowym
znaki są już uwzględnione w ich
orientacjach). Każda wewnętrzna
$(n-1)$-ściana należy dokładnie
do dwóch $n$-sympleksów.
Orientacje indukowane na ich
wspólnej ścianie są przeciwne,
więc odpowiadające współczynniki
w $\partial c$ się znoszą.
W przypadku $\partial M=\varnothing$
otrzymujemy $\partial c=0$.
W przypadku brzegu pozostają tylko
ściany leżące w $\partial M$,
więc $\partial c\in C_{n-1}(\partial M)$:
$c$ jest cyklem względnym.

W punkcie $x$ we wnętrzu
$n$-sympleksu obraz $c$ w
grupie lokalnej jest generatorem
o znaku określonym orientacją.
Jeśli $x$ leży na ścianie,
mała gwiazda wokół $x$ jest
kulą złożoną z przyległych
sympleksów; znoszenie się
wewnętrznych ścian pokazuje,
że ta sama suma daje generator
również tam. W przypadku brzegu powyższy warunek lokalnego
generatora nakładamy na punkty \emph{wnętrza}.
Dla $x\in\partial M$ zwykła grupa
$H_n(M,M\setminus\{x\};\Z)$ jest zerowa:
półkula i półkula z usuniętym środkiem płaskiej ściany
są ściągalne (dla $n\geq1$). Nie daje więc lokalnej
orientacji brzegu. Zamiast tego brzeg klasy względnej
spełnia $\partial[M,\partial M]=[\partial M]$
przy orientacji „zewnętrzna normalna najpierw”.

Wyjaśnijmy jedyność w przypadku
zamkniętym. W każdym $n$-cyklu
$z=\sum a_\sigma\sigma$ warunek
$\partial z=0$ przy wspólnej
$(n-1)$-ścianie mówi, że
współczynniki $a_\sigma$ po obu
jej stronach są równe
(po uwzględnieniu orientacji).
Graf sąsiedztwa najwyższych
sympleksów jest spójny,
gdy $M$ jest spójna i $n\geq1$:
droga między ich wnętrzami
może po małej perturbacji
omijać szkielet wymiaru
$n-2$ i kolejno przecinać
ściany kowymiaru jeden.
Wszystkie współczynniki są
więc jednakowe. Ponieważ
$C_{n+1}(K)=0$, nie ma
nietrywialnych symplicjalnych
brzegów w stopniu $n$.
Stąd $H_n(M;\Z)=\Z[c]$. Dla $n=0$ spójna
zamknięta rozmaitość jest punktem i wynik jest bezpośredni.
W przypadku względnym to samo porównanie współczynników
przez ściany wewnętrzne daje jedyność klasy względnej.
Homologia symplicjalna jest
zgodna z singularną na
triangulacji, jak pokazano
w Rozdziale~\ref{ch:homologia-przestrzeni}.
Jeśli porzucimy znaki i pracujemy
nad $\mathbb F_2$, dwie ściany
sumują się do $0$ bez wybierania
orientacji. Ten sam argument
daje klasę modulo $2$.
\end{proof}

\begin{przyklad}
Na $S^2$ suma zorientowanych
trójkątów daje $[S^2]$.
Na torusie analogiczna suma
daje $[T^2]$. Dla wstęgi
Möbiusa nie istnieje
klasa $[M,\partial M]$ nad $\Z$
o lokalnych współczynnikach
równych $+1$ w każdym punkcie:
po obejściu rdzenia orientacja
zmienia znak. Klasa modulo $2$
istnieje. Poprawną wersję
całkowitą uzyskamy z lokalnymi
współczynnikami w
Sekcji~\ref{sec:lokalne-wsp-16}.
\end{przyklad}

\section{Dualność Poincarégo przez komórki dualne}
\label{sec:dualnosc-komorkowa-16}

Niech $\sigma$ będzie $k$-sympleksem
triangulacji $K$ zamkniętej
$n$-rozmaitości. Po podziale
barycentrycznym bierzemy wszystkie
małe sympleksy odpowiadające
łańcuchom ścian
$\sigma=\sigma_k\subsetneq
\sigma_{k+1}\subsetneq\cdots
\subsetneq\sigma_n$.
Ich suma, z odpowiednimi
ścianami, tworzy
$(n-k)$-wymiarową
\emph{komórkę dualną} $D(\sigma)$.
Komórka $D(\sigma)$ przecina
$\sigma$ w barycentrum i nie
przecina wnętrza innych
$k$-sympleksów. W triangulacji zgodnej ze strukturą gładką link
$\sigma$ jest sferą PL (dla najwyższych sympleksów
komórki dualne są punktami), więc
$D(\sigma)$ jest domkniętym
dyskiem, a jego brzeg składa
się z komórek dualnych do
$(k+1)$-sympleksów
zawierających $\sigma$.

\begin{figure}[htbp]
\centering
\begin{tikzpicture}[>=Latex,font=\small,line cap=round,scale=1.08]
 \coordinate (o) at (0,0);
 \coordinate (a) at (-2,-1.65);
 \coordinate (b) at (2,-1.65);
 \coordinate (c) at (2,1.65);
 \coordinate (d) at (-2,1.65);
 \fill[manifoldgold!18]
  (-1,-.825)--(0,-1.10)--(1,-.825)--
  ({4/3},0)--(1,.825)--(0,1.10)--
  (-1,.825)--({-4/3},0)--cycle;
 \draw[manifoldblue!65,thick] (a)--(b)--(c)--(d)--cycle;
 \foreach \v in {a,b,c,d}
  \draw[manifoldblue!65,thick] (o)--(\v);
 \draw[manifoldgold!85!black,very thick]
  (-1,-.825)--(0,-1.10)--(1,-.825)--
  ({4/3},0)--(1,.825)--(0,1.10)--
  (-1,.825)--({-4/3},0)--cycle;
 \foreach \x/\y in {0/-1.10,0/1.10}
  \fill[manifoldgold!85!black] (\x,\y) circle (1.4pt);
 \fill[manifoldgold!85!black] ({4/3},0) circle (1.4pt);
 \fill[manifoldgold!85!black] ({-4/3},0) circle (1.4pt);
 \fill[manifoldred] (o) circle (2.5pt)
  node[below right=2pt] {$\sigma^0$};
 \node[manifoldgold!65!black] at (0,.65) {$D(\sigma^0)$};
 \node[manifoldblue,anchor=south] at (0,1.78) {triangulacja};
\end{tikzpicture}
\caption{Niebieskie trójkąty stanowią fragment
triangulacji powierzchni. Złoty dysk jest
komórką dualną do środkowego wierzchołka.
Analogicznie krawędź ma dualną krawędź
przecinającą ją poprzecznie.}
\label{fig:dualna-komorka-16}
\end{figure}

\begin{twierdzenie}[Dualność Poincarégo]
\label{thm:poincare-dualnosc-16}
Dla zamkniętej zorientowanej
rozmaitości $M^n$ i dowolnych
abelowych grup współczynników $G$ mnożenie
kapowe przez klasę fundamentalną
jest izomorfizmem
\[
 \operatorname{PD}:H^k(M;G)
       \stackrel{\cong}{\longrightarrow}
       H_{n-k}(M;G),\qquad
 \alpha\longmapsto[M]\frown\alpha.
\]
Dla zwartej zorientowanej
rozmaitości z brzegiem
odpowiednie izomorfizmy to
\[
 H^k(M;G)\cong H_{n-k}(M,\partial M;G),
 \qquad
 H^k(M,\partial M;G)\cong H_{n-k}(M;G).
\]
Przy $G=\mathbb F_2$ orientacja
nie jest potrzebna.
\end{twierdzenie}
\begin{proof}
Najpierw niech $\partial M=\varnothing$.
Wybieramy dualny rozkład komórkowy
$K^\vee$. Ustalmy orientacje sympleksów. Dla $G=\Z$ niech
$\sigma^\#$ będzie kołańcuchem równym $1$ na $\sigma$
i $0$ na pozostałych sympleksach tego wymiaru.
Orientujemy $D(\sigma)$ tak, aby kolejność
$(T\sigma,TD(\sigma))$ dawała orientację $M$.
Dla dowolnej grupy abelowej przenosimy po prostu
współczynnik $g\in G$ przy $\sigma^\#$ na tę komórkę.
Daje to izomorfizmy grup
\[
 \mathcal D_k:C^k(K;G)\longrightarrow C_{n-k}(K^\vee;G).
\]
Rachunek brzegu ma konkretny znak. Jeśli $\tau^{k+1}$
zawiera $\sigma^k$ i $u$ wskazuje od $\sigma$ do wnętrza
$\tau$, to konwencja brzegu daje
$\operatorname{or}(\tau)=[\tau:\sigma](-u,\operatorname{or}(\sigma))$.
Po przeniesieniu $u$ za $k$ wektorów otrzymujemy
$[\tau:\sigma](-1)^{k+1}(\operatorname{or}(\sigma),u)$.
W komórce dualnej kierunek $u$ wskazuje na zewnątrz
ku ścianie $D(\tau)$. Porównanie orientacji $M$ daje
\begin{equation}\label{eq:dualny-lancuch-16}
 \partial\mathcal D_k=(-1)^{k+1}\mathcal D_{k+1}\dd.
\end{equation}
Mimo znaku izomorfizmy $\mathcal D_k$ przenoszą
kocykle na cykle i kobrzegi na brzegi, zatem dają
$H^k(K;G)\cong H_{n-k}(K^\vee;G)$.
Jeśli chcemy dosłownego izomorfizmu kompleksów
z różniczkami $\dd$ i $\partial$, należy użyć
$(-1)^{k(k+1)/2}\mathcal D_k$. Zachowujemy jednak
nieprzeskalowane $\mathcal D_k$, zgodne ze znakiem
iloczynu kapowego w~\eqref{eq:kapowy-brzeg-16}.

Wyjaśnijmy utożsamienie z $[M]\frown\alpha$.
W podziale barycentrycznym obie konstrukcje dla
$\sigma^\#$ mają nośnik w zamkniętej gwieździe $\sigma$.
W stopniu $k=n$ dają punkty z tym samym współczynnikiem
orientacji. Różnicę można wypełnić drogą w tej gwieździe.
Schodząc kolejno z $k$ do $k-1$, z obu wzorów na brzeg
wynika, że różnica po odjęciu już zbudowanej homotopii
na kościanach jest cyklem w tej samej gwieździe.
Gwiazda jest ściągalna, więc cykl można wypełnić.
To indukcja analogiczna do dowodu przemienności
iloczynu kubkowego; daje homotopię po sprowadzeniu
obu konstrukcji do wspólnego podziału. Wyznaczają
więc tę samą mapę na kohomologii, czyli $\operatorname{PD}$.
Oba rozkłady komórkowe obliczają homologię $M$.

Gdy $\partial M\ne\varnothing$,
dualna komórka przy brzegu
ma dodatkową ścianę
\emph{leżącą w brzegu}.
Po przejściu do łańcuchów
względnych ta ściana znika;
daje to pierwszy z podanych
izomorfizmów. Z kolei
kołańcuchy względne znikają
na sympleksach brzegowych,
więc ich dualne komórki
nie mają obciętej ściany:
tworzą podkompleks dualny $L$ z komórek
$D(\sigma)$ dla $\sigma\not\subset\partial M$.
Ten podkompleks oblicza homologię $M$.
Istotnie, $M\setminus\partial M$ cofa się na $L$:
w każdym barycentrycznym sympleksie zerujemy wagi
wierzchołków odpowiadających ścianom brzegowym
i normalizujemy pozostałe. Ich suma jest dodatnia
poza brzegiem, a homotopia przez liniową zmianę wag
jest zgodna na ścianach. Włączenie wnętrza do $M$
jest równoważnością homotopijną, bo kołnierz pozwala
odsunąć brzeg do wnętrza. Stąd drugi izomorfizm
z homologii bezwzględnej.
W każdym przypadku
sprawdzenie~\eqref{eq:dualny-lancuch-16}
jest tym samym lokalnym
rachunkiem incydencji;
kołnierz wynikający z twierdzenia o kołnierzu brzegu
zapewnia standardowy model
półprzestrzenny przy brzegu.
Nad $\mathbb F_2$ wszystkie
znaki orientacji znikają.
\end{proof}

\begin{wniosek}[Parowanie i charakterystyka Eulera]
\label{cor:pd-parowanie-16}
Dla zamkniętej zorientowanej
$M^n$ i ciała $\mathbb F$
parowanie
\[
 H^k(M;\mathbb F)\times
 H^{n-k}(M;\mathbb F)
 \longrightarrow\mathbb F,\qquad
 (\alpha,\beta)\longmapsto
      \langle\alpha\smile\beta,[M]\rangle
\]
jest niezdegenerowane. Zatem
$b_k=b_{n-k}$. Jeśli $n$
jest nieparzyste, $\chi(M)=0$.
\end{wniosek}
\begin{proof}
Dualność identyfikuje
$H^k$ z $H_{n-k}$.
Nad ciałem każda niezerowa
klasa homologii jest wykrywana
przez pewną klasę kohomologii
z $H^{n-k}$, bo dual przestrzeni
wektorowej rozdziela punkty.
Definicja iloczynu kapowego
i kubkowego utożsamia tę
ewaluację z podaną liczbą.
Stąd parowanie jest
niezdegenerowane i w skończonym
wymiarze obie strony mają
równe wymiary. Jeśli $n$
jest nieparzyste, w sumie
$\chi(M)=\sum_k(-1)^kb_k$
wyrazy dla $k$ i $n-k$
mają przeciwne znaki,
więc sumują się do zera.
\end{proof}

\begin{przyklad}[Torus i iloczyn dwóch sfer]
Na $T^2$ klasy
$a,b\in H^1(T^2;\Z)$,
odpowiadające dwóm kierunkom,
spełniają
$\langle a\smile b,[T^2]\rangle=1$
po ustaleniu orientacji.
Na $S^2\times S^2$ dwie
klasy $A,B\in H_2(-;\Z)$
reprezentowane przez
$S^2\times\{*\}$ i
$\{*\}\times S^2$
mają macierz przecięć
$\begin{psmallmatrix}0&1\\1&0\end{psmallmatrix}$.
Samych rang $H_2\cong\Z^2$
nie wystarcza do odtworzenia
tej macierzy.
\end{przyklad}

\section{Liczba przecięcia: znak i niezmienniczość}
\label{sec:przeciecia-16}

Niech $A^p,B^{n-p}\subset M^n$ będą
zamkniętymi zorientowanymi
podrozmaitościami zamkniętej
zorientowanej $M$. Po małej
perturbacji, możliwej dzięki twierdzeniu o aproksymacji transwersalnej,
możemy założyć $A\pitchfork B$.
Przecięcie ma wymiar zero,
a ze zwartości jest skończone.

\begin{definicja}[Znak i liczba przecięcia]
\label{def:liczba-przeciecia-16}
Dla $x\in A\cap B$ kładziemy
$\epsilon_x=+1$, jeżeli
uporządkowana baza
$T_xA$ \emph{a następnie}
$T_xB$ daje orientację $T_xM$,
i $\epsilon_x=-1$ w przeciwnym
przypadku. \emph{Liczba
przecięcia} to
\[
 I(A,B)=\sum_{x\in A\cap B}\epsilon_x.
\]
Bez orientacji liczymy
$\#(A\cap B)\pmod 2$.
\end{definicja}

\begin{figure}[htbp]
\centering
\begin{tikzpicture}[>=Latex,font=\small,line cap=round,scale=1.04]
 \fill[manifoldblue!5] (-2.1,-1.65) rectangle (2.1,1.65);
 \draw[manifoldblue!45,thick]
   (-2.1,-1.65) rectangle (2.1,1.65);
 \draw[manifoldblue,very thick]
   (-2.1,0)--(2.1,0);
 \draw[manifoldblue,very thick,->]
   (1.25,0)--(1.75,0);
 \draw[manifoldgold!85!black,very thick]
   (0,-1.65)--(0,1.65);
 \draw[manifoldgold!85!black,very thick,->]
   (0,.60)--(0,1.10);
 \fill[manifoldred] (0,0) circle (2.6pt);
 \node[manifoldred,anchor=south west] at (.10,.08) {$x,\;\epsilon_x=+1$};
 \node[manifoldblue,anchor=north] at (1.35,-.12) {$A$};
 \node[manifoldgold!75!black,anchor=west] at (.12,1.15) {$B$};
 \node[anchor=south] at (0,1.72)
       {przeciwległe brzegi są sklejone};
\end{tikzpicture}
\caption{Dwa podstawowe okręgi torusa
przedstawionego jako kwadrat z parami
sklejonych boków. W punkcie $x$ kierunek
poziomy, po nim pionowy, daje dodatnią
orientację. Liczba przecięcia wynosi $1$.}
\label{fig:przeciecie-torus-16}
\end{figure}

\begin{twierdzenie}[Przecięcie i iloczyn kubkowy]
\label{thm:przeciecie-kubek-16}
Niech $\alpha_A=\operatorname{PD}^{-1}
(i_{A*}[A])\in H^{n-p}(M;\Z)$
i $\alpha_B=\operatorname{PD}^{-1}
(i_{B*}[B])\in H^p(M;\Z)$.
Przy konwencji orientowania przecięcia
z Definicji~\ref{def:liczba-przeciecia-16}
zachodzi
\begin{equation}\label{eq:przeciecie-kubek-16}
 I(A,B)=
 \langle\alpha_A\smile\alpha_B,[M]\rangle.
\end{equation}
Liczba zależy więc tylko od klas
homologii obu podrozmaitości;
jest niezmienna przy izotopii.
Jeżeli jest niezerowa, żadna
perturbacja w tych klasach nie
może rozdzielić $A$ i $B$.
\end{twierdzenie}
\begin{proof}
Użyjemy klasy Thoma z Twierdzenia~\ref{thm:euler-thom-16}.
Jej zamieszczony niżej dowód jest komórkowy i nie
korzysta z obecnego twierdzenia ani z klas $w_i$.
Orientujemy normalne warunkiem „normalna, potem styczna”.
Wycięcie w otoczeniu tubularnym daje klasy
$u_A\in H^{n-p}(M,M\setminus A;\Z)$ i
$u_B\in H^p(M,M\setminus B;\Z)$.
Ich obrazy w kohomologii bezwzględnej są $\alpha_A,\alpha_B$:
lokalnie oba opisy przypisują $+1$ dodatnio zorientowanemu
normalnemu dyskowi, co jest konstrukcją dualnej klasy.
Względny iloczyn ma nośnik na przecięciu, tzn.
\[
 u_A\smile u_B\in H^n(M,M\setminus(A\cap B);\Z).
\]
Wynika to z wersji względnej wzoru kubkowego:
po dostatecznym podziale sympleksy poza przecięciem
leżą w jednym z dwóch otwartych dopełnień,
gdzie odpowiedni czynnik znika.
Wycięcie rozkłada ostatnią grupę na sumę grup
$H^n(B_x,B_x\setminus\{x\};\Z)\cong\Z$
dla małych rozłącznych kul wokół punktów przecięcia.
Ewaluacja na $[M]$ jest sumą tych lokalnych liczb.
Jest to argument całkowity, nie tylko rachunek form
rzeczywistych, który mógłby zgubić torsję.

Oto kontrola znaku.
Niech $q=n-p$, a
$(a_1,\ldots,a_p)$
i $(b_1,\ldots,b_q)$
będą dodatnimi
bazami $T_xA$ i $T_xB$.
Załóżmy najpierw,
że $(a,b)$ orientuje
$T_xM$ dodatnio.
Normalną do $A$
orientujemy warunkiem
$(\nu_A,a)>0$,
a normalną do $B$
warunkiem
$(\nu_B,b)>0$.
Wówczas odpowiadająca
pierwszej normalnej
$q$-forma $\omega_A$
ma na bloku $b$
wartość $(-1)^{pq}$,
a $p$-forma $\omega_B$
drugiej normalnej
ma na bloku $a$
wartość $+1$.
We wzorze na
$(\omega_A\wedge\omega_B)(a,b)$
przestawienie bloku
$b$ przed $a$
daje drugi czynnik
$(-1)^{pq}$.
Iloczyn dwóch znaków
wynosi $+1$.
Odwrócenie orientacji
jednego z $A,B$
zmienia zarówno
$\epsilon_x$, jak
i tę wartość na $-1$.
Zatem lokalna ewaluacja
iloczynu jest $\epsilon_x$.
Po zsumowaniu po rozłącznych
małych otoczeniach punktów
przecięcia otrzymujemy
\eqref{eq:przeciecie-kubek-16}.
W powyższym rachunku formy $\omega_A,\omega_B$
opisują orientacje liniowych bloków; całkowite
generatory są ustalone przez klasy Thoma.
Poza tymi kulami nie ma wkładu dzięki względnemu
nośnikowi iloczynu. Ten sam dowód modulo $2$
działa bez orientacji. Dla transwersalnych przecięć
dodatniego wymiaru identyfikuje iloczyn klas Thoma
z klasą Thoma sumy normalnych; daje analogiczną
interpretację iloczynu jako klasy dualnej do przecięcia.

Prawa strona zależy jedynie
od $i_{A*}[A]$ i $i_{B*}[B]$,
ponieważ dualność Poincarégo
i iloczyn kubkowy są
określone na klasach.
Izotopijne zanurzenia
indukują homotopijne
odwzorowania w $M$,
więc dają tę samą
klasę homologii.
Gdyby obrazy dało się
rozłączyć, suma po
pustym przecięciu
wynosiłaby $0$.
\end{proof}

\begin{przyklad}[Krzywe na torusie]
W standardowej bazie
$H_1(T^2;\Z)=\Z a\oplus\Z b$
krzywe o klasach
$pa+qb$ i $ra+sb$
mają przecięcie algebraiczne
\[
 I\bigl((p,q),(r,s)\bigr)
 =\det\begin{pmatrix}p&r\\q&s\end{pmatrix}
 =ps-qr.
\]
Równość wynika z dwuliniowości
iloczynu i z
$I(a,b)=1$, $I(b,a)=-1$,
$I(a,a)=I(b,b)=0$.
Na przykład $(1,1)$ i $(1,-1)$
mają liczbę przecięcia $-2$:
co najmniej dwa przecięcia,
licząc bez znaków, są
nieuniknione w tych klasach.
Liczba $0$ nie gwarantuje
rozłączności konkretnych
przedstawicieli: dodatnie
i ujemne przecięcia mogą
się znosić.
\end{przyklad}

\begin{uwaga}[Samoprzecięcie]
Dwie kopie tej samej
podrozmaitości nie są
transwersalne, gdy się
pokrywają. Jedną z nich
przesuwamy w małym
otoczeniu tubularnym
za pomocą przekroju
wiązki normalnej.
Jeżeli przekrój ma zera,
przesunięta kopia
przecina pierwotną
właśnie w tych zerach.
Ta obserwacja prowadzi
do klasy Eulera
w Sekcji~\ref{sec:klasy-eulera-16}.
\end{uwaga}

\section{Lokalne współczynniki i brak orientacji}
\label{sec:lokalne-wsp-16}

Zwykła grupa współczynników
$\Z$ przyjmuje, że liczba
$+1$ oznacza to samo w
każdym punkcie. Na rozmaitości
nieorientowalnej lokalny
generator orientacji wraca
po pewnej pętli jako $-1$.
Trzeba więc przenosić
współczynniki wraz z drogą.

\begin{definicja}[System lokalny i system orientacji]
\label{def:system-lokalny-16}
\emph{System lokalny} grup abelowych
$\mathcal L$ przypisuje punktowi
$x$ grupę $\mathcal L_x$,
a drodze $\gamma:x\to y$
izomorfizm
$T_\gamma:\mathcal L_x\to\mathcal L_y$,
niezmienny przy homotopii
drogi z ustalonymi końcami
i zgodny ze składaniem dróg.
Transport wzdłuż drogi stałej jest identycznością.
Dla rozmaitości bez brzegu \emph{system orientacji}
$\mathcal O_M$ ma włókno
\[
 (\mathcal O_M)_x
  =H_n(M,M\setminus\{x\};\Z)
  \cong\Z;
\]
transport lokalnej orientacji
wzdłuż drogi jest mnożeniem
przez $+1$ albo $-1$. Przy brzegu definiujemy
system przez orientacje $T_xM$, równoważnie przez
przedłużenie systemu z wnętrza w kołnierzu.
Nie stosujemy tam powyższego wzoru na lokalną homologię.
\end{definicja}

W łańcuchach z $\mathcal L$
sympleks singularny ma
współczynnik w grupie
nad jego pierwszym
wierzchołkiem. Przy operatorze
brzegu współczynnik ściany
pomijającej pierwszy
wierzchołek przenosimy
wzdłuż pierwszej krawędzi
do nowego wierzchołka;
dla pozostałych ścian
punkt bazowy się nie
zmienia. Znaki przemienne
są takie jak w zwykłym
operatorze $\partial$.
Dwie drogi transportu
przy dwukrotnym usunięciu
wierzchołków są homotopijne
w obrębie sympleksu,
więc odpowiadające
wyrazy w $\partial^2$
znoszą się parami.
Tak powstaje
$H_*(M;\mathcal L)$.

\begin{przyklad}[Skręcone współczynniki na okręgu]
Niech $\mathcal L$ nad $S^1$
ma włókno $\Z$ i transport
po pełnym okrążeniu
równy $-1$. W komórkowym
modelu okręgu jest
jedna $0$-komórka
i jedna $1$-komórka.
Przy zwykłych
współczynnikach końce
krawędzi dają
$\partial_1=1-1=0$.
Tutaj jeden koniec
wraca ze znakiem
$-1$, więc
$\partial_1=(-1)-1=-2$.
Stąd
\[
 H_1(S^1;\mathcal L)=0,\qquad
 H_0(S^1;\mathcal L)=\Z/2.
\]
Po przejściu do
$\mathbb F_2$ macierz
$-2$ staje się zerem
i dostajemy zwykłe
grupy okręgu.
\end{przyklad}

\begin{twierdzenie}[Dualność ze skręconymi współczynnikami]
\label{thm:pd-skrecenie-16}
Dla spójnej zamkniętej
$n$-rozmaitości, także
nieorientowalnej, istnieje
kanoniczna skręcona klasa
fundamentalna
$[M]_{\mathcal O}\in
H_n(M;\mathcal O_M)$.
Dla stałej abelowej grupy
współczynników $G$
otrzymujemy izomorfizm
\[
 H^k(M;G)\cong
 H_{n-k}(M;\mathcal O_M\otimes G).
\]
Dla zwartej rozmaitości z brzegiem używamy
klasy $[M,\partial M]_{\mathcal O}$ i otrzymujemy
obie wersje względne:
\[
 H^k(M;G)\cong H_{n-k}(M,\partial M;\mathcal O_M\otimes G),
 \qquad
 H^k(M,\partial M;G)\cong H_{n-k}(M;\mathcal O_M\otimes G).
\]
Po zredukowaniu
współczynników modulo $2$
system orientacji staje
się stały.
\end{twierdzenie}
\begin{proof}
W każdej małej mapie
wybieramy generator
lokalnej homologii.
Gdy przechodzimy
do sąsiedniej mapy,
ewentualna zmiana
orientacji o $-1$
zmienia jednocześnie
współczynnik w
$\mathcal O_M$ o $-1$.
Iloczyn tych dwóch
znaków jest $+1$.
Suma wszystkich
najwyższych sympleksów
z tak dobranymi
lokalnymi współczynnikami
ma zatem zerowy brzeg
na każdej wspólnej
ścianie i daje
$[M]_{\mathcal O}$.
W małej kuli jej obraz
jest lokalnym generatorem.

W dowodzie
Twierdzenia~\ref{thm:poincare-dualnosc-16}
orientacja $M$ była
potrzebna do orientowania
komórek $D(\sigma)$.
Teraz orientujemy
$D(\sigma)$ za pomocą
lokalnego generatora
$(\mathcal O_M)_x$.
Przejście do sąsiedniej
komórki wnosi dokładnie
znak transportu
w $\mathcal O_M$.
Rachunek incydencji
\eqref{eq:dualny-lancuch-16}
pozostaje więc prawdziwy
z tymi współczynnikami.
Znowu dostajemy
izomorfizmy przenoszące jądra i obrazy
(ze znakiem zależnym od stopnia), a więc ich homologie.
Na brzegu przechodzimy
do łańcuchów względnych
tak samo jak w dowodzie
wersji zorientowanej.
W $\mathbb F_2$
mamy $-1=+1$,
więc monodromia znika.
\end{proof}

\begin{przyklad}[Wstęga Möbiusa i płaszczyzna rzutowa]
Dla wstęgi Möbiusa $W$
nie istnieje zwykła
całkowita klasa
$[W,\partial W]$
generująca orientację
we wszystkich punktach,
ale istnieje
$[W,\partial W]_{\mathcal O}
\in H_2(W,\partial W;\mathcal O_W)
\cong\Z$.
Dla $\R P^2$ mamy
$H_2(\R P^2;\Z)=0$,
za to
$H_2(\R P^2;\mathcal O)\cong\Z$
i $H_2(\R P^2;\mathbb F_2)
\cong\mathbb F_2$.
Lokalne współczynniki
odzyskują orientację
bez twierdzenia,
że przestrzeń nagle
stała się orientowalna.
\end{przyklad}

\section{Kołańcuchy z lokalnymi współczynnikami: pełna konstrukcja}
\label{sec:lokalne-kolancuchy-v4}

Powyższa definicja transportu pozwala zdefiniować homologię
ze skręconymi współczynnikami. Do dualności i ciągu
spektralnego fibracji potrzebna jest również \emph{kohomologia}
ze skręconymi współczynnikami. Ustalamy przestrzeń $X$ i
system lokalny grup abelowych $\mathcal L$ z transportami
$T_\gamma:\mathcal L_x\to\mathcal L_y$.

\begin{definicja}[Skręcony kompleks singularny]
Dla singularnego $k$-sympleksu $\sigma:\Delta^k\to X$
piszemy $v_j=\sigma(e_j)$ i $\sigma_j$ dla jego $j$-tej
ściany. Łańcuch stopnia $k$ jest skończoną sumą symboli
$a\sigma$, $a\in\mathcal L_{v_0}$. Operator brzegu jest dany
przez
\[
\partial(a\sigma)=T_{[v_0,v_1]}(a)\,\sigma_0
+\sum_{j=1}^k(-1)^j a\,\sigma_j.
\]
Przez $C^k(X;\mathcal L)$ oznaczamy \emph{wszystkie} rodziny
$c(\sigma)\in\mathcal L_{v_0}$ (bez warunku skończonego nośnika)
przypisane singularnym $k$-sympleksom $\sigma$.
Koróżniczka to
\[
(\delta c)(\sigma)=T_{[v_0,v_1]}^{-1}c(\sigma_0)
+\sum_{j=1}^{k+1}(-1)^j c(\sigma_j),
\qquad\sigma:\Delta^{k+1}\to X.
\]
Definiujemy $H^k(X;\mathcal L)=\ker\delta^k/\operatorname{im}
\delta^{k-1}$ oraz analogicznie $H_k(X;\mathcal L)$.
\end{definicja}

\begin{stwierdzenie}
\label{prop:local-delta-v4}
Dla wszystkich $k\ge0$ zachodzą $\partial^2=0$ i
$\delta^2=0$. Każde ciągłe $f:Y\to X$ indukuje
$f^*:H^k(X;\mathcal L)\to H^k(Y;f^*\mathcal L)$.
\end{stwierdzenie}
\begin{proof}
W rozwinięciu podwójnego brzegu każda ściana
kowymiaru $2$ występuje dwa razy: po usunięciu
wierzchołków $i<j$ w kolejności $i,j$ i $j,i$.
Odpowiadające znaki są przeciwne. Gdy oba usunięcia
zmieniają punkt bazowy współczynnika, transporty
wzdłuż dróg $[v_0,v_1]*[v_1,v_2]$ i $[v_0,v_2]$
są identyczne, ponieważ krawędzie leżą na brzegu
$2$-sympleksu i drogi są w nim homotopijne z ustalonymi
końcami. Dlatego również współczynniki zgadzają się.
To dowodzi $\partial^2=0$.

Dla $\delta^2$ te same pary usuniętych wierzchołków
występują z przeciwnymi znakami. W pierwszej ścianie
korzystamy z odwrotnego transportu, a równość
$T_{[v_0,v_1]}^{-1}T_{[v_1,v_2]}^{-1}
=T_{[v_0,v_2]}^{-1}$ jest konsekwencją poprzedniej
homotopii trójkątnej. Pozostałe wyrazy znoszą się tak
jak w zwykłych kołańcuchach. Dla ciągłego $f$
kładziemy $(f^*c)(\tau)=c(f\circ\tau)$;
transporty $f^*\mathcal L$ zostały zdefiniowane
przez transporty w $\mathcal L$ wzdłuż $f\circ\gamma$.
Dlatego $\delta f^*=f^*\delta$.
\end{proof}

\begin{stwierdzenie}[Monodromia w stopniu zerowym]
Niech $X$ będzie łukowo spójna, $x_0\in X$,
$A=\mathcal L_{x_0}$ i
$\rho:\pi_1(X,x_0)\to\operatorname{Aut}(A)$ monodromią.
Wówczas $H^0(X;\mathcal L)\cong A^{\pi_1(X,x_0)}$
oraz $H_0(X;\mathcal L)\cong A_{\pi_1(X,x_0)}$,
gdzie pierwsza grupa składa się z elementów niezmienniczych,
a druga jest ilorazem przez relacje $\rho(g)a-a$.
\end{stwierdzenie}
\begin{proof}
Warunek $\delta c=0$ na dowolnej drodze od $x$ do $y$
stanowi dokładnie równanie $c(y)=T_\gamma c(x)$.
Wartość $c(x_0)$ wyznacza całą rodzinę, a jej
niezależność od drogi jest równoważna niezmienniczości
względem pętli. Homologia zerowa jest generowana przez
symbole $[x,a]$ modulo brzeg dróg,
$[y,T_\gamma a]-[x,a]$. Po sprowadzeniu wszystkich
punktów do $x_0$ pozostają jedynie relacje
$\rho(g)a-a$.
\end{proof}

\begin{przyklad}
Dla okręgu i systemu $\mathcal L$ o monodromii
$-\id$ na $\Z$ komórkowe różniczki i koróżniczki są
mnożeniem przez $-2$. Stąd
$H_0(S^1;\mathcal L)=\Z/2$,
$H_1(S^1;\mathcal L)=0$,
$H^0(S^1;\mathcal L)=0$ oraz
$H^1(S^1;\mathcal L)=\Z/2$.
\end{przyklad}

\section{Dualność Poincarégo--Lefschetza dla rozmaitości z brzegiem}
\label{sec:pl-v4}
\begin{twierdzenie}[Dwie wersje dualności]
\label{thm:pl-v4}
Niech $M$ będzie zwartą gładką $n$-rozmaitością z brzegiem
$\partial M$ (brzeg może być pusty),
$\mathcal L$ systemem lokalnym grup abelowych, a
$\mathcal O_M$ systemem orientacji. Istnieje kanoniczna
klasa $[M,\partial M]_{\mathcal O}\in
H_n(M,\partial M;\mathcal O_M)$, której iloczyn kapowy
daje, dla $0\le k\le n$, naturalne izomorfizmy
\[
H^k(M;\mathcal L)\cong
H_{n-k}(M,\partial M;\mathcal L\otimes\mathcal O_M),\qquad
H^k(M,\partial M;\mathcal L)\cong
H_{n-k}(M;\mathcal L\otimes\mathcal O_M).
\]
Dla zorientowanego $M$ i stałego systemu $G$
można usunąć czynnik $\mathcal O_M$ po wyborze orientacji.
\end{twierdzenie}
\begin{proof}
\emph{Krok 1: klasa fundamentalna.}
Wybierzmy skończoną triangulację gładkiej rozmaitości,
w której $\partial M$ jest podkompleksem. Każdemu
$n$-sympleksowi $\sigma$ nadajemy orientację i dobieramy
współczynnik z włókna $\mathcal O_M$ tak, aby lokalnie
para tych orientacji odpowiadała kanonicznej lokalnej
jednostce. Na wspólnej ścianie dwóch najwyższych
sympleksów ściany indukowane mają przeciwne orientacje,
a transport współczynnika koryguje ewentualne odwrócenie
orientacji rozmaitości. Zatem suma sympleksów ma brzeg
wyłącznie w $\partial M$. Lokalnie, w małej kuli we wnętrzu,
jej obraz jest kanonicznym generatorem z lokalnym
systemem orientacji. Tak otrzymujemy klasę względną.

\emph{Krok 2: dualny podział komórkowy.}
Po barycentrycznym podziale triangulacji, dla $k$-simpleksu
$\sigma$ bierzemy sumę sympleksów, których najniższą
ścianą jest $\sigma$. Otrzymana komórka dualna ma wymiar
$n-k$. Każdy łańcuch takich małych sympleksów znajduje się
w gwieździe $\sigma$; gwiazda jest ściągalna, więc
transport $\mathcal L$ między jej punktami nie zależy od
wybranej w niej drogi. Orientacja $\sigma$ wraz z lokalną
orientacją $M$ określa orientację dualnej komórki.
Na komórce wspólnej dwóm dualnym komórkom znak incydencji
jest, po stałym przeskalowaniu zależnym od stopnia,
transpozycją znaku incydencji pierwotnych sympleksów.
To daje izomorfizm kompleksu komórkowych kołańcuchów
w stopniu $k$ z kompleksem łańcuchów dualnych w stopniu
$n-k$, z różniczką odwracającą stopień.

\emph{Krok 3: pierwszy izomorfizm.}
Dla $C^k(M;\mathcal L)$ dopuszczamy również kołańcuchy
na sympleksach brzegu. Dualne komórki przy brzegu mają
odpowiednie fragmenty granicy w $\partial M$.
Po przejściu do kompleksu względnego
$C_{n-k}(M,\partial M;\mathcal L\otimes\mathcal O_M)$
fragmenty te znikają. Dopasowanie dualnych komórek
z kołańcuchami jest bijekcją baz, a zgodność znaków
incydencji z poprzedniego kroku zapewnia zgodność z
różniczkami. Przechodzimy do homologii.

\emph{Krok 4: drugi izomorfizm.}
Względne kołańcuchy $C^k(M,\partial M;\mathcal L)$
zerują się na sympleksach brzegowych. Odpowiadają im
wewnętrzne komórki dualne. Ich suma tworzy podkompleks
komórkowy, którego realizacja jest homotopijnie równoważna
$\operatorname{int}M$ (w barycentrycznych współrzędnych
odpychamy składową brzegową). Kołnierz brzegu daje
równoważność homotopijną $\operatorname{int}M\simeq M$.
Znów otrzymujemy zgodność kompleksów i drugą wersję dualności.

\emph{Krok 5: identyfikacja z iloczynem kapowym.}
Wzór $c\frown\alpha$ z początku rozdziału przecina
pierwotny i dualny sympleks w jednym barycentrum,
odtwarzając generator komórki dualnej wraz z jego
współczynnikiem i znakiem. W przypadku kolejnych podziałów
barycentrycznych ewentualna różnica map łańcuchowych
jest wypełniana indukcyjnie w ściągalnej gwieździe
odpowiedniego sympleksu. Wskutek tego obie mapy
indukują ten sam homomorfizm na homologii.
Zgodność z długimi ciągami pary wynika ze wzoru na
brzeg iloczynu kapowego i
$\partial[M,\partial M]_{\mathcal O}=[\partial M]_{\mathcal O}$,
z odpowiednią konwencją orientacji brzegu.
\end{proof}

\begin{przyklad}[Dysk i wstęga Möbiusa]
Dla $D^n$ mamy
$H^0(D^n;\Z)=\Z$, $H^k(D^n;\Z)=0$ dla $k>0$;
pierwszy izomorfizm daje
$H_n(D^n,S^{n-1};\Z)=\Z$. Drugi daje
$H^n(D^n,S^{n-1};\Z)\cong H_0(D^n;\Z)=\Z$.
Dla wstęgi Möbiusa $W$ istnieje
$[W,\partial W]_{\mathcal O}\in H_2(W,\partial W;\mathcal O_W)$,
choć $H_2(W,\partial W;\Z)=0$. To pokazuje, że czynnik
$\mathcal O_M$ jest konieczny, a nie dekoracyjny.
\end{przyklad}

\section{Dualność Aleksandra: kohomologia \v{C}echa i dopełnienia}
\label{sec:alexander-v4}
\begin{twierdzenie}[Dualność Aleksandra]
\label{thm:alexander-v4}
Niech $n\ge2$, $A\subsetneq S^n$ będzie niepustym zbiorem
zwartym, a $G$ grupą abelową. Wtedy dla każdego $i\ge0$
zachodzi naturalny izomorfizm
\[
\widetilde H_i(S^n\setminus A;G)
\cong\widetilde{\check H}^{\,n-i-1}(A;G),
\]
przy czym w ujemnym stopniu prawa strona jest zerowa.
Jeżeli $A$ jest zwartym wielościanem, można zastąpić
kohomologię \v{C}echa kohomologią singularną.
\end{twierdzenie}
\begin{proof}
\emph{Krok 1: regularne otoczenie.}
Najpierw niech $N\subset S^n$ będzie zwartą
$n$-wymiarową podrozmaitością z brzegiem,
$N\ne S^n$, i niech $M=S^n\setminus\operatorname{int}N$.
Dualność Poincarégo--Lefschetza
(Twierdzenie~\ref{thm:pl-v4}) daje
$H^q(N;G)\cong H_{n-q}(N,\partial N;G)$.
Wybieramy kołnierz $\partial N$ i stosujemy wycinanie,
aby otrzymać
$H_{n-q}(N,\partial N;G)\cong H_{n-q}(S^n,M;G)$.
Wycinanie wolno zastosować po niewielkim odsunięciu
względem kołnierza; samo formalne wykreślenie
$S^n\setminus N$ nie spełnia automatycznie warunku
$\overline U\subset\operatorname{int}A$ z twierdzenia o wycinaniu.

\emph{Krok 2: długi ciąg sfery.}
Długi ciąg pary $(S^n,M)$ i
$\widetilde H_j(S^n;G)=0$ dla $j\ne n$
pokazują w stopniach $0\le i\le n-2$, że
$\widetilde H_i(M;G)\cong H_{i+1}(S^n,M;G)
\cong H^{n-i-1}(N;G)$.
Gdy $i=n-1$, uwzględniamy fragment
$H_n(S^n;G)\to H_n(S^n,M;G)\to H_{n-1}(M;G)\to0$.
Po dualności pierwszy homomorfizm jest diagonalną mapą
$G\to H^0(N;G)$ przypisującą wszystkim składowym
jednakowy współczynnik. Jego kokernel to
$\widetilde H^0(N;G)$.
Daje to także przypadek $i=n-1$.
Dla $i>n$ grupy homologii $M$ znikają ze względu na
wymiar $n$. Również $H_n(M;G)=0$: każda składowa zwartej
rozmaitości $M$ ma niepusty brzeg, a dualność
Poincarégo--Lefschetza daje
$H_n(M;G)\cong H^0(M,\partial M;G)=0$.
Zatem dla $i\ge n$ obie strony są zerowe.
Zatem
$\widetilde H_i(S^n\setminus\operatorname{int}N;G)
\cong\widetilde H^{n-i-1}(N;G)$.

\emph{Krok 3: dowolny zwarty $A$.}
Wybierzmy malejące regularne otoczenia
$N_1\supset\operatorname{int}N_2\supset N_2\supset\cdots$
o gładkich brzegach, takie że
$\bigcap_jN_j=A$ oraz wnętrza $N_j$ tworzą bazę
otoczeń $A$. Istnieją, ponieważ w sferze można dobierać
skończone unie małych kul, a następnie wygładzać ich
brzegi z zachowaniem zawierań. Niech
$M_j=S^n\setminus\operatorname{int}N_j$.
Mamy $S^n\setminus A=\bigcup_j\operatorname{int}M_j$.
Każdy singularny cykl i każde skończone wypełnienie
cyklu brzegiem mają zwarty obraz; zwartość i otwartość
rosnącej rodziny $\operatorname{int}M_j$ oznaczają, że
mieszczą się w pewnym etapie. Homologia singularna
spełnia zatem
$\widetilde H_i(S^n\setminus A;G)
\cong\varinjlim_j\widetilde H_i(M_j;G)$
(włączenia $\operatorname{int}M_j\hookrightarrow M_j$
są równoważnościami homotopijnymi dzięki kołnierzom).

Izomorfizmy z kroku 2 są naturalne względem zmniejszania
otoczeń, więc przechodząc do granicy prostej otrzymujemy
$\varinjlim_j\widetilde H^{n-i-1}(N_j;G)$.
Na mocy opisu kohomologii \v{C}echa przez otoczenia
(Twierdzenie~\ref{thm:cech-neighbourhood-v4}) jest to
$\widetilde{\check H}^{\,n-i-1}(A;G)$.
Dla wielościanów stosujemy
Stwierdzenie~\ref{prop:cech-poly-v4}.
\end{proof}

\begin{wniosek}[Jordan--Brouwer]
Dla każdego osadzenia $S^{n-1}\hookrightarrow S^n$, $n\ge2$,
dopełnienie ma dokładnie dwie składowe łukowe.
\end{wniosek}
\begin{proof}
Obraz jest zwarty i homeomorficzny z $S^{n-1}$;
jego kohomologia \v{C}echa jest zwykłą kohomologią sfery.
Dla $i=0$ twierdzenie daje
$\widetilde H_0(S^n\setminus S^{n-1};\Z)\cong\Z$.
Dopełnienie jest otwarte w rozmaitości, a więc lokalnie
łukowo spójne. Jeśli ma $c$ składowych, to
$\widetilde H_0\cong\Z^{c-1}$ dla skończonego $c$;
w przypadku nieskończonym ranga byłaby nieskończona.
Zatem $c=2$.
\end{proof}

\begin{przyklad}[Dopełnienie węzła]
Dla dowolnego węzła $K\cong S^1\subset S^3$ otrzymujemy
$H_1(S^3\setminus K;\Z)\cong\Z$ oraz zanik wyższych
dodatnich homologii. Homologia dopełnienia nie odróżnia
więc splątania; subtelniejsza jest jego grupa podstawowa.
\end{przyklad}

\begin{uwaga}[Granice twierdzenia]
Nie wolno bez uzasadnienia zastępować
$\widetilde{\check H}^{\,n-i-1}(A;G)$ przez zwykłą
kohomologię singularną $A$, jeśli $A$ jest dzikim zbiorem
zwartym. Nie wolno też w miejsce dopełnienia w $S^n$
bez zmiany warunków wstawić po prostu $\R^n\setminus A$:
punkt w nieskończoności może zmieniać klasy homologii.
\end{uwaga}

\section{Wiązki nad sferami i funkcje sklejające}
\label{sec:clutching-16}

Sferę $S^m$ rozkładamy na dwa
domknięte dyski $D^m_+$ i $D^m_-$,
sklejone wzdłuż równika
$S^{m-1}$. Wiązka wektorowa
nad każdym dyskiem jest
trywialna, ponieważ dysk
jest ściągalny. Cała informacja
o wiązce leży więc w
\emph{funkcji sklejającej}
na równiku.

\begin{twierdzenie}[Klasyfikacja przez równik]
\label{thm:clutching-16}
Dla $m\geq2$ klasy izomorfizmu zorientowanych
rzeczywistych wiązek
rzędu $r\geq1$ nad $S^m$
odpowiadają klasom homotopii
\[
 [S^{m-1},SO(r)]
 \cong\pi_{m-1}(SO(r)).
\]
Funkcja $g:S^{m-1}\to SO(r)$
skleja dwie wiązki
trywialne przez
$(x,v)_+\sim(x,g(x)v)_-$.
Izomorfizmy w tej klasyfikacji zachowują orientację włókien.
Dla niezorientowanych
wiązek nad $S^1$ pozostają
dwie klasy w każdym dodatnim
rzędzie: wiązka trywialna
i wiązka z jednym skręceniem
typu Möbiusa.
\end{twierdzenie}
\begin{proof}
Wybierzmy metrykę na włóknach
i dodatnio zorientowane ramy
nad $D^m_+$ i $D^m_-$.
Na ich wspólnym równiku
ramy różnią się macierzą
$g(x)\in SO(r)$. Otrzymujemy
opisane sklejenie, więc
każda wiązka jest takiej
postaci. Wybieramy ramy gładkie, na przykład przez transport
wzdłuż promieni dysków. Wtedy $g$ jest gładka;
klasy izomorfizmu
topologicznego i gładkiego
na sferze są tu zgodne
przez gładkie przybliżenie.

Zmiana ramy na górnym
lub dolnym dysku mnoży
$g$ z lewej lub prawej
przez ograniczenie
$h_\pm:D^m_\pm\to SO(r)$.
Każde takie ograniczenie
jest homotopijne do stałej:
ściągamy dysk do punktu.
Ponieważ $SO(r)$ jest
spójna dla $r\geq2$,
stałe macierze można
dalej połączyć z
jednostkową. Zatem
izomorfizm wiązek
nie zmienia klasy
homotopii $g$.
Dla $r=1$ grupa $SO(1)$
jest trywialna i wniosek
również jest prawdziwy.

Odwrotnie, gdy $g_t$ jest homotopią funkcji
sklejających, wygładzamy ją względnie końców i sklejamy
wiązki trywialne nad
$D^m_\pm\times[0,1]$
za pomocą $g_t$.
Powstaje wiązka nad
$S^m\times[0,1]$.
Wybierzmy na niej
metrykę i koneksję
(lokalne zwykłe pochodne
można skleić podziałem
jedności). Transport
równoległy w kierunku
od $t=0$ do $t=1$
zależy gładko od punktu
sfery i daje izomorfizm
wiązek na obu końcach.
To dowodzi, że homotopijne
funkcje sklejające
dają izomorficzne wiązki.
Stałość punktu bazowego
nie zmienia klasyfikacji dla spójnej grupy $SO(r)$.
Nie tylko mapę, ale i homotopię $g_t$ można
znormalizować: $\widetilde g_t(x)=g_t(x_0)^{-1}g_t(x)$
stale przyjmuje wartość $I$ w $x_0$. To pokazuje
injektywność i surjektywność przejścia od swobodnych
do opartych klas homotopii.

Nad $S^1$ równik to
dwa punkty $S^0$.
Po zmianie ramy w
jednym z dwóch łuków
pozostaje tylko informacja,
czy transport dookoła
okręgu ma wyznacznik
$+1$ czy $-1$.
Składowa $GL_r(\R)$
o wyznaczniku dodatnim
jest spójna, podobnie
ujemna; otrzymujemy
dokładnie dwie klasy.
\end{proof}

\begin{figure}[htbp]
\centering
\begin{tikzpicture}[>=Latex,font=\small,line cap=round,scale=1.08]
 \shade[inner color=white,outer color=manifoldblue!43]
    (0,0) circle (1.55);
 \fill[white,opacity=.38] (-1.55,0)
     arc[start angle=180,end angle=360,x radius=1.55,y radius=.41]
     --cycle;
 \draw[manifoldblue!75,very thick]
     (-1.55,0)
     arc[start angle=180,end angle=360,x radius=1.55,y radius=.41];
 \draw[manifoldblue!55,dashed]
     (-1.55,0)
     arc[start angle=180,end angle=0,x radius=1.55,y radius=.41];
 \node[manifoldblue!80!black] at (0,1.05) {$D^m_+$};
 \node[manifoldblue!80!black] at (0,-1.05) {$D^m_-$};
 \draw[->,manifoldgold!85!black,very thick]
   (2.20,.54) .. controls (2.85,.55) and (2.85,-.55) .. (2.20,-.54);
 \node[anchor=west,manifoldgold!75!black] at (2.8,.03)
     {$g:S^{m-1}\to SO(r)$};
 \node[anchor=west] at (-1.30,-1.88)
     {dwie trywializacje spotykają się na równiku};
\end{tikzpicture}
\caption{Sfera jako dwa dyski z równikiem.
W każdej połowie wiązka jest iloczynem;
macierz $g(x)$ mówi, jak zidentyfikować
włókna po obu stronach równika.}
\label{fig:clutching-16}
\end{figure}

\begin{przyklad}[Dlaczego $TS^2$ jest skręcona]
Zorientowane wiązki płaszczyznowe
nad $S^2$ opisuje
$\pi_1(SO(2))=\pi_1(S^1)=\Z$.
Dla wiązki $TS^2$ wybierzmy
stereograficzne zespolone
współrzędne $z$ i $w$
na półkulach z przejściem
$w=1/z$ w pasie równika.
Pochodna przejścia
$\dd w/\dd z=-z^{-2}$
obraca ramę styczną
dwukrotnie, gdy $z$
okrąża równik raz.
Funkcja sklejająca ma
więc liczbę nawinięć
$\pm2$ (znak zależy
od wyboru orientacji
równika). Wiązka $TS^2$
nie jest trywialna:
na sferze nie istnieje
globalna rama dwóch
niezależnych pól stycznych.
\end{przyklad}

\begin{przyklad}[Grupa obrotów w wymiarze trzy]
Dla zorientowanych wiązek
rzędu $3$ nad $S^2$
otrzymujemy
$\pi_1(SO(3))\cong\Z/2$:
są dwie klasy.
Opis podwójnego nakrycia
$\operatorname{Spin}(3)\to SO(3)$
z teorii nakryć grup obrotów
pokazuje, skąd bierze się
dwójka. Pętla macierzy
obrotu reprezentująca
klasę niezerową podnosi
się do drogi kończącej
w drugim punkcie włókna
nakrycia. Takiej
funkcji sklejającej
nie można podnieść
do zamkniętej pętli
w $\operatorname{Spin}(3)$;
wiązka ma niezerowe
$w_2\in H^2(S^2;\mathbb F_2)$.
Wiązka trywialna ma
$w_2=0$.
\end{przyklad}

\section{Klasy charakterystyczne: przeszkody w kohomologii}
\label{sec:klasy-char-16}

\begin{definicja}[Idea klasy charakterystycznej]
Klasą charakterystyczną
wiązki $E\to M$ nazywamy
przypisanie $E\mapsto c(E)
\in H^k(M;G)$ zgodne
z cofnięciem:
$c(f^*E)=f^*c(E)$.
Jeżeli wiązka jest trywialna,
jej klasy dodatniego stopnia
są zerowe, bo jest cofnięciem
wiązki nad punktem.
Odwrotna implikacja
nie musi zachodzić:
klasy charakterystyczne
nie rozróżniają wszystkich
wiązek.
\end{definicja}

Najbardziej elementarna
przeszkoda dotyczy
orientacji. Nie trzeba
jej definiować przez
ogólną teorię przestrzeni
klasyfikujących.

\begin{stwierdzenie}[Pierwsza klasa Stiefela--Whitneya]
\label{prop:w1-orientacja-16}
Niech $E\to M$ będzie
rzeczywistą wiązką
rzędu $r$. Wybierzmy
trywializacje nad dobrym
pokryciem $\{U_i\}$ i
funkcje przejścia
$g_{ij}:U_i\cap U_j\to GL_r(\R)$.
Znaki
$\epsilon_{ij}=\operatorname{sgn}
\det g_{ij}\in\{\pm1\}$
wyznaczają klasę
$w_1(E)\in H^1(M;\mathbb F_2)$.
Mamy $w_1(E)=0$
wtedy i tylko wtedy,
gdy $E$ jest orientowalna.
Ponadto
$w_1(E\oplus F)=w_1(E)+w_1(F)$.
\end{stwierdzenie}
\begin{proof}
Na potrójnym przecięciu
$g_{ij}g_{jk}=g_{ik}$,
a więc
$\epsilon_{ij}\epsilon_{jk}
=\epsilon_{ik}$.
Po utożsamieniu $-1$
z $1\in\mathbb F_2$
jest to równanie jednokocyklu. Aby utożsamić jego
klasę z kohomologią singularną, można obyć się
bez dodatkowego formalizmu: mnożymy znaki przejść
wzdłuż pętli. Warunek na potrójnych przecięciach
zapewnia niezmienniczość przy podziale i homotopii
pętli, a więc homomorfizm $\pi_1(M)\to\mathbb F_2$.
Utożsamienia $H_1(M;\Z)=\pi_1(M)_{\mathrm{ab}}$
i $H^1(M;\mathbb F_2)=\operatorname{Hom}(H_1(M;\Z),\mathbb F_2)$
dają żądaną klasę na każdej składowej. Gdy homomorfizm
znika, transport orientacji z ustalonego punktu nie
zależy od drogi i nadaje zgodne orientacje lokalnym ramom.
Zmiana
ramy na $U_i$ mnoży
każdy znak $\epsilon_{ij}$
przez znak nowej ramy
na $U_i$ i na $U_j$;
odpowiada to dodaniu
kobrzegu, więc klasa
nie zależy od ram.
Jeśli $w_1(E)=0$,
kocykl znaków jest
kobrzegiem: zmieniamy
orientacje lokalnych
ram tak, aby wszystkie
macierze przejścia
miały dodatni wyznacznik.
Ramy nadają wtedy
spójną orientację
wiązce. Odwrotnie
spójna orientacja
pozwala od początku
wybrać takie ramy
i daje kocykel zerowy.
Macierz przejścia
dla $E\oplus F$
jest blokowa,
$\det(g_{ij}\oplus h_{ij})
=\det(g_{ij})\det(h_{ij})$;
znaki mnożą się, a
w $\mathbb F_2$
ich klasy się dodają.
\end{proof}

\begin{przyklad}
Wiązka Möbiusa nad $S^1$
ma monodromię $-1$,
więc $w_1$ jest
niezerowym elementem
$H^1(S^1;\mathbb F_2)$.
Trywialna linia ma $w_1=0$.
Dla rozmaitości
$w_1(TM)$ jest dokładnie
przeszkodą do jej
orientowalności.
Jest to również
monodromia systemu
$\mathcal O_M$
z Sekcji~\ref{sec:lokalne-wsp-16}.
\end{przyklad}

\begin{twierdzenie}[Pełna rodzina $w_i$]
\label{thm:wi-klasy-16}
Dla wiązki rzeczywistej
$E$ rzędu $r$ istnieją
klasy Stiefela--Whitneya
\[
 w_i(E)\in H^i(M;\mathbb F_2),
 \quad 0\leq i\leq r,\qquad
 w_0=1,\quad w_i=0\ (i>r),
\]
naturalne przy cofnięciu.
Dla sumy wiązek
\[
 w(E\oplus F)=w(E)\smile w(F),
 \qquad w(E)=1+w_1(E)+\cdots+w_r(E).
\]
W szczególności dla zanurzenia
$S\subset M$ z twierdzenia o otoczeniu tubularnym
zachodzi
\begin{equation}\label{eq:w-normalne-16}
 w(TM|_S)=w(TS)\smile w(\nu(S,M)).
\end{equation}
\end{twierdzenie}
\begin{proof}
Dla $r=0$ przyjmujemy $w(E)=1$; dalej $r\geq1$.
Korzystamy z niezależnej konstrukcji klasy Thoma
i klasy $e_2$ poniżej. Dowód tej konstrukcji nie
zakłada istnienia $w_i$; zależności wyjaśnia uwaga
po Twierdzeniu~\ref{thm:euler-thom-16}.

Uzasadnijmy najpierw potrzebny rachunek we włóknie.
Na $\R P^m$ niech $x=w_1(\lambda)$ dla linii
tautologicznej. Ograniczenie ustalonego funkcjonału
liniowego do $\lambda$ jest przekrojem $\lambda^*$,
którego zera tworzą hiperpłaszczyznę $\R P^{m-1}$.
Parzystość przecięć pętli z tą hiperpłaszczyzną
jest zmianą znaku jej podniesienia do jednostkowych
wektorów linii: funkcjonał zmienia znak przy każdym
poprzecznym przejściu przez zero. Jest to dokładnie
monodromia $w_1$, więc hiperpłaszczyzna jest dualna do $x$.
$j$ niezależnych hiperpłaszczyzn przecina się
transwersalnie, a w $\R P^j$ ich przecięcie jest
jednym punktem. Modułowa wersja rachunku przecięć daje
$\langle x^j,[\R P^j]_2\rangle=1$.
Z obliczenia komórkowego grup $\R P^m$ wiemy, że
każda grupa $H^j(-;\mathbb F_2)$ dla $0\leq j\leq m$
jest jednowymiarowa. Zatem $1,x,\ldots,x^m$ są bazą,
a $x^{m+1}=0$ z przyczyn wymiarowych.

Podajemy konstrukcję algebraiczną klas ogólnej wiązki.
Niech $P(E)$ będzie wiązką
prostych we włóknach $E$,
$\lambda\to P(E)$ wiązką
tautologiczną, a
$x=w_1(\lambda)\in
H^1(P(E);\mathbb F_2)$.
Włóknem $P(E)$ jest
$\R P^{r-1}$; klasy
$1,x,\ldots,x^{r-1}$
po ograniczeniu do
każdego włókna tworzą
bazę jego kohomologii.
Stąd
\[
 H^*(P(E);\mathbb F_2)
  =\bigoplus_{j=0}^{r-1}
     H^*(M;\mathbb F_2)\,x^j
\]
jako moduł nad $H^*(M)$.
To twierdzenie o module
można sprawdzić komórka
po komórce: nad otwartą
komórką $E$ jest
trywialna, więc
$P(E)$ jest iloczynem
z $\R P^{r-1}$.
W względnym ciągu
dokładnym dołączania
komórki mapy
$1,x,\ldots,x^{r-1}$
są izomorfizmami na
względnych grupach
każdego stopnia.
Indukcja i lemat
o pięciu odwzorowaniach
dają powyższy rozkład
na całej skończonej
triangulacji. Dla
niezwartej bazy
ten sam rachunek
przeprowadza się
na filtracji przez
szkielety lokalnie
skończonego kompleksu:
relatywne grupy przy
dołączaniu wszystkich
komórek danego stopnia
są produktami grup
dla pojedynczych komórek,
a mapy pozostają
izomorfizmami.

Klasa $x^r$ ma wobec
tej bazy jednoznaczny
rozkład. Zapisujemy
go w postaci
\[
 x^r+
   \pi^*w_1(E)x^{r-1}
     +\cdots+\pi^*w_r(E)=0.
\]
Współczynnik $w_i$
ma stopień $i$,
bo każdy wyraz ma
całkowity stopień $r$.
Jednoznaczność rozkładu
i zgodność konstrukcji
$P(E)$ z cofnięciem
dają naturalność.
Przy rozszczepieniu
$E=L_1\oplus\cdots\oplus L_r$
relacja włóknowa
ma czynniki
$\prod_j(x+w_1(L_j))$.
Oto bezpośrednie
sprawdzenie. Nad $P(E)$
tautologiczna linia
$\lambda$ wchodzi
do $\pi^*E$, więc
homomorfizm inkluzji
jest nigdzie
nieznikającym
przekrojem wiązki
$\operatorname{Hom}(\lambda,\pi^*E)
=\bigoplus_j(\lambda^*\otimes\pi^*L_j)$.
Jej klasa Eulera modulo $2$, oznaczana niezależnie
przez $e_2$, jest zatem zerowa
na mocy
Twierdzenia~\ref{thm:euler-thom-16}
i uwagi o redukcji
modulo $2$ w
Sekcji~\ref{sec:klasy-eulera-16}.
Klasa Eulera modulo $2$
linii jest jej
$w_1$, a
$w_1(\lambda^*\otimes L_j)
=x+w_1(L_j)$
z rachunku znaków
funkcji przejścia.
Multiplikatywność
klasy Eulera sumy
daje właśnie
$\prod_j(x+w_1(L_j))=0$.
Porównanie współczynników
daje
$w(E)=\prod_j(1+w_1(L_j))$.

Każdą wiązkę można
rozszczepić po cofnięciu
do wiązki pełnych flag:
po kolei wybieramy
prostą w $E$ i w
ortogonalnym do niej
dopełnieniu. W każdym
takim kroku powyższy
rozkład modułowy
zawiera składnik $1$,
więc cofnięcie
na kohomologiach
jest injektywne.
Wzór iloczynowy
sprawdzony po
rozszczepieniu
schodzi zatem na $M$
i daje twierdzenie o sumie Whitneya. Dla dwóch
wiązek najpierw rozszczepiamy $E$, potem cofniętą $F$;
oba cofnięcia są injektywne, a obie wiązki rozszczepione.
Współczynnik stopnia jeden jest $\sum_jw_1(L_j)$,
czyli wcześniej zdefiniowany znak wyznacznika przejść.
Nowa definicja $w_1$ jest więc zgodna z poprzednią.
Wreszcie
$TM|_S\cong TS\oplus\nu$
z rozszczepienia przestrzeni stycznych wzdłuż podrozmaitości;
podstawienie do wzoru
sumy daje~\eqref{eq:w-normalne-16}.
\end{proof}

\begin{uwaga}[Gdzie już pojawiło się $w_2$?]
W teorii spinowej
skonstruowaliśmy
$w_2(E)$ dla zorientowanej
wiązki jako przeszkodę
podniesienia funkcji
przejścia z $SO(r)$
do $\operatorname{Spin}(r)$.
Wyjaśnijmy zgodność z klasą stopnia dwa z
Twierdzenia~\ref{thm:wi-klasy-16}. Wystarczy ograniczyć
się do $2$-szkieletu bazy: restrykcja $H^2(M;\mathbb F_2)$
do tego szkieletu jest injektywna, ponieważ grupy
kołańcuchów i kobrzegi w stopniach $1,2$ są te same.
Na kompleksie wymiaru dwa wiązka rzędu co najmniej trzy
ma niezerowy przekrój: wybieramy go na wierzchołkach,
przedłużamy po krawędziach przez spójność sfery włókna,
a po $2$-komórkach przez jej jednospójność.
Powtarzając, rozszczepiamy zorientowaną wiązkę na
trywialną część i zorientowaną płaszczyznę $L$.
Wzór sumy daje $w_2(E)=w_2(L)=e(L)\bmod2$.
Wybierzmy zorientowaną ramę $L$ na $1$-szkielecie;
jest to możliwe dzięki spójności $SO(2)$.
Na $2$-komórce, względem tej ustalonej ramy brzegowej,
ostatnia liczba jest parzystością liczby nawinięć
funkcji $SO(2)$, bo to
stopień przeszkody przedłużenia niezerowego przekroju.
Ta sama parzystość jest przeszkodą podniesienia pętli
do $\operatorname{Spin}(2)\to SO(2)$, nakrycia $z\mapsto z^2$.
Po dodaniu kierunków trywialnych obrót o $2\pi$
nadal podnosi się do drogi od $1$ do $-1$ w grupie Spin.
Oba kocykle na $2$-komórkach są więc równe,
co dowodzi zgodności klas. Dla rzędu jeden obie są zerowe.
Z warunku
\eqref{eq:w-normalne-16}
wynika, że jeśli
$\nu(S,M)$ ma obramowanie,
to $w(\nu)=1$ i
$w(TM)|_S=w(TS)$.
To prosty test konieczny
dla niektórych zanurzeń
obramowanych.
\end{uwaga}

\begin{uwaga}[Klasy zespolone i Pontriagina]
Dla zespolonej wiązki
rzędu $r$ analogiczna
konstrukcja z
$\C P^{r-1}$
i klasą pierwszą
tautologicznej linii
o stopniu $2$
daje klasy Cherna
$c_i(E)\in H^{2i}(M;\Z)$.
Dla wiązki rzeczywistej
klasy Pontriagina
definiuje się przez
$p_i(E)=(-1)^i
c_{2i}(E\otimes_\R\C)
\in H^{4i}(M;\Z)$.
Ich naturalność i
zgodność ze stabilnym
dołożeniem wiązki
trywialnej wynikają
z konstrukcji projektowej.
W tym rozdziale
szczegółowo stosujemy
$w_1,w_2$ i klasę
Eulera; pełny rachunek
klas $c_i,p_i$
oraz konstrukcja
Cherna--Weila
wykraczają poza
niezbędne przygotowanie
do chirurgii.
\end{uwaga}

\section{Klasa Eulera, zera pól i samoprzecięcie}
\label{sec:klasy-eulera-16}

\begin{twierdzenie}[Klasa Thoma i klasa Eulera]
\label{thm:euler-thom-16}
Niech $E\to M$ będzie
zorientowaną rzeczywistą
wiązką rzędu $r$ z metryką
we włóknach. Dla jej
wiązki dysków $D(E)$
i wiązki sfer $S(E)$
istnieje klasa
\[
 u_E\in H^r(D(E),S(E);\Z),
\]
której ograniczenie do
każdego włókna
$(D^r,S^{r-1})$ jest
dodatnim generatorem.
Iloczyn z $u_E$ daje
izomorfizm Thoma
\[
 H^j(M;\Z)\stackrel{\cong}{\longrightarrow}
 H^{j+r}(D(E),S(E);\Z).
\]
Jeśli $s_0:M\to D(E)$
jest przekrojem zerowym,
to \emph{klasa Eulera}
$e(E)=s_0^*u_E\in H^r(M;\Z)$
jest naturalna przy
cofnięciu i spełnia
$e(E\oplus F)=e(E)\smile e(F)$
dla zorientowanych wiązek.
Istnienie nigdzie nieznikającego
przekroju pociąga $e(E)=0$.
\end{twierdzenie}

W zapisie $s_0^*u_E$ najpierw zapominamy o względności,
przechodząc z $H^r(D(E),S(E))$ do $H^r(D(E))$,
a potem cofamy klasę przez $s_0$. Równoważnym modelem
klasy Thoma jest $H^r(E,E\setminus0)$: wycięcie
i radialne cofnięcie włókien utożsamiają go z modelem
dyskowym. Ten model pozwala cofać klasę przez przekrój
jako klasę pary $(M,M\setminus Z(s))$.

\begin{proof}
Nad komórką bazy wiązka jest trywialna, zatem odpowiednia
para w przestrzeni całkowitej jest iloczynem tej komórki
z parą $(D^r,S^{r-1})$.
Włóknowa grupa względna
w stopniu $r$ jest $\Z$;
orientacja określa jej
dodatni generator.
Na przecięciu sąsiednich
trywializacji macierz
przejścia ma dodatni
wyznacznik, więc przenosi
ten generator w niego
samego. Sama zgodność lokalnych generatorów nie jest jeszcze
dowodem istnienia globalnej klasy. Skonstruujmy ją
w kompleksie komórkowym pary $(D(E),S(E))$.
Nad każdą $d$-komórką bazy dołączamy względną
komórkę $D^d\times D^r$ wymiaru $d+r$;
część $D^d\times S^{r-1}$ leży już w $S(E)$,
a $\partial D^d\times D^r$ nad wcześniejszym szkieletem.
Orientujemy iloczyn kolejnością „baza, włókno”.
Liczby incydencji są tymi z bazy pomnożonymi
przez stopień transportu orientacji włókna, czyli $+1$.
Kompleks względny kołańcuchów jest więc kompleksem
bazy przesuniętym o $r$. Stały kocykl $1$ w stopniu
zero bazy wyznacza kocykl stopnia $r$ pary,
którego klasa to $u_E$. Jest to jedyna klasa
o wymaganych ograniczeniach: pod tym przesunięciem
odpowiada funkcji na składowych bazy o wartości $1$.
Dla niezwartej bazy używamy produktów współczynników
po komórkach; indukcja po szkieletach ma skończenie
wiele etapów, gdyż rozmaitość ma skończony wymiar.

Teraz sprawdzamy, że mnożenie przez tę już
skonstruowaną klasę daje deklarowany izomorfizm:
na dołączanej $d$-komórce
relatywne grupy pary
wiązki są przesunięciem
o $r$ relatywnych grup
$(D^d,S^{d-1})$.
Mnożenie przez lokalny
generator jest tam
izomorfizmem. W długich
ciągach dokładnych
dla starego szkieletu
i nowej komórki wszystkie
kwadraty są naturalne;
indukcja i lemat o
pięciu odwzorowaniach
dają izomorfizm w
każdym stopniu.
Mnożenie ma dokładnie postać
$a\mapsto\pi^*a\smile u_E$; na każdej komórce
jest opisanym przesunięciem przez włóknowy generator.
Dla rzędu zero przyjmujemy $S(E)=\varnothing$
i $u_E=e(E)=1$.

Cofnięcie wiązki przenosi
orientacje i włóknowe
generatory, więc przenosi
$u_E$ na klasę Thoma
cofniętej wiązki. Po
ograniczeniu do
przekroju zerowego
daje to naturalność
$e(f^*E)=f^*e(E)$.
Dla sumy $E\oplus F$ możemy użyć dysku w normie
maksimum, czyli $D(E)\times_M D(F)$. Jego brzeg to
$(S(E)\times_M D(F))\cup(D(E)\times_M S(F))$.
Cofamy obie klasy Thoma przez rzuty i bierzemy
ich względny iloczyn kubkowy. W każdym włóknie
jest on dodatnim generatorem w orientacji „$E$, potem $F$”,
więc z jedyności jest klasą Thoma sumy.
Zmiana normy przez radialny homeomorfizm nie zmienia
klasy. Cofnięcie do zera daje
$e(E\oplus F)=e(E)\smile e(F)$.

Jeśli $s$ nigdzie
nie znika, po unormowaniu
jest odwzorowaniem
$M\to S(E)$.
Włączenie $S(E)\subset D(E)$
powoduje $s^*u_E=0$,
bo względna klasa
na $S(E)$ znika.
Przekroje $s_0$ i $s$
są homotopijne jako
odwzorowania $M\to D(E)$
przez $t\mapsto ts$.
Ich cofnięcia klasy
$u_E$ są więc równe;
$e(E)=0$.
\end{proof}

\begin{uwaga}[Wersja modulo $2$ i brak błędnego koła]
Powyższy dowód nie używa klas Stiefela--Whitneya.
Z $\mathbb F_2$ działa dla każdej wiązki, ponieważ
zmiana orientacji włókna nie zmienia generatora.
Oznaczmy otrzymaną klasę Eulera modulo $2$ przez $e_2(E)$.
Jest naturalna, multiplikatywna i znika przy istnieniu
niezerowego przekroju, z dokładnie tych samych dowodów.

Dla linii $L$ mamy $e_2(L)=w_1(L)$. Wybierzmy
niezerowe wartości przekroju w wierzchołkach triangulacji
i przedłużmy go ogólnie na krawędzie. Parzystość zer
na krawędzi wynosi $1$ dokładnie wtedy, gdy wartości
na końcach mają przeciwne znaki po transporcie
w trywializacji tej krawędzi. Jest to kocykl zmiany
orientacji, czyli $w_1(L)$; zmiana wartości przy
wierzchołku dodaje kobrzeg. Z drugiej strony te same
parzystości są lokalnymi stopniami cofnięcia klasy Thoma,
więc reprezentują $e_2(L)$.

Te własności $e_2$ wystarczają do dowodu wzoru
w Twierdzeniu~\ref{thm:wi-klasy-16}; nie zakłada on
z góry równości $e_2=w_r$. Dopiero po rozszczepieniu
$E$ na linie otrzymujemy
$e_2(E)=\prod_jw_1(L_j)=w_r(E)$.
Injektywność cofnięcia do wiązki flag przenosi
tę równość na bazę. Dla wiązki zorientowanej
redukcja całkowitej klasy Thoma jest klasą Thoma modulo $2$,
więc również $e(E)\bmod2=w_r(E)$.
\end{uwaga}

\begin{twierdzenie}[Zera przekroju reprezentują klasę Eulera]
\label{thm:euler-zeros-16}
Niech $M^n$ będzie
zamknięta i zorientowana,
a $E\to M$ zorientowana
rzędu $r\leq n$.
Dla przekroju $s$
transwersalnego do zera
zbiór $Z(s)$ jest
zorientowaną
$(n-r)$-podrozmaitością i
\[
 \operatorname{PD}^{-1}[Z(s)]=e(E).
\]
Jeśli $r=n$, zera są
skończone i
\begin{equation}\label{eq:euler-zeros-16}
 \langle e(E),[M]\rangle
   =\sum_{x\in Z(s)}
       \operatorname{sgn}\det(\D_xs).
\end{equation}
W bazach zgodnych z
orientacją znak nie
zależy od wyboru
trywializacji.
\end{twierdzenie}
\begin{proof}
Regularność wynika z twierdzenia o przekrojach transwersalnych.
Pochodna pionowa w zerze daje izomorfizm
$\nu(Z(s),M)\cong E|_{Z(s)}$. Nadajemy normalnej
orientację przeniesioną z $E$, a $Z(s)$ orientujemy
warunkiem „normalna, potem styczna daje orientację $M$”.

W modelu $H^r(E,E\setminus0)$ cofnięcie przez $s$
jest klasą $s^*u_E\in H^r(M,M\setminus Z(s))$.
W małej trywializacji ograniczenie $s$ do normalnego
dysku ma w zerze odwracalną różniczkę o dodatnim
wyznaczniku, więc jego lokalny stopień jest $+1$.
Wycięcie i otoczenie tubularne identyfikują
$s^*u_E$ z klasą Thoma normalnej do $Z(s)$.
Jej obraz w kohomologii bezwzględnej jest klasą
dualną do $Z(s)$: na każdej normalnej komórce
ocenia $+1$, zgodnie z orientacjami w dowodzie dualności.
Liniowa homotopia $s\simeq s_0$ w $E$ pokazuje,
że ten sam obraz jest $e(E)$.

Dla $r=n$ każdy punkt zerowy jest zorientowaną
$0$-komórką o współczynniku
$\operatorname{sgn}\det\D_xs$. Zwartość i izolacja
zer dają ich skończoność. Ewaluacja na klasie
fundamentalnej sumuje te współczynniki,
co daje~\eqref{eq:euler-zeros-16}.
\end{proof}

\begin{wniosek}[Samoprzecięcie]
\label{cor:samoprzeciecie-16}
Niech $M^{2k}$ będzie zorientowana, a
$S^k\subset M$ zamknięta i zorientowana.
W tym wzorze nadajemy wiązce normalnej orientację
przez kolejność \emph{styczna do $S$, potem normalna},
zgodną z orientacją $M$. Wówczas
to liczba przecięcia
$S$ z małą przesuniętą
kopią wynosi
\[
 I(S,S)=\langle e(\nu(S,M)),[S]\rangle.
\]
\end{wniosek}
\begin{proof}
W otoczeniu tubularnym
zapisujemy $S$ jako
przekrój zerowy $\nu$.
Wybierzmy mały przekrój
$s$ transwersalny do
zera i przesuńmy $S$
do jego wykresu.
Wykres przecina
przekrój zerowy
dokładnie dla $s(x)=0$.
W bazach „styczna, normalna” kolumny baz
przekroju zerowego i wykresu tworzą macierz
$\begin{psmallmatrix}I&I\\0&\D_xs\end{psmallmatrix}$.
Jej wyznacznik to $\det\D_xs$, więc znak przecięcia
jest lokalnym stopniem $s$. Ten rachunek ustala
potrzebną orientację normalnej; nie można jej
odwrócić dowolnie, zachowując ten sam zapis wzoru.
Suma znaków to
$\langle e(\nu),[S]\rangle$
z~\eqref{eq:euler-zeros-16}.
\end{proof}

\begin{przyklad}
Dla $S^2\times\{p\}
\subset S^2\times S^2$
wiązka normalna jest
trywialna: jej włókno
to stała przestrzeń
$T_pS^2$. Istnieje
niezerowy stały
przekrój, więc
$I(S,S)=0$.
Dwie różne sfery
$S^2\times\{p\}$
i $\{q\}\times S^2$
przecinają się
jednak w jednym
punkcie z liczbą
przecięcia $1$.
\end{przyklad}

\section{Twierdzenie Poincarégo--Hopfa}
\label{sec:poincare-hopf-16}

\begin{definicja}[Indeks zera pola]
Jeżeli pole $X$ na
$n$-rozmaitości ma
izolowane zero $x$,
wybierzmy małą kulę w mapie wokół $x$
oraz trywializację stycznej przez różniczkę tej mapy.
\emph{Indeks} $\operatorname{ind}_x X$
jest stopniem mapy
$X/|X|:S^{n-1}\to S^{n-1}$, dla $n\geq1$.
Zmiana mapy zmienia jednocześnie orientację dziedziny
i współrzędnych wektora, więc dwa znaki się mnożą
do $+1$. Zmienną macierz przejścia na małej kuli
można zdeformować do jej wartości w środku, nie
przechodząc przez macierze osobliwe. Indeks nie
zależy więc od mapy ani od lokalnej orientacji.
Dla $n=1$ stopień liczymy na $\widetilde H_0(S^0)$;
dla $n=0$ umownie każdy punkt zerowy ma indeks $1$.
Jeśli $\D_xX$ jest
odwracalna,
$\operatorname{ind}_xX
=\operatorname{sgn}\det\D_xX$:
na małej sferze
$X(y)=\D_xX(y-x)+o(|y-x|)$
jest homotopijne
do odwracalnej
mapy liniowej.
\end{definicja}

\begin{twierdzenie}[Poincaré--Hopf]
\label{thm:poincare-hopf-16}
Jeśli $M$ jest
zamkniętą
zorientowaną rozmaitością,
a pole $X$ ma
izolowane zera, to
\begin{equation}\label{eq:poincare-hopf-16}
 \sum_{x:X(x)=0}
     \operatorname{ind}_x X
 =\langle e(TM),[M]\rangle
 =\chi(M).
\end{equation}
W szczególności
$\chi(M)\ne0$
uniemożliwia istnienie
nigdzie nieznikającego
pola stycznego.
\end{twierdzenie}
\begin{proof}
Dla $n=0$ obie strony liczą punkty $M$.
Dalej przyjmujemy $n\geq1$.
\emph{Od indeksów do klasy Eulera.}
Przy zerze nieosobliwym
pierwsza równość jest
dokładnie
Twierdzeniem~\ref{thm:euler-zeros-16}
dla $E=TM$. Dla zera
izolowanego, lecz
osobliwego, wybieramy
kulę bez zer na
brzegu. Mała perturbacja
pola wewnątrz kuli,
niezmieniająca go
przy brzegu, może być
transwersalna do
przekroju zerowego
na mocy twierdzenia o transwersalności.
Stopień mapy
$X/|X|$ na brzegu
nie zmienia się
w tej homotopii;
równa się sumie
znaków nowo
powstałych zer
wewnątrz. Zwartość
$M$ pozwala wykonać
ten zabieg w
skończenie wielu
rozłącznych kulach.

\emph{Dlaczego suma jest $\chi(M)$?}
Przypomnijmy i uzasadnijmy
potrzebną postać
zasady Lefschetza.
Dla gładkiego
$f:M\to M$ o
skończenie wielu
nieosobliwych
punktach stałych
zachodzi
\begin{equation}\label{eq:lefschetz-lokalny-16}
 \sum_{f(x)=x}
    \operatorname{sgn}\det(I-\D_xf)
 =\sum_j(-1)^j
    \operatorname{tr}
      (f_*:H_j(M;\Q)\to H_j(M;\Q)).
\end{equation}
Nadajemy grafowi $\Gamma_f$ i przekątnej $\Delta$
orientacje przeniesione z $M$, a $M\times M$ orientację
produktową. W punkcie stałym, w tej kolejności płatów,
macierz baz stycznych wynosi
\[
 \begin{pmatrix}I&I\\\D_xf&I\end{pmatrix},
 \qquad \det=\det(I-\D_xf).
\]
Lewa strona~\eqref{eq:lefschetz-lokalny-16}
jest więc $I(\Gamma_f,\Delta)$.

Obliczymy tę liczbę bez pomijania znaku globalnego.
Wybierzmy bazę $a_{ij}$ przestrzeni $H^i(M;\Q)$
i bazę dualną $b_{ij}\in H^{n-i}(M;\Q)$, tzn.
$\langle a_{ij}\smile b_{i\ell},[M]\rangle=\delta_{j\ell}$.
Przez $u\times v$ oznaczamy
$\pr_1^*u\smile\pr_2^*v$. Te iloczyny tworzą bazę
kohomologii produktu nad $\Q$. W tym szczególnym
przypadku formułę Künnetha można uzasadnić bez torsji:
komórki produktu mają brzeg
$\partial(e^p\times e^q)=\partial e^p\times e^q+
(-1)^p e^p\times\partial e^q$.
Ich kompleks jest tensorem kompleksów czynników.
Nad ciałem każdy skończony kompleks rozkłada się na
homologię z zerową różniczką i dwuwyrazowe kompleksy
izomorfizmów; tensor z każdym z tych ostatnich
jest ściągalny. Pozostaje tensor homologii,
a po dualizacji podana baza kohomologii.

Klasa dualna przekątnej ma postać
\[
 \alpha_\Delta=\sum_{i,j}(-1)^{ni}\,b_{ij}\times a_{ij}.
\]
Sprawdzamy ją przez ewaluację na bazowych iloczynach:
z definicji dualności
$\langle\alpha_\Delta\smile(u\times v),[M\times M]\rangle
=\langle u\smile v,[M]\rangle$.
Dla $|u|=i$, $|v|=n-i$ znak przestawienia czynników
w iloczynie $(b_{ij}\times a_{ij})\smile(u\times v)$
wynosi $(-1)^{i^2}$, a zamiana kolejności $b_{ij},u$
daje $(-1)^{i(n-i)}$. Razem z $(-1)^{ni}$ jest to
$+1$. Dualność baz daje żądaną ewaluację; w innych
stopniach wkład jest zerowy. Niezdegenerowanie
parowania w produkcie dowodzi wzoru na $\alpha_\Delta$.

Teraz ograniczamy tę klasę do grafu, gdzie
$\pr_1$ jest identycznością, a $\pr_2=f$:
\begin{align*}
 I(\Gamma_f,\Delta)
 &=\langle\alpha_\Delta,[\Gamma_f]\rangle\\
 &=\sum_{i,j}(-1)^{ni}
       \langle b_{ij}\smile f^*a_{ij},[M]\rangle\\
 &=\sum_i(-1)^i\operatorname{tr}
       (f^*:H^i(M;\Q)\to H^i(M;\Q)).
\end{align*}
Ostatni znak wynika z
$ni+i(n-i)\equiv i\pmod2$. Sumowanie po $j$
wybiera elementy diagonalne macierzy $f^*$.
Ewaluacja identyfikuje $f^*$ z dualnym odwzorowaniem
do $f_*$, więc ich ślady są równe. Otrzymujemy
\eqref{eq:lefschetz-lokalny-16}.

Teraz przyjmijmy, że
zera $X$ są nieosobliwe,
i niech $f=\Phi_{-t}$
będzie przepływem
$X$ przez bardzo
mały czas $-t<0$.
Punkty zerowe są
punktami stałymi.
Poza małymi kulami
otaczającymi zera
norma $X$ ma
dodatnie minimum,
więc z rozwinięcia
$\Phi_{-t}(y)=y-tX(y)+O(t^2)$
jednolitego na tym
zbiorze nie ma
innych punktów
stałych dla małego $t$.
W mapie przy zerze wyrażenie
$F(y,t)=(y-\Phi_{-t}(y))/t$ przedłuża się gładko
do $t=0$ przez $F(y,0)=X(y)$; wynika to ze wzoru
całkowego na przyrost przepływu. Ponieważ
$\D_yF(x,0)=\D_xX$ jest odwracalna, twierdzenie
o funkcji uwikłanej z parametrem $t$ daje jedyne
rozwiązanie w jednym ustalonym otoczeniu $x$.
Jest nim stale $x$. Zmniejszamy kule do tych
otoczeń, a ich zwarte dopełnienie obejmuje także
pozostałe pierścienie wokół zer. Poprzednie
jednolite oszacowanie wyklucza tam punkty stałe.
W punkcie zerowym
\[
 I-\D_x\Phi_{-t}
   =t\D_xX+O(t^2),
\]
zatem znak wyznacznika
to $\operatorname{ind}_xX$.
Przepływ jest
homotopijny do
identyczności,
więc $f_*$ jest
identycznością
na wszystkich
$H_j(M;\Q)$.
Prawa strona
\eqref{eq:lefschetz-lokalny-16}
wynosi
$\sum_j(-1)^j b_j=\chi(M)$.
Otrzymujemy drugą
równość~\eqref{eq:poincare-hopf-16}.
Perturbacja z początku
dowodu rozszerza ją
na zera osobliwe.
\end{proof}

\begin{figure}[htbp]
\centering
\begin{tikzpicture}[>=Latex,font=\small,line cap=round,scale=1.05]
 \shade[inner color=white,outer color=manifoldblue!43] (0,0) circle (1.55);
 \foreach \c in {-.88,-.48,0,.48,.88}
  \draw[manifoldblue!55,thin]
   plot[domain=8:172,samples=65,variable=\t]
    ({1.55*\c*sin(\t)},{1.55*cos(\t)});
 \foreach \c in {-.88,-.48,0,.48,.88}
  \draw[->,manifoldgold!85!black,thick]
   plot[domain=115:62,samples=25,variable=\t]
    ({1.55*\c*sin(\t)},{1.55*cos(\t)});
 \fill[manifoldred] (0,1.55) circle (2.7pt);
 \fill[manifoldred] (0,-1.55) circle (2.7pt);
 \node[anchor=west,manifoldred] at (.16,1.55) {$+1$};
 \node[anchor=west,manifoldred] at (.16,-1.55) {$+1$};
 \node[anchor=west] at (1.9,.30) {$h(x,y,z)=z$};
 \node[anchor=west] at (1.9,-.12) {$X=e_z-z(x,y,z)$};
\end{tikzpicture}
\caption{Gradient wysokości na $S^2$
ma dwa zera: na biegunie dolnym
i górnym. Oba mają indeks $+1$,
więc suma indeksów i $\chi(S^2)$ są równe $2$.
Złote łuki pokazują rzuty trajektorii gradientu
w metryce standardowej: biegną po południkach od minimum
do maksimum. Promień rysunku jest tylko skalą graficzną.}
\label{fig:poincare-hopf-16}
\end{figure}

\begin{przyklad}[Sfera i torus]
Na $S^2$ gradient
funkcji wysokości
ma minimum i maksimum.
Lokalne macierze
pochodnej pola są
odpowiednio dodatnio
i ujemnie określone
w dwóch wymiarach,
ale oba wyznaczniki
są dodatnie:
$1+1=\chi(S^2)=2$.
Na torusie $T^2$
istnieje nigdzie
nieznikające pole
styczne wzdłuż
jednego czynnika
$S^1$, więc
$e(TT^2)=0$ i
$\chi(T^2)=0$.
To dopowiada topologiczną
treść wzoru
Gaussa--Bonneta z twierdzenia Gaussa--Bonneta.
\end{przyklad}

\section{Forma przecięcia i przygotowanie do chirurgii}
\label{sec:forma-przeciecia-16}

Gdy $n=2k$, dwie
$k$-podrozmaitości
przecinają się
w punktach.
Ich liczba przecięcia
tworzy na środkowej
homologii formę
dwuliniową. Część
torsyjna nie jest
wykrywana przez
wartości całkowite
tej formy, więc
najpierw ją usuwamy.

\begin{twierdzenie}[Forma przecięcia]
\label{thm:forma-przeciecia-16}
Dla zamkniętej
zorientowanej
$M^{2k}$ parowanie
\[
 Q_M:
 \frac{H_k(M;\Z)}{\mathrm{Tor}}
 \times
 \frac{H_k(M;\Z)}{\mathrm{Tor}}
   \longrightarrow\Z,
 \qquad
 Q_M(a,b)=
 \langle\operatorname{PD}^{-1}(a)
 \smile\operatorname{PD}^{-1}(b),[M]\rangle
\]
jest \emph{unimodularne}, czyli doskonałe nad $\Z$:
macierz w dowolnej bazie
części wolnej ma wyznacznik
$\pm1$. Ponadto
$Q_M(a,b)=(-1)^kQ_M(b,a)$.
W wymiarze $4m$ forma
jest symetryczna, a
w wymiarze $4m+2$
antysymetryczna.
\end{twierdzenie}
\begin{proof}
Z dualności Poincarégo
otrzymujemy
$H_k(M;\Z)\cong H^k(M;\Z)$.
Twierdzenie o
współczynnikach
uniwersalnych daje
ciąg dokładny
\[
0\longrightarrow
 \operatorname{Ext}
     (H_{k-1}(M),\Z)
 \longrightarrow H^k(M;\Z)
 \longrightarrow
 \operatorname{Hom}(H_k(M),\Z)
 \longrightarrow0.
\]
Grupy homologii
zwartej triangulacji
są skończenie
generowane.
Składnik $\operatorname{Ext}$
jest więc torsyjny,
a grupa $\operatorname{Hom}$
jest dualną kratą
części wolnej $H_k$.
Po podzieleniu przez
torsję otrzymujemy
izomorfizm
\[
 H_k(M)/\mathrm{Tor}
 \ \cong\
 \operatorname{Hom}
    (H_k(M)/\mathrm{Tor},\Z).
\]
Dla wybranej konwencji kapowej ten izomorfizm
wysyła $a$ na funkcjonał
\[
 b\longmapsto\langle\operatorname{PD}^{-1}(a),b\rangle
 =Q_M(b,a)=(-1)^kQ_M(a,b).
\]
Wspólny czynnik $(-1)^k$ nie zmienia odwracalności.
Klasy torsyjne nie wpływają na $Q_M$, ponieważ
ich całkowita ewaluacja jest torsyjną liczbą całkowitą,
a więc zerem. Definicja na ilorazie jest zatem jednoznaczna.
Macierz odwzorowania
między wolną grupą
abelową a jej dualną
jest odwracalna nad
$\Z$, stąd jej
wyznacznik to $\pm1$.
Dla wszystkich klas, także niereprezentowanych
przez zanurzone podrozmaitości, przemienność
stopniowana iloczynu kubkowego daje znak
$(-1)^{k^2}=(-1)^k$. To dowodzi obu symetrii.
Nie zakładamy, że każda całkowita klasa homologii
ma reprezentanta będącego podrozmaitością.
\end{proof}

\begin{przyklad}
Na $S^2\times S^2$
w bazie
$[S^2\times\{p\}]$,
$[\{q\}\times S^2]$
forma ma macierz
\[
 Q=
 \begin{pmatrix}
  0&1\\1&0
 \end{pmatrix},
 \qquad\det Q=-1.
\]
Na torusie
$H_1(T^2)=\Z^2$
i w bazie dwóch
podstawowych okręgów
otrzymujemy
$\begin{psmallmatrix}
0&1\\-1&0
\end{psmallmatrix}$.
Druga forma jest
antysymetryczna,
bo oba cykle mają
wymiar nieparzysty.
\end{przyklad}

\begin{definicja}[Elementarna operacja chirurgiczna]
\label{def:chirurgia-zarys-16}
Niech $0\leq k\leq n-1$
i niech dane będzie
\emph{obramowane}
zanurzenie
$S^k\times D^{\,n-k}
\hookrightarrow M^n$.
Elementarna
\emph{chirurgia indeksu $k$} (w przyjętej tutaj
konwencji: wymiar operowanej sfery; uchwyt w jej
śladzie ma indeks $k+1$) polega na usunięciu
wnętrza tego
otoczenia i
wklejeniu
$D^{k+1}\times
S^{\,n-k-1}$
wzdłuż wspólnego brzegu
\[
 \partial(S^k\times D^{n-k})
 =S^k\times S^{n-k-1}
 =\partial(D^{k+1}\times S^{n-k-1}).
\]
Dla ustalonego obramowanego zanurzenia sklejamy
punkty o tych samych współrzędnych produktowych na
$S^k\times S^{n-k-1}$; jest to określony dyfeomorfizm
brzegów. Jego dowolna dodatkowa zmiana jest dodatkową
daną i może zmienić wynik. Kołnierze określają gładką
strukturę przy szwie. Jeśli chcemy zachować orientację,
identyfikacja musi odwracać orientacje indukowane
na brzegach obu sklejanych części.
\end{definicja}

To wyjaśnia rolę
wcześniejszych tematów.
Twierdzenie o otoczeniu
tubularnym zapewnia
obszar, który można
wyciąć. Obramowanie
z definicji obramowania wiązki normalnej
pozwala zapisać go
jako \emph{iloczyn}
$S^k\times D^{n-k}$.
Klasy charakterystyczne
mogą wykazać, że
normalna wiązka
nie ma takiego
obramowania.
Forma przecięcia
kontroluje, jak
cykle środkowego
wymiaru są ze sobą
powiązane; po
operacji trzeba
obliczyć ją na nowo.
Szczegółowa teoria
sklejania, uchwytów
i zmian homologii
będzie osobnym
rozdziałem.

\begin{przyklad}[Cięcie sfery w dwóch punktach]
Na $S^2$ wybierzmy
dwa rozłączne
małe dyski:
$S^0\times D^2$
to właśnie ich suma.
Po wycięciu wnętrz
pozostaje walec
$S^1\times[0,1]$.
W chirurgii
$k=0,n=2$
dołączamy
$D^1\times S^1$,
drugi walec,
wzdłuż dwóch
brzegowych okręgów.
Wybieramy sklejenia zgodne z orientacją powstałej
powierzchni, czyli odwracające orientacje jej brzegów.
Wtedy dwa walce składają się w zorientowaną zamkniętą
powierzchnię rodzaju jeden, czyli torus. Bez tego
warunku sklejenie może dać butelkę Kleina:
po obejściu podłużnego okręgu kierunek poprzeczny
wraca wówczas z przeciwną orientacją. Sama
charakterystyka Eulera nie rozróżnia tych wyników.
Na poziomie
charakterystyki
Eulera: ze sfery
o $\chi=2$ usunięcie
dwóch dysków
odejmuje $2$,
a dołączenie walca nie zmienia wyniku:
$\chi(A\cup B)=\chi(A)+\chi(B)-\chi(A\cap B)$,
przy czym zarówno walec, jak i dwa okręgi sklejenia
mają charakterystykę $0$.
Torus ma więc
$\chi=0$, zgodnie
z bezpośrednim
obliczeniem.
\end{przyklad}

\par\smallskip
{\footnotesize\noindent\emph{Dalsza lektura:}
\href{https://pi.math.cornell.edu/~hatcher/AT/AT.pdf}{A.\ Hatcher, \emph{Algebraic Topology}, dualność i iloczyn kubkowy},
\href{https://pi.math.cornell.edu/~hatcher/VBKT/VB.pdf}{A.\ Hatcher, \emph{Vector Bundles and K-Theory}, sklejanie i klasy charakterystyczne},
\href{https://math.stanford.edu/~ralph/bookR4.pdf}{R.\ Cohen, \emph{Bundles, Manifolds, and Homotopy}, przecięcia i klasa Thoma}.\par}


% Iteracja V: ciagi spektralne i ciag Serre'a.
\chapter{Ciągi spektralne: od filtracji do fibracji Serre'a}
\label{ch:ciagi-spektralne-v5}

\section{Motywacja: dlaczego jedna długa sekwencja dokładna nie wystarcza?}
Dla fibracji $F\to E\xrightarrow{p}B$ długi ciąg grup homotopii
porównuje $\pi_*(F)$, $\pi_*(E)$ i $\pi_*(B)$, ale nawet
znajomość wszystkich grup $H_*(F)$ i $H_*(B)$ nie pozwala po prostu
pomnożyć ich i otrzymać $H_*(E)$. W fibracji Hopfa
$S^1\to S^3\to S^2$ naiwne mnożenie homologii dawałoby
homologię $S^1\times S^2$, a nie $S^3$. Musi istnieć
mechanizm, który \emph{usuwa} niektóre klasy przy sklejaniu.

Takim mechanizmem jest ciąg spektralny: na każdej \emph{stronie}
$E^r$ widzimy grupy podobne do kolejnych przybliżeń homologii,
a różniczka $d^r$ mierzy obstrukcję przeniesienia klasy do
następnego przybliżenia. Słowo ,,spektralny'' ma tu tradycyjne
znaczenie i nie oznacza widma operatora.

Najpierw wyprowadzimy konstrukcję z filtracji \emph{dowolnego}
kompleksu łańcuchowego. Dopiero potem zastosujemy ją do
filtracji przestrzeni całkowitej wiązki przez przeciwobrazy
szkieletów podstawy. Wszystkie homologie do końca tego rozdziału
mają współczynniki całkowite, chyba że zaznaczono inaczej.

\section{Filtracja kompleksu i pierwsza strona}
\begin{definicja}[Filtracja rosnąca kompleksu]
Niech $(C_*,\partial)$ będzie kompleksem łańcuchowym grup abelowych
w stopniach nieujemnych. \emph{Filtracją rosnącą} nazywamy
podkompleksy
$\cdots\subseteq F_{p-1}C_*\subseteq F_pC_*\subseteq\cdots$,
takie że $\partial(F_pC_n)\subseteq F_pC_{n-1}$.
Filtracja jest \emph{ograniczona w każdym stopniu}, jeżeli
dla każdego $n$ istnieją $a_n\le b_n$ z
$F_{a_n-1}C_n=0$ oraz $F_{b_n}C_n=C_n$.
Dla $p<0$ w przykładach pierwszoćwiartkowych przyjmujemy $F_pC=0$.
\end{definicja}
\begin{przyklad}[Filtracja przez szkielety]
Jeżeli $X$ jest CW-kompleksem, to jego szkielety
$X^{-1}=\varnothing\subset X^0\subset X^1\subset\cdots$
dają filtrację komórkowego kompleksu $C_*^{\mathrm{CW}}(X)$.
W każdej grupie łańcuchów granica komórki $n$-wymiarowej
leży w szkielecie $n-1$.
\end{przyklad}
\begin{definicja}[Strona zerowa i pierwsza]
Dla filtrowanego kompleksu $C$ definiujemy
\[
E^0_{p,q}=F_pC_{p+q}/F_{p-1}C_{p+q},\qquad
 d^0:E^0_{p,q}\to E^0_{p,q-1},\quad [c]\mapsto[\partial c].
\]
Jest to dobrze określone, a $(d^0)^2=0$.
\emph{Pierwsza strona} to $E^1_{p,q}=H_{p+q}(F_pC/F_{p-1}C)$.
Grupa $E^1_{p,q}$ jest grupą homologii ilorazu kompleksów,
a nie ilorazem $H(F_pC)/H(F_{p-1}C)$.
\end{definicja}
\begin{przyklad}[Elementarna filtracja]
Niech $C_2=\Z x$, $C_1=\Z y$, $\partial x=y$,
pozostałe grupy są zerowe, oraz $F_0C=\Z y$,
$F_1C=F_0C$, $F_2C=C$. Zwykła homologia $C$ znika,
ale pierwsza strona ma $E^1_{2,0}=\Z$ i $E^1_{0,1}=\Z$.
Różniczka $d^1$ jest zerowa ze względów stopniowania,
więc oba elementy przeżywają do strony drugiej.
Dopiero $d^2:E^2_{2,0}\to E^2_{0,1}$ jest izomorfizmem.
To najprostszy przykład klasy, która ,,ginie dopiero na drugiej stronie''.
\end{przyklad}

\section{Konstrukcja wszystkich różniczek i dowód strony następnej}
\begin{definicja}[Cykle i brzegi rzędu $r$]
Niech $r\ge1$ i $n=p+q$. Określamy
\[
Z^r_{p,q}=\{x\in F_pC_n:\partial x\in F_{p-r}C_{n-1}\},
\quad
B^r_{p,q}=Z^r_{p,q}\cap\bigl(F_{p-1}C_n+\partial(F_{p+r-1}C_{n+1})\bigr).
\]
Kładziemy $E^r_{p,q}=Z^r_{p,q}/B^r_{p,q}$.
Dla $r=1$ identyfikujemy ten iloraz z poprzednio zdefiniowaną
pierwszą stroną.
\end{definicja}
\begin{lemat}[Poprawność różniczki]
\label{lem:dr-v5}
Dla $r\ge1$ wzór $d^r[x]=[\partial x]$ wyznacza
homomorfizm
$d^r:E^r_{p,q}\to E^r_{p-r,q+r-1}$,
a $d^r\circ d^r=0$.
\end{lemat}
\begin{proof}
Niech $x\in Z^r_{p,q}$. Wtedy $\partial x\in F_{p-r}C_{n-1}$,
a $\partial^2x=0$, zatem $\partial x$ jest cyklem
rzędu $r$ w miejscu $(p-r,q+r-1)$. Załóżmy, że zmienimy
reprezentanta na $x'=x+y+\partial z$, gdzie
$y\in F_{p-1}C_n$, $z\in F_{p+r-1}C_{n+1}$,
a $x'\in Z^r_{p,q}$. Wówczas
$\partial x'-\partial x=\partial y$,
co jest brzegiem rzędu $r$ w miejscu docelowym, bo
$y\in F_{(p-r)+r-1}C_n=F_{p-1}C_n$.
Własność $(d^r)^2=0$ pochodzi z $\partial^2=0$.
\end{proof}
\begin{twierdzenie}[Strony sukcesywne]
\label{thm:pages-v5}
Dla każdego $r\ge1$ istnieje naturalny izomorfizm
$E^{r+1}_{p,q}\cong H_{p,q}(E^r,d^r)$, czyli
$E^{r+1}_{p,q}\cong\ker(d^r:E^r_{p,q}\to E^r_{p-r,q+r-1})/
\operatorname{im}(d^r:E^r_{p+r,q-r+1}\to E^r_{p,q})$.
\end{twierdzenie}
\begin{proof}
Wyjaśnijmy osobno przechodzenie do jądra i następnie
ilorazowanie przez obraz. Warunek $d^r[x]=0$ mówi, że
$\partial x$ można zapisać jako
$\partial x=y+\partial z$, gdzie
$y\in F_{p-r-1}C_{n-1}$ i $z\in F_{p-1}C_n$.
Zastąpmy $x$ przez $x-z$ (wolno, gdyż $z\in F_{p-1}$).
Wtedy $\partial(x-z)=y\in F_{p-r-1}C_{n-1}$,
więc $x-z\in Z^{r+1}_{p,q}$. Odwrotnie, cykl rzędu $r+1$
jest cyklem rzędu $r$ zabijanym przez $d^r$.

Ustalmy teraz dwie reprezentacje $x,x'\in Z^{r+1}_{p,q}$.
Różnica ich klas na stronie $r$ jest obrazem $d^r$
wtedy i tylko wtedy, gdy modulo $F_{p-1}C_n$ jest
brzegiem elementu z $F_{p+r}C_{n+1}$, którego brzeg
leży w $F_pC_n$; ostatni warunek jest dokładnie
warunkiem bycia cyklem rzędu $r$ w miejscu $(p+r,q-r+1)$.
Otrzymujemy więc relację
$x-x'\in F_{p-1}C_n+\partial(F_{p+r}C_{n+1})$,
z uwzględnieniem warunku bycia w $Z^{r+1}$.
Jest to definicja mianownika $B^{r+1}_{p,q}$.
Opisane przejście jest zgodne z dodawaniem i daje izomorfizm.
\end{proof}
\begin{uwaga}[Indeksy nie są dekoracją]
Różniczka w \emph{homologicznym} ciągu spektralnym ma
bistopień $(-r,r-1)$. Dla \emph{kohomologicznego} ciągu
spektralnego przy standardowej konwencji ma bistopień
$(r,1-r)$. W dalszym rachunku każdorazowo zaznaczamy,
którą konwencję stosujemy.
\end{uwaga}

\section{Zbieżność i problem rozszerzeń}
\begin{definicja}
Filtracja kompleksu indukuje filtrację homologii przez
$F_pH_n(C)=\operatorname{im}(H_n(F_pC)\to H_n(C))$.
Piszemy $\operatorname{gr}_pH_n(C)=F_pH_n(C)/F_{p-1}H_n(C)$.
\end{definicja}
\begin{twierdzenie}[Zbieżność przy filtracji ograniczonej]
\label{thm:convergence-v5}
Niech $F_pC$ będzie filtracją ograniczoną w każdym stopniu.
Wtedy dla każdego $(p,q)$ grupy $E^r_{p,q}$ stabilizują się,
a ich wartość stabilna $E^\infty_{p,q}$ jest naturalnie
izomorficzna z $\operatorname{gr}_pH_{p+q}(C)$.
\end{twierdzenie}
\begin{proof}
Dla ustalonego $n=p+q$, gdy $r$ jest dostatecznie duże,
$F_{p-r}C_{n-1}=0$ oraz $F_{p+r-1}C_{n+1}=C_{n+1}$.
Dlatego licznik w definicji $E^r$ staje się
$Z_n(C)\cap F_pC_n$, zaś mianownik
$\bigl(Z_n(C)\cap F_pC_n\bigr)\cap
(F_{p-1}C_n+\partial C_{n+1})$.
Przejście od cyklu $z$ do $[z]\in H_n(C)$
ma jądro równe cyklom będącym granicami w $C$.
Po odjęciu cyklu z $F_{p-1}C_n$ opisany iloraz jest
więc dokładnie obrazem homologii filtru $F_p$ modulo
obraz z poprzedniego filtru, czyli
$F_pH_n(C)/F_{p-1}H_n(C)$. Ponieważ filtry w danym
stopniu są ograniczone, nie występuje problem nieskończonych
kroków ani dopełniania granic.
\end{proof}
\begin{przyklad}[Nie wolno sumować $E^\infty$ bez uzasadnienia]
Jeżeli w pewnym stopniu filtracja ma dwa kolejne ilorazy
$E^\infty_{0,n}\cong\Z/2$ i
$E^\infty_{1,n-1}\cong\Z/2$, to grupa $H_n$ może być
izomorficzna zarówno z $\Z/4$, jak i z
$\Z/2\oplus\Z/2$: obie mają filtrację z tymi ilorazami.
Dane $E^\infty$ same nie rozwiązują \emph{problemu rozszerzeń}.
Jeśli natomiast wszystkie występujące ilorazy są grupami wolnymi,
to kolejne krótkie ciągi dokładne rozszczepiają się jako ciągi
grup abelowych, choć zwykle niekanonicznie.
\end{przyklad}
\begin{uwaga}
Przy filtracjach nieograniczonych trzeba osobno sformułować
warunki zbieżności i kompletności. Nie używamy
Twierdzenia~\ref{thm:convergence-v5} bez tych hipotez.
\end{uwaga}

\section{Dokładna para: skąd pochodzą strony z punktu widzenia par}
\begin{definicja}[Dokładna para]
Dokładną parą grup abelowych nazywamy układ
$D\xrightarrow{i}D\xrightarrow{j}E\xrightarrow{k}D$
z $\operatorname{im}i=\ker j$,
$\operatorname{im}j=\ker k$,
$\operatorname{im}k=\ker i$.
Wtedy $d=j\circ k:E\to E$ spełnia $d^2=0$.
\end{definicja}
\begin{proof}
Ponieważ $kj=0$ z dokładności w środkowym wierzchołku,
$d^2=j(kj)k=0$.
\end{proof}
\begin{stwierdzenie}
Filtracja kompleksu $F_pC$ wyznacza dokładną parę
$D^1_{p,q}=H_{p+q}(F_pC)$ oraz
$E^1_{p,q}=H_{p+q}(F_pC,F_{p-1}C)$.
Mapy $i,j,k$ pochodzą odpowiednio z włączenia filtru,
przejścia do ilorazu i homomorfizmu łączącego
w długim ciągu dokładnym pary kompleksów.
Przy tej konstrukcji $d^1=j\circ k$, a kolejne pary
pochodne dają strony $E^r$ z poprzedniej sekcji.
\end{stwierdzenie}
\begin{proof}
Dla każdego $p$ stosujemy krótki ciąg dokładny kompleksów
$0\to F_{p-1}C\to F_pC\to F_pC/F_{p-1}C\to0$.
Z dokładności jego długiego ciągu wynika kolejno
$\operatorname{im}i=\ker j$, $\operatorname{im}j=\ker k$
i $\operatorname{im}k=\ker i$, gdy uwzględnimy przesunięcie
indeksu $p\mapsto p-1$ mapy $k$.
Na poziomie cyklu $x\in F_pC$ z $\partial x\in F_{p-1}C$
homomorfizm łączący wysyła $[x]$ w $[\partial x]$;
a następnie $j$ bierze jego klasę modulo $F_{p-2}C$.
To jest dokładnie definicja $d^1$.
\end{proof}

\section{Fibracja nad kompleksem CW i lokalne współczynniki}
Długi ciąg homotopii nie wykrywa torsji i iloczynów w $H_*(E)$.
Zastosujemy właśnie zbudowany mechanizm do fibracji.
Najpierw podajemy szczegółowy dowód dla \emph{lokalnie trywialnych}
wiązek nad skończeniewymiarowymi CW-kompleksami. Ich rzutowania
są fibracjami Hurewicza w zwykłych sytuacjach CW. Istnieje też
pełne twierdzenie dla fibracji Serre'a; wymaga jednak osobnego
argumentu porównującego homologię względną preobrazów komórek
z homologią ich włókien i nie utożsamiamy go z lokalną trywialnością.

\begin{lemat}[Fibracja nad dyskiem]
\label{lem:serre-disk-v5}
Niech $p:E\to D^m$ będzie fibracją Hurewicza, $b_0\in D^m$
i $F=p^{-1}(b_0)$. Wówczas $p$ jest homotopijnie
trywialna \emph{nad podstawą}: istnieją mapy
$E\rightleftarrows D^m\times F$ nad $D^m$,
których złożenia są homotopijne (nad $D^m$) do identyczności.
\end{lemat}
\begin{proof}
Wybierzmy kontrakcję podstawy $h_t:D^m\to D^m$,
$h_0=\id$, $h_1\equiv b_0$. Podnosimy homotopię
$H(e,t)=h_t(p(e))$ z podniesienia początkowego $\id_E$.
Otrzymujemy $L:E\times I\to E$ z $pL(e,t)=h_t(p(e))$
i $L(e,0)=e$. W szczególności $r(e)=L(e,1)$ leży w $F$.
Mapa $g:E\to D^m\times F$ dana przez $e\mapsto(p(e),r(e))$
jest więc ciągła. Własność podnoszenia dla
$Z=D^m\times F$ zastosowana do homotopii
$K((b,x),t)=h_{1-t}(b)$ i mapy początkowej
$(b,x)\mapsto x\in F\subset E$ dostarcza mapy
$f:D^m\times F\to E$ nad $D^m$ przy $t=1$.
Złączenie podstawowych kontrakcji i ich homotopii
względem końców daje homotopie włókniste
$fg\simeq\id_E$ i $gf\simeq\id_{D^m\times F}$:
wymagana homotopia podniesień otrzymywana jest przez
własność Hurewicza zastosowaną do rodziny parametryzowanej
przez $E\times I$ i odpowiednio $F\times D^m\times I$.
Można również bezpośrednio sprawdzić ten wniosek przez
podnoszenie liniowej homotopii między odpowiednimi
złożeniami dróg w wypukłym $D^m$. Ważne, że homotopie
pozostają nad podstawą, a nie tylko w przestrzeni całkowitej.
\end{proof}

\begin{definicja}[Monodromia na homologii włókna]
Niech $p:E\to B$ będzie fibracją Hurewicza o łukowo spójnej
podstawie, $b_0\in B$, $F=p^{-1}(b_0)$.
Transport wzdłuż drogi $\gamma:b_0\to b_1$ daje określoną
z dokładnością do homotopii równoważność
$T_\gamma:F_{b_0}\to F_{b_1}$.
W każdym stopniu $q$ otrzymujemy system lokalny
$\mathcal H_q(F)$ o włóknach $H_q(F_b;\Z)$ i transporcie
$(T_\gamma)_*$. Przy zamkniętej pętli działa
$\pi_1(B,b_0)$ na $H_q(F;\Z)$.
\end{definicja}

\begin{lemat}[Względna homologia preobrazów szkieletów]
\label{lem:serre-cells-v5}
Niech $p:E\to B$ będzie lokalnie trywialną wiązką nad
CW-kompleksem $B$, a $E_j=p^{-1}(B^j)$.
Dla $j\ge0$ i $k\in\Z$ istnieje izomorfizm
\[
H_k(E_j,E_{j-1};\Z)\cong
\bigoplus_{\alpha\in\operatorname{Cells}_j(B)}H_{k-j}(F_\alpha;\Z),
\]
gdzie $F_\alpha$ jest włóknem nad wnętrzem komórki
$e^j_\alpha$, a $H_l(F_\alpha)=0$ dla $l<0$.
Znak identyfikacji zależy od wyboru orientacji komórki.
\end{lemat}
\begin{proof}
CW-kompleks $B^j$ jest otrzymany z $B^{j-1}$ przez
dołączenie dysków charakterystycznych
$\chi_\alpha:D^j_\alpha\to B^j$. Ich otwarte wnętrza
są parami rozłączne. Cofnięcie lokalnie trywialnej wiązki
na dysk $D^j_\alpha$ jest trywialne (wiązka nad
parazwartą przestrzenią ściągalną jest trywialna), więc
nad nim mamy parę
$(D^j\times F_\alpha,S^{j-1}\times F_\alpha)$.
Jej brzeg ma kołnierz: para jest kofibracyjna.

Relatywne wycinanie dla komórkowego przylegania daje
izomorfizm z sumą bezpośrednią względnych homologii
kolejnych par dyskowych. Dla skończonej rodziny komórek
jest to zwykłe wycinanie po pogrubieniu brzegu kołnierzem;
w przypadku nieskończenie wielu komórek stosujemy fakt,
że każdy singularny łańcuch jest skończony, a jego obraz
zwarty, więc obliczenie przechodzi do granicy prostej
po skończonych podkompleksach. Stąd otrzymujemy sumę
prostą (a nie iloczyn).

Grupy $H_l(D^j,S^{j-1};\Z)$ wynoszą $\Z$ dla $l=j$
i zero dla pozostałych $l$. Z relatywnego wzoru
Künnetha wynika zatem
$H_k(D^j\times F,S^{j-1}\times F;\Z)\cong H_{k-j}(F)$;
człon $\operatorname{Tor}$ znika z powodu wolności
homologii pary $(D^j,S^{j-1})$.
Sumując po komórkach, otrzymujemy tezę.
\end{proof}

\begin{twierdzenie}[Homologiczny ciąg spektralny Serre'a]
\label{thm:serre-v5}
Niech $p:E\to B$ będzie lokalnie trywialną wiązką o typie
CW-kompleksu jako włóknie, $B$ spójnym i skończeniewymiarowym
CW-kompleksem, $b_0\in B$, $F=p^{-1}(b_0)$.
Istnieje naturalny pierwszoćwiartkowy ciąg spektralny
\[
E^2_{p,q}\cong H_p(B;\mathcal H_q(F)),\qquad
 d^r:E^r_{p,q}\longrightarrow E^r_{p-r,q+r-1},
\]
zbieżny do $H_{p+q}(E;\Z)$ w sensie skończonej w każdym
stopniu filtracji $F_pH_{p+q}(E)$ z
$E^\infty_{p,q}\cong F_pH_{p+q}(E)/F_{p-1}H_{p+q}(E)$.
Jeżeli monodromia na $H_q(F;\Z)$ jest trywialna, można
zastąpić $\mathcal H_q(F)$ stałym systemem współczynników.
Dla fibracji Serre'a o odpowiednim typie CW analogiczna
konkluzja pozostaje prawdziwa; nie traktujemy jej jednak
jako udowodnionej w pełnej ogólności poniższym argumentem.
\end{twierdzenie}
\begin{proof}
\emph{Krok 1: filtrujemy przestrzeń całkowitą.}
Niech $B^p$ oznacza szkielet $p$-wymiarowy $B$, przyjmijmy
$B^{-1}=\varnothing$ i $E_p=p^{-1}(B^p)$.
Filtrujemy singularny kompleks $C_*(E)$ podkompleksami
$F_pC_*(E)=\operatorname{im}(C_*(E_p)\to C_*(E))$.
Każda singularna mapa sympleksu do $E$ ma zwarty obraz,
a jej rzut do CW-kompleksu jest zawarty w skończonym
podkompleksie $B$. Jeżeli $m=\dim B<\infty$, to
$E_m=E$, więc filtracja $C_*(E)$ jest rzeczywiście
ograniczona w każdym stopniu: $F_{-1}C_*=0$ i
$F_mC_*=C_*$. Dzięki temu wolno zastosować
Twierdzenie~\ref{thm:convergence-v5} bez ukrytych
założeń o zupełności filtracji.

\emph{Krok 2: obliczamy wyrazy strony pierwszej.}
Mamy $E^1_{p,q}=H_{p+q}(E_p,E_{p-1};\Z)$.
Wybierzmy dla każdej zorientowanej $p$-komórki $e^p_\alpha$
mapę charakterystyczną
$\chi_\alpha:(D^p,S^{p-1})\to(B^p,B^{p-1})$.
Lemat~\ref{lem:serre-cells-v5} i lokalne trywializacje nad komórkami dają
\[
H_{p+q}(E_p,E_{p-1};\Z)\cong
\bigoplus_{\alpha\in\operatorname{Cells}_p(B)}
H_{p+q}(D^p\times F,S^{p-1}\times F;\Z).
\]
Ponieważ $H_k(D^p,S^{p-1};\Z)$ jest $\Z$ dla $k=p$
i zerem poza tym stopniem, względna postać
Künnetha (bez członu $\operatorname{Tor}$, gdyż odpowiednia
grupa jest wolna) identyfikuje każdą summandę z
$H_q(F;\Z)$. Zmiana trywializacji działa na tej grupie
przez monodromię. Dla $q<0$ grupy są zerowe.

\emph{Krok 3: identyfikujemy $d^1$.}
Homomorfizm łączący pary $(E_p,E_{p-1})$ bierze granicę
względnego cyklu nad komórką i ogranicza ją do brzegu komórki.
Dla dwóch sąsiednich komórek suma współczynników
zawiera znak liczby incydencji ich orientacji, a
współczynnik z włókna trzeba przenieść wzdłuż drogi
łączącej wybrane punkty bazowe. Jest to dokładnie
różniczka komórkowych łańcuchów z lokalnym systemem
$\mathcal H_q(F)$, którego konstrukcję podaliśmy w
Sekcji~\ref{sec:lokalne-kolancuchy-v4}.
Stąd $E^2_{p,q}=H_p(B;\mathcal H_q(F))$.

\emph{Krok 4: różniczki wyższe i zbieżność.}
Z par kolejnych filtrów otrzymujemy dokładną parę,
a z niej wyższe różniczki i strony, jak w
Twierdzeniu~\ref{thm:pages-v5}.
Dla ustalonego stopnia całkowitego $n$ w pierwszej ćwiartce
możliwe są jedynie indeksy $0\le p\le n$,
zatem filtracja homologii ma skończoną długość.
Ponieważ $E_m=E$, filtracja łańcuchów jest skończona;
nie ma ukrytych klas pochodzących z nieskończonego
dopełnienia filtracji. Standardowa
identyfikacja cykli i brzegów stabilnych z dowodu
Twierdzenia~\ref{thm:convergence-v5} daje opis $E^\infty$.
Naturalność pochodzi z naturalności wszystkich etapów:
przeciwobrazów szkieletów, homomorfizmów łączących i
transportu wzdłuż dróg.
\end{proof}

\begin{uwaga}[Uwaga o zakresie dowodu]
Krok 2 korzysta z filtracyjnej wersji wycinania dla
CW-kompleksów, a nie ze zwykłego wycinania wzdłuż
domkniętych komórek bez sprawdzenia jego hipotez.
Dla wiązek lokalnie trywialnych odpowiedni izomorfizm
otrzymuje się najpierw po zastąpieniu $E_{p-1}$ małym
otwartym otoczeniem deformacyjnie retrahującym się
na $E_{p-1}$, a następnie przez wycinanie nad rozłącznymi
otwartymi $p$-komórkami. Uogólnienie na wszystkie fibracje
Serre'a wymaga odrębnego twierdzenia porównawczego.
\end{uwaga}

\subsection*{Od wiązek do ogólnych fibracji Serre'a}
Powyższy dowód przeprowadziliśmy dla lokalnie trywialnych wiązek.
Teraz uzasadnimy rozszerzenie na fibracje Serre'a bez zakładania
lokalnej trywialności. Rozróżnienie jest ważne: fibracja Serre'a
ma własność podnoszenia homotopii względem dysków, ale nie
musi posiadać lokalnych kart trywializujących.

\begin{lemat}[Porównanie z fibracją przestrzeni dróg]
\label{lem:serre-mappingpath-v6}
Niech $p:E\to B$ będzie surjektywną fibracją Serre'a, $B$ przestrzenią
łukowo spójną, a jej włókna niech mają typ homotopijny CW-kompleksów.
Zdefiniujmy
\[
E^h=\{(e,\gamma)\in E\times B^I:\gamma(0)=p(e)\},\qquad
p^h(e,\gamma)=\gamma(1).
\]
Przestrzeń dróg $B^I$ wyposażamy w topologię zwarto-otwartą,
a $I=[0,1]$. Wtedy $p^h$ jest fibracją Hurewicza, a mapa
$j:E\to E^h$, $j(e)=(e,c_{p(e)})$, jest mapą nad $B$.
Dla każdego $b\in B$ mapa włókien
$j_b:F_b\to (p^h)^{-1}(b)$ indukuje izomorfizmy
wszystkich grup homologii singularnej.
Co więcej, dla dowolnego podkompleksu CW $A\subseteq B$
mapa $j:p^{-1}(A)\to(p^h)^{-1}(A)$ indukuje
izomorfizmy homologii, a dla pary podkompleksów
$A'\subseteq A$ również homologii względnej.
\end{lemat}
\begin{proof}
Najpierw wyjaśnijmy własność Hurewicza. Dla homotopii
$h:Z\times I\to B$ oraz podniesienia początkowego
$\widetilde h_0(z)=(e_z,\gamma_z)$, gdzie
$\gamma_z(1)=h(z,0)$, pozostawiamy $e_z$ bez zmiany,
a drogę $\gamma_z$ przedłużamy od $h(z,0)$ do $h(z,t)$
fragmentem $s\mapsto h(z,st)$. Reparametryzując oba
odcinki liniowo, uzyskujemy ciągłe podniesienie $h$.

Mapa $E\to E^h$, $e\mapsto(e,c_{p(e)})$, jest
równoważnością homotopijną \emph{całych przestrzeni}: rzut
$r:E^h\to E$, $(e,\gamma)\mapsto e$, spełnia $rj=\id$,
a $jr\simeq\id$ przez skracanie drogi $\gamma$ od jej końca.
Ta homotopia nie zachowuje jednak rzutowania na $B$;
dlatego nie wystarcza do porównania poszczególnych włókien.

Dla $b\in B$ porównujemy długie ciągi homotopii
fibracji Serre'a $F_b\to E\to B$ i fibracji Hurewicza
$(p^h)^{-1}(b)\to E^h\to B$. W podstawie występuje
identyczność, a w przestrzeni całkowitej $j$ jest
równoważnością homotopijną. Z lematu pięciu w dodatnich
stopniach oraz z dokładności na poziomie zbiorów składowych
wynika, że $j_b$ jest słabą równoważnością homotopijną.
Słaba równoważność dowolnych przestrzeni indukuje
izomorfizm homologii singularnej: wynika to z twierdzenia
Whiteheada zastosowanego do ich kanonicznych przybliżeń CW
przez realizacje geometryczne kompleksów singularnych.

Teraz niech $A\subseteq B$ będzie podkompleksem.
Cofnięcia obu fibracji nad $A$ są fibracjami Serre'a;
porównanie ich długich ciągów homotopii na każdej składowej
$A$ daje słabą równoważność na przestrzeniach całkowitych
nad $A$. Stąd izomorfizm homologii. Jeżeli
$A'\subseteq A$, stosujemy do otrzymanych map par
długie ciągi homologii i lemat pięciu, co daje ostatnią tezę.
\end{proof}

\begin{lemat}[Lokalny rachunek nad komórką dla fibracji Hurewicza]
\label{lem:serre-cell-hurewicz-v6}
Niech $q:Y\to B$ będzie fibracją Hurewicza nad CW-kompleksem
$B$, a $Y_p=q^{-1}(B^p)$. Dla każdej zorientowanej
$p$-komórki $e^p_\alpha$ wybierzmy punkt $b_\alpha$ w jej
wnętrzu. Wtedy
\[
H_{p+q}(Y_p,Y_{p-1};\Z)\cong
\bigoplus_{\alpha\in\operatorname{Cells}_p(B)}
H_q(q^{-1}(b_\alpha);\Z),\qquad q\ge0,
\]
a grupy względne w stopniu $p+q$ z $q<0$ znikają.
Po tych identyfikacjach homomorfizm łączący
$H_{p+q}(Y_p,Y_{p-1})\to H_{p+q-1}(Y_{p-1},Y_{p-2})$
jest różniczką komórkową z transportem współczynników
$H_q(q^{-1}(b))$ wzdłuż dróg w podstawie.
\end{lemat}
\begin{proof}
Dla każdej mapy charakterystycznej
$\chi_\alpha:(D^p,S^{p-1})\to(B^p,B^{p-1})$
cofamy $q$ do fibracji nad dyskiem. Lemat
\ref{lem:serre-disk-v5} daje równoważność włóknistą
z produktem $D^p\times F_\alpha$; ograniczenie na
$S^{p-1}$ daje równoważność par. Dlatego homologia
względna cofnięcia wynosi
$H_{p+q}(D^p\times F_\alpha,S^{p-1}\times F_\alpha)
\cong H_q(F_\alpha)$, co wynika z względnego wzoru
Künnetha, gdyż $H_p(D^p,S^{p-1})=\Z$.

Wyjaśnijmy przejście od pojedynczej komórki do całego
szkieletu. Wybierzmy w $B^p$ otoczenie $U$ szkieletu
$B^{p-1}$ deformacyjnie retrahujące się na $B^{p-1}$
względem tego szkieletu. Podnosząc tę deformację
z własności Hurewicza otrzymujemy równoważność
$q^{-1}(B^{p-1})\hookrightarrow q^{-1}(U)$.
Po wycięciu $q^{-1}(U)$ pozostają rozłączne części
pochodzące z wnętrz komórek. Zwarta nośność pojedynczego
łańcucha singularnego redukuje dowolny rachunek do
skończonej liczby komórek i daje sumę prostą.
Użyte wycinanie jest wycinaniem \emph{par po pogrubieniu}
skieletu, a nie nieuzasadnionym wycinaniem domkniętych
brzegów komórek.

Aby obliczyć $d^1$, wybieramy orientacje komórek i drogi
łączące punkty bazowe. Na ścianie dysku
charakterystycznego transportujemy klasę $H_q(F)$ do
włókna nad sąsiednią $(p-1)$-komórką. Wkład tej ściany
jest jej liczbą incydencji pomnożoną przez transport.
Dlaczego nie powstają inne współczynniki?
Na małym dysku w ścianie możemy przez homotopię mapy
charakterystycznej zastąpić jej zachowanie mapą o lokalnym
stopniu $+1$ lub $-1$; naturalność homomorfizmu
łączącego pary dyskowej mnoży wówczas generator
$H_p(D^p,S^{p-1})$ przez odpowiedni znak, pozostawiając
transportowaną klasę włókna. Sumowanie po ścianach
jest dokładnie definicją granicy komórkowej z lokalnymi
współczynnikami. Zależność od wyboru dróg zmienia
jedynie identyfikacje włókien systemu lokalnego,
nie sam kompleks.
\end{proof}

\begin{twierdzenie}[Ciąg spektralny Serre'a dla fibracji]
\label{thm:serre-general-v6}
Niech $p:E\to B$ będzie surjektywną fibracją Serre'a,
$B$ łukowo spójnym, skończeniewymiarowym CW-kompleksem,
a włókna niech mają typ homotopijny CW-kompleksów.
Dla wybranego $b_0\in B$ oznaczmy $F=p^{-1}(b_0)$.
Wtedy istnieje naturalny homologiczny ciąg spektralny
\[
E^2_{s,t}\cong H_s(B;\mathcal H_t(F;\Z)),\qquad
 d^r:E^r_{s,t}\longrightarrow E^r_{s-r,t+r-1},
\]
który silnie zbiega do $H_{s+t}(E;\Z)$, w tym sensie, że
$E^\infty_{s,t}\cong F_sH_{s+t}(E)/F_{s-1}H_{s+t}(E)$,
a filtracja jest skończona w każdym stopniu.
System $\mathcal H_t(F;\Z)$ wyznacza monodromia fibracji
na homologii włókien. Nie wymagamy lokalnej trywialności.
\end{twierdzenie}
\begin{proof}
Zastępujemy $p$ fibracją przestrzeni dróg $p^h:E^h\to B$
z Lematu~\ref{lem:serre-mappingpath-v6}. Rzutowanie $p^h$
ma własność podnoszenia Hurewicza, więc Lemat
\ref{lem:serre-cell-hurewicz-v6} oblicza stronę
$E^1_{s,t}=H_{s+t}((E^h)_s,(E^h)_{s-1})$
i identyfikuje $d^1$ z granicą komórkową z lokalnymi
współczynnikami $H_t((p^h)^{-1}(b))$.

Dla każdego szkieletu $B^s$ porównanie $j:E\to E^h$
indukuje izomorfizm na homologii jego przeciwobrazu;
na mocy długich ciągów par daje więc izomorfizm
$H_* (E_s,E_{s-1})\to H_*((E^h)_s,(E^h)_{s-1})$.
Są to izomorfizmy \emph{zgodne z homomorfizmami łączącymi},
ponieważ pochodzą z jednej mapy filtrowanych przestrzeni.
Dają zatem izomorfizm dokładnych par i wszystkich
stron ich ciągów spektralnych. Na włóknach $j$ indukuje
izomorfizmy homologii, więc identyfikuje systemy
lokalnych współczynników. Otrzymujemy
$E^2_{s,t}=H_s(B;\mathcal H_t(F;\Z))$.

Dla $d=\dim B<\infty$ filtracja
$E_{-1}=\varnothing\subseteq E_0\subseteq\cdots
\subseteq E_d=E$ jest ograniczona \emph{na poziomie
kompleksów łańcuchowych}. Stosujemy zatem
Twierdzenie~\ref{thm:convergence-v5}: daje ono stabilizację,
identyfikację kolejnych ilorazów i brak problemu
warunkowej zbieżności przy nieskończonej filtracji.
Naturalność zachodzi dla map fibracji respektujących
wybraną filtrację szkieletową; jej niezależność od wyboru
modelu CW wynika z porównania przez wspólne
przybliżenia komórkowe.
\end{proof}
\begin{uwaga}
Dla bazy o nieskończonym wymiarze także istnieje klasyczny
ciąg Serre'a w szerokiej ogólności. Przejście od
skończeniewymiarowych filtracji do sytuacji nieskończonej
wymaga osobnej analizy zbieżności; nie jest wnioskiem
z samego twierdzenia o filtracji ograniczonej.
\end{uwaga}


\section{Wersja kohomologiczna i transgresja}
\begin{twierdzenie}[Kohomologiczny ciąg spektralny]
Jeżeli $F\to E\to B$ jest lokalnie trywialną wiązką
nad spójnym skończeniewymiarowym CW-kompleksem $B$, to istnieje ciąg spektralny
kohomologiczny, dla którego
\[
E_2^{p,q}\cong H^p(B;\mathcal H^q(F)),\qquad
 d_r:E_r^{p,q}\longrightarrow E_r^{p+r,q-r+1},
\]
a $E_\infty^{p,q}$ są kolejnymi ilorazami odpowiedniej
filtracji $H^{p+q}(E;\Z)$. W przypadku trywialnej monodromii
system $\mathcal H^q(F)$ jest stały. Dla orientowalnej
wiązki sferowej najważniejszą różniczką jest transgresja.
\end{twierdzenie}
\begin{proof}
Zamiast filtrować kompleks łańcuchowy, filtrujemy kompleks
kołańcuchów względnych względem przeciwobrazów szkieletów.
Otrzymana dokładna para ma różniczkę o indeksach
$(r,1-r)$; lokalne obliczenie grup względnych jest dualne do
rachunku z dowodu Twierdzenia~\ref{thm:serre-v5}.
Dokładność w kolejnych krokach zapewnia długi ciąg
kohomologii pary. Ponieważ szkielet $B^m=B$ dla pewnego $m$, filtracja
kołańcuchów jest ograniczona i zbieżność wynika z
twierdzenia o filtracji ograniczonej.
\end{proof}
\begin{definicja}[Klasa transgresywna]
W kohomologicznym ciągu spektralnym klasę
$u\in E_2^{0,q}$ nazywamy \emph{transgresywną}, jeśli
przeżywa do strony $E_{q+1}^{0,q}$. Wówczas
$d_{q+1}u\in E_{q+1}^{q+1,0}$ jest jej transgresją.
Nie każdy element $H^q(F)$ jest transgresywny:
wcześniejsza różniczka może go unicestwić.
\end{definicja}
\begin{uwaga}[Iloczyn i reguła Leibniza]
Jeśli $B$ jest jednospójna, włókno ma odpowiedni typ CW,
a współczynniki należą do przemiennego pierścienia,
ciąg kohomologiczny ma strukturę multiplikatywną;
dla jednorodnego $a\in E_r^{p,q}$ i $b$ zachodzi
$d_r(ab)=d_r(a)b+(-1)^{p+q}a\,d_r(b)$.
Wymaga to filtracyjnego odwzorowania diagonalnego
i konstrukcji iloczynu na stronie zerowej; jest ono zgodne
z iloczynem kubkowym całkowitej kohomologii. Nie wolno
zakładać identycznego wzoru przy lokalnym systemie o
nieprzemiennym działaniu bez zdefiniowania odpowiedniego
sparowania systemów współczynników.
\end{uwaga}

\begin{wniosek}[Kohomologiczna wersja dla fibracji Serre'a]
Przy hipotezach Twierdzenia~\ref{thm:serre-general-v6}
ciąg spektralny z poprzedniego twierdzenia istnieje także
dla fibracji Serre'a. Jego strona druga ma postać
$E_2^{p,q}=H^p(B;\mathcal H^q(F;\Z))$ i zbiega
do $H^{p+q}(E;\Z)$ przez skończoną filtrację.
\end{wniosek}
\begin{proof}
Porównanie z fibracją przestrzeni dróg z
Lematu~\ref{lem:serre-mappingpath-v6} daje izomorfizm
kohomologii wszystkich par preobrazów szkieletów.
Po odwróceniu strzałek w ich długich ciągach dokładnych
otrzymujemy izomorfizm kohomologicznych dokładnych par.
Obliczenie komórkowe jest dualne do
Lematu~\ref{lem:serre-cell-hurewicz-v6}; przy nieskończonym
zbiorze komórek w danym wymiarze sumę łańcuchów
zastępuje iloczyn kołańcuchów, jak wymaga definicja
kohomologii komórkowej. Skończeniewymiarowość bazy
zapewnia skończoną filtrację i zbieżność.
\end{proof}

\section{Przykład podstawowy: fibracja Hopfa i niezerowa \texorpdfstring{$d_2$}{d2}}
Rozważmy wiązkę $S^1\to S^3\to S^2$ z
Twierdzenia~\ref{thm:hopf-v4}. Podstawa jest jednospójna,
więc monodromia jest trywialna. Na stronie drugiej
mamy tylko cztery niezerowe pozycje:
$E^2_{0,0}=E^2_{2,0}=E^2_{0,1}=E^2_{2,1}=\Z$.

\begin{center}
\begin{tikzpicture}[x=1.25cm,y=1.0cm,>=Stealth]
\draw[->] (-.35,0)--(3.3,0) node[right]{$p$};
\draw[->] (0,-.25)--(0,2.3) node[above]{$q$};
\foreach \i in {0,1,2,3}{\draw(\i,.07)--(\i,-.07) node[below] {\small\i};}
\foreach \i in {0,1,2}{\draw(.07,\i)--(-.07,\i) node[left] {\small\i};}
\foreach \x/\y in {0/0,2/0,0/1,2/1}{\node[draw,rounded corners,fill=blue!7,minimum width=.65cm,minimum height=.43cm] at (\x,\y){$\Z$};}
\draw[->,very thick,red!70!black] (1.75,.13)--(.25,.88) node[midway,above] {$d^2$};
\node[align=left] at (2.45,1.64){\footnotesize $E^2_{p,q}$\\[-1pt]\footnotesize fibracji Hopfa};
\end{tikzpicture}
\end{center}

\begin{twierdzenie}
W homologicznej sekwencji Hopfa homomorfizm
$d^2:E^2_{2,0}\to E^2_{0,1}$ jest izomorfizmem
mnożenia przez $\pm1$. W konsekwencji na stronie
$E^\infty$ pozostają tylko grupy w pozycjach $(0,0)$
i $(2,1)$, obie izomorficzne z $\Z$.
\end{twierdzenie}
\begin{proof}
Na stronie $E^2$ nie ma niezerowych różniczek poza
wskazaną: różniczka ma bistopień $(-r,r-1)$,
a podstawa ma tylko stopnie $0$ i $2$.
Pozycja $(2,0)$ odpowiada stopniowi całkowitemu $2$,
a $(0,1)$ stopniowi $1$. Ponieważ
$H_1(S^3;\Z)=H_2(S^3;\Z)=0$, obie pozycje muszą
zniknąć na stronie $E^\infty$; jedyną możliwością
jest $d^2$ o zerowym jądrze i zerowym kokernelu.
Jako endomorfizm $\Z$ jest więc mnożeniem przez $\pm1$.
Pozostałe dwa wyrazy nie mogą zostać zabite ani nie mogą
być obrazami różniczki z pierwszej ćwiartki.
Na stopniach $0$ i $3$ nie ma problemu rozszerzeń,
co odtwarza homologię $S^3$.
\end{proof}
\begin{uwaga}
W kohomologii ta sama fibracja daje
$d_2:E_2^{0,1}\to E_2^{2,0}$, również izomorfizm.
Klasa generatora $H^1(S^1;\Z)$ transgresuje do generatora
$H^2(S^2;\Z)$. Nie należy zamieniać kierunków strzałek
homologicznej i kohomologicznej różniczki.
\end{uwaga}

\section{Przykład z nietrywialną monodromią: butelka Kleina}
Niech $K$ będzie odwzorowaniem sklejenia okręgu przez
odbicie, czyli butelką Kleina opisaną w rozdziale o fibracjach.
Mamy wiązkę $S^1\to K\to S^1$. Działanie monodromii
na $H_0(S^1;\Z)=\Z$ jest trywialne, a na
$H_1(S^1;\Z)=\Z$ jest mnożeniem przez $-1$.
Dla współczynników $\Z_-$ otrzymujemy komórkowy kompleks
$\Z\xrightarrow{-2}\Z$, a dla stałych $\Z$
kompleks $\Z\xrightarrow{0}\Z$. Zatem
\[
E^2_{0,0}=\Z,\quad E^2_{1,0}=\Z,\quad
E^2_{0,1}=\Z/2,\quad E^2_{1,1}=0.
\]
Ponieważ $\dim S^1=1$, wszystkie $d^r$ dla $r\ge2$
mają zerowe cele. Dostajemy $H_0(K;\Z)=\Z$,
$H_2(K;\Z)=0$ i krótką sekwencję
$0\to\Z/2\to H_1(K;\Z)\to\Z\to0$.
Rozszczepia się ona, bo $\Z$ jest wolna, zatem
$H_1(K;\Z)\cong\Z\oplus\Z/2$.
Wyrzucenie systemów lokalnych ze wzoru $E^2$ dałoby
błędny wynik $H_1(K)\cong\Z^2$.

\section{Wiązki sferowe, klasa Eulera i ciąg Gysina}
\begin{definicja}[Wiązka sferowa zorientowanej wiązki]
Niech $\xi\to B$ będzie orientowalną rzeczywistą wiązką
wektorową rzędu $m+1$, wyposażoną w metrykę włóknową.
Niech $D(\xi)$ i $S(\xi)$ oznaczają wiązki odpowiednio
jednostkowych dysków i jednostkowych sfer. Mamy
wiązki $S^m\to S(\xi)\xrightarrow{\pi}B$ i
$D^{m+1}\to D(\xi)\to B$.
\end{definicja}
\begin{twierdzenie}[Izomorfizm Thoma i klasa Eulera]
\label{thm:thom-v5}
Niech $B$ będzie skończonym CW-kompleksem, a
$\xi\to B$ zorientowaną rzeczywistą wiązką wektorową
rzędu $r\ge1$ z metryką włóknową. Oznaczmy przez
$D(\xi)$ i $S(\xi)$ podwiązki jednostkowych dysków
oraz sfer, a przez $\pi:D(\xi)\to B$ rzutowanie.
Wówczas istnieje dokładnie jedna klasa
$u_\xi\in H^r(D(\xi),S(\xi);\Z)$, której ograniczenie
do każdej pary włókien $(D^r_b,S^{r-1}_b)$ jest
jej dodatnim generatorem wyznaczonym przez orientację
$\xi_b$. Dla każdego $k\ge0$ homomorfizm
\[
T_\xi:H^k(B;\Z)\longrightarrow
H^{k+r}(D(\xi),S(\xi);\Z),\qquad
\alpha\longmapsto\pi^*\alpha\smile u_\xi
\]
jest izomorfizmem. Obraz $u_\xi$ przez
$H^r(D(\xi),S(\xi))\to H^r(D(\xi))\cong H^r(B)$
nazywamy klasą Eulera $e(\xi)$.
\end{twierdzenie}
\begin{proof}
\emph{Krok 1: filtracja pary wiązek.}
Dla szkieletu $B^p$ oznaczmy
$D_p=D(\xi)|_{B^p}$, $S_p=S(\xi)|_{B^p}$.
Relatywny kompleks kołańcuchów pary $(D,S)$
filtrujemy przez jądra ograniczeń do
$(D_{p-1},S_{p-1})$. Różniczki kolejnych stron
tej filtracji pochodzą z długich ciągów dokładnych
trójek $(D_p,D_{p-1}\cup S_p,S_p)$; całkowita
kohomologia jest $H^*(D,S)$, ponieważ $B$ ma
skończony wymiar.

\emph{Krok 2: wyrazy względne nad komórkami.}
Dla zorientowanej $p$-komórki $e^p_\alpha$ cofnięcie
wiązki na dysk charakterystyczny jest trywialne.
Wycinanie po pogrubieniu brzegu dysku i względny
iloczyn Künnetha dają
\[
\begin{aligned}
&H^{p+q}(D_p,D_{p-1}\cup S_p;\Z)\\
&\qquad\cong
\bigoplus_{\alpha\in\operatorname{Cells}_p(B)}
H^p(D^p,S^{p-1};\Z)\otimes
H^q(D^r,S^{r-1};\Z).
\end{aligned}
\]
Wszystkie grupy $H^q(D^r,S^{r-1};\Z)$ znikają
poza $q=r$, gdzie są równe $\Z$. Orientacja $\xi$
sprawia, że transport tego generatora wokół każdej
pętli jest tożsamością. Otrzymany \emph{jedyny}
niezerowy wiersz strony pierwszej jest więc
komórkowym kompleksem kołańcuchów
$C^p_{\mathrm{CW}}(B;\Z)$, przesuniętym w stopniu o $r$.
Wobec tego
\[
E_2^{p,q}=\begin{cases}H^p(B;\Z),&q=r,\\0,&q\ne r.
\end{cases}
\]
Żadna różniczka $d_s$ dla $s\ge2$ nie może łączyć
niezerowych wierszy. Ciąg spektralny stabilizuje się
na $E_2$. Dla każdego całkowitego stopnia $k+r$
pozostaje dokładnie jeden możliwy iloraz filtracji,
zatem nie powstaje także problem rozszerzeń.
W rezultacie otrzymujemy naturalne \emph{addytywne}
izomorfizmy
$\Theta_k:H^k(B;\Z)\xrightarrow{\cong}H^{k+r}(D,S;\Z)$.

\emph{Krok 3: konstrukcja globalnej klasy.}
Niech $1_B\in H^0(B;\Z)$ będzie klasą funkcji
przyjmującej wartość $1$ na każdej składowej bazy.
Definiujemy $u_\xi=\Theta_0(1_B)$.
Po ograniczeniu filtracji do dowolnej komórki
powraca wybrany dodatni generator względnej kohomologii
włókna. Ponieważ ograniczenia do włókien są zgodne
z transportem indukowanym przez orientację,
klasa $u_\xi$ ma żądaną własność. Jeżeli druga klasa
ma te same ograniczenia, jej różnica odpowiada
poprzez $\Theta_0$ klasie z $H^0(B)$ równej zeru
na każdej składowej, jest więc zerowa.

\emph{Krok 4: identyfikacja izomorfizmu z iloczynem.}
Zdefiniujmy $T_\xi(\alpha)=\pi^*\alpha\smile u_\xi$.
Filtracja komórkowa jest zgodna z iloczynem kubkowym:
na ilorazie odpowiadającym komórce $e^p$ mapa
$T_\xi$ jest zewnętrznym iloczynem
\[
H^p(D^p,S^{p-1};\Z)\longrightarrow
H^{p+r}(D^p\times D^r,
 (S^{p-1}\times D^r)\cup(D^p\times S^{r-1});\Z)
\]
z dodatnim generatorem $H^r(D^r,S^{r-1};\Z)$.
Obie grupy są izomorficzne z $\Z$, a mapa wysyła
generator na generator, jest więc izomorfizmem.
Przez naturalność iloczynu otrzymujemy izomorfizm
na wszystkich składnikach strony pierwszej oraz
przemienność z różniczką $d_1$.
Porównanie filtracji skończonych (indukcyjnie z lematu
pięciu dla kolejnych krótkich ciągów filtracji)
pokazuje, że $T_\xi$ jest izomorfizmem w każdym stopniu.
Ponieważ $T_\xi(1_B)=u_\xi$, jego działanie zgadza się
z wyżej skonstruowanym izomorfizmem $\Theta_k$.
\end{proof}
\begin{uwaga}[Wiązka nieorientowalna i współczynniki lokalne]
Jeżeli $\xi$ nie jest orientowalna, drugi krok nie daje
stałego wiersza $\Z$: otrzymujemy lokalny system orientacji
$\mathcal O_\xi$ i odpowiednią wersję izomorfizmu Thoma
z lokalnymi współczynnikami. Dlatego w zwykłym
całkowitoliczbowym twierdzeniu powyżej orientowalność
nie jest założeniem zbędnym.
\end{uwaga}


\begin{twierdzenie}[Ciąg Gysina]
\label{thm:gysin-v5}
Dla zorientowanej wiązki sferowej nad skończonym CW-kompleksem
jak powyżej istnieje
długi ciąg dokładny
\[
\begin{aligned}
\cdots &\to H^{k-m-1}(B;\Z)
\xrightarrow{\smile e(\xi)}H^k(B;\Z)
\xrightarrow{\pi^*}H^k(S(\xi);\Z)\\
&\xrightarrow{\pi_!}H^{k-m}(B;\Z)
\xrightarrow{\smile e(\xi)}H^{k+1}(B;\Z)\to\cdots.
\end{aligned}
\]
Jeżeli $m\ge1$, to w kohomologicznym ciągu spektralnym
Serre'a generator $H^m(S^m;\Z)$, jeśli jest transgresywny,
przechodzi pod $d_{m+1}$ w klasę $e(\xi)$
(z dokładnością do znaku). Dla $m=0$ włókno $S^0$
nie jest spójne i powyższy opis transgresji nie ma
tej postaci; sam ciąg Gysina nadal wynika z pary
$(D(\xi),S(\xi))$.
\end{twierdzenie}
\begin{proof}
Stosujemy długi ciąg dokładny pary
$(D(\xi),S(\xi))$:
$H^k(D,S)\to H^k(D)\to H^k(S)\to H^{k+1}(D,S)$.
Ponieważ $D(\xi)$ deformacyjnie retrahuje się na przekrój
zerowy, mamy $H^k(D)\cong H^k(B)$.
Z Twierdzenia~\ref{thm:thom-v5} mamy
$H^k(D,S)\cong H^{k-m-1}(B)$ oraz
$H^{k+1}(D,S)\cong H^{k-m}(B)$.
Mapa $H^k(D,S)\to H^k(D)$ jest przez naturalność
iloczynu kubkowego mnożeniem przez obraz klasy Thoma,
czyli przez $e(\xi)$. Homomorfizm $\pi_!$ jest
połączeniem mapy łączącej i odwrotności izomorfizmu Thoma.
To daje wszystkie strzałki oraz dokładność.

Dla $m\ge1$ w ciągu spektralnym wiązki sferowej możliwe
są jedynie wiersze $q=0,m$. Na pierwszej stronie, na której
różniczka może łączyć te wiersze, jej stopień wynosi $m+1$.
Porównanie homomorfizmu łączącego pary dyskowej
z homomorfizmem z dokładnej pary filtracji identyfikuje
transgresję generatora włókna z obrazem klasy Thoma.
Zgodność znaku zależy od przyjętych orientacji komórek.
\end{proof}
\begin{przyklad}
Dla trywialnej wiązki $\xi$ istnieje nigdzie niezerowy
przekrój, więc $e(\xi)=0$, a ciąg Gysina rozszczepia
poszczególne grupy (niekoniecznie w sposób kanoniczny).
Dla wiązki Hopfa $S^1\to S^3\to S^2$ klasa Eulera
odpowiedniej zorientowanej wiązki rzeczywistej rzędu $2$
ma wartość $\pm1$ razy generator $H^2(S^2;\Z)$,
co jest równoważne izomorfizmowi $d_2$ w poprzednim rachunku.
\end{przyklad}

\section{Zastosowanie do przestrzeni pętli i ograniczenia metody}
Dla $n\ge2$ mamy fibrację przestrzeni dróg
$\Omega S^n\to PS^n\to S^n$, której przestrzeń całkowita
jest ściągalna. W tym przykładzie korzystamy z \emph{ogólnej}
wersji ciągu spektralnego Serre'a dla fibracji Hurewicza:
fibracja przestrzeni dróg nie została tutaj wykazana jako
lokalnie trywialna wiązka, zatem jej użycie wykracza poza
wariant konstrukcyjnie udowodniony w tym rozdziale.
Poniższy rachunek jest poprawnym zastosowaniem tego
ogólnego twierdzenia, a nie jego dodatkowym dowodem. Ponieważ $S^n$ jest jednospójna,
lokalne współczynniki są trywialne. Na stronie drugiej
występują tylko kolumny $p=0,n$.
Różniczka mogąca unicestwić klasę
$E^2_{n,0}=H_n(S^n;\Z)\cong\Z$ ma stopień $r=n$
i cel $E^n_{0,n-1}$. Wszystkie wcześniejsze różniczki
z tej pozycji są zerowe z powodów wymiarowych,
a pozycja $(n,0)$ nie może być trafiona przez różniczkę
z pierwszej ćwiartki o $p>n$. Skoro $H_n(PS^n)=0$,
$d^n:E^n_{n,0}\to E^n_{0,n-1}$ musi być injekcją.
Również $H_{n-1}(PS^n)=0$, a $(0,n-1)$ może być
trafiona tylko w ten sposób, więc ta mapa jest surjekcją.
Zatem $H_{n-1}(\Omega S^n;\Z)\cong\Z$.
Jest to zgodne z izomorfizmem
$\pi_{n-1}(\Omega S^n)\cong\pi_n(S^n)$, ale wynika tu
z rachunku homologicznego.

\begin{uwaga}[Czego jeszcze nie udowodniliśmy?]
Niniejszy rozdział daje pełną algebraiczną konstrukcję
stron i różniczek filtrowanego kompleksu oraz rachunki
ciągu Serre'a dla wiązek lokalnie trywialnych.
Wersja dla wszystkich fibracji Serre'a, pełny dowód
wielo- i naturalności transgresji, ogólne twierdzenia o
silnej zbieżności filtracji nieograniczonych i struktura
multiplikatywna z dowolnymi skręconymi współczynnikami
wymagają osobnych lematów technicznych. Nie są automatycznymi
konsekwencjami samego zapisu $E_2^{p,q}=H^p(B;H^q(F))$.
\end{uwaga}

\section{Podsumowanie metody obliczeniowej}
Przy konkretnej fibracji ustalamy najpierw system lokalny
$\mathcal H_q(F)$ i rzeczywiste działanie monodromii.
Następnie obliczamy $H_p(B;\mathcal H_q(F))$ oraz
rozrysowujemy niezerowe pozycje strony $E^2$.
Sprawdzamy, które różniczki mogą być niezerowe z uwagi
na ich bistopień i znane grupy przestrzeni całkowitej.
Dopiero po obliczeniu $E^\infty$ rozwiązujemy problem
rozszerzeń. Wersja kohomologiczna dostarcza dodatkowo
iloczynów oraz transgresji, której podstawowym przykładem
jest klasa Eulera zorientowanej wiązki sferowej.
\clearpage


\chapter{Przestrzenie klasyfikujące i teoria przeszkód}
\label{ch:klasyfikacja-v6}

\section{Dlaczego potrzebujemy przestrzeni klasyfikujących?}
W rozdziałach o homologii i kohomologii przypisaliśmy przestrzeni $X$
grupy $H^n(X;A)$. Czy grupy te można opisać za pomocą
\emph{odwzorowań} z $X$ do jednej przestrzeni, zależnej tylko od
$A$ i $n$? Podobnie: czy wszystkie wiązki danego typu nad $X$
można uzyskać jako cofnięcia jednej wiązki uniwersalnej?
Na oba pytania odpowiedź jest twierdząca po sprecyzowaniu
hipotez. Wprowadzimy przestrzenie Eilenberga--Mac Lane'a,
przestrzenie klasyfikujące grup i teorię przeszkód, która
wyjaśnia, dlaczego obiekty globalne mogą nie istnieć
pomimo lokalnego istnienia.

Zakładamy w tym rozdziale, że przestrzenie służące za
podstawy klasyfikacji są CW-kompleksami. Gdy mówimy
o ogólnych wiązkach głównych nad przestrzenią $X$,
zakładamy, że $X$ jest parazwarta, a wiązka jest
\emph{numerowalna} (istnieje podział jedności podporządkowany
pokryciu trywializującemu). Dla CW-kompleksów
warunek ten nie stanowi dodatkowego ograniczenia
w standardowej teorii wiązek lokalnie trywialnych.

\section{Przestrzenie Eilenberga--Mac Lane'a}
\begin{definicja}
Niech $n\ge1$. Dla $n=1$ niech $G$ będzie dowolną grupą;
dla $n\ge2$ niech $G$ będzie grupą abelową.
\emph{Przestrzenią Eilenberga--Mac Lane'a typu $K(G,n)$}
nazywamy łukowo spójny CW-kompleks $K$ z wyróżnionym
punktem, dla którego
\[
\pi_j(K)=\begin{cases}G,&j=n,\\0,&j\ge1,\ j\ne n.
\end{cases}
\]
Przy identyfikacji $\pi_n(K)\cong G$ ustalamy
izomorfizm, a nie tylko istnienie izomorfizmu.
\end{definicja}
\begin{przyklad}
Mamy $S^1\simeq K(\Z,1)$, a nieskończeniewymiarowa
przestrzeń rzutowa rzeczywista $\R P^\infty$ jest
$K(\Z/2,1)$: jej dwukrotne nakrycie $S^\infty$
jest ściągalne. Pokażemy także, że
$\C P^\infty\simeq K(\Z,2)$.
Sfera $S^n$ dla $n\ge2$ \emph{nie} jest przestrzenią
$K(\Z,n)$, gdyż posiada niezerowe grupy homotopii
w wyższych wymiarach.
\end{przyklad}

\begin{twierdzenie}[Istnienie przestrzeni $K(G,n)$]
\label{thm:KGn-existence-v6}
Dla każdej grupy $G$ istnieje CW-kompleks $K(G,1)$.
Dla każdej grupy abelowej $A$ i $n\ge2$ istnieje
CW-kompleks $K(A,n)$. Każdy taki obiekt jest jednoznaczny
z dokładnością do równoważności homotopijnej zachowującej
ustaloną identyfikację odpowiedniej grupy homotopii.
\end{twierdzenie}
\begin{proof}
Najpierw skonstruujmy $K(G,1)$. Wybierzmy prezentację
grupy $G$ przez generatory i relacje, niekoniecznie skończoną.
Zbudujmy spójny dwuwymiarowy CW-kompleks $X_2$:
jeden wierzchołek, jedna $1$-komórka dla każdego generatora
i jedna $2$-komórka przyklejona wzdłuż słowa odpowiadającego
każdej relacji. Twierdzenie van Kampena daje $\pi_1(X_2)\cong G$.
Następnie dla wszystkich generatorów grupy $\pi_2(X_2)$
doklejamy $3$-komórki wzdłuż reprezentujących je map
$S^2\to X_2$. W ten sposób zabijamy $\pi_2$, nie zmieniając
$\pi_1$. Powtarzamy indukcyjnie: na etapie $k$ doklejamy
$(k+1)$-komórki zabijające $\pi_k$, bez zmiany $\pi_j$
dla $j<k$. W słabej topologii granicznego CW-kompleksu
każda mapa sfery ma obraz homotopijnie reprezentowalny
w skończonym etapie odpowiedniej filtracji komórkowej;
stąd otrzymujemy $\pi_k=0$ dla $k\ge2$.

Dla grupy abelowej $A$ i $n\ge2$ wybieramy prezentację
$0\to R\to F\to A\to0$, w której $R$ i $F$ są wolne.
Bierzemy bukiet sfer $S^n$ indeksowany bazą $F$;
reprezentanci relacji $R\subset F\cong\pi_n(\bigvee S^n)$
wyznaczają mapy dołączające $(n+1)$-komórki.
Powstały kompleks jest $(n-1)$-spójny i ma
$H_n\cong A$, zatem z twierdzenia Hurewicza
$\pi_n\cong A$. Dołączając kolejno komórki wymiarów
$n+2,n+3,\ldots$ zabijamy pozostałe grupy homotopii,
nie zmieniając $\pi_n$. Otrzymujemy $K(A,n)$.

Aby wyjaśnić jednoznaczność, zauważmy, że dla dwóch takich
kompleksów $K,L$ homomorfizm $\pi_n(K)\to\pi_n(L)$
można zrealizować odwzorowaniem komórkowym: na
$n$-szkielecie podajemy reprezentanty grupy $\pi_n(L)$,
a przeszkody rozszerzania na wyższe komórki znikają,
ponieważ $\pi_j(L)=0$ dla $j>n$. Dla $n=1$ realizację
homomorfizmu $\pi_1$ uzyskujemy z prezentacji przez
wypełnienie relacji, a dalsze komórki dołączamy, gdyż
$\pi_j(L)=0$ dla $j>1$. W ten sposób konstruujemy mapę
indukującą zadany izomorfizm $\pi_n$ (odpowiednio $\pi_1$).
Jest to słaba równoważność homotopijna, a więc na mocy
twierdzenia Whiteheada równoważność homotopijna.
\end{proof}

\section{Kohomologia jako zbiór klas homotopii}
Niech $[X,Y]$ oznacza zbiór klas homotopii ciągłych
odwzorowań $X\to Y$. Dla $n\ge1$ i grupy abelowej $A$
zbiór $[X,K(A,n)]$ ma naturalną strukturę grupy abelowej,
choć nie jest ona oczywista z samej definicji homotopii.

\begin{twierdzenie}[Reprezentacja kohomologii singularnej]
\label{thm:rep-coho-v6}
Dla dowolnego CW-kompleksu $X$, grupy abelowej $A$ i
$n\ge1$ istnieje naturalna bijekcja
\[
[X,K(A,n)]\cong H^n(X;A).
\]
Jest ona izomorfizmem grup przy naturalnej strukturze
na zbiorze klas homotopii po lewej stronie.
Dla $n\ge2$ tę samą bijekcję można sformułować
w wersji z wyróżnionym punktem dla spójnych CW-kompleksów.
\end{twierdzenie}
\begin{proof}
Pokażemy szczegółowo mechanizm komórkowy dla $n\ge2$;
przypadek $n=1$ otrzymuje się analogicznie z $K(A,1)$
lub przez klasyfikację homomorfizmów do grupy abelowej $A$.

\emph{Krok 1: kołańcuch z odwzorowania.}
Ponieważ $K(A,n)$ jest $(n-1)$-spójna, za pomocą
przybliżenia komórkowego i kolejnych homotopii na komórkach
możemy doprowadzić mapę $f:X\to K(A,n)$ do postaci
stałej na $X^{n-1}$. Na każdej zorientowanej $n$-komórce
$e^n_\alpha$ mapa ilorazowa
$D^n/\partial D^n\cong S^n\to K(A,n)$
wyznacza element $a_\alpha\in\pi_n(K(A,n))=A$.
Przypisanie $e^n_\alpha\mapsto a_\alpha$ definiuje
kołańcuch komórkowy $c_f\in C^n_{\mathrm{CW}}(X;A)$.

\emph{Krok 2: warunek kocyklu.}
Niech $e^{n+1}$ będzie komórką dołączoną mapą
$\varphi:S^n\to X^n$. Obstrukcja rozszerzenia mapy
z brzegu na jej dysk leży w $\pi_n(K(A,n))=A$.
Po zastosowaniu komórkowego przybliżenia $\varphi$
obstrukcja jest sumą elementów $a_\alpha$ ważonych
liczbami incydencji $[e^{n+1}:e^n_\alpha]$.
Jest to dokładnie $(\delta c_f)(e^{n+1})$.
Ponieważ $f$ istnieje na całym $X$, otrzymujemy
$\delta c_f=0$.

\emph{Krok 3: realizacja dowolnego kocyklu.}
Odwrotnie, dla $c\in Z^n_{\mathrm{CW}}(X;A)$
ustalamy mapę stałą na $X^{n-1}$ i wybieramy na
każdej $n$-komórce reprezentanta elementu $c(e^n)$.
Warunek $\delta c=0$ jest dokładnie warunkiem
wypełnienia map na granicach wszystkich $(n+1)$-komórek.
Na komórkach wyższych wymiarów przeszkody leżą
w $\pi_j(K(A,n))$ dla $j>n$ i dlatego znikają.
Powstaje mapa $f_c:X\to K(A,n)$.

\emph{Krok 4: homotopie i kobrzegi.}
Jeżeli $c'=c+\delta b$, gdzie
$b\in C^{n-1}_{\mathrm{CW}}(X;A)$, konstruujemy
homotopię na $X\times I$ przez nadanie na pryzmatach
$(n-1)$-komórek elementów $b(e^{n-1})\in A$.
Warunek możliwości rozszerzenia przez pryzmaty
$n$-komórek ma postać
$c'-c-\delta b=0$; powyżej tych wymiarów obstrukcje
znikają. Odwrotnie każda homotopia $f\simeq g$ po
komórkowym uporządkowaniu względem $X\times\partial I$
wyznacza kołańcuch $b$ z $c_g-c_f=\delta b$.
Zatem klasy homotopii odpowiadają dokładnie klasom
$Z^n/B^n=H^n_{\mathrm{CW}}(X;A)$, które są naturalnie
izomorficzne z kohomologią singularną.
Dodawanie klas odpowiada dodawaniu reprezentujących
je kocykli, co nadaje lewej stronie strukturę grupową.
\end{proof}

\begin{przyklad}[Dlaczego samą grupą nie da się zastąpić pierścienia]
Niech $h\in H^2(\C P^2;\Z)$ będzie generatorem.
Z twierdzenia reprezentacyjnego odpowiada mu mapa
$f:\C P^2\to K(\Z,2)$. Iloczyn
$h\smile h\in H^4(\C P^2;\Z)$ jest niezerowy,
natomiast na $S^2\vee S^4$ iloczyn dwóch klas
stopnia $2$ jest zerowy. Grupy addytywne homologii
obu przestrzeni są jednak izomorficzne. Stąd pełny
pierścień kohomologii zawiera informację niewidoczną
w samych grupach.
\end{przyklad}

\section{Przestrzeń klasyfikująca grupy dyskretnej}
\begin{definicja}[Wiązka główna]
Niech $G$ będzie grupą topologiczną. \emph{Główną
wiązką $G$} nad $B$ nazywamy lokalnie trywialne
rzutowanie $P\to B$ z ciągłym, wolnym prawym działaniem
$G$ na $P$, które lokalnie identyfikuje $P|_U$
z $U\times G$ i działanie z prawym mnożeniem na $G$.
Izomorfizmy wiązek mają zachowywać zarówno rzutowanie,
jak i działanie $G$.
\end{definicja}

\begin{twierdzenie}[Konstrukcja $EG\to BG$ dla grupy dyskretnej]
\label{thm:bar-BG-v6}
Dla dowolnej grupy dyskretnej $G$ istnieje
ściągalny CW-kompleks $EG$ z wolnym działaniem $G$,
dla którego iloraz $BG=EG/G$ jest $K(G,1)$.
Rzutowanie $EG\to BG$ jest uniwersalną główną
wiązką $G$ wśród wiązek nad CW-kompleksami.
\end{twierdzenie}
\begin{proof}
Konstrukcję można podać bez odwołania do
nieskończonych przestrzeni Grassmanna. Bierzemy
realizację geometryczną zbioru symplicjalnego
złożonego z ciągów $(g_0,\ldots,g_k)\in G^{k+1}$:
opuszczanie jednego wyrazu definiuje ściany,
a powtarzanie wyrazu degeneracje. Równoczesne
lewe mnożenie wszystkich współrzędnych daje wolne
działanie $G$ na realizacji $EG$. Operator dodający
z przodu element neutralny,
$(g_0,\ldots,g_k)\mapsto(1,g_0,\ldots,g_k)$,
jest dodatkową degeneracją; odpowiadające mu stożkowanie
w realizacji daje kontrakcję $EG$.
Rzutowanie $EG\to BG$ jest nakryciem o grupie
transformacji pokrycia $G$ i ma ściągalną przestrzeń
całkowitą, zatem $\pi_1(BG)\cong G$ i
$\pi_j(BG)=0$ dla $j>1$. Jest więc modelem $K(G,1)$.

Dla spójnego CW-kompleksu $X$ główne wiązki
grupy dyskretnej $G$ odpowiadają homomorfizmom
$\pi_1(X)\to G$ z dokładnością do sprzężenia:
homomorfizm jest monodromią nakrycia, a nakrycie
odtwarzamy jako iloraz $\widetilde X\times G$
przez działanie grupy $\pi_1(X)$ wyznaczone tym
homomorfizmem. Z drugiej strony
$[X,BG]$ jest tym samym zbiorem klas homomorfizmów,
ponieważ $BG=K(G,1)$. Zatem każdą wiązkę otrzymujemy
jako cofnięcie $EG\to BG$; zgodność z homotopiami
map klasyfikujących wynika z niezmienniczości
monodromii względem homotopii.
Dla niespójnego $X$ konstrukcję stosujemy oddzielnie
do każdej składowej.
\end{proof}

\begin{uwaga}[Nienaturalność wyboru monodromii]
Po ustaleniu punktów w bazie i we włóknie otrzymuje się
konkretny homomorfizm $\pi_1(X)\to G$.
Bez wyboru punktu we włóknie zmiana trywializacji
sprzęga ten homomorfizm przez element $G$.
Dlatego w klasyfikacji \emph{niepunktowanych} wiązek
występują klasy sprzężoności.
\end{uwaga}

\section{Grupy topologiczne, wiązki uniwersalne i przestrzenie Grassmanna}
Dla grupy topologicznej $G$ model $BG$ można otrzymać
z nieskończonego złączenia $EG=G*G*G*\cdots$.
Punkt w tym złączeniu ma postać skończonej formalnej
kombinacji $\sum_i t_i g_i$, gdzie $t_i\ge0$,
$\sum_i t_i=1$, a współrzędną $g_i$ pomijamy,
gdy $t_i=0$. Grupa $G$ działa równocześnie prawym
mnożeniem na wszystkich niezerowych współrzędnych.

\begin{lemat}[Ściągalność modelu uniwersalnego]
Nieskończone złączenie $EG=G*G*\cdots$ jest ściągalne,
a diagonalne działanie $G$ jest wolne.
\end{lemat}
\begin{proof}
Każdy punkt ma co najmniej jedną współrzędną
o dodatniej wadze, na której prawostronna translacja
jest wolna; zatem całe działanie jest wolne.
Dla ściągalności wygodny jest równoważny
\emph{symplicjalny model} przestrzeni $EG$ z
$k$-sympleksami $G^{k+1}$, którego operatory ścian
opuszczają współrzędne, a degeneracje je powtarzają.
Operacja
$(g_0,\ldots,g_k)\mapsto(1_G,g_0,\ldots,g_k)$
jest dodatkową degeneracją, a odpowiadający jej
stożek na realizacji daje homotopię identyczności
do mapy stałej. Przejście pomiędzy tym modelem
a nieskończonym złączeniem zachowuje typ homotopijny
(oba są standardowymi modelami uniwersalnej
numerowalnej wiązki); to wystarcza do wniosku.
\end{proof}

\begin{twierdzenie}[Klasyfikacja numerowalnych wiązek głównych]
\label{thm:BG-general-v6}
Niech $G$ będzie grupą topologiczną, $X$ przestrzenią
parazwartą. Klasy izomorfizmu numerowalnych głównych
wiązek $G$ nad $X$ pozostają w naturalnej bijekcji z
$[X,BG]$, gdzie $BG=EG/G$ dla uniwersalnego
numerowalnego modelu $EG\to BG$.
\end{twierdzenie}
\begin{proof}
Pokażmy najpierw, skąd bierze się mapa klasyfikująca.
Niech $P\to X$ ma pokrycie trywializujące $\{U_i\}$
z lokalnymi przekrojami $s_i$, a $\{\lambda_i\}$ będzie
podporządkowanym mu lokalnie skończonym podziałem jedności.
Dla $e\in P_x$ piszemy $e=s_i(x)g_i(e)$ na $U_i$.
Przypisanie
\[
\Phi(e)=\sum_i\lambda_i(x)g_i(e)\in EG
\]
jest ciągłe (lokalnie występuje skończenie wiele składników)
i $G$-ekwawariantne: $\Phi(eh)=\Phi(e)h$.
Przechodząc do ilorazu otrzymujemy mapę
$f:X\to BG$ oraz izomorfizm wiązek
$P\cong f^*(EG)$. To sprawdza się w każdej trywializacji:
współrzędna $g_i(e)$ przekształca się przez prawą
translację, więc na włóknach otrzymujemy bijekcję
$G$-przestrzeni zgodną z topologią.

Dwa podziały jedności i dwa wybory lokalnych przekrojów
prowadzą do homotopijnych map klasyfikujących:
umieszczamy wybrane konstrukcje w rozłącznych
zbiorach współrzędnych nieskończonego złączenia
i interpolujemy współczynniki przez $(1-t)$ oraz $t$.
Analogiczna interpolacja pokazuje, że izomorficzne
wiązki wyznaczają homotopijne mapy. Odwrotnie,
cofnięcia wiązki uniwersalnej wzdłuż map homotopijnych
są izomorficzne, ponieważ numerowalne wiązki nad
$X\times I$ mają izomorficzne restrykcje do końców
cylindra. Zbieżność konstrukcji nie wymaga nieskończonej
sumy w otoczeniu pojedynczego punktu, gdyż użyliśmy
lokalnie skończonego podziału jedności.
\end{proof}

\begin{center}
\begin{tikzpicture}[>=Stealth,node distance=1.7cm]
\node (P) {$P$};
\node (EG) [right=3.4cm of P] {$EG$};
\node (X) [below=1.4cm of P] {$X$};
\node (BG) [below=1.4cm of EG] {$BG$};
\draw[->,thick] (P)--(EG) node[midway,above]{$\Phi$};
\draw[->,thick] (P)--(X) node[midway,left]{$\pi$};
\draw[->,thick] (EG)--(BG) node[midway,right]{$\pi_{\rm univ}$};
\draw[->,thick] (X)--(BG) node[midway,below]{$f$};
\end{tikzpicture}
\end{center}
Powyższy diagram jest przemienny, a górna mapa
$\Phi$ jest $G$-ekwawariantna. Taką własność ma
mapa klasyfikująca skonstruowana w dowodzie.

\begin{wniosek}[Klasyfikacja wiązek wektorowych]
Niech $X$ będzie parazwartą przestrzenią, $r\ge1$.
Izomorfizmy rzeczywistych wiązek wektorowych rzędu $r$
nad $X$ odpowiadają klasom homotopii map
$X\to BO(r)$, a zespolonych rzędu $r$ -- klasom map
$X\to BU(r)$. Standardowe modele mają postać
granic przestrzeni Grassmanna skończeniewymiarowych.
\end{wniosek}
\begin{proof}
Metryka włóknowa na rzeczywistej wiązce redukuje
grupę strukturalną do $O(r)$. Jej wiązka ram
ortonormalnych jest główną wiązką $O(r)$,
a wiązkę pierwotną odzyskujemy przez dołączenie
reprezentacji standardowej na $\R^r$.
Stosujemy Twierdzenie~\ref{thm:BG-general-v6}.
W przypadku zespolonym wybieramy metrykę hermitowską
i zastępujemy $O(r)$ przez $U(r)$.
\end{proof}

\begin{przyklad}[Zespolone wiązki liniowe i $\C P^\infty$]
Grupa $U(1)=S^1$ działa swobodnie na
$S^\infty=\bigcup_{N\ge0}S^{2N+1}$, który jest
ściągalny. Iloraz to $\C P^\infty$, zatem
$BU(1)\simeq\C P^\infty$. Z długiego ciągu homotopii
fibracji $S^1\to S^\infty\to\C P^\infty$ otrzymujemy
$\pi_2(\C P^\infty)\cong\Z$ i
$\pi_j(\C P^\infty)=0$ dla $j\ne2$, $j\ge1$.
Stąd $\C P^\infty=K(\Z,2)$ i dla CW-kompleksu $X$
\[
\{\text{zespolone wiązki liniowe nad }X\}/\cong
\quad\longleftrightarrow\quad H^2(X;\Z).
\]
Odpowiadająca klasie kohomologii liczba lub klasa to
pierwsza klasa Cherna $c_1(L)$. W szczególności
nad $S^2$ wiązki liniowe klasyfikuje liczba całkowita
$\langle c_1(L),[S^2]\rangle$.
\end{przyklad}

\section{Teoria przeszkód: rozszerzanie map i przekrojów}
\begin{definicja}[Przeszkoda rozszerzania]
Niech $(X,A)$ będzie parą CW, a $Y$ przestrzenią
łukowo spójną. Załóżmy, że mapa $f$ została określona
na $A\cup X^n$. Dla $(n+1)$-komórki
$e_\alpha^{n+1}$ z mapą charakterystyczną
$\chi_\alpha:D^{n+1}\to X$ ograniczenie
$f\circ\chi_\alpha|_{S^n}$ reprezentuje klasę
$[f\circ\chi_\alpha|_{S^n}]\in\pi_n(Y)$,
po przetransportowaniu punktu bazowego.
Rodzinę tych klas nazywamy \emph{kołańcuchem przeszkody}
rozszerzenia na komórki stopnia $n+1$.
Dla $n\ge2$ jest to kołańcuch o wartościach
w systemie lokalnym otrzymanym z działania
$\pi_1(Y)$ na $\pi_n(Y)$, cofniętym przez
$f_*:\pi_1(X)\to\pi_1(Y)$ (mapa $f$ jest już
określona na $2$-szkielecie). Transporty wymagają
ustalenia dróg łączących punkty bazowe. Jeżeli
to działanie jest trywialne, współczynniki są stałe.
Dla pary $(X,A)$ w rachunku względnym uwzględniamy
te same transporty w $X$.
\end{definicja}

\begin{center}
\begin{tikzpicture}[>=Stealth]
 \draw[thick,fill=blue!4] (0,0) circle (.88);
 \draw[very thick,red!70!black] (0,0) circle (.88);
 \node at (0,0){$D^{n+1}$};
 \node[red!70!black,below] at (0,-.95){$S^n$};
 \draw[->,thick] (1.2,.15)--(3.45,.15)
   node[midway,above]{$f\circ\chi|_{S^n}$};
 \draw[thick,fill=blue!4] (4.4,0) ellipse (.7 and .85);
 \node at (4.4,0){$Y$};
 \node[align=center] at (2.3,-1.78){\footnotesize Przedłużenie na dysk istnieje dokładnie wtedy,\\
 \footnotesize gdy klasa mapy na brzegu w $\pi_n(Y)$ jest zerowa.};
\end{tikzpicture}
\end{center}

\begin{twierdzenie}[Pierwsza przeszkoda dla sekcji]
\label{thm:obstruction-section-v6}
Niech $p:E\to B$ będzie lokalnie trywialną wiązką nad
spójnym CW-kompleksem, której włókno $F$ jest
$(n-1)$-spójne dla pewnego $n\ge2$.
Przypuśćmy, że skonstruowano przekrój $s$ nad
$n$-szkieletem $B^n$. Dołączenie każdej $(n+1)$-komórki
wyznacza element $\pi_n(F)$, a elementy te tworzą
kocykl komórkowy $o_{n+1}(s)$ o wartościach
w systemie lokalnym $\mathcal\pi_n(F)$.
Przy ustalonym przekroju nad $(n-1)$-szkieletem
jego klasa w $H^{n+1}(B;\mathcal\pi_n(F))$ nie zależy
od jego rozszerzenia na $n$-szkielet. Ponieważ
włókno jest $(n-1)$-spójne, dwa przekroje nad
$(n-1)$-szkieletem można połączyć odpowiednią
homotopią; wynikająca stąd klasa jest więc
niezależna od wyborów (po identyfikacji systemów).
Nazywa się ją \emph{pierwszą przeszkodą} istnienia przekroju. Znikanie tej klasy pozwala
skonstruować przekrój na $B^{n+1}$ (po ewentualnej
zmianie przekroju na $B^n$), lecz nie zapewnia
jeszcze jego istnienia na całym $B$, jeżeli
wyższe grupy $\pi_j(F)$ są niezerowe.
\end{twierdzenie}
\begin{proof}
Cofnięcie wiązki na dysk charakterystyczny każdej
$(n+1)$-komórki jest trywialne. Przekrój już istniejący
na jej brzegu wyznacza więc mapę $S^n\to F$.
Można ją przedłużyć na dysk dokładnie wtedy, gdy
jej klasa w $\pi_n(F)$ znika.
Dla różnych wyborów trywializacji identyfikacje
$\pi_n(F_b)$ różnią się transportem w systemie lokalnym;
zatem wszystkie te klasy są współczynnikami tego samego
kołańcucha komórkowego $o_{n+1}(s)$.

Pokażmy dokładniej, skąd pochodzi warunek kocyklu.
Trzeba uwzględnić nie tylko liczby incydencji,
lecz również działanie monodromii na $\pi_n(F)$.
Niech $e^{n+2}$ będzie komórką z mapą
charakterystyczną $\chi:D^{n+2}\to B$.
Cofnięta wiązka nad dyskiem jest trywialna.
Na podziale komórkowym $\partial D^{n+2}$ przekrój
wyznaczony na $n$-szkielecie daje elementy
$\pi_n(F)$ na granicach $(n+1)$-komórek.
Każdą z nich transportujemy do wspólnego włókna
nad wybranym wierzchołkiem $\partial D^{n+2}$.
Suma ze znakami incydencji jest dokładnie wartością
$\delta o_{n+1}(s)$ na $e^{n+2}$.

Aby wykazać jej znikanie, używamy relatywnego
kompleksu homotopii pary
$(D^{n+2},\partial D^{n+2})$ oraz naturalności
homomorfizmów brzegowych dla map charakterystycznych.
Element odpowiadający obstrukcji jest obrazem
brzegu odpowiedniej klasy względnej, a następne
zastosowanie homomorfizmu brzegowego daje zero
($\partial^2=0$ w łańcuchowym opisie komórkowym
tej pary). W lokalnej trywializacji wszystkie
transporty są identyczne, a po powrocie do $B$
odzyskujemy je przez powyższą identyfikację systemu
lokalnego. Stąd $\delta o_{n+1}(s)=0$.
Jest to \emph{komórkowy lemat o kocyklu przeszkody};
samo przytoczenie $\partial^2=0$ bez opisania
transportów byłoby niewystarczające.

Jeżeli zmienimy przekrój na komórkach wymiaru $n$,
zachowując go na $(n-1)$-szkielecie, różnice
opisują klasy map sferowych w $\pi_n(F)$; współczynniki
tej zmiany tworzą $n$-kołańcuch $b$, a nowa przeszkoda jest
$o_{n+1}(s)+\delta b$. Dlatego klasa kohomologii
nie zależy od takiej zmiany. Gdy $[o_{n+1}(s)]=0$,
wybieramy $b$ z $\delta b=-o_{n+1}(s)$,
modyfikujemy przekrój na $n$-komórkach i otrzymujemy
zerowe przeszkody dla wszystkich $(n+1)$-komórek.
Podnoszenie homotopii i rozszerzanie komórka po komórce
dają przekrój na $B^{n+1}$. Na dalszych szkieletach
problem się powtarza dla kolejnych $\pi_j(F)$.
\end{proof}

\begin{uwaga}[Zakres hipotez]
Gdy $n=1$ grupa $\pi_1(F)$ może nie być abelowa i
pierwsza przeszkoda wymaga sformułowania w teorii
nieabelowych kocykli. Dla $n\ge2$ odpowiednie grupy
są abelowe. Dla fibracji niebędących wiązkami lokalnie
trywialnymi potrzebna jest wersja komórkowej teorii
przeszkód z podnoszeniem homotopii; nie wyprowadzamy
jej tu przez pozorną lokalną trywializację.
\end{uwaga}

\begin{wniosek}[Klasa Eulera i przekrój nigdzie niezerowy]
Niech $\xi\to B$ będzie orientowaną rzeczywistą wiązką
wektorową rzędu $r\ge2$ nad CW-kompleksem $B$.
Pierwsza przeszkoda skonstruowania przekroju jednostkowej
wiązki sfer $S(\xi)\to B$ leży w $H^r(B;\Z)$ i
jest równa klasie Eulera $e(\xi)$, z ustaloną konwencją
orientacji. W szczególności $e(\xi)\ne0$ wyklucza
przekrój nigdzie niezerowy.
\end{wniosek}
\begin{proof}
Włókno $S^{r-1}$ jest $(r-2)$-spójne i
$\pi_{r-1}(S^{r-1})\cong\Z$. Dla $r\ge3$
możemy zastosować Twierdzenie~\ref{thm:obstruction-section-v6}
z $n=r-1\ge2$. Przypadek $r=2$ rozpatrujemy
osobno: zorientowana wiązka płaszczyzn redukuje
grupę strukturalną do $SO(2)\cong U(1)$,
a jej przekrój jednostkowy jest równoważny
trywializacji tej wiązki. Jej pierwszą przeszkodą
jest klasa $c_1=e\in H^2(B;\Z)$, zgodnie z klasyfikacją
wiązki liniowej zespolonej przez $BU(1)$.
Dla $r\ge3$ przeszkoda na $r$-komórce jest
stopniem mapy jej brzegu do sfery włókna. W konstrukcji Gysina ten sam stopień
otrzymujemy, ograniczając klasę Thoma do dysku
charakterystycznego i przechodząc do zwykłej
kohomologii; naturalność homomorfizmu łączącego
identyfikuje więc kołańcuch przeszkody z reprezentantem
klasy Eulera. Brak przekroju jest wówczas konieczny.
Nie twierdzimy odwrotności bez dodatkowych hipotez:
po zniknięciu klasy Eulera mogą istnieć wyższe przeszkody.
\end{proof}

\begin{przyklad}[Włosy na sferze]
Dla sfery $S^{2m}$ mamy
$\langle e(TS^{2m}),[S^{2m}]\rangle
=\chi(S^{2m})=2$ na mocy twierdzenia
Poincarégo--Hopfa. Ponieważ klasa Eulera jest niezerowa,
żadna parzystowymiarowa sfera nie ma globalnego
nigdzie niezerowego pola stycznego. Dla $S^{2m+1}$
charakterystyka Eulera znika; w tym przypadku pole
nigdzie niezerowe można dodatkowo skonstruować
jawnie, stosując na $\R^{2m+2}\cong\C^{m+1}$
wektor $x\mapsto ix$, styczny do sfery i niezerowy.
\end{przyklad}

\begin{przyklad}[Znikanie klasy Eulera nie wystarcza]
Istnieje orientowana rzeczywista wiązka wektorowa rzędu
$3$ nad $S^4$, która nie ma przekroju nigdzie niezerowego,
chociaż jej klasa Eulera jest równa zeru.
Istotnie, grupa $\pi_3(SO(3))\cong\Z$ jest niezerowa:
przykładowo dwukrotne nakrycie $S^3\to SO(3)$
przez jednostkowe kwaterniony indukuje izomorfizm
w trzeciej grupie homotopii. Wybierzmy wiązkę $\xi$
sklejoną na równiku $S^3\subset S^4$ funkcją
reprezentującą niezerowy element tej grupy.
Z klasyfikacji przez funkcje sklejające wynika,
że $\xi$ jest nietrywialna. Tymczasem
$e(\xi)\in H^3(S^4;\Z)=0$.
Gdyby $\xi$ miała przekrój nigdzie niezerowy,
rozkładałaby się jako
$\xi\cong\underline\R\oplus\eta$,
gdzie $\eta$ jest orientowaną rzeczywistą wiązką
rzędu $2$. Ale $\eta$ musi być trywialna, gdyż
$\pi_3(SO(2))=\pi_3(S^1)=0$, a więc wszystkie
takie wiązki nad $S^4$ są trywialne.
Otrzymalibyśmy sprzeczność. Jest to konkretny przykład,
w którym \emph{pierwsza} przeszkoda znika,
lecz przeszkoda wyższego rzędu pozostaje.
\end{przyklad}

\section{Funkcje sklejające i wiązki nad sferami}
\begin{twierdzenie}[Konstrukcja przez równik]
Niech $G$ będzie spójną łukowo grupą topologiczną,
a $n\ge2$. Główne wiązki $G$ nad $S^n$ powstają
przez sklejenie dwóch trywialnych wiązek nad
półkulami za pomocą funkcji
$g:S^{n-1}\to G$ na równiku. Klasy takich wiązek
odpowiadają elementom $\pi_{n-1}(G)$, z naturalną
konwencją dotyczącą punktów bazowych; dla niepunktowanych
funkcji sklejających odpowiednią relację homotopii
należy uwzględnić w klasyfikacji.
\end{twierdzenie}
\begin{proof}
Półkule $D^n_\pm$ są ściągalne, więc każda główna wiązka
nad nimi jest trywialna. Dwie trywializacje różnią się
na przekroju równikowym przez mapę do $G$.
Sklejając $(D^n_+\times G)$ z $(D^n_-\times G)$
według relacji $(x,h)_+\sim(x,g(x)h)_-$ otrzymujemy
wiązkę. Zmiana trywializacji nad półkulą działa na
g przez mnożenie z lewej lub prawej strony funkcją,
która rozszerza się na dysk, a zatem jest homotopijna
do stałej. Gdy $G$ jest spójna łukowo, odpowiednie
stałe są połączone drogami i otrzymujemy klasyfikację
przez homotopię funkcji z równika do $G$.
Dla $n\ge2$ i spójnej łukowo grupy $G$ jest to
$\pi_{n-1}(G)$ (w szczególności dla $n=2$ jest to
abelowa grupa $\pi_1(G)$).
\end{proof}
\begin{przyklad}
Dla $G=U(1)=S^1$ oraz $n=2$ funkcja
$g_k(e^{it})=e^{ikt}$ ma stopień $k$.
Sklejona zespolona wiązka liniowa ma
$\langle c_1(L_k),[S^2]\rangle=k$ z przyjętą
orientacją i konwencją klejenia. Dla $k=0$ jest trywialna,
a dla $k=\pm1$ otrzymujemy tautologiczną wiązkę
nad $\C P^1$ albo jej dualną (zależnie od znaku).
\end{przyklad}


\section{Uzupełnienia audytu: konstrukcje i granice hipotez}
\label{sec:audyt-VI-VII}
\begin{lemat}[Komórki przy konstrukcji $K(G,n)$]
Niech $A$ będzie CW-kompleksem i niech $Y=A\cup_f e^{m}$
dla $m\ge3$. Dla każdego $j\le m-2$ włączenie
$A\hookrightarrow Y$ indukuje izomorfizm
$\pi_j(A)\to\pi_j(Y)$. Dla $j=m-1$ indukuje
surjekcję, której jądro jest normalnie generowane
przez klasę mapy dołączającej (w wymiarach $m-1\ge2$
z naturalnym uwzględnieniem działania $\pi_1$).
\end{lemat}
\begin{proof}
Para $(Y,A)$ jest, przez przybliżenie komórkowe
względem $A$, $(m-1)$-spójna: klasy względne w stopniach
$m-1$ i niższych można reprezentować mapami dysków,
których obraz po przybliżeniu komórkowym unika wnętrza
jedynej nowej $m$-komórki. W długim ciągu
$\pi_j(A)\to\pi_j(Y)\to\pi_j(Y,A)$
otrzymujemy zatem izomorfizmy do wymiaru $m-2$.
W stopniu $m-1$ jądro mapy włączenia jest obrazem
$\partial:\pi_m(Y,A)\to\pi_{m-1}(A)$,
a relatywny generator ma jako brzeg klasę mapy $f$.
Dla kilku komórek stosujemy argument kolejno,
a następnie przechodzimy do CW-granicy; zwartość
sfer gwarantuje faktoryzację map przez pewien etap.
\end{proof}

\begin{uwaga}[Co jest rzeczywiście udowodnione]
Lemat wyjaśnia konstrukcję istnienia $K(G,n)$.
W klasyfikacji \emph{wszystkich} wiązek głównych
występuje oddzielne twierdzenie o uniwersalnym
numerowalnym modelu $EG\to BG$: jego użycie wymaga
parazwartości i własności Hausdorffa podstawy
oraz kontrolowania lokalnie skończonych podziałów
jedności. Wzór na mapę klasyfikującą pokazuje
surjektywność klasyfikacji i jej naturalność;
pełna homotopijna jednoznaczność zawiera dodatkowo
argument o wiązkach nad cylindrem $X\times I$.
Nie zastępujemy tego argumentu samym stwierdzeniem,
że dwie mapy są punktowo podobne.
\end{uwaga}

\begin{przyklad}[Rozszerzenie na komórkę w dwóch etapach]
Rozważmy przestrzeń Moore'a
$M(\Z/m,n)=S^n\cup_{m}e^{n+1}$,
$m\ge2$, $n\ge2$. Mapa $S^n\to K(A,n)$
reprezentująca $a\in A$ przedłuża się na
$M(\Z/m,n)$ dokładnie wtedy, gdy $ma=0$.
Istotnie, po złożeniu z mapą dołączającą stopnia $m$
otrzymujemy $ma\in\pi_n(K(A,n))=A$, a przedłużenie
istnieje dokładnie dla klasy zerowej. Wynik jest zgodny
z $H^n(M(\Z/m,n);A)\cong A[m]$.
\end{przyklad}

\section{Podsumowanie i granice obecnej iteracji}
Uzyskaliśmy trzy mechanizmy służące różnym celom:
$K(A,n)$ reprezentuje kohomologię,
$BG$ klasyfikuje wiązki, a kołańcuchy przeszkód
wykrywają możliwość rozszerzania map i przekrojów.
Użyte konstrukcje są zgodne z poprzednim rozdziałem:
przestrzenie klasyfikujące i klasy Eulera można badać
ciągami spektralnymi fibracji, a pierwsza przeszkoda
jest obliczana w kohomologii z lokalnymi współczynnikami.

Wieże Postnikowa, klasy $k$-niezmienników,
operacje Steenroda i szczegółowa teoria przeszkód
z wieloma niezerowymi grupami homotopii włókna
stanowią naturalny następny etap. Nie utożsamiamy
pierwszej przeszkody z całą teorią przeszkód.


% Rozdzial glowny iteracji VII. Do wstawienia przed \appendix.
\chapter{Wieże Postnikowa, iloczyny Whiteheada i operacje Steenroda}
\label{ch:postnikow-v7}

\section{Dlaczego grupy homotopii nie wystarczają?}
Niech $X$ będzie łukowo spójnym CW-kompleksem. Grupy
$\pi_n(X)$ opisują odwzorowania sfer, ale zapisanie samej
rodziny tych grup nie mówi jeszcze, jak odwzorowania
różnych wymiarów oddziałują ze sobą. Właśnie tę brakującą
informację zachowują \emph{niezmienniki Postnikowa}. Z kolei
\emph{iloczyny Whiteheada} mierzą oddziaływanie elementów
grup homotopii, a \emph{operacje Steenroda} wykrywają
informację niewidoczną w addytywnych grupach kohomologii
czy nawet w samych kwadratach iloczynu kubkowego.

Podamy najpierw konstrukcje dla \emph{jednospójnych}
CW-kompleksów z wybranym punktem bazowym. Jednospójność
oznacza tu $\pi_0(X)=*$ i $\pi_1(X)=0$. Dzięki temu
na wyższych grupach homotopii nie występuje nietrywialne
działanie grupy podstawowej. Dopiero w ostatniej sekcji
wyjaśnimy, co zmienia się w przypadku ogólnym.

\begin{definicja}
Niech $n\ge1$. \emph{$n$-typem homotopijnym} nazywamy
łukowo spójny CW-kompleks $Y$, dla którego $\pi_j(Y)=0$
dla każdego $j>n$. Mapa $f:X\to Y$ jest
\emph{$n$-równoważnością} w tym rozdziale, jeśli
$f_*:\pi_j(X)\to\pi_j(Y)$ jest izomorfizmem dla
$1\le j\le n$ (przy zgodnym wyborze punktów bazowych).
\emph{$n$-etapem Postnikowa} przestrzeni $X$ nazywamy
parę $(P_nX,p_n)$, gdzie $P_nX$ jest $n$-typem,
a $p_n:X\to P_nX$ jest $n$-równoważnością.
\end{definicja}

\begin{uwaga}
Zdefiniowane tu pojęcie $n$-równoważności jest słabsze
od spotykanego w literaturze określenia \emph{$n$-spójna mapa},
które dodatkowo narzuca warunek surjektywności w wymiarze
$n+1$ lub warunek na włókno homotopijne, zależnie od
konwencji. Nie będziemy mieszać tych terminów.
\end{uwaga}

\begin{twierdzenie}[Istnienie etapów Postnikowa]
\label{thm:postnikow-istence-v7}
Dla każdego łukowo spójnego CW-kompleksu $X$ i każdej
liczby $n\ge1$ istnieje etap $(P_nX,p_n)$.
Można dobrać etapy w wieżę odwzorowań
$\cdots\to P_{n+1}X\to P_nX\to\cdots\to P_1X$
zgodnych z odwzorowaniami $p_n$ w kategorii homotopijnej.
\end{twierdzenie}
\begin{proof}
Do $X$ dołączamy komórki wymiaru co najmniej $n+2$.
Na początku wybieramy reprezentantów wszystkich elementów
$\pi_{n+1}(X)$ i dołączamy do ich map
$S^{n+1}\to X$ komórki wymiaru $n+2$.
Długi ciąg homotopii pary oraz względne twierdzenie
o przybliżeniu komórkowym pokazują, że nowe komórki
nie zmieniają $\pi_j$ dla $j\le n$, natomiast zabijają
wybrane elementy $\pi_{n+1}$. Jeżeli na tym etapie
pozostałyby elementy $\pi_{n+1}$, dołączamy komórki
wzdłuż ich reprezentantów; równoważnie można od razu
zabrać reprezentantów całej grupy. Następnie zabijamy
$\pi_{n+2}$ komórkami wymiaru $n+3$, i tak dalej.

Niech $Y$ będzie sumą kolejnych podkompleksów w topologii
CW. Dla ustalonego $j>n$ wszystkie elementy $\pi_j$
któregoś etapu giną po dołączeniu komórek $(j+1)$-wymiarowych,
a późniejsze komórki o wymiarze co najmniej $j+2$
nie odtwarzają $\pi_j$. Każda mapa $S^j\to Y$ jest
homotopijna do mapy trafiającej do skończonego podkompleksu
$Y$, a więc do pewnego etapu konstrukcji (twierdzenie
o zwartościowym charakterze map do CW-kompleksu).
Otrzymujemy $\pi_j(Y)=0$. Ponieważ dołączaliśmy tylko
komórki wymiaru co najmniej $n+2$, dla $j\le n$
mapa $X\to Y$ indukuje izomorfizm na $\pi_j$.
Przyjmujemy $P_nX=Y$.

Aby uzyskać mapy $P_{n+1}X\to P_nX$, przedłużamy
$p_n:X\to P_nX$ do $P_{n+1}X$ komórka po komórce.
Komórki dołączone do $X$ w konstrukcji $P_{n+1}X$
mają wymiary co najmniej $n+3$, a zatem przeszkody
przedłużenia leżą w grupach $\pi_j(P_nX)$ z $j\ge n+2$,
które są zerowe. Złożenie z $p_{n+1}$ jest homotopijne
z $p_n$. Wybory przedłużeń nie są kanoniczne na poziomie
pojedynczych przestrzeni, ale wyznaczają poprawną wieżę
w kategorii homotopijnej.
\end{proof}

\begin{przyklad}[Dwa pierwsze etapy sfery]
Dla $X=S^2$ mamy $P_1S^2\simeq *$, a
$P_2S^2\simeq K(\Z,2)\simeq\C P^\infty$.
Nie wynika stąd, że $S^2\simeq\C P^\infty$:
pierwsza przestrzeń ma niezerową grupę $\pi_3\cong\Z$,
a druga już nie. Konstrukcja $P_3S^2$ musi więc
zawierać informację, której nie ma w samym $K(\Z,2)$.
\end{przyklad}

\section{Włókna kolejnych etapów i niezmienniki \texorpdfstring{$k$}{k}}
\begin{definicja}[Włókno homotopijne]
Niech $f:B\to Y$ będzie mapą przestrzeni z punktami
bazowymi i $y_0\in Y$. \emph{Włóknem homotopijnym} $f$
nad $y_0$ nazywamy przestrzeń
\[
\operatorname{hofib}(f)=\{(b,\gamma):b\in B,
\ \gamma:[0,1]\to Y,\ \gamma(0)=f(b),\ \gamma(1)=y_0\},
\]
z topologią podprzestrzeni $B\times Y^{[0,1]}$,
gdzie $Y^{[0,1]}$ ma topologię zwarto-otwartą.
\end{definicja}

\begin{lemat}[Włókno mapy do przestrzeni Eilenberga--Mac Lane'a]
\label{lem:hofib-K-v7}
Dla $n\ge1$ i grupy abelowej $A$ włókno homotopijne
odwzorowania $k:B\to K(A,n+1)$ jest włóknem fibracji
nad $B$ o typie homotopijnym $K(A,n)$.
Dla spójnej $B$ działanie monodromii na $A$ jest trywialne.
\end{lemat}
\begin{proof}
Rzutowanie $\operatorname{hofib}(k)\to B$ jest cofnięciem
fibracji dróg $P_{y_0}K(A,n+1)\to K(A,n+1)$,
gdzie $P_{y_0}Y$ składa się z dróg kończących się w $y_0$.
Włókno jest przestrzenią pętli $\Omega K(A,n+1)$.
Z $\pi_j(\Omega Y)=\pi_{j+1}(Y)$ wynika, że ma ona
jedyną niezerową grupę homotopii $\pi_n=A$,
więc jest słabo równoważna $K(A,n)$; po wybraniu
modeli CW daje to równoważność homotopijną.
Działanie $\pi_1(B)$ otrzymujemy przez cofnięcie
działania $\pi_1(K(A,n+1))$, a ta grupa jest zerowa
dla $n\ge1$.
\end{proof}

\begin{twierdzenie}[Rozkład Postnikowa dla przestrzeni jednospójnej]
\label{thm:postnikov-k-v7}
Niech $X$ będzie jednospójnym CW-kompleksem, $n\ge2$,
a $A_n=\pi_n(X)$ (grupa abelowa). Etapy Postnikowa
można dobrać wraz z fibracjami homotopijnymi
\[
K(A_n,n)\longrightarrow P_nX\xrightarrow{q_n}P_{n-1}X.
\]
Fibracja ta jest, z dokładnością do równoważności
włóknowej, włóknem homotopijnym pewnej mapy
\(k_n:P_{n-1}X\to K(A_n,n+1)\).
Klasa $[k_n]\in H^{n+1}(P_{n-1}X;A_n)$
nazywa się \emph{$n$-tym niezmiennikiem Postnikowa}.
Zależy od wyboru identyfikacji grup homotopii,
a bez tych identyfikacji jest określona z dokładnością
do działania odpowiednich automorfizmów.
\end{twierdzenie}
\begin{proof}
Rozpatrzmy mapę $q_n:P_nX\to P_{n-1}X$,
odpowiadającą zgodnym etapom. Zastępujemy ją fibracją
przestrzeni dróg, nie zmieniając typu homotopijnego
przestrzeni całkowitej. Z długiego ciągu homotopii
wynika, że jej włókno homotopijne $F_n$ ma
$\pi_j(F_n)=0$ dla $j\ne n$ oraz
$\pi_n(F_n)\cong\pi_n(X)=A_n$: dla $j<n$
mapa $q_n$ indukuje izomorfizm, a dla $j>n$
obie grupy w wieży są zerowe. Zatem
$F_n\simeq K(A_n,n)$.

Wyjaśnienia wymaga, dlaczego taka fibracja jest określona
przez jedną klasę kohomologii. Na kolejnych komórkach
jednospójnej podstawy $B=P_{n-1}X$ wybieramy lokalne
podniesienia do $E=P_nX$. Ponieważ włókno
$K(A_n,n)$ ma zerowe grupy $\pi_j$ dla $j<n$,
wybory można uzgodnić do wymiaru $n$. Na
$(n+1)$-komórkach pojawiają się elementy $A_n$;
jak w teorii przeszkód tworzą one kocykl i klasę
$k\in H^{n+1}(B;A_n)$. Wyższe przeszkody
na zgodność wyborów nie występują, ponieważ
$\pi_j(K(A_n,n))=0$ dla $j>n$.

Jest to podstawowy przypadek klasyfikacji \emph{głównych
fibracji $K(A,n)$}. Można go również zobaczyć
bezpośrednio: wybieramy mapę $k:B\to K(A_n,n+1)$
reprezentującą klasę i konstruujemy jej włókno
homotopijne zgodnie z Lematem~\ref{lem:hofib-K-v7}.
Na każdym szkielecie obie fibracje mają ten sam
kołańcuch przeszkody, a indukcja komórkowa daje
ich równoważność włóknową. Ważna jest tu \emph{prostota}
struktury homotopijnej: $\pi_1(B)=0$ i dlatego
nie występują skręcone współczynniki ani niezależne
działania monodromii na $A_n$.
\end{proof}

\begin{uwaga}[Granica powyższego dowodu]
Identyfikacja włókna przez długi ciąg homotopii jest
pełnym rachunkiem. Natomiast klasyfikacja wszystkich
głównych fibracji $K(A,n)$ przez $H^{n+1}$ korzysta
z twierdzenia o klasyfikacji takich fibracji, którego
konstrukcję komórkową podaliśmy w argumentacji.
Nie twierdzimy, że \emph{dowolna} fibracja o włóknie
$K(A,n)$ jest główna: bez dodatkowych danych może
wystąpić nietrywialne działanie grupy automorfizmów
włókna. Wieża Postnikowa przestrzeni jednospójnej
ma wymagany typ główny.
\end{uwaga}

\begin{center}
\begin{tikzpicture}[>=Stealth, every node/.style={font=\small}]
\node (X) at (-3.6,1.8) {$X$};
\node (P4) at (-1.5,1.8) {$P_4X$};
\node (P3) at (0.7,1.8) {$P_3X$};
\node (P2) at (3.0,1.8) {$P_2X$};
\draw[->,thick] (X)--(P4) node[midway,above]{$p_4$};
\draw[->,thick] (P4)--(P3) node[midway,above]{$q_4$};
\draw[->,thick] (P3)--(P2) node[midway,above]{$q_3$};
\node (K4) at (-1.5,.35) {$K(\pi_4X,4)$};
\node (K3) at (.7,.35) {$K(\pi_3X,3)$};
\draw[->,thick] (K4)--(P4);
\draw[->,thick] (K3)--(P3);
\node (K5) at (-1.5,-1.0) {$K(\pi_4X,5)$};
\node (K4b) at (.7,-1.0) {$K(\pi_3X,4)$};
\draw[->,thick,dashed] (P3) to[out=-125,in=10] (K5)
 node[midway,left] {$k_4$};
\draw[->,thick,dashed] (P2) to[out=-130,in=10] (K4b)
 node[midway,right] {$k_3$};
\end{tikzpicture}
\end{center}
\noindent Schemat wieży: włókno mapy $q_n$ ma typ $K(\pi_nX,n)$,
a klasa $k_n$ jest reprezentowana przez mapę z $P_{n-1}X$
do $K(\pi_nX,n+1)$. Przerywane strzałki oznaczają
mapy klasyfikujące, a nie części ciągu fibracji.

\section{Niezmiennik \texorpdfstring{$k_3$}{k3} dla \texorpdfstring{$S^2$}{S2}}
\begin{twierdzenie}[Pierwszy nietrywialny etap sfery]
\label{thm:s2-k-v7}
Niech $u\in H^2(K(\Z,2);\Z)$ będzie generatorem,
wybranym tak, że $H^*(K(\Z,2);\Z)\cong\Z[u]$.
Dla $X=S^2$ niezmiennik $k_3$ odpowiada klasie
$\pm u^2\in H^4(K(\Z,2);\Z)\cong\Z$.
W szczególności $P_3S^2$ jest włóknem homotopijnym
mapy reprezentującej $\pm u^2$.
\end{twierdzenie}
\begin{proof}
Mamy $\pi_2(S^2)=\Z$ i $\pi_3(S^2)=\Z$,
stąd $P_2S^2\simeq K(\Z,2)$, zaś włókno
$P_3S^2\to P_2S^2$ ma typ $K(\Z,3)$.
Kohomologiczny ciąg spektralny Serre'a tej fibracji
z całkowitymi współczynnikami ma w niskich stopniach
$E_2^{p,q}=H^p(K(\Z,2);H^q(K(\Z,3);\Z))$.
Dla $p+q\le4$ wiersze $q=0$ i $q=3$
są jedynymi mogącymi uczestniczyć w różniczce
zaczynającej się w $E_4^{0,3}=\Z$;
jej celem jest $E_4^{4,0}=\Z\cdot u^2$.

Wybierzmy model $P_3S^2$ skonstruowany z $S^2$
przez dołączanie komórek wymiaru co najmniej $5$.
Nie ma wtedy nowych względnych komórek w stopniu $4$,
więc
$H^4(P_3S^2,S^2;\Z)=0$ i
$H^5(P_3S^2,S^2;\Z)$ może być niezerowa.
Z długiego ciągu pary wynika, że
$H^4(P_3S^2;\Z)=0$, ponieważ
$H^4(S^2;\Z)=0$ i $H^4(P_3S^2,S^2;\Z)=0$.
W konsekwencji klasa $u^2$ musi zginąć przed
stroną $E_\infty$, a jedyną różniczką mogącą
ją zabić jest $d_4:E_4^{0,3}\to E_4^{4,0}$.
Jej obraz musi być całym $\Z$, zatem jest to
mnożenie przez $\pm1$. Z definicji transgresji
klasa $k_3$ jest obrazem generatora włókna
przez tę różniczkę, co daje $k_3=\pm u^2$.
\end{proof}

\begin{uwaga}
W dowodzie wystarczy silniejszy fakt komórkowy:
$P_3S^2$ uzyskujemy z $S^2$ przez dołączanie komórek
w wymiarach co najmniej $5$. Nie wynika z tego
izomorfizm $H^5$, ani ogólne izomorfizmy wyższej
kohomologii. Informacja o \emph{pierwszej} niezerowej
różniczce w ciągu Serre'a jest tym, co rozróżnia
$S^2$ od $K(\Z,2)$.
\end{uwaga}


\begin{przyklad}[\texorpdfstring{$S^2$}{S2} i
\texorpdfstring{$\C P^2$}{CP2}: jednakowe $\pi_2$,
ale różne $k_3$]
Obie przestrzenie są jednospójne i mają
$\pi_2\cong\Z$, więc ich drugi etap Postnikowa
to $K(\Z,2)$. Natomiast z fibracji Hopfa
$S^1\to S^5\to\C P^2$ wynika
$\pi_3(\C P^2)\cong\pi_3(S^5)=0$.
Zatem $P_3\C P^2\simeq K(\Z,2)$, a odpowiadający
niezmiennik $k_3$ znika, inaczej niż dla $S^2$.
Same drugie grupy homotopii nie odróżniają więc tych
przestrzeni; pierwszy niezerowy niezmiennik Postnikowa
dostarcza dodatkowej informacji.
\end{przyklad}

\section{Iloczyny Whiteheada i komórka górna iloczynu sfer}
\begin{definicja}[Iloczyn Whiteheada]
Niech $X$ będzie jednospójną przestrzenią z punktem
bazowym, $p,q\ge2$, a
$\alpha\in\pi_p(X)$, $\beta\in\pi_q(X)$.
Oznaczmy przez $\iota_p:S^p\to S^p\vee S^q$ i
$\iota_q:S^q\to S^p\vee S^q$ włączenia składników.
Produkt $S^p\times S^q$ ma strukturę CW, w której
$(S^p\times\{*\})\cup(\{*\}\times S^q)=S^p\vee S^q$
jest szkieletem bez ostatniej komórki wymiaru $p+q$.
Niech $\omega:S^{p+q-1}\to S^p\vee S^q$
będzie mapą dołączającą tej komórki, z określoną
konwencją orientacji. \emph{Iloczynem Whiteheada}
nazywamy
$[\alpha,\beta]=(\alpha\vee\beta)_*[\omega]
\in\pi_{p+q-1}(X)$, gdzie
$\alpha\vee\beta:S^p\vee S^q\to X$.
\end{definicja}

\begin{uwaga}
Zmiana konwencji orientacji może zmienić globalny znak
iloczynu. W dalszych rachunkach, w których znak ma
znaczenie, podajemy wynik z dokładnością do znaku.
Nie mylimy iloczynu Whiteheada z komutatorem
zwykłej grupy podstawowej: dla $p,q\ge2$ wynik
należy do grupy w stopniu $p+q-1$.
\end{uwaga}

\begin{twierdzenie}[Przeszkoda do przedłużenia na iloczyn]
\label{thm:whitehead-product-v7}
Niech $p,q\ge2$, $X$ będzie przestrzenią z punktem
bazowym, a $\alpha:S^p\to X$ i $\beta:S^q\to X$
odwzorowaniami punktowanymi. Mapa
$\alpha\vee\beta:S^p\vee S^q\to X$ przedłuża się
na $S^p\times S^q$ wtedy i tylko wtedy, gdy
$[\alpha,\beta]=0$ w $\pi_{p+q-1}(X)$.
\end{twierdzenie}
\begin{proof}
Przestrzeń $S^p\times S^q$ powstaje z bukietu
$S^p\vee S^q$ przez dołączenie jednej komórki
$D^{p+q}$ wzdłuż mapy $\omega$.
Przedłużenie odwzorowania $f=\alpha\vee\beta$
przez tę komórkę jest równoważne przedłużeniu
$f\circ\omega:S^{p+q-1}\to X$ na dysk.
Taka mapa przedłuża się na dysk dokładnie wtedy,
gdy jest homotopijna do mapy stałej (z uwzględnieniem
punktu bazowego), czyli gdy reprezentowana klasa
$[\alpha,\beta]$ jest zerowa.
\end{proof}

\begin{definicja}[Niezmiennik Hopfa przez komórkę dołączaną]
Niech $n\ge2$ będzie parzyste i
$f:S^{2n-1}\to S^n$ będzie mapą punktowaną.
Niech $C_f=S^n\cup_f e^{2n}$ będzie jej stożkiem
odwzorowania. Wybierzmy generatory
$a\in H^n(C_f;\Z)$ oraz $b\in H^{2n}(C_f;\Z)$.
Istnieje jednoznaczne $H(f)\in\Z$ takie, że
$a\smile a=H(f)b$. Liczbę $H(f)$ nazywamy
\emph{niezmiennikiem Hopfa} $f$ względem tych wyborów.
Przy ustalonej orientacji komórki górnej jej wartość
jest dobrze określona; zmiana generatora $b$ odwraca znak.
\end{definicja}

\begin{twierdzenie}[Kwadrat Whiteheada na \texorpdfstring{$S^2$}{S2}]
\label{thm:whitehead-square-v7}
Dla generatora $\iota_2\in\pi_2(S^2)$ i generatora
Hopfa $\eta\in\pi_3(S^2)\cong\Z$ zachodzi
$[\iota_2,\iota_2]=\pm2\eta$.
\end{twierdzenie}
\begin{proof}
Rozważmy CW-rozkład $S^2\times S^2$ z komórkami
w wymiarach $0,2,2,4$. Niech $x,y$ będą klasami
stopnia $2$ pochodzącymi z dwóch rzutowań, a $u$
generatorem stopnia $4$. W pierścieniu kohomologii
mamy $x^2=y^2=0$ i $x\smile y=u$,
po ustaleniu orientacji. Zatem
$(x+y)^2=2u$.

Niech $\omega:S^3\to S^2\vee S^2$ oznacza mapę
dołączającą komórkę czterowymiarową, a
$\nabla:S^2\vee S^2\to S^2$ mapę składającą oba
składniki przez identyczność. Stożek
$C_{\nabla\circ\omega}=S^2\cup_{\nabla\circ\omega}e^4$
jest skonstruowany z tej samej komórki górnej.
Z własności ilorazu dołączającego istnieje mapa
$F:S^2\times S^2\to C_{\nabla\circ\omega}$,
która na bukiecie jest równa $\nabla$, a na komórce
czterowymiarowej ma stopień $1$. Jeżeli $a,b$
są generatorami odpowiednio $H^2(C_{\nabla\circ\omega})$
i $H^4(C_{\nabla\circ\omega})$, to
$F^*a=x+y$ i $F^*b=\pm u$.
Zatem $F^*(a^2)=2u$, skąd, przez izomorfizm
$F^*:H^4(C_{\nabla\circ\omega})\to H^4(S^2\times S^2)$,
otrzymujemy $H(\nabla\circ\omega)=\pm2$.
Ponieważ niezmiennik Hopfa
$H:\pi_3(S^2)\to\Z$ jest izomorfizmem, wysyłającym
$\eta$ na $\pm1$, stąd $[\iota_2,\iota_2]=\pm2\eta$.
\end{proof}

\begin{przyklad}[Dwa przyklejenia tej samej komórki]
Rozważmy $S^2\vee S^4$ oraz
$\C P^2\simeq S^2\cup_{\eta}e^4$.
Obie przestrzenie mają $H_0=H_2=H_4=\Z$ i zerowe
pozostałe grupy homologii zredukowanej, ale
$H^*(S^2\vee S^4;\Z)$ ma zerowy kwadrat generatora
stopnia $2$, a w $H^*(\C P^2;\Z)$ jego kwadrat
jest generatorem $H^4$. Ta różnica jest ujęta
przez niezmiennik Hopfa mapy dołączającej.
\end{przyklad}

\section{Podnoszenie przez wieżę Postnikowa i wyższe przeszkody}
\begin{twierdzenie}[Przeszkoda podniesienia]
\label{thm:postnikov-lift-v7}
Niech $B$ będzie jednospójnym CW-kompleksem, $A$
grupą abelową, $n\ge2$, a
$k:B\to K(A,n+1)$ mapą. Oznaczmy przez
$p:E=\operatorname{hofib}(k)\to B$ odpowiadającą
fibrację. Dla CW-kompleksu $Z$ i ciągłej mapy
$f:Z\to B$ istnieje mapa $\widetilde f:Z\to E$
ze złożeniem $p\circ\widetilde f=f$ wtedy i tylko
wtedy, gdy $f^*[k]=0$ w $H^{n+1}(Z;A)$.
\end{twierdzenie}
\begin{proof}
Z definicji $E$ punkt nad $b\in B$ jest parą
$(b,\gamma)$, gdzie $\gamma$ jest drogą od $k(b)$
do wyróżnionego punktu $K(A,n+1)$. Podniesienie
$\widetilde f(z)=(f(z),\gamma_z)$ jest ciągłe
dokładnie wtedy, gdy drogi $\gamma_z$ zależą
ciągle od $z$. To jest homotopia od $k\circ f$
do mapy stałej. A zatem podniesienie istnieje
dokładnie wtedy, gdy $[k\circ f]=0$.
Na mocy twierdzenia o reprezentacji kohomologii
$[Z,K(A,n+1)]\cong H^{n+1}(Z;A)$, więc ten
warunek to $f^*[k]=0$.
\end{proof}

\begin{przyklad}[Pierwsza wyższa przeszkoda dla sfery]
Niech $u\in H^2(K(\Z,2);\Z)$ będzie generatorem,
a $Z$ dowolnym CW-kompleksem. Klasa
$c\in H^2(Z;\Z)$ reprezentuje mapę
$f:Z\to P_2S^2\simeq K(\Z,2)$. Istnienie jej
podniesienia do $P_3S^2$ wymaga i wystarcza,
aby $c\smile c=0$ w $H^4(Z;\Z)$, gdyż
$k_3=\pm u^2$. \emph{Nie} jest to samo co
istnienie mapy $Z\to S^2$ indukującej $c$:
w dalszych etapach Postnikowa mogą pojawić się
nowe przeszkody. Jeżeli jednak $\dim Z\le3$,
przeszkoda w stopniu $4$ znika z powodów
wymiarowych, a dalsze również nie występują.
\end{przyklad}

\begin{uwaga}[Uwaga o wyborze podniesienia]
Jeżeli $f^*[k]=0$, to zbiór homotopii
zerujących $k\circ f$, \emph{z ustalonym $f$}
i z homotopiami włóknowymi, jest torsorem
nad $[Z,\Omega K(A,n+1)]\cong H^n(Z;A)$
(po wyborze jednego zerowania). Jeżeli dodatkowo
identyfikuje się podniesienia nad mapami $f$
zmieniającymi się przez homotopię, może pojawić się
działanie grupy pętli przestrzeni odwzorowań do $B$;
wtedy bez dodatkowych hipotez nie należy twierdzić,
że klasy wszystkich podniesień tworzą taki torsor.
\end{uwaga}

\section{Steenrod: od diagonalnej do działań na kohomologii}
\subsection*{Dlaczego zwykły iloczyn kubkowy nie wystarczy?}
Dla dowolnej przestrzeni $X$ z iloczynem kubkowym
znamy mnożenie $H^p(X;\F_2)\otimes H^q(X;\F_2)
\to H^{p+q}(X;\F_2)$, gdzie $\F_2=\Z/2\Z$.
Jednak nie każda operacja kohomologiczna jest
uzyskiwana przez pomnożenie dwóch niezależnie
wybranych klas. Interesują nas \emph{naturalne}
odwzorowania zwiększające stopień
$H^n(X;\F_2)\to H^{n+i}(X;\F_2)$,
zdefiniowane dla wszystkich przestrzeni $X$.

\begin{definicja}[Wyższe diagonalne i iloczyny $\smile_r$]
Niech $C_*(X;\F_2)$ będzie unormowanym singularnym
kompleksem łańcuchowym, $\delta$ jego koróżniczką
na $C^*(X;\F_2)$, a $T(a\otimes b)=b\otimes a$
zamianą czynników tensora (nad $\F_2$ znaki
Koszula znikają). \emph{Wyższymi przybliżeniami diagonalnej}
nazywamy naturalne odwzorowania stopnia $r$
\(\Delta_r:C_m(X;\F_2)\to
(C_*(X;\F_2)\otimes C_*(X;\F_2))_{m+r}\),
$r\ge0$, dla których $\Delta_0$ jest diagonalną
Alexandra--Whitneya i
\[
\partial\Delta_r+\Delta_r\partial=(1+T)\Delta_{r-1}
\quad(r\ge1).
\]
Dla jednorodnych kołańcuchów $a,b$ definiujemy
$\smile_r$ przez
$(a\smile_r b)(c)=(a\otimes b)(\Delta_r c)$,
a $\smile_{-1}$ określamy jako zero.
Wtedy $a\smile_r b$ ma stopień
$\deg a+\deg b-r$.
\end{definicja}

\begin{lemat}[Istnienie i tożsamość koróżniczki]
\label{lem:cup-i-v7}
Istnieją naturalne operatory $\Delta_r$. Korzystamy
z klasycznej unormowanej konstrukcji Steenroda; można je
wybrać tak, by $\smile_0$ zgadzał się z iloczynem
Alexandra--Whitneya, a $a\smile_{\deg a}a=a$
dla kołańcucha o wartościach w $\F_2$ w odpowiednim
stopniu. Dla jednorodnych kołańcuchów zachodzi
\[
\delta(a\smile_r b)=(\delta a)\smile_r b
+a\smile_r(\delta b)+a\smile_{r-1}b+b\smile_{r-1}a.
\]
\end{lemat}
\begin{proof}
Opiszmy konstrukcję komórkową, nie tylko odwołanie
do istnienia. Na singularnym sympleksie $\sigma:\Delta^m\to X$
operatory wyznacza się na standardowym sympleksie
$\Delta^m$, a następnie przenosi przez $\sigma$.
Dla $r=0$ bierzemy zwykły wzór Alexandra--Whitneya:
\(\Delta_0[v_0\dots v_m]=
\sum_{i=0}^m[v_0\dots v_i]\otimes[v_i\dots v_m]\).

Załóżmy, że $\Delta_0,\ldots,\Delta_{r-1}$
zostały już określone naturalnie. Dla każdego
$\Delta^m$ chcemy rozwiązać równanie
$\partial\Delta_r+\Delta_r\partial=(1+T)\Delta_{r-1}$.
Na brzegu sympleksu $\Delta_r$ jest już znana
z indukcji po wymiarze $m$. Żądany brzeg
$\Delta_r[\Delta^m]$ to suma
$(1+T)\Delta_{r-1}[\Delta^m]
+\Delta_r(\partial[\Delta^m])$.
Jest ona cyklem: stosując $\partial$ i hipotezę
indukcyjną otrzymujemy $(1+T)^2\Delta_{r-2}=0$
ponieważ $T^2=1$ w charakterystyce $2$.
Grupa docelowa jest tensorowym kompleksem
łańcuchowym standardowego sympleksu z sobą;
jest on acykliczny w dodatnich stopniach, gdyż
$\Delta^m$ jest ściągalny, a łańcuchy są wolne.
Żądany cykl leży w stopniu $m+r-1>0$
poza krańcowym przypadkiem $m=0,r=1$,
który ustala się bezpośrednio. Można go więc
wypełnić, wybierając $\Delta_r[\Delta^m]$.
To daje jednocześnie indukcję po $m$ i $r$.
Przenoszenie operatorów ze standardowych sympleksów przez
odwzorowania singularne zapewnia naturalność względem
ciągłych odwzorowań. \emph{Trzeba przy tym zaznaczyć
istotne ograniczenie tego argumentu:} sama acykliczność
modeli daje rozwiązania relacji homotopii, lecz nie dowodzi
jeszcze warunku normalizacji
$a\smile_n a=a$ dla dowolnego kołańcucha stopnia $n$.
Do uzyskania tej normalizacji wybieramy standardowe
\emph{unormowane wyższe diagonalne Steenroda}; ich
jawna konstrukcja wymaga dodatkowych warunków na wypełnienia
(np. wzorów przedziałowych na uporządkowanych sympleksach).
W dalszym ciągu normalizacja jest zatem przywoływaną
własnością tej konkretnej klasycznej konstrukcji, a nie
wnioskiem z samego rozwiązywania równania na $\Delta_r$.

Ostatnią tożsamość otrzymujemy, wartościując relację
na $\Delta_r$ przez $a\otimes b$. Człony
z $\partial$ dają koróżniczki obu czynników,
a człon $(1+T)\Delta_{r-1}$ daje
$a\smile_{r-1}b+b\smile_{r-1}a$.
\end{proof}

\begin{definicja}[Kwadraty Steenroda]
Niech $X$ będzie przestrzenią topologiczną,
$n\ge0$, $i\ge0$. Dla kocyklu
$a\in Z^n(X;\F_2)$ definiujemy
\[
\operatorname{Sq}^i[a]=
\begin{cases}
[a\smile_{n-i}a]\in H^{n+i}(X;\F_2),&0\le i\le n,\\
0,&i>n.
\end{cases}
\]
Operacje $\operatorname{Sq}^i$ nazywamy
\emph{kwadratami Steenroda}.
\end{definicja}

\begin{twierdzenie}[Poprawność definicji i własności podstawowe]
\label{thm:steenrod-basic-v7}
Dla każdej przestrzeni $X$ operacje
$\operatorname{Sq}^i:H^n(X;\F_2)\to H^{n+i}(X;\F_2)$
są dobrze określonymi naturalnymi homomorfizmami
przestrzeni wektorowych nad $\F_2$. Ponadto:
$\operatorname{Sq}^0=\id$,
$\operatorname{Sq}^n(x)=x\smile x$ dla $x\in H^n$,
a $\operatorname{Sq}^i(x)=0$ dla $i>\deg x$.
\end{twierdzenie}
\begin{proof}
Oznaczmy $r=n-i$. Dla kocyklu $a$ tożsamość
Lematu~\ref{lem:cup-i-v7} daje
$\delta(a\smile_r a)=2a\smile_{r-1}a=0$.
Zatem przedstawiona formuła wyznacza klasę
kohomologii. Niech $a'=a+\delta b$, gdzie
$b\in C^{n-1}(X;\F_2)$. Po rozwinięciu
$a'\smile_r a'$ różnica w stosunku do
$a\smile_r a$ jest sumą
\[
a\smile_r\delta b+\delta b\smile_r a
+\delta b\smile_r\delta b.
\]
Jest to koróżniczka kołańcucha
\(a\smile_r b+b\smile_r a
+b\smile_r\delta b+b\smile_{r-1}b\):
pierwsze dwa składniki otrzymujemy z tożsamości
dla $a\smile_r b+b\smile_r a$,
a pozostałe z tożsamości dla
$b\smile_r\delta b+b\smile_{r-1}b$.
Gdy $r=0$, składnik z $\smile_{-1}$ jest zerowy.
Zatem wynik nie zależy od reprezentanta.

Dla kocykli $a,c\in Z^n$ wyrazy mieszane
$a\smile_r c+c\smile_r a$ są koróżniczką
$a\smile_{r+1}c$. Wobec tego
$\operatorname{Sq}^i([a]+[c])=
\operatorname{Sq}^i[a]+\operatorname{Sq}^i[c]$.
Naturalność wynika z naturalności $\Delta_r$.
Normalizacja $a\smile_n a=a$ daje
$\operatorname{Sq}^0=\id$, zaś dla $i=n$
mamy $r=0$, czyli zwykły iloczyn kubkowy.
Przypadek $i>n$ jest częścią definicji.
\end{proof}

\begin{twierdzenie}[Wzór Cartana i relacje Adema]
\label{thm:steenrod-cartan-v7}
Dla $x\in H^p(X;\F_2)$, $y\in H^q(X;\F_2)$
oraz $k\ge0$ zachodzi \emph{wzór Cartana}
\[
\operatorname{Sq}^k(x\smile y)=
\sum_{i+j=k}\operatorname{Sq}^i(x)\smile\operatorname{Sq}^j(y).
\]
Ponadto dla $0<a<2b$ zachodzą \emph{relacje Adema}
\[
\operatorname{Sq}^a\operatorname{Sq}^b
=\sum_{j=0}^{\lfloor a/2\rfloor}
\binom{b-j-1}{a-2j}\operatorname{Sq}^{a+b-j}\operatorname{Sq}^j,
\]
gdzie współczynniki dwumianowe liczymy modulo $2$,
a iloczyn operatorów oznacza ich złożenie.
\end{twierdzenie}
\begin{uwaga}[Szkic i status dowodów]
Wzór Cartana wynika z porównania wyższych diagonalnych
dla tensora kompleksów z tensorami diagonalnych
obu czynników: składniki stopnia $k$ powstają ze
wszystkich rozkładów $k=i+j$. Zgodność wymaga jednak
jawnych homotopii między diagonalnymi ekwiwariantnymi.
Relacje Adema wymagają jeszcze wyższych homotopii
dla \emph{iterowanych} diagonalnych. W obecnej iteracji
podajemy oba standardowe twierdzenia i ich zastosowania,
ale \emph{nie deklarujemy pełnego dowodu Cartana ani
Adema}. Stanowią one odrębny etap konstrukcyjny
(teoria spójnych diagonalnych i algebry Steenroda).
Natomiast dowody istnienia kwadratów, ich naturalności,
addytywności, niestabilności oraz górnego kwadratu
zostały podane wyżej.
\end{uwaga}

\section{Przykłady operacji Steenroda}

\begin{przyklad}[Pierwszy kwadrat na trójkącie]
Niech $K$ będzie kompleksem symplicjalnym i
$a\in Z^1(K;\F_2)$. Dla uporządkowanej
zorientowanej $2$-komórki $[v_0,v_1,v_2]$
zwykły iloczyn kubkowy jest dany przez
$\bigl(a\smile a\bigr)[v_0,v_1,v_2]
=a[v_0,v_1]a[v_1,v_2]$. Zatem
$\operatorname{Sq}^1[a]=[a\smile a]$.
To wzór wykonawczy: wystarczy obliczyć wartości
kocyklu na krawędziach. Dla $a$ stopnia $1$
operacje wyższe niż $\operatorname{Sq}^1$ znikają
na mocy niestabilności.
\end{przyklad}

\begin{przyklad}[\texorpdfstring{$\R P^\infty$}{RP infinity}]
Niech $t\in H^1(\R P^\infty;\F_2)$ będzie generatorem.
Mamy $H^*(\R P^\infty;\F_2)\cong\F_2[t]$.
Z własności podstawowych wynika
$\operatorname{Sq}^0(t)=t$,
$\operatorname{Sq}^1(t)=t^2$ i
$\operatorname{Sq}^i(t)=0$ dla $i>1$.
Stosując wzór Cartana do $t^m$, otrzymujemy
\[
\operatorname{Sq}^i(t^m)=\binom mi t^{m+i}.
\]
W $\R P^N$ ten sam wzór obowiązuje, jeśli
$m+i\le N$; w większych stopniach prawa strona
znika ze względów wymiarowych.
\end{przyklad}

\begin{przyklad}[\texorpdfstring{$\C P^\infty$}{CP infinity}]
Niech $u\in H^2(\C P^\infty;\F_2)$ będzie redukcją
modulo $2$ uniwersalnej klasy Cherna. Ponieważ
grupy w nieparzystych stopniach są zerowe,
$\operatorname{Sq}^1(u)=0$, a z własności górnego
kwadratu $\operatorname{Sq}^2(u)=u^2$.
Wzór Cartana daje
$\operatorname{Sq}^{2j}(u^m)=\binom mj u^{m+j}$,
a $\operatorname{Sq}^{2j+1}(u^m)=0$.
\end{przyklad}

\begin{przyklad}[Zastosowanie do cofnięć]
Jeśli $f:X\to\R P^\infty$ klasyfikuje rzeczywistą
wiązkę liniową $L$ i $w_1(L)=f^*t$, to naturalność
kwadratów Steenroda daje
$\operatorname{Sq}^1(w_1(L))=w_1(L)^2$.
Dla przestrzeni o większej liczbie generatorów
operacje potrafią odróżnić mapy, które na poziomie
addytywnym dają te same grupy kohomologii.
\end{przyklad}

\section{Przygotowanie do dalszej teorii}
\begin{uwaga}[Przypadek niejednospójny]
Jeżeli $\pi_1(X)$ nie jest zerowa, jej działanie
na $\pi_n(X)$ dla $n\ge2$ musi wejść do danych
wieży Postnikowa. Odpowiednie niezmienniki $k$
należą do grup kohomologii z \emph{lokalnymi
współczynnikami}. W szczególności nie wolno
zastępować ich bez uzasadnienia grupami
$H^{n+1}(P_{n-1}X;\pi_nX)$ ze stałymi współczynnikami.
\end{uwaga}

\begin{uwaga}[Nieskończone wieże i zbieżność]
Skończone etapy Postnikowa opisują kolejno coraz
więcej grup homotopii. Przejście do granicy
homotopijnej całej wieży jest odrębnym zagadnieniem:
w ogólnej sytuacji występują warunki kompletności
i możliwe człony $\varprojlim^1$. W tym rozdziale
nie utożsamiamy automatycznie $X$ ze zwykłą
granicą odwrotną $\varprojlim P_nX$.
\end{uwaga}

\begin{uwaga}[Miejsce geometrii algebraicznej]
W późniejszym połączeniu z geometrią algebraiczną
ciągi spektralne Serre'a będą punktem odniesienia
dla ciągu Leraya, a operacje Steenroda dla analogicznych
operacji w kohomologii motywicznej. Same
konstrukcje analityczne, etalne i snopowe wymagają
jednak osobnych hipotez i nie są utożsamiane z
operacjami zdefiniowanymi dla topologicznych
CW-kompleksów.
\end{uwaga}


% Iteracja VIII, rozdział główny 15. Do umieszczenia przed rozdziałem o stabilizacji.
\chapter{Operacje kohomologiczne: naturalność, iloczyny i relacje}
\label{ch:operacje-v8}

\section{Dlaczego potrzebujemy operacji, a nie tylko pierścienia kohomologii?}
Iloczyn kubkowy porównuje dwie klasy tej samej przestrzeni. Kwadraty
Steenroda są silniejszym narzędziem: każda operacja
$\operatorname{Sq}^i:H^n(X;\F_2)\to H^{n+i}(X;\F_2)$
jest naturalna względem $X$ i zależy od jednej klasy.
Jeżeli dwie przestrzenie mają izomorficzne pierścienie kohomologii,
nie wynika jeszcze, że identyfikacja zachowuje wszystkie operacje
kohomologiczne. Szczególnie wyrazisty przykład otrzymamy po zawieszeniu
$\C P^2$, kiedy wszystkie iloczyny dodatnich stopni znikną, ale
$\operatorname{Sq}^2$ pozostanie nietrywialne.

W poprzednim rozdziale skonstruowaliśmy operacje za pomocą wyższych
diagonalnych. Należy przy tym oddzielić dwa twierdzenia:
acykliczność kompleksu modeli pozwala rozwiązać relacje homotopii
diagonalnych, ale \emph{nie wymusza} wybranego unormowania ani
wszystkich zgodności iterowanych diagonalnych. Stosując klasyczne
\emph{unormowane operacje Steenroda}, wykorzystujemy także te dodatkowe
zgodności. Właśnie one są potrzebne do wzoru Cartana i relacji Adema.

\section{Operacje naturalne i kobrzeg Bocksteina}
\begin{definicja}[Operacja kohomologiczna]
Niech $G,H$ będą grupami abelowymi, a $n,m\ge0$.
\emph{Naturalną operacją kohomologiczną} stopnia $(n,m)$ nazywamy
rodzinę funkcji
$\theta_X:H^n(X;G)\to H^m(X;H)$, indeksowaną
przestrzeniami topologicznymi $X$, taką że dla każdego ciągłego
$f:X\to Y$ i $u\in H^n(Y;G)$ zachodzi
$f^*\theta_Y(u)=\theta_X(f^*u)$.
Nie wymagamy addytywności, o ile nie jest jawnie wskazana.
Jeżeli $m=n+i$, mówimy, że operacja \emph{podwyższa stopień o $i$}.
\end{definicja}

\begin{przyklad}[Dlaczego mnożenie nie jest zawsze addytywne]
Przyporządkowanie $u\mapsto u\smile u$ jest naturalne. Dla klas
całkowitych stopnia parzystego nie musi być addytywne:
$(u+v)^2=u^2+2uv+v^2$. Przy współczynnikach w $\F_2$ składniki mieszane
znoszą się, ponieważ iloczyn jest przemienny w charakterystyce $2$.
Natomiast addytywność wszystkich $\operatorname{Sq}^i$
nie wynika z samej naturalności: jest osobną własnością
standardowej konstrukcji operacji Steenroda.
\end{przyklad}

\begin{definicja}[Kobrzeg Bocksteina]
Dla krótkiego ciągu dokładnego grup abelowych
$0\to A\xrightarrow{\iota}B\xrightarrow{q}C\to0$
\emph{homomorfizmem Bocksteina} nazywamy łączący homomorfizm
$\beta:H^n(X;C)\to H^{n+1}(X;A)$
w długim ciągu kohomologii współczynników.
Dla $0\to\F_2\xrightarrow{\,\cdot2\,}\Z/4\Z\to\F_2\to0$
piszemy $\beta_2:H^n(X;\F_2)\to H^{n+1}(X;\F_2)$.
\end{definicja}

\begin{stwierdzenie}[Wzór wykonawczy dla Bocksteina]
\label{prop:bockstein-v8}
Niech $a\in Z^n(X;\F_2)$ i niech $\widetilde a\in C^n(X;\Z/4\Z)$
będzie dowolnym podniesieniem kołańcucha $a$. Wtedy
$\delta\widetilde a$ przyjmuje wartości w podgrupie
$2(\Z/4\Z)$, a klasa kołańcucha
$(\delta\widetilde a)/2$ zidentyfikowanego z kołańcuchem
$\F_2$ jest równa $\beta_2[a]$.
\end{stwierdzenie}
\begin{proof}
Ponieważ $q(\delta\widetilde a)=\delta(q\widetilde a)=\delta a=0$,
wartości $\delta\widetilde a$ należą do jądra $q=2(\Z/4\Z)$.
Wybierzmy $b\in C^{n+1}(X;\F_2)$ spełniający
$\iota b=\delta\widetilde a$. Z równości
$0=\delta^2\widetilde a=\iota\delta b$ i różnowartościowości
$\iota$ dostajemy $\delta b=0$.
Jeżeli zmienimy podniesienie o $\iota c$, to $b$ zmieni się
o $\delta c$; jeśli $a$ zmieni się o kobrzeg, wynik zmieni się
co najwyżej o kobrzeg. Jest to dokładnie konstrukcja mapy łączącej
w długim ciągu pochodzącym z krótkiego ciągu kompleksów kołańcuchów.
\end{proof}

\begin{twierdzenie}[Identyfikacja pierwszego kwadratu]
\label{thm:sq1-bock-v8}
Dla każdej przestrzeni topologicznej $X$ i każdego $n\ge0$
zachodzi $\operatorname{Sq}^1=\beta_2$ jako homomorfizm
$H^n(X;\F_2)\to H^{n+1}(X;\F_2)$.
W szczególności $\operatorname{Sq}^1\operatorname{Sq}^1=0$.
\end{twierdzenie}
\begin{proof}
Pierwsza równość jest \emph{standardową tożsamością unormowanej
konstrukcji $\smile_i$}. Jej kołańcuchową weryfikację otrzymuje się,
porównując wzór z poprzedniego rozdziału
$\operatorname{Sq}^1[a]=[a\smile_{n-1}a]$
z kołańcuchem $(\delta\widetilde a)/2$ z
Stwierdzenia~\ref{prop:bockstein-v8}. Porównanie wymaga jawnego
unormowanego modelu $\smile_i$; sama relacja
$\partial\Delta_i+\Delta_i\partial=(1+T)\Delta_{i-1}$
nie jest do niego wystarczająca. Drugą równość można udowodnić bez tego porównania.
Niech $\rho_2:H^n(X;\Z)\to H^n(X;\F_2)$
było redukcją modulo $2$, a
$\beta_{\Z}:H^n(X;\F_2)\to H^{n+1}(X;\Z)$
homomorfizmem łączącym ciągu
$0\to\Z\xrightarrow{\cdot2}\Z\to\F_2\to0$.
Wprost z kołańcuchowej definicji obydwu homomorfizmów
łączących mamy $\beta_2=\rho_2\circ\beta_{\Z}$.
Dokładność długiego ciągu daje
$\beta_{\Z}\circ\rho_2=0$. Dlatego
$\beta_2^2=\rho_2\beta_{\Z}\rho_2\beta_{\Z}=0$.
\end{proof}

\begin{uwaga}[Zakres poprzedniego dowodu]
W identyfikacji $\operatorname{Sq}^1=\beta_2$ wykorzystujemy
znany fakt o \emph{konkretnych} unormowanych diagonalnych, a nie
przypisujemy go abstrakcyjnemu istnieniu dowolnej rodziny homotopii.
W pozostałych obliczeniach niezależność od modelu jest rozumiana
na poziomie operacji kohomologicznych, nie równości kołańcuchów.
\end{uwaga}

\section{Całkowita operacja i rachunki z użyciem wzoru Cartana}
\begin{definicja}[Całkowity kwadrat Steenroda]
Dla przestrzeni $X$ i klasy jednorodnej $x\in H^n(X;\F_2)$
przyjmujemy $\operatorname{Sq}(x)=\sum_{i=0}^n\operatorname{Sq}^i(x)$.
Przedłużamy ten wzór addytywnie na $H^*(X;\F_2)$.
Jest to odwzorowanie grup addytywnych; składowa stopnia $n+i$
operacji $\operatorname{Sq}$ ma stopień $i$.
\end{definicja}

\begin{twierdzenie}[Postać multiplikatywna wzoru Cartana]
\label{thm:cartan-v8}
Dla dowolnej przestrzeni $X$ całkowity kwadrat jest
homomorfizmem pierścieni z jedynką:
$\operatorname{Sq}(x\smile y)=\operatorname{Sq}(x)\smile\operatorname{Sq}(y)$.
Równość składowych stopnia $\deg x+\deg y+k$ jest
wzorem Cartana z Twierdzenia~\ref{thm:steenrod-cartan-v7}.
\end{twierdzenie}
\begin{proof}
Dokładnie rozdzielamy krok formalny od geometrycznego.
Korzystamy z następującego standardowego twierdzenia
pomocniczego o \emph{zewnętrznym kwadracie}: istnieje
naturalna konstrukcja klas ekwiwariantnych, multiplikatywna
względem iloczynu kubkowego, której ograniczenie do
przekątnej daje unormowane kwadraty Steenroda.
Nie dowodzimy tu tego twierdzenia ekwiwariantnego;
pokazujemy natomiast szczegółowo, jak wynika z niego Cartan.
W klasycznej konstrukcji wybiera się wolną, ściągalną
przestrzeń $E\Sigma_2$, gdzie
$\Sigma_2=\{1,\tau\}$ działa przez zamianę dwóch czynników.
Oznaczamy $B\Sigma_2:=E\Sigma_2/\Sigma_2$;
ma ona typ homotopijny $\R P^\infty$.
Następnie tworzy się przestrzeń orbit
$E\Sigma_2\times_{\Sigma_2}(X\times X)$.
Po ograniczeniu do diagonalnego odwzorowania
$X\to X\times X$ otrzymuje się klasę w
$H^{2n}(B\Sigma_2\times X;\F_2)$ o postaci
$\sum_{i=0}^n t^{n-i}\otimes\operatorname{Sq}^i(x)$,
przy czym $t\in H^1(B\Sigma_2;\F_2)$ jest generatorem.
Naturalność i formuła Künnetha nad ciałem zapewniają jednoznaczność
współczynników tej ekspansji.

Mnożenie zewnętrznych potęg jest zgodne z iloczynem klas,
po czym ograniczenie do przekątnej zachowuje iloczyny.
Dla $x,y$ daje to iloczyn wielomianów
$\bigl(\sum_i t^{n-i}\otimes\operatorname{Sq}^i(x)\bigr)
 \bigl(\sum_j t^{m-j}\otimes\operatorname{Sq}^j(y)\bigr)$.
Porównanie współczynników przy $t^{n+m-k}$
pozwala otrzymać żądany wzór.
\end{proof}
\begin{uwaga}[O rzeczywistym zakresie argumentu]
W dowodzie użyto \emph{twierdzenia o zewnętrznych potęgach}:
istnienia klasy i jej multiplikatywności, a także jej
identyfikacji z wcześniejszymi kołańcuchowymi kwadratami.
Te trzy własności wymagają technicznego dowodu ekwiwariantnego
na poziomie rezolwenty grupy $\Sigma_2$; opisaliśmy krok rachunkowy,
a nie pełny niezależny dowód istnienia potęgi. Daje to przejrzyste
wyprowadzenie Cartana \emph{z jawnie wskazanej konstrukcji},
bez ukrywania potrzebnego faktu pomocniczego.
\end{uwaga}

\begin{przyklad}[Rachunek na iloczynie rzeczywistych przestrzeni rzutowych]
Niech $X=\R P^{a}\times\R P^{b}$ i niech
$u,v\in H^1(X;\F_2)$ będą klasami z odpowiednich czynników.
Dla $0\le i\le a$, $0\le j\le b$ wzór Cartana daje
$\operatorname{Sq}^{k}(u^iv^j)=
\sum_{r+s=k}\binom ir\binom js u^{i+r}v^{j+s}$,
przy czym jednomiany przekraczające relacje
$u^{a+1}=v^{b+1}=0$ przyjmujemy za zero.
W szczególności dla $i=j=1$:
$\operatorname{Sq}^1(uv)=u^2v+uv^2$,
$\operatorname{Sq}^2(uv)=u^2v^2$.
\end{przyklad}

\section{Relacje Adema: co rzeczywiście wynika z przykładów?}
\begin{twierdzenie}[Relacje Adema w charakterystyce dwa]
\label{thm:adem-v8}
Dla liczb całkowitych $0<a<2b$ i każdej przestrzeni $X$
na $H^*(X;\F_2)$ zachodzi
\[
\operatorname{Sq}^a\operatorname{Sq}^b
 =\sum_{j=0}^{\lfloor a/2\rfloor}
 \binom{b-j-1}{a-2j}
 \operatorname{Sq}^{a+b-j}\operatorname{Sq}^{j},
\]
gdzie dwumiany redukujemy modulo $2$.
\end{twierdzenie}
\begin{uwaga}[Status dowodu]
Jest to twierdzenie klasyczne. Pełny dowód kołańcuchowy
wymaga porównania dwóch iterowanych operacji z działaniem
większej grupy symetrii (np. w modelu rozszerzonych potęg),
a więc danych nieobecnych w samej relacji dla $\smile_i$.
Nie przyjmujemy rachunku na jednym pierścieniu
$H^*(\R P^\infty;\F_2)$ za dowód \emph{uniwersalny}.
Natomiast poniżej rozpisujemy konsekwencje i sprawdzamy ich
zgodność z wykonawczymi formułami.
\end{uwaga}
\begin{przyklad}[Dwie pierwsze relacje]
Dla $(a,b)=(1,1)$ jedyny dwumian jest
$\binom{0}{1}=0$, stąd
$\operatorname{Sq}^1\operatorname{Sq}^1=0$.
Jest to także konsekwencja tożsamości
$\operatorname{Sq}^1=\beta_2$.
Dla $(a,b)=(2,2)$ mamy
$\binom{1}{2}=0$, $\binom{0}{0}=1$, więc
$\operatorname{Sq}^2\operatorname{Sq}^2=
\operatorname{Sq}^3\operatorname{Sq}^1$.
Na klasie $t^m\in H^m(\R P^\infty;\F_2)$
obie strony działają jako mnożenie przez
$\binom m2\binom{m+2}2$ oraz
$m\binom{m+1}3$, liczone modulo $2$.
Zgodność współczynników można sprawdzić z twierdzenia Lucasa,
ale sam ten rachunek nie stanowi jeszcze dowodu
równości wszystkich naturalnych operacji.
\end{przyklad}

\begin{twierdzenie}[Postać dopuszczalna jednomianów]
Każde złożenie dodatnich kwadratów Steenroda można,
wykorzystując relacje Adema, przedstawić jako sumę
złożeń
$\operatorname{Sq}^{i_1}\cdots\operatorname{Sq}^{i_r}$
o indeksach $i_j\ge2i_{j+1}$. Takie złożenia nazywamy
\emph{dopuszczalnymi}. Twierdzenie dotyczy rozpięcia
algebry Steenroda przez jednomiany dopuszczalne;
ich liniowa niezależność wymaga odrębnego argumentu.
\end{twierdzenie}
\begin{proof}
Jeżeli $i_j<2i_{j+1}$, relacja Adema zastępuje
parę $\operatorname{Sq}^{i_j}\operatorname{Sq}^{i_{j+1}}$
sumą złożeń $\operatorname{Sq}^{i_j+i_{j+1}-k}\operatorname{Sq}^{k}$,
gdzie $k\le\lfloor i_j/2\rfloor$.
Lewa liczba jest większa od prawej co najmniej dwukrotnie,
a indeksy zachowują sumę.
Ściśle uzasadnijmy zakończenie tej procedury. Dla słowa
$\operatorname{Sq}^{i_1}\cdots\operatorname{Sq}^{i_\ell}$
o dodatnich indeksach określmy sumy końcowe
$s_j=i_j+\cdots+i_\ell$ dla $1\le j\le\ell$.
W pierwszej kolejności porównujemy długość słów, a przy
tej samej długości porównujemy leksykograficznie ciągi
$(s_\ell,s_{\ell-1},\ldots,s_1)$. W podstawieniu
pary $(a,b)$ przez $(a+b-j,j)$, gdzie $j>0$,
wszystkie sumy końcowe na prawo od tej pary pozostają
niezmienione, zaś suma zaczynająca się od drugiego
wyrazu pary maleje z $b+S$ do $j+S$, ponieważ
$j\le\lfloor a/2\rfloor<b$; $S$ jest sumą
indeksów leżących jeszcze dalej na prawo.
Każde takie podstawienie ściśle zmniejsza więc
wskazany ciąg leksykograficznie. Gdy $j=0$, otrzymujemy
$\operatorname{Sq}^{a+b}\operatorname{Sq}^0=
\operatorname{Sq}^{a+b}$, toteż długość słowa maleje.
Ponieważ całkowita suma indeksów nie zmienia się,
a długość nigdy nie rośnie, rozpatrujemy tylko skończenie
wiele możliwych słów. Nie może zatem istnieć nieskończony
ciąg redukcji.
\end{proof}

\section{Są takie przestrzenie, których pierścień kohomologii nie rozróżnia}
\begin{lemat}[Iloczyny na zawieszeniu]
\label{lem:susp-cupzero-v8}
Jeżeli $X$ jest dobrze bazowanym CW-kompleksem,
to iloczyn dwóch klas dodatnich stopni w
$\widetilde H^*(\Sigma X;R)$ jest zerowy
(dla dowolnego pierścienia współczynników $R$).
\end{lemat}
\begin{proof}
Zawieszenie $\Sigma X$ jest przestrzenią ko-H: posiada
odwzorowanie uszczypnięcia
$\Sigma X\to\Sigma X\vee\Sigma X$.
Złożenie z włączeniem bukietu do iloczynu jest
homotopijne do przekątnej
$\Sigma X\to\Sigma X\times\Sigma X$.
Gdy $a,b$ mają dodatnie stopnie, ich zewnętrzny
iloczyn ograniczony do bukietu jest zerowy,
ponieważ na każdym składniku bukietu jeden czynnik
odwzorowuje się na punkt bazowy. Ponieważ iloczyn kubkowy
jest przeciwobrazem zewnętrznego po przekątnej,
$a\smile b=0$.
\end{proof}

\begin{twierdzenie}[Zawieszona przestrzeń rzutowa nie jest bukietem sfer]
\label{thm:sigma-cp2-v8}
Niech $Y=\Sigma\C P^2$. Wtedy $Y$ i $S^3\vee S^5$
mają izomorficzne całkowite grupy homologii i izomorficzne
pierścienie kohomologii (ze współczynnikami $\Z$ i $\F_2$),
ale nie są homotopijnie równoważne.
\end{twierdzenie}
\begin{proof}
Rozkład komórkowy $\C P^2\simeq S^2\cup_\eta e^4$
daje $Y\simeq S^3\cup_{\Sigma\eta}e^5$.
Obie porównywane przestrzenie mają po jednej komórce
niezerowej w stopniach $3$ i $5$, więc z różniczek
komórkowych otrzymujemy
$\widetilde H_3=\widetilde H_5\cong\Z$ oraz pozostałe
zredukowane grupy homologii równe zeru.
Iloczyny dodatnich stopni są zerowe na obu przestrzeniach:
na $Y$ z Lematu~\ref{lem:susp-cupzero-v8}, a na bukiecie
sfer bezpośrednio.

Niech $u\in H^2(\C P^2;\F_2)$ będzie redukcją klasy
hiperpłaszczyzny, $u^2\ne0\in H^4(\C P^2;\F_2)$.
Z poprzedniego rozdziału wiemy, że
$\operatorname{Sq}^2(u)=u^2$.
Operacje Steenroda są \emph{stabilne}, to znaczy
komutują z izomorfizmem zawieszenia
$\sigma:\widetilde H^q(X;\F_2)\to
\widetilde H^{q+1}(\Sigma X;\F_2)$.
Dlatego
$\operatorname{Sq}^2(\sigma u)=\sigma(u^2)\ne0$
w $H^5(Y;\F_2)$.
Tymczasem na $S^3\vee S^5$ klasa stopnia $3$
pochodzi z $S^3$, a
$\operatorname{Sq}^2:H^3(S^3;\F_2)\to H^5(S^3;\F_2)$
jest zerowa, więc na bukiecie również.
Równoważność homotopijna zachowywałaby operacje
naturalne, zatem nie może istnieć.
\end{proof}

\begin{uwaga}[Dlaczego stabilność kwadratów jest istotna]
Stabilność $\operatorname{Sq}^i$ można sprawdzać na
unormowanym modelu wyższych diagonalnych dla par
$(CX,X)$ i na łączącej w ciągu kohomologii względnej.
Pełna zgodność łańcuchowa wymaga tego samego wyboru
spójnych diagonalnych, który pojawia się w formule Cartana.
Używamy tutaj standardowego twierdzenia o stabilności,
a nie wyprowadzamy go wyłącznie z naturalności operacji.
\end{uwaga}

\section{Wiązki wektorowe i charakterystyczny wzór Thoma}
\begin{twierdzenie}[Kwadraty Steenroda klasy Thoma]
\label{thm:thom-wu-v8}
Niech $\xi\to B$ będzie rzeczywistą wiązką wektorową
rzędu $r$ nad CW-kompleksem, a
$U\in H^r(D(\xi),S(\xi);\F_2)$ jej mod-$2$ klasą Thoma.
Dla $0\le i\le r$ istnieją jednoznaczne klasy
$w_i(\xi)\in H^i(B;\F_2)$ takie, że
$\operatorname{Sq}^i(U)=p^*w_i(\xi)\smile U$.
Są to klasy Stiefela--Whitneya; $w_0=1$,
a $w_i=0$ dla $i>r$.
\end{twierdzenie}
\begin{proof}
Względny izomorfizm Thoma
$H^i(B;\F_2)\xrightarrow{\cong}
H^{r+i}(D(\xi),S(\xi);\F_2)$,
$x\mapsto p^*x\smile U$,
pozwala jednoznacznie zdefiniować $w_i$ jako przeciwobraz
$\operatorname{Sq}^i(U)$. Z własności
$\operatorname{Sq}^0=\id$ dostajemy $w_0=1$.
Z niestabilności operacji $\operatorname{Sq}^i(U)=0$
dla $i>r$, więc $w_i=0$ dla $i>r$.
Naturalność klas wynika z naturalności izomorfizmu
Thoma i operacji Steenroda.
\end{proof}

\begin{wniosek}[Wzór Whitney'a]
Dla dwóch wiązek rzeczywistych $\xi,\eta\to B$
zachodzi $w(\xi\oplus\eta)=w(\xi)\smile w(\eta)$,
gdzie $w=\sum_iw_i$.
\end{wniosek}
\begin{proof}
Klasa Thoma sumy prostej odpowiada zewnętrznemu
iloczynowi klas Thoma obu wiązek po przekątnej podstawy.
Stosujemy wzór Cartana do tego iloczynu:
$\operatorname{Sq}(U_{\xi\oplus\eta})
 =\operatorname{Sq}(U_\xi)\smile\operatorname{Sq}(U_\eta)$.
Następnie wyciągamy czynniki $w(\xi)$ i $w(\eta)$
oraz korzystamy z różnowartościowości izomorfizmu Thoma.
\end{proof}

\begin{przyklad}[Wiązka tautologiczna nad $\R P^n$]
Dla rzeczywistej wiązki liniowej tautologicznej
$\gamma^1\to\R P^n$ z generatorem
$t\in H^1(\R P^n;\F_2)$ mamy $w(\gamma^1)=1+t$.
Ponieważ $\gamma^1\oplus(\gamma^1)^\perp$ jest trywialna,
ze wzoru Whitney'a otrzymujemy
$w((\gamma^1)^\perp)=(1+t)^{-1}=1+t+\cdots+t^n$.
Równości zachodzą w pierścieniu
$\F_2[t]/(t^{n+1})$; nie zakładamy zbieżności
nieskończonego szeregu formalnego.
\end{przyklad}

\section{Co wymaga dalszej techniki?}
Sam wzór Cartana jest w pełni obliczalny na klasach
pierścieni generowanych przez elementy małego stopnia;
relacje Adema pozwalają upraszczać iterowane kwadraty.
Pełny niezależny dowód ogólnych relacji Adema wymaga
budowy spójnych działań symetrycznych na iterowanych
produktach i zostaje wyraźnie oddzielony od podanych
rachunków. W następnym rozdziale stabilność operacji
będzie kluczem do niezerowości zawieszonej mapy Hopfa.


% Iteracja VIII, rozdział główny 16. Do umieszczenia przed \appendix.
\chapter{Zawieszenie, twierdzenie Freudenthala i homotopia stabilna}
\label{ch:stabilna-v8}

\section{Motywacja: co pozostaje po wielokrotnym zawieszaniu?}
Dla każdej liczby $n\ge1$ wiemy, że
$\pi_n(S^n)\cong\Z$. Dla grup
$\pi_{n+k}(S^n)$ przy ustalonym $k\ge1$ wyniki
są trudniejsze, lecz po dostatecznym zwiększeniu $n$
grupy przestają się zmieniać. Ta stabilizacja jest
konsekwencją twierdzenia Freudenthala.

W tym rozdziale konsekwentnie używamy \emph{zredukowanego}
zawieszenia
$\Sigma X=S^1\wedge X$ dla CW-kompleksu z punktem
bazowym. Dla $X=S^n$ jest ono homeomorficzne z $S^{n+1}$.
Przyjmujemy standardowe punkty bazowe włączone jako
podkompleksy, aby ilorazy i pary były kofibracjami.

\begin{definicja}[Zawieszenie odwzorowania]
Dla punktowanego odwzorowania $f:X\to Y$
definiujemy $\Sigma f=\id_{S^1}\wedge f:
\Sigma X\to\Sigma Y$.
Jeżeli $f:(S^q,*)\to(X,x_0)$, to
$\Sigma f:S^{q+1}\to\Sigma X$ wyznacza
\emph{homomorfizm zawieszenia}
$E:\pi_q(X)\to\pi_{q+1}(\Sigma X)$.
Dla $q=1$ mapę rozumiemy jako homomorfizm z grupy,
niekoniecznie abelowej, do abelowej grupy $\pi_2$.
\end{definicja}

\begin{stwierdzenie}[Naturalność]
Jeżeli $f:X\to Y$ jest punktowane, to dla $q\ge1$
zachodzi $(\Sigma f)_*\circ E_X=E_Y\circ f_*$.
\end{stwierdzenie}
\begin{proof}
Dla reprezentanta $\alpha:S^q\to X$ oba złożenia są
reprezentowane przez to samo odwzorowanie
$\Sigma(f\circ\alpha)=(\Sigma f)\circ(\Sigma\alpha)$.
\end{proof}

\section{Łączność przestrzeni i dokładność homotopijna}
\begin{definicja}[$r$-spójność]
Niech $r\ge0$. Punktowaną przestrzeń $X$ nazywamy
\emph{$r$-spójną}, jeśli jest łukowo spójna i
$\pi_i(X)=0$ dla $1\le i\le r$.
Parę punktowaną $(Y,A)$ nazywamy $r$-spójną,
jeśli $\pi_i(Y,A)=0$ dla $1\le i\le r$
oraz $\pi_0(A)\to\pi_0(Y)$ jest surjekcją.
Dla $r=0$ oznacza to tylko warunek na składowe.
\end{definicja}

\begin{lemat}[Stożek i względne grupy homotopii]
\label{lem:cone-relative-v8}
Niech $X$ będzie $r$-spójnym CW-kompleksem,
$r\ge0$. Wtedy para $(CX,X)$ jest $(r+1)$-spójna,
a dla $q\ge2$ odwzorowanie brzegowe
$\partial:\pi_{q+1}(CX,X)\to\pi_q(X)$
jest izomorfizmem.
\end{lemat}
\begin{proof}
Stożek jest ściągalny, więc $\pi_q(CX)=0$ dla
$q\ge1$. Fragment długiego ciągu homotopii pary ma postać
$\pi_{q+1}(CX)\to\pi_{q+1}(CX,X)
\xrightarrow{\partial}\pi_q(X)\to\pi_q(CX)$.
Skrajne grupy są zerowe, stąd izomorfizm.
Dla $q\le r$ zeruje się też $\pi_{q+1}(CX,X)$;
warunki w niskich stopniach wynikają z łączności $X$.
\end{proof}

\begin{twierdzenie}[Ekscyzja homotopijna Blakersa--Masseya]
\label{thm:bm-v8}
Niech $Z=A\cup B$ będzie triadą CW, gdzie
$C=A\cap B$ jest niepuste i łukowo spójne.
Jeżeli pary $(A,C)$ i $(B,C)$ są odpowiednio
$m$-spójna i $n$-spójna, $m,n\ge1$, to homomorfizm
$\pi_q(A,C)\to\pi_q(Z,B)$ jest izomorfizmem dla
$q<m+n$ i surjekcją dla $q=m+n$.
\end{twierdzenie}
\begin{uwaga}[Status dowodu ekscyzji homotopijnej]
Jest to głębokie twierdzenie i \emph{nie} wynika ze
zwykłej ekscyzji homologicznej. Pełny dowód można
przeprowadzić przez redukcję do triad CW oraz
usuwanie przecięć obrazów kostek z komórkami
(ogólna pozycja). Wykorzystujemy je jako
wcześniej jawnie sformułowany wynik pomocniczy;
dowód Freudenthala poniżej jest pełny \emph{względem}
tego twierdzenia, nie stanowi niezależnego dowodu
Blakersa--Masseya.
\end{uwaga}

\begin{twierdzenie}[Freudenthal]
\label{thm:freudenthal-v8}
Niech $X$ będzie $r$-spójnym, dobrze bazowanym
CW-kompleksem, gdzie $r\ge0$. Wówczas
$E:\pi_q(X)\to\pi_{q+1}(\Sigma X)$
jest izomorfizmem dla $1\le q<2r+1$ oraz
surjekcją dla $q=2r+1$.
\end{twierdzenie}
\begin{proof}
Przedstawmy zawieszenie jako sumę dwóch
zredukowanych stożków $C_+X$ i $C_-X$,
których przecięcie ma typ homotopijny $X$.
Oba stożki są ściągalne. Pary $(C_+X,X)$
i $(C_-X,X)$ są $(r+1)$-spójne na mocy
Lematu~\ref{lem:cone-relative-v8}.
Ekscyzja homotopijna dla tej triady daje
$\pi_{q+1}(C_+X,X)\to
\pi_{q+1}(\Sigma X,C_-X)$
izomorfizm, gdy $q+1<2(r+1)$, i surjekcję,
gdy $q+1=2(r+1)$.

Z długich ciągów homotopii par i ściągalności
stożków otrzymujemy naturalne identyfikacje
$\pi_{q+1}(C_+X,X)\cong\pi_q(X)$ oraz
$\pi_{q+1}(\Sigma X,C_-X)\cong
\pi_{q+1}(\Sigma X)$ dla $q\ge2$.
Homomorfizm pomiędzy tymi dwiema grupami
przyporządkowuje reprezentantowi $f:S^q\to X$
mapę $\Sigma f:S^{q+1}\to\Sigma X$,
co widać przez potraktowanie reprezentanta względnego
jako mapy $(D^{q+1},S^q)\to(C_+X,X)$
i sklejenie drugiego stożka.
Zatem wskazana ekscyzja jest dokładnie homomorfizmem $E$.
Daje to podany zakres izomorfizmu i surjekcji.
Pozostał przypadek $q=1$. Jeżeli $r\ge1$, obie
grupy są zerowe: $X$ jest jednospójna, a $\Sigma X$
jest co najmniej $2$-spójna. Istotnie, $\Sigma X$ jest
jednospójna, ma $H_2(\Sigma X;\Z)\cong H_1(X;\Z)=0$,
a więc $\pi_2(\Sigma X)=0$ z twierdzenia Hurewicza. Jeżeli $r=0$, przestrzeń $X$
jest łukowo spójna i $\Sigma X$ jest jednospójna.
Twierdzenie Hurewicza, izomorfizm zawieszenia w homologii
oraz twierdzenie o abelianizacji grupy podstawowej dają
naturalne identyfikacje
\[
\pi_2(\Sigma X)\cong H_2(\Sigma X;\Z)
\cong H_1(X;\Z)\cong\pi_1(X)_{\mathrm{ab}}.
\]
Aby sprawdzić, że to właśnie $E$, weźmy pętlę
$\alpha:S^1\to X$; obraz klasy $E[\alpha]$ przez
homomorfizm Hurewicza jest zawieszeniem cyklu
reprezentującego $[\alpha]$. Wynika to z naturalności
homomorfizmu Hurewicza oraz izomorfizmu zawieszenia dla
$S^1$ i $X$. Zatem $E$ jest odwzorowaniem abelianizacji,
w szczególności surjekcją, czego wymaga teza.
\end{proof}

\begin{center}
\begin{tikzpicture}[>=Stealth,font=\small]
\node[draw,rounded corners,minimum width=30mm,minimum height=11mm] (X) at (0,0) {$X$};
\node[draw,rounded corners,minimum width=30mm,minimum height=11mm] (SX) at (4.3,0) {$\Sigma X$};
\node[draw,rounded corners,minimum width=30mm,minimum height=11mm] (SSX) at (8.6,0) {$\Sigma^2 X$};
\draw[->,thick] (X)-- node[above]{zawieszenie} (SX);
\draw[->,thick] (SX)-- node[above]{zawieszenie} (SSX);
\node[font=\footnotesize] at (4.3,-1.05)
{Po dostatecznie wielu krokach grupy $\pi_{q+k}(\Sigma^kX)$ stabilizują się.};
\end{tikzpicture}
\end{center}

\begin{przyklad}[Ostry brzeg zakresu]
Sfera $S^2$ jest $1$-spójna, więc
$E:\pi_2(S^2)\to\pi_3(S^3)$ jest izomorfizmem,
ale dla $q=3$ twierdzenie gwarantuje już tylko surjekcję
$E:\pi_3(S^2)\to\pi_4(S^3)$.
Pokażemy dalej, że jest to redukcja $\Z\to\Z/2$.
W szczególności zastąpienie w twierdzeniu
warunku $q<2r+1$ przez $q\le2r+1$
byłoby nieprawdziwe.
\end{przyklad}

\section{Stabilne grupy homotopii sfer}
\begin{definicja}[Stabilna grupa homotopii]
Dla dobrze bazowanego CW-kompleksu $X$ i $k\in\Z$
definiujemy, o ile stopnie po pewnym zawieszeniu są dodatnie,
$\pi_k^{\mathrm s}(X)=
\varinjlim_{j}\pi_{k+j}(\Sigma^jX)$,
gdzie morfizmami układu są homomorfizmy zawieszenia.
W szczególności dla $k\ge0$ kładziemy
$\pi_k^{\mathrm s}=\pi_k^{\mathrm s}(S^0)
=\varinjlim_{n}\pi_{n+k}(S^n)$.
Grupy $\pi_k^{\mathrm s}$ nazywamy \emph{stabilnymi
pniami homotopii sfer} (dla $k\ge0$).
\end{definicja}

\begin{twierdzenie}[Stabilizacja w ustalonym stopniu]
\label{thm:stable-range-v8}
Niech $X$ będzie niepustym, dobrze bazowanym
CW-kompleksem. Dla każdego ustalonego $k\ge0$
odwzorowania
$\pi_{k+j}(\Sigma^j X)\to
\pi_{k+j+1}(\Sigma^{j+1}X)$
są izomorfizmami dla dostatecznie dużych $j$.
Dla sfer $S^n$ i $k\ge0$ wystarcza $n\ge k+2$.
\end{twierdzenie}
\begin{proof}
Po pierwszym zawieszeniu $\Sigma X$ jest łukowo spójne,
a dla $j\ge1$ przestrzeń $\Sigma^jX$ jest co najmniej
$(j-1)$-spójna, gdy $X$ jest niepustym CW-kompleksem.
Stosujemy Twierdzenie~\ref{thm:freudenthal-v8}
z $r=j-1$ do grupy o indeksie $q=k+j$.
Wymagamy $k+j<2(j-1)+1=2j-1$, czyli $j>k+1$.
Dla $X=S^0$ otrzymujemy $\Sigma^nS^0\cong S^n$
i ten sam rachunek daje $n\ge k+2$.
\end{proof}

\begin{wniosek}[Zerowy stabilny pień sfer]
$\pi_0^{\mathrm s}\cong\Z$.
\end{wniosek}
\begin{proof}
Dla $n\ge2$ mamy $\pi_n(S^n)\cong\Z$,
gdzie izomorfizm wyznacza stopień odwzorowania.
Zawieszenie zachowuje stopień, co można sprawdzić
na homologii fundamentalnej i izomorfizmie zawieszenia
$\widetilde H_n(S^n)\cong\widetilde H_{n+1}(S^{n+1})$.
Dlatego system stabilizujący jest układem identyczności
na $\Z$.
\end{proof}

\section{Niezerowy stabilny element Hopfa}
\begin{lemat}[Rozszczepienie zawieszonego iloczynu]
\label{lem:sigma-product-v8}
Dla punktowanych CW-kompleksów $X,Y$ istnieje naturalna
równoważność homotopijna
$\Sigma(X\times Y)\simeq
\Sigma X\vee\Sigma Y\vee\Sigma(X\wedge Y)$.
W szczególności zawieszenie mapy dołączającej
$S^{p+q-1}\to S^p\vee S^q$ górnej komórki
$S^p\times S^q$, dla $p,q\ge2$, jest nulhomotopijne.
\end{lemat}
\begin{proof}
Mamy kofibrację $X\vee Y\hookrightarrow X\times Y
\to X\wedge Y$.
Po zawieszeniu włączenie
$\Sigma(X\vee Y)\to\Sigma(X\times Y)$
ma lewy homotopijny odwrotny: tworzą go zawieszone
projekcje na $X$ i $Y$ połączone operacją uszczypnięcia
na zawieszeniu. To rozszczepia zawieszoną kofibrację
w kategorii homotopijnej. W otrzymanym rozkładzie
trzeci składnik ma typ $\Sigma(X\wedge Y)$,
gdyż iloraz pary ma ten typ i zawieszone włączenie
jest rozszczepione. Dla dwóch sfer część pochodząca
z górnej komórki jest więc oddzielnym składnikiem
$S^{p+q+1}$, a zawieszona mapa dołączająca jest
nulhomotopijna.
\end{proof}

\begin{twierdzenie}[Pierwszy stabilny pień]
\label{thm:stable-eta-v8}
Niech $\eta:S^3\to S^2$ będzie mapą Hopfa o
$H(\eta)=1$. Wówczas zawieszenie $E\eta$
reprezentuje niezerowy element rzędu $2$ w
$\pi_4(S^3)$, a
$\pi_1^{\mathrm s}\cong\Z/2\Z$.
\end{twierdzenie}
\begin{proof}
Wiemy z poprzednich rozdziałów, że
$\pi_3(S^2)\cong\Z$ z generatorem $\eta$.
Z twierdzenia Freudenthala otrzymujemy surjekcję
$E:\pi_3(S^2)\twoheadrightarrow\pi_4(S^3)$.
Z Twierdzenia~\ref{thm:whitehead-square-v7} mamy
$[\iota_2,\iota_2]=\pm2\eta$.
Lemat~\ref{lem:sigma-product-v8} mówi, że zawieszenie
iloczynu Whiteheada jest zerowe, więc $2E\eta=0$.
Wynika stąd, że $\pi_4(S^3)$ jest albo zerowa,
albo izomorficzna z $\Z/2\Z$.

Załóżmy, że $E\eta=0$. Wtedy stożek mapy
$\Sigma\eta:S^4\to S^3$ ma typ
$S^3\vee S^5$. Z drugiej strony zawieszenie zachowuje
konstrukcję stożka odwzorowania, zatem
$C_{\Sigma\eta}\simeq\Sigma C_\eta
\simeq\Sigma\C P^2$.
Twierdzenie~\ref{thm:sigma-cp2-v8} udowodniło,
że te przestrzenie nie mogą być równoważne:
$\operatorname{Sq}^2$ jest nietrywialne na pierwszej
i zerowe na drugiej. Sprzeczność dowodzi $E\eta\ne0$.
Stąd $\pi_4(S^3)\cong\Z/2\Z$.
Dla $n\ge3$ grupy $\pi_{n+1}(S^n)$ są już w zakresie
stabilnym z Twierdzenia~\ref{thm:stable-range-v8},
więc wyznaczają $\pi_1^{\mathrm s}\cong\Z/2\Z$.
\end{proof}

\begin{uwaga}[Co wykazał ten dowód?]
Wykorzystaliśmy tylko: homotopię fibracji Hopfa,
równanie Whiteheada, rozszczepienie zawieszonego
iloczynu, Freudenthala oraz stabilność kwadratów
Steenroda. Nie potrzebowaliśmy tabel wszystkich
niestabilnych grup homotopii sfer. Najtrudniejszy
zewnętrzny składnik argumentu to stabilność operacji
Steenroda, jawnie wskazana w poprzednim rozdziale.
\end{uwaga}

\section{Stabilne odwzorowania między skończonymi kompleksami}
\begin{definicja}[Klasa stabilna odwzorowania]
Dla skończonych punktowanych CW-kompleksów $X,Y$
przyjmujemy
$[X,Y]_{\mathrm s}:=\varinjlim_{j\ge0}
[\Sigma^jX,\Sigma^jY]_*,$
gdzie $[-,-]_*$ oznacza punktowane klasy homotopii.
Nie zakładamy, że jest to już grupa dla $j=0$;
po dostatecznym zawieszeniu klasy tworzą grupy
przez strukturę ko-H przestrzeni źródłowej.
\end{definicja}

\begin{stwierdzenie}[Stabilizacja klas map]
Niech $X,Y$ będą skończonymi punktowanymi
CW-kompleksami, a $Y$ niech będzie łukowo spójna.
Dla dostatecznie dużych $j$ zawieszenie
$[\Sigma^jX,\Sigma^jY]_*
\to[\Sigma^{j+1}X,\Sigma^{j+1}Y]_*$
jest bijekcją.
\end{stwierdzenie}
\begin{proof}
Po $j$ zawieszeniach przestrzeń $\Sigma^jY$
jest co najmniej $j$-spójna. Względna wersja
twierdzenia Freudenthala (dla odwzorowań z kompleksów
o ustalonym skończonym wymiarze) zapewnia
bijektywność zawieszenia, gdy wymiar $\Sigma^jX$
leży poniżej progu $2j+1$.
Jeżeli $\dim X=d$, to $\dim\Sigma^jX=d+j$,
więc nierówność spełnia się dla dostatecznie dużego $j$.
\end{proof}
\begin{uwaga}
Przywołana względna wersja Freudenthala jest
osobnym twierdzeniem o przestrzeniach odwzorowań
albo konsekwencją twierdzenia o ekscyzji
homotopijnej. W powyższym dowodzie jest wskazana
jawnie; nie wynika z punktowego sprawdzenia
izomorfizmów $\pi_q$ bez argumentu przez szkielety $X$.
\end{uwaga}

\section{Widma: dlaczego stabilizacji nie opisuje pojedyncza przestrzeń?}
\begin{definicja}[Widmo sekwencyjne]
\emph{Widmem sekwencyjnym} nazywamy rodzinę
punktowanych przestrzeni $E_0,E_1,\ldots$ i
odwzorowań struktury $\sigma_n:\Sigma E_n\to E_{n+1}$.
Jeżeli dołączone odwzorowania $E_n\to\Omega E_{n+1}$
są słabymi równoważnościami dla wszystkich $n$,
mówimy o \emph{$\Omega$-widmie}.
Dla $k\in\Z$ stabilne grupy homotopii widma
definiujemy przez
$\pi_k(E)=\varinjlim_{n\ge\max(0,1-k)}\pi_{n+k}(E_n)$,
z mapami indukowanymi przez zawieszenie oraz struktury.
\end{definicja}

\begin{przyklad}[Widmo sferyczne]
Widmo $\mathbb S$ ma $\mathbb S_n=S^n$ i mapy
struktury $\Sigma S^n\cong S^{n+1}$.
Jego grupy homotopii to stabilne pnie sfer:
$\pi_k(\mathbb S)=\pi_k^{\mathrm s}$ dla $k\ge0$.
To widmo nie jest $\Omega$-widmem, gdyż
$S^n\not\simeq\Omega S^{n+1}$ w ogólności.
\end{przyklad}

\begin{przyklad}[Widmo Eilenberga--Mac Lane'a]
Niech $A$ będzie grupą abelową. Wybierzmy
$HA_n=K(A,n)$ dla $n\ge0$, gdzie $K(A,0)$
traktujemy jako przestrzeń dyskretną z wyróżnionym zerem.
Standardowa równoważność
$K(A,n)\simeq\Omega K(A,n+1)$ wyznacza mapy
struktury $\Sigma K(A,n)\to K(A,n+1)$.
Wtedy $HA$ jest $\Omega$-widmem w sensie słabej równoważności,
a jego grupy stabilne wynoszą
$\pi_k(HA)=A$ dla $k=0$ i $0$ dla $k\ne0$.
Dla $k=0$ otrzymujemy na każdym poziomie $n\ge1$
$\pi_n(K(A,n))\cong A$, zaś dla $k\ne0$
wszystkie późniejsze grupy z definicji $K(A,n)$ są zerowe.
\end{przyklad}

\begin{twierdzenie}[Kohomologia jako stabilna teoria]
\label{thm:HA-rep-v8}
Dla punktowanego CW-kompleksu $X$, grupy abelowej $A$
i liczby całkowitej $m\ge1$ zachodzi naturalna bijekcja
$\widetilde H^m(X;A)\cong[X,K(A,m)]_*$.
Po zawieszeniu otrzymujemy zgodne identyfikacje
$\widetilde H^{m+j}(\Sigma^jX;A)
\cong[\Sigma^jX,K(A,m+j)]_*$.
Wszystkie te identyfikacje są zgodne z izomorfizmem
zawieszenia w zredukowanej kohomologii.
\end{twierdzenie}
\begin{proof}
Pierwsza bijekcja jest twierdzeniem o reprezentowalności
kohomologii przez $K(A,m)$ z wcześniejszego rozdziału.
Zredukowany izomorfizm zawieszenia wynika z długiego
ciągu dokładnego pary $(CX,X)$ i ściągalności stożka.
Po stronie przestrzeni Eilenberga--Mac Lane'a
przechodzi on przez adjunkcję $\Sigma\dashv\Omega$ i
słabą równoważność
$K(A,m)\simeq\Omega K(A,m+1)$.
Naturalność wynika z naturalności tych dwóch konstrukcji.
\end{proof}

\begin{uwaga}[Od kohomologii do teorii reprezentowanych]
Pełne stwierdzenie spektralne ma postać
$\widetilde H^m(X;A)\cong[\Sigma^\infty X,\Sigma^mHA]$
w stabilnej kategorii homotopii. Tutaj
$\Sigma^\infty X$ oznacza \emph{widmo zawieszeń}
(o $n$-tym poziomie $\Sigma^nX$), a nawias
$[E,F]$ zbiór stabilnych klas odwzorowań widm. Jego dowód wymaga
konstrukcji odwzorowań widm, przejścia do modeli
homotopijnych oraz twierdzenia o reprezentowalności.
W obecnym skrypcie podajemy dokładną wersję dla
$K(A,m)$, a zapis spektralny wskazuje kierunek
dalszego uogólnienia.
\end{uwaga}

\section{Czego nie można odczytać z samej tabeli grup?}
Znajomość wszystkich grup homologii $H_*(X)$ nie
wyznacza stabilnego typu homotopijnego: zawieszona
$\C P^2$ i bukiet sfer mają identyczne homologie,
ale różnią się działaniem $\operatorname{Sq}^2$.
Podobnie same grupy homotopii mogą nie wyznaczać
niestabilnego typu homotopijnego: potrzebne są
niezmienniki Postnikowa.

Kolejnymi krokami są ciąg spektralny Adamsa,
wykorzystujący algebrę Steenroda i funktory $\operatorname{Ext}$,
oraz twierdzenie Pontriagina--Thoma wiążące
stabilną homotopię z kobordyzmami. Dopiero po ich
przygotowaniu sensowne jest systematyczne
obliczanie dalszych stabilnych pni sfer.

% Iteracja IX: odrębny rozdział do wstawienia tuż przed \appendix.
\chapter{Konstrukcja Pontriagina--Thoma i podstawy teorii kobordyzmów}
\label{ch:pt-v9}

\section{Motywacja: jak odwzorowanie w sferę koduje podrozmaitość?}
W poprzednim rozdziale zdefiniowaliśmy stabilne grupy homotopii sfer
za pomocą zawieszeń. Ich elementy mają interpretację geometryczną:
można je reprezentować przez zwarte podrozmaitości o \emph{uramowanej}
wiązce normalnej. Przejście od podrozmaitości do odwzorowania
polega na ściągnięciu do jednego punktu wszystkiego poza jej
otoczeniem tubularnym. Jest to \emph{konstrukcja Pontriagina--Thoma}.
Po zastąpieniu uramowania dowolną wiązką normalną otrzymujemy
przestrzenie Thoma i widma $MO$, $MSO$; służą one do opisu
kobordyzmów nieorientowanych i zorientowanych. Odróżniamy przy tym
trzy poziomy: twierdzenie dla uramowanych podrozmaitości,
jego wersję dla wiązek normalnych i twierdzenia obliczeniowe
o pierścieniach kobordyzmu.

Pracujemy w kategorii \emph{gładkich}, Hausdorffowskich,
drugoprzeliczalnych rozmaitości. Rozmaitość \emph{zamknięta}
oznacza tutaj zwartą rozmaitość bez brzegu; dopuszczamy
rozmaitość pustą. Podstawy o podrozmaitościach, transwersalności
i otoczeniach tubularnych przypomniano w dodatku. Rozkłady
kobordyzmów na uchwyty są przedmiotem osobnego skryptu z geometrii
różniczkowej; tutaj badamy przede wszystkim \emph{niezmienniki}
kobordyzmu.

\section{Włókno wartości regularnej i uramowanie normalne}
\begin{definicja}[Wiązka normalna i uramowanie]
Niech $M^n\subseteq Q^{n+k}$ będzie domkniętą, gładko
zanurzoną podrozmaitością, $k\ge1$. Jej \emph{wiązką normalną}
nazywamy rzeczywistą wiązkę wektorową rzędu $k$,
$\nu_{M/Q}:=TQ|_M/TM$. Po wyborze metryki można ją
utożsamić z prostopadłym dopełnieniem $TM$ w $TQ|_M$.
\emph{Uramowaniem normalnym} nazywamy izomorfizm wiązek
$\varphi:\nu_{M/Q}\to M\times\R^k$, czyli ciągłą
(gładką) rodzinę uporządkowanych baz we włóknach.
\emph{Podrozmaitość uramowana} oznacza parę $(M,\varphi)$.
Zmiana uramowania nie musi zachowywać jego klasy kobordyzmu.
\end{definicja}

\begin{lemat}[Wartość regularna i uramowanie]
\label{lem:pt-regular-v9}
Niech $Q^{n+k}$ będzie zwartą gładką rozmaitością,
$f:Q\to S^k$ odwzorowaniem gładkim, a $q\in S^k$
wartością regularną $f$. Wtedy $M=f^{-1}(q)$
jest zamkniętą podrozmaitością wymiaru $n$.
Po wyborze uporządkowanej bazy przestrzeni $T_qS^k$
różniczka $f$ wyznacza kanoniczne \emph{względem tego
wyboru} uramowanie $\nu_{M/Q}$.
\end{lemat}
\begin{proof}
Dla $x\in M$ surjektywność
$\D f_x:T_xQ\to T_qS^k$ daje, z twierdzenia o wartości
regularnej, gładką podrozmaitość $M$ z
$T_xM=\ker\D f_x$. Zatem $\D f_x$ indukuje
izomorfizm ilorazów
$T_xQ/T_xM\to T_qS^k$; izomorfizmy zależą gładko od
$x$ i składają się na izomorfizm wiązki
$\nu_{M/Q}\cong M\times T_qS^k$.
Ustalona baza $T_qS^k\cong\R^k$ kończy konstrukcję.
Zwartość $Q$ zapewnia zwartość zbioru $M$, a jego
domkniętość wynika z ciągłości $f$.
\end{proof}

\begin{uwaga}[Przedmiot uramowania]
Uramowujemy \emph{wiązkę normalną}, niekoniecznie
wiązkę styczną. Istnienie stabilnego uramowania normalnego
jest równoważne istnieniu stabilnego uramowania stycznego,
lecz konkretne uramowania należy porównać za pomocą relacji
$TM\oplus\nu_{M/\R^{n+k}}\cong\underline\R^{n+k}$.
W przestrzeni $S^{n+k}$ używamy wykresu
$S^{n+k}\setminus\{\infty\}\cong\R^{n+k}$;
w ten sposób ramowanie normalne nie zależy od trywialności
całej wiązki stycznej sfery.
\end{uwaga}

\section{Przestrzeń Thoma i odwzorowanie ściągające}
\begin{definicja}[Przestrzeń Thoma]
Niech $\xi\to B$ będzie rzeczywistą wiązką wektorową
rzędu $k$ z metryką włóknową. Oznaczmy przez $D(\xi)$
odpowiednio wiązkę dysków jednostkowych, a przez $S(\xi)$
wiązkę sfer jednostkowych. \emph{Przestrzenią Thoma} jest
przestrzeń punktowana
$\operatorname{Th}(\xi):=D(\xi)/S(\xi)$,
w której całą $S(\xi)$ utożsamiamy do jednego punktu.
Dla wiązki trywialnej
$\underline\R^k_B\to B$ mamy naturalny homeomorfizm
$\operatorname{Th}(\underline\R^k_B)\cong B_+\wedge S^k$;
$B_+$ oznacza $B$ z dodanym rozłącznym punktem bazowym.
\end{definicja}

\begin{lemat}[Odwzorowanie ściągające]
\label{lem:pt-collapse-v9}
Niech $i:M^n\hookrightarrow Q^{n+k}$ będzie gładkim
osadzeniem zwartej podrozmaitości we wnętrzu zwartej
rozmaitości $Q$ (ewentualnie z brzegiem), z
$M\cap\partial Q=\varnothing$. Wybierzmy otoczenie tubularne
$U\cong \operatorname{int}D(\nu_{M/Q})$ z domknięciem
we wnętrzu $Q$. Wtedy istnieje ciągłe punktowane
odwzorowanie
\[
c_{M,Q}: Q/\partial Q\longrightarrow
\operatorname{Th}(\nu_{M/Q}),
\]
przypisujące wszystkim punktom spoza $U$ punkt bazowy,
a punktom w $U$ ich współrzędne tubularne. Jeśli
$\partial Q=\varnothing$, przez $Q/\partial Q$ rozumiemy
$Q_+$ z oddzielnym punktem bazowym. Klasa homotopii
odwzorowania nie zależy od dostatecznie małego
wyboru otoczenia ani od metryki tubularnej.
\end{lemat}
\begin{proof}
Wybierzmy identyfikację domkniętego otoczenia
tubularnego z $D(\nu)$. Na jego wnętrzu mapę określa
rzut do ilorazu $D(\nu)/S(\nu)$; na dopełnieniu
wnętrza przyjmujemy punkt bazowy. Na brzegu tubularnym
obie formuły dają ten sam punkt, zatem z lematu o
sklejaniu odwzorowań otrzymujemy ciągłość.
Otoczenie można wybrać poza $\partial Q$, więc
odwzorowanie jest stałe na $\partial Q$ i przechodzi
na wskazany iloraz. Dwa dostatecznie małe wybory
porównujemy po zmniejszeniu do wspólnego otoczenia;
rodzina promieniowych przeskalowań włókien daje
homotopię odwzorowań ściągających. Każde dwa
otoczenia tubularne są izotopijne po zwężeniu,
co jest standardową częścią twierdzenia o otoczeniach
tubularnych. To uzasadnia niezależność klasy homotopii.
\end{proof}

Jeśli $\nu_{M/Q}$ ma uramowanie $\varphi$, otrzymujemy
złożenie
$c_{M,Q}:Q_+\to\operatorname{Th}(\nu_{M/Q})
\cong M_+\wedge S^k\to S^k$, gdzie ostatnia strzałka
zapomina punkt $M$ (na współczynniku $M_+$).
Dla $Q=S^{n+k}$ z wybranym punktem $\infty$ poza $M$
uzyskujemy mapę punktowaną $S^{n+k}\to S^k$.

\begin{center}
\begin{tikzpicture}[>=Stealth,font=\small]
\begin{scope}[xshift=0cm]
\draw[thick] (0,0) ellipse (1.50 and 0.91);
\draw[manifoldblue,thick,fill=manifoldblue!10] (0,0) ellipse (1.12 and 0.46);
\draw[manifoldred,very thick] (-0.74,-0.02) .. controls (-0.28,0.27) and (0.18,-0.24) .. (0.76,0.06);
\node[font=\footnotesize] at (0,1.20) {$S^{n+k}$};
\node[font=\footnotesize,manifoldblue] at (0,-0.70) {$U\cong D(\nu)$};
\node[font=\footnotesize,manifoldred] at (0,0.33) {$M$};
\end{scope}
\draw[->,thick] (1.65,0)--node[above,align=center,font=\footnotesize] {ściągnięcie} (4.25,0);
\begin{scope}[xshift=5.85cm]
\draw[thick,fill=manifoldblue!8] (0,0) circle (0.92);
\draw[thick,dashed] (0,0) circle (0.52);
\fill (0.88,0.14) circle (1.5pt);
\node[font=\footnotesize,anchor=west] at (1.00,0.12) {$*$};
\node[font=\footnotesize] at (0,1.20) {$\operatorname{Th}(\nu)$};
\node[font=\footnotesize] at (0,-1.18) {brzeg włókien $\mapsto *$};
\end{scope}
\end{tikzpicture}
\end{center}
\begin{center}\small Schemat idei: wnętrze rurki opisuje współrzędne normalne,
zaś jej zewnętrze zostaje utożsamione z punktem bazowym.\end{center}

\section{Twierdzenie Pontriagina dla podrozmaitości uramowanych}
\begin{definicja}[Kobordyzm uramowany w sferze]
Ustalmy $n\ge0$, $k\ge1$. Dwa zamknięte uramowane
$n$-wymiarowe podrozmaitości $M_0,M_1$ w
$S^{n+k}\setminus\{\infty\}$ są \emph{uramowanie kobordantne}
(oznaczamy $M_0\sim_{\mathrm{fr}}M_1$), jeśli istnieje
zwarta $(n+1)$-wymiarowa podrozmaitość
$W\subset(S^{n+k}\setminus\{\infty\})\times[0,1]$
prostopadła do końcowych poziomów w otoczeniach brzegów,
taka że
$\partial W=M_0\times\{0\}\sqcup M_1\times\{1\}$,
a jej wiązka normalna rzędu $k$ ma uramowanie
ograniczające się do danych uramowań końcowych.
Utożsamiamy klasy różniące się izotopią uramowaną.
Zbiór tych klas oznaczamy $\mathcal F_n^k$.
\end{definicja}

\begin{twierdzenie}[Pontriagin--Thom w ustalonym kowymiarze]
\label{thm:pt-framed-v9}
Dla $n\ge0$ i $k\ge1$ konstrukcja ściągająca
wyznacza bijekcję
\[
\operatorname{PT}:\mathcal F_n^k\xrightarrow{\cong}
[S^{n+k},S^k]_*=\pi_{n+k}(S^k).
\]
Jest zgodna ze zwiększaniem kowymiaru przez dodanie
trywialnego kierunku normalnego i z zawieszeniem
odwzorowań. W zakresie, w którym grupy homotopii
są abelowe, suma odpowiada rozłącznej sumie
uramowanych podrozmaitości umieszczonych w rozłącznych
kulach.
\end{twierdzenie}
\begin{proof}[Dowód z użyciem transwersalności i twierdzenia tubularnego]
\emph{Krok 1: od podrozmaitości do odwzorowania.}
Na mocy Lematu~\ref{lem:pt-collapse-v9} otoczenie tubularne
$M$ daje mapę $S^{n+k}\to\operatorname{Th}(\nu_{M/S^{n+k}})$.
Uramowanie utożsamia tę przestrzeń z $M_+\wedge S^k$;
po zapomnieniu parametru $M$ otrzymujemy mapę do $S^k$.
Uramowany kobordyzm $W$ ma otoczenie tubularne w
$S^{n+k}\times[0,1]$, które można wybrać produktowe
przy końcach przedziału. Jego ściągnięcie jest zatem
homotopią między mapami obu brzegów.
Wynik zależy tylko od klasy kobordyzmu.

\emph{Krok 2: od mapy do podrozmaitości.}
Niech $f:S^{n+k}\to S^k$ będzie odwzorowaniem
punktowanym. Z twierdzenia o aproksymacji gładkiej
możemy zastąpić je homotopijnym odwzorowaniem gładkim,
stałym blisko $\infty$. Wybierzmy wartość regularną
$q\ne*$ (Sard); jeśli $q$ nie należy do obrazu $f$,
wówczas $M=\varnothing$, a $f$ jest nullhomotopijne, gdyż
jego obraz mieści się w $S^k\setminus\{q\}\cong\R^k$.
W pozostałym przypadku $M=f^{-1}(q)$ jest niepustą, zwartą
$n$-podrozmaitością poza $\infty$.
Lemat~\ref{lem:pt-regular-v9} daje uramowanie normalne
przez wybór dodatnio zorientowanej bazy $T_qS^k$.
Jeśli $f_0,f_1$ są homotopijne przez homotopię punktowaną,
przybliżamy homotopię względnie do jej końców,
aby była poprzeczna do $q$. Przeciwobraz $q$
w $S^{n+k}\times[0,1]$ jest wtedy uramowanym
kobordyzmem z iloczynowym zachowaniem końców.
Tak powstaje odwzorowanie w stronę odwrotną.

\emph{Krok 3: odwrotność obu konstrukcji.}
Dla podrozmaitości z otoczeniem tubularnym mapa
ściągająca ma jako przeciwobraz punktu $q$ dokładnie
zerowy przekrój wiązki normalnej; różniczka w kierunku
normalnym odtwarza zadane uramowanie. W drugą stronę
wybierzmy małą domkniętą kulę $D$ wokół wartości
regularnej $q$, nie zawierającą punktu bazowego.
Po jej zmniejszeniu $f^{-1}(D)\to D$ jest właściwą
submersją (dziedzina i sfera są zwarte, a wszystkie
punkty blisko $q$ są regularne). Twierdzenie Ehresmanna (dla właściwej submersji,
z wersją nad brzegiem dysku) daje lokalnie trywialną
wiązkę nad $D$. Ponieważ $D$ jest ściągalny,
otrzymujemy jej trywializację $f^{-1}(D)\cong M\times D$.
Wybierając trywializację, która wzdłuż włókna centralnego
ma różniczkę normalną zgodną z $\mathrm{D}f$, uzyskujemy
identyfikację z otoczeniem tubularnym zachowującą
uramowanie po przejściu do izotopii małych dysków.
Ten ostatni krok nie wynika z samej zwartości: korzysta
z jednoznaczności otoczeń tubularnych z dokładnością
do izotopii.
Wybierzmy odwzorowanie stopnia $1$,
$S^k\to D/\partial D\cong S^k$, które wysyła
do punktu bazowego dopełnienie wnętrza $D$.
Jest ono homotopijne do identyczności (po naturalnym
utożsamieniu orientacji). Złożenie z $f$ jest więc
homotopijne do $f$, a w wybranej trywializacji
jest właśnie mapą ściągającą otoczenie tubularne.
Obie konstrukcje są wzajemnie odwrotne.

\emph{Krok 4: zawieszenie i dodawanie.}
Nowy trywialny kierunek normalny zastępuje
$D^k/S^{k-1}$ przez $D^{k+1}/S^k$;
odpowiada to zawieszeniu mapy ściągającej.
Dla dwóch rozłącznych osadzeń w kulach mapy
ściągające sumują się poprzez mapę uszczypnięcia
$S^{n+k}\to S^{n+k}\vee S^{n+k}$. To daje
zgodność z operacjami grupowymi.
\end{proof}

\begin{przyklad}[Stopień mapy jako liczba punktów]
Niech $f:S^1\to S^1$ będzie gładkim odwzorowaniem
stopnia $d$, a $q\in S^1$ jego wartością regularną.
Włókno $f^{-1}(q)$ jest skończonym zbiorem punktów.
Uramowanie normalne w każdym punkcie odpowiada znakowi
lokalnej pochodnej w zorientowanych współrzędnych.
Po zsumowaniu ze znakami otrzymujemy
$d=\sum_{x\in f^{-1}(q)}\operatorname{sgn}\D f_x$.
Tak stopień, wcześniej definiowany przez homologię,
staje się liczbą algebraiczną punktów uramowanej
zerowymiarowej podrozmaitości.
\end{przyklad}

\begin{wniosek}[Homotopia stabilna i kobordyzm uramowany]
\label{cor:pt-stable-v9}
Niech $\Omega_n^{\mathrm{fr}}:=\varinjlim_k\mathcal F_n^k$,
gdzie mapy przejścia dodają jeden trywialny kierunek
normalny. Dla każdego $n\ge0$ istnieje naturalny
izomorfizm grup
$\Omega_n^{\mathrm{fr}}\cong\pi_n^{\mathrm s}$.
\end{wniosek}
\begin{proof}
Twierdzenie~\ref{thm:pt-framed-v9} identyfikuje wszystkie
poziomy tego układu z $\pi_{n+k}(S^k)$,
a mapy przejścia z zawieszeniami. Po wzięciu
tej samej granicy prostej otrzymujemy dokładnie
definicję stabilnego pnia z poprzedniego rozdziału.
\end{proof}

\begin{przyklad}[Włókno Hopfa i pierwszy stabilny pień]
Odwzorowanie Hopfa $\eta:S^3\to S^2$ ma włókno
regularne będące okręgiem $S^1\subset S^3$, z
uramowaniem normalnym pochodzącym od $\D\eta$.
Z poprzedniego rozdziału wiemy, że $\Sigma\eta$
reprezentuje niezerowy element $\pi_1^{\mathrm s}\cong\Z/2\Z$.
Dlatego otrzymany uramowany okrąg nie jest brzegiem
uramowanego kobordyzmu, choć sam okrąg bez dodatkowej
struktury jest brzegiem dysku. Uramowanie zawiera
informację niewidoczną dla zwykłej homologii.
\end{przyklad}

\section{Związek konstrukcji z dualnością Poincarégo i klasą Eulera}
Pokażemy, dlaczego przestrzeń Thoma jest właściwym celem
odwzorowania ściągającego również wtedy, gdy normalna
wiązka nie jest trywialna. Jest to zarazem powrót do
dualności Poincarégo i teorii przecięć z wcześniejszych rozdziałów.

\begin{twierdzenie}[Ściągnięcie klasy Thoma daje dualność podrozmaitości]
\label{thm:pt-pd-v9}
Niech $Q^m$ będzie zamkniętą, spójną, zorientowaną
rozmaitością gładką, a $M^{m-k}\subset Q$ zamkniętą,
zorientowaną podrozmaitością bez brzegu i o kowymiarze
$k\ge1$. Wyposażmy $\nu=\nu_{M/Q}$ w orientację
wyznaczoną z orientacji $Q$ i $M$.
Niech $U\in\widetilde H^k(\operatorname{Th}(\nu);\Z)$
będzie całkowitą klasą Thoma, z normalizacją lokalną
zgodną z tą orientacją. Wówczas
\[
c_{M,Q}^*U=\operatorname{PD}_Q(i_*[M])
\quad\text{w }H^k(Q;\Z),
\]
gdzie $i:M\hookrightarrow Q$ jest włączeniem, a
$\operatorname{PD}_Q:H_{m-k}(Q;\Z)\to H^k(Q;\Z)$
izomorfizmem dualności Poincarégo.
Bez założeń orientowalności analogiczny wzór zachodzi
ze współczynnikami $\F_2$.
\end{twierdzenie}
\begin{proof}
Niech $\sigma:\Delta^k\to Q$ będzie gładkim
sympleksem singularnym poprzecznym do $M$ i mającym
brzeg rozłączny z $M$. Jego obraz przecina $M$ w
skończenie wielu punktach: transwersalność daje
zbiór dyskretny, a zwartość dziedziny zapewnia
skończoność. Przy każdym punkcie przecięcia możemy
wybrać współrzędne $(y,z)\in\R^{m-k}\times\R^k$
na $Q$ z $M=\{z=0\}$. Klasa Thoma w tych
współrzędnych jest generatorem względnej kohomologii
$H^k(D^k,S^{k-1};\Z)$. Wartość
$\langle c_{M,Q}^*U,\sigma\rangle$ jest więc
sumą lokalnych stopni składowej normalnej
$z\circ\sigma$ przy jej zerach; są to dokładnie
znaki poprzecznych przecięć $\sigma$ z $M$.
Wyjaśnijmy również, dlaczego otrzymujemy równość
\emph{klas całkowitych}, a nie tylko jednakowe
sparowania z cyklami wolnymi od torsji. Klasa Thoma $U$
przez ekscyzję odpowiada klasie względnej
$u\in H^k(Q,Q\setminus M;\Z)$. W łańcuchowej
konstrukcji iloczynu zewnętrznego i cap-produktu
z klasą fundamentalną lokalny rachunek ma postać
\[
[D^{m-k}\times D^k]\frown
\operatorname{pr}_{D^k}^*U_k
=[D^{m-k}\times\{0\}],
\]
gdzie $U_k$ oznacza dodatni generator
$H^k(D^k,S^{k-1};\Z)$, a klasy po obu stronach
rozumiemy w odpowiednich grupach względnych.
Istotnie, na podziale symplicjalnym zgodnym z
otoczeniem tubularnym jedynie sympleksy przecinające
$M$ wnoszą niezerowy wkład do iloczynu cap;
lokalny wkład wynosi właśnie powyższy generator.
Po zsumowaniu lokalnych fragmentów, przejściu przez
ekscyzję i zapomnieniu podpory w $Q\setminus M$
dostajemy tożsamość \emph{w homologii}
$[Q]\frown c_{M,Q}^*U=i_*[M]$.
Ponieważ $[Q]\frown(-)$ jest izomorfizmem dualności
Poincarégo, otrzymujemy żądaną równość klas
kohomologii również w obecności torsji. Współczynniki
$\F_2$ usuwają potrzebę wyboru orientacji i znaków.
\end{proof}

\begin{wniosek}[Zera przekroju i klasa Eulera]
\label{cor:euler-zero-v9}
Niech $\xi\to Q$ będzie zorientowaną rzeczywistą wiązką
wektorową rzędu $k\ge1$ nad zamkniętą, zorientowaną
rozmaitością $Q^m$. Jeżeli gładki przekrój
$s:Q\to\xi$ jest poprzeczny do przekroju zerowego,
to $Z(s)=\{x:s(x)=0\}$ jest zamkniętą, zorientowaną
podrozmaitością wymiaru $m-k$ i
$\operatorname{PD}_Q([Z(s)])=e(\xi)$ w $H^k(Q;\Z)$.
\end{wniosek}
\begin{proof}
Transwersalność przekroju do przekroju zerowego
daje podrozmaitość wymiaru $m-k$, a pochodna przekroju
w kierunkach normalnych identyfikuje
$\nu_{Z(s)/Q}$ z $\xi|_{Z(s)}$ i przenosi orientację.
Definicja klasy Eulera jako przeciwobrazu klasy Thoma
przez przekrój zerowy i homotopia $t\mapsto t s$
utożsamiają ją z lokalną klasą zer przekroju.
Twierdzenie~\ref{thm:pt-pd-v9} identyfikuje tę klasę
z dualną Poincarégo do $[Z(s)]$. Jeśli $k>m$,
transwersalny przekrój nie ma zer; wówczas obie strony
są zerowe, ponieważ $H^k(Q;\Z)=0$.
\end{proof}

\begin{przyklad}[Dlaczego znak ma znaczenie]
Dla zorientowanej wiązki liniowej nad okręgiem
każdy przekrój poprzeczny do zera ma zera o
lokalnych indeksach $+1$ lub $-1$, których suma
jest zerowa. Odpowiada to $H^1(S^1;\Z)\cong\Z$,
ale klasa Eulera zorientowanej rzeczywistej wiązki
rzędu $1$ jest zerowa: orientowana wiązka liniowa
jest trywialna. Bez orientacji można natomiast otrzymać
wiązki Möbiusa i niezerową klasę $w_1$.
\end{przyklad}

\section{Kobordyzmy nieorientowane i zorientowane}
\begin{definicja}[Kobordyzm gładki]
Niech $M_0,M_1$ będą zamkniętymi gładkimi
$n$-rozmaitościami. \emph{Kobordyzmem (nieorientowanym)}
między nimi nazywamy zwartą gładką
$(n+1)$-rozmaitość $W$ wraz z dyfeomorfizmem
$\partial W\cong M_0\sqcup M_1$.
W kategorii zorientowanej wymagamy orientacji wszystkich
obiektów i izomorfizmu zorientowanych brzegów
$\partial W\cong (-M_0)\sqcup M_1$, gdzie
$-M_0$ ma odwróconą orientację, a $\partial W$ ma
orientację indukowaną konwencją normalnej zewnętrznej.
Zbiory klas oznaczamy odpowiednio
$\Omega_n^O$ i $\Omega_n^{SO}$. Dodawanie klas
jest zadane sumą rozłączną, a mnożenie produktem
kartezjańskim rozmaitości.
\end{definicja}

\begin{stwierdzenie}[Struktury algebraiczne]
\label{prop:bordism-group-v9}
Dla każdego $n\ge0$ zbiory $\Omega_n^O$ i
$\Omega_n^{SO}$ są grupami abelowymi. W grupie
nieorientowanej zachodzi $2[M]=0$, a w grupie
zorientowanej $[-M]=-[M]$. Sumy
$\Omega_*^O:=\bigoplus_{n\ge0}\Omega_n^O$
i $\Omega_*^{SO}:=\bigoplus_{n\ge0}\Omega_n^{SO}$
są pierścieniami stopniowanymi, przy czym
w przypadku zorientowanym mnożenie jest
przemienne z regułą znaku $(-1)^{mn}$ w stopniach
$m,n$.
\end{stwierdzenie}
\begin{proof}
Sumowanie rozłączne jest łączne i przemienne z
dokładnością do dyfeomorfizmu; pusty zbiór jest
zerem. Walec $M\times[0,1]$ ma dwa egzemplarze
$M$ jako brzeg, co daje $[M]+[M]=0$ bez orientacji.
Z orientacją brzeg walca ma postać $(-M)\sqcup M$,
więc $[-M]+[M]=0$. Jeśli $M\sim M'$ przez $W$,
to dla zamkniętego $N$ produkt $W\times N$ jest
kobordyzmem między $M\times N$ i $M'\times N$.
Zatem mnożenie nie zależy od reprezentantów.
Zamiana czynników $M^m\times N^n\to N^n\times M^m$
zmienia orientację o $(-1)^{mn}$, co wyjaśnia
regułę przemienności. Pozostałe aksjomaty wynikają
z łączności i rozdzielności produktu względem
sumy rozłącznej.
\end{proof}

\begin{przyklad}[Pierwsze grupy kobordyzmu]
\label{ex:bordism-low-v9}
Wymiar zero daje $\Omega_0^O\cong\Z/2\Z$
(parzystość liczby punktów) oraz
$\Omega_0^{SO}\cong\Z$ (liczba punktów zorientowanych,
ze znakami). Każda zwarta $1$-rozmaitość ma parzystą
liczbę punktów brzegowych, bo jej składowe z brzegiem są
odcinkami. I odwrotnie, parę punktów można połączyć
odcinkiem, zorientowanym od punktu ujemnego do dodatniego
w przypadku zorientowanym. To dowodzi obu izomorfizmów. Każda zwarta rozmaitość jednowymiarowa
bez brzegu jest sumą okręgów, a okrąg jest brzegiem
dysku, więc $\Omega_1^O=\Omega_1^{SO}=0$.
Każda zamknięta zorientowana powierzchnia jest brzegiem
trójwymiarowego ciała z uchwytami, dlatego
$\Omega_2^{SO}=0$. Z klasyfikacji zamkniętych
powierzchni otrzymamy dalej
$\Omega_2^O\cong\Z/2\Z$, generowaną przez
$\R P^2$. W dwóch ostatnich wnioskach używamy
klasyfikacji powierzchni jako wyniku znanego z
topologii rozmaitości, a nie twierdzenia dowiedzionego
tutaj od początku.
\end{przyklad}

\section{Klasy charakterystyczne jako przeszkody kobordyzmu}
\begin{definicja}[Liczba Stiefela--Whitneya]
Niech $M^n$ będzie zamkniętą gładką $n$-rozmaitością,
a $w_j(TM)\in H^j(M;\F_2)$ klasami Stiefela--Whitneya.
Dla nieujemnych indeksów $i_1,\ldots,i_s$ o sumie
$i_1+\cdots+i_s=n$ liczbę
\[
\bigl\langle w_{i_1}(TM)\smile\cdots\smile
w_{i_s}(TM),[M]_2\bigr\rangle\in\F_2
\]
nazywamy \emph{liczbą Stiefela--Whitneya}.
Tutaj $[M]_2\in H_n(M;\F_2)$ oznacza klasę fundamentalną;
nie wymaga ona orientowalności $M$.
\end{definicja}

\begin{twierdzenie}[Zanik liczb charakterystycznych na brzegu]
\label{thm:sw-bdy-v9}
Jeśli zamknięta $n$-rozmaitość $M$ jest dyfeomorficzna
z brzegiem zwartej $(n+1)$-rozmaitości $W$, to
wszystkie liczby Stiefela--Whitneya $M$ są zerowe.
Jeżeli $M$ i $W$ są zorientowane, analogicznie
zanikają liczby Pontriagina złożone z klas
$p_j(TM)\in H^{4j}(M;\Z)$ o łącznym stopniu $n$.
\end{twierdzenie}
\begin{proof}
Niech $i:M\hookrightarrow W$ oznacza włączenie brzegu.
Z twierdzenia o kołnierzu otrzymujemy
$TW|_M\cong TM\oplus\underline\R$.
Stabilność klas charakterystycznych daje
$i^*w_j(TW)=w_j(TM)$ oraz analogicznie dla $p_j$.
Jeżeli $P$ jest dowolnym jednomianem charakterystycznym
stopnia $n$, to
$\langle P(TM),[M]\rangle=
\langle i^*P(TW),[M]\rangle=
\langle P(TW),i_*[M]\rangle$.
Ale $i_*[M]=0$ w $H_n(W;\F_2)$:
wynika to z długiego ciągu homologii pary $(W,M)$,
ponieważ $[M]$ jest obrazem względnej klasy fundamentalnej
$[W,M]\in H_{n+1}(W,M;\F_2)$ przez homomorfizm brzegowy.
W przypadku zorientowanym ten sam argument prowadzimy
ze współczynnikami $\Z$, używając zgodności
orientacji brzegu. Każde rozważane sparowanie wynosi zero.
\end{proof}

\begin{przyklad}[Dlaczego $\R P^2$ nie jest brzegiem]
Niech $t\in H^1(\R P^2;\F_2)$ będzie generatorem.
W punkcie odpowiadającym prostej $\ell\subset\R^{n+1}$
przestrzeń styczna $T_{[\ell]}\R P^n$ utożsamia się
z $\operatorname{Hom}(\ell,\ell^\perp)$. Z rozkładu
$\R^{n+1}=\ell\oplus\ell^\perp$ wynika izomorfizm
wiązek
$\underline\R\oplus T\R P^n
\cong\operatorname{Hom}(\gamma^1,\underline\R^{n+1})
\cong(n+1)(\gamma^1)^*$.
Metryka euklidesowa utożsamia $(\gamma^1)^*$ z
$\gamma^1$. Stąd
$T\R P^n\oplus\underline\R\cong(n+1)\gamma^1$,
a więc $w(T\R P^2)=(1+t)^3=1+t+t^2$.
Stąd $\langle w_1^2,[\R P^2]_2\rangle=1$ i
$\langle w_2,[\R P^2]_2\rangle=1$, ponieważ
$t^2$ generuje $H^2(\R P^2;\F_2)$.
Na mocy Twierdzenia~\ref{thm:sw-bdy-v9}
$\R P^2$ nie jest brzegiem zwartej $3$-rozmaitości.
Warto przy tym pamiętać, że $\R P^2$ nie jest
orientowalna: argument dotyczy \emph{nieorientowanego}
kobordyzmu.
\end{przyklad}

\begin{twierdzenie}[Obliczenie $\Omega_2^O$]
\label{thm:omega2-v9}
Istnieje izomorfizm
$\Omega_2^O\cong\Z/2\Z$ zadany przez
$[M]\mapsto\chi(M)\bmod2$; generatorem jest
klasa $\R P^2$.
\end{twierdzenie}
\begin{proof}
Dla zamkniętej powierzchni $M$ równość
$\langle w_2(TM),[M]_2\rangle=\chi(M)\bmod2$
wynika z twierdzenia Poincarégo--Hopfa modulo $2$.
Na mocy Twierdzenia~\ref{thm:sw-bdy-v9} odwzorowanie
jest więc niezmiennikiem kobordyzmu.
Każda zamknięta spójna powierzchnia jest orientowalną
powierzchnią $\Sigma_g$ albo sumą spójną $h$ egzemplarzy
$\R P^2$. Powierzchnie $\Sigma_g$ są brzegami ciał
z uchwytami. Z kolei weźmy
$(M\sqcup N)\times[0,1]$ i dołączmy do jego górnej
części trójwymiarowy uchwyt $D^1\times D^2$
wzdłuż dwóch rozłącznych dysków $D^2$,
po jednym na każdej powierzchni. Nowa górna
powierzchnia jest sumą spójną $M\#N$,
dolna pozostaje $M\sqcup N$. W ten sposób
otrzymujemy kobordyzm dowodzący
$[M\#N]=[M]+[N]$ w $\Omega_2^O$.
Wynika stąd, że każda klasa jest wielokrotnością
$[\R P^2]$. Równanie $2[\R P^2]=0$ pochodzi
ze Stwierdzenia~\ref{prop:bordism-group-v9}, a liczba
$\chi(\R P^2)\bmod2=1$ pokazuje, że generator
jest niezerowy. Dostajemy żądany izomorfizm.
\end{proof}

\begin{uwaga}[Zakres twierdzenia Thoma o liczbach charakterystycznych]
W dalszej teorii dowodzi się głębszego twierdzenia:
\emph{dwie zamknięte gładkie $n$-rozmaitości są
nieorientowanie kobordantne wtedy i tylko wtedy, gdy
wszystkie odpowiadające im liczby Stiefela--Whitneya
są równe}. Kierunek konieczny udowodniliśmy powyżej;
kierunek dostateczny opiera się na obliczeniu pierścienia
$\Omega_*^O$ metodami Thoma i nie jest tu dowiedziony.
\end{uwaga}

\section{Od uramowania do widm \texorpdfstring{$MO$ i $MSO$}{MO i MSO}}
\begin{definicja}[Przestrzenie Grassmanna i wiązki tautologiczne]
Dla $k\ge1$, $N\ge0$ oznaczmy przez
$\mathrm{Gr}_k(\R^{k+N})$ przestrzeń $k$-wymiarowych
podprzestrzeni liniowych $\R^{k+N}$. Jej punktowi $V$
odpowiada włókno $V$ rzeczywistej wiązki tautologicznej
$\gamma_{k,N}$. Granica włączeń
$\mathrm{Gr}_k(\R^{k+N})\hookrightarrow
\mathrm{Gr}_k(\R^{k+N+1})$ daje model $BO(k)$
z uniwersalną wiązką $\gamma_k$. Zastąpienie
$\mathrm{Gr}_k$ przestrzenią zorientowanych
$k$-płaszczyzn daje $BSO(k)$ i $\gamma_k^{SO}$.
Przestrzenie Thoma tych wiązek oznaczamy
$MO(k):=\operatorname{Th}(\gamma_k)$ oraz
$MSO(k):=\operatorname{Th}(\gamma_k^{SO})$.
\end{definicja}

\begin{stwierdzenie}[Klasyfikacja wiązki normalnej]
\label{prop:normal-class-v9}
Jeśli $M^n\subset\R^{n+k}$ jest zamkniętą
podrozmaitością, to jej wiązka normalna
$\nu_M$ jest klasyfikowana przez odwzorowanie
$c_\nu:M\to BO(k)$; w szczególności
$\nu_M\cong c_\nu^*\gamma_k$.
Jeśli $M$ jest zorientowana, to $\nu_M$ jest
orientowalna i można wybrać podniesienie mapy
klasyfikującej do $BSO(k)$, zgodne z orientacją
$TM\oplus\nu_M\cong\underline\R^{n+k}$.
\end{stwierdzenie}
\begin{proof}
Wybierzmy metrykę euklidesową. Włókno normalne
$\nu_{M,x}$ jest $k$-wymiarową podprzestrzenią
$\R^{n+k}$, więc funkcja $x\mapsto\nu_{M,x}$
jest ciągłą (a nawet gładką) mapą do
$\mathrm{Gr}_k(\R^{n+k})$. Jej cofnięciem
$\gamma_{k,n}$ jest z definicji właśnie $\nu_M$.
Wybrana orientacja $M$ i standardowa orientacja
$\R^{n+k}$ wyznaczają orientację każdego włókna
normalnego przez zasadę sumy bezpośredniej.
Stąd otrzymujemy mapę do zorientowanej
przestrzeni Grassmanna.
\end{proof}

\begin{twierdzenie}[Pontriagin--Thom dla kobordyzmu]
\label{thm:pt-MO-v9}
Dla każdego $n\ge0$ istnieją naturalne izomorfizmy
\[
\Omega_n^O\cong\pi_n(MO):=
\varinjlim_{k\to\infty}\pi_{n+k}(MO(k)),
\qquad
\Omega_n^{SO}\cong\pi_n(MSO):=
\varinjlim_{k\to\infty}\pi_{n+k}(MSO(k)).
\]
Mapy przejścia po prawej pochodzą z odwzorowań
$\Sigma MO(k)\to MO(k+1)$ (odpowiednio
$\Sigma MSO(k)\to MSO(k+1)$), wyznaczonych
przez dodanie trywialnej prostej do wiązki.
Po prawej występują grupy homotopii \emph{widm},
a nie grupy homotopii pojedynczej przestrzeni $BO$
czy $BSO$.
\end{twierdzenie}
\begin{proof}
\emph{Od rozmaitości do mapy.}
Osadźmy zamkniętą gładką $n$-rozmaitość
$M$ w $\R^{n+k}$ dla dostatecznie dużego $k$
(twierdzenie Whitneya). Osadzenie ma otoczenie
tubularne oraz wiązkę normalną $\nu_M$.
Po uzwartowieniu jednopunktowym
$\R^{n+k}\cup\{\infty\}=S^{n+k}$ otrzymujemy mapę ściągającą
$S^{n+k}\to\operatorname{Th}(\nu_M)$.
Wybór mapy klasyfikującej
$c_\nu:M\to BO(k)$ daje odwzorowanie przestrzeni
Thoma $\operatorname{Th}(\nu_M)\to MO(k)$,
a złożenie definiuje klasę w $\pi_{n+k}(MO(k))$.
Dla rozmaitości zorientowanej wybieramy orientację normalną
wyznaczoną przez standardową orientację przestrzeni
$\R^{n+k}$ i orientację $M$; wówczas mapą docelową jest
$MSO(k)$. Nie chodzi tu o trywializację wiązki normalnej.
Zmiana osadzenia w dostatecznie dużym wymiarze nie
zmienia stabilnej klasy. Dokładniej, jeśli
$i_0,i_1:M\hookrightarrow\R^{n+k}$ są dwoma osadzeniami,
po przejściu do $\R^{2(n+k)+1}$ można zestawić
$i_0$ i $i_1$ w rozłącznych ortogonalnych składnikach
i połączyć przez rodzinę osadzeń
$x\mapsto((1-t)i_0(x),t i_1(x),t)$; dla $t\in(0,1)$
różnowartościowość i iniektywność różniczki wynikają
z odpowiednich własności $i_0$ i $i_1$.
Po zaokrągleniu końców rodzina wyznacza izotopię
w wymiarze stabilnym. Zachowanie orientacji
normalnej wynika z ciągłości orientacji w rodzinie.
Identyfikacja nadmiarowych kierunków normalnych
jest zgodna z mapami strukturalnymi widma.
Jeśli $M_0$ i $M_1$ są kobordantne przez $W$,
możemy po zwiększeniu wymiaru osadzić $W$ w
$\R^{n+k}\times[0,1]$, produktowo przy brzegu.
Ściągnięcie otoczenia tubularnego $W$ daje homotopię
między klasami z $M_0$ i $M_1$. Tym samym
otrzymujemy dobrze określony homomorfizm z
odpowiedniej grupy kobordyzmu.

\emph{Od mapy do rozmaitości.}
Klasa $\pi_{n+k}(MO(k))$ jest reprezentowana
przez mapę z \emph{kompaktowej} sfery. W modelu
CW przestrzeni $MO(k)$ można, po homotopii
komórkowej, sprowadzić tę mapę do przestrzeni
Thoma pewnej skończenie wymiarowej wiązki
$\gamma_{k,N}$ nad $\mathrm{Gr}_k(\R^{k+N})$.
Nie wolno traktować całej przestrzeni Thoma jako
rozmaitości gładkiej w punkcie bazowym: jest ona
ilorazem, a więc ten krok wymaga lokalizacji.
Wybieramy otwarte otoczenie przekroju zerowego
we wnętrzu $D(\gamma_{k,N})$; na jego przeciwobrazie
mapa do przestrzeni Thoma przyjmuje wartości
w gładkiej całkowitej przestrzeni wiązki.
Stosujemy względne twierdzenie o aproksymacji
i transwersalności na zwartym obszarze zawierającym
przeciwobraz przekroju zerowego, pozostawiając
odwzorowanie niezmienione w sąsiedztwie punktu bazowego.
Po tej operacji wybieramy reprezentanta poprzecznego
do przekroju zerowego. Przeciwobraz sekcji zerowej
jest zamkniętą gładką
$n$-podrozmaitością $M\subset S^{n+k}$, a jej
wiązka normalna jest izomorficzna z cofnięciem
$\gamma_{k,N}$. W wersji zorientowanej analogiczna
konstrukcja zachowuje orientację normalną.
Homotopię map można uczynić poprzeczną do sekcji
zerowej względnie do końców; jej przeciwobraz jest
kobordyzmem otrzymanych rozmaitości.

\emph{Sprawdzenie odwrotności i stabilizacji.}
Dla map ściągających przeciwobraz sekcji zerowej
to początkowa rozmaitość wraz z daną wiązką normalną.
Odwrotnie, w otoczeniu regularnego przeciwobrazu
przybliżona mapa do przestrzeni Thoma daje współrzędne
normalne; ściągnięcie otoczenia tubularnego odtwarza
jej klasę homotopii, jak w dowodzie
Twierdzenia~\ref{thm:pt-framed-v9}.
Dodanie trywialnego kierunku normalnego jest zgodne
z $\Sigma MO(k)\to MO(k+1)$. Przechodząc do granic
prostych, uzyskujemy obustronne izomorfizmy.
\end{proof}

\begin{uwaga}[Co zostało użyte jako narzędzie pomocnicze?]
Dowód zawiera konstrukcję obu odwzorowań i argumenty,
że są wzajemnie odwrotne. Wymaga jednak standardowych
wyników geometrii różniczkowej: osadzania Whitneya,
otoczeń tubularnych, aproksymacji gładkiej,
transwersalności względnej i stabilizacji osadzeń.
Ich pełnych dowodów nie powtarzamy; znajdują się
w osobnej części książki o rozmaitościach. Podobnie
techniczny argument faktoryzacji przez skończony
etap Grassmanna jest faktem o filtracji CW.
Nie jest to dowód obliczenia całych pierścieni
$\Omega_*^O$ i $\Omega_*^{SO}$.
\end{uwaga}


\section{Uzupełnienia audytu: zakres transwersalności i konstrukcje stabilne}
\begin{uwaga}[Wartość regularna poza obrazem]
Dla mapy $f:S^{n+k}\to S^k$ wartość regularna może nie należeć
do obrazu. Jeżeli $q\notin f(S^{n+k})$, przeciwobraz $f^{-1}(q)$
jest pustą podrozmaitością, a mapa $f$ jest homotopijna ze stałą,
ponieważ $S^k\setminus\{q\}$ jest ściągalna. Rozróżnienie
to jest niezbędne, aby konstrukcja odwrotna Pontriagina--Thoma
obejmowała również klasę zerową.
\end{uwaga}

\begin{lemat}[Podstawowy przypadek aproksymacji dla przestrzeni Thoma]
\label{lem:thom-transversal-v9-audyt}
Niech $\xi\to B$ będzie gładką wiązką rzędu $k$ nad zwartą
gładką rozmaitością $B$, a $h:S^N\to\operatorname{Th}(\xi)$
ciągłym odwzorowaniem punktowanym. Wtedy $h$ jest
homotopijne, względnie do małego otoczenia punktu bazowego,
z odwzorowaniem, którego przeciwobraz $B\subset\operatorname{Th}(\xi)$
przekroju zerowego jest gładką podrozmaitością
kowymiaru $k$ z normalną wiązką izomorficzną do
cofnięcia $\xi$.
\end{lemat}
\begin{proof}
Wybierzmy mniejsze otwarte otoczenie $V$ przekroju zerowego
w $D(\xi)$, którego domknięcie nie przecina $S(\xi)$.
Jego obraz w ilorazie jest gładką otwartą podrozmaitością
$\operatorname{Th}(\xi)$, ponieważ iloraz nie utożsamia
żadnych punktów $V$. Zbiór $h^{-1}(B)$ jest domknięty,
więc zwarty. Ma otoczenie $U\subset S^N$, którego
obraz mieści się w $V$. Twierdzenie o względnej gładkiej
aproksymacji pozwala wygładzić $h|_U$ w mniejszym
otoczeniu $h^{-1}(B)$, zaś twierdzenie o transwersalności
pozwala wybrać małe zaburzenie poprzeczne do przekroju
zerowego. Zaburzenie może mieć podporę zwartą w $U$:
mapa pozostaje zatem niezmieniona poza $U$, w tym
w otoczeniu punktu bazowego. Twierdzenie o przeciwobrazie
transwersalnym daje podrozmaitość kowymiaru $k$.
Różniczka poprzecznego odwzorowania identyfikuje jej
normalną wiązkę z cofnięciem normalnej przekroju
zerowego, a ta ostatnia jest kanonicznie izomorficzna
z $\xi$. Cały argument używa zwykłej transwersalności
\emph{tylko na gładkim obszarze przestrzeni Thoma},
nie na jej punkcie bazowym.
\end{proof}

\begin{uwaga}[Podstawa i wybór grupy strukturalnej]
Orientacja $M$ nie daje uramowania $\nu_M$;
daje jedynie redukcję grupy strukturalnej normalnej
wiązki z $O(k)$ do $SO(k)$ przy wyborze orientacji otoczenia.
Dlatego klasyfikacja przez $MSO$ nie jest szczególnym
przypadkiem klasyfikacji uramowanej bez zapomnienia
dodatkowych danych. Z kolei dla ogólnej przestrzeni
klasyfikującej $BO(k)$ wymaga się numerycznych wiązek;
na zwartych rozmaitościach użytych tutaj warunek ten
jest spełniony.
\end{uwaga}

\section{Po co nam widma Thoma?}
Zbierzmy konstrukcje w jednym diagramie pojęciowym:
\[
\begin{gathered}
\text{uramowanie normalne} \quad\longleftrightarrow\quad
\pi_*^{\mathrm s}\cong\Omega_*^{\mathrm{fr}},\\[3pt]
\text{wiązka normalna bez orientacji}\quad\longleftrightarrow\quad
\pi_*(MO)\cong\Omega_*^O,\\[3pt]
\text{zorientowana wiązka normalna}\quad\longleftrightarrow\quad
\pi_*(MSO)\cong\Omega_*^{SO}.
\end{gathered}
\]
Widzimy, że \emph{różne rodzaje struktur normalnych}
prowadzą do różnych teorii kobordyzmu i różnych widm.
Tę samą ideę można zastosować do struktur zespolonych;
powstaje widmo $MU$ i teoria kobordyzmu zespolonego
$\Omega_*^U\cong\pi_*(MU)$. W tym przypadku chodzi
o stabilną strukturę zespoloną na wiązce normalnej;
nie wymagamy, aby sama rozmaitość miała strukturę
rozmaitości zespolonej. Podobnie struktury spinowe
prowadzą do widma $MSpin$, ale wymagają wyjaśnienia
odpowiednich przeszkód kohomologicznych.

\begin{uwaga}[Późniejsze połączenie z geometrią algebraiczną]
Kobordyzm zespolony topologiczny $MU$ jest szczególnie
ważny dla topologii przestrzeni zespolonych. Nie należy
go jednak \emph{utożsamiać bez dodatkowej konstrukcji}
z algebraiczną teorią kobordyzmu na rozmaitościach
algebraicznych (np. teorią Levine'a--Morela).
Porównanie wymaga map realizacji i hipotez dotyczących
ciała bazowego, gładkości oraz użytej teorii orientowanej.
\end{uwaga}

\section{Podsumowanie i dalsza droga}
Otrzymaliśmy rzeczywiste przejście między językiem
homotopii i geometrią podrozmaitości: wartości regularne
tworzą podrozmaitości uramowane, a otoczenia tubularne
zamieniają je z powrotem w odwzorowania do sfer.
Przestrzenie Thoma pozwalają ten mechanizm uogólnić
na kobordyzmy bez ramowania. Wersje $MO$ i $MSO$
dają geometryczną interpretację grup homotopii widm,
a klasy charakterystyczne wykrywają przeszkody do
zerowania klas kobordyzmu.

W dalszych etapach naturalne są: szczegółowe obliczenie
pierścienia $\Omega_*^O$, teoria klas Stiefela--Whitneya
jako pełnych niezmienników kobordyzmu, orientowalność
zespolona widm $MU$ i współczynniki formalnych praw
grupowych. Kolejne dwa rozdziały poświęcamy homologii Morse'a oraz
twierdzeniu Lefschetza o punktach stałych.

% Literatura do rozdziału: R. Cohen, Bundles, Manifolds and Homotopy,
% rozdział o Pontriaginie--Thomie; R. E. Stong, Notes on Cobordism Theory.

% Iteracja X, rozdział XVIII. Wymaga preambuły wersji IX.
\chapter{Homologia Morse'a: od punktów krytycznych do homologii singularnej}
\label{ch:morse-v10}

\section{Dlaczego rachunek punktów krytycznych wykrywa topologię?}
W teorii singularnej grupy łańcuchów są zwykle nieskończenie generowane,
ponieważ generatory stanowią wszystkie sympleksy singularne. Dla zwartej
gładkiej rozmaitości można skonstruować znacznie mniejszy kompleks:
jego generatorami będą punkty krytyczne jednej funkcji, a różniczkę
wyznaczymy, licząc odpowiednio zorientowane trajektorie gradientowe.
Zaskakujący jest fakt, że jego homologia nie zależy od wybranej
funkcji ani metryki: odzyskujemy zwykłą homologię singularną.
Przedstawimy dwa obrazy tej samej konstrukcji: geometryczny (trajektorie)
i komórkowy (doczepianie uchwytów). Nie będziemy utożsamiać pełnego
twierdzenia o sklejaniu trajektorii z prostą konsekwencją twierdzenia
Sarda: wprowadza ono osobny, głęboki argument analityczny.

\section{Funkcja Morse'a, indeks i lokalny model}
\begin{definicja}[Punkt niezdegenerowany i indeks Morse'a]
Niech $M$ będzie gładką rozmaitością bez brzegu o wymiarze $n$,
a $f\in C^\infty(M,\R)$. Punkt $p\in M$ nazywamy \emph{krytycznym}, jeśli $\dd f_p=0$.
W dowolnym układzie współrzędnych wokół $p$ określmy w takim
punkcie formę $\operatorname{Hess}_p f$ przez macierz drugich
pochodnych funkcji zapisanej we współrzędnych.
Jest ona niezależna od mapy jako forma dwuliniowa na $T_pM$, gdyż
składniki zależne od pochodnych pierwszego rzędu znikają dla $\dd f_p=0$.
Punkt krytyczny jest \emph{niezdegenerowany}, gdy hesjan nie ma jądra.
\emph{Indeksem} $\lambda(p)$ nazywamy wymiar maksymalnej podprzestrzeni,
na której $\operatorname{Hess}_pf$ jest ujemnie określony.
Funkcja jest \emph{funkcją Morse'a}, jeżeli wszystkie jej punkty
krytyczne są niezdegenerowane.
Zbiór punktów krytycznych oznaczamy przez
$\operatorname{Crit}(f):=\{p\in M:\dd f_p=0\}$.
\end{definicja}

\begin{twierdzenie}[Lemat Morse'a — używane twierdzenie geometryczne]
\label{thm:morse-lemma-v10}
Jeśli $p$ jest niezdegenerowanym punktem krytycznym funkcji
$f\in C^\infty(M,\R)$ o indeksie $\lambda$, to istnieją gładkie
współrzędne $(x_1,\ldots,x_n)$, w których $p$ ma współrzędne $0$ oraz
\[
 f(x)=f(p)-\sum_{i=1}^{\lambda}x_i^2+
                \sum_{i=\lambda+1}^{n}x_i^2.
\]
W szczególności niezdegenerowane punkty krytyczne są izolowane.
\end{twierdzenie}
\begin{proof}
Przyjmijmy lokalne współrzędne z $p=0$. Ponieważ
$\mathrm df_0=0$, dwukrotne zastosowanie wzoru Newtona--Leibniza daje
\[
f(x)-f(0)=\int_{t\in[0,1]}(1-t)\,\D^2f(tx)[x,x] \,\dd t
=x^\top A(x)x,
\quad A(x)=\int_{t\in[0,1]}(1-t)\,\operatorname{Hess}f(tx)\,\dd t.
\]
Macierz $A(x)$ jest symetryczna i gładka oraz
$A(0)=\tfrac12\operatorname{Hess}_0f$ jest odwracalna.
Po jednej stałej liniowej zmianie współrzędnych
możemy założyć $A(0)=J=\operatorname{diag}(-I_\lambda,I_{n-\lambda})$.
Rozpiszmy $A(x)$ blokowo względem
$\R^\lambda\oplus\R^{n-\lambda}$; blok górny jest
ujemnie określony w małym otoczeniu zera.
Gładka eliminacja bloków mieszanych prowadzi do
formy blokowo diagonalnej, której drugi blok jest
dodatnio określonym dopełnieniem Schura.
Na stożkach macierzy dodatnio określonych operacja
pierwiastka kwadratowego jest gładka.
Normalizując w ten sposób oba bloki, znajdujemy
gładką rodzinę macierzy odwracalnych $L(x)$ z
$A(x)=L(x)^\top J L(x)$.
Wobec tego $y=L(x)x$ ma w zerze różniczkę $L(0)$,
a zatem jest lokalnym dyfeomorfizmem na mocy twierdzenia
o funkcji odwrotnej. W nowych współrzędnych
$f=f(0)-\sum_{i=1}^\lambda y_i^2+\sum_{i=\lambda+1}^n y_i^2$.
Ta postać pokazuje też, że $p$ jest izolowanym punktem krytycznym.
\end{proof}

\begin{wniosek}[Skończoność punktów krytycznych]
Jeżeli $M$ jest zwartą rozmaitością bez brzegu, a $f:M\to\R$ jest
funkcją Morse'a, to zbiór $\operatorname{Crit}(f)$ jest skończony.
\end{wniosek}
\begin{proof}
Zbiór krytyczny jest domknięty, gdyż jest miejscem zerowym ciągłej
sekcji $\dd f$ wiązki kostycznej, a zatem zwarty. Z poprzedniego
twierdzenia jest dyskretny. Zwarta przestrzeń dyskretna jest skończona.
\end{proof}

\section{Przepływ gradientowy, rozmaitości stabilne i transwersalność}
\begin{definicja}[Para Morse'a--Smale'a]
Niech $M$ będzie zwartą gładką rozmaitością bez brzegu,
$f:M\to\R$ funkcją Morse'a, a $g$ metryką riemannowską.
Pole $X=-\operatorname{grad}_g f$ ma zupełny przepływ
$\Phi:\R\times M\to M$. Dla $p\in\operatorname{Crit}(f)$ definiujemy
\begin{align*}
 W^u(p)&=\{x\in M:\lim_{t\to-\infty}\Phi_t(x)=p\},\\
 W^s(p)&=\{x\in M:\lim_{t\to+\infty}\Phi_t(x)=p\}.
\end{align*}
Para $(f,g)$ jest \emph{parą Morse'a--Smale'a}, jeżeli
$W^u(p)$ przecina $W^s(q)$ transwersalnie dla wszystkich
$p,q\in\operatorname{Crit}(f)$. Transwersalność rozumiemy w każdym
punkcie przecięcia, jako równość $T_xW^u(p)+T_xW^s(q)=T_xM$.
\end{definicja}

\begin{lemat}[Monotoniczność energii i wymiary]
\label{lem:morse-energy-v10}
Dla każdej trajektorii $t\mapsto\Phi_t(x)$ zachodzi
$\frac{\dd}{\dd t}f(\Phi_t(x))=-\|\operatorname{grad}_g f(\Phi_t(x))\|_g^2$.
Rozmaitości $W^u(p)$ i $W^s(p)$ są zanurzonymi rozmaitościami,
obie są dyfeomorficzne odpowiednio z $\R^{\lambda(p)}$ oraz
$\R^{n-\lambda(p)}$. Jeśli trajektoria nietrywialna łączy
$p$ z $q$, to $f(p)>f(q)$.
\end{lemat}
\begin{proof}
Pierwsza równość jest regułą łańcuchową i definicją gradientu.
Jeżeli całka z $\|\operatorname{grad}f\|^2$ na dowolnym odcinku
nietrywialnej trajektorii byłaby zerowa, gradient byłby na nim zerowy
i trajektoria byłaby stała, sprzeczność. Stąd ścisła nierówność.
W otoczeniu $p$ liniaryzacja pola ma postać
$-\operatorname{Hess}_pf$ po identyfikacji za pomocą $g$; ma
$\lambda(p)$ dodatnich i $n-\lambda(p)$ ujemnych wartości własnych.
Twierdzenie o rozmaitościach stabilnej i niestabilnej daje wskazane
wymiary i gładkie zanurzenia. Globalne rozmaitości otrzymuje się
przez obrazy lokalnych przy przepływie; w przypadku gradientowym
są dyfeomorficzne z przestrzeniami euklidesowymi odpowiednich wymiarów.
To ostatnie korzysta ze standardowego twierdzenia o rozmaitościach
stabilnych, które traktujemy jako narzędzie z geometrii dynamiki.
\end{proof}

Dla $p\ne q$ oznaczamy przez
$\widetilde{\mathcal M}(p,q)=W^u(p)\cap W^s(q)$
przestrzeń \emph{sparametryzowanych} trajektorii. Translacja czasu
$\tau\cdot x=\Phi_\tau(x)$ wyznacza wolne działanie $\R$.
Iloraz $\mathcal M(p,q)=\widetilde{\mathcal M}(p,q)/\R$
nazywamy przestrzenią \emph{trajektorii niesparametryzowanych}.

\begin{stwierdzenie}[Wymiar przestrzeni trajektorii]
\label{prop:morse-dimension-v10}
Dla pary Morse'a--Smale'a i $p\ne q$ przestrzeń
$\mathcal M(p,q)$ jest gładką rozmaitością wymiaru
$\lambda(p)-\lambda(q)-1$, o ile jest niepusta.
\end{stwierdzenie}
\begin{proof}
Transwersalność daje
$\dim(W^u(p)\cap W^s(q))=\lambda(p)+(n-\lambda(q))-n
=\lambda(p)-\lambda(q)$.
W każdym punkcie przecięcia pole $X$ jest niezerowe i styczne do
przecięcia, więc jego przepływ daje lokalny model
$(t,y)\mapsto\Phi_t(y)$ dla działania translacji.
Globalny przekrój tego działania otrzymujemy na przecięciu z
regularnym poziomem $f=c$ dla $f(q)<c<f(p)$:
na każdej trajektorii $f$ jest ściśle malejąca, więc przecina
poziom $c$ dokładnie raz, transwersalnie.
Ten przekrój jest dyfeomorficzny z ilorazem i ma wymiar
$\lambda(p)-\lambda(q)-1$.
\end{proof}

\section{Zliczanie trajektorii i kompleks Morse'a}
Wybierzmy orientację każdej niestabilnej rozmaitości $W^u(p)$.
Wzdłuż $W^s(q)$ wybór orientacji $W^u(q)$ indukuje orientację wiązki
normalnej do $W^s(q)$, zgodnie z rozkładem stabilnym i niestabilnym
przy $q$ i transportem przez przepływ. Gdy $\lambda(p)-\lambda(q)=1$,
w punkcie $x\in W^u(p)\cap W^s(q)$ transwersalność pozwala porównać
orientację $T_xW^u(p)$ z orientacjami normalną do $W^s(q)$
i dodatniego kierunku przepływu. Tak otrzymujemy znak
$\varepsilon(\gamma)\in\{-1,+1\}$ niesparametryzowanej
trajektorii $\gamma\in\mathcal M(p,q)$. To ujęcie nie wymaga
orientowalności całej rozmaitości $M$. Przy zmianie orientacji
$W^u(p)$ zmieniają się znaki trajektorii wychodzących z $p$;
przy zmianie orientacji $W^u(q)$ zmieniają się znaki trajektorii
dochodzących do $q$. Obie zmiany odpowiadają zmianie znaków
odpowiednich generatorów kompleksu Morse'a.

\begin{twierdzenie}[Kompaktyfikacja trajektorii — twierdzenie geometryczne]
\label{thm:morse-compact-v10}
Niech $M$ będzie zwartą rozmaitością bez brzegu,
a $(f,g)$ parą Morse'a--Smale'a. Wtedy:
\begin{enumerate}[label=(\roman*)]
\item jeśli $\lambda(p)-\lambda(q)=1$, zbiór $\mathcal M(p,q)$
      jest skończony;
\item jeśli $\lambda(p)-\lambda(r)=2$, przestrzeń
      $\mathcal M(p,r)$ dopuszcza zwartą kompaktyfikację
      $\overline{\mathcal M}(p,r)$ o strukturze zorientowanej
      jednowymiarowej rozmaitości z brzegiem, przy czym
      \[
       \partial\overline{\mathcal M}(p,r)=
       \coprod_{\lambda(q)=\lambda(p)-1}
       \mathcal M(p,q)\times\mathcal M(q,r),
      \]
      a znak punktu brzegowego odpowiada iloczynowi znaków
      dwóch sklejanych trajektorii, z jednym wspólnym znakiem
      wynikającym z przyjętej orientacji brzegu.
\end{enumerate}
\end{twierdzenie}
\begin{proof}
Wyjaśnijmy dokładnie zakres użytego wyniku. Z równania energii
\[
\int_{t\in\R}\|\operatorname{grad}f(\gamma(t))\|_g^2\,\dd t
=f(p)-f(q)
\]
wynika ograniczenie energii trajektorii.
Przy zwartości $M$ ciąg trajektorii można, po dobraniu translacji
czasu i wybraniu podciągów, badać przez zwartość rozwiązań ODE
na przedziałach ograniczonych. Utrata zwartości może polegać na tym,
że trajektoria spędza coraz dłuższy czas w otoczeniu pośredniego
punktu krytycznego $q$. Powstaje wówczas trajektoria
\emph{łamana} $p\to q\to r$. Ścisły spadek $f$
i transwersalność ograniczają liczbę łamań; gdy różnica
indeksów wynosi $1$, łamanie jest niemożliwe, a dla różnicy
$2$ pojawia się dokładnie jeden punkt pośredni.
Aby stwierdzić, że każdemu takiemu złamaniu odpowiada
\emph{jeden koniec gładkiej jednowymiarowej rozmaitości}
$\mathcal M(p,r)$, potrzebne jest twierdzenie o sklejaniu
trajektorii i zgodności orientacji. Jego dowód wykorzystuje
lokalną postać przepływu przy punktach hiperbolicznych
oraz oszacowania dla sklejonych rozwiązań i nie jest
konsekwencją samej zwartości ODE. Przyjmujemy ten
konkretny wynik pomocniczy; wszystkie dalsze użycia
odnoszą się dokładnie do punktów (i)--(ii), bez
niejawnego założenia silniejszej kompaktyfikacji.
W szczególności nie twierdzimy, że samo ograniczenie energii
dowodzi istnienia współrzędnych sklejania lub zgodności
orientacji brzegowych. Te części są treścią przyjętego
twierdzenia geometrycznego.
\end{proof}

\begin{definicja}[Kompleks Morse'a]
\label{def:morse-complex-v10}
Dla pary Morse'a--Smale'a $(f,g)$ na zwartej rozmaitości bez brzegu
niech $CM_k(f,g)$ będzie wolną grupą abelową z bazą
$\{p\in\operatorname{Crit}(f):\lambda(p)=k\}$.
Określamy homomorfizm $\partial_k:CM_k(f,g)\to CM_{k-1}(f,g)$,
$\partial_kp=\sum_{\lambda(q)=k-1}n(p,q)q$, gdzie
$n(p,q)=\sum_{\gamma\in\mathcal M(p,q)}\varepsilon(\gamma)\in\Z$.
Skończoność sum wynika z Twierdzenia~\ref{thm:morse-compact-v10}.
\end{definicja}

\begin{twierdzenie}[Różniczka Morse'a spełnia $\partial^2=0$]
\label{thm:morse-d2-v10}
Dla każdej pary Morse'a--Smale'a $(f,g)$ na zwartej gładkiej
rozmaitości bez brzegu zachodzi $\partial_{k-1}\partial_k=0$
dla wszystkich $k\in\Z$.
\end{twierdzenie}
\begin{proof}
Współczynnik przy krytycznym punkcie $r$ o indeksie $k-2$
w $\partial_{k-1}\partial_kp$ wynosi
$\sum_{\lambda(q)=k-1}n(p,q)n(q,r)$.
Każdy składnik tego wyrażenia odpowiada punktowi brzegowemu
$\overline{\mathcal M}(p,r)$ z odpowiednim znakiem, według
Twierdzenia~\ref{thm:morse-compact-v10}(ii). Zorientowana
zwarta jednowymiarowa rozmaitość jest skończoną sumą odcinków
i okręgów; końce każdego odcinka mają przeciwne znaki,
a okręgi nie mają brzegu. Suma znaków jej punktów brzegowych
jest więc zerowa. Zatem wskazany współczynnik znika dla
każdego $r$, co dowodzi $\partial_{k-1}\partial_k=0$.
\end{proof}

Definiujemy $HM_k(f,g)=\ker\partial_k/\operatorname{im}\partial_{k+1}$.
Istotą teorii jest wykazanie, że grupy te nie są nowymi
niezmiennikami zależnymi od pary $(f,g)$.

\section{Od uchwytów do homologii singularnej}
\begin{twierdzenie}[Zmiana podpoziomu w punkcie krytycznym]
\label{thm:morse-handle-v10}
Niech $M$ będzie zwartą gładką rozmaitością bez brzegu,
$f:M\to\R$ funkcją Morse'a, a $a<b$ liczbami regularnymi
takimi, że $f^{-1}([a,b])$ zawiera dokładnie jeden punkt
krytyczny $p$ o indeksie $k$.
Wtedy podpoziom $M^b=f^{-1}(( -\infty,b])$ powstaje,
z dokładnością do homotopii, z $M^a$ przez dołączenie
jednej $k$-komórki. Ponadto
$H_j(M^b,M^a;\Z)\cong\Z$ dla $j=k$ i jest zerem
dla $j\ne k$.
\end{twierdzenie}
\begin{proof}
Z lematu Morse'a wokół $p$ mamy $f-f(p)=-\|u\|^2+\|v\|^2$
dla $u\in\R^k$, $v\in\R^{n-k}$. Uchwyt
$D^k\times D^{n-k}$ doczepiamy wzdłuż
$S^{k-1}\times D^{n-k}$, ponieważ tylko w tych kierunkach
występuje spadek wartości $f$ przy przechodzeniu
przez poziom krytyczny. Na obszarze bez punktów krytycznych, dla dowolnie ustalonej
metryki $g$, pole $-\operatorname{grad}_g f/\|\operatorname{grad}_g f\|_g^2$
spełnia $\dd f(X)=-1$; jego przepływ identyfikuje regularne
podpoziomy po pominięciu małych otoczeń punktów krytycznych.
Czas przepływu jest skończony, a zwartość $f^{-1}([a,b])$
zapobiega ucieczce trajektorii przez nieskończoność.
Daje to standardowe twierdzenie o dołączaniu uchwytu
(przyjmujemy jego lokalno-globalną wersję z teorii Morse'a).
Para $(D^k\times D^{n-k},S^{k-1}\times D^{n-k})$
deformacyjnie zwija się do $(D^k,S^{k-1})$.
Względna homologia tej pary jest $\Z$ w stopniu $k$
i zerem w pozostałych stopniach. Ekscyzja przenosi
ten rachunek na $(M^b,M^a)$.
\end{proof}

\begin{lemat}[Kilka punktów krytycznych o tej samej wartości]
\label{lem:morse-same-level-audit}
Niech $M$ będzie zwartą gładką rozmaitością bez brzegu,
$f:M\to\R$ funkcją Morse'a, a $a<c<b$ liczbami takimi, że
$a,b$ są regularne i jedyną wartością krytyczną w $[a,b]$ jest $c$.
Niech $p_1,\dots,p_r$ będą wszystkimi punktami krytycznymi na
poziomie $c$ i niech $k_i=\lambda(p_i)$. Wtedy $M^b$ powstaje z $M^a$,
z dokładnością do homotopii par, przez jednoczesne dołączenie
$r$ komórek wymiarów $k_1,\dots,k_r$. W szczególności
\[
 H_j(M^b,M^a;\Z)\cong
 \bigoplus_{\{i:k_i=j\}}\Z.
\]
\end{lemat}
\begin{proof}
Ponieważ punkty $p_i$ są skończenie liczne i izolowane, istnieją
parami rozłączne małe otoczenia Morse'a $U_i$ zawarte w
$f^{-1}((a,b))$. W każdym $U_i$ lokalny rachunek z dowodu
Twierdzenia~\ref{thm:morse-handle-v10} dostarcza uchwytu
$D^{k_i}\times D^{n-k_i}$ z częścią przyczepiającą
$S^{k_i-1}\times D^{n-k_i}$. Poza sumą tych otoczeń
nie ma punktów krytycznych na $f^{-1}([a,b])$;
znormalizowany przepływ gradientowy pozwala przesunąć
poziomy regularne do okolic $U_i$. Dlatego uchwyty można
doczepić jednocześnie, a ich części przyczepiające są
parami rozłączne. Względna ekscyzja redukuje
$H_j(M^b,M^a;\Z)$ do sumy prostej homologii par
$(D^{k_i}\times D^{n-k_i},S^{k_i-1}\times D^{n-k_i})$.
Każda wnosi jedną kopię $\Z$ w stopniu $k_i$.
Argument wykorzystuje tę samą geometryczną wersję twierdzenia
o doczepianiu uchwytów co Twierdzenie~\ref{thm:morse-handle-v10}.
\end{proof}

\begin{twierdzenie}[Zgodność homologii Morse'a i singularnej]
\label{thm:morse-singular-v10}
Dla zwartej gładkiej rozmaitości $M$ bez brzegu i pary
Morse'a--Smale'a $(f,g)$ istnieją izomorfizmy
$HM_k(f,g)\cong H_k(M;\Z)$ dla wszystkich $k\ge0$.
Są naturalne względem map kontynuacji pomiędzy
parami Morse'a--Smale'a, po ustaleniu zwyczajowych
identyfikacji z homologią singularną.
\end{twierdzenie}
\begin{proof}
Dla funkcji Morse'a wybierzmy przedziały rozdzielające
\emph{różne wartości krytyczne}, a niekoniecznie poszczególne
punkty krytyczne: kilka punktów może mieć tę samą wartość
(np. dwa siodła funkcji $\cos\theta+\cos\varphi$ na torusie).
Lemat~\ref{lem:morse-same-level-audit} umożliwia doczepienie
wszystkich uchwytów danego poziomu jednocześnie. Otrzymujemy
filtrację przez podpoziomy i strukturę CW, w której na każdy
punkt krytyczny $p$ przypada komórka wymiaru $\lambda(p)$.
Zatem grupa $k$-łańcuchów komórkowych jest wolna na punktach
krytycznych o indeksie $k$.
W transwersalnym ustawieniu Morse'a--Smale'a liczba
incydencji komórki niestabilnej $p$ z komórką
niższego indeksu $q$ jest stopniem mapy przyłączającej,
a stopień ten można obliczyć przez zorientowane przecięcia
$W^u(p)$ z przekrojem poprzecznym do $W^s(q)$.
Każde takie przecięcie odpowiada jednej trajektorii
$p\to q$, a znak przecięcia jest dokładnie
$\varepsilon(\gamma)$. Otrzymujemy więc
$\partial^{\mathrm{cell}}p=\sum_qn(p,q)q=\partial p$.
W tym przejściu używa się \emph{twierdzenia o zgodności
liczb incydencji z liczeniem trajektorii}; wynika ono
z konstrukcji map przyłączających za pomocą przepływu
oraz transwersalnego przecięcia z regularnym poziomem.
Jest to standardowe twierdzenie geometryczne, którego
pełny dowód wymaga kontroli deformacji brzegów
uchwytów w pobliżu wielu punktów krytycznych; nie
zastępujemy go bezpodstawnym stwierdzeniem, że
same grupy łańcuchów są z definicji identyczne.
Homologia komórkowa jest izomorficzna z singularną,
co kończy argument po zastosowaniu przytoczonego
twierdzenia o liczbach incydencji. Naturalność
wynika z map kontynuacji, które opisujemy dalej.
\end{proof}

\begin{uwaga}[Kontynuacja, niezmienniczość i zakres dowodu]
Dla dwóch par Morse'a--Smale'a $(f_0,g_0)$ i $(f_1,g_1)$
wybiera się gładką rodzinę danych i liczy trajektorie równania
$\dot\gamma=-\operatorname{grad}_{g_s}f_s(\gamma)$,
przy parametrze $s\in\R$ stałym dla $s\ll0$ i $s\gg0$.
Dla regularnej rodziny parametrów daje to mapę łańcuchową
$CM_*(f_0,g_0)\to CM_*(f_1,g_1)$; kompaktyfikacja
jednowymiarowych przestrzeni trajektorii dowodzi zgodności
z różniczką. Porównując dwie homotopie rodzin,
otrzymuje się homotopię łańcuchową; odwrócenie
kontynuacji daje izomorfizm w homologii.
Pełne niezależne przeprowadzenie tej konstrukcji
wymaga analitycznego twierdzenia o sklejaniu
(zależnego od parametru). W naszym skrypcie
niezmienniczość wynika również z izomorfizmu
z homologią singularną w Twierdzeniu~\ref{thm:morse-singular-v10}.
\end{uwaga}

\section{Nierówności Morse'a i charakterystyka Eulera}
\begin{twierdzenie}[Silne nierówności Morse'a]
\label{thm:morse-inequalities-v10}
Niech $M$ będzie zwartą rozmaitością bez brzegu, $(f,g)$
parą Morse'a--Smale'a, $m_k$ liczbą punktów krytycznych
indeksu $k$ i $b_k=\dim_\Q H_k(M;\Q)$.
Istnieje wielomian $Q(t)\in\Z[t]$ o nieujemnych współczynnikach taki, że
\[
\sum_{k=0}^{n}m_kt^k-\sum_{k=0}^{n}b_kt^k=(1+t)Q(t).
\]
W szczególności $m_k\ge b_k$,
$\sum_{j=0}^k(-1)^{k-j}m_j\ge
\sum_{j=0}^k(-1)^{k-j}b_j$ i
$\sum_k(-1)^km_k=\chi(M)$.
\end{twierdzenie}
\begin{proof}
Tensorujmy kompleks Morse'a z $\Q$. Niech
$r_k=\operatorname{rank}_\Q\partial_k$; przyjmujemy $r_0=r_{n+1}=0$.
Z wymiarów $\ker\partial_k$ i ilorazu homologii wynika
$m_k=b_k+r_k+r_{k+1}$. Pomnóżmy przez $t^k$
i zsumujmy po $k$. Otrzymujemy wzór z
$Q(t)=\sum_{k=0}^{n-1}r_{k+1}t^k$, którego współczynniki
są nieujemnymi liczbami całkowitymi.
Porównanie współczynników daje słabe nierówności.
Naprzemienne sumowanie do $k$ daje różnicę $r_{k+1}\ge0$;
podstawienie $t=-1$ daje wzór na charakterystykę Eulera.
\end{proof}

\section{Przykłady obliczeń}
\begin{przyklad}[Sfera i funkcja wysokości]
Dla $S^n\subset\R^{n+1}$, $n\ge1$, funkcja
$f(x)=x_{n+1}$ ma dwa punkty krytyczne: minimum indeksu $0$
i maksimum indeksu $n$. Dlatego $CM_0\cong CM_n\cong\Z$
i pozostałe grupy łańcuchowe są zerowe.
Dla $n>1$ różniczka automatycznie znika z przyczyn
stopniowania. Dla $n=1$ dwie trajektorie od maksimum
do minimum mają przeciwne znaki i znoszą się.
W obu przypadkach odzyskujemy $H_0(S^n;\Z)\cong\Z$
i $H_n(S^n;\Z)\cong\Z$.
\end{przyklad}

\begin{przyklad}[Cztery punkty krytyczne torusa]
Na $T^2=S^1\times S^1$ rozważmy
$f(\theta,\varphi)=\cos\theta+\cos\varphi$
z płaską metryką produktową. Punkty $(\pi,\pi)$,
$(0,\pi)$, $(\pi,0)$, $(0,0)$ mają indeksy odpowiednio
$0,1,1,2$. Stąd $CM_0\cong CM_2\cong\Z$ i $CM_1\cong\Z^2$.
Ponieważ znana homologia torusa ma rangi $1,2,1$,
nierówności Morse'a są tu równościami; wymusza to
$\partial_1=\partial_2=0$ po tensorowaniu z $\Q$,
a ponieważ grupy łańcuchowe są wolne, różniczki są zerowe
również nad $\Z$. Nie zakładaliśmy więc bez uzasadnienia,
że symetria automatycznie znosi wszystkie znaki.
\end{przyklad}

\begin{przyklad}[Niedoskonała funkcja na okręgu]
Dla $f(\theta)=\cos(2\theta)$ na $S^1$ mamy dwa maksima
$p_1,p_2$ i dwa minima $q_1,q_2$. Po wybraniu orientacji
niestabilnych kierunków różniczka może mieć postać
$\partial p_1=q_1-q_2$ oraz $\partial p_2=q_2-q_1$.
Macierz $\partial_1$ ma rząd $1$ nad $\Z$ i obraz
$\Z(q_1-q_2)$. Zatem $HM_1\cong\Z$ oraz
$HM_0\cong\Z$, mimo że funkcja posiada więcej punktów
krytycznych niż funkcja wysokości na okręgu.
\end{przyklad}

\begin{center}
\begin{tikzpicture}[>=Stealth,scale=0.93,font=\small]
\draw[->] (-3.55,0)--(3.58,0) node[right] {$\theta$};
\draw[very thick,manifoldblue,domain=-3.14159:3.14159,samples=140]
 plot (\x,{0.85*cos(deg(2*\x))});
\draw[manifoldred,->,thick] (-0.40,0.59)--(-0.75,0.06);
\draw[manifoldred,->,thick] ( 0.40,0.59)--( 0.75,0.06);
\fill[manifoldblue] (0,0.85) circle (1.7pt);
\fill[manifoldblue] (-1.5708,-0.85) circle (1.7pt);
\fill[manifoldblue] ( 1.5708,-0.85) circle (1.7pt);
\draw[manifoldblue,fill=white] (-3.14159,.85) circle (1.7pt);
\draw[manifoldblue,fill=white] ( 3.14159,.85) circle (1.7pt);
\node[above] at (0,1.02) {maksimum};
\node[below] at (-1.5708,-1.02) {minimum};
\node[below] at ( 1.5708,-1.02) {minimum};
\node[below] at (-3.14159,-0.08) {$-\pi$};
\node[below] at ( 3.14159,-0.08) {$\pi$};
\node[anchor=west] at (-3.46,1.53) {$f(\theta)=\cos(2\theta)$};
\end{tikzpicture}
\end{center}
\begin{center}\small Wykres $f(\theta)=\cos(2\theta)$ dla
$\theta\in[-\pi,\pi]$; końce przedziału oznaczają ten sam
punkt okręgu. Pełne punkty oznaczają trzy widoczne ekstrema
we wnętrzu przedziału, a punkty puste -- jeden wspólny
punkt maksymalny po utożsamieniu końców. Strzałki wskazują
kierunek ujemnego przepływu gradientowego.\end{center}

\section{Związek z teorią uchwytów i z topologią algebraiczną}
Kompleks Morse'a pozwala obliczać niezmienniki topologiczne
przez skończony rachunek algebraiczny, a geometria przepływu
wyjaśnia, skąd biorą się współczynniki różniczki.
W rozdziale o Pontriaginie--Thomie punkty regularne
i podrozmaitości reprezentowały klasy stabilnej homotopii;
tutaj liczby trajektorii reprezentują mapy brzegowe w
kompleksie łańcuchowym. Przy późniejszym połączeniu z geometrią algebraiczną
pożyteczna będzie analogia z przestrzeniami modułów i
zorientowanymi liczbami przecięć, lecz do przeniesienia tej
konstrukcji potrzebne są osobne twierdzenia o zwartości,
transwersalności i orientacji przestrzeni modułów.


% Iteracja X, rozdział XIX. Wymaga preambuły wersji IX.
\chapter{Stopień, indeks punktu stałego i twierdzenie Lefschetza}
\label{ch:lefschetz-v10}

\section{Dlaczego ślad działania na homologii wykrywa punkty stałe?}
Niech $X$ będzie skończonym kompleksem CW, a $f:X\to X$
odwzorowaniem ciągłym. Zwykły stopień mierzy działanie
$f$ na najwyższej homologii sfery. Dla przestrzeni, które
mają niezerową homologię w wielu stopniach, właściwym
uogólnieniem jest naprzemienna suma śladów działań
$f_*$ na wszystkich grupach homologii. Wynik może wymusić
istnienie punktu stałego, chociaż nie daje jego współrzędnych.
Rozdział oddziela twierdzenie dla dowolnych skończonych
kompleksów CW od dokładniejszego wzoru indeksowego na
gładkich rozmaitościach.

\section{Liczba Lefschetza i formuła śladu na łańcuchach}
\begin{definicja}[Liczba Lefschetza]
Niech $X$ będzie skończonym kompleksem CW, a $f:X\to X$
odwzorowaniem ciągłym. Ponieważ $H_k(X;\Q)$ są
skończeniewymiarowe i znikają dla $k>\dim X$, definiujemy
\[
 L(f)=\sum_{k\ge0}(-1)^k\operatorname{tr}
 \bigl(f_*:H_k(X;\Q)\to H_k(X;\Q)\bigr)\in\Z.
\]
Całkowitość wynika z tego, że $f_*$ zachowuje obrazy
wolnych części grup $H_k(X;\Z)$; macierze całkowite
mają całkowite ślady. Liczba $L(f)$ nie zależy od
wyboru bazy homologii.
\end{definicja}

\begin{lemat}[Lemat śladu Hopfa]
\label{lem:hopf-trace-v10}
Niech $(C_*,\partial)$ będzie ograniczonym kompleksem
skończeniewymiarowych przestrzeni wektorowych nad ciałem
$K$, a $F_*:C_*\to C_*$ odwzorowaniem łańcuchowym.
Wówczas
\[
 \sum_k(-1)^k\operatorname{tr}(F_k\mid C_k)
 =\sum_k(-1)^k\operatorname{tr}(H_k(F)\mid H_k(C)).
\]
\end{lemat}
\begin{proof}
Niech $Z_k=\ker\partial_k$ oraz $B_k=\operatorname{im}\partial_{k+1}$.
Wszystkie podprzestrzenie są niezmiennicze względem $F$.
Dokładne ciągi
$0\to Z_k\to C_k\to B_{k-1}\to0$ oraz
$0\to B_k\to Z_k\to H_k(C)\to0$
rozszczepiają się jako ciągi przestrzeni wektorowych.
Dla endomorfizmu zachowującego podprzestrzeń ślad jest
sumą śladu ograniczenia i śladu odwzorowania indukowanego
na ilorazie. Otrzymujemy
$\operatorname{tr}(F_k\mid C_k)=
 \operatorname{tr}(H_k(F))+
 \operatorname{tr}(F_k\mid B_k)+
 \operatorname{tr}(F_{k-1}\mid B_{k-1})$.
W naprzemiennej sumie dwa ostatnie wyrazy znoszą się
po przesunięciu indeksu. To daje tezę.
\end{proof}

\begin{wniosek}[Homotopie, deformacje i charakterystyka Eulera]
Jeśli $f,g:X\to X$ są homotopijne, to $L(f)=L(g)$.
Ponadto $L(\id_X)=\chi(X)$, gdzie
$\chi(X)=\sum_k(-1)^k\dim_\Q H_k(X;\Q)$.
Jeżeli $X$ jest niepusty i kontraktylny, to $L(f)=1$
dla każdego ciągłego $f:X\to X$.
\end{wniosek}
\begin{proof}
Homotopijne odwzorowania indukują te same homomorfizmy
w homologii. Dla $\id_X$ wszystkie ślady są wymiarami.
Dla przestrzeni kontraktylnej jedyną niezerową grupą
homologii racjonalnej jest $H_0(X;\Q)\cong\Q$,
a ciągła mapa przestrzeni spójnej działa na niej identycznie.
\end{proof}

\section{Twierdzenie Lefschetza dla skończonych kompleksów}
\begin{twierdzenie}[Lefschetz — istnienie punktu stałego]
\label{thm:lefschetz-v10}
Jeśli $X$ jest skończonym kompleksem CW, $f:X\to X$
ciągłym odwzorowaniem i $L(f)\ne0$, to istnieje
$x\in X$ takie, że $f(x)=x$.
\end{twierdzenie}
\begin{proof}
Najpierw rozważmy skończony wielokompleks $X=|K|$.
Zapiszemy dowód kontrapozycji, a następnie wyprowadzimy
wersję dla dowolnego skończonego CW-kompleksu przez
osadzenie jako retraktu otoczenia w przestrzeni euklidesowej.

Załóżmy zatem dla sprzeczności, że $X=|K|$ i $f$
nie ma punktów stałych. Umieśćmy $|K|$ w przestrzeni
euklidesowej. Dzięki zwartości
$\delta=\min_{x\in X}\|f(x)-x\|>0$.
Przez dostatecznie drobny podział barycentryczny
kompleksu docelowego i wyjściowego możemy założyć,
że każdy sympleks $\sigma$ tego samego drobnego podziału
ma średnicę mniejszą niż $\delta/4$ (jednostajna
ciągłość $f$ nie jest tu potrzebna do oszacowania odległości
$f(\sigma)$ od $\sigma$). Rzeczywiście
$\operatorname{dist}(f(\sigma),\sigma)>3\delta/4$
(odległość dowolnego $f(x)$ od dowolnego $y\in\sigma$
wynosi co najmniej $\delta-\|x-y\|$).
Niech $K$ oznacza ustalony drobny podział obu kopii
wielościanu $X$. Dobierzmy następny podział $K'$ dziedziny
i aproksymację symplicjalną $g:K'\to K$ mapy $f$.
Lemat o aproksymacji symplicjalnej zapewnia homotopię
$|g|\simeq f$ oraz, przy wystarczająco drobnej triangulacji
\emph{celu}, oszacowanie
$\sup_{x\in X}\||g|(x)-f(x)\|<\delta/4$.
Istotne jest rozróżnienie dwóch podziałów: $K'$ służy
do przybliżania mapy, a drobny podział $K$ kontroluje
zarówno średnice sympleksów w bazie grup łańcuchów,
jak i błąd aproksymacji.

Niech $S:C_*(K;\Q)\to C_*(K';\Q)$ oznacza odwzorowanie
subdywizji. Złożenie $F=g_\#\circ S$ jest
odwzorowaniem łańcuchowym i indukuje w homologii
$f_*$. Każdy sympleks występujący w $S(\sigma)$
leży geometrycznie w $\sigma$. Jego obraz przez $g$
leży w odległości co najwyżej $\delta/4$ od
$f(\sigma)$, a zatem nie może być sympleksem $\sigma$,
ponieważ $f(\sigma)$ jest oddalone od $\sigma$
o ponad $3\delta/4$. W macierzy $F_k$ względem
bazy zorientowanych $k$-sympleksów wszystkie wyrazy
diagonalne są więc zerowe. Z Lematu~\ref{lem:hopf-trace-v10}
wynika $L(f)=\sum_k(-1)^k\operatorname{tr}F_k=0$,
wbrew założeniu. To dowodzi tezy dla wielościanów.
Pozostaje przypadek skończonego kompleksu CW $X$.
Korzystamy z klasycznego twierdzenia o otoczeniowych
retraktach. Przestrzeń metryzowalna $X$ jest \emph{absolutnym
retraktorem otoczeniowym} (ANR), gdy dla każdego
osadzenia $X$ jako domkniętego podzbioru przestrzeni
metryzowalnej $Y$ jest retraktem pewnego otwartego
otoczenia w $Y$. Skończony CW-kompleks jest zwartą
metryzowalną przestrzenią typu ANR; można go więc osadzić
jako domknięty podzbiór $i:X\hookrightarrow\R^N$
tak, że istnieje otwarte otoczenie $U$ i retrakcja
$r:U\to X$, $r\circ i=\id_X$. To jest osobne twierdzenie
o ANR, niewynikające z samego twierdzenia o aproksymacji
symplicjalnej.
Zwartość pozwala znaleźć skończony wielościan
$P\subset U$ zawierający $i(X)$: bierzemy skończenie
wiele dostatecznie małych sympleksów triangulacji
$\R^N$, których suma zawiera $i(X)$ i leży w $U$.
Niech $j:X\to P$ będzie tym osadzeniem, a
$R=r|_P:P\to X$ retrakcją. Definiujemy
$F=j\circ f\circ R:P\to P$.
Gdyby $F(y)=y$, to $y\in j(X)$ i z $R\circ j=\id_X$
wynikałby punkt stały $f$, sprzeczność.
Tak więc $F$ nie ma punktów stałych i z przypadku
wielościanowego $L(F)=0$.
Dla każdego $k$ zachodzi $F_*=j_* f_*R_*$,
a z cykliczności śladu także dla map między różnymi
skończeniewymiarowymi przestrzeniami
$\operatorname{tr}(j_*f_*R_*)=
\operatorname{tr}(f_*R_*j_*)=\operatorname{tr}(f_*)$.
Zatem $L(f)=L(F)=0$, co kończy dowód.
\end{proof}

\begin{wniosek}[Twierdzenie Brouwera]
Każde ciągłe odwzorowanie $f:D^n\to D^n$, $n\ge1$,
ma punkt stały.
\end{wniosek}
\begin{proof}
Dysk $D^n$ jest skończonym wielościanem z dokładnością
do homeomorfizmu i jest kontraktylny, więc
$L(f)=1$. Stosujemy Twierdzenie~\ref{thm:lefschetz-v10}.
\end{proof}

\begin{uwaga}[Właściwy zakres wniosku]
Odwrotność twierdzenia Lefschetza jest fałszywa:
$L(f)=0$ nie oznacza istnienia odwzorowania bez punktów
stałych w danej klasie homotopii ani nie oznacza, że
$f$ nie ma punktów stałych. Przykładowo identyczność
na okręgu ma $L(\id_{S^1})=\chi(S^1)=0$, ale każdy
punkt jest stały. Obrót okręgu przez kąt niezerowy
ma także $L=0$ i nie ma punktów stałych.
\end{uwaga}

\section{Indeks izolowanego punktu stałego}
\begin{definicja}[Indeks lokalny]
\label{def:fixed-index-v10}
Niech $M^n$ będzie gładką rozmaitością bez brzegu,
$n\ge1$, $f:M\to M$ odwzorowaniem ciągłym oraz
$p\in M$ izolowanym punktem stałym. Wybierzmy mapę
$\varphi:U\to\R^n$ wokół $p$ i małą kulę $B$
o domknięciu zawartym w $\varphi(U)$ taką, że
$\varphi\circ f\circ\varphi^{-1}$ jest określona
na $\overline B$ i $p$ jest jedynym punktem stałym
w $\varphi^{-1}(\overline B)$. Dla $F=\varphi f\varphi^{-1}$
określamy
$\operatorname{ind}_p(f)=\deg(\id-F,B,0)$,
tj. stopień Brouwera odwzorowania
$x\mapsto x-F(x)$ względem zera; równoważnie jest to
stopień mapy
$\partial B\to S^{n-1}$, $x\mapsto(x-F(x))/\|x-F(x)\|$.
\end{definicja}

\begin{stwierdzenie}[Właściwości indeksu]
Indeks jest niezależny od wystarczająco małej mapy i kuli,
addytywny przy rozdzieleniu izolowanych punktów stałych
oraz niezmienniczy przy homotopii, podczas której
na brzegu rozważanego otoczenia nie pojawiają się punkty
stałe. Jeśli $f$ jest różniczkowalne w $p$
i $\det(I-\D f_p)\ne0$, to
$\operatorname{ind}_p(f)=\operatorname{sgn}\det(I-\D f_p)$.
\end{stwierdzenie}
\begin{proof}
Niezależność od otoczenia, addytywność i niezmienniczość
są własnościami stopnia Brouwera, wynikającymi
z ekscyzji, względnej homologii par $(B,\partial B)$
i homotopijnej niezmienniczości stopnia.
Dla ostatniej równości w lokalnych współrzędnych
$F(x)=F(0)+\D F_0x+o(\|x\|)$ oraz $F(0)=0$.
Jeśli $A=I-\D F_0$ jest odwracalna, wybieramy
$\varepsilon>0$ tak, by na sferze $\|x\|=\varepsilon$
reszta $o(\|x\|)$ miała normę mniejszą niż
$\tfrac12\min_{\|v\|=1}\|Av\|\,\varepsilon$.
Homotopia $Ax-t(F(x)-\D F_0x)$, $0\le t\le1$,
nie ma zer na brzegu kuli. Stopień $A$ wynosi
$\operatorname{sgn}\det A$, co daje wzór.
Niezależność od współrzędnych dla dowolnej mapy ciągłej
wymaga czegoś więcej niż sprzężenia różniczek $f$.
Jeśli $h$ jest przejściem między dwoma układami współrzędnych
z $h(0)=0$, to w małym otoczeniu $0$ zachodzi
$h(x)-h(F(x))=(\D h_0+E(x))(x-F(x))$,
gdzie $\|E(x)\|\to0$ dla $x\to0$. Wynika to ze wzoru
na przyrost $h$ wzdłuż odcinka łączącego $x$ z $F(x)$;
w dostatecznie małej kuli odcinek leży w dziedzinie $h$.
Na jej brzegu macierze $\D h_0+tE(x)$, $0\le t\le1$,
są odwracalne. Stopień pola różnicowego zmienia się przy
przejściu do nowych współrzędnych o dwa czynniki
$\operatorname{sgn}\det\D h_0$: od zmiany współrzędnych
dziedziny oraz od mnożenia różnicy przez $\D h_0$.
Ich iloczyn wynosi $1$, stąd indeks nie zależy od mapy.
Także niezależność od małej kuli wynika z ekscyzji,
ponieważ jedynym zerem $\id-F$ pozostaje punkt $0$.
\end{proof}

\section{Formuła Lefschetza--Hopfa przez przecięcie grafu z diagonalną}
\begin{twierdzenie}[Formuła indeksowa]
\label{thm:lefschetz-hopf-v10}
Niech $M^n$ będzie zamkniętą gładką rozmaitością,
$n\ge1$, a $f:M\to M$ gładkim odwzorowaniem,
którego wszystkie punkty stałe są nieosobliwe,
tzn. $\det(I-\D f_p)\ne0$ dla każdego $f(p)=p$.
Wówczas punktów stałych jest skończenie wiele i
\[
 L(f)=\sum_{p:f(p)=p}\operatorname{sgn}\det(I-\D f_p).
\]
\end{twierdzenie}
\begin{proof}
Punkt stały $p$ odpowiada przecięciu grafu
$\Gamma_f=\{(x,f(x)):x\in M\}$ z diagonalną
$\Delta=\{(x,x):x\in M\}$ w $M\times M$.
Warunek $\det(I-\D f_p)\ne0$ jest dokładnie warunkiem
transwersalności obu $n$-wymiarowych podrozmaitości
w rozmaitości $2n$-wymiarowej. Przecięcia są więc
izolowane, a ponieważ $M$ jest zwarta, jest ich
skończenie wiele. Lokalny znak przecięcia jest
$\operatorname{sgn}\det(I-\D f_p)$: normalny kierunek
do diagonali identyfikujemy z różnicą pierwszego
i drugiego wektora stycznego.

Wyjaśnijmy równość globalnej sumy przecięć ze śladem
bez przemilczenia znaków. Najpierw rozpatrzmy spójną zorientowaną rozmaitość $M$.
Jeśli $M$ jest niespójna, to $f$ może przenosić jej spójne
składowe w inne składowe. W rozkładzie
$H^*(M;\Q)=\bigoplus_{\alpha}H^*(M_\alpha;\Q)$
śladowi $f^*$ odpowiadają tylko bloki diagonalne, a więc
te składowe $M_\alpha$, dla których $f(M_\alpha)\subseteq M_\alpha$.
Pozostałe składowe nie zawierają punktów stałych i nie
wnoszą wkładu do śladu. Dlatego argument dla przypadku
spójnego wystarczy na składowych zachowywanych przez $f$.
Dla rozmaitości nieorientowalnej potrzebna będzie odrębna
wersja dualności z lokalnym systemem orientacji;
nie wynika ona z rachunku dla zwykłej klasy fundamentalnej.
Niech $a_{p,j}$ będzie bazą $H^p(M;\Q)$ i $b_{p,j}$
bazą $H^{n-p}(M;\Q)$ dualną w sensie
$\langle a_{p,i}\smile b_{p,j},[M]\rangle=\delta_{ij}$.
Przez $\pi_1,\pi_2:M\times M\to M$ oznaczmy rzutowania.
Klasa diagonalna ma postać
\[
 \operatorname{PD}_{M\times M}[\Delta]
 =\sum_{p=0}^{n}\sum_j(-1)^{n-p}
 \pi_1^*a_{p,j}\smile\pi_2^*b_{p,j}.
\]
Wzór wynika z własności
$\langle\operatorname{PD}[\Delta]\smile
 \pi_1^*\alpha\smile\pi_2^*\beta,[M\times M]\rangle
 =\langle\alpha\smile\beta,[M]\rangle$
i reguły znaków Künnetha.
Dla mapy grafu $j_f:M\to M\times M$,
$j_f(x)=(x,f(x))$, otrzymujemy
\begin{align*}
 \langle j_f^*\operatorname{PD}[\Delta],[M]\rangle
 &=\sum_{p,j}(-1)^{n-p}
 \langle a_{p,j}\smile f^*b_{p,j},[M]\rangle\\
 &=\sum_{q=0}^n(-1)^q
 \operatorname{tr}(f^*\mid H^q(M;\Q))=L(f).
\end{align*}
W drugim przejściu podstawiliśmy $q=n-p$;
macierze $f^*$ i $f_*$ mają te same ślady nad $\Q$.
Lewa strona jest globalną algebraiczną liczbą
przecięć grafu z diagonalną, a transwersalność
pozwala policzyć ją jako sumę znaków lokalnych.
Dla nieorientowalnego $M$ końcowa równość również jest
prawdziwa. Wymaga jednak odpowiednika klasy diagonalnej
z lokalnymi współczynnikami i sparowania skręconej
klasy Thoma normalnej do $\Delta$ z klasą fundamentalną
współczynnikową. Tę wersję twierdzenia indeksowego
przyjmujemy tutaj jako odrębny standardowy wynik;
powyższy rachunek bazami i wzorem Künnetha dowodzi
formuły bezpośrednio w przypadku orientowalnym.
Lokalna definicja indeksu pozostaje poprawna również
bez globalnego wyboru orientacji $M$.
\end{proof}

\begin{center}
\begin{tikzpicture}[>=Stealth,font=\small,scale=.95]
\draw[->] (-0.2,0)--(5.6,0) node[below] {$x$};
\draw[->] (0,-0.2)--(0,4.5) node[left] {$y$};
\draw[very thick,manifoldblue] (0,0)--(4.1,4.1)
 node[above right] {$\Delta=\{y=x\}$};
\draw[very thick,manifoldred,domain=0:4.2,samples=80]
 plot (\x,{2.6-0.5*\x+0.1*\x*\x});
\node[manifoldred,anchor=west] at (4.4,2.6) {$\Gamma_f$};
\fill[manifoldgreen] (2.0,2.0) circle (1.8pt);
\draw[dashed,manifoldgreen] (2.0,0)--(2.0,2.0);
\node[anchor=north] at (2.0,-.12) {punkt stały};
\end{tikzpicture}
\end{center}
\begin{center}\small Schemat lokalny: punkty stałe są przecięciami
wykresu $y=f(x)$ z diagonalną. Rysunek ilustruje mechanizm,
nie jest wykresem dowolnej mapy rozmaitości.\end{center}

\section{Obliczenia: sfery, torusy i kontrprzykłady}
\begin{przyklad}[Sfera i stopień]
Jeśli $f:S^n\to S^n$ jest ciągłe, $n\ge1$,
to $H_0(S^n;\Q)\cong\Q$ i $H_n(S^n;\Q)\cong\Q$,
a na drugiej grupie $f_*$ jest mnożeniem przez $\deg f$.
Wobec tego $L(f)=1+(-1)^n\deg f$.
Dla odwzorowania antypodalnego $a(x)=-x$ mamy
$\deg a=(-1)^{n+1}$, więc $L(a)=0$, zgodnie z tym,
że $a$ nie ma punktów stałych.
Z drugiej strony każda mapa $S^{2m}\to S^{2m}$
stopnia różnego od $-1$ musi mieć punkt stały.
\end{przyklad}

\begin{przyklad}[Torus i całkowita macierz]
Niech $A\in M_d(\Z)$ i $f_A:T^d=\R^d/\Z^d\to T^d$
będzie mapą $[x]\mapsto[Ax]$. Ponieważ
$H_k(T^d;\Q)\cong\bigwedge^k\Q^d$,
mapa $f_{A*}$ działa jako $\bigwedge^k A$.
Z tożsamości wyznacznikowej
$\det(I-tA)=\sum_{k=0}^d(-t)^k
 \operatorname{tr}(\bigwedge^kA)$ wynika
$L(f_A)=\det(I-A)$.
Jeśli $\det(I-A)\ne0$, liczba rozwiązań $[Ax]=[x]$
wynosi $|\det(I-A)|$; wszystkie są nieosobliwe,
a ich indeksy mają znak $\operatorname{sgn}\det(I-A)$.
Formuła Lefschetza--Hopfa daje więc tożsamość
$L(f_A)=\det(I-A)$ również jako rachunek lokalny.
\end{przyklad}

\begin{przyklad}[Rotacja okręgu]
Dla rotacji $R_\alpha:S^1\to S^1$ z $\alpha\notin2\pi\Z$
mapa jest homotopijna do identyczności i ma
$L(R_\alpha)=1-1=0$; nie ma też punktu stałego.
Dla $\alpha=0$ ta sama liczba Lefschetza wynosi zero,
choć wszystkie punkty są stałe. Dowodzi to, że
warunek $L(f)\ne0$ jest wystarczający, lecz niekonieczny,
dla istnienia punktu stałego.
\end{przyklad}

\begin{uwaga}[Zakres dowodu formuły indeksowej]
Dowód formuły Lefschetza--Hopfa w przypadku zorientowanym
jest redukcją do dualności Poincarégo i klasy diagonalnej,
nie nowym dowodem dualności. W przypadku nieorientowalnym
korzystamy z wersji indeksowej uwzględniającej lokalne
współczynniki. Dla skończonych CW-kompleksów samą
implikację $L(f)\ne0\Rightarrow\operatorname{Fix}(f)\ne\varnothing$
udowodniliśmy niezależnie, stosując aproksymację
symplicjalną i własność ANR.
\end{uwaga}

\section{Relacje z teorią Morse'a i dalszymi częściami książki}
Dla funkcji Morse'a liczba $\chi(M)=\sum_p(-1)^{\lambda(p)}$
jest naprzemienną sumą generatorów kompleksu Morse'a.
Z drugiej strony $L(\id_M)=\chi(M)$ jest śladem działania
identyczności na homologii. To dwa geometrycznie odmienne
wyprowadzenia tego samego niezmiennika. Połączeniem obu
jest twierdzenie Poincarégo--Hopfa: suma lokalnych indeksów
zer pola wektorowego jest równa $\chi(M)$.
Nie utożsamiamy tego twierdzenia z przypadkiem mapy
$\id_M$ w formule indeksowej, gdyż identyczność nie
ma izolowanych punktów stałych. Należałoby najpierw
wybrać odpowiednie zaburzenie identyczności i porównać
lokalne stopnie; taki argument należy do teorii indeksu
pól wektorowych.

% Źródła orientacyjne: A. Hatcher, Algebraic Topology;
% J. Milnor, Morse Theory; M. Schwarz, Morse Homology;
% R. Brown, Lefschetz Fixed Point Theorem; D. Kim, The Lefschetz fixed-point formula.


% Iteracja XI: szczegółowy audyt części analitycznej homologii Morse'a.
\chapter{Zwartość i sklejanie trajektorii Morse'a: podstawy analityczne}
\label{ch:morse-analytic-v11}

\section{Cel i hipotezy twierdzeń}
W rozdziale~\ref{ch:morse-v10} skonstruowaliśmy różniczkę
Morse'a i użyliśmy kompaktyfikacji przestrzeni trajektorii.
Nie należy jednak utożsamiać trzech odmiennych faktów:
(i) istnienia granic ciągów trajektorii, być może łamanych,
(ii) faktu, że każda trajektoria łamana jest granicą trajektorii
niełamanych, oraz (iii) zgodności znaków orientacji na
brzegu kompaktyfikacji. Pierwszy fakt wynika w dużej mierze
z ODE, zwartości i hipotezy Morse'a--Smale'a; drugi wymaga
argumentu sklejania, a trzeci reguły dla orientacji
przestrzeni przecięć. Rozdzielimy je i podamy dokładnie
hipotezy każdego kroku.

W całym rozdziale $M$ oznacza \emph{zamkniętą} (zwartą, bez brzegu)
gładką rozmaitość, $f:M\to\R$ funkcję Morse'a, $g$ metrykę
riemannowską, a $X=-\operatorname{grad}_g f$ jej ujemne pole
gradientowe. Przez $\Phi_t$ oznaczamy zupełny przepływ $X$.
Gdy potrzebujemy wymiarów przestrzeni moduli, dodatkowo zakładamy
warunek Morse'a--Smale'a. Nie wymagamy orientowalności $M$;
orientacje wybieramy oddzielnie na rozmaitościach niestabilnych.

\section{Równanie energii i granice trajektorii}
\begin{lemat}[Równanie energii i punkty graniczne]
\label{lem:morse-energy-analytic-v11}
Niech $\gamma:\R\to M$ będzie całkową krzywą pola $X$.
Wówczas dla $a<b$ zachodzi
\[
 f(\gamma(a))-f(\gamma(b))=
 \int_{t\in[a,b]}\|\operatorname{grad}_g f(\gamma(t))\|_g^2\,\dd t.
\]
Ponadto zbiory graniczne $\omega(\gamma)$ i $\alpha(\gamma)$ są
niepustymi, spójnymi podzbiorami $\operatorname{Crit}(f)$;
w konsekwencji są pojedynczymi punktami krytycznymi.
\end{lemat}
\begin{proof}
Pierwsza tożsamość wynika z
$\tfrac{\dd}{\dd t}f(\gamma(t))=-\|\nabla f(\gamma(t))\|^2$.
Monotoniczność i zwartość $M$ dają skończone granice
$\ell_+=\lim_{t\to+\infty}f(\gamma(t))$ oraz
$\ell_-=\lim_{t\to-\infty}f(\gamma(t))$.
Zbiór $\omega(\gamma)$ jest przecięciem malejącej rodziny
niepustych, zwartych i spójnych zbiorów
$\overline{\gamma([T,\infty))}$, więc jest niepusty i spójny.
Jest niezmienniczy względem przepływu. Jeżeli punkt $y$ tego
zbioru byłby niekrytyczny, to
$\tfrac{\dd}{\dd t}f(\Phi_t(y))|_{t=0}<0$.
Tymczasem na całym zbiorze $\omega(\gamma)$ funkcja $f$
ma stałą wartość $\ell_+$, a niezmienniczość daje sprzeczność.
Zatem $\omega(\gamma)\subseteq\operatorname{Crit}(f)$.
Zbiór krytyczny jest skończony, a spójny podzbiór zbioru
skończonego jest jednopunktowy. Rozumowanie dla $\alpha(\gamma)$
jest analogiczne, po zmianie kierunku czasu. Ponieważ zbiory
graniczne są jednopunktowe, mamy faktyczną zbieżność krzywej
przy $t\to\pm\infty$.
\end{proof}

\begin{lemat}[Jednostajne ograniczenie czasu z dala od zer]
\label{lem:morse-transit-v11}
Niech $U$ będzie otwartym otoczeniem wszystkich punktów
krytycznych $f$ oraz $K=M\setminus U$.
Jeżeli $K\ne\varnothing$, to istnieje $m>0$ takie, że
$\|\nabla f(x)\|_g\ge m$ dla $x\in K$.
Łączny czas, jaki dowolna niestała trajektoria gradientowa
spędza w $K$, nie przekracza
$\bigl(\max_M f-\min_M f\bigr)/m^2$.
\end{lemat}
\begin{proof}
Zwartość $K$ i brak punktów krytycznych dają
$m=\min_K\|\nabla f\|_g>0$.
Dla dowolnego mierzalnego podzbioru $A\subset\R$, dla którego
$\gamma(A)\subseteq K$, równanie energii daje
$m^2\lambda_1(A)\le\int_{t\in A}\|\nabla f(\gamma(t))\|_g^2\,\dd t
\le\max_Mf-\min_Mf$.
Tutaj $\lambda_1$ oznacza jednowymiarową miarę Lebesgue'a.
W szczególności nierówność zachodzi dla całego zbioru czasów,
w których krzywa przebywa w $K$.
\end{proof}

\section{Granice ciągów rozwiązań i trajektorie łamane}
\begin{definicja}[Trajektoria łamana]
\label{def:broken-morse-v11}
Dla punktów krytycznych $p,q$ \emph{trajektorią łamaną}
od $p$ do $q$ nazywamy skończony ciąg
$([\gamma_1],\ldots,[\gamma_\ell])$,
gdzie $\gamma_j$ jest niestałą krzywą całkową $X$,
$\gamma_1(-\infty)=p$, $\gamma_\ell(+\infty)=q$ oraz
$\gamma_j(+\infty)=\gamma_{j+1}(-\infty)$.
Nawias oznacza klasę z dokładnością do translacji parametru czasu.
\end{definicja}

\begin{twierdzenie}[Zwartość z dokładnością do łamania]
\label{thm:morse-broken-compact-v11}
Niech $p\ne q$ będą punktami krytycznymi $f$ i niech
$([\gamma_j])_{j\in\N}$ będzie ciągiem elementów
$\mathcal M(p,q)$.
Istnieje podciąg oraz trajektoria łamana
$([\eta_1],\ldots,[\eta_\ell])$ od $p$ do $q$ takie, że
każdą $\eta_i$ można otrzymać jako granicę $C^\infty$
na zwartych przedziałach odpowiednio \emph{przesuniętych w czasie}
krzywych $\gamma_j$. Ponadto wartości $f$ na kolejnych
punktach krytycznych są ściśle malejące, więc
$\ell\le\#\operatorname{Crit}(f)-1$.
Jeśli $(f,g)$ jest parą Morse'a--Smale'a, to
$\ell\le\lambda(p)-\lambda(q)$.
\end{twierdzenie}
\begin{proof}
Zbiór punktów krytycznych jest skończony. Wybierzmy
skończony zbiór regularnych wartości $c_1>\cdots>c_N$
tak, aby między każdymi dwiema kolejnymi różnymi wartościami
krytycznymi leżało jedno $c_i$ (niektóre $c_i$ mogą leżeć
poza przedziałem $(f(q),f(p))$ i wtedy je pomijamy).
Aby kontrolować \emph{także końce} trajektorii, wybieramy
odpowiednio blisko wartości $f(p)$ dodatkowy regularny
poziom $c_+$, położony poniżej $f(p)$, a blisko $f(q)$
poziom $c_-$, położony powyżej $f(q)$.
Można je wybrać tak, by każdy punkt przecięcia trajektorii
wychodzącej z $p$ z $f^{-1}(c_+)$ należał do zwartej
sfery przekrojowej lokalnej rozmaitości niestabilnej
$W^u_{\mathrm{loc}}(p)$, a analogicznie dla $c_-$
i $W^s_{\mathrm{loc}}(q)$. Rzeczywiście w lokalnych
współrzędnych stabilnych i niestabilnych funkcja $f$
ma na tych rozmaitościach odpowiednio niezdegenerowane
maksimum i minimum; poziomy blisko wartości krytycznej
są zwartymi sferami, a na brzegu wybranego małego
otoczenia spadek (względnie wzrost w czasie wstecznym)
ma dodatnie ograniczenie od dołu. Dlatego trajektoria
musi przeciąć taki poziom przed opuszczeniem otoczenia.
Ten wybór wyklucza utratę końcowego punktu $p$ lub $q$
na rzecz innego punktu krytycznego o tej samej wartości.

Na każdej niestałej trajektorii od $p$ do $q$ funkcja
$f\circ\gamma_j$ jest ściśle malejąca i przyjmuje każdą
wartość $c_i\in(f(q),f(p))$ dokładnie raz; do rodziny
przekrojów włączamy także $c_+$ i $c_-$, jeżeli nie były
już wybrane. Niech
$x_{j,i}$ będzie odpowiednim punktem przecięcia krzywej
z zwartą regularną hiperpłaszczyzną $f^{-1}(c_i)$.
Po wybraniu podciągu wszystkie $x_{j,i}$ zbiegają do pewnych
$x_i\in f^{-1}(c_i)$.

Każda krzywa $t\mapsto\Phi_t(x_{j,i})$ zbiega na zwartych
przedziałach do $\eta_i(t):=\Phi_t(x_i)$, nawet w topologii
$C^\infty$, ponieważ przepływ jest gładki, a $M$ zwarta.
Granice w obu kierunkach czasu są punktami krytycznymi
na mocy Lematu~\ref{lem:morse-energy-analytic-v11}.
Zachodzą dwa przypadki dla sąsiednich przekrojów $c_i,c_{i+1}$:
czas przejścia krzywej $\gamma_j$ między nimi ma podciąg
ograniczony albo dąży do $+\infty$.
W pierwszym przypadku $x_i$ oraz $x_{i+1}$ należą do tej
samej trajektorii granicznej, gdyż przejście jest realizowane
przepływem przez skończony czas.
W drugim przypadku graniczne odcinki są rozdzielone
punktem krytycznym: gdyby graniczna trajektoria nie
zbliżała się do krytycznego punktu w obszarze między
poziomami, Lemat~\ref{lem:morse-transit-v11}, użyty do
tego obszaru z usuniętymi małymi otoczeniami punktów
krytycznych, ograniczałby czasy przejścia.
Ponieważ przepływ nie może zwiększać $f$, a odstępy czasu
między dwoma sąsiednimi przekrojami są jedynym miejscem
możliwej utraty zwartości, przy sklejaniu
sąsiednich granicznych odcinków ich końce muszą być tym
samym punktem krytycznym; dwie różne krytyczne wartości
pośrednie zostałyby rozdzielone przez dodatkowy poziom
$c_i$. W rezultacie otrzymujemy skończony ciąg trajektorii
łączących $p$ z $q$.

Każdy nietrywialny odcinek powoduje ścisły spadek $f$,
co daje ograniczenie liczby odcinków przez liczbę
punktów krytycznych minus jeden.
Gdy para jest Morse'a--Smale'a, z wymiaru
$\dim\mathcal M(a,b)=\lambda(a)-\lambda(b)-1\ge0$
wynika $\lambda(a)>\lambda(b)$ dla każdej niepustej
przestrzeni trajektorii $\mathcal M(a,b)$.
Sumując różnice indeksów otrzymujemy
$\ell\le\lambda(p)-\lambda(q)$.
\end{proof}

\begin{wniosek}[Skończoność trajektorii o różnicy indeksów jeden]
\label{cor:morse-rigid-finite-v11}
Niech $(f,g)$ będzie parą Morse'a--Smale'a.
Jeśli $\lambda(p)-\lambda(q)=1$, to
$\mathcal M(p,q)$ jest skończonym zbiorem.
\end{wniosek}
\begin{proof}
Z twierdzenia wynika, że każdy ciąg w $\mathcal M(p,q)$
ma podciąg zbieżny z dokładnością do łamania.
Dwa lub więcej niezerowych odcinków wymagałoby spadku indeksu
co najmniej dwa, więc granica pozostaje niełamana.
Wybór przekroju z regularnym poziomem $f=c$ identyfikuje
zbieżność modulo przesunięcia czasu ze zwykłą zbieżnością
punktów na przekroju. Zatem $\mathcal M(p,q)$ jest zwarta.
Z warunku Morse'a--Smale'a jest zerowymiarową rozmaitością,
czyli przestrzenią dyskretną; zwarta przestrzeń dyskretna
jest skończona.
\end{proof}

\begin{uwaga}[Co uzyskaliśmy bez twierdzenia o sklejaniu?]
Zwartość z dokładnością do łamania i skończoność
$\mathcal M(p,q)$ przy różnicy indeksów jeden są dowiedzione
bez używania konstrukcji gluing (przy przyjęciu
standardowego twierdzenia o rozmaitościach stabilnych
i gładkiej teorii ODE). Natomiast z tych dwóch
faktów \emph{nie wynika}, że każda trajektoria łamana
jest końcem jednowymiarowej przestrzeni trajektorii.
Dopiero twierdzenie o sklejaniu nadaje temu końcowi
strukturę przedziału i pozwala liczyć znaki brzegowe.
\end{uwaga}

\section{Analiza w otoczeniu punktu hiperbolicznego}
\begin{lemat}[Wykładniczy zanik na rozmaitości stabilnej]
\label{lem:morse-exponential-v11}
Niech $q$ będzie niezdegenerowanym punktem krytycznym $f$.
Istnieją otoczenie $V$ punktu $q$ oraz liczby $C,\mu>0$
takie, że każda trajektoria z lokalnej rozmaitości stabilnej,
która od czasu $0$ pozostaje w $V$, spełnia
$\operatorname{dist}_g(\gamma(t),q)\le
C e^{-\mu t}\operatorname{dist}_g(\gamma(0),q)$ dla $t\ge0$.
Analogiczna nierówność dla $t\le0$ zachodzi na lokalnej
rozmaitości niestabilnej.
\end{lemat}
\begin{proof}
W lokalnych współrzędnych przystosowanych do rozmaitości
stabilnej równanie na tej rozmaitości ma postać
$\dot v=A_sv+R(v)$, gdzie wszystkie wartości własne
$A_s$ mają ujemne części rzeczywiste oraz $R(0)=0$,
$\D R(0)=0$. Istnieje dodatnio określona macierz $P$
spełniająca równanie Lapunowa
$A_s^\top P+PA_s=-I$; można ją otrzymać przez
$P=\int_{t\in[0,\infty)}e^{tA_s^\top}e^{tA_s}\,\dd t$.
Dla $Q(v)=v^\top Pv$ mamy
$\dot Q=-\|v\|^2+2v^\top PR(v)$.
Po zmniejszeniu otoczenia składnik ostatni ma moduł
co najwyżej $\frac12\|v\|^2$, zatem
$\dot Q\le-\mu_0 Q$ dla pewnego $\mu_0>0$.
Równoważność norm i nierówność Gronwalla kończą dowód.
Lokalne równanie na rozmaitości niestabilnej rozpatrujemy
po odwróceniu kierunku czasu. W dowodzie użyliśmy istnienia
gładkich lokalnych rozmaitości stabilnych z twierdzenia
o rozmaitościach stabilnych.
\end{proof}

\begin{uwaga}[Dlaczego to nie wystarcza do sklejania?]
Zanik rozwiązań ustalonych na $W^s(q)$ i $W^u(q)$ nie daje
jeszcze rozwiązania, które przechodzi \emph{blisko} $q$,
ale nie należy do żadnej z tych rozmaitości. Aby je
skonstruować, trzeba kontrolować równanie dla małych
odchyleń w obu kierunkach czasu jednocześnie.
Kluczowe są oszacowania dla operatora
$\mathcal D_\gamma=\tfrac{\dd}{\dd t}-\D X(\gamma(t))$
w odpowiednich ważonych przestrzeniach funkcyjnych.
\end{uwaga}

\section{Sklejanie i orientacje jednowymiarowych przestrzeni moduli}
\begin{twierdzenie}[Model gluing w różnicy indeksów dwa]
\label{thm:morse-gluing-v11}
Niech $(f,g)$ będzie parą Morse'a--Smale'a, a punkty
krytyczne $p,q,r$ niech spełniają
$\lambda(p)-\lambda(q)=\lambda(q)-\lambda(r)=1$.
Dla każdej pary trajektorii
$([\gamma_-],[\gamma_+])\in
\mathcal M(p,q)\times\mathcal M(q,r)$
istnieje $T_0>0$ i gładka mapa
\[
\operatorname{Glue}_{\gamma_-,\gamma_+}:
(T_0,\infty)\longrightarrow\mathcal M(p,r)
\]
taka, że jej obrazy zbiegają do danej trajektorii łamanej
w sensie Twierdzenia~\ref{thm:morse-broken-compact-v11}.
Dla dostatecznie dużego $T$ mapa jest lokalnie różnowartościowa,
a każda dostatecznie bliska trajektoria niełamana należy do
jej obrazu. Parametr $e^{-T}$ opisuje topologicznie zbliżanie się do
brzegu; dla struktury gładkiego kołnierza należy dobrać
parametr sklejania zgodny z uzyskanymi oszacowaniami.
Nie twierdzimy bez dodatkowego argumentu, że sama funkcja
$e^{-T}$ jest zawsze gładką współrzędną na brzegu.
\end{twierdzenie}
\begin{proof}[Dowód warunkowy: wymagany lemat analityczny]
Przesuwamy parametryzacje $\gamma_-$ i $\gamma_+$ tak,
aby na wybranych przekrojach poprzecznych do przepływu
przy $q$ miały ustalone czasy przejścia. Wstawiamy między
ich końcowe fragmenty odcinek długości w przybliżeniu $2T$.
Wykładniczy zanik z Lematu~\ref{lem:morse-exponential-v11}
pozwala sklejać krzywe za pomocą gładkich funkcji odcięcia
w przybliżone rozwiązanie $u_T$, którego defekt
$\dot u_T-X(u_T)$ jest rzędu $O(e^{-\mu T})$.

Aby uzyskać rzeczywiste rozwiązanie, rozważamy w lokalnej
trywializacji operator
$\mathcal D_T=\tfrac{\dd}{\dd t}-\D X(u_T(t))$.
Warunek Morse'a--Smale'a daje surjektywność odpowiednich
operatorów linearyzacji wzdłuż $\gamma_-$ i $\gamma_+$.
Niezbędny jest następujący \emph{lemat analityczny o prawym
odwrotnym}: po wybraniu wag $0<\delta<\mu$ mniejszych od luk
spektralnych hesjanu w punktach końcowych i po
usunięciu kierunku translacji czasu (w odpowiednim
przekroju afinicznym) rozpatrujemy operator między
ważonymi przestrzeniami Sobolewa $W^{1,2}_\delta$
i $L^2_\delta$. Dopiero po wydzieleniu kierunku
parametru sklejania wymagamy istnienia prawego odwrotnego
$Q_T$ i stałej $C$ niezależnej od dużego $T$, takich że
$\mathcal D_TQ_T=\id$, $\|Q_T\|\le C$, w odpowiednich
normach ważonych. Buduje się ją przez sklejanie prawych
odwrotności dla $\gamma_-$ i $\gamma_+$ oraz poprawienie
błędu operatorem Neumanna. Ten lemat wymaga kontroli
jednowymiarowego kierunku związanym z długością szyjki;
nie wynika automatycznie z oddzielnej surjektywności
operatorów dla $\gamma_-$ i $\gamma_+$.

Pisząc poprawkę jako $u_T+\xi$, sprowadzamy równanie do
$\xi=-Q_T\bigl(\dot u_T-X(u_T)+N_T(\xi)\bigr)$,
gdzie $\|N_T(\xi)-N_T(\eta)\|\le
C(\|\xi\|+\|\eta\|)\|\xi-\eta\|$ dla małych norm.
Jednostajny szacunek na $Q_T$ i zanik defektu zapewniają,
że prawa strona jest kontrakcją w kuli promienia
$O(e^{-\mu T})$. Stąd otrzymujemy rodzinę prawdziwych
trajektorii i jej gładką zależność od $T$.
Argument \emph{jednoznaczności} w kołnierzu opiera się na
tej samej konstrukcji odwrotności oraz na analizie
przekrojów: każda bliska trajektoria ma jednoznaczną
normalizację czasu i długość szyjki.

Zatem główne równanie korekcyjne i oszacowania są opisane
jawnie, ale sam lemat o jednostajnym prawym odwrotnym
przyjmujemy tutaj jako standardowe twierdzenie analityczne.
Nie należy uważać powyższego schematu za jego pełny
dowód niezależny.
\end{proof}

\begin{center}
\begin{tikzpicture}[>=Stealth,font=\small,x=1cm,y=1cm]
\coordinate (P) at (-4,0.5); \coordinate (Q) at (0,0); \coordinate (R) at (4,-0.5);
\fill[manifoldblue] (P) circle (2pt);
\fill[manifoldblue] (Q) circle (2pt);
\fill[manifoldblue] (R) circle (2pt);
\draw[manifoldblue,very thick,-{Stealth[length=2mm]}] (P) .. controls (-2.7,.75) and (-1.6,.15) .. (-.25,.05);
\draw[manifoldblue,very thick,-{Stealth[length=2mm]}] (.25,-.06) .. controls (1.7,-.45) and (2.8,-.85) .. (R);
\node[above] at (P) {$p$};\node[above] at (Q) {$q$};\node[above] at (R) {$r$};
\draw[densely dotted,manifoldgold,thick] (-.9,-.65) rectangle (.9,.65);
\node[below] at (0,-.78) {otoczenie $q$; czas szyjki $2T$};
\node[above] at (-2.25,.55) {$\gamma_-$};
\node[above] at (2.3,-.25) {$\gamma_+$};
\end{tikzpicture}
\par\smallskip\small Trajektoria łamana $p\to q\to r$.
Sklejanie zastępuje dwa odcinki jedną trajektorią
przechodzącą blisko $q$ przez długi czas.
\end{center}

\begin{stwierdzenie}[Orientacje i różniczka kwadratowa]
\label{prop:morse-d2-analytic-v11}
Przy założeniach Twierdzenia~\ref{thm:morse-gluing-v11}
i zgodnych orientacjach rozmaitości niestabilnych,
zorientowane końce jednowymiarowej kompaktyfikacji
$\overline{\mathcal M}(p,r)$ są oznaczone iloczynami
$\varepsilon(\gamma_-)\varepsilon(\gamma_+)$
z jednym, wspólnym znakiem konwencji brzegowej.
W konsekwencji dla kompleksu Morse'a $\partial^2=0$,
o ile zastosujemy twierdzenie o zgodności orientacji
przy sklejaniu (niezależnie przywołane poniżej).
\end{stwierdzenie}
\begin{proof}
W pobliżu końca gluing orientacja jednowymiarowego
moduli jest orientacją parametru długości szyjki,
a znak pozostałych czynników pochodzi z orientacji
transwersalnych przecięć $W^u(p)$ z $W^s(q)$ i
$W^u(q)$ z $W^s(r)$.
Zgodność znaku tej operacji z regułą brzegu należy do
standardowej \emph{orientacyjnej części twierdzenia
o sklejaniu}. Po jej zastosowaniu każdy punkt
\(\mathcal M(p,q)\times\mathcal M(q,r)\)
wnosi do współczynnika $r$ w $\partial^2p$ swój
zorientowany znak brzegowy (pomnożony przez stały
znak globalny). Zorientowana zwarta
jednowymiarowa rozmaitość ma sumę znaków punktów
brzegowych równą zeru. Stąd każdy współczynnik
$\partial^2p$ znika. Pełny dowód reguły orientacji
wymaga porównania orientacji wyznaczników operatorów
linearyzacji i pozostaje odrębnym krokiem od
samej zwartości ODE.
\end{proof}

\section{Dokładny bilans: co jest, a czego jeszcze nie dowiedziono?}
Dzięki Lematowi~\ref{lem:morse-energy-analytic-v11},
Twierdzeniu~\ref{thm:morse-broken-compact-v11} i
Wnioskowi~\ref{cor:morse-rigid-finite-v11} mamy już
pełne uzasadnienie skończoności liczby sztywnych trajektorii,
przy użyciu jedynie standardowej teorii gładkich przepływów
i rozmaitości stabilnych. Dla jednowymiarowych moduli
pozostały dwa jasno nazwane narzędzia analityczne:
jednostajne prawe odwrotności dla operatorów gluing i
zgodność orientacji na brzegu. Wnioski homologiczne są
więc dowiedzione \emph{względem tych twierdzeń pomocniczych},
a nie przez sam rachunek energii. Jest to precyzyjniejszy
status niż stwierdzenie, że ``kompaktyfikacja wynika
ze zwartości rozwiązań ODE''.


% Rozdział do wstawienia do głównego skryptu; korzysta z jego preambuły.
\chapter{Indeks pól wektorowych i twierdzenie Poincarégo--Hopfa z brzegiem}
\label{ch:poincare-hopf-v11}

\section{Dlaczego nie powtarzamy dowodu dla rozmaitości zamkniętych?}
Twierdzenie Poincarégo--Hopfa dla zamkniętych rozmaitości
zorientowanych, wraz z definicją lokalnego indeksu pola,
zostało już udowodnione w
Sekcji~\ref{sec:poincare-hopf-16}.
Z kolei Twierdzenie~\ref{thm:lefschetz-hopf-v10} utożsamia
liczbę Lefschetza z sumą lokalnych indeksów punktów stałych.
W tym rozdziale dokładnie porównamy oba rodzaje indeksów,
a potem potraktujemy rozmaitości z brzegiem.
Nie wolno bez komentarza stosować wzoru dla rozmaitości zamkniętych
do pola, które znika lub jest styczne do brzegu: pojawia się wtedy
dodatkowy wkład brzegowy.

\section{Od zera pola do punktu stałego}
\begin{lemat}[Lokalne porównanie indeksów]
\label{lem:field-fixed-v11}
Niech $M$ będzie gładką rozmaitością wymiaru $n\ge1$,
$X$ gładkim polem wektorowym, a $p\in M$ jego izolowanym zerem.
Wybierzmy metrykę riemannowską i otoczenie $U$ punktu $p$
o zwartym domknięciu, zawarte w jednej mapie, takie że
$X$ nie znika na $\partial U$. W razie potrzeby przedłużamy metrykę i pole do otoczenia
$M$ w rozmaitości bez brzegu, po czym dla dostatecznie
małego $t>0$ określamy lokalnie
$f_t(x)=\exp_x(-tX(x))$.
Definicja indeksu dotyczy wyłącznie wewnętrznego
zera $p$, więc nie zależy od wyboru przedłużenia.
Wówczas $f_t$ ma w $U$ jedyny punkt stały $p$, po odpowiednim
zmniejszeniu $U$, i
$\operatorname{ind}_p(f_t)=\operatorname{ind}_p(X)$.
Jeśli zero jest nieosobliwe, oba indeksy są równe
$\operatorname{sgn}\det(\D_pX)$.
\end{lemat}
\begin{proof}
Przyjmijmy mapę $\varphi$ z $\varphi(p)=0$.
Zwartość $\overline U$ pozwala wybrać $t>0$ tak małe,
aby dla każdego $x\in\overline U$ wektor $-tX(x)$ leżał
w obszarze, w którym $\exp_x$ jest różnowartościowe.
Stąd $f_t(x)=x$ wtedy i tylko wtedy, gdy $X(x)=0$.
Oznaczmy $F_t=\varphi\circ f_t\circ\varphi^{-1}$ i
$V=\mathrm d\varphi(X)$.
Na zwartym $\varphi(\overline U)$ zachodzi jednostajnie
$F_t(y)=y-tV(y)+O(t^2)$, toteż
$t^{-1}(y-F_t(y))\to V(y)$ jednostajnie.
Na $\varphi(\partial U)$ wektor $V$ jest oddalony od zera.
Wobec tego dla małego $t$ homotopia
$(1-s)V(y)+s\,t^{-1}(y-F_t(y))$ nie ma zer na brzegu.
Niezmienniczość stopnia daje
$\deg(\id-F_t,U,0)=\deg(V,U,0)$.
Jeśli $\D_pX$ jest odwracalna, lokalny stopień pola
wynosi znak jej wyznacznika. Dodatkowo
$\D_pf_t=I-t\D_pX$ (po identyfikacji przestrzeni stycznej),
więc znak $\det(I-\D_pf_t)$ jest ten sam.
\end{proof}

\begin{wniosek}[Dowód Poincarégo--Hopfa przez Lefschetza]
Niech $M$ będzie zamkniętą, zorientowaną, gładką
$n$-rozmaitością, $n\ge1$, a $X$ gładkim polem z izolowanymi
zerami. Wtedy $\sum_{X(p)=0}\operatorname{ind}_p(X)=\chi(M)$.
\end{wniosek}
\begin{proof}
Dla małego $t$ odwzorowanie $f_t=\exp(-tX)$ jest globalnie
określone i homotopijne z identycznością poprzez
$s\mapsto\exp(-stX)$. Jego punkty stałe są dokładnie zerami $X$.
Dla zer nieosobliwych Twierdzenie~\ref{thm:lefschetz-hopf-v10}
oraz Lemat~\ref{lem:field-fixed-v11} dają
$\sum_p\operatorname{ind}_p(X)=L(f_t)=\chi(M)$.
W przypadku izolowanych zer osobliwych stosujemy małe,
lokalne perturbacje pola, utrzymując je bez zmian na brzegach
rozłącznych kul otaczających zera. Transwersalność do
przekroju zerowego zapewnia nieosobliwość nowych zer,
a homotopijna niezmienniczość stopnia gwarantuje, że suma
ich indeksów w każdej kuli równa się staremu indeksowi.
\end{proof}

\section{Wersja z brzegiem: pole skierowane na zewnątrz}
\begin{definicja}[Pole skierowane ściśle na zewnątrz]
Niech $M$ będzie zwartą gładką rozmaitością z brzegiem,
a $\nu$ zewnętrznym jednostkowym normalnym polem
wzdłuż $\partial M$ dla ustalonej metryki.
Pole styczne $X$ na $M$ jest \emph{ściśle skierowane na zewnątrz},
jeśli $\langle X(x),\nu(x)\rangle>0$ dla każdego
$x\in\partial M$. Warunek nie zależy od wybranej metryki:
wyraża dodatniość składowej normalnej w koorientacji brzegu.
\end{definicja}

\begin{twierdzenie}[Poincaré--Hopf dla rozmaitości z brzegiem]
\label{thm:PH-boundary-v11}
Niech $M$ będzie zwartą zorientowaną gładką
$n$-rozmaitością z brzegiem, $n\ge1$.
Jeśli gładkie pole $X\in\mathfrak X(M)$ jest ściśle skierowane
na zewnątrz na $\partial M$ i ma izolowane zera we wnętrzu,
to zbiór tych zer jest skończony oraz
\[
\sum_{p\in\operatorname{int}M:\,X(p)=0}\operatorname{ind}_p(X)=\chi(M).
\]
\end{twierdzenie}
\begin{proof}
Zera nie mogą zbliżać się do brzegu, bo pole ma tam
niezerową składową normalną. Zatem tworzą zwarty dyskretny
podzbiór wnętrza i jest ich skończenie wiele.
Najpierw zakładamy, że są nieosobliwe.
Wybierzmy metrykę i pole przedłużone do otoczenia
$M$ w jej podwojeniu (lub innej rozmaitości bez brzegu).
Tam definiujemy mapę wykładniczą. Przedłużenie nie
musi zachowywać warunku skierowania na zewnątrz
poza $M$, ponieważ rozważamy punkty $x\in M$.
Dla dostatecznie małego $t>0$ odwzorowanie
$f_t(x)=\exp_x(-tX(x))$ jest określone na całym $M$
i przyjmuje wartości w $M$: wewnątrz wynika to z małości
$t$, a przy brzegu wektor $-X$ jest ściśle skierowany
do wnętrza; przez zwartość warunek jest jednostajny
na kołnierzu. Mamy $f_t\simeq\id_M$ przez tę samą
rodzinę z parametrem $st$, a punkty stałe $f_t$ są
wyłącznie wewnętrznymi zerami $X$.

Dla ścisłego zastosowania twierdzenia indeksowego
przechodzimy do podwojenia $DM$ rozmaitości $M$.
Niech $j:M\hookrightarrow DM$ oznacza pierwszą kopię,
a $r:DM\to M$ odwzorowanie składające obie kopie
na pierwszą. Spełnia ono $r\circ j=\id_M$.
Odwzorowanie $F=j\circ f_t\circ r:DM\to DM$ ma dokładnie
te same punkty stałe co $f_t$; wszystkie leżą we wnętrzu
pierwszej kopii. Na ich otoczeniach $F$ jest gładkie.
Po wygładzeniu $F$ poza tymi otoczeniami, względnie do nich,
możemy zastosować formułę Lefschetza--Hopfa na zamkniętej
rozmaitości $DM$. Wygładzenie nie tworzy nowych punktów
stałych, gdyż na zwartym dopełnieniu małych otoczeń
odległość punktu od jego obrazu ma dodatnie minimum.

Cykliczność śladu dla $j_*,r_*$ i $r_*j_*=\id$ daje
$L(F)=L(f_t)=L(\id_M)=\chi(M)$.
Lokalne indeksy $F$ i $f_t$ są identyczne, zaś
Lemat~\ref{lem:field-fixed-v11} identyfikuje je z indeksami
pola $X$. Otrzymujemy żądaną równość.
Dla izolowanych zer osobliwych perturbujemy $X$ w
skończenie wielu rozłącznych kulach we wnętrzu,
zachowując warunek na brzegu i indeksy lokalne,
po czym stosujemy przypadek nieosobliwy.
\end{proof}

\begin{przyklad}[Dysk, kula i pole radialne]
Na $D^n\subset\R^n$ pole $X(x)=x$ jest ściśle zewnętrzne
na brzegu i ma jedno zero w $0$, którego indeks wynosi $+1$.
Stąd $\chi(D^n)=1$. Na kuli z usuniętym wnętrzem mniejszej
współśrodkowej kuli pole $X(x)=x$ nie jest zewnętrzne
na wewnętrznej składowej brzegu; nie można więc bez
zmiany hipotez stosować powyższego twierdzenia.
\end{przyklad}

\begin{przyklad}[Pole gradientowe funkcji Morse'a]
Niech $M$ będzie zwartą rozmaitością bez brzegu,
a $f:M\to\R$ funkcją Morse'a. Dla $X=\operatorname{grad}f$
każde zero $p$ jest nieosobliwe i
$\operatorname{ind}_p(X)=(-1)^{\lambda(p)}$.
Dlatego twierdzenie Poincarégo--Hopfa daje
$\chi(M)=\sum_p(-1)^{\lambda(p)}$.
Jest to ta sama równość, którą otrzymaliśmy algebraicznie
z kompleksu Morse'a, ale interpretacja lokalna jest inna.
\end{przyklad}

\section{Co się zmienia, gdy pole jest styczne do brzegu?}
\begin{uwaga}
Jeśli $X$ jest tylko styczne do $\partial M$ albo jego
składowa normalna zmienia znak, suma indeksów zer wewnętrznych
nie musi równać się $\chi(M)$. Trzeba uwzględnić warunek
brzegowy, np. stopień części normalizowanej pola na brzegu
lub odpowiedni indeks brzegowy. Nie podajemy niepoprawnego
wzoru bez ustalenia konwencji dla tych indeksów.
\end{uwaga}


% Iteracja XI: samodzielne wprowadzenie do nośników zwartych i Borel--Moore'a.


% Rozdział do wstawienia do głównego skryptu; korzysta z jego preambuły.
\chapter{Kohomologia z nośnikami zwartymi i homologia Borel--Moore'a}
\label{ch:compact-bm-v11}

\section{Motywacja: dlaczego zwykła dualność nie wystarcza?}
Dla zwartej zorientowanej $n$-rozmaitości $M$ znamy izomorfizm
$H^k(M;\Z)\cong H_{n-k}(M;\Z)$. Na prostej $\R$ nie można
bez zmian napisać tego samego: $H^0(\R;\Z)=\Z$, lecz
$H_1(\R;\Z)=0$. Intuicja jest geometryczna.
Otwarta prosta ma klasę fundamentalną reprezentowaną przez
\emph{nieskończoną, ale lokalnie skończoną} sumę kolejnych
zorientowanych odcinków. Nie stanowi ona zwykłego skończonego
łańcucha singularnego. Dualnie potrzebujemy form i kołańcuchów,
które mają nośnik zwarty i mogą być całkowane po takiej sumie.
Otrzymamy dwa różne izomorfizmy, których nie wolno mylić:
\[
H_c^k(M;G)\cong H_{n-k}(M;G),\qquad
H^k(M;G)\cong H_{n-k}^{\mathrm{BM}}(M;G)
\]
(dla rozmaitości zorientowanych, bez brzegu; w przypadku
nieorientowalnym należy skręcić współczynniki systemem orientacji).

\section{Kohomologia z nośnikami zwartymi}
\begin{definicja}[Definicja przez pary względne]
\label{def:compact-support-cohom-v11}
Niech $X$ będzie lokalnie zwartą przestrzenią Hausdorffa,
a $G$ grupą abelową. Przez $\mathcal K(X)$ oznaczmy
skierowany zbiór zwartych podzbiorów $X$ uporządkowany
przez inkluzję. Definiujemy
\[
H_c^k(X;G):=\varinjlim_{K\in\mathcal K(X)}
H^k(X,X\setminus K;G).
\]
Jeżeli $K\subseteq L$, to zawieranie par
$(X,X\setminus L)\to(X,X\setminus K)$ daje przez
kontrawariantność kohomologii morfizm
$H^k(X,X\setminus K;G)\to H^k(X,X\setminus L;G)$.
To właśnie są strzałki układu skierowanego.
\end{definicja}

\begin{uwaga}[Nośnik kołańcucha a nośnik klasy]
Klasa $H_c^k$ nie oznacza dowolnego kołańcucha singularnego
zerowego poza skończoną listą sympleksów. Definicja względna
mówi, że klasa może być reprezentowana \emph{względem dopełnienia}
pewnego zwartego $K$. Dla lokalnie skończonych triangulacji
rozmaitości porównamy ją z kołańcuchami komórkowymi o skończonym
nośniku. W rozdziale o kohomologii snopów (Rozdział~\ref{ch:snopy-verdier-v12}) rozróżniamy
ściśle: nośnik sekcji snopu, nośnik względnej klasy
singularnej oraz lokalnie skończony nośnik łańcucha Borel--Moore'a. Ta równoważność wymaga twierdzenia o porównaniu
konstrukcji komórkowej i singularnej.
\end{uwaga}

\begin{stwierdzenie}[Przestrzeń zwarta i mapa właściwa]
\label{prop:compact-proper-v11}
Jeśli $X$ jest zwarta, to $H_c^k(X;G)\cong H^k(X;G)$.
Jeżeli $f:X\to Y$ jest ciągłym odwzorowaniem \emph{właściwym},
tj. przeciwobraz każdego zwartego podzbioru $Y$ jest zwarty,
to $f$ indukuje homomorfizm
$f^*:H_c^k(Y;G)\to H_c^k(X;G)$.
\end{stwierdzenie}
\begin{proof}
Dla zwartego $X$ zbiór $K=X$ jest największym elementem
$\mathcal K(X)$; granica prosta równa się
$H^k(X,X\setminus X;G)=H^k(X;G)$.
Dla mapy właściwej i zwartego $K\subseteq Y$ mamy
$L=f^{-1}(K)$ zwarty, a mapę par
$f:(X,X\setminus L)\to(Y,Y\setminus K)$.
Daje ona morfizm w kohomologii względnej. Naturalność
względem inkluzji zbiorów zwartych pozwala przejść
do granic prostych. Bez właściwości $f$ przeciwobraz
$K$ może nie być zwarty, więc konstrukcja nie musi działać.
\end{proof}

\begin{przyklad}[Dlaczego $H_c^0(\R)=0$?]
Na spójnej niezwartej rozmaitości bez brzegu
każda lokalnie stała funkcja o nośniku zwartym jest zerowa.
Z tego powodu $H_c^0(\R;\Z)=0$, podczas gdy
$H^0(\R;\Z)=\Z$. Na przestrzeni dyskretnej $\N$
klasa w stopniu zerowym może być niezerowa tylko na
skończenie wielu punktach, zatem
$H_c^0(\N;\Z)\cong\bigoplus_{m\in\N}\Z$.
\end{przyklad}

\section{Kompaktowo podparte formy różniczkowe}
\begin{definicja}[Kompleks de Rhama z nośnikami zwartymi]
Dla gładkiej rozmaitości $M$ niech $\Omega_c^k(M)$ będzie
przestrzenią gładkich $k$-form $\omega$ o zwartym nośniku
$\operatorname{supp}\omega=\overline{\{x:\omega_x\ne0\}}$.
Ponieważ $\operatorname{supp}(\dd\omega)\subseteq
\operatorname{supp}(\omega)$, różniczka zewnętrzna definiuje
kompleks $(\Omega_c^*(M),\dd)$. Jego kohomologię oznaczamy
$H_{c,\mathrm{dR}}^k(M)$.
\end{definicja}

\begin{lemat}[Całkowanie na prostej]
\label{lem:compact-R-v11}
Odwzorowanie $\omega=f(x)\,\dd x\mapsto
\int_{x\in\R}f(x)\,\dd x$ indukuje izomorfizmy
$H_{c,\mathrm{dR}}^1(\R)\cong\R$ oraz
$H_{c,\mathrm{dR}}^0(\R)=0$.
\end{lemat}
\begin{proof}
Jeśli $f\in C_c^\infty(\R)$, to $f\,\dd x$ jest zamknięta.
Jeżeli $f\,\dd x=\dd F$ dla pewnej
$F\in C_c^\infty(\R)$, to
$\int_{x\in\R}f(x)\,\dd x=F(+\infty)-F(-\infty)=0$.
Odwrotnie, gdy ta całka jest zerowa, funkcja
$F(x)=\int_{u\in(-\infty,x]} f(u)\,\dd u$ ma nośnik
zwarty: na lewo od nośnika $f$ i na prawo od niego
jest równa zeru. Zatem $\dd F=f\,\dd x$.
Istnieje $f$ o całce $1$, co dowodzi surjektywności.
Dla form stopnia $0$ warunek $\dd F=0$ oznacza
stałość $F$, a jedyna stała funkcja o nośniku zwartym
na $\R$ jest zerowa.
\end{proof}

\begin{twierdzenie}[Lokalny rachunek na $\R^n$]
\label{thm:compact-Rn-v11}
Dla $n\ge1$ zachodzi
$H_{c,\mathrm{dR}}^k(\R^n)=0$ przy $k\ne n$ oraz
$H_{c,\mathrm{dR}}^n(\R^n)\cong\R$.
W stopniu $n$ izomorfizm jest dany przez całkowanie form.
\end{twierdzenie}
\begin{proof}
Przedstawimy indukcję wykorzystującą $\R^n=\R^{n-1}\times\R$.
Wybierzmy $\rho\in C_c^\infty(\R)$ o
$\int_{t\in\R}\rho(t)\,\dd t=1$.
Rozłóżmy formę w sposób jednoznaczny:
$\omega=\alpha(x,t)+\dd t\wedge\beta(x,t)$,
gdzie $\alpha$ i $\beta$ nie zawierają $\dd t$.
Operator całkowania po włóknie określamy wzorem
$I(\omega)=\int_{t\in\R}\beta(x,t)\,\dd t$.
Jest to forma stopnia o jeden mniejszego, o nośniku zwartym,
na $\R^{n-1}$. Zapisując $\dd_x$ dla różniczki względem $x$,
mamy
$\dd\omega=\dd_x\alpha+\dd t\wedge
(\partial_t\alpha-\dd_x\beta)$, a więc
$I\dd=-\dd_x I$. Definiujemy też
$s(\eta)=\rho(t)\,\dd t\wedge\eta$;
wówczas $I\circ s=\id$ i $\dd s=-s\dd_x$.

Dla formy poziomej $\beta$ niech
$P\beta(x,t)=\int_{u\in(-\infty,t]}\beta(x,u)\,\dd u-
\bigl(\int_{u\in\R}\beta(x,u)\,\dd u\bigr)R(t)$,
gdzie $R(t)=\int_{u\in(-\infty,t]}\rho(u)\,\dd u$.
Wynik ma nośnik zwarty, gdyż poza dużym przedziałem
oba wyrazy się znoszą, również jednostajnie względem $x$.
Określmy $H(\omega)=P\beta$. Mamy
$\partial_t P\beta=\beta-\rho\,I\omega$ i
$P(\partial_t\alpha)=\alpha$ dla form $\alpha$ o
nośniku zwartym, a $P\dd_x=\dd_xP$.
Stąd przez bezpośrednie podstawienie
$\dd H\omega+H\dd\omega=
\alpha+\dd t\wedge(\beta-\rho\,I\omega)
=\omega-sI\omega$.
Jest to jawna homotopia łańcuchowa; mimo że $I$ i $s$
zmieniają stopień o jeden, indukują odwrotne
izomorfizmy kohomologii po przesunięciu stopnia. Zatem
$H_{c,\mathrm{dR}}^k(\R^n)\cong
H_{c,\mathrm{dR}}^{k-1}(\R^{n-1})$.
Początek indukcji daje Lemat~\ref{lem:compact-R-v11}.
\end{proof}

\begin{twierdzenie}[De Rham z nośnikami zwartymi]
\label{thm:compact-deRham-v11}
Dla gładkiej drugoprzeliczalnej rozmaitości $M$ całkowanie
form po gładkich sympleksach singularnych wyznacza
naturalne izomorfizmy
$H_{c,\mathrm{dR}}^k(M)\cong H_c^k(M;\R)$.
\end{twierdzenie}
\begin{proof}[Argument wykorzystujący globalne twierdzenie porównawcze]
Wyjaśniamy zakres porównania. Na otwartym zbiorze
współrzędnych izomorfizm otrzymuje się ze zwykłego
lematu Poincarégo i wersji z nośnikami zwartymi
z Twierdzenia~\ref{thm:compact-Rn-v11}.
Dla otwartego pokrycia $M=U\cup V$ istnieje krótki ciąg
kompleksów
$0\to\Omega_c^*(U\cap V)\to
\Omega_c^*(U)\oplus\Omega_c^*(V)\to
\Omega_c^*(M)\to0$:
pierwsza mapa to $(\alpha,-\alpha)$ przez przedłużenie
zerem, druga jest sumą, a surjektywność wynika z podziału
jedności na zwartym nośniku formy. Prowadzi to do ciągu
Mayera--Vietorisa. Analogiczny długi ciąg Mayera--Vietorisa istnieje
w singularnej kohomologii z nośnikami zwartymi,
ale nie wynika z samej definicji przez pary:
trzeba najpierw zastosować ekscyzję i odpowiednią
konstrukcję z nośnikami (lub rezolwenty snopowe).
Porównanie map łączących i lemat pięciu przenoszą
izomorfizm z otwartych kostek na ich skończone sumy,
a następnie na wszystkie rozmaitości przy użyciu
lokalnie skończonego pokrycia i wyczerpania przez
zwarte podzbiory. Ostatni krok można sformułować jako
standardowe twierdzenie o rozdzielczości kompleksu
stałych snopów przez formy różniczkowe; nie utożsamiamy
go z pojedynczym zastosowaniem lematu Poincarégo.
\end{proof}

\section{Łańcuchy lokalnie skończone i homologia Borel--Moore'a}
\begin{definicja}[Kompleks lokalnie skończonych łańcuchów]
\label{def:BM-chain-v11}
Niech $X$ będzie lokalnie zwartą przestrzenią Hausdorffa.
\emph{Lokalnie skończonym singularnym $k$-łańcuchem}
o współczynnikach w $G$ nazywamy formalną sumę
$c=\sum_{\sigma:\Delta^k\to X}a_\sigma\sigma$,
w której dla każdego zwartego $K\subset X$ tylko
skończenie wiele sympleksów o $a_\sigma\ne0$
ma obraz przecinający $K$. Zbiór takich łańcuchów
oznaczamy $C_k^{\mathrm{lf}}(X;G)$.
Różniczkę brzegową definiujemy na składnikach jak
w homologii singularnej; jest dobrze określona,
ponieważ obraz ściany sympleksu zawiera się w obrazie
samego sympleksu. Homologię tego kompleksu oznaczamy
$H_k^{\mathrm{BM}}(X;G)$ i nazywamy
\emph{homologią Borel--Moore'a} (lub homologią lokalnie skończoną).
\end{definicja}

\begin{stwierdzenie}[Zwartość i właściwość]
\label{prop:BM-proper-v11}
Jeżeli $X$ jest zwarta, to $H_k^{\mathrm{BM}}(X;G)=H_k(X;G)$.
Jeżeli $f:X\to Y$ jest właściwym odwzorowaniem ciągłym
między lokalnie zwartymi przestrzeniami Hausdorffa,
to indukuje $f_*:H_k^{\mathrm{BM}}(X;G)\to H_k^{\mathrm{BM}}(Y;G)$.
\end{stwierdzenie}
\begin{proof}
W przestrzeni zwartej warunek lokalnej skończoności
zastosowany do $K=X$ wymusza skończoność całej sumy.
Jeśli $c$ jest lokalnie skończony, a $L\subset Y$ jest zwarty,
to $f^{-1}(L)$ jest zwarty. Tylko skończenie wiele
sympleksów z $c$ przecina ten przeciwobraz, toteż
ich obrazy przez $f$ przecinają $L$ jedynie skończenie
wiele razy; po zsumowaniu współczynników jednakowych
obrazów nadal otrzymujemy rodzinę lokalnie skończoną.
Mapa łańcuchowa jest zatem określona
poprawnie. Przy niewłaściwych mapach ta własność
może zawieść; np. mapa stała $\R\to\{*\}$
wysyłałaby nieskończony lokalnie skończony łańcuch
na nieskończoną sumę w punkcie.
\end{proof}

\begin{center}
\begin{tikzpicture}[>=Stealth,font=\small,x=0.8cm,y=0.8cm]
\fill[black!7] (-1.05,-.55) rectangle (1.05,.50);
\draw[thick] (-4.4,0)--(4.4,0);
\foreach \m in {-4,-3,-2,-1,0,1,2,3,4}
 {\fill[manifoldblue] (\m,0) circle (1.5pt);}
\foreach \m in {-4,-3,-2,-1,0,1,2,3}
 {\draw[->,manifoldblue,semithick] (\m+.12,.12)--(\m+.86,.12);}
\node[anchor=east] at (-4.5,.09) {$\cdots$};
\node[anchor=west] at (4.5,.09) {$\cdots$};
\node[below] at (-1,-.10) {$-1$};
\node[below] at (0,-.10) {$0$};
\node[below] at (1,-.10) {$1$};
\node[above] at (0,.53) {$K=[-1,1]$};
\end{tikzpicture}
\par\smallskip\small Schemat lokalnie skończonego łańcucha na prostej:
każdy zwarty odcinek przecina tylko skończenie wiele składników.
\end{center}

\begin{przyklad}[Nieskończony cykl na prostej]
\label{ex:BM-line-v11}
Niech $\sigma_m:[0,1]\to\R$ będzie dane przez
$\sigma_m(t)=m+t$, $m\in\Z$. Suma
$c=\sum_{m\in\Z}\sigma_m$ jest lokalnie skończona;
jej brzeg znika, ponieważ przy każdym wierzchołku
$m\in\Z$ wkład końca $\sigma_{m-1}$ znosi wkład
początku $\sigma_m$. Nie jest to skończony łańcuch
singularny. Reprezentuje generator
$H_1^{\mathrm{BM}}(\R;\Z)\cong\Z$.
Natomiast łańcuch $\sum_{m\ge0}\sigma_m$ ma brzeg
$-[0]$, więc $[0]$ znika w $H_0^{\mathrm{BM}}(\R;\Z)$.
Wartości grup zostaną uzasadnione przez dualność poniżej.
\end{przyklad}

\begin{przyklad}[Dlaczego jednopunktowe uzwarcenie bywa mylące?]
Dla nieskończonej przestrzeni dyskretnej $X=\N$
mamy $H_0^{\mathrm{BM}}(X;\Z)=\prod_{m\in\N}\Z$,
ponieważ lokalnie skończona suma punktów może mieć
dowolnie wiele niezerowych współczynników, a nie
ma nietrywialnych granic w stopniu zero.
Jednocześnie $H_0(X;\Z)=\bigoplus_{m\in\N}\Z$.
Nie wolno więc bez dodatkowych warunków utożsamiać
lokalnie skończonej homologii singularnej z
$H_*(X^+,\{\infty\};\Z)$: dla uzwarcenia
jednopunktowego przeliczalnej przestrzeni dyskretnej
ostatnia grupa w stopniu zero jest sumą prostą,
a nie produktem. Takie utożsamienie wymaga starannego
wyboru teorii homologii i hipotez o przestrzeni.
\end{przyklad}

\section{Dualność dla rozmaitości niezwartych}
\begin{definicja}[Lokalnie skończona klasa fundamentalna]
Niech $M$ będzie spójną zorientowaną gładką
$n$-rozmaitością bez brzegu. Triangulacja lokalnie
skończona, zorientowana zgodnie z orientacją $M$,
wyznacza lokalnie skończony $n$-cykl złożony z
orientowanych sympleksów najwyższego wymiaru.
Jego klasa $[M]_{\mathrm{BM}}\in H_n^{\mathrm{BM}}(M;\Z)$
nazywa się klasą fundamentalną Borel--Moore'a.
Jej lokalny obraz w $H_n(M,M\setminus\{x\};\Z)$
jest lokalnym generatorem zgodnym z orientacją.
\end{definicja}

\begin{twierdzenie}[Dualność Poincarégo w wersji niezwartej]
\label{thm:noncompact-PD-v11}
Niech $M$ będzie gładką, drugoprzeliczalną,
zorientowaną $n$-rozmaitością bez brzegu,
a $G$ dowolną grupą abelową. Iloczyn z klasą
$[M]_{\mathrm{BM}}$ indukuje naturalne izomorfizmy
\[
H_c^k(M;G)\cong H_{n-k}(M;G),\qquad
H^k(M;G)\cong H_{n-k}^{\mathrm{BM}}(M;G)
\]
dla każdego $k\in\Z$. W przypadku rozmaitości
nieorientowalnej w grupach homologii po prawej
występuje system orientacji $\mathcal O_M\otimes_{\Z}G$.
\end{twierdzenie}
\begin{proof}
Przedstawimy model komórkowy realizujący oba izomorfizmy.
Gładka rozmaitość drugoprzeliczalna ma lokalnie skończoną
triangulację $K$. Po odpowiednim podziale barycentrycznym
każdemu zorientowanemu $k$-sympleksowi $\sigma$ odpowiada
blok dualny (komórka z podziału barycentrycznego) $D(\sigma)$ wymiaru $n-k$, złożona z
sympleksów barycentrycznych łańcuchów
$\sigma=\sigma_k<\sigma_{k+1}<\cdots<\sigma_n$.
Lokalna skończoność $K$ sprawia, że $D(\sigma)$ jest
skończonym łańcuchem i że rodzina komórek dualnych
jest lokalnie skończona.

Wybór orientacji $M$ wraz z orientacją $\sigma$
ustala orientację każdego bloku $D(\sigma)$.
W wymiarze zerowym blok dualny odpowiada
$n$-sympleksowi; dla lokalnie skończonej triangulacji
bloki są skończone, choć może ich być nieskończenie wiele.
Reguła orientacji brzegu implikuje,
że współczynnik $D(\tau)$ w $\partial D(\sigma)$
jest, z globalnie ustaloną konwencją znaków,
współczynnikiem $\sigma$ w $\delta\tau$.
W ten sposób operator przypisujący
$\sigma^*\mapsto D(\sigma)$ (gdzie $\sigma^*$ jest
kołańcuchem bazowym) daje izomorfizm kompleksów
z odpowiednim przesunięciem stopni między
komórkowymi kołańcuchami $K$ i dualnymi łańcuchami.

Zwykłe kołańcuchy komórkowe są \emph{dowolnymi}
rodzinami współczynników, tj.
$C^k(K;G)=\prod_{\sigma\in K^{(k)}}G$.
Odpowiadają one lokalnie skończonym łańcuchom
komórek dualnych. Kołańcuchy o nośniku zwartym
mają natomiast tylko skończenie wiele niezerowych
współczynników, czyli tworzą
$C_c^k(K;G)=\bigoplus_{\sigma\in K^{(k)}}G$;
odpowiadają zwykłym skończonym łańcuchom dualnym.
Komórkowo-singularne twierdzenie porównawcze,
stosowane do lokalnie skończonej triangulacji
i do łańcuchów lokalnie skończonych, pozwala
przejść od tych dwóch izomorfizmów kompleksów
do grup wymienionych w tezie. Konstrukcja jest
lokalna i odpowiada iloczynowi z klasą fundamentalną.

W przypadku nieorientowalnym orientacje komórek dualnych
nie sklejają się przy przejściu między mapami.
Ich zmiany znaków są dokładnie funkcjami przejścia
lokalnego systemu $\mathcal O_M$; ten system
musi wystąpić po stronie homologii.
\end{proof}

\begin{uwaga}[Zakres dowodu komórkowego]
Użyliśmy dwóch standardowych wyników topologii
rozmaitości: triangulowalności \emph{gładkich}
rozmaitości i zgodności homologii lokalnie
skończonej komórkowej z singularną. Nie rozszerzamy
tego argumentu bez dodatkowych hipotez na dowolne
niezwarte przestrzenie topologiczne ani na
niektóre nietriangulowalne rozmaitości topologiczne.
\end{uwaga}

\section{Przykłady i obliczenia grup}
\begin{przyklad}[Przestrzeń euklidesowa]
Dla $n\ge1$ znamy $H^0(\R^n;\Z)=\Z$ i
$H^k(\R^n;\Z)=0$ dla $k>0$.
Z dualności otrzymujemy
\[
H_k^{\mathrm{BM}}(\R^n;\Z)=
\begin{cases}\Z,&k=n,\\0,&k\ne n,\end{cases}
\qquad
H_c^k(\R^n;\Z)=
\begin{cases}\Z,&k=n,\\0,&k\ne n.\end{cases}
\]
Odtwarzamy w ten sposób rachunek form zwartych
z Twierdzenia~\ref{thm:compact-Rn-v11},
ale teraz również z całkowitymi współczynnikami.
\end{przyklad}

\begin{przyklad}[Otwarty cylinder]
Niech $M=S^1\times\R$, z orientacją produktową.
Ma typ homotopijny $S^1$, zatem $H_0(M)=H_1(M)=\Z$,
a pozostałe grupy homologii są zerowe.
Z dualności w wymiarze $2$ otrzymujemy
$H_c^0(M)=0$, $H_c^1(M)=\Z$, $H_c^2(M)=\Z$ oraz
$H_0^{\mathrm{BM}}(M)=0$, $H_1^{\mathrm{BM}}(M)=\Z$,
$H_2^{\mathrm{BM}}(M)=\Z$.
W wersji de Rhama generator stopnia $1$ reprezentuje
$\rho(t)\,\dd t$, a stopnia $2$ forma
$\dd\theta\wedge\rho(t)\,\dd t$, gdzie
$\int_{t\in\R}\rho(t)\,\dd t=1$ i $\dd\theta$ oznacza
standardową globalną formę kątową na okręgu.
\end{przyklad}

\begin{przyklad}[Wnętrze zwartej rozmaitości z brzegiem]
Niech $W$ będzie zwartą gładką $n$-rozmaitością
z brzegiem i niech $M=\operatorname{int}W$.
Kołnierz brzegu i ekscyzja, w połączeniu z
porównaniem komórkowym, dają
$H_c^k(M;G)\cong H^k(W,\partial W;G)$ oraz
$H_k^{\mathrm{BM}}(M;G)\cong H_k(W,\partial W;G)$.
Dla $W=D^n$ uzyskujemy
$H_n^{\mathrm{BM}}(\R^n;\Z)=\Z$,
choć zwykła homologia $H_n(\R^n;\Z)$ jest zerowa.
\end{przyklad}

\section{Związek z dualnością Poincarégo--Lefschetza}
Jeśli $M=\operatorname{int}W$ dla zwartej rozmaitości
z brzegiem, to izomorfizmy
$H_c^k(M)\cong H^k(W,\partial W)$ oraz
$H_{n-k}(M)\cong H_{n-k}(W)$
zamieniają pierwszą dualność z
Twierdzenia~\ref{thm:noncompact-PD-v11}
w dualność Poincarégo--Lefschetza dla pary
$(W,\partial W)$. Druga dualność odpowiada
$H^k(W)\cong H_{n-k}(W,\partial W)$.
W ten sposób oba równania nie są osobnymi
przypadkowymi analogiami: są dwiema wersjami
tego samego mechanizmu geometrycznego.

\section{Wskazówki przed połączeniem z geometrią algebraiczną}
Dla gładkiej zespolonej rozmaitości algebraicznej
wymiaru zespolonego $d$ występuje naturalna orientacja
rzeczywistej $2d$-rozmaitości. Grupy
$H_{2d}^{\mathrm{BM}}$ są właściwym miejscem dla klas
fundamentalnych niezwartych podrozmaitości, a
$H_c^*(\cdot)$ jest potrzebna do dualności i map
właściwych. W następnym wspólnym skrypcie można
połączyć te konstrukcje z cyklami algebraicznymi,
przecięciami i kohomologią snopów; utożsamienie
cykli algebraicznych z wszystkimi klasami topologicznymi
nie jest jednak prawdziwe bez dodatkowych warunków.


\chapter{Kohomologia snopów z nośnikami i podstawy dualności Verdiera}
\label{ch:snopy-verdier-v12}

\section{Dlaczego teoria snopów jest potrzebna?}
W rozdziale~\ref{ch:compact-bm-v11} spotkaliśmy dwie postacie dualności
Poincarégo dla niezwartych rozmaitości: jedną z kohomologią o nośnikach
zwartych, drugą z homologią Borel--Moore'a. Ich elementarne dowody
wykorzystywały lokalnie skończone triangulacje. Tymczasem w geometrii
algebraicznej często pracujemy na przestrzeniach osobliwych, a zamiast
jednej grupy współczynników pojawiają się dane zmieniające się lokalnie.
Potrzebujemy konstrukcji, która jednocześnie pamięta (i) współczynniki
lokalne, (ii) warunki na nośniki oraz (iii) sposób, w jaki lokalne
dane sklejają się w globalne.

Snop jest właśnie uporządkowanym systemem takich danych. Kohomologia
snopów mierzy przeszkody w ich globalnym sklejaniu, zaś dualność Verdiera
jest snopowym uogólnieniem dualności Poincarégo. Rozpoczniemy od
konstrukcji elementarnych, a twierdzenia wymagające kategorii pochodnych
wyraźnie oddzielimy od wyników, których dowody przedstawimy tutaj.

W całym rozdziale $k$ jest ustalonym ciałem, chyba że jawnie podamy inny
pierścień współczynników. Przestrzenie są Hausdorffa i lokalnie zwarte,
a w twierdzeniach porównawczych i o nośnikach dodatkowo parazwarte.
Kohomologia jest indeksowana \emph{kohomologicznie}: różniczka podnosi
stopień o jeden. Zapis $\mathcal F$ oznacza snop $k$-przestrzeni wektorowych.

\section{Przedsnopy, snopy i włókna snopu}
\begin{definicja}[Przedsnop i warunek snopowy]
\label{def:sheaf-v12}
Niech $X$ będzie przestrzenią topologiczną. Przedsnopem
$k$-przestrzeni wektorowych $\mathcal F$ na $X$ nazywamy przyporządkowanie
każdemu otwartemu $U\subseteq X$ przestrzeni $\mathcal F(U)$ i każdemu
zawieraniu $V\subseteq U$ liniowej mapy ograniczenia
$\rho^U_V:\mathcal F(U)\to\mathcal F(V)$, przy czym
$\rho^U_U=\id$ i $\rho^V_W\rho^U_V=\rho^U_W$ dla $W\subseteq V\subseteq U$.
Przedsnop jest \emph{snopem}, jeśli dla każdego otwartego $U$ oraz
pokrycia $U=\bigcup_iU_i$ zachodzą:
\begin{enumerate}[label=(\roman*)]
\item jeśli $s,t\in\mathcal F(U)$ mają równe ograniczenia do każdego $U_i$,
  to $s=t$ (jednoznaczność);
\item jeśli $s_i\in\mathcal F(U_i)$ i
  $s_i|_{U_i\cap U_j}=s_j|_{U_i\cap U_j}$ dla wszystkich $i,j$,
  istnieje $s\in\mathcal F(U)$ o $s|_{U_i}=s_i$ (istnienie sklejenia).
\end{enumerate}
\end{definicja}

\begin{przyklad}[Snopy funkcji, form i rozwiązań lokalnych]
Na gładkiej rozmaitości $M$ przyporządkowania
$U\mapsto C^\infty(U)$ oraz $U\mapsto\Omega^q(U)$ są snopami,
podobnie jak $U\mapsto\mathcal O(U)$ dla holomorficznych funkcji na
zespolonej rozmaitości. Dla ustalonego operatora różniczkowego $P$
przestrzenie $\{f\in C^\infty(U):Pf=0\}$ również tworzą snop
(jeśli operator jest zgodny z ograniczeniami).
Natomiast przyporządkowanie $U\mapsto H^q(U;k)$ dla $q>0$ nie jest
na ogół snopem: globalna klasa kohomologii może znikać na wszystkich
zbiorach drobnego pokrycia, nie będąc zerowa globalnie.
Na przykład generator $H^1(S^1;k)$ znika po ograniczeniu
do każdej z dwóch dostatecznie dużych otwartych łuków.
\end{przyklad}

\begin{definicja}[Włókno snopu w punkcie]
Dla snopu $\mathcal F$ i $x\in X$ jego \emph{włóknem} (stalkiem)
nazywamy granicę prostą
$\mathcal F_x=\varinjlim_{x\in U}\mathcal F(U)$, gdzie przejścia
są ograniczeniami, a $U$ przebiegają otwarte otoczenia $x$.
Element $\mathcal F_x$ jest klasą pary $(U,s)$, $x\in U$,
z utożsamieniem dwóch par, jeśli sekcje zgadzają się
na pewnym mniejszym otoczeniu $x$.
\end{definicja}

\begin{lemat}[Dokładność sprawdza się na włóknach]
\label{lem:stalk-exact-v12}
Ciąg morfizmów snopów $k$-przestrzeni wektorowych
$\mathcal A\xrightarrow{u}\mathcal B\xrightarrow{v}\mathcal C$
jest dokładny (czyli $\operatorname{im}u=\ker v$ jako snopy)
wtedy i tylko wtedy, gdy ciąg
$\mathcal A_x\xrightarrow{u_x}\mathcal B_x\xrightarrow{v_x}\mathcal C_x$
jest dokładny dla każdego $x\in X$.
\end{lemat}
\begin{proof}
Warunek jądra jest lokalny, ponieważ dla sekcji $b$ równość $v(b)=0$
można sprawdzać na wszystkich jej kiełkach. Obraz \emph{snopowy}
morfizmu $u$ składa się z sekcji, które lokalnie należą do obrazu $u$.
Zatem kiełek $b_x$ należy do $\operatorname{im}u_x$ dokładnie wtedy,
gdy po zmniejszeniu otoczenia $x$ sekcja $b$ ma lokalny przeciwobraz.
To równoważność dokładności w punkcie i dokładności snopowej.
Nie wynika stąd, że ciąg przestrzeni sekcji globalnych jest dokładny:
lokalne przeciwobrazy mogą nie sklejać się globalnie.
\end{proof}

\begin{przyklad}[Snop stały a system lokalny]
Na spójnej przestrzeni lokalnie łukowo spójnej snop $k_X$
przyporządkowuje otwartemu $U$ \emph{lokalnie stałe} funkcje $U\to k$;
nie należy go mylić z przedsnopem $U\mapsto k$ nawet wtedy,
gdy $U$ jest niespójne. Każdy system lokalny o włóknie $k^r$
można otrzymać z reprezentacji
$\pi_1(X,x_0)\to\operatorname{GL}_r(k)$,
jeśli $X$ jest spójna, lokalnie łukowo spójna i półlokalnie jednospójna.
\end{przyklad}

\section{Nośnik sekcji i przedłużenie przez zero}
\begin{definicja}[Sekcja o nośniku zwartym]
\label{def:gamma-c-v12}
Niech $X$ będzie lokalnie zwarta i Hausdorffa, a $\mathcal F$
niech będzie snopem na $X$. Dla $s\in\mathcal F(X)$ określmy
$\operatorname{supp}(s)=\{x\in X:s_x\ne0\}$
(zbiór ten jest już domknięty na mocy warunku snopowego).
Przestrzeń sekcji o nośniku zwartym oznaczamy
\[\Gamma_c(X,\mathcal F)
  :=\{s\in\mathcal F(X):\operatorname{supp}(s)\text{ jest zwarty}\}.\]
Dla domkniętego $Z\subseteq X$ analogicznie oznaczamy
$\Gamma_Z(X,\mathcal F)=\{s:\operatorname{supp}(s)\subseteq Z\}$.
\end{definicja}

\begin{uwaga}[Skąd domkniętość nośnika?]
Zbiór punktów $x$ o $s_x=0$ jest otwarty: zerowy kiełek oznacza,
że na pewnym otoczeniu $x$ sekcja $s$ jest zerowa. Jego dopełnienie
jest więc domknięte. Dla snopu funkcji jest to zwykłe domknięcie
zbioru, na którym funkcja jest różna od zera.
\end{uwaga}

\begin{definicja}[Przedłużenie przez zero]
\label{def:jbang-v12}
Niech $j:U\hookrightarrow X$ będzie zanurzeniem \emph{otwartym}
i $\mathcal F$ snopem na $U$. Przez $j_!\mathcal F$ oznaczamy snop,
który ma na otwartym $V\subseteq X$ następujące sekcje:
\[(j_!\mathcal F)(V)=\{s\in\mathcal F(V\cap U):
\operatorname{supp}_{V\cap U}(s)\text{ jest domknięty w }V\}.\]
W szczególności sekcje te dają się przedłużyć przez zero
poza $U$. Na włóknach:
\[
(j_!\mathcal F)_x\cong
\begin{cases}\mathcal F_x,&x\in U,\\0,&x\notin U.\end{cases}
\]
Natomiast obraz bezpośredni $j_*\mathcal F$ definiuje wzór
$(j_*\mathcal F)(V)=\mathcal F(V\cap U)$ dla otwartych $V\subseteq X$.
\end{definicja}

\begin{przyklad}[Dlaczego $j_!$ różni się od $j_*$?]
Niech $X=\R$, $U=(0,1)$ i $\mathcal F=k_U$. Dla $x=0$ mamy
$(j_!k_U)_0=0$, ale $(j_*k_U)_0\cong k$, gdyż dla małego
$\varepsilon>0$ funkcje lokalnie stałe na $(0,\varepsilon)$
mają jedną swobodną wartość. Istotne jest więc, czy
informację dopuszczamy \emph{aż do brzegu}, czy znikającą przy nim.
\end{przyklad}

\begin{center}
\begin{tikzpicture}[x=1cm,y=1cm,font=\small]
\draw[-{Stealth[length=2mm]},thick] (-3.9,0)--(4.0,0);
\node[below] at (-3.4,0) {$X=\R$};
\draw[line width=2.3pt,manifoldblue] (-1.9,0)--(2.0,0);
\draw[fill=white,thick] (-1.9,0) circle (2pt);
\draw[fill=white,thick] (2.0,0) circle (2pt);
\node[above,manifoldblue] at (0,0.12) {$U=(0,1)$ (schematycznie)};
\node[below] at (-1.9,-0.14) {$0$};
\node[below] at (2.0,-0.14) {$1$};
\node[below,align=center] at (0,-0.70) {Na brzegu: $(j_!k_U)_x=0$,\\ ale $(j_*k_U)_x\cong k$.};
\end{tikzpicture}
\par\smallskip\small Przedłużenie przez zero zachowuje dane wewnątrz $U$,
ale usuwa ich kiełki na brzegu otwartego podzbioru.
\end{center}

\begin{stwierdzenie}[Sekcje na otwartym podzbiorze]
\label{prop:jbang-sections-v12}
Przy założeniach definicji~\ref{def:jbang-v12}, jeśli $X$ jest lokalnie
zwarta i Hausdorffa, to zachodzi naturalna identyfikacja
$\Gamma_c(X,j_!\mathcal F)\cong\Gamma_c(U,\mathcal F)$.
\end{stwierdzenie}
\begin{proof}
Sekcja $j_!\mathcal F$ jest na $X\setminus U$ lokalnie zerowa,
a więc odpowiada sekcji $\mathcal F$ na $U$ z zerowym przedłużeniem.
Jeśli jej nośnik jest zwarty w $X$, to leży w $U$ i jest
zwarty również w $U$; w drugą stronę nośnik zwarty w $U$
jest zwarty w $X$ (z ciągłości włożenia), a przedłużenie
przez zero określa sekcję $j_!\mathcal F$.
Konstrukcje są odwrotne i naturalne.
\end{proof}

\section{Kohomologia snopów i kohomologia z nośnikami}
\begin{definicja}[Funktory pochodne sekcji]
\label{def:sheaf-cohom-v12}
Dla snopu $\mathcal F$ na $X$ wybieramy rezolwentę iniektywną
$0\to\mathcal F\to\mathcal I^0\to\mathcal I^1\to\cdots$
w kategorii snopów $k$-przestrzeni wektorowych. Definiujemy
\[
H^q(X,\mathcal F)=H^q\bigl(\Gamma(X,\mathcal I^\bullet)\bigr),
\quad H_c^q(X,\mathcal F)
=H^q\bigl(\Gamma_c(X,\mathcal I^\bullet)\bigr),
\quad H_Z^q(X,\mathcal F)
=H^q\bigl(\Gamma_Z(X,\mathcal I^\bullet)\bigr).
\]
Są to odpowiednio kohomologia snopu, kohomologia z nośnikami
zwartymi oraz kohomologia lokalna z nośnikiem w $Z$.
Ich niezależność od wyboru rezolwenty wynika z twierdzenia
porównawczego dla rezolwent iniektywnych.
\end{definicja}

\begin{uwaga}[Co zakładamy z algebry homologicznej?]
Istnienie dostatecznej liczby snopów iniektywnych i twierdzenie
o niezależności rezolwenty są standardowymi rezultatami
algebry homologicznej w kategorii snopów. Konstrukcja powyższa
nie \emph{dowodzi} ich istnienia; szczegóły kategorii
pochodnej rozwiniemy przed dualnością Verdiera.
\end{uwaga}

\begin{lemat}[Ciąg Mayera--Vietorisa ze zwartymi nośnikami]
\label{lem:sheaf-mv-v12}
Niech $X$ będzie lokalnie zwartą przestrzenią Hausdorffa,
$X=U\cup V$ pokryciem otwartym, a $\mathcal F$ snopem
$k$-przestrzeni wektorowych na $X$. Zachodzi naturalny
długi ciąg dokładny
\[
\cdots\longrightarrow H_c^q(U\cap V,\mathcal F)
\longrightarrow H_c^q(U,\mathcal F)\oplus H_c^q(V,\mathcal F)
\longrightarrow H_c^q(X,\mathcal F)
\longrightarrow H_c^{q+1}(U\cap V,\mathcal F)\longrightarrow\cdots,
\]
gdzie wszędzie po lewej stosujemy ograniczenie snopu do
wskazanego otwartego podzbioru.
\end{lemat}
\begin{proof}
Niech $a:U\cap V\hookrightarrow X$, $j:U\hookrightarrow X$
i $l:V\hookrightarrow X$ oznaczają zanurzenia otwarte.
Sprawdzenie na włóknach (Lemat~\ref{lem:stalk-exact-v12}) daje krótki
ciąg dokładny snopów
\[
0\longrightarrow a_!(\mathcal F|_{U\cap V})
\xrightarrow{\,s\mapsto(s,-s)\,}
 j_!(\mathcal F|_U)\oplus l_!(\mathcal F|_V)
\xrightarrow{\,(s,t)\mapsto s+t\,}\mathcal F\longrightarrow0.
\]
Na przecięciu jest to $0\to F\to F\oplus F\to F\to0$,
a poza przecięciem identyczność na odpowiednim składniku.
Ponieważ przedłużenie przez zero przy otwartym zanurzeniu
jest dokładne (co także wynika ze sprawdzenia na włóknach),
ciąg pozostaje dokładny po zastosowaniu rezolwent w kategorii
pochodnej. Następnie stosujemy funktor $R\Gamma_c(X,-)$.
Przez Stwierdzenie~\ref{prop:jbang-sections-v12} i dokładność
$j_!$ identyfikujemy otrzymane człony z
$R\Gamma_c(U\cap V,-)$, $R\Gamma_c(U,-)$ i $R\Gamma_c(V,-)$.
Długi ciąg kohomologii trójkąta dokładnego daje tezę.
\end{proof}

\begin{przyklad}[Obliczenie $H_c^*(\R;k)$]
Niech $U=(-\infty,1)$, $V=(0,\infty)$, więc $U\cap V=(0,1)$
i $\R=U\cup V$. Każdy z tych trzech zbiorów jest homeomorficzny
z $\R$, dlatego w tym przykładzie ciąg Mayera--Vietorisa
sam nie oblicza nieznanej grupy. Zamiast kołowego rozumowania
używamy definicji względnej z
Definicji~\ref{def:compact-support-cohom-v11}:
$H_c^q(\R;k)\cong H^q([-1,1],\{-1,1\};k)$
przez wycięcie i wyczerpanie przedziałami. Ciąg względny daje
$H_c^0(\R;k)=0$ i $H_c^1(\R;k)\cong k$.
Poniżej odzyskamy ten sam wynik z rezolwenty de Rhama.
\end{przyklad}

\section{Od kohomologii \v{C}echa do kohomologii snopów}
\label{sec:cech-sheaves-v12}
Nie należy mylić \emph{homologii \v{C}echa} z rozdziału o
nerwach pokryć z \emph{kohomologią \v{C}echa snopu}, której
współczynniki stanowią sekcje snopu na przecięciach otwartych.

\begin{definicja}[Kołańcuchy \v{C}echa ze współczynnikami w snopie]
Niech $\mathcal U=(U_i)_{i\in I}$ będzie uporządkowanym pokryciem
otwartym $X$, a $\mathcal F$ snopem abelowym. Definiujemy
\[
\check C^q(\mathcal U,\mathcal F)
=\prod_{i_0<\cdots<i_q}\mathcal F(U_{i_0}\cap\cdots\cap U_{i_q}),
\qquad
(\delta c)_{i_0\cdots i_{q+1}}
=\sum_{j=0}^{q+1}(-1)^j
c_{i_0\cdots\widehat{i_j}\cdots i_{q+1}}
\big|_{U_{i_0}\cap\cdots\cap U_{i_{q+1}}}.
\]
Niepuste przecięcia są przyjmowane w indeksie; dla pustego
przecięcia przestrzeń sekcji jest zerowa. Kohomologię tego
kompleksu oznaczamy $\check H^q(\mathcal U,\mathcal F)$.
\end{definicja}

\begin{stwierdzenie}[Pierwsza przeszkoda sklejania]
$\check H^0(\mathcal U,\mathcal F)=\mathcal F(X)$,
gdy $\mathcal F$ jest snopem i $\mathcal U$ pokrywa $X$.
\end{stwierdzenie}
\begin{proof}
Zero-kocykl to rodzina sekcji $(s_i)$ spełniająca
$s_j|_{U_i\cap U_j}=s_i|_{U_i\cap U_j}$.
Warunek snopowy daje dokładnie jedną sekcję na całym $X$.
W standardowym kompleksie \v{C}echa nie ma kołańcuchów stopnia $-1$,
zatem w stopniu zerowym nie bierzemy ilorazu przez żadne kobrzegi.
\end{proof}

\begin{twierdzenie}[Kiedy \v{C}ech oblicza kohomologię snopów?]
\label{thm:leray-sheaf-v12}
Jeżeli dla każdego niepustego skończonego przecięcia
$U_{i_0}\cap\cdots\cap U_{i_q}$ oraz każdego $r>0$ mamy
$H^r(U_{i_0}\cap\cdots\cap U_{i_q},\mathcal F)=0$,
to naturalna mapa
$\check H^m(\mathcal U,\mathcal F)\to H^m(X,\mathcal F)$
jest izomorfizmem dla każdego $m\ge0$.
\end{twierdzenie}
\begin{proof}[Podwójny kompleks i dwie filtracje]
Niech $\mathcal F\to\mathcal I^\bullet$ będzie rezolwentą
iniektywną. Dla uporządkowanego pokrycia $\mathcal U$
wprowadzamy dwukompleks
$C^{p,q}=\prod_{i_0<\cdots<i_p}
\Gamma(U_{i_0\cdots i_p},\mathcal I^q)$ dla $p,q\ge0$,
z różniczką całkowitą
$D=\delta+(-1)^p d_{\mathcal I}$ na $C^{p,q}$.
Obie filtracje są w pierwszej ćwiartce, zatem w każdym
stopniu całkowitym występuje skończenie wiele par $(p,q)$;
standardowe twierdzenie o zbieżności ma tu zastosowanie.

Najpierw liczymy kohomologię w kierunku $p$. Snop
iniektywny jest wiotki (flabby): ograniczenie jego sekcji
do zbioru otwartego jest surjektywne. Dla wiotkiego snopu
rozszerzony kompleks \v{C}echa danego pokrycia ma
zanik kohomologii w dodatnich stopniach. Istotne jest,
że chodzi tutaj o \emph{wiotkość} snopów $\mathcal I^q$,
a nie o sam warunek snopowy, który zapewnia tylko
dokładność w stopniu zerowym. Wobec tego pierwsza
filtracja daje kohomologię kompleksu
$\Gamma(X,\mathcal I^\bullet)$, czyli $H^*(X,\mathcal F)$.

W drugim kierunku, ponieważ iniektywne snopy są
acykliczne na dowolnym otwartym podzbiorze,
$H^q(\Gamma(U_{i_0\cdots i_p},\mathcal I^\bullet))
=H^q(U_{i_0\cdots i_p},\mathcal F)$.
Iloczyny przestrzeni wektorowych są dokładne,
więc $E_1^{p,q}=\prod_{i_0<\cdots<i_p}
H^q(U_{i_0\cdots i_p},\mathcal F)$.
Z hipotezy wynika zanik dla $q>0$.
Wiersz $q=0$ jest dokładnie kompleksem
$\check C^\bullet(\mathcal U,\mathcal F)$.
Dlatego drugi ciąg spektralny degeneruje po obliczeniu
kohomologii tego wiersza. Obie filtracje obliczają
kohomologię tego samego totalnego kompleksu;
indukowana mapa porównawcza jest zatem izomorfizmem.
Użyliśmy przy tym standardowego lematu o wiotkości
snopów iniektywnych, wyraźnie oddzielonego od rachunku
z dwiema filtracjami.
\end{proof}

\begin{przyklad}[Okrąg i jego lokalny system]
Niech $L$ będzie systemem lokalnym o włóknie $V=k^r$
na $S^1$, o monodromii $A\in\operatorname{GL}(V)$.
Jednokomórkowy model okręgu z odpowiednim skręceniem
współczynników ma kompleks kołańcuchów
$V\xrightarrow{A-I}V$ w stopniach $0$ i $1$.
Dlatego
$H^0(S^1;L)=\ker(A-I)$ oraz
$H^1(S^1;L)=\operatorname{coker}(A-I)$.
Na przykład dla $V=\R$ i $A=-1$ obie grupy są zerowe;
przy $A=1$ obie wynoszą $\R$.
To pokazuje, dlaczego znajomość samych włókien $L_x$
nie wystarcza: trzeba również znać ich transport.
\end{przyklad}

\section{Rezolwenta de Rhama i związki z poprzednim rozdziałem}
\begin{twierdzenie}[Twierdzenie de Rhama ze zwartymi nośnikami — dowód snopowy]
\label{thm:de-rham-sheaf-v12}
Niech $M$ będzie gładką drugoprzeliczalną rozmaitością,
bez brzegu. Wówczas naturalne całkowanie form wyznacza
$H^q(\Omega_c^\bullet(M),\dd)\cong H_c^q(M;\R)$,
gdzie prawa strona jest singularną kohomologią z nośnikami
zwartymi z Definicji~\ref{def:compact-support-cohom-v11}.
\end{twierdzenie}
\begin{proof}
Rozpatrzmy ciąg snopów
$0\to\R_M\to\Omega_M^0\xrightarrow{\dd}\Omega_M^1
\xrightarrow{\dd}\cdots$. Jest on dokładny na włóknach:
lokalny lemat Poincarégo dowodzi dokładności w dodatnich
stopniach, a w stopniu zero jądrem $\dd$ są funkcje lokalnie
stałe. Dla każdej formy i każdego otwartego pokrycia
istnieje gładki podział jedności podporządkowany
lokalnie skończonemu uszczegółowieniu; mnożenie form
przez jego funkcje daje operatory lokalizujące.
Dlatego snopy $\Omega_M^q$ są \emph{drobne} (fine).

Dla przestrzeni lokalnie zwartych i parazwartych snopy
drobne są acykliczne względem $\Gamma_c$. Wyjaśnijmy
mechanizm tego standardowego lematu: w kompleksie
\v{C}echa sekcji o nośnikach zwartych podział jedności
pozwala zbudować homotopię kontrahującą dodatnie
stopnie. Każda sekcja o nośniku zwartym spotyka tylko
skończenie wiele elementów lokalnie skończonego pokrycia,
więc operatory zachowują warunek zwartości nośnika;
na przestrzeni parazwartej porównanie rezolwenty
\v{C}echa z rezolwentą iniektywną daje zanik wyższych
funktorów pochodnych. W efekcie kompleks
$\Gamma_c(M,\Omega_M^\bullet)=\Omega_c^\bullet(M)$
oblicza $R\Gamma_c(M,\R_M)$.

Dla rozmaitości lokalnie kontraktowalnej porównanie
kohomologii stałego snopu z kohomologią singularną
(respektujące nośniki) utożsamia
$H_c^q(M,\R_M)$ z $H_c^q(M;\R)$.
Trzeba doprecyzować, co znaczy tutaj „całkowanie form”
w przypadku nośników zwartych. Dla ustalonego kompaktu
$K\subset M$ całkowanie po gładkich sympleksach daje
kołańcuch reprezentujący klasę względną
$H^q(M,M\setminus K;\R)$: forma o nośniku zawartym
w $K$ zeruje się na sympleksach o obrazie w $M\setminus K$.
Po przejściu do granicy prostej po kompaktach $K$
otrzymujemy mapę do $H_c^q(M;\R)$.
Jej zgodność z izomorfizmem snopowym wynika z lokalnego
lematu Poincarégo i naturalności konstrukcji względem
ograniczeń i ciągów Mayera--Vietorisa.
Stosujemy jako jawne twierdzenie porównawcze
identyfikację pochodnej kohomologii stałego snopu
z singularną na rozmaitościach; samo stwierdzenie,
że formy całkujemy po sympleksach, nie jest niezależnym
dowodem tego twierdzenia.
\end{proof}

\section{Kategorie pochodne: tylko niezbędny aparat}
\begin{definicja}[Kompleks i przesunięcie]
Dla kompleksu kohomologicznego $(A^\bullet,d)$ ustalamy
$A[n]^q=A^{q+n}$ oraz $d_{A[n]}=(-1)^n d_A$.
Morfizm kompleksów jest \emph{quasi-izomorfizmem}, gdy
indukuje izomorfizmy wszystkich grup kohomologii.
Kategoria pochodna $D(X)$ jest otrzymywana z kategorii
kompleksów snopów przez formalne odwrócenie wszystkich
quasi-izomorfizmów.
\end{definicja}

\begin{przyklad}[Odczytywanie przesunięcia]
Jeżeli $V$ jest przestrzenią wektorową skoncentrowaną
w stopniu $0$, to $V[n]$ ma jedyną niezerową kohomologię
w stopniu $-n$. W szczególności $\R_M[n]$ nie jest
zwykłym snopem ustawionym w stopniu $n$, tylko
kompleksem przesuniętym \emph{w lewo}.
\end{przyklad}

\begin{definicja}[Pochodne funktory Hom]
Przez $R\!\operatorname{Hom}_k(A^\bullet,B^\bullet)$
oznaczamy pochodny kompleks Hom, a przez
$R\mathcal Hom(\mathcal F^\bullet,\mathcal G^\bullet)$
pochodny snop wewnętrznych homomorfizmów.
Definiuje się je za pomocą stosownych rezolwent
iniektywnych; ich kohomologie mierzą nie tylko
homomorfizmy, lecz także klasy rozszerzeń.
\end{definicja}

\section{Snop dualizujący i twierdzenie Verdiera}
\begin{definicja}[Snop dualizujący i dualność Verdiera]
\label{def:verdier-v12}
Niech $a:X\to\{*\}$ będzie mapą do punktu.
W formalizmie sześciu operacji dla odpowiedniej kategorii
przestrzeni (np. lokalnie zwartych, lokalnie kontraktowalnych,
o skończonym wymiarze kohomologicznym), $a^!$ oznacza
prawy adjunkt do pochodnego obrazu właściwego $Ra_!$.
Kompleks
$\omega_X^\bullet:=a^!k$ nazywamy \emph{kompleksem dualizującym}.
Dla kompleksu $\mathcal F^\bullet$ definiujemy
\[
\mathbb D_X\mathcal F^\bullet
:=R\mathcal Hom(\mathcal F^\bullet,\omega_X^\bullet).
\]
Przez $D^b_c(X;k)$ oznaczamy kategorię ograniczonych
kompleksów snopów o konstruktybilnych grupach kohomologii,
gdy $X$ jest odpowiednio stratyfikowalna.
\end{definicja}

\begin{twierdzenie}[Globalna dualność Verdiera — wersja nad ciałem]
\label{thm:verdier-global-v12}
Niech $X$ będzie lokalnie zwartą, lokalnie kontraktowalną,
parazwartą przestrzenią o skończonym wymiarze kohomologicznym,
na której istnieje standardowy formalizm $Ra_!\dashv a^!$.
Dla $\mathcal F^\bullet\in D^b_c(X;k)$ zachodzi
naturalny quasi-izomorfizm
\[
R\Gamma\bigl(X,\mathbb D_X\mathcal F^\bullet\bigr)
\simeq R\!\operatorname{Hom}_k
\bigl(R\Gamma_c(X,\mathcal F^\bullet),k\bigr).
\]
W szczególności, gdy $R\Gamma_c(X,\mathcal F^\bullet)$
ma skończeniewymiarowe grupy kohomologii,
$H^{-q}(X,\mathbb D_X\mathcal F^\bullet)
\cong H_c^q(X,\mathcal F^\bullet)^\vee$,
gdzie $V^\vee=\operatorname{Hom}_k(V,k)$.
\end{twierdzenie}
\begin{proof}[Wskazanie dokładnego miejsca użycia adjunkcji]
Z definicji kompleksu dualizującego i z adjunkcji
$Ra_!\dashv a^!$ wynika
$R\!\operatorname{Hom}_X(\mathcal F^\bullet,a^!k)
\simeq R\!\operatorname{Hom}_k(Ra_!\mathcal F^\bullet,k)$.
Lewą stronę identyfikujemy z
$R\Gamma(X,R\mathcal Hom(\mathcal F^\bullet,a^!k))$,
a $Ra_!\mathcal F^\bullet$ z
$R\Gamma_c(X,\mathcal F^\bullet)$.
Po podstawieniu dostajemy zapisaną równoważność.
Dowód istnienia $a^!$ i konstrukcji adjunkcji jest
osobnym, nietrywialnym twierdzeniem teorii pochodnych
snopów, którego tutaj nie wyprowadzamy.
\end{proof}

\begin{twierdzenie}[Postać dualizująca na rozmaitości]
\label{thm:verdier-manifold-v12}
Jeżeli $M$ jest topologiczną $n$-rozmaitością bez brzegu,
orientowalną nad $k$, to
$\omega_M^\bullet\simeq k_M[n]$; bez wyboru orientacji
$\omega_M^\bullet\simeq\operatorname{or}_{M,k}[n]$,
gdzie $\operatorname{or}_{M,k}$ jest lokalnym systemem orientacji.
Dlatego dla $\mathcal F=k_M$ globalna dualność Verdiera
daje dualność Poincarégo z nośnikami zwartymi.
\end{twierdzenie}
\begin{proof}[Wyprowadzenie z lokalnego modelu]
Ustalmy punkt $x\in M$ i kulę współrzędnych $B$.
Lokalna kohomologia z nośnikiem w $x$ spełnia
$H^q(B,B\setminus\{x\};k)\cong k$ dla $q=n$
i znika w innych stopniach. Orientacja lokalna to
wybór generatora tej jednowymiarowej grupy; generatory
w sąsiednich kartach różnią się znakiem stopnia
lokalnej zmiany współrzędnych (albo działaniem $-1$
w ciele $k$). Konstruuje to system $\operatorname{or}_{M,k}$.

Sam rachunek lokalnej kohomologii \emph{nie wystarcza}
do zbudowania globalnego quasi-izomorfizmu. Potrzebujemy
ponadto twierdzenia o lokalizacji dla $a^!$:
po ograniczeniu do $B$ kompleks dualizujący jest
kompleksem dualizującym na $B$, a jego funkcjonalne
odwzorowania przejścia realizują właśnie otrzymany
system orientacji. Z taką identyfikacją lokalną i
zgodnością na przecięciach sklejamy quasi-izomorfizm
$\omega_M^\bullet\simeq\operatorname{or}_{M,k}[n]$.
Twierdzenie o lokalizacji $a^!$ jest tu jawnie
przyjętym składnikiem formalizmu sześciu operacji.
\end{proof}

\begin{przyklad}[Prosta rzeczywista]
Dla $X=\R$ i $\mathcal F=k_\R$ mamy
$\mathbb D_\R(k_\R)\simeq k_\R[1]$.
Stąd $H^{-1}(\R,\mathbb D k_\R)\cong H^0(\R,k)=k$,
co zgadza się z dualną przestrzenią
$H_c^1(\R;k)^\vee\cong k$.
Dla $q=0$ obie strony są zerowe.
\end{przyklad}

\begin{przyklad}[Węzeł krzywej zespolonej]
\label{ex:node-duality-v12}
Niech $X=\{(z,w)\in\C^2:zw=0\}\cap B_\varepsilon$,
gdzie $B_\varepsilon$ jest małą otwartą kulą wokół $0$.
Otoczenie węzła jest topologicznie stożkiem na
dwóch rozłącznych okręgach (geometryczny link w $S^3$
jest splotem Hopfa), a $X\setminus\{0\}$ deformacyjnie zwija się
do $S^1\sqcup S^1$. Ponieważ $X$ jest kontraktowalna,
ciąg kohomologii pary daje
\[
H^q(X,X\setminus\{0\};k)\cong
\widetilde H^{q-1}(S^1\sqcup S^1;k)
=\begin{cases}
k,&q=1,\\k^2,&q=2,\\0,&q\notin\{1,2\}.
\end{cases}
\]
Natomiast w punkcie gładkim zespolonej krzywej
lokalna kohomologia pary jest niezerowa tylko w stopniu $2$,
gdzie wynosi $k$. Dlatego \emph{stały kompleks przesunięty}
$k_X[2]$ nie może odgrywać roli kompleksu dualizującego
w takiej samej postaci jak na rozmaitości gładkiej.
To pierwszy sygnał pojawienia się w geometrii
algebraicznej kompleksów dualizujących i snopów perwersyjnych.
\end{przyklad}

\section{Sześć operacji i właściwe odwzorowania}
Na poziomie pochodnym porządkujemy operacje
$f^{-1}$ (odwrotny obraz snopu), $Rf_*$ (obraz bezpośredni),
$Rf_!$ (obraz właściwy), $f^!$ (nadzwyczajny odwrotny obraz),
$\otimes^L$ i $R\mathcal Hom$.
Dla snopów \emph{wektorowych nad ciałem} zwykłe $f^{-1}$
i $\otimes$ są już dokładne, dlatego w tej kategorii $f^{-1}$ nie wymaga
pochodnego poprawiania, a $\otimes^L=\otimes$.

\begin{stwierdzenie}[Obraz właściwy i dualność]
\label{prop:verdier-proper-v12}
Jeśli $f:X\to Y$ jest właściwym ciągłym odwzorowaniem
między przestrzeniami spełniającymi hipotezy
Twierdzenia~\ref{thm:verdier-global-v12}, to
$Rf_!\cong Rf_*$. Dla konstruktybilnego
$\mathcal F^\bullet$ zachodzi ponadto naturalna
równoważność
\[
Rf_*\,\mathbb D_X\mathcal F^\bullet
\simeq\mathbb D_Y\,Rf_*\mathcal F^\bullet.
\]
\end{stwierdzenie}
\begin{proof}[Argument z formalizmu sześciu operacji]
Właściwość mapy oznacza, że warunek właściwości
nośnika sekcji na przeciwobrazie dowolnego otwartego
zbioru jest automatyczny; więc $f_!=f_*$ i tak samo
po przejściu do funktorów pochodnych. Zgodność z dualnością
otrzymujemy z twierdzenia o dualności względnej
$Rf_*R\mathcal Hom_X(\mathcal F,f^!\mathcal G)
\simeq R\mathcal Hom_Y(Rf_*\mathcal F,\mathcal G)$
razem z $(a_Y\circ f)^!\simeq f^!a_Y^!$.
Jest to zastosowanie twierdzenia o dualności względnej,
a nie jego niezależny dowód.
\end{proof}

\begin{uwaga}[Co wchodzi do geometrii algebraicznej?]
Dla przestrzeni analitycznej $X^{\mathrm{an}}$ mamy snop
stały $k_X$ i jego kohomologię topologiczną. Mamy też
snop funkcji holomorficznych $\mathcal O_X$ oraz dla
schematów snopy koherentne w topologii Zariskiego.
Nie należy bez hipotez identyfikować
$H^*(X^{\mathrm{an}},k)$, $H^*(X,\mathcal O_X)$
i kohomologii etalnej. Należą do różnych teorii
współczynników i, w tym ostatnim przypadku,
topologii Grothendiecka. Dualność Verdiera prowadzi
od lokalnych systemów do kompleksów konstruktybilnych,
a dalej do kompleksów przecięcia i twierdzenia o rozkładzie.
Tych zaawansowanych wyników nie uznajemy jeszcze za udowodnione.
\end{uwaga}

\section{Ciąg lokalizacji dla podzbioru zamkniętego}
\begin{twierdzenie}[Lokalizacja nośników]
\label{thm:localization-v12}
Niech $X$ będzie lokalnie zwartą przestrzenią Hausdorffa,
$Z\subseteq X$ zbiorem domkniętym, $U=X\setminus Z$,
a $j:U\hookrightarrow X$, $i:Z\hookrightarrow X$
kanonicznymi włożeniami. Dla snopu $\mathcal F$
$k$-przestrzeni wektorowych na $X$ istnieje naturalny
ciąg dokładny
\[
\cdots\to H_c^q(U,\mathcal F|_U)\to H_c^q(X,\mathcal F)
\to H_c^q(Z,\mathcal F|_Z)
\to H_c^{q+1}(U,\mathcal F|_U)\to\cdots.
\]
\end{twierdzenie}
\begin{proof}
Rozpatrzmy ciąg snopów na $X$:
\[
0\longrightarrow j_!(\mathcal F|_U)
\longrightarrow\mathcal F\longrightarrow i_*(\mathcal F|_Z)
\longrightarrow0.
\]
Sprawdźmy dokładność, nie ograniczając się do sekcji globalnych.
Dla $x\in U$ na włóknach otrzymujemy
$0\to\mathcal F_x\xrightarrow{\id}\mathcal F_x\to0$.
Dla $x\in Z$ otrzymujemy
$0\to0\to\mathcal F_x\xrightarrow{\id}\mathcal F_x\to0$,
ponieważ ograniczenie snopu do $Z$ nie zmienia
kiełków przy punktach $Z$. Zatem ciąg jest dokładny
na mocy Lematu~\ref{lem:stalk-exact-v12}.

Pochodny funktor sekcji o nośnikach zwartych
przekształca ten ciąg w długi ciąg kohomologii.
Pierwszy człon identyfikujemy ze wspomnianą
kohomologią $U$ przez dokładność $j_!$.
Ponieważ $i$ jest zanurzeniem domkniętym i właściwym,
$\Gamma_c(X,i_*\mathcal G)=\Gamma_c(Z,\mathcal G)$;
identyfikacja pozostaje prawdziwa po zastosowaniu
funktorów pochodnych. Otrzymujemy wzór z tezy.
\end{proof}

\begin{przyklad}[Usunięcie punktu z prostej]
Dla $X=\R$, $Z=\{0\}$ i $U=\R\setminus\{0\}$ mamy
$H_c^0(Z;k)=k$, $H_c^1(U;k)=k^2$ oraz
$H_c^1(\R;k)=k$. Istotny fragment ciągu lokalizacji to
\[
0\longrightarrow k\xrightarrow{\Delta}k^2
\xrightarrow{\,(a,b)\mapsto b-a\,}k\longrightarrow0,
\]
przy odpowiednim wyborze generatorów orientacji
obu półprostych. Łączący homomorfizm $\Delta$
jest przekątny: usunięcie jednego punktu tworzy dwie
niezależne klasy o nośnikach zwartych w stopniu pierwszym.
\end{przyklad}

\section{Przykład z geometrii algebraicznej: kohomologia \texorpdfstring{$\mathcal O(m)$}{O(m)} na \texorpdfstring{$\mathbb P^1$}{P1}}
\label{sec:p1-sheaf-v12}
Nie poprzestawajmy na formalizmie: snopy pozwalają
wykonać rachunek, którego nie opisuje zwykła homologia
singularna. Niech $\mathbb P^1=\mathbb P^1(\C)$,
$U_0=\{[x:y]:x\ne0\}$, $U_\infty=\{[x:y]:y\ne0\}$.
Oznaczmy $z=y/x$ na $U_0$, $w=x/y$ na $U_\infty$.
Na przecięciu $w=1/z$.

\begin{definicja}[Snop $\mathcal O(m)$]
Dla $m\in\Z$ wiązkę liniową $\mathcal O(m)$
trywializujemy na $U_0$ i $U_\infty$ przez
lokalne ramy $e_0$ i $e_\infty$ z regułą przejścia
$e_\infty=z^m e_0$ na $U_0\cap U_\infty$.
Snop $\mathcal O(m)$ oznacza snop \emph{holomorficznych}
przekrojów tej wiązki.
\end{definicja}

\begin{twierdzenie}[Kohomologia skręceń na prostej rzutowej]
\label{thm:p1-twist-v12}
Dla każdego $m\in\Z$ zachodzą izomorfizmy
\[
H^0(\mathbb P^1,\mathcal O(m))\cong
\begin{cases}\C^{m+1},&m\ge0,\\0,&m<0,\end{cases}
\qquad
H^1(\mathbb P^1,\mathcal O(m))\cong
\begin{cases}\C^{-m-1},&m\le-2,\\0,&m\ge-1.\end{cases}
\]
Ponadto $H^q(\mathbb P^1,\mathcal O(m))=0$ dla $q\ge2$.
\end{twierdzenie}
\begin{proof}
Zaczniemy od kompleksu \v{C}echa dla pokrycia
$\{U_0,U_\infty\}$, identyfikując przekroje z funkcjami
holomorficznymi w odpowiednich trywializacjach.
Zgodność przekrojów na przecięciu ma postać
$f_0(z)=z^m f_\infty(1/z)$.
Zatem kompleks ma dwa niezerowe wyrazy:
\[
0\longrightarrow
\mathcal O(\C)\oplus\mathcal O(\C)
\xrightarrow{\,\delta\,}\mathcal O(\C^*)
\longrightarrow0,
\qquad
\delta(f_0,f_\infty)=f_0(z)-z^m f_\infty(1/z).
\]
Jego jądro liczy globalne przekroje: funkcja
$f_0\in\mathcal O(\C)$ może przedłużyć się zgodnie
z ramą w nieskończoności wtedy i tylko wtedy,
gdy jest wielomianem stopnia co najwyżej $m$ przy $m\ge0$;
przy $m<0$ pozostaje jedynie przekrój zerowy.

Obliczmy teraz kokernele. Każda funkcja holomorficzna
na $\C^*$ ma rozwinięcie Laurenta
$h(z)=\sum_{j\in\Z}a_jz^j$, zbieżne lokalnie jednostajnie
na $\C^*$. Część $\sum_{j\ge0}a_jz^j$ jest funkcją
całkowitą na $\C$ i może zostać usunięta przez
odjęcie $f_0$. Część $\sum_{j\le m}a_jz^j$
ma postać $z^m f_\infty(1/z)$, gdzie
$f_\infty(w)=\sum_{j\le m}a_jw^{m-j}$ jest
funkcją całkowitą w $w$; można więc usunąć również
wszystkie potęgi $j\le m$. Gdy $m\ge-1$, te dwa
zbiory wykładników pokrywają całą $\Z$.
Gdy $m\le-2$, pozostają dokładnie klasy
$z^{m+1},z^{m+2},\ldots,z^{-1}$, których liczba wynosi
$-m-1$. Stąd obliczenie $H^1$ kompleksu \v{C}echa.

Aby utożsamić je z \emph{pochodną} kohomologią snopu,
stosujemy Twierdzenie~\ref{thm:leray-sheaf-v12}:
$U_0\cong\C$, $U_\infty\cong\C$ i
$U_0\cap U_\infty\cong\C^*$ są przestrzeniami Steina,
a kohomologie wyższych stopni snopów holomorficznych
przekrojów wiązek wektorowych znikają na nich
na mocy \emph{twierdzenia B Cartana}. To ostatnie
jest tutaj wyraźnie przywołanym wynikiem analizy
zespolonej, a nie ukrytą konsekwencją samej definicji
snopu. Pokrycie jest zatem acykliczne i rachunek
\v{C}echa daje wszystkie grupy kohomologii.
\end{proof}

\begin{uwaga}[Dwa różne rodzaje kohomologii]
Na $\mathbb P^1$ singularna kohomologia ze stałymi
współczynnikami ma $H^0\cong H^2\cong\Z$ i
$H^1=0$, niezależnie od $m$. Natomiast
$H^1(\mathbb P^1,\mathcal O(-3))\cong\C^2$.
To nie jest sprzeczność: grupy badają ten sam
zbiór topologiczny, ale \emph{inne współczynniki snopowe}.
\end{uwaga}


\section{Podsumowanie, ograniczenia i dalszy kierunek}
W tym rozdziale udowodniliśmy bez aparatu pochodnego
warunek snopowy, sprawdzanie dokładności na włóknach,
obliczenie $H^0$ metodą \v{C}echa, opis przedłużenia przez zero
i elementarne przykłady. Dla kohomologii z nośnikami
wyprowadziliśmy ciąg Mayera--Vietorisa z dokładnego
ciągu snopów. Rezolwenta de Rhama tłumaczy,
dlaczego lokalny lemat Poincarégo daje globalne
kohomologie z nośnikami zwartymi.

Dualność Verdiera jest tu sformułowana precyzyjnie
\emph{względem} istnienia funktora $a^!$
i standardowego formalizmu sześciu operacji.
Nie rozwinęliśmy niezależnego dowodu jego istnienia,
pełnej teorii konstruktybilności ani
pełnego dowodu względnej dualności. Takie rozdzielenie
statusu dowodów będzie ważne przy scalaniu skryptu
z teorią snopów i geometrią algebraiczną.


\chapter{Homologia przecięć, snopy perwersyjne i twierdzenie o rozkładzie}
\label{ch:ih-perverse-xiii}

\section{Dlaczego zwykła homologia nie wystarcza dla osobliwości?}
W rozdziale~\ref{ch:snopy-verdier-v12} otrzymaliśmy dualność Verdiera,
ale dla przestrzeni osobliwej kompleks dualizujący nie musi być
przesunięciem snopu stałego. Przykładowo dla węzła $zw=0$
lokalne grupy $H^q(X,X\setminus\{0\};k)$ są niezerowe w dwóch
stopniach, a w punkcie gładkim krzywej jedynie w stopniu $2$.
Podobna trudność pojawia się w dualności Poincarégo: zwykłe grupy
homologii zwartej przestrzeni osobliwej nie muszą układać się w
niedegenerującą parę stopni dopełniających się do wymiaru.

Ideą Goresky'ego--MacPhersona jest ograniczenie łańcuchów, które
mogą przechodzić przez warstwy osobliwe. Liczbę dopuszczalnych
wymiarów przecięcia reguluje funkcja zwana
\emph{perwersyjnością} (ang. \emph{perversity}).
Współczesna konstrukcja snopowa opisuje tę samą teorię przez
\emph{kompleks przecięcia} $IC_X$, który jest szczególnym
przykładem snopu perwersyjnego. Umożliwia to sformułowanie
ważnego twierdzenia o rozkładzie obrazów właściwych.

W części łańcuchowej pracujemy nad ustalonym ciałem $k$,
zazwyczaj $\Q$. W zastosowaniach algebraicznych
przyjmujemy $k=\Q$ i topologię klasyczną zespolonych
rozmaitości algebraicznych. Tam, gdzie twierdzenie wymaga
zaawansowanej teorii pochodnej lub Hodge'a, wskazujemy je
jawnie jako wynik przyjęty, zamiast zastępować dowód samą
formułą.

\section{Pseudorozmaitości warstwowe i funkcje perwersyjności}
\begin{definicja}[Warstwowa pseudorozmaitość PL]
\label{def:pseudomanifold-xiii}
Niech $n\ge2$. Przez $n$-wymiarową pseudorozmaitość
warstwową typu PL rozumiemy parazwartą przestrzeń $X$
z filtracją domkniętymi podzbiorami
$X=X_n\supseteq X_{n-1}=X_{n-2}\supseteq\cdots\supseteq X_0
\supseteq X_{-1}=\varnothing$, zgodną z lokalnymi
triangulacjami, taką że każda różnica
$S_j=X_j\setminus X_{j-1}$ jest (być może pustą)
$j$-rozmaitością PL, a każdy punkt $x\in S_j$ dla $j<n$ posiada
warstwowe otoczenie homeomorficzne z
$\R^j\times cL_x$; tutaj $L_x$ jest zwartą
$(n-j-1)$-wymiarową przestrzenią warstwową,
a $cL_x=(L_x\times[0,1))/(L_x\times\{0\})$
ma naturalną filtrację stożkową. W punktach warstwy
głównej $S_n$ wymagamy otoczeń homeomorficznych
z $\R^n$; nie stosujemy tam wzoru na link
wymiaru $-1$. Dla $j<n$ wymiar linku może wynosić
co najmniej $1$ wobec braku warstw kowymiaru $1$. Zakładamy gęstość
$X\setminus X_{n-2}$ w $X$ i brak warstw kowymiaru $1$.
Przestrzeń nazywamy \emph{normalną}, jeżeli linki
wszystkich warstw osobliwych są spójne. W rozdziale
rozważamy tylko przestrzenie z lokalnie skończoną
triangulacją zgodną z filtracją.
\end{definicja}
\begin{uwaga}
Dla skończonego kompleksu PL (a także dla stratyfikacji
Whitneya zespolonej rozmaitości algebraicznej po przejściu
do odpowiedniego modelu topologicznego) warunki powyższe
zapewniają lokalne modele stożkowe. W ogólnej topologii
istnienie takich modeli nie jest automatyczne.
Dla $S\subset X$ o wymiarze $n-c$ liczbę $c$ nazywamy
\emph{kowymiarem rzeczywistym} warstwy $S$.
\end{uwaga}

\begin{definicja}[Perwersyjność Goresky'ego--MacPhersona]
\label{def:perversity-xiii}
Perwersyjnością klasyczną na $X$ nazywamy funkcję
$\bar p:\{2,\ldots,n\}\to\Z$ spełniającą
$\bar p(2)=0$ oraz
$\bar p(c+1)\in\{\bar p(c),\bar p(c)+1\}$
dla $2\le c<n$. Perwersyjność górna $\bar t$
oraz dwie środkowe dane są przez
\[
\bar t(c)=c-2,\qquad
\bar m(c)=\left\lfloor\frac{c-2}{2}\right\rfloor,
\qquad
\bar n(c)=\left\lceil\frac{c-2}{2}\right\rceil.
\]
Dwie perwersyjności $\bar p,\bar q$ są
\emph{komplementarne}, gdy $\bar p(c)+\bar q(c)=c-2$
dla wszystkich $2\le c\le n$.
\end{definicja}
\begin{przyklad}[Środkowa perwersyjność]
Dla kowymiarów $c=2,3,4,5$ wartości $\bar m(c)$
są kolejno $0,0,1,1$, a wartości $\bar n(c)$
wynoszą $0,1,1,2$. Dla parzystych kowymiarów
$\bar m(c)=\bar n(c)$; dla nieparzystych mogą się
różnić. Ta drobna pozornie rozbieżność będzie istotna
w przykładzie zawieszenia torusa.
\end{przyklad}

\section{Łańcuchy dopuszczalne i operator brzegu}
\begin{definicja}[Dopuszczalność]
\label{def:allowability-xiii}
Niech $X$ będzie $n$-wymiarową pseudorozmaitością PL
z ustaloną triangulacją zgodną z filtracją.
Dla łańcucha PL $\xi$ wymiaru $i$ oznaczmy przez
$|\xi|$ sumę nośników sympleksów o niezerowych
współczynnikach. Łańcuch $\xi$ jest
\emph{$\bar p$-dopuszczalny}, gdy dla każdego $c\ge2$
\[
\dim\bigl(|\xi|\cap X_{n-c}\bigr)
\le i-c+\bar p(c)
\]
oraz te same nierówności zachodzą dla $\partial\xi$
ze stopniem $i-1$. Przyjmujemy
$\dim\varnothing=-\infty$. Dopuszczamy skończone
subdywizje zgodne z filtracją i przechodzimy do
identyfikacji wynikających z subdywizji.
Dla grup $IH_*(cL;k)$ i innych niezwartych przestrzeni
korzystamy z \emph{łańcuchów skończonych}
(o nośniku zwartym), nie z lokalnie skończonych.
Wzory dla dwóch wariantów nie są zamienne.
Kompleks takich łańcuchów oznaczamy
$I^{\bar p}C_*(X;k)$, a jego homologię przez
$IH_i^{\bar p}(X;k)$.
\end{definicja}
\begin{uwaga}[Dlaczego warunek na brzeg jest osobny?]
Nierówność dla $|\xi|$ nie gwarantuje nierówności
dla $|\partial\xi|$: brzeg może leżeć w warstwie
osobliwej, mimo że sam łańcuch ma wyższy wymiar.
Dlatego w definicji jawnie wymagamy dopuszczalności
obu łańcuchów. Wtedy $\partial^2=0$ jest zwykłą
tożsamością symplicjalną i
$I^{\bar p}C_*(X;k)$ rzeczywiście jest kompleksem.
Dla przestrzeni nielokalnie skończonej powyższy zapis
odnosi się do łańcuchów o nośniku zwartym;
wersja lokalnie skończona wymaga osobnego warunku
lokalnej skończoności. Nie utożsamiamy tych wersji.
\end{uwaga}
\begin{przyklad}[Łańcuch przechodzący przez osobliwość]
Jeśli $X$ ma wymiar $4$, a punkt $v$ jest warstwą
kowymiaru $4$, to dla $\bar m(4)=1$ i $i=1$
warunek ma postać $\dim(|\xi|\cap\{v\})\le-2$.
Zatem dopuszczalny łańcuch jednowymiarowy nie
przechodzi przez $v$. Dla $i=3$ prawa strona wynosi
$0$, więc spotkanie $v$ przez łańcuch jest dozwolone,
lecz trzeba dodatkowo sprawdzić dopuszczalność jego brzegu.
\end{przyklad}

\section{Wzór na stożek: pierwszy istotny rachunek}
\begin{twierdzenie}[Wzór stożkowy]
\label{thm:cone-ih-xiii}
Niech $L$ będzie zwartą $(n-1)$-wymiarową
warstwową pseudorozmaitością PL, $n\ge2$, a $cL$
jej otwartym stożkiem, którego wierzchołek jest
warstwą kowymiaru $n$. Niech $\bar p$ będzie klasyczną
perwersyjnością, a $k$ ustalonym ciałem. Wówczas
\[
IH_i^{\bar p}(cL;k)\cong
\begin{cases}
IH_i^{\bar p}(L;k),&0\le i<n-1-\bar p(n),\\
0,&i\ge n-1-\bar p(n),
\end{cases}
\]
przy czym po prawej dla $L$ stosujemy perwersyjność
ograniczoną do kowymiarów występujących w $L$.
\end{twierdzenie}
\begin{proof}[Dowód po przyjęciu lematu o dopuszczalnym stożkowaniu PL]
Po usunięciu wierzchołka $v$ mamy
$cL\setminus\{v\}\cong L\times(0,1)$.
Przesunięcie wzdłuż czynnika $(0,1)$ stanowi
warstwową homotopię do poziomu $L\times\{1/2\}$.
Standardowy operator pryzmatu dla tej homotopii
zachowuje nierówności dopuszczalności, gdyż nie
zmienia kowymiaru spotykanej warstwy; w konsekwencji
$IH_i(cL\setminus\{v\};k)\cong IH_i(L;k)$.

Oznaczmy $r=n-1-\bar p(n)$.
Jeżeli $i<r$, dopuszczalny $i$-łańcuch nie może
spotkać $v$, bo nierówność w wierzchołku miałaby
postać $0\le i-n+\bar p(n)<-1$.
Również dopuszczalne homologiczne relacje między
cyklami tego stopnia można zastąpić relacjami poza
wierzchołkiem: ich wymiar $i+1\le r$ nadal wymusza
uniknięcie $v$. Zatem w tym zakresie homologie są
te same co dla $L\times(0,1)$.

Dla $i\ge r$ weźmy dopuszczalny cykl $z$.
Wybierzmy wierzchołek $v$ i utwórzmy geometryczny
stożek $v*z$, po wcześniejszej subdywizji, jeśli
trzeba. Jego brzeg spełnia
$\partial(v*z)=z-v*(\partial z)=z$.
W wierzchołku warunek dopuszczalności dla wymiaru
$i+1$ to $0\le i+1-n+\bar p(n)$, co zachodzi
właśnie dla $i\ge r$. Dla pozostałych warstw
stożkowanie podwyższa jednocześnie wymiar łańcucha
i wymiar przecięcia z daną warstwą co najwyżej
o jeden. Warunki dla $z$ gwarantują więc warunki
dla $v*z$. Jeśli $z$ spotyka $v$, stosujemy następujący
lemat PL: po wspólnej subdywizji łańcuchów w
stożku, triangulacji zgodnej z filtracją oraz
operacji ``stożka łańcuchowego'' istnieje operator
$C_v$, określony na dopuszczalnych cyklach w
stopniu $i\ge r$, taki że
$\partial C_v(z)=z$, a każdy sympleks $C_v(z)$
spełnia warunki dopuszczalności. \emph{Szkic weryfikacji:}
na sympleksach nierozpiętych na $v$ jest to zwykły
stożek; na sympleksach zawierających $v$ stosuje się
normalizację kompleksu symplicjalnego i usuwa
sympleksy zdegenerowane. Wspólna subdywizja
zapewnia, że operatory brzegu oraz te redukcje
są ze sobą zgodne. Warunek w wierzchołku to
$0\le i+1-n+\bar p(n)$, a dla innych warstw
oba wymiary rosną co najwyżej o jeden.
Zatem w wysokich stopniach dopuszczalne cykle
są brzegami.

\emph{Granica dowodu:} zgodność normalizacji z
filtracją oraz pełna konstrukcja operatora $C_v$
dla arbitralnej dopuszczalnej triangulacji są
treścią standardowego lematu o stożkowaniu PL,
który tutaj przyjmujemy. Wzór na próg i
kontrola nierówności są uzasadnione jawnie,
lecz nie traktujemy powyższego szkicu jako
niezależnego dowodu całej teorii łańcuchów
przecięć.
\end{proof}
\begin{przyklad}[Stożek na torusie]
\label{ex:torus-cone-xiii}
Niech $L=T^2$ i $X=cT^2$, więc $\dim_\R X=3$.
Dla $\bar m(3)=0$ wartość graniczna wynosi
$r=3-1-0=2$. Współczynniki $k$ przyjmujemy w ciele.
Z wzoru stożkowego otrzymujemy
$IH_0^{\bar m}(X;k)\cong k$,
$IH_1^{\bar m}(X;k)\cong k^2$ i
$IH_i^{\bar m}(X;k)=0$ dla $i\ge2$.
Dla górnej perwersyjności $\bar n(3)=1$ próg to $r=1$:
pozostaje tylko $IH_0^{\bar n}(X;k)\cong k$.
To pierwszy przykład, w którym dwie środkowe
perwersyjności wykrywają różne dane.
\end{przyklad}

\section{Dualność przecięcia i przykład zawieszenia torusa}
\begin{twierdzenie}[Dualność Goresky'ego--MacPhersona]
\label{thm:ih-duality-xiii}
Niech $X$ będzie zwartą, zorientowaną, normalną
$n$-wymiarową pseudorozmaitością PL bez brzegu,
a $k$ ciałem. Dla komplementarnych perwersyjności
klasycznych $\bar p+\bar q=\bar t$ istnieje
naturalne, niedegenerujące sparowanie
$IH_i^{\bar p}(X;k)\otimes_k
 IH_{n-i}^{\bar q}(X;k)\to k$.
Jeżeli wszystkie występujące warstwy mają parzyste
kowymiar rzeczywisty, to dla środkowej
perwersyjności mamy $\bar m=\bar n$ na warstwach,
a więc otrzymujemy samodualność środkowej teorii.
\end{twierdzenie}
\begin{proof}[Zakres przedstawionego argumentu]
Na najwyższej warstwie parowanie powstaje z przecięcia
dopuszczalnych łańcuchów w położeniu ogólnym.
Warunki $\bar p(c)+\bar q(c)=c-2$ gwarantują,
że przecięcia i ich brzegi nie przechodzą przez
warstwy osobliwe w wymiarach, w których mogłyby
zniszczyć zgodność z operatorem brzegu.
Niedegeneracja jest istotą głębokiego twierdzenia
i wymaga indukcji po warstwach, wzoru stożkowego,
lokalnych klas orientacji i odpowiedniej wersji
dualności Verdiera dla kompleksu przecięcia.
Nie uzyskujemy jej z samego liczenia wymiarów:
przyjmujemy ją tutaj jako twierdzenie Goresky'ego--MacPhersona.
\end{proof}

\begin{przyklad}[Dlaczego dolna perwersyjność nie zawsze jest samodualna?]
\label{ex:suspension-torus-xiii}
Niech $X=\Sigma T^2$ będzie zawieszeniem torusa.
Jest to zwarta, zorientowana pseudorozmaitość wymiaru $3$
z dwoma punktami osobliwymi, których linkiem jest $T^2$.
Pokryjmy $X$ dwoma otwartymi stożkami $U,V$; ich
przecięcie $W=U\cap V$ jest warstwowo homotopijnie
równoważne $T^2$. Zatem
$IH_i^{\bar p}(W;k)\cong H_i(T^2;k)$, czyli
w stopniach $0,1,2$ otrzymujemy odpowiednio $k,k^2,k$.

Dla dolnej perwersyjności $\bar m(3)=0$ wzór
stożkowy daje
$IH_0^{\bar m}(U;k)=k$,
$IH_1^{\bar m}(U;k)=k^2$ oraz
$IH_i^{\bar m}(U;k)=0$ dla $i\ge2$;
to samo dotyczy $V$. W długim ciągu
Mayera--Vietorisa mamy, po opuszczeniu
symbolu $\bar m$:
\[
0\longrightarrow IH_3(X)\longrightarrow
H_2(T^2;k)\longrightarrow0,
\]
więc $IH_3(X)\cong k$. Kolejny istotny
fragment ma postać
\[
0\longrightarrow IH_2(X)\longrightarrow
k^2\xrightarrow{\,u\mapsto(u,-u)\,}k^2\oplus k^2
\longrightarrow IH_1(X)\longrightarrow
k\xrightarrow{\,t\mapsto(t,-t)\,}k\oplus k.
\]
Obie zapisane mapy są różnowartościowe.
Pierwsza ma kokernele izomorficzne z $k^2$,
zatem $IH_2(X)=0$ i $IH_1(X)\cong k^2$.
Na końcu tego ciągu otrzymujemy $IH_0(X)\cong k$.

Dla górnej perwersyjności $\bar n(3)=1$ wzór
stożkowy daje dla $U,V$ jedynie grupę
$IH_0^{\bar n}=k$, a wszystkie wyższe znikają.
Wtedy ciąg Mayera--Vietorisa redukuje się
do izomorfizmów
$IH_3^{\bar n}(X)\cong H_2(T^2;k)=k$ oraz
$IH_2^{\bar n}(X)\cong H_1(T^2;k)=k^2$;
w stopniach $1,0$ otrzymujemy dokładny ciąg
\[
0\longrightarrow IH_1^{\bar n}(X)
\longrightarrow k\xrightarrow{\,t\mapsto(t,-t)\,}
k\oplus k\longrightarrow IH_0^{\bar n}(X)
\longrightarrow0.
\]
Stąd $IH_1^{\bar n}(X)=0$ i
$IH_0^{\bar n}(X)\cong k$.
Podsumowując,
\[
IH_i^{\bar m}(X;k)=(k,k^2,0,k),\qquad
IH_i^{\bar n}(X;k)=(k,0,k^2,k)
\quad (i=0,1,2,3).
\]
Niedegenerujące parowanie z
Twierdzenia~\ref{thm:ih-duality-xiii} łączy
$IH_1^{\bar m}$ z $IH_2^{\bar n}$,
a nie z $IH_2^{\bar m}=0$.
Przyczyną jest nieparzysty kowymiar osobliwości;
warunek Witta nie jest tutaj spełniony.
\emph{Zależność dowodowa:} użyliśmy standardowego
ciągu Mayera--Vietorisa dla dopuszczalnych kompleksów
łańcuchów; jego dowód wymaga przejścia do
kompleksu snopowego lub wersji podziału jedności
na łańcuchach PL i jest przyjętym twierdzeniem
pomocniczym tego rozdziału.
\end{przyklad}

\section{Kompleksy konstruktybilne, podpory i podpory nadzwyczajne}
Przenosimy się teraz do zespolonej geometrii
algebraicznej. Niech $X$ będzie \emph{zespoloną odmianą
algebraiczną} skończonego typu (nie zakładamy jej
gładkości) z ustaloną skończoną
stratyfikacją, której warstwy są gładkie i spójne.
Przez $D_c^b(X;k)$ oznaczamy ograniczoną kategorię
pochodną kompleksów snopów $k$-przestrzeni wektorowych,
których snopy kohomologii są lokalnymi systemami
skończonego rzędu na każdej warstwie.
Dla punktu $x\in X$ niech $i_x:\{x\}\hookrightarrow X$.

\begin{definicja}[Warunki podpory i kopodpory]
\label{def:perverse-support-xiii}
Kompleks $K\in D_c^b(X;k)$ spełnia
\emph{warunek podpory}, jeżeli dla każdego punktu $x$
leżącego w warstwie $S$ wymiaru zespolonego $s$
$H^j(i_x^*K)=0$ dla $j>-s$.
Spełnia \emph{warunek kopodpory}, jeżeli
$H^j(i_x^!K)=0$ dla $j<s$.
Kompleks spełniający oba warunki nazywamy
\emph{snopem perwersyjnym} (w środkowej konwencji
kohomologicznej). Kategorię takich obiektów oznaczamy
$\operatorname{Perv}(X;k)$.
\end{definicja}
\begin{uwaga}[Dlaczego przy kopodporze występuje $i_x^!$?]
Kiełek $i_x^*K$ sprawdza, co widać w sąsiedztwie
punktu $x$, a $i_x^!K$ rejestruje klasy lokalne
\emph{z nośnikiem} w $x$. Występują tu dokładnie
obie operacje omawiane przy dualności Verdiera.
W istocie kopodpora $K$ jest podporą
$\mathbb D_XK$ po odpowiedniej dualizacji.
\end{uwaga}

\begin{stwierdzenie}[Podstawowy przykład: rozmaitość gładka]
\label{prop:smooth-perverse-xiii}
Niech $M$ będzie gładką zespoloną rozmaitością
wymiaru zespolonego $d$, a $L$ lokalnym systemem
skończeniewymiarowych $k$-przestrzeni wektorowych
na $M$. Wtedy $L[d]$ jest snopem perwersyjnym.
\end{stwierdzenie}
\begin{proof}
Dla punktu $x\in M$ zwykły kiełek $i_x^*L[d]$
ma jedyną niezerową kohomologię w stopniu $-d$.
Z drugiej strony lokalna kohomologia pary
$(B^{2d},B^{2d}\setminus\{0\})$ ma jedyną niezerową
grupę w stopniu $2d$. Dlatego $i_x^!L[d]$
ma jedyną niezerową kohomologię w stopniu
$2d-d=d$. Obie nierówności z
Definicji~\ref{def:perverse-support-xiii}
są spełnione. Ponieważ konstrukcja zachowuje
lokalne trywializacje $L$, wynik nie zależy od
wybranego punktu.
\end{proof}

\begin{twierdzenie}[Serce struktury $t$]
\label{thm:heart-perverse-xiii}
Dla zespolonej rozmaitości algebraicznej $X$
i ciała $k$ warunki podpory i kopodpory wyznaczają
ograniczoną strukturę $t$ na $D_c^b(X;k)$.
Jej sercem jest abelowa kategoria
$\operatorname{Perv}(X;k)$; istnieją funktory
kohomologii perwersyjnej $\,{}^p\!H^j$.
Dualność Verdiera zachowuje to serce i zamienia
warunek podpory z warunkiem kopodpory.
\end{twierdzenie}
\begin{proof}[Status dowodu]
Pełny dowód wymaga konstrukcji perwersyjnych
truncacji, sklejenia struktur $t$ wzdłuż podziału
na warstwy oraz sprawdzenia dokładności funktorów
odnośnych do włożenia otwartego i domkniętego.
Nie ma go w poprzednim rozdziale, dlatego
przyjmujemy tutaj twierdzenie o strukturze $t$
jako jawny wynik podstawowy. Nierówności w
Definicji~\ref{def:perverse-support-xiii}
są pełną definicją testowaną na punktach,
a nie skrótowym określeniem „złego stopnia”.
\end{proof}

\section{Przedłużenie pośrednie i kompleks przecięcia}
\begin{definicja}[Przedłużenie pośrednie]
\label{def:middle-extension-xiii}
Niech $j:U\hookrightarrow X$ będzie otwartym włożeniem,
a $P\in\operatorname{Perv}(U;k)$.
Przedłużeniem pośrednim $j_{!*}P$ nazywamy obraz
morfizmu w kategorii abelowej $\operatorname{Perv}(X;k)$:
\[
{}^p\!H^0(Rj_!P)\longrightarrow {}^p\!H^0(Rj_*P).
\]
Jest to jedyny perwersyjny kompleks przedłużający $P$,
który nie ma nietrywialnego podobiektu ani ilorazu
podpartego w $X\setminus U$.
\end{definicja}
\begin{definicja}[Kompleks przecięcia]
\label{def:ic-xiii}
Niech $X$ będzie zredukowaną, czystowymiarową zespoloną
odmianą algebraiczną wymiaru zespolonego $d$, a
$j:U=X_{\mathrm{reg}}\hookrightarrow X$ włożeniem
gładkiej części. Jej \emph{kompleksem przecięcia}
ze współczynnikami w $k$ nazywamy
$IC_X(k)=j_{!*}(k_U[d])$.
Dla dowolnego $i\in\Z$ definiujemy
\[
IH^i(X;k):=\mathbb H^{i-d}(X,IC_X(k)),
\]
gdzie $\mathbb H$ oznacza hiperkohomologię kompleksu snopów.
Dla gładkiego $X$ otrzymujemy $IC_X(k)=k_X[d]$
i $IH^i(X;k)=H^i(X;k)$.
\end{definicja}
\begin{uwaga}[Dwie konwencje stopni]
Często autorzy oznaczają przez $IC_X$ kompleks
\emph{nieprzesunięty}, którego hiperkohomologia
jest $IH^i$. U nas $IC_X$ jest \emph{perwersyjny}
i zawiera przesunięcie $[d]$; stąd konieczne jest
$\mathbb H^{i-d}$, a nie $\mathbb H^i$.
Nie należy mylić tego z przesunięciem rzeczywistego
wymiaru $2d$ występującym w dualności Verdiera.
\end{uwaga}

\begin{twierdzenie}[Porównanie wersji łańcuchowej i snopowej]
\label{thm:ic-comparison-xiii}
Niech $X$ będzie zwartą zespoloną odmianą
algebraiczną czystego wymiaru $d$,
stratyfikowaną zgodnie z lokalnym modelem stożkowym,
a współczynniki należą do ciała $k$.
Wówczas hiperkohomologia perwersyjnego kompleksu
$IC_X(k)$ oblicza środkową kohomologię przecięć
Goresky'ego--MacPhersona, z uwzględnieniem
konwencji przesunięcia w Definicji~\ref{def:ic-xiii}.
Ponadto $\mathbb D_XIC_X(k)\simeq IC_X(k^\vee)$,
po naturalnym utożsamieniu dualnego systemu
współczynników i wyborze normalizacji.
\end{twierdzenie}
\begin{proof}[Zakres]
Podstawą porównania jest lokalny wzór stożkowy
i aksjomaty charakteryzujące kompleks przecięcia:
warunek normalizacji na części gładkiej,
odpowiednie ograniczenia podpór oraz reguła
przedłużenia przez kolejne warstwy osobliwe.
Dowód przeprowadza się indukcyjnie po stratach;
nie wynika on automatycznie z definicji $j_{!*}$.
Przyjmujemy go jako twierdzenie porównawcze.
Samodualność korzysta z komplementarności
środkowych perwersyjności, które na przestrzeniach
zespolonych zgadzają się na kowymiarach warstw
\emph{zespolonych} (czyli parzystych kowymiarach
rzeczywistych). Nie twierdzimy tego dla dowolnej
stratyfikacji rzeczywistej; w szczególności
przykład $\Sigma T^2$ nie jest modelem
algebraicznym o tej własności.
\end{proof}

\section{Obliczenie: krzywa węzłowa i jej normalizacja}
\label{sec:nodal-perverse-xiii}
Rozpatrzmy najpierw lokalnie
$X=\{zw=0\}\cap B_\varepsilon\subset\C^2$.
Normalizacja $\nu:\widetilde X\to X$ rozdziela
obie gałęzie; $\widetilde X$ jest sumą rozłączną
dwóch małych dysków. Jest to odwzorowanie skończone,
a zatem właściwe, o zerowym wymiarze włókien.

\begin{stwierdzenie}[Kompleks przecięcia węzła]
\label{prop:node-ic-xiii}
Dla $X$ powyżej i dowolnego ciała $k$ zachodzi
naturalny quasi-izomorfizm
$IC_X(k)\simeq\nu_*k_{\widetilde X}[1]$.
W szczególności w węźle $x=0$ mamy
$H^{-1}(i_x^*IC_X(k))\cong k^2$,
a wszystkie pozostałe grupy kohomologii kiełka
są zerowe.
\end{stwierdzenie}
\begin{proof}
Ponieważ $\nu$ jest skończona, jej włókna są
skończone i $R^q\nu_*k_{\widetilde X}=0$ dla $q>0$;
w naszym lokalnym modelu wynika to również
bezpośrednio z faktu, że przeciwobrazy dostatecznie
małych otoczeń są sumą dwóch dysków.
Kompleks $K=\nu_*k_{\widetilde X}[1]$ ogranicza się
do $k_U[1]$ na części gładkiej $U$.
W punkcie $x$ jego zwykły kiełek ma kohomologię
$k^2$ wyłącznie w stopniu $-1$. Natomiast
$H^q(i_x^!K)$ obliczamy przez lokalne pary
dysków i ich nakłutych wersji: niezerowa grupa
$k^2$ pojawia się wyłącznie dla $q=1$.
Takie dane spełniają warunki podpory i kopodpory;
w szczególności $H^0(i_x^*K)=H^0(i_x^!K)=0$.
Przy użyciu adjunkcji dla włożenia punktu
ostatnie dwa zaniki oznaczają, że $K$ nie ma
podobiektów ani ilorazów perwersyjnych
podpartych wyłącznie w $x$. Z jednoznaczności
przedłużenia pośredniego otrzymujemy tezę.
\end{proof}

\begin{przyklad}[Krzywa rzutowa: kohomologia zwykła i przecięć]
Niech $C$ będzie nieredukowalną zespoloną
krzywą rzutową z jednym zwyczajnym węzłem,
której normalizacja $\nu:\mathbb P^1\to C$
jest izomorfizmem poza dwiema różnymi preobrazami
węzła (np. odpowiednia wymierna krzywa sześcienna).
Topologicznie $C$ powstaje z $S^2$ przez
utożsamienie dwóch punktów, więc
$H^0(C;k)\cong H^2(C;k)\cong k$,
$H^1(C;k)\cong k$.
Z drugiej strony
$IC_C(k)=\nu_*k_{\mathbb P^1}[1]$ i właściwość $\nu$
dają $IH^i(C;k)\cong H^i(\mathbb P^1;k)$.
A zatem $IH^1(C;k)=0$ przy zachowaniu
$IH^0(C;k)=IH^2(C;k)=k$.
Nowa teoria nie „ignoruje” osobliwości: rozdziela
lokalne gałęzie i odzyskuje prawidłową dualność.
\end{przyklad}

\begin{center}
\begin{tikzpicture}[scale=0.9,font=\small,>=Stealth]
\draw[manifoldblue,thick] (-3.8,0) ellipse (0.62 and 0.74);
\draw[manifoldblue,thick] (-2.3,0) ellipse (0.62 and 0.74);
\fill[manifoldred] (-3.1,0) circle (2.3pt);
\node[above] at (-3.1,0.18) {dwie gałęzie};
\draw[->,thick] (-1.35,0)--(0.7,0)
 node[midway,above] {$\nu$};
\draw[manifoldblue,thick] (1.65,0) ellipse (1.02 and 0.8);
\draw[manifoldblue,thick] (2.2,0.1) ellipse (1.02 and 0.8);
\fill[manifoldred] (1.92,0.71) circle (2.3pt);
\node[above] at (2.3,1.0) {węzeł};
\node[below] at (-3, -1.1) {$\widetilde X$ (schematycznie)};
\node[below] at (2, -1.1) {$X$ (schematycznie)};
\end{tikzpicture}
\par\smallskip\small Normalizacja oddziela lokalne gałęzie; rysunek
pokazuje ideę, a nie dokładny obraz algebraiczny krzywej.
\end{center}

\section{Twierdzenie o rozkładzie i małe odwzorowania}
\begin{definicja}[Odwzorowanie małe i półmałe]
Niech $f:Y\to X$ będzie właściwym morfizmem
zespolonych odmian algebraicznych czystego wymiaru $d$,
generycznie skończonym i dominującym. Dla $r\ge1$
niech $X_r=\{x\in X:\dim_\C f^{-1}(x)\ge r\}$.
Odwzorowanie jest \emph{półmałe}, jeśli
$\operatorname{codim}_\C X_r\ge2r$ dla każdego $r\ge1$,
a \emph{małe}, jeśli wszędzie zachodzi
nierówność ostra. Przyjmujemy
$\operatorname{codim}\varnothing=+\infty$.
\end{definicja}
\begin{uwaga}
Przy powyższej definicji już założyliśmy, że $f$ jest
dominujący i generycznie skończony pomiędzy
czystowymiarowymi przestrzeniami tego samego wymiaru;
dla ogólnych morfizmów półmałość formułuje się przez
stratyfikację i nierówność
$2\dim f^{-1}(x)\le\dim Y-\dim S$
dla $x$ w warstwie $S$.
\end{uwaga}

\begin{twierdzenie}[Twierdzenie o rozkładzie --- wersja projektowa]
\label{thm:decomposition-xiii}
Niech $f:Y\to X$ będzie projektowym morfizmem
zespolonych odmian algebraicznych, a $Y$
niech będzie gładka i czystowymiarowa o wymiarze
zespolonym $d$. W kategorii $D_c^b(X;\Q)$ istnieje
(niekanoniczny na ogół) izomorfizm
\[
Rf_*\Q_Y[d]\simeq
\bigoplus_{j\in\Z}{}^p\!H^j(Rf_*\Q_Y[d])[-j].
\]
Każdy snop perwersyjny po prawej jest półprosty,
a więc rozkłada się jako skończona suma kompleksów
$IC_{\overline S}(L)$, przedłużonych przez zero
do $X$, gdzie $S$ przebiega odpowiednie gładkie,
lokalnie domknięte warstwy, a $L$ są
półprostymi lokalnymi systemami $\Q$-przestrzeni
wektorowych na $S$. W konsekwencji pochodny obraz
właściwy rozpada się na przesunięte kompleksy przecięcia.
\end{twierdzenie}
\begin{proof}[Jawnie używany wynik głęboki]
Nie dowodzimy tu twierdzenia o rozkładzie BBDG.
Jego klasyczne dowody wymagają czystości i wag,
polaryzowanych struktur Hodge'a albo innych
silnych twierdzeń geometrycznych.
Z konstrukcji perwersyjnej struktury $t$ wynika
istnienie snopów ${}^p\!H^j$, ale \emph{nie}
wynika z niej rozszczepienie kompleksu na sumę
swoich kohomologii perwersyjnych.
Właśnie to rozszczepienie i półprostota są
nietrywialną treścią przyjętego twierdzenia.
\end{proof}
\begin{wniosek}[Rozdzielczości małe]
\label{cor:small-map-ic-xiii}
Niech $f:Y\to X$ będzie właściwą \emph{małą}
rozdzielczością osobliwości, tj. biracjonalnym
morfizmem z gładkiej czystowymiarowej
$Y$ na nieredukowalną $X$, który jest izomorfizmem
nad gęstą częścią gładką i spełnia nierówności
Definicji powyżej. Wówczas
$Rf_*\Q_Y[d]\simeq IC_X(\Q)$.
\end{wniosek}
\begin{proof}[Wyjaśnienie zastosowanego kryterium]
Małość eliminuje nietrywialne składniki
perwersyjne podparte w mniejszych stratach;
na części otwartej $Rf_*\Q_Y[d]$ ogranicza się
do $\Q_U[d]$. Przyjmujemy standardowe twierdzenie
o małych mapach (wynikające z aksjomatów
kompleksu przecięcia), które daje jego
identyfikację z przedłużeniem pośrednim.
Małość nie oznacza jedynie, że włókna są spójne:
potrzebne są jawne nierówności wymiarów.
\end{proof}

\section{Przykład geometryczny: rozdmuchanie płaszczyzny}
\label{sec:blowup-decomp-xiii}
Niech $f:\widetilde{\C^2}\to\C^2$ będzie
rozdmuchaniem początku. Jest to projektowe,
a więc właściwe odwzorowanie; poza $0$ jest
izomorfizmem, natomiast $f^{-1}(0)\cong\mathbb P^1$.
W szczególności $f$ jest półmałe, lecz nie małe:
$\operatorname{codim}_\C\{0\}=2=2\dim_\C f^{-1}(0)$.

\begin{stwierdzenie}[Rozkład dla rozdmuchania]
\label{prop:blowup-decomp-xiii}
W kategorii $D_c^b(\C^2;\Q)$ mamy
\[
Rf_*\Q_{\widetilde{\C^2}}[2]
\simeq\Q_{\C^2}[2]\oplus\Q_{\{0\}}.
\]
Dla rozdmuchania punktu na $\mathbb P^2$ otrzymujemy
$\dim_\Q H^0=\dim_\Q H^4=1$ oraz
$\dim_\Q H^2=2$ dla rozdmuchanej powierzchni.
\end{stwierdzenie}
\begin{proof}
Twierdzenie o zmianie bazy dla właściwego odwzorowania
utożsamia kiełek $R^qf_*\Q$ w punkcie $x$
z $H^q(f^{-1}(x);\Q)$. Dla $x\ne0$ włókno jest
punktem, a dla $x=0$ jest $\mathbb P^1$.
Wynika stąd $R^0f_*\Q\cong\Q_{\C^2}$,
$R^2f_*\Q\cong\Q_{\{0\}}$ i zanik pozostałych.
Nie wolno z tego \emph{samego} wnioskować, że
$Rf_*\Q$ jest sumą swoich snopów kohomologii:
morfizmy łączące w kategorii pochodnej mogą
być nietrywialne. Rozszczepienie podane w tezie
wynika z Twierdzenia~\ref{thm:decomposition-xiii};
przesunięcie o $[2]$ lokuje drugi składnik
w stopniu $0$. Po zastosowaniu $R\Gamma$
do analogicznego wzoru dla rozdmuchania $\mathbb P^2$
dodatkowa klasa w stopniu $2$ zwiększa drugą
liczbę Bettiego o jeden, przy zachowaniu zerowego
i czwartego stopnia.
\end{proof}
\begin{uwaga}[Skąd potrzeba półprostoty?]
Dane o wszystkich $R^qf_*\Q$ nie określają
jeszcze klasy izomorfizmu $Rf_*\Q$ w kategorii
pochodnej. Twierdzenie o rozkładzie kontroluje
również klasy rozszerzeń między tymi składnikami.
Właśnie dlatego jest tak silnym narzędziem
w geometrii algebraicznej.
\end{uwaga}

\section{Względne twarde twierdzenie Lefschetza}
\begin{twierdzenie}[Względne twarde Lefschetz]
\label{thm:relative-hard-lefschetz-xiii}
Niech $f:Y\to X$ będzie projektowym morfizmem
zespolonych odmian algebraicznych, $Y$
gładką czystowymiarową o wymiarze $d$,
a $\eta\in H^2(Y;\Q)$ klasą pierwszego Cherna
wiązki liniowej względnie ample względem $f$.
Wtedy dla każdego $r\ge0$ mnożenie przez
$\eta^r$ wyznacza izomorfizm
\[
\eta^r:{}^p\!H^{-r}(Rf_*\Q_Y[d])
\longrightarrow{}^p\!H^{r}(Rf_*\Q_Y[d]).
\]
\end{twierdzenie}
\begin{proof}[Status]
Jest to niezależne głębokie twierdzenie, a nie
formalny wniosek z definicji snopu perwersyjnego.
Pełny dowód wymaga aparatu Hodge'a lub teorii wag.
Warto jednak zauważyć analogię ze zwykłym twardym
twierdzeniem Lefschetza dla gładkich odmian
projektowych: tu stopnie kohomologii zastąpione
są stopniami perwersyjnymi.
\end{proof}
\begin{przyklad}[Projekcja z włóknem $\mathbb P^1$]
Niech $f:\mathbb P^1\times\mathbb P^1\to\mathbb P^1$
będzie pierwszą projekcją. Wtedy
$Rf_*\Q_{\mathbb P^1\times\mathbb P^1}[2]
\simeq\Q_{\mathbb P^1}[2]\oplus\Q_{\mathbb P^1}$.
Ponieważ $P=\Q_{\mathbb P^1}[1]$ jest perwersyjny,
dwie składowe są $P[1]$ i $P[-1]$, zatem
${}^p\!H^{-1}\cong P\cong{}^p\!H^{1}$.
Klasa $\eta$ pochodząca z drugiego czynnika
przenosi generator $H^0(\mathbb P^1;\Q)$ na
generator $H^2(\mathbb P^1;\Q)$ w każdym włóknie;
stąd w tym przypadku izomorfizm względnego
Lefschetza można sprawdzić bezpośrednio.
\end{przyklad}

\section{Granice aktualnego rozdziału i pomost do geometrii algebraicznej}
Uzupełniliśmy schemat
\[
\text{łańcuchy dopuszczalne}
\longrightarrow IH^*(X)
\longleftrightarrow IC_X
\longrightarrow\operatorname{Perv}(X)
\longrightarrow Rf_*IC_Y.
\]
Łańcuchowa konstrukcja i rachunki stożkowe
mają szczegółowe uzasadnienia. Pojęcie snopu
perwersyjnego jest zdefiniowane przez
sprawdzalne warunki na kiełkach i kopiełkach;
przykłady rozmaitości gładkiej oraz węzła
zostały obliczone. Twierdzenie o rozkładzie,
względne twarde Lefschetz i pełna równoważność
konstrukcji łańcuchowej ze snopową są
\emph{jawnie przyjętymi} głębokimi twierdzeniami.

Przy późniejszym scalaniu z geometrią algebraiczną
warto dołączyć rozdziały o strukturach Hodge'a,
monodromii i cyklach zanikających, dualności
Grothendiecka oraz o \emph{$t$-strukturach} w kategoriach
pochodnych. Nie utożsamiamy tu kohomologii
singularnej i przecięć ze snopami koherentnymi
ani z kohomologią etalną bez dodatkowych twierdzeń.



\section{Uwaga redakcyjna po audycie rozdziału XIII}
\label{sec:xiii-status-po-audycie}
\begin{uwaga}[Zakres samodzielności dowodów]
Rachunek wzoru stożkowego wymaga standardowego lematu
PL o stożkowaniu dopuszczalnych łańcuchów; jawny
rachunek nierówności nie zastępuje konstrukcji
tych operatorów. Formuła Mayera--Vietorisa dla
homologii przecięć wymaga odpowiedniej wersji
kompleksów łańcuchów dopuszczalnych i ich
niezmienniczości przy przejściu do pokryć otwartych.
Z kolei porównanie $IC_X$ z teorią łańcuchową,
istnienie perwersyjnej struktury $t$, dualność
Goresky'ego--MacPhersona, twierdzenie o rozkładzie
i względne twarde Lefschetz są \\emph{jawnie
przyjętymi} głębokimi wynikami. W dalszym
rozdziale rozwijamy aparat kategorii pochodnych
i teorii Hodge'a potrzebny do precyzyjnego
sformułowania tych zależności; nie ogłaszamy
przez to automatycznie pełnych dowodów tych twierdzeń.
\end{uwaga}
\chapter{Kategorie pochodne, struktury Hodge'a, cykle zanikające i monodromia}
\label{ch:xiv-hodge-monodromia}

\section{Motywacja i kolejność konstrukcji}
W rozdziale~\ref{ch:ih-perverse-xiii} pojawiły się trzy zjawiska,
których nie da się opisać samymi grupami kohomologii:
obiekt $Rf_*K$ nie jest na ogół określony przez swoje snopy
$R^qf_*K$, kompleks przecięcia jest definiowany przez
warunki na kiełkach i kopiełkach, a w degenerującej rodzinie
rozmaitości istotne są nie tylko grupy $H^q$, ale także sposób,
w jaki zmieniają się one przy okrążaniu punktu krytycznego.

Najpierw rozwiniemy minimum rachunku w kategoriach pochodnych,
następnie strukturę Hodge'a w przypadku gładkim i jej wersję
mieszaną dla prostego przykładu osobliwego. Na koniec
zdefiniujemy cykle pobliskie oraz zanikające i obliczymy
monodromię dla osobliwości kwadratowych.
Wszystkie współczynniki, jeżeli nie zaznaczono inaczej, należą
do $\Q$; przy lokalnych obliczeniach analitycznych jawnie
używamy też $\C$ lub $\Z$.

\section{Kompleksy, odwzorowania stożkowe i trójkąty wyróżnione}
\begin{definicja}[Kompleks kohomologiczny]
Niech $\mathcal A$ będzie kategorią abelową. Kompleks
kohomologiczny $A^\bullet$ jest rodziną obiektów
$A^i\in\mathcal A$ i morfizmów
$d_A^i:A^i\to A^{i+1}$ spełniających
$d_A^{i+1}d_A^i=0$. Dla $m\in\Z$ definiujemy
przesunięcie $A[m]$ przez
$A[m]^i=A^{i+m}$ oraz $d_{A[m]}^i=(-1)^m d_A^{i+m}$.
Kohomologię oznaczamy
$H^i(A^\bullet)=\ker d_A^i/\operatorname{im}d_A^{i-1}$.
Morfizm kompleksów $f:A^\bullet\to B^\bullet$ jest
\emph{quasi-izomorfizmem}, gdy $H^i(f)$ jest izomorfizmem
dla każdego $i$.
\end{definicja}

\begin{definicja}[Stożek morfizmu]
Dla morfizmu kompleksów $f:A^\bullet\to B^\bullet$
definiujemy $\operatorname{Cone}(f)^i=B^i\oplus A^{i+1}$ oraz
\[
d_{\operatorname{Cone}(f)}(b,a)=(d_Bb+f(a),-d_Aa).
\]
Istnieje naturalny ciąg morfizmów kompleksów
$B^\bullet\to\operatorname{Cone}(f)\to A^\bullet[1]$,
a ciąg $A^\bullet\xrightarrow f B^\bullet\to
\operatorname{Cone}(f)\to A^\bullet[1]$
nazywamy standardowym trójkątem stożkowym.
\end{definicja}
\begin{lemat}[Stożek jest kompleksem]
\label{lem:cone-complex-xiv}
Dla każdego morfizmu kompleksów $f:A^\bullet\to B^\bullet$
operator $d_{\operatorname{Cone}(f)}$ spełnia $d^2=0$.
\end{lemat}
\begin{proof}
Dla $(b,a)\in B^i\oplus A^{i+1}$ otrzymujemy
$d^2(b,a)=(d_B^2b+d_Bf(a)-f(d_Aa),d_A^2a)$.
Pierwszy i ostatni składnik znikają, a środkowe znoszą się,
gdyż $f$ jest morfizmem kompleksów, tj. $d_Bf=fd_A$.
\end{proof}

\begin{twierdzenie}[Długi ciąg stożka]
\label{thm:cone-les-xiv}
Dla $f:A^\bullet\to B^\bullet$ istnieje naturalny długi
ciąg dokładny
\[
\cdots\to H^i(A)\xrightarrow{H^i(f)}H^i(B)
\to H^i(\operatorname{Cone}(f))\to H^{i+1}(A)\to\cdots.
\]
W szczególności $f$ jest quasi-izomorfizmem wtedy i tylko wtedy,
gdy $\operatorname{Cone}(f)$ jest acykliczny.
\end{twierdzenie}
\begin{proof}
Rzut na drugi czynnik daje krótki ciąg dokładny kompleksów
$0\to B^\bullet\to\operatorname{Cone}(f)\to A^\bullet[1]\to0$.
Z konstrukcji homomorfizmu łączącego otrzymujemy
$H^{i+1}(A)\to H^{i+1}(B)$ równe, z przyjętą konwencją
znaku, $H^{i+1}(f)$: istotnie reprezentant $(0,a)$
klasy z $H^i(A[1])$ ma różniczkę $(f(a),0)$.
Długi ciąg krótkiego ciągu kompleksów jest dokładnie
ciągiem z tezy. Jeśli $H^i(f)$ są izomorfizmami,
dokładność wymusza $H^i(\operatorname{Cone}(f))=0$.
Odwrotnie, gdy kohomologie stożka znikają, te same
fragmenty ciągu wymuszają jednocześnie różnowartościowość
i surjektywność wszystkich $H^i(f)$.
\end{proof}

\begin{przyklad}[Niebanalne zjawisko zapisane przez stożek]
Dla morfizmu kompleksów $\Z\xrightarrow{\cdot m}\Z$
(skupionych w stopniu $0$) stożek ma w stopniach
$-1,0$ grupy $\Z,\Z$ i różniczkę będącą mnożeniem
przez $m$. Dla $m\neq0$ jego jedyną niezerową kohomologią
jest $H^0\cong\Z/m\Z$. W ten sposób obiekt ilorazowy
można odtworzyć z morfizmu, nie wybierając rozszczepienia.
\end{przyklad}

\begin{definicja}[Kategoria pochodna]
Kategorię homotopii $K(\mathcal A)$ otrzymujemy z kategorii
kompleksów w $\mathcal A$ przez utożsamienie morfizmów
łańcuchowo homotopijnych. Jej \emph{lokalizację względem
quasi-izomorfizmów} oznaczamy $D(\mathcal A)$ i nazywamy
kategorią pochodną. $D^b(\mathcal A)$ oznacza pełną
podkategorię obiektów o ograniczonej kohomologii.
W tej kategorii obrazy trójkątów stożkowych są
\emph{trójkątami wyróżnionymi}.
\end{definicja}
\begin{uwaga}[Co pozostaje poza rachunkiem elementarnym?]
Sama definicja lokalizacji nie zapewnia automatycznie
istnienia zbiorów morfizmów o właściwej wielkości ani
formalizmu funktorów pochodnych. Dla rozważanych dalej
kategorii snopów stosujemy standardowe twierdzenia
o rezolwentach iniektywnych i kompleksach $K$-iniektywnych.
Nie utożsamiamy również każdego trójkąta wyróżnionego
z dosłownym krótkim ciągiem snopów.
\end{uwaga}

\section{Funktory pochodne i dlaczego nie wystarczają snopy \texorpdfstring{$R^qf_*$}{Rq f*}}
\begin{definicja}[Pochodny obraz prosty i hiperkohomologia]
Niech $f:X\to Y$ będzie ciągłym odwzorowaniem
lokalnie zwartych przestrzeni, a $\mathcal F^\bullet$
kompleksem snopów $k$-przestrzeni wektorowych na $X$.
Po wybraniu odpowiedniej rezolwenty iniektywnej
$\mathcal F^\bullet\to I^\bullet$ definiujemy
$Rf_*\mathcal F^\bullet=f_*I^\bullet$ w kategorii pochodnej
oraz $R\Gamma(X,\mathcal F^\bullet)$ w podobny sposób.
Niech $\mathbb H^q(X,\mathcal F^\bullet)=
H^q(R\Gamma(X,\mathcal F^\bullet))$.
Dla pojedynczego snopu $\mathcal F$ mamy
$R^qf_*\mathcal F=H^q(Rf_*\mathcal F)$.
\end{definicja}
\begin{twierdzenie}[Ciąg spektralny Leraya]
\label{thm:leray-derived-xiv}
Niech $f:X\to Y$ będzie ciągłym odwzorowaniem
przestrzeni topologicznych, a $\mathcal F$ snopem grup
abelowych na $X$. Dla zwykłej kohomologii snopowej
i nieujemnych funktorów pochodnych otrzymujemy
pierwszoćwiartkowy, zbieżny do kohomologii w każdym
stopniu całkowitym ciąg spektralny
$E_2^{p,q}=H^p(Y,R^qf_*\mathcal F)
\Rightarrow H^{p+q}(X,\mathcal F)$.
\end{twierdzenie}
\begin{proof}[Źródło twierdzenia]
Kompozycja $\Gamma(Y,-)\circ f_*$ jest równa
$\Gamma(X,-)$. Twierdzenie o ciągu spektralnym
złożenia funktorów pochodnych (wraz z warunkiem
acykliczności obrazów rezolwent iniektywnych)
daje wskazaną stronę $E_2$ i jej zbieżność.
Sama konstrukcja stożka nie zastępuje tego
twierdzenia o rezolwentach.
\end{proof}
\begin{uwaga}[Dane kohomologiczne nie określają obiektu]
W kategorii pochodnej nie jest na ogół prawdą, że
$K\simeq\bigoplus_qH^q(K)[-q]$.
Taki rozkład może być blokowany przez nietrywialne
morfizmy łączące w trójkątach truncacji.
Dlatego obliczenie samych $R^qf_*\mathcal F$ nie dowodzi
rozkładu $Rf_*\mathcal F$; to właśnie istotna treść
twierdzenia z rozdziału~\ref{ch:ih-perverse-xiii}.
\end{uwaga}

\begin{przyklad}[Te same snopy kohomologii, różne obiekty pochodne]
\label{ex:derived-not-split-xiv}
Niech $X=S^2$, a $\Q_X$ będzie snopem stałym.
Standardowa identyfikacja
$\operatorname{Hom}_{D(X;\Q)}(\Q_X,\Q_X[2])
\cong H^2(S^2;\Q)\cong\Q$
pozwala wybrać niezerowy morfizm
$e:\Q_X\to\Q_X[2]$.
Zdefiniujmy $K_e=\operatorname{Cone}(e)[-2]$.
Z długiego ciągu stożka wynika, że jedynymi
niezerowymi snopami kohomologii $K_e$ są
$H^0(K_e)\cong\Q_X$ i $H^1(K_e)\cong\Q_X$.
Dla $e=0$ otrzymujemy
$K_0\simeq\Q_X\oplus\Q_X[-1]$.
Dla $e\ne0$ trójkąt
$\Q_X\to K_e\to\Q_X[-1]\xrightarrow{e[-1]}\Q_X[1]$
nie jest rozszczepiony, ponieważ jego morfizm łączący
jest niezerowy. Jest to trójkąt standardowej truncacji
$K_e$, więc $K_e$ nie jest izomorficzny z sumą
$\Q_X\oplus\Q_X[-1]$.
Widzimy dwa obiekty z izomorficznymi snopami
kohomologii, lecz różnymi klasami izomorfizmu
w kategorii pochodnej.
\end{przyklad}

\section{Struktury Hodge'a: od rozkładu form do filtracji}
\begin{definicja}[Czysta struktura Hodge'a]
Niech $m\in\Z$. \emph{Racjonalną czystą strukturą
Hodge'a wagi $m$} na skończeniewymiarowej przestrzeni
$V_{\Q}$ nazywamy rozkład
$V_{\C}=V_{\Q}\otimes_{\Q}\C=\bigoplus_{p+q=m}V^{p,q}$,
w którym $\overline{V^{p,q}}=V^{q,p}$; kreska oznacza
sprzężenie względem struktury racjonalnej.
Filtracją Hodge'a nazywamy
$F^pV_{\C}=\bigoplus_{r\ge p}V^{r,m-r}$.
\end{definicja}
\begin{stwierdzenie}[Rozkład a dwie filtracje]
\label{prop:hodge-filtration-xiv}
Dla czystej struktury Hodge'a wagi $m$ zachodzi
$V_{\C}=F^pV_{\C}\oplus\overline{F^{m-p+1}V_{\C}}$
na każdym $p\in\Z$. Odwrotnie, malejąca skończona
filtracja $F^\bullet V_{\C}$ spełniająca te równości
określa jednoznacznie strukturę Hodge'a wagi $m$ przez
$V^{p,q}=F^pV_{\C}\cap\overline{F^qV_{\C}}$
dla $p+q=m$.
\end{stwierdzenie}
\begin{proof}
W pierwszą stronę $F^p$ składa się bezpośrednia suma
wszystkich $V^{r,m-r}$ z $r\ge p$. Z kolei
$\overline{F^{m-p+1}}$ jest sumą składników
$V^{r,m-r}$ o $r\le p-1$. Składniki są rozłączne
i razem wyczerpują wszystkie pary $r+(m-r)=m$.
W drugą stronę dwie przeciwstawne filtracje
wyznaczają odpowiednią rozkładową gradację:
$F^p\cap\overline{F^{m-p}}$ jest jej $p$-tym składnikiem,
a suma tych przecięć jest bezpośrednia, co wynika
indukcyjnie z równości dla kolejnych $p$.
Zastosowanie sprzężenia zamienia $p$ i $m-p$.
Jest to standardowa rekonstrukcja z dwóch filtracji przeciwnych;
w każdym stopniu używamy rozkładu z kolejnych
ilorazów filtracji, a nie dowolnego rozszczepienia.
\end{proof}

\begin{twierdzenie}[Rozkład Hodge'a dla rozmaitości zwartej Kählerowskiej]
\label{thm:kahler-hodge-xiv}
Niech $M$ będzie zwartą rozmaitością Kählerowską
wymiaru zespolonego $d$. Dla każdego $k\ge0$
istnieje naturalny rozkład
\[
H^k(M;\C)=\bigoplus_{p+q=k}H^{p,q}(M),
\qquad H^{p,q}(M)\cong H^q(M,\Omega_M^p),
\]
gdzie $\Omega_M^p$ jest snopem holomorficznych $p$-form.
Ponadto $\overline{H^{p,q}(M)}=H^{q,p}(M)$.
Jeżeli $M$ jest gładką zespoloną odmianą projektową,
to $H^k(M;\Q)$ z tym rozkładem jest czystą
strukturą Hodge'a wagi $k$.
\end{twierdzenie}
\begin{proof}[Dowód z jawnie wskazanymi twierdzeniami analitycznymi]
Na rozmaitości zespolonej formy zespolone rozkładają
się według typu: $\Omega^k(M;\C)=
\bigoplus_{p+q=k}\Omega^{p,q}(M)$.
Operator zewnętrzny ma postać
$d=\partial+\bar\partial$,
gdzie $\partial$ i $\bar\partial$ zwiększają
odpowiednio pierwszy i drugi stopień.

Z twierdzenia Hodge'a dla form różniczkowych
każda klasa de Rhama ma dokładnie jednego
reprezentanta harmonicznego dla ustalonej metryki.
Z tożsamości Kählerowskich (przyjęty tu wynik
analityczny) otrzymujemy
$\Delta_d=2\Delta_{\bar\partial}$.
Laplasjan $\Delta_{\bar\partial}$ zachowuje typ
$(p,q)$, więc rozkład form harmonicznych według
typu daje rozkład kohomologii de Rhama.
Twierdzenie Dolbeaulta identyfikuje przestrzeń
form $\bar\partial$-harmonicznych typu $(p,q)$
z $H^q(M,\Omega_M^p)$.
Sprzężenie form harmonicznych zamienia typy
$(p,q)$ i $(q,p)$. Dla rozmaitości projektowej
istnienie metryki Kählerowskiej jest klasycznym
wynikiem geometrycznym. Zatem wszystkie warunki
definicji struktury Hodge'a są spełnione.
\end{proof}

\begin{przyklad}[Prosta rzutowa i krzywa eliptyczna]
Na $\mathbb P^1$ niezerowe grupy Hodge'a to
$H^{0,0}(\mathbb P^1)\cong\C$ oraz
$H^{1,1}(\mathbb P^1)\cong\C$; w szczególności
$H^1(\mathbb P^1;\C)=0$.
Jeżeli $E=\C/\Lambda$ jest krzywą eliptyczną,
to $H^1(E;\C)=H^{1,0}(E)\oplus H^{0,1}(E)$,
oba składniki są jednowymiarowe: reprezentują je
odpowiednio klasy form $\dd z$ i $\dd\bar z$.
W stopniu drugim mamy $H^{1,1}(E)\cong\C$.
Stąd $h^{p,q}=h^{q,p}$, a $b_1(E)=2$.
\end{przyklad}

\begin{wniosek}[Parzystość nieparzystych liczb Bettiego]
Jeżeli $M$ jest zwartą rozmaitością Kählerowską,
to $b_{2r+1}(M)$ jest parzyste dla każdego $r\ge0$.
\end{wniosek}
\begin{proof}
Dla $p+q=2r+1$ nie zachodzi $p=q$, więc
składniki $H^{p,q}$ i $H^{q,p}$ tworzą rozłączne
pary tej samej liczby wymiarów. Sumowanie po
wszystkich parach daje parzysty wymiar $H^{2r+1}(M;\C)$.
\end{proof}

\section{Ciąg spektralny Hodge'a--de Rhama i jego zakres}
Niech $X$ będzie gładką zespoloną odmianą algebraiczną.
Filtracja ``głupia'' $F^p\Omega_X^\bullet=
\Omega_X^{\ge p}$ prowadzi do ciągu spektralnego
\[
E_1^{p,q}=H^q(X,\Omega_X^p)
\Longrightarrow\mathbb H^{p+q}(X,\Omega_X^\bullet).
\]
W przypadku gładkim projektowym jest to algebraiczna
wersja ciągu Hodge'a--de Rhama, a jego zbieżna
hiperkohomologia odpowiada kohomologii singularnej
$H^{p+q}(X^{\mathrm{an}};\C)$ po wykorzystaniu
twierdzenia porównawczego Grothendiecka.
\begin{twierdzenie}[Degeneracja na stronie $E_1$]
Niech $X$ będzie gładką właściwą odmianą algebraiczną
nad $\C$. Wtedy ciąg spektralny Hodge'a--de Rhama
degeneruje na stronie $E_1$.
\end{twierdzenie}
\begin{proof}[Zakres]
Dla odmian projektowych nad $\C$ jest to konsekwencja
rozkładu Hodge'a i zgodności filtracji Dolbeaulta.
W ogólnym gładkim właściwym przypadku stosuje się
odpowiednie twierdzenie analityczne lub algebraiczne.
Nie utożsamiamy tej degeneracji z dowolnym
ciągiem spektralnym dla przestrzeni osobliwych.
\end{proof}

\section{Mieszane struktury Hodge'a i osobliwa krzywa węzłowa}
\begin{definicja}[Mieszana struktura Hodge'a]
Mieszaną strukturą Hodge'a na skończeniewymiarowym
$V_{\Q}$ nazywamy rosnącą, skończoną filtrację
$W_\bullet V_{\Q}$ (filtrację wag) oraz malejącą,
skończoną filtrację $F^\bullet V_{\C}$ takie, że
na $\operatorname{Gr}_m^WV=W_mV/W_{m-1}V$
filtracja indukowana przez $F$ określa czystą
strukturę Hodge'a wagi $m$ dla każdego $m$.
\end{definicja}
\begin{twierdzenie}[Deligne: zakres użycia]
Dla dowolnej zespolonej odmiany algebraicznej $X$
skończonego typu grupy $H^k(X;\Q)$ są naturalnie
wyposażone w mieszane struktury Hodge'a, funktorialne
dla morfizmów algebraicznych. Dla $X$ gładkiej
projektowej struktura na $H^k$ jest czysta wagi $k$.
\end{twierdzenie}
\begin{proof}[Status]
To głębokie twierdzenie Deligne'a. Jego dowód
wymaga zwartyfikacji, rozdzielczości osobliwości
oraz mieszanych kompleksów Hodge'a. W niniejszym
rozdziale wyprowadzamy natomiast konkretny rachunek
wag w przykładzie, przyjmując funktorialność
struktur mieszanych i ścisłość ich morfizmów.
\end{proof}

\begin{przyklad}[Waga zerowa w $H^1$ krzywej z węzłem]
\label{ex:nodal-mhs-xiv}
Niech $C$ będzie nieredukowalną rzutową krzywą
wymierną z jednym zwyczajnym węzłem, a
$\nu:\widetilde C\cong\mathbb P^1\to C$ jej normalizacją.
Preobraz węzła składa się z dwóch punktów $a,b$.
Istnieje krótki ciąg snopów stałych na $C$:
\[
0\to\Q_C\to\nu_*\Q_{\widetilde C}
\xrightarrow{\,\mathrm{ev}_a-\mathrm{ev}_b\,}
\Q_{\{x\}}\to0.
\]
Dokładność w $x$ wynika z mapy $\Q\to\Q^2$
przyporządkowującej stałej $c$ parę $(c,c)$,
której kokernelem jest $\Q$ przez różnicę wartości.
Poza węzłem $\nu$ jest izomorfizmem.
W długim ciągu kohomologii $H^0(C;\Q)\to
H^0(\widetilde C;\Q)$ jest izomorfizmem,
a różnica wartości globalnej funkcji stałej
wynosi zero. Dlatego
$H^1(C;\Q)\cong\Q$,
$H^2(C;\Q)\cong\Q$.

Z funktorialności mieszanego Hodge'a otrzymujemy,
że $H^1(C;\Q)$ jest izomorficzne z czystą
strukturą $H^0(\{x\};\Q)$ wagi $0$.
W szczególności $H^1(C;\Q)$ ma wagę $0$, a nie $1$.
Tymczasem kompleks przecięcia normalizuje tę
anomalię w innym sensie:
$IH^1(C;\Q)=H^1(\mathbb P^1;\Q)=0$,
co obliczyliśmy już w rozdziale~\ref{ch:ih-perverse-xiii}.
To są dwie różne konstrukcje, a nie dwie nazwy
tego samego niezmiennika.
\end{przyklad}

\section{Rodziny, lokalne systemy i monodromia}
\begin{definicja}[Monodromia rodziny]
Niech $p:E\to B$ będzie lokalnie trywialną wiązką
z włóknem $F$, $b\in B$. Transport po pętlach
w $B$ wyznacza (po wybraniu identyfikacji włókna)
reprezentację
$\rho_q:\pi_1(B,b)\to\operatorname{Aut}(H^q(F;\Q))$,
zwaną reprezentacją monodromii kohomologicznej.
\end{definicja}
\begin{twierdzenie}[Ehresmann: wersja używana]
Jeżeli $p:E\to B$ jest surjektywną, właściwą gładką
submersją między rozmaitościami gładkimi bez brzegu,
to $p$ jest lokalnie trywialną wiązką gładką.
W szczególności $R^qp_*\Q_E$ jest lokalnym
systemem o włóknie $H^q(F;\Q)$.
\end{twierdzenie}
\begin{proof}[Szkic z precyzyjnym użyciem właściwości]
Metryka na $E$ pozwala wybrać gładki rozkład
$TE=\ker(dp)\oplus H$. Dla pola wektorowego
na małym dysku w $B$ podnosimy je poziomo do $H$.
Właściwość mapy zapewnia, że nad odpowiednio
małym zwartym otoczeniem podstawy trajektorie
podniesionych pól nie uciekają w skończonym czasie.
Przepływy budują dyfeomorfizmy lokalnych włókien.
Pełny dowód uwzględnia zgodność przy zmianie pól;
przyjmujemy tutaj standardową wersję twierdzenia
Ehresmanna.
\end{proof}

\begin{center}
\begin{tikzpicture}[scale=0.85,>=Stealth,font=\small]
\draw[manifoldblue,thick] (0,0) circle (1.1);
\fill[manifoldred] (0,0) circle (2.1pt);
\node[below right] at (0,0) {$0$};
\draw[manifoldgreen,very thick,->] (1.1,0)
 arc[start angle=0,end angle=305,radius=1.1];
\node[right] at (1.25,0.75) {$\gamma$};
\node at (0,-1.42) {$\Delta^*$};
\draw[manifoldblue,thick] (4.1,0) ellipse (0.95 and 0.44);
\draw[manifoldblue,thick] (4.1,0.2) ellipse (0.95 and 0.44);
\node[above] at (4.1,1.0) {$F_t$};
\draw[->,thick] (1.5,0)--(2.8,0)
 node[midway,above] {transport};
\node at (4.1,-1.42) {$T$};
\end{tikzpicture}
\par\medskip\small Transport wzdłuż zamkniętej pętli w podstawie
wyznacza działanie monodromii na kohomologii włókna.
\end{center}

\begin{uwaga}[Punkt krytyczny]
Dla $f:X\to\Delta$, gdzie $\Delta\subset\C$ jest dyskiem,
submersja może zawieść nad $0$. Wtedy
$R^qf_*\Q$ może nie być lokalnym systemem na całym dysku,
choć jego ograniczenie do $\Delta^*=\Delta\setminus\{0\}$
może być lokalnie stałe. Monodromia pętli otaczającej
$0$ opisuje część informacji traconej przy specjalizacji.
\end{uwaga}

\section{Włókna Milnora, cykle pobliskie i zanikające}
\begin{definicja}[Włókno Milnora]
Niech $f:(\C^{n+1},0)\to(\C,0)$ będzie holomorficznym
germem, którego jedynym punktem krytycznym w
pewnym otoczeniu jest $0$, i niech $f(0)=0$.
Dla dostatecznie małych $0<|t|\ll\delta$ definiujemy
lokalne włókno Milnora
$F_t=f^{-1}(t)\cap B_\delta^{2n+2}$.
Z twierdzenia Milnora wynika, że dla odpowiednio
wybranych $t,\delta$ jego typ homotopijny nie
zależy od parametrów, a rodzina nad małym
przebitym dyskiem jest lokalnie trywialna.
\end{definicja}

\begin{definicja}[Cykle pobliskie i zanikające]
\label{def:nearby-vanishing-xiv}
Niech $f:X\to\Delta$ będzie analitycznym odwzorowaniem,
$X_0=f^{-1}(0)$, $i:X_0\hookrightarrow X$, a
$\widetilde{\Delta^*}\to\Delta^*$ uniwersalnym nakryciem.
Niech $\widetilde X^*=X\times_\Delta\widetilde{\Delta^*}$
z odwzorowaniem $\widetilde j:\widetilde X^*\to X$.
Dla $K\in D_c^b(X;\Q)$ definiujemy
\emph{kompleks cykli pobliskich}
$R\psi_f K=i^*R\widetilde j_*\widetilde j^*K$.
Istnieje naturalny morfizm specjalizacji
$\mathrm{sp}:i^*K\to R\psi_fK$.
\emph{Kompleks cykli zanikających} definiujemy jako
$R\phi_fK=\operatorname{Cone}(\mathrm{sp})$.
Wtedy istnieje trójkąt wyróżniony
\[
i^*K\longrightarrow R\psi_fK
\longrightarrow R\phi_fK\longrightarrow i^*K[1].
\]
Przy tej konwencji nie wprowadzamy dodatkowego
przesunięcia $[-1]$ do definicji $R\phi_f$.
\end{definicja}
\begin{uwaga}[Precyzja konstrukcji]
W definicji cykli pobliskich $\widetilde j$
nie jest zwykłym włożeniem otwartym:
uwzględnia uniwersalne nakrycie $\Delta^*$.
Działanie grupy jego automorfizmów
$\pi_1(\Delta^*)\cong\Z$ daje naturalny
operator monodromii $T$ na $R\psi_fK$ i $R\phi_fK$.
Otrzymanie tych funktorów w kategorii
konstruktybilnych kompleksów oraz zgodność
z wybranym nakryciem jest standardowym twierdzeniem
teorii snopów, z którego tu korzystamy.
\end{uwaga}

\begin{stwierdzenie}[Włókno Milnora i cykle zanikające]
\label{prop:milnor-stalk-xiv}
Jeśli $X$ jest małą kulą w $\C^{n+1}$,
$f:X\to\Delta$ ma izolowany punkt krytyczny $0$,
a $K=\Q_X$, to
$H^q((R\psi_fK)_0)\cong H^q(F_t;\Q)$.
Ponieważ $X_0\cap B_\delta$ jest lokalnie
ściągalne do $0$, dla $q>0$ zachodzi
$H^q((R\phi_fK)_0)\cong\widetilde H^q(F_t;\Q)$;
w stopniu $0$ analogiczne stwierdzenie dotyczy
kohomologii zredukowanej, o ile $F_t$ jest niepuste.
\end{stwierdzenie}
\begin{proof}
Pierwszy izomorfizm jest twierdzeniem o kiełkach
cykli pobliskich; używa lokalnej trywialności
fibracji Milnora i właściwych warunków
konstruktybilności. Następnie bierzemy
ciąg długi trójkąta z
Definicji~\ref{def:nearby-vanishing-xiv}.
Kiełek $i^*\Q_X$ w $0$ ma kohomologię $\Q$
w stopniu $0$ i zero w pozostałych stopniach.
Morfizm $H^0(i^*K)\to H^0(R\psi_fK)$
wysyła $1$ na funkcję stałą $1$ na każdej
składowej włókna. Dokładność daje właśnie
kohomologię zredukowaną włókna Milnora,
z uwzględnieniem przesunięć w trójkącie.
\end{proof}

\section{Dwa jawne rachunki monodromii}
\begin{przyklad}[Węzeł płaski $f(z,w)=zw$]
\label{ex:node-monodromy-xiv}
Dla $f:\C^2\to\C$, $f(z,w)=zw$, włókno
$F_t$ dla $t\ne0$ jest homotopijnie równoważne
okręgowi $S^1$: z równania $w=t/z$ wynika,
że w małej kuli $z$ przebiega pierścień.
Stąd $\widetilde H^1(F_t;\Q)\cong\Q$
i inne zredukowane grupy znikają.
Zatem jedyną niezerową kohomologią
$(R\phi_f\Q_{\C^2})_0$ jest $H^1\cong\Q$.

Monodromię można zobaczyć bezpośrednio:
przy przejściu $t\mapsto te^{i\theta}$
transport $(z,w)\mapsto(e^{i\theta}z,w)$
zachowuje równanie włókna i normę modułu $z$;
po pełnym obrocie $\theta=2\pi$ jest identycznością.
Indukowana monodromia na $H^1(F_t;\Q)$ to więc $T=1$.
Ten \emph{lokalny} wynik nie oznacza, że monodromia
na całym $H^1$ gładkiej krzywej degenerującej
do krzywej węzłowej musi być identycznością:
może ona być niebanalną transwekcją.
\end{przyklad}

\begin{twierdzenie}[Osobliwość kwadratowa $A_1$: pełne obliczenie]
\label{thm:a1-monodromy-xiv}
Niech $n\ge1$ oraz
$f:\C^{n+1}\to\C$,
$f(z_0,\dots,z_n)=\sum_{j=0}^nz_j^2$.
Wówczas lokalne włókno Milnora ma typ homotopijny $S^n$,
a monodromia geometryczna wywołuje na
$\widetilde H^n(F_t;\Z)\cong\Z$ mnożenie przez
$(-1)^{n+1}$.
\end{twierdzenie}
\begin{proof}
Po zastąpieniu $t\ne0$ przez $1$ za pomocą
skalowania zespolonego zbadamy afiniczną
kwadrykę $Q=f^{-1}(1)$. Zapiszmy $z=x+iy$,
gdzie $x,y\in\R^{n+1}$. Równanie $f(z)=1$
jest równoważne układowi
$\lVert x\rVert^2-\lVert y\rVert^2=1$,
$\langle x,y\rangle=0$.
Dla $u\in S^n$ i $p\in T_u^*S^n$ utożsamionego
z wektorem $p\perp u$, wzór
\[
(u,p)\longmapsto\sqrt{1+\lVert p\rVert^2}\,u+ip
\]
jest dyfeomorfizmem $T^*S^n\to Q$.
Odwrotność ma postać
$u=x/\lVert x\rVert$, $p=y$;
równania kwadryki zapewniają $p\perp u$.
Skalowanie $p\mapsto sp$ dla $0\le s\le1$
przekształca się w deformacyjną retrakcję
$Q\to S^n$ (przekrój zerowy).
Dla włókna lokalnego po przeskalowaniu $t$
uzyskujemy część kwadryki $Q$ ograniczoną
odpowiednio dużą kulą. Retrakcja zmniejsza
normę $\lVert z\rVert$, zatem pozostaje w tej
części i wykazuje jej ten sam typ homotopijny.

Ponieważ $f$ jest jednorodny stopnia $2$,
pętli $t(s)=e^{2\pi is}$ odpowiada transport
$z\mapsto e^{\pi is}z$, a po pełnym okrążeniu
otrzymujemy $z\mapsto-z$. Na retrakcie
$S^n\subset Q$ jest to odwzorowanie
antypodyczne stopnia $(-1)^{n+1}$.
Ponieważ retrakcja jest równoważnością
homotopijną, taki jest również operator
monodromii na $\widetilde H^n(F_t;\Z)$.
\end{proof}
\begin{przyklad}[Zmiana parzystości]
Dla $n=1$, odpowiadającego zwykłemu węzłowi
płaskiemu, monodromia działa na $H^1(S^1;\Z)$
jak $+1$. Dla $n=2$ działa na
$H^2(S^2;\Z)$ jak $-1$. Zmiana znaku wynika z
parzystości wymiaru sfery, a nie ze zmiany
liczby cykli zanikających.
\end{przyklad}

\begin{uwaga}[Różne znaczenia monodromii przy węźle]
Są tu trzy różne obiekty: monodromia lokalnego włókna Milnora,
monodromia całego $H^1$ gładkiej krzywej oraz monodromia
po zmianie parametru rodziny. Pierwsza dla $zw=t$ działa
identycznie na $H^1$ pierścienia; druga może być
nietrywialną transwekcją; po podstawieniu $t=s^m$
trzecia jest $m$-tą potęgą monodromii pierwotnej.
Nie wolno więc wnioskować z lokalnej identyczności, że
monodromia globalnej rodziny jest trywialna.
\end{uwaga}

\section{Formuła Picarda--Lefschetza i jej zakres}
\begin{twierdzenie}[Picard--Lefschetz dla rodzin krzywych]
\label{thm:picard-lefschetz-xiv}
Niech $\pi:\mathcal C\to\Delta$ będzie właściwą
holomorficzną rodziną, której włókna nad
$\Delta^*$ są gładkimi zwartymi spójnymi
krzywymi zespolonymi, a włókno nad $0$ posiada
jedną zwyczajną osobliwość węzłową.
Niech $\delta\in H_1(C_t;\Z)$ będzie
zorientowanym cyklem zanikającym.
Po ustaleniu orientacji pętli i konwencji
sparowania przecięć monodromia na $H_1(C_t;\Z)$
ma postać transwekcji
\[
T(\gamma)=\gamma+\langle\gamma,\delta\rangle\,\delta.
\]
Znak przy ostatnim składniku zależy od
przyjętej konwencji orientacji.
\end{twierdzenie}
\begin{proof}[Status i rachunek lokalnego modelu]
Pełna identyfikacja monodromii globalnej z
skręceniem Dehna wokół $\delta$ jest twierdzeniem
Picarda--Lefschetza. Po jego przyjęciu wzór na
homologię wynika z policzenia wpływu skręcenia
na krzywą $\gamma$: każde poprzeczne przecięcie
z $\delta$ dokłada jeden obieg wokół $\delta$,
a znaki tych obiegów sumują się do
$\langle\gamma,\delta\rangle$.
Nie sprowadzamy tego globalnego rachunku do
stwierdzenia o monodromii lokalnego pierścienia
$\{zw=t\}$, gdzie działanie na $H_1$ jest identyczne.
\end{proof}
\begin{przyklad}[Macierz monodromii]
Załóżmy, że $C_t$ ma rodzaj $1$ oraz
$H_1(C_t;\Z)=\Z a\oplus\Z b$, przy czym
$\delta=a$ i $\langle b,a\rangle=1$.
Wówczas $T(a)=a$, $T(b)=b+a$.
W bazie $(a,b)$ macierz monodromii jest
$\bigl(\begin{smallmatrix}1&1\\0&1\end{smallmatrix}\bigr)$.
Jej nilpotentna część $N=T-I$ spełnia
$N\ne0$ oraz $N^2=0$; ilustruje to różnicę
między monodromią unipotentną a identyczną.
\end{przyklad}

\section{Powrót do snopów perwersyjnych i ograniczenia teorii}
\begin{twierdzenie}[Dokładność perwersyjna cykli]
Dla algebraicznego odwzorowania $f:X\to\C$ i
konstruktybilnego kompleksu perwersyjnego $P$
funktory $R\psi_f[-1]$ i $R\phi_f[-1]$
są $t$-dokładne dla środkowych perwersyjnych
struktur $t$ (przy standardowej konwencji
funkcji cykli i odpowiedniej kategorii
konstruktybilnych kompleksów).
\end{twierdzenie}
\begin{proof}[Status]
Wynik wymaga pełnej teorii funktorów cykli
pobliskich i zanikających oraz sklejenia
perwersyjnych struktur $t$. Przyjmujemy go
jako twierdzenie; powyższe rachunki dla
$A_1$ same w sobie nie dają jego ogólnego dowodu.
\end{proof}
\begin{uwaga}[Zastosowania i dalsza granica]
Cykle zanikające rejestrują topologię zmieniającą
się podczas degeneracji; w szczególności w
wymiarze zespolonym $n$ dla izolowanych punktów
krytycznych są skupione w wymiarze środkowym
(co wynika z twierdzenia Milnora o bukiecie sfer).
Pełny formalizm mieszanych modułów Hodge'a,
teoremat o quasi-unipotentności monodromii
oraz dowody twierdzenia o rozkładzie przez
czystość wag pozostają poza zakresem tego rozdziału.
Użycie słów ``twierdzenie'' i ``dowód'' w
poprzednich rozdziałach nie usuwa tych zależności.
\end{uwaga}

\section{Kontrola zależności i przygotowanie do geometrii algebraicznej}
W tym rozdziale bez odwoływania się do głębokich twierdzeń
dowiedliśmy tożsamości
$\dd^2=0$ dla stożka, jego długiego ciągu,
równoważności ``quasi-izomorfizm $\Leftrightarrow$
stożek acykliczny'', rachunku filtracji
czystej struktury Hodge'a oraz typu homotopijnego
i monodromii osobliwości kwadratowej. Pozostałe
twierdzenia oznaczyliśmy jako użyte narzędzia.
Przy scalaniu z geometrią algebraiczną
należy współdzielić definicje snopów, hiper-
kohomologii i funktorów pochodnych, a nie
powielać je w obu skryptach.

\chapter{Wariacje struktur Hodge'a, koneksja Gaussa--Manina i degeneracje}
\label{ch:xv-vhs-degeneracje}

\section{Dlaczego struktura Hodge'a musi się zmieniać?}
W rozdziale~\ref{ch:xiv-hodge-monodromia} wyposażyliśmy
$H^k(X;\Q)$ gładkiej odmiany projektowej $X$ w strukturę Hodge'a
oraz obliczyliśmy monodromię izolowanej osobliwości kwadratowej.
Teraz rozważymy \emph{rodzinę} $f:\mathcal X\to S$ odmian i zapytamy,
jak przy zmianie parametru $s\in S$ porównywać podprzestrzenie
$H^{p,q}(X_s)\subset H^k(X_s;\C)$.
Grupy kohomologii są lokalnie identyczne jako przestrzenie wektorowe,
lecz filtracja Hodge'a może się zmieniać. To rozróżnienie prowadzi
do koneksji Gaussa--Manina, równań Picarda--Fuchsa i granicznych
struktur mieszanych.

\begin{definicja}[Gładka właściwa rodzina projektowa]
Niech $S$ będzie spójną gładką odmianą algebraiczną nad $\C$
(rozpatrywaną także jako rozmaitość analityczna),
$\mathcal X$ gładką odmianą algebraiczną nad $\C$,
a $f:\mathcal X\to S$ gładkim, surjektywnym morfizmem
projektowym. Odpowiadające mu odwzorowanie analityczne
jest właściwą submersją holomorficzną. Włóknem w $s\in S$ nazywamy
$X_s=f^{-1}(s)$. Każde $X_s$ jest wówczas zwartą gładką
odmianą projektową. W tym rozdziale właśnie taką rodzinę
nazywamy krótko \emph{gładką rodziną projektową}.
\end{definicja}
\begin{uwaga}[Po co wszystkie hipotezy?]
Właściwość i warunek submersji pozwalają zastosować twierdzenie
Ehresmanna, a projektowość zapewnia rozkład Hodge'a oraz
jego polaryzację. Na bazie $S$ nie musi istnieć globalny wybór
baz kohomologii włókien: przeszkodą jest monodromia.
\end{uwaga}

\section{Lokalny system, płaska wiązka i koneksja Gaussa--Manina}
\begin{definicja}[Lokalny system kohomologii]
Niech $f:\mathcal X\to S$ będzie gładką rodziną projektową,
$k\ge0$. Oznaczmy przez
$\mathbb V_{\Z}:=R^kf_*\Z_{\mathcal X}/\mathrm{tors}$
lokalny system wolnych grup abelowych skończonego rzędu,
przez $\mathbb V_{\C}=\mathbb V_{\Z}\otimes_{\Z}\C$
jego zespolenie, a przez
$\mathcal H:=\mathbb V_{\C}\otimes_{\C_S}\mathcal O_S$
skończonego rzędu holomorficzną wiązkę wektorową.
Dokładniej, iloraz przez torsję bierze się włókno po włóknie;
z twierdzenia Ehresmanna otrzymujemy lokalny system
$\mathbb V_{\Z}$ o włóknie $H^k(X_s;\Z)/\mathrm{tors}$.
\end{definicja}
\begin{stwierdzenie}[Płaska koneksja z lokalnego systemu]
\label{prop:xv-gm-local}
Dla każdego zespolonego lokalnego systemu $\mathbb V$ na
zespolonej rozmaitości $S$ wiązka
$\mathcal H=\mathbb V\otimes_{\C_S}\mathcal O_S$
ma kanoniczną koneksję holomorficzną
$\nabla:\mathcal H\to\Omega_S^1\otimes_{\mathcal O_S}\mathcal H$
charakteryzowaną przez $\nabla(v\otimes g)=v\otimes\dd g$
dla lokalnie stałego $v$. Jest ona płaska i jej lokalnymi
przekrojami poziomymi są dokładnie przekroje $\mathbb V$.
\end{stwierdzenie}
\begin{proof}
Wybierzmy pokrycie $S$ zbiorami po prostu spójnymi $U_\alpha$,
na których $\mathbb V$ jest stały, i lokalne bazy
$e_{\alpha,1},\dots,e_{\alpha,r}$. Dla
$\sigma=\sum_ig_ie_{\alpha,i}$ przyjmujemy
$\nabla\sigma=\sum_i\dd g_i\otimes e_{\alpha,i}$.
Na przecięciu $U_\alpha\cap U_\beta$ macierz przejścia
między bazami ma lokalnie stałe współczynniki, a więc
jej różniczka znika. Zatem wzór nie zależy od wyboru bazy.
Reguła Leibniza wynika z $\dd(hg)=h\dd g+g\dd h$,
natomiast $\nabla^2=0$ z $\dd^2=0$ w płaskich bazach.
Warunek $\nabla\sigma=0$ oznacza $\dd g_i=0$ dla każdego $i$,
czyli lokalną stałość wszystkich $g_i$.
\end{proof}
\begin{definicja}[Koneksja Gaussa--Manina]
Koneksję z poprzedniego stwierdzenia, zastosowaną do
$\mathbb V_{\C}=R^kf_*\C_{\mathcal X}$, nazywamy
\emph{koneksją Gaussa--Manina} rodziny $f$ w stopniu $k$.
\end{definicja}
\begin{stwierdzenie}[Różniczkowanie okresów]
\label{prop:xv-period-derivative}
Niech $U\subset S$ będzie dyskiem współrzędnych z parametrem $t$,
$\gamma_t\in H_k(X_t;\Z)$ będzie rodziną cykli poziomą
w sensie lokalnego systemu homologii, a
$\omega(t)\in\mathcal H|_U$ różniczkowalną sekcją.
Wówczas dla funkcji okresowej
$\Pi(t)=\langle\gamma_t,\omega(t)\rangle$
zachodzi
$\Pi'(t)=\langle\gamma_t,\nabla_{\partial/\partial t}\omega(t)\rangle$.
\end{stwierdzenie}
\begin{proof}
Na $U$ możemy jednocześnie wybrać płaskie bazy homologii
oraz kohomologii. W tych bazach $\gamma_t$ ma stałe
współrzędne, a $\omega(t)$ ma współrzędne zmienne.
Różniczkujemy zwykłą skończoną sumę iloczynów współrzędnych.
Z definicji koneksji współczynniki pochodne są współrzędnymi
$\nabla_{\partial/\partial t}\omega$. Zależność od bazy
znika, ponieważ parowanie jest naturalne.
\end{proof}

\section{Filtracja Hodge'a i warunek transwersalności Griffithsa}
\begin{definicja}[Wariacja czystych struktur Hodge'a]
Niech $m\in\Z$ i $S$ będzie zespoloną rozmaitością.
\emph{Całkowitą wariacją czystych struktur Hodge'a wagi $m$}
nazywamy lokalny system wolnych grup abelowych
$\mathbb V_{\Z}$ skończonego rzędu oraz malejącą
filtrację $F^\bullet\mathcal H$ przez holomorficzne
podwiązki $\mathcal H=\mathbb V_{\Z}\otimes_{\Z}\mathcal O_S$
o stałych rzędach, spełniającą:
\begin{enumerate}[label=\textup{(\roman*)},leftmargin=*]
\item dla każdego $s\in S$ filtracja na
$\mathbb V_{\Z,s}\otimes\C$ wyznacza czystą strukturę Hodge'a wagi $m$;
\item dla płaskiej koneksji $\nabla$ zachodzi
$\nabla(F^p\mathcal H)\subseteq
F^{p-1}\mathcal H\otimes\Omega_S^1$ dla każdego $p\in\Z$.
\end{enumerate}
Ostatni warunek nazywamy \emph{transwersalnością Griffithsa}.
Polaryzowana wariacja zawiera ponadto płaskie, całkowite
parowanie polaryzujące na wszystkich włóknach.
\end{definicja}
\begin{uwaga}
Warunek transwersalności nie stwierdza, że
$\nabla F^p\subset F^p$. To ostatnie oznaczałoby
stałość podwiązki Hodge'a względem transportu i jest
ogólnie fałszywe. Ponadto samo pojęcie lokalnego systemu
nie wystarcza do zdefiniowania wariacji Hodge'a: potrzeba
holomorficznej filtracji i jej zgodności z koneksją.
\end{uwaga}
\begin{twierdzenie}[Wariacja Hodge'a z rodziny projektowej]
\label{thm:xv-vhs-family}
Dla gładkiej rodziny projektowej $f:\mathcal X\to S$
i każdego $k\ge0$ lokalny system
$R^kf_*\Z_{\mathcal X}/\mathrm{tors}$ wraz z filtracją
pochodzącą od względnego kompleksu de Rhama
$\Omega^\bullet_{\mathcal X/S}$ tworzy polaryzowaną
wariację czystych struktur Hodge'a wagi $k$.
\end{twierdzenie}
\begin{proof}[Wyodrębnienie części elementarnej i wyniku używanego]
Z twierdzenia Ehresmanna otrzymujemy lokalny system,
a ze Stwierdzenia~\ref{prop:xv-gm-local} płaską koneksję.
Względna filtracja form przez ich stopień holomorficzny
wyznacza podprzestrzenie $F^pH^k(X_s;\C)$.
Stałość ich wymiarów i fakt, że tworzą podwiązki holomorficzne,
wynikają z teorii kohomologii koherentnej i zmiany podstawy
dla gładkich rodzin projektowych. W każdym włóknie
rozkład Hodge'a daje wymagane czyste struktury.

Pozostaje warunek (ii). Lokalny dowód jego postaci
na poziomie kompleksów można uzyskać tak:
wybieramy holomorficzne współrzędne $(t,z_1,\dots,z_d)$
w otoczeniu punktu całkowitej przestrzeni, w których
$f(t,z)=t$. W różniczce absolutnej reprezentanta
względnej formy typu $p$ składniki zawierające $\dd t$
powstają przez różniczkowanie współczynników i przez
porównanie lokalnych podniesień form względnych.
Po wydzieleniu czynnika $\dd t$ mają stopień
względny co najmniej $p-1$. Ten lokalny rachunek dowodzi
właściwego ograniczenia stopnia, lecz jeszcze nie pokazuje,
że otrzymana koneksja jest niezależna od reprezentanta klasy.
Do tego potrzebna jest zgodność różniczki absolutnej
z filtracją względnego kompleksu de Rhama i przejście
przez hiperkohomologię; jest to nietrywialne twierdzenie
porównawcze. Po jego zastosowaniu otrzymuje się
$\nabla F^p\subseteq F^{p-1}\otimes\Omega_S^1$.
Niezależność od wybranych podniesień i sklejanie
tego rachunku na kompleksie \v{C}echa--de Rhama
jest standardowym twierdzeniem o konstrukcji koneksji
Gaussa--Manina; korzystamy z niego w tym miejscu.
Polaryzacja pochodzi z formy przecięcia i klasy
hiperpłaszczyzny przez twierdzenie Hodge'a--Riemanna.
Nie przedstawiamy tego ostatniego jako wyniku
dowiedzionego wyłącznie powyższym rachunkiem.
\end{proof}

\section{Okresy krzywej eliptycznej i półpłaszczyzna okresów}
\begin{stwierdzenie}[Parametr okresów]
Niech $E$ będzie gładką krzywą eliptyczną, $a,b$ zorientowaną
bazą $H_1(E;\Z)$ o $\langle a,b\rangle=1$, a
$0\ne\omega\in H^{1,0}(E)$. Okresy
$A=\int_a\omega$, $B=\int_b\omega$ spełniają $A\ne0$ i
$\tau=B/A\in\mathbb H=\{z\in\C:\operatorname{Im}z>0\}$
przy zgodnym wyborze orientacji formy przecięcia.
\end{stwierdzenie}
\begin{proof}
Gdyby $A=0$, wszystkie okresy formy $\omega$ musiałyby
być rzeczywiście wielokrotnościami $B$, co przeczy
jej dodatniemu polaryzowanemu parowaniu.
Ściślej, relacja Riemanna dla parowania na $H^1(E)$ daje
$\mathrm i\int_E\omega\wedge\overline\omega
=\mathrm i(A\overline B-B\overline A)=2\operatorname{Im}(\overline A B)>0$.
Lewa strona jest dodatnia, ponieważ w lokalnej
współrzędnej $\omega=h(z)\,\dd z$ i
$\mathrm i\,\omega\wedge\overline\omega$
jest dodatnią formą objętościową (z przyjętą orientacją).
Stąd $\operatorname{Im}(\overline AB)>0$, a więc $A\ne0$
i $\operatorname{Im}(B/A)>0$.
\end{proof}
\begin{uwaga}[Zmiana bazy cykli]
Po zamianie bazy $(a,b)$ na inną dodatnio zorientowaną
bazę przez macierz z $\mathrm{SL}_2(\Z)$ parametr
$\tau$ zmienia się ułamkowo liniowo. Nie jest więc
funkcją jednowartościową na dowolnej bazie parametrów:
po obejściu pętli może ulec działaniu monodromii.
\end{uwaga}

\section{Równania Picarda--Fuchsa: od płaskiej koneksji do ODE}
\begin{stwierdzenie}[Równanie spełniane przez okresy]
\label{prop:xv-pf-general}
Niech $S$ będzie gładką nierozkładalną krzywą algebraiczną,
$f:\mathcal X\to S$ gładką rodziną projektową,
$\mathcal H$ algebraiczną wiązką de Rhama rzędu $r$
z koneksją Gaussa--Manina, a $\omega$ jej przekrojem
wymiernym. Ustalmy niestałą funkcję wymierną $t\in\C(S)$
i ograniczmy się do otwartego podzbioru, na którym $t$
jest regularna oraz $\dd t\neq0$. Wtedy $t$ jest
lokalną współrzędną analityczną, a różniczkowanie
$\partial_t$ zachowuje ciało $\C(S)$. Istnieją
funkcje wymierne $a_0,\dots,a_r\in\C(S)$, nie wszystkie zerowe,
takie że każda lokalna funkcja okresowa
$\Pi(t)=\int_{\gamma_t}\omega_t$ spełnia
$\sum_{j=0}^ra_j(t)\Pi^{(j)}(t)=0$
w dziedzinie, na której współczynniki są określone.
\end{stwierdzenie}
\begin{proof}
Nad ciałem funkcji wymiernych $K=\C(S)$ przestrzeń
generycznego włókna $\mathcal H\otimes K$ ma wymiar $r$.
Ponieważ $t\in K$ jest niestałe, różniczkowanie względem
$t$ jest derywacją ciała $K$ (na krzywej algebraicznej nad
$\C$ mamy $\Omega^1_{K/\C}=K\,\dd t$).
Algebraiczna koneksja przenosi tę derywację na $\mathcal H\otimes K$.
Wektory $\omega,\nabla_{\partial_t}\omega,\dots,
\nabla_{\partial_t}^{r}\omega$ są liniowo zależne nad $K$.
Bierzemy nietrywialną relację z $a_j\in K$.
Parujemy ją z lokalnie poziomym cyklem $\gamma_t$ i
korzystamy $r$ razy ze
Stwierdzenia~\ref{prop:xv-period-derivative}.
Daje to żądane równanie. Algebraiczność koneksji
i przekroju jest w założeniach istotna:
bez niej otrzymalibyśmy współczynniki meromorficzne,
a niekoniecznie wymierne na $S$.
\end{proof}

\section{Rodzina Legendre'a: rachunek okresu od początku}
\label{sec:xv-legendre}
Niech $S=\mathbb P^1\setminus\{0,1,\infty\}$ i rozważmy
rodzinę gładkich krzywych eliptycznych będących
uzupełnieniami rzutowymi krzywych afinicznych
\[
E_\lambda:\qquad y^2=x(x-1)(x-\lambda),\qquad \lambda\in S.
\]
Jej wyróżnik w postaci Weierstrassa wynosi
$\Delta(\lambda)=16\lambda^2(1-\lambda)^2$.
Rzeczywiście wyróżnik wielomianu po prawej stronie,
którego pierwiastki to $0,1,\lambda$, jest
$\lambda^2(1-\lambda)^2$, a wyróżnik równania
Weierstrassa $y^2=P_3(x)$ jest szesnastokrotnością tej liczby.
Forma $\omega_\lambda=\dd x/y$ jest holomorficzna i
nigdzie niezerowa na krzywej rzutowej $E_\lambda$.

\begin{stwierdzenie}[Rzeczywisty okres]
\label{prop:xv-legendre-period}
Dla $0<\lambda<1$ niech $a_\lambda$ będzie cyklem,
którego rzut na oś $x$ dwukrotnie przebiega odcinek
$[0,\lambda]$ z odpowiednio wybraną orientacją.
Wówczas okres ma postać
\[
\Pi_a(\lambda)=\int_{a_\lambda}\omega_\lambda
=2\int_0^\lambda\frac{\dd x}
{\sqrt{x(1-x)(\lambda-x)}}
=2\pi\sum_{n=0}^\infty
\frac{\binom{2n}{n}^2}{16^n}\lambda^n.
\]
Szereg zbiega bezwzględnie dla $|\lambda|<1$.
\end{stwierdzenie}
\begin{proof}
Na wnętrzu odcinka $[0,\lambda]$ prawa strona równania
krzywej jest dodatnia, a dwa arkusze nakrycia podwójnego
mają przeciwne znaki $y$. Całka po zamkniętym cyklu
jest więc dwa razy całką po jednym arkuszu.
Podstawienie $x=\lambda u$ daje
$\Pi_a(\lambda)=2\int_0^1
[u(1-u)(1-\lambda u)]^{-1/2}\,\dd u$.
Dla $|\lambda|<1$ rozwijamy
$(1-\lambda u)^{-1/2}=
\sum_{n\ge0}\binom{2n}{n}\lambda^nu^n/4^n$.
Na dowolnym zwartym poddysku $|\lambda|\le r<1$
szereg jest jednostajnie zdominowany przez
stałą pomnożoną przez $[u(1-u)]^{-1/2}$,
więc możemy całkować wyraz po wyrazie.
Z podstawienia $u=\sin^2\theta$ i klasycznej
rekurencji całek trygonometrycznych dostajemy
$\int_0^1 u^{n-1/2}(1-u)^{-1/2}\,\dd u
=\pi\binom{2n}{n}/4^n$.
Po pomnożeniu współczynników otrzymujemy tezę.
\end{proof}

\begin{twierdzenie}[Równanie Picarda--Fuchsa rodziny Legendre'a]
\label{thm:xv-legendre-pf}
Każdy lokalny okres formy $\omega_\lambda$ w rodzinie
$E_\lambda$ spełnia równanie
\[
\lambda(1-\lambda)\Pi''(\lambda)
+(1-2\lambda)\Pi'(\lambda)
-\frac14\Pi(\lambda)=0.
\]
\end{twierdzenie}
\begin{proof}
Niech $a_n=\binom{2n}{n}^2/16^n$.
Bezpośrednio z ilorazu kolejnych współczynników wynika
$(n+1)^2a_{n+1}=(n+\frac12)^2a_n$.
Jeśli $u(\lambda)=\sum_{n\ge0}a_n\lambda^n$,
to współczynnik przy $\lambda^n$ w
$\lambda(1-\lambda)u''+(1-2\lambda)u'-u/4$
jest równy
$(n+1)^2a_{n+1}-(n+\frac12)^2a_n=0$.
Wykazaliśmy więc równanie dla niezerowego okresu
$\Pi_a=2\pi u$. Dla ogólnego okresu odwołujemy się
do własności rzędu dwa koneksji Gaussa--Manina i
do związku równania hipergeometrycznego z rodziną
Legendre'a: drugi okres, dla cyklu nad odcinkiem $[\lambda,1]$,
po podstawieniu $x=1-(1-\lambda)s$ ma postać
$2\pi\mathrm i\,u(1-\lambda)$, z wyborem orientacji cyklu
odpowiadającym znakowi $+\mathrm i$.
Dokładniej, całka po drugim cyklu (z wybraną gałęzią
pierwiastka na $[\lambda,1]$) daje, po podstawieniu
$x=1-(1-\lambda)s$, właśnie funkcję
$2\pi\mathrm i\,u(1-\lambda)$, z dokładnością do
ustalonego znaku wynikającego z orientacji cyklu.
Ponieważ operator
$\lambda(1-\lambda)D^2+(1-2\lambda)D-1/4$
jest niezmienniczy względem podstawienia
$\lambda\mapsto1-\lambda$, obie niezależne
funkcje spełniają to równanie; ich niezależność
wynika z odmiennego zachowania przy $\lambda=0$
(stałość pierwszego okresu i logarytmiczna
rozbieżność drugiego, wykazana poniżej).
Każdy okres jest ich lokalną kombinacją liniową
ze stałymi współczynnikami.
\end{proof}
\begin{uwaga}[Kontrola normalizacji]
To jest równanie hipergeometryczne
${}_2F_1(\frac12,\frac12;1;\lambda)$.
Ważny jest znak członu zerowego: stoi przed nim
$-1/4$, a nie $+1/4$. Inna normalizacja okresu
może mnożyć rozwiązanie przez stałą, nie zmienia jednak
samego równania.
\end{uwaga}

\begin{stwierdzenie}[Logarytm i monodromia Legendre'a]
\label{prop:xv-legendre-T2}
Niech $u(\lambda)={}_2F_1(\frac12,\frac12;1;\lambda)$
oraz $v(\lambda)=\mathrm i\,u(1-\lambda)$, początkowo
na $0<\lambda<1$, z kontynuacją analityczną.
W sąsiedztwie $\lambda=0$ zachodzi
$u(\lambda)=1+O(\lambda)$ oraz
$u(1-\lambda)=-\pi^{-1}u(\lambda)\log\lambda+h(\lambda)$
z funkcją holomorficzną $h$ w otoczeniu zera.
Po dodatnim okrążeniu $\lambda=0$ otrzymujemy
$u\mapsto u$ oraz $v\mapsto v+2u$.
W szczególności w bazie rozwiązań $(u,v)$
monodromia działa na kolumnę $(u,v)^t$ przez macierz
$\left(\begin{smallmatrix}1&0\\2&1\end{smallmatrix}\right)$.
Działanie na bazie przestrzeni rozwiązań $(u,v)$
ma macierz transponowaną; przy porównywaniu z
monodromią homologii trzeba ponadto ustalić,
czy stosujemy transport cykli, czy jego odwrotność.
\end{stwierdzenie}
\begin{proof}
Równanie Picarda--Fuchsa ma w $\lambda=0$
regularny punkt osobliwy z podwójnym wykładnikiem
indycjalnym $0$. Metoda Frobeniusa daje drugie
rozwiązanie postaci $c\,u(\lambda)\log\lambda+h(\lambda)$,
gdzie $h$ jest holomorficzne.
Aby wyznaczyć $c$, korzystamy z całki
$u(1-\lambda)=\pi^{-1}\int_0^1
[s(1-s)(1-(1-\lambda)s)]^{-1/2}\,\dd s$.
Dla $s=1-z$ i małego $z$ mianownik ma część dominującą
$\sqrt{z(z+\lambda)}$. Całka
$\int_0^\varepsilon [z(z+\lambda)]^{-1/2}\,\dd z$
ma asymptotykę $-\log\lambda+O(1)$.
Stąd współczynnik wiodący części logarytmicznej wynosi
$c=-1/\pi$, ponieważ $u(0)=1$.
Tożsamość z całym członem $u(\lambda)\log\lambda$
(a nie tylko wyrazem $\log\lambda$) i holomorficzność
reszty $h$ pochodzą z metody Frobeniusa dla regularnego
punktu osobliwego. Samo oszacowanie całki nie dowodziłoby
tej silniejszej postaci.
Przy dodatnim obrocie logarytm zmienia się o $2\pi\mathrm i$,
zatem $u(1-\lambda)\mapsto u(1-\lambda)-2\mathrm i\,u(\lambda)$
i $v\mapsto v+2u$.
\end{proof}
\begin{uwaga}[Dlaczego pojawia się $2$, a nie $1$?]
W standardowym modelu jednokrotnego wygładzenia węzła
$zw=t$ globalna monodromia krzywej jest transwekcją.
Rodzina Legendre'a w parametrze $\lambda$ nie jest tym samym
modelem jednoparametrowym: przy $\lambda=0$
$\Delta(\lambda)$ ma zero rzędu dwa, a odpowiedni
parametr wygładzający lokalnie jest jednostką pomnożoną
przez $\lambda^2$. Dlatego monodromia na integralnej
homologii jest kwadratem transwekcji,
zależnie od orientacji pętli sprzężonym z
$\left(\begin{smallmatrix}1&2\\0&1\end{smallmatrix}\right)$
lub z jego odwrotnością.
Można to sprawdzić bez odwoływania się do samego
wyróżnika. W otoczeniu $x=y=\lambda=0$ wybieramy
holomorficzną gałąź $\sqrt{1-x}$ i wprowadzamy
$Y=y/\sqrt{1-x}$ oraz $X=x-\lambda/2$.
Wtedy równanie rodziny przechodzi dokładnie w
$Y^2+X^2=\lambda^2/4$.
Po podstawieniu $U=Y+\mathrm iX$, $V=Y-\mathrm iX$
otrzymujemy $UV=\lambda^2/4$. Jest to zwykłe
wygładzenie węzła po bazowej zmianie parametru
$t=\lambda^2/4$. Zatem okrążenie $\lambda=0$
odpowiada dwukrotnemu okrążeniu $t=0$.
Sama obecność zwyczajnego węzła we włóknie
centralnym nie wyznacza bez znajomości całkowitej
rodziny wykładnika potęgi transwekcji.
\end{uwaga}

\begin{center}
\begin{tikzpicture}[>=Stealth,scale=0.89,font=\small]
\draw[->] (-2.3,0)--(2.3,0);
\node[below] at (2.00,-0.10) {$\Re\lambda$};
\draw[->] (0,-1.65)--(0,1.65) node[above] {$\Im\lambda$};
\draw[manifoldblue,thick] (0,0) circle (1.05);
\fill[manifoldred] (0,0) circle (2.4pt);
\node[below right] at (0,0) {$0$};
\draw[manifoldgreen,very thick,->] (1.05,0)
 arc[start angle=0,end angle=320,radius=1.05];
\node[above right] at (0.85,0.85) {$\gamma$};
\node at (0,-2.02) {pętla w $\Delta^*$};
\draw[->,thick] (2.55,0)--(4.05,0);
\node[above] at (3.3,0.10) {okresy};
\draw[->] (4.5,0)--(8.35,0);
\node[below] at (8.05,-0.10) {$\Re\tau$};
\draw[->] (5.75,-0.35)--(5.75,2.75) node[above] {$\Im\tau$};
\draw[manifoldblue,densely dotted,thick] (4.5,0)--(8.1,0);
\fill[manifoldblue] (5.55,1.3) circle (2.2pt);
\fill[manifoldgreen] (7.55,1.3) circle (2.2pt);
\draw[manifoldgreen,thick,->] (5.75,1.3)--(7.35,1.3);
\node[below] at (5.55,1.3) {$\tau$};
\node[above] at (7.55,1.3) {$\tau+2$};
\node at (6.35,-1.02) {górna półpłaszczyzna};
\end{tikzpicture}
\par\smallskip\small Schemat: okrążenie $\lambda=0$ w rodzinie Legendre'a
zmienia lokalnie wybraną gałąź stosunku okresów o $2$.
Wykres nie przedstawia globalnej mapy jednoznacznej.
\end{center}

\section{Quasi-unipotencja, filtracja monodromiczna i graniczna struktura Hodge'a}
\begin{definicja}[Unipotencja i logarytm monodromii]
Niech $V_\Q$ będzie skończeniewymiarową przestrzenią
wektorową, a $T\in\operatorname{GL}(V_\Q)$.
Operator $T$ nazywamy \emph{unipotentnym}, jeżeli
$(T-I)^r=0$ dla pewnego $r\ge1$, a \emph{quasi-unipotentnym},
jeżeli pewna jego dodatnia potęga jest unipotentna.
Jeżeli $T$ jest unipotentny, definiujemy nilpotentny
operator $N=\log T:=\sum_{j\ge1}(-1)^{j+1}(T-I)^j/j$.
Suma jest skończona.
\end{definicja}
\begin{twierdzenie}[Twierdzenie o lokalnej monodromii: zakres]
Niech $f:\mathcal X\to\Delta^*$ będzie gładką właściwą
rodziną projektową, która pochodzi z algebraicznej rodziny
nad przebitym dyskiem. Wówczas lokalna monodromia na
$H^k(X_t;\Q)$ jest quasi-unipotentna: jej wartości własne
są pierwiastkami z jedności. Po skończonym rozgałęzionym
podstawieniu $t=s^m$ otrzymujemy unipotentną monodromię.
\end{twierdzenie}
\begin{proof}[Status]
Jest to twierdzenie o monodromii regularnej osobliwości
w geometrycznej wariacji Hodge'a. W ogólności nie wynika
z faktu istnienia płaskiej koneksji; konieczny jest warunek
geometryczności oraz twierdzenie o regularności.
Przyjmujemy je jako wynik zewnętrzny. Rachunek Legendre'a
powyżej dowodzi unipotencji bezpośrednio dla tej rodziny.
\end{proof}

\begin{definicja}[Filtracja monodromiczna w przypadku $N^2=0$]
Niech $V_\Q$ będzie dwuwymiarowa, $N\ne0$ i $N^2=0$.
Filtracją wag monodromii \emph{z centrum $1$} nazywamy
filtrację $0=W_{-1}\subset W_0\subset W_1\subset W_2=V_\Q$
zadaną przez $W_0=\operatorname{im}N$ oraz
$W_1=\ker N$. W rozważanym przypadku
$W_0=W_1$ jest jednowymiarowa.
\end{definicja}
\begin{lemat}[Jednoznaczność filtracji w rzędzie dwa]
\label{lem:xv-weights-rank2}
Dla $V_\Q$ wymiaru dwa i $N\ne0$, $N^2=0$,
filtracja z poprzedniej definicji jest jedyną filtracją
$W_\bullet$ spełniającą $NW_j\subset W_{j-2}$ oraz
taką, że $N:\operatorname{Gr}_2^W V\to
\operatorname{Gr}_0^W V$ jest izomorfizmem.
\end{lemat}
\begin{proof}
Nilpotentny niezerowy operator na przestrzeni
wymiaru dwa ma rząd jeden, a z $N^2=0$ wynika
$\operatorname{im}N\subset\ker N$. Obie przestrzenie
mają wymiar jeden, zatem są równe.
Jeśli $NW_2\subset W_0$ i indukowana mapa między
skrajnymi ilorazami ma być izomorfizmem, to
$W_0$ musi zawierać $\operatorname{im}N$ i mieć wymiar
jeden, więc $W_0=\operatorname{im}N$.
Ponieważ $NW_1=0$ i $W_1$ ma wymiar jeden,
$W_1=\ker N$. Odwrotnie, te wybory oczywiście
spełniają wskazane warunki, a mapa na ilorazach
jest niezerowa między jednowymiarowymi przestrzeniami,
więc jest izomorfizmem.
\end{proof}

\begin{twierdzenie}[Graniczna mieszana struktura Hodge'a: wersja używana]
\label{thm:xv-lmhs}
Niech $(\mathbb V_\Z,F^\bullet)$ będzie polaryzowaną
wariacją struktur Hodge'a wagi $k$ nad $\Delta^*$,
pochodzącą od algebraicznej rodziny, i przypuśćmy,
że monodromia $T$ jest unipotentna.
Na uniwersalnym nakryciu $t=e^{2\pi\mathrm iz}$
można utożsamić wszystkie włókna $\mathbb V$ z ustaloną
przestrzenią $V_\Q$, a filtrację Hodge'a zapisać jako
$F^\bullet(z)$. Przy odpowiedniej konwencji działania
monodromii granica
$F^\bullet_\infty=\lim_{\operatorname{Im}z\to\infty}
\exp(-zN)F^\bullet(z)$, $N=\log T$,
istnieje w przestrzeni flag (po użyciu kanonicznego
przedłużenia) i wraz z przesuniętą do centrum $k$
filtracją monodromiczną $W(N)[-k]$ tworzy graniczną
mieszaną strukturę Hodge'a.
\end{twierdzenie}
\begin{proof}[Status i sens wzoru]
Wartość $F^\bullet(z)$ zależy od gałęzi logarytmu,
a po zastąpieniu $z$ przez $z+1$ jest przekształcana
przez $T=\exp N$. Wtedy
$\exp(-(z+1)N)\exp(N)F^\bullet(z)
=\exp(-zN)F^\bullet(z)$: skorygowana filtracja
jest jednowartościowa. Jest to bezpośrednio sprawdzony
fragment twierdzenia. Jej przedłużenie przez $t=0$,
istnienie granicy i warunki mieszanej struktury Hodge'a
wymagają twierdzenia Schmida o orbicie nilpotentnej
(oraz odpowiedniej teorii przedłużeń). Nie wynikają
z samego rachunku z eksponentą i tutaj korzystamy
z tego głębokiego twierdzenia.
\end{proof}

\begin{przyklad}[Dwie różne struktury przy degeneracji krzywej węzłowej]
\label{ex:xv-central-limit}
Niech gładka krzywa eliptyczna degeneruje do
nieredukowalnej krzywej wymiernej z jednym węzłem,
a lokalna monodromia na $V=H^1(E_t;\Q)$ jest
nieidentycznościowa i unipotentna.
Wtedy $\dim_\Q V=2$, $N\ne0$, $N^2=0$.
Według Lematu~\ref{lem:xv-weights-rank2}
$\operatorname{Gr}_0^WV\cong\Q$,
$\operatorname{Gr}_1^WV=0$,
$\operatorname{Gr}_2^WV\cong\Q(-1)$ jako czyste
struktury Hodge'a (ostatnia identyfikacja wynika z
polaryzacji i twierdzenia o granicznej strukturze).
Tymczasem dla \emph{singularnego włókna centralnego}
$H^1(E_0;\Q)\cong\Q$ ma jedynie wagę $0$,
jak wykazaliśmy w Przykładzie~\ref{ex:nodal-mhs-xiv}.
Zatem $H^1(E_0)$ i graniczna mieszana struktura
Hodge'a gładkiego włókna nie są tym samym obiektem:
nie zgadzają się nawet ich wymiary.
\end{przyklad}
\begin{uwaga}[Tate'owskie przesunięcie wagi]
Symbol $\Q(-1)$ oznacza jednowymiarową czystą
strukturę Hodge'a wagi $2$, typu $(1,1)$.
W przykładzie powyżej $N$ obniża wagę o dwa;
po uwzględnieniu skręcenia Tate'a otrzymujemy
izomorfizm $N:\operatorname{Gr}_2^WV\to
\operatorname{Gr}_0^WV(-1)$.
\end{uwaga}

\section{Granice twierdzeń i dalsza droga do geometrii algebraicznej}
W tym rozdziale w pełni udowodniliśmy istnienie płaskiej
koneksji z lokalnego systemu, wzór na różniczkowanie okresów,
ogólne kryterium istnienia równania Picarda--Fuchsa
przy wyborze algebraicznego parametru lokalnego,
rozwinięcie szeregu okresu Legendre'a, jego równanie i
monodromię oraz lemat liniowo-algebraiczny o filtracji
monodromicznej. Przyjęliśmy natomiast klasyczne twierdzenia
o wariacjach Hodge'a pochodzących z rodzin projektowych,
regularności i quasi-unipotencji monodromii oraz istnieniu
granicznej mieszanej struktury Hodge'a.
Ich niezależne dowody wymagają analizy asymptotycznej
koneksji z osobliwościami regularnymi, teorii modułów
Hodge'a i głębszych twierdzeń algebraiczno-geometrycznych.
Takie rozdzielenie zapobiega kołowemu używaniu
w teorii snopów perwersyjnych i twierdzenia
o rozkładzie.

\chapter{Sześć operacji snopowych, dualność Verdiera i cykle specjalizacji}
\label{ch:xvi-szesc-operacji}

\section{Motywacja: dlaczego zwykły obraz bezpośredni nie wystarcza?}
W Rozdziale~\ref{ch:snopy-verdier-v12} skonstruowaliśmy snopy,
przedłużenie przez zero i kohomologię z nośnikami zwartymi, a w
Rozdziale~\ref{ch:xiv-hodge-monodromia} napotkaliśmy cykle pobliskie.
Potrzebujemy teraz języka, w którym operacje te tworzą jeden spójny
rachunek. Problem jest geometryczny: jeśli $f:X\to Y$ nie jest właściwe,
przesuwanie informacji z $X$ na $Y$ może zachować sekcje rozciągające się
do nieskończoności albo odrzucić te, których nośnik nie jest właściwy
nad $Y$. Różnicę kodują odpowiednio $Rf_*$ i $Rf_!$.

Ustalamy ciało $k$ i korzystamy z kompleksów snopów $k$-przestrzeni
wektorowych, numerowanych kohomologicznie. Dla przestrzeni $X$
piszemy $D(X)=D(\operatorname{Sh}(X;k))$ i $D^b_c(X)$ dla
ograniczonych kompleksów o konstruktybilnych snopach kohomologii.
W twierdzeniach konstruktybilnych $X$ i $Y$ oznaczają
analityfikacje zespolonych odmian algebraicznych o skończonej
stratyfikacji algebraicznej, a $f:X\\to Y$ jest analityfikacją
morfizmu algebraicznego. Na takich przestrzeniach posługujemy się
kategorią $D^b_c(X;k)$ kompleksów o konstruktybilnych snopach
kohomologii. W przykładach topologicznych można szerzej dopuścić
dobre, lokalnie zwarte przestrzenie Hausdorffa o skończonym wymiarze
kohomologicznym, pod warunkiem sprawdzenia właściwych hipotez
konkretnego twierdzenia. Kategoria $D(X)$ wszystkich kompleksów nie
musi zachowywać ograniczoności ani konstruktybilności przy dowolnej
mapie $f$. Twierdzenia formalizmu sześciu operacji przyjmujemy w
sprecyzowanej kategorii, a dowody przypadków szczególnych zaznaczamy.
Dla kompleksu $A$ piszemy $\\mathcal H^q(Rf_*A)$ na oznaczenie
jego $q$-tego snopu kohomologii; symbol $R^qf_*\\mathcal F$
stosujemy do snopu $\\mathcal F$ skoncentrowanego w stopniu zerowym. Współczynniki w ciele upraszczają
iloczyn tensorowy, lecz nie eliminują konieczności pochodnego
wewnętrznego Homu.

\begin{definicja}[Przesunięcie i stożek]
Dla kompleksu kohomologicznego $K$ definiujemy $K[1]^n=K^{n+1}$
i $d_{K[1]}=-d_K$. Jeśli $u:A\to B$ jest morfizmem kompleksów,
jego stożek $\operatorname{Cone}(u)$ ma
$\operatorname{Cone}(u)^n=B^n\oplus A^{n+1}$ i
$d(b,a)=(d_Bb+u(a),-d_Aa)$. W kategorii pochodnej otrzymujemy
trójkąt wyróżniony
$A\xrightarrow{u}B\to\operatorname{Cone}(u)\to A[1]$.
\end{definicja}
\begin{lemat}[Długi ciąg kohomologii stożka]
Dla każdego morfizmu kompleksów $u:A\to B$ istnieje naturalny
ciąg dokładny
$\cdots\to H^n(A)\to H^n(B)\to
H^n(\operatorname{Cone}(u))\to H^{n+1}(A)\to\cdots$.
\end{lemat}
\begin{proof}
Rzut $\operatorname{Cone}(u)^n\to A^{n+1}$ daje krótki ciąg
dokładny kompleksów $0\to B\to\operatorname{Cone}(u)\to A[1]\to0$.
Morfizm łączący na kohomologii jest indukowany przez $u$ (z
konwencją znaku zależną od przyjętej definicji stożka). Standardowa
konstrukcja długiego ciągu dokładnego daje tezę.
\end{proof}

\section{Sześć funktorów i ich adjunkcje}
\begin{definicja}[Dwie pary funktorów związanych z odwzorowaniem]
Niech $f:X\to Y$ będzie mapą w rozważanej kategorii przestrzeni.
Przez $f^*=f^{-1}$ oznaczamy dokładny funktor obrazu odwrotnego
snopów $k$-przestrzeni wektorowych, a przez $Rf_*$ jego prawy
pochodny obraz bezpośredni. Dla snopu $\mathcal F$ obraz bezpośredni
jest dany przez $(f_*\mathcal F)(V)=\mathcal F(f^{-1}V)$.
Przez $f_!\mathcal F$ oznaczamy snop sekcji o nośniku właściwym
nad otwartym podzbiorem bazy: na $V\subseteq Y$ rozważamy
sekcje nad $f^{-1}V$, których nośnik odwzorowuje się właściwie na $V$,
i dokonujemy snopifikacji, jeśli jest potrzebna. Prawy funktor
pochodny oznaczamy $Rf_!$. W tych samych założeniach $f^!$
jest prawym adjunktem do $Rf_!$ na odpowiednich kategoriach
pochodnych.
\end{definicja}
\begin{uwaga}[Cztery operacje geometryczne]
Zapis $f^*$ nie jest tutaj rozszerzaniem skalarów po pierścieniach
funkcji; pracujemy ze snopami $k$-przestrzeni wektorowych.
Dla $f$ właściwego zachodzi $Rf_!\simeq Rf_*$. Zwykły
$f_*$ jest lewostronnie dokładny, zaś $f_!$ i $Rf_!$
nie muszą zachowywać krótkich ciągów dokładnych bez przejścia
na kategorię pochodną.
\end{uwaga}
\begin{definicja}[Tensor i wewnętrzny Hom]
Dla $A,B\in D(X)$ przez $A\otimes_k^L B$ oznaczamy
pochodny iloczyn tensorowy snopów, a przez
$R\mathcal Hom_X(A,B)$ pochodny \emph{wewnętrzny Hom},
charakteryzowany przez adjunkcję
\[
\operatorname{Hom}_{D(X)}(A\otimes_k^L B,C)
\cong\operatorname{Hom}_{D(X)}
\bigl(A,R\mathcal Hom_X(B,C)\bigr).
\]
W naszym przypadku każdy $k$-moduł jest płaski, toteż
$\otimes_k^L$ można obliczać zwykłym tensorem kompleksów.
Natomiast $R\mathcal Hom_X$ jest funktorem pochodnym w kategorii
snopów: zwykły wewnętrzny Hom nie wystarcza.
\end{definicja}
\begin{twierdzenie}[Podstawowe adjunkcje: sformułowanie formalizmu]
\label{thm:xvi-adjunctions}
Niech $f:X\to Y$ będzie morfizmem, dla którego obowiązuje
formalizm sześciu operacji. Dla $A\in D(X)$ i $B\in D(Y)$
istnieją naturalne izomorfizmy
\[
\operatorname{Hom}_{D(X)}(f^*B,A)
\cong\operatorname{Hom}_{D(Y)}(B,Rf_*A),
\quad
\operatorname{Hom}_{D(Y)}(Rf_!A,B)
\cong\operatorname{Hom}_{D(X)}(A,f^!B).
\]
Razem z $\otimes^L$ i $R\mathcal Hom$ tworzą one sześć
operacji $f^*,Rf_*,Rf_!,f^!,\otimes^L,R\mathcal Hom$.
\end{twierdzenie}
\begin{proof}[Część dowiedziona i część przyjęta]
Pierwsza adjunkcja w stopniu zerowym wynika z definicji obrazu
odwrotnego: morfizm $f^{-1}\mathcal G\to\mathcal F$
odpowiada morfizmowi $\mathcal G\to f_*\mathcal F$.
Ponieważ $f^{-1}$ jest dokładny, adjunkcja przechodzi na
kompleksy przez zastosowanie rezolwent injektywnych do
argumentu $\mathcal F$.
Druga adjunkcja \emph{definiuje} szczególny obraz odwrotny $f^!$;
świadomego dowodu jego istnienia dla ogólnych map nie można
zastąpić manipulacją symbolami. Istnienie i własności $f^!$
w przyjętej kategorii stanowią tu standardowe twierdzenie
formalizmu dualności Verdiera.
\end{proof}

\section{Otwarte i zamknięte włożenia: jawne obliczenia}
Niech $i:Z\hookrightarrow X$ będzie włożeniem zamkniętym,
a $j:U=X\setminus Z\hookrightarrow X$ otwartym.
Dla snopu $\mathcal F$ na $X$ definiujemy snop sekcji
z nośnikiem w $Z$ przez
$(\underline\Gamma_Z\mathcal F)(V)=
\{s\in\mathcal F(V):\operatorname{supp}(s)\subseteq Z\cap V\}$
na otwartych $V\subseteq X$. Odróżniamy go od przestrzeni
sekcji globalnych $\Gamma_Z(X,\mathcal F)$.
Wówczas dla $A\in D(X)$ zachodzi
$i_*i^!A\simeq R\underline\Gamma_Z(A)$;
jest to pochodny kompleks snopów z nośnikiem w $Z$.

\begin{lemat}[Pierwszy ciąg lokalizacyjny na poziomie snopów]
Dla każdego snopu $\mathcal F$ na $X$ mamy dokładny ciąg
\[
0\longrightarrow j_!j^*\mathcal F
\longrightarrow\mathcal F\longrightarrow i_*i^*\mathcal F
\longrightarrow0.
\]
\end{lemat}
\begin{proof}
Wystarczy sprawdzić dokładność na kiełkach.
Dla $x\in U$ kiełki trzech pierwszych wyrazów to
$\mathcal F_x\to\mathcal F_x\to0$.
Dla $x\in Z$ są to $0\to\mathcal F_x\to\mathcal F_x$.
W obu przypadkach otrzymujemy ciąg dokładny.
\end{proof}
\begin{twierdzenie}[Dwa trójkąty lokalizacyjne]
\label{thm:xvi-localization}
Dla $A\in D^b_c(X)$ zachodzą naturalne trójkąty wyróżnione
\[
 j_!j^*A\to A\to i_*i^*A\to (j_!j^*A)[1],
\]
\[
 i_*i^!A\to A\to Rj_*j^*A\to (i_*i^!A)[1].
\]
\end{twierdzenie}
\begin{proof}[Pierwszy trójkąt i status drugiego]
Pierwszy trójkąt pochodzi z poprzedniego krótkiego ciągu dokładnego
zastosowanego składnikowo do kompleksu; funktory $j_!,j^*,i_*,i^*$
są tutaj dokładne. Dla drugiego wybierzmy rezolwentę
injektywną $A\to I^\bullet$. Snopy injektywne są wiotkie
(\emph{flabby}), czyli ograniczenie sekcji do dowolnego
mniejszego zbioru otwartego jest surjekcją.
Dla każdego otwartego $V\subseteq X$ mamy zatem krótki ciąg
$0\to\underline\Gamma_Z(I^q)(V)
\to I^q(V)\to I^q(V\cap U)\to0$.
Po zebraniu stopni i przejściu na snopy otrzymujemy
krótki ciąg kompleksów
$0\to\underline\Gamma_Z(I^\bullet)
\to I^\bullet\to j_*j^*I^\bullet\to0$.
Ponieważ $j^*$ ma dokładny lewy adjunkt $j_!$,
zachowuje snopy injektywne; dlatego
$j_*j^*I^\bullet$ oblicza $Rj_*j^*A$.
Pierwszy kompleks jest modelem dla
$R\underline\Gamma_Z(A)\simeq i_*i^!A$.
Powstały trójkąt ma zatem żądaną postać.
Korzystaliśmy tutaj z istnienia rezolwent injektywnych
oraz identyfikacji pochodnych sekcji z nośnikiem
z $i_*i^!$, której konstrukcja wynika z adjunkcji.
\end{proof}
\begin{przyklad}[Zwykły i wyjątkowy obraz punktu na prostej]
\label{ex:xvi-point-line}
Niech $X=\R$, $Z=\{0\}$, $U=\R\setminus\{0\}$,
$A=k_{\R}$. Wówczas $i^*A\simeq k_Z$, ale
$i^!A\simeq k_Z[-1]$. Aby otrzymać drugą równość,
obliczamy lokalną kohomologię:
$H^q(\R,\R\setminus\{0\};k)$ jest równa $k$ dla $q=1$
i zerowa w pozostałych stopniach. Rzeczywiście, długi ciąg
pary i $H^0(\R\setminus\{0\};k)\cong k\oplus k$
pokazują, że $H^1$ jest ilorazem $k^2/\Delta k\cong k$.
Daje to $i^!k_{\R}\simeq k_Z[-1]$.
\end{przyklad}
\begin{przyklad}[Dlaczego $j_!\ne j_*$?]
Dla $j:(0,1)\hookrightarrow\R$ i $A=k_{(0,1)}$
kiełek $(j_!A)_0$ jest zerowy, natomiast
$(j_*A)_0\cong k$. Z drugiego trójkąta wynika również,
że $Rj_*$ może mieć niezerowe wyższe snopy kohomologii,
jeśli lokalne dopełnienie jest topologicznie nietrywialne.
W tym konkretnym przykładzie wyższe $R^qj_*A$ znikają.
\end{przyklad}


\section{Jak rozpoznać $f^!$? Trzy podstawowe przypadki}
\begin{stwierdzenie}[Włożenie otwarte]
\label{prop:xvi-open-shriek}
Dla włożenia otwartego $j:U\hookrightarrow X$ zachodzi
$j^!\simeq j^*$ na kategorii pochodnej snopów $k$-przestrzeni
wektorowych.
\end{stwierdzenie}
\begin{proof}
Dla snopu $A$ na $U$ i snopu $B$ na $X$ naturalne
przyporządkowanie morfizmów daje
$\operatorname{Hom}_X(j_!A,B)\cong
\operatorname{Hom}_U(A,j^*B)$.
Przyporządkowanie jest zgodne ze złożeniami i wynika
z definicji sekcji przedłużonych przez zero.
Ponieważ $j_!$ i $j^*$ są dokładne (co można sprawdzić
na kiełkach), ta sama adjunkcja obowiązuje dla kompleksów.
Natomiast $j^!$ jest z definicji prawym adjunktem do $Rj_!=j_!$.
Jednoznaczność adjunktu daje $j^!\simeq j^*$.
\end{proof}

\begin{twierdzenie}[Gładka czystość: wersja analityczno-topologiczna]
\label{thm:xvi-smooth-purity}
Niech $f:M\to N$ będzie gładką submersją rzeczywistych
rozmaitości, której włókna mają stały wymiar $d$, i niech
$\operatorname{or}_{M/N}$ będzie systemem ich orientacji.
Dla konstruktybilnego kompleksu $B\in D^b_c(N)$ zachodzi
naturalny izomorfizm
\[
 f^!B\simeq f^*B\otimes_k\operatorname{or}_{M/N}[d].
\]
Dla orientowanej submersji system orientacji jest trywialny;
w szczególności dla $a:\R^d\to\{*\}$ mamy
$a^!k\simeq k_{\R^d}[d]$.
\end{twierdzenie}
\begin{proof}[Rachunek lokalny i zaznaczenie zależności]
Z twierdzenia o lokalnej postaci submersji dostajemy
lokalne współrzędne, w których $f$ jest projekcją
$V\times\R^d\to V$. Zorientowane włókno posiada lokalną
klasę fundamentalną w
$H_c^d(\R^d;k)\cong k$, więc przesunięcie o $d$
i skręcenie systemem orientacji są konieczne, aby
adjunkcja $Rf_!\dashv f^!$ zachodziła w modelu
produktowym. Na przecięciach kart zmiana lokalnej
orientacji mnoży generator przez $+1$ lub $-1$,
co składa się na $\operatorname{or}_{M/N}$.
Zgodność lokalnych identyfikacji z ogólną adjunkcją i
ich kanoniczne sklejenie są częścią twierdzenia o
czystości (purity) w dualności Verdiera. Nie wywodzimy
tej globalnej części wyłącznie z lokalnego rachunku grup.
\end{proof}
\begin{przyklad}[Ta sama przestrzeń, różne obrazy do punktu]
Dla $a:\R\to\{*\}$ mamy
$Ra_*k_{\R}\simeq k$ oraz $Ra_!k_{\R}\simeq k[-1]$.
Pierwsza to zwykła kohomologia $\R$, druga to
kohomologia z nośnikami zwartymi; jest ona skoncentrowana
w stopniu $1$. Ponadto $a^!k\simeq k_{\R}[1]$.
Te trzy wzory pokazują jednocześnie, dlaczego
$f_*$ nie zastępuje $f_!$ i dlaczego $f^!$
nie powinno się utożsamiać z $f^*$.
\end{przyklad}

\section{Kohomologia z nośnikiem zwartym i obiekt dualizujący}
\begin{definicja}[Kompleks dualizujący]
Dla $a:X\to\{*\}$ definiujemy $\omega_X=a^!k$,
a dla $A\in D^b_c(X)$ dual Verdiera
$\mathbb D_XA=R\mathcal Hom_X(A,\omega_X)$.
\end{definicja}
\begin{twierdzenie}[Dualność Verdiera dla ciała]
\label{thm:xvi-verdier}
Dla lokalnie zwartej, odpowiednio skończeniewymiarowej
przestrzeni $X$ i $A\in D^b_c(X)$ zachodzi naturalny
quasi-izomorfizm
\[
R\Gamma(X,\mathbb D_X A)
\simeq R\operatorname{Hom}_k(R\Gamma_c(X,A),k).
\]
Jeśli $X$ jest rzeczywistą rozmaitością wymiaru $n$,
$\omega_X\simeq\operatorname{or}_X[n]$, gdzie
$\operatorname{or}_X$ jest systemem orientacji.
\end{twierdzenie}
\begin{proof}[Dowód wskazanego modelu i wyraźna zależność]
Dla $X=\R^n$ i $A=k_X$ obie strony mają jedną
jednowymiarową grupę kohomologii w stopniu $-n$:
$R\Gamma_c(\R^n;k)\simeq k[-n]$, zatem
$R\operatorname{Hom}_k(k[-n],k)\simeq k[n]$.
Również $R\Gamma(\R^n,k_X[n])\simeq k[n]$.
Dla ogólnej rozmaitości te lokalne identyfikacje skleja
system orientacji. Istnienie kanonicznego quasi-izomorfizmu,
funktorialność i twierdzenie dla osobliwych $X$ należą do
pełnej dualności Verdiera; nie wynikają z porównania wymiarów
w przykładzie euklidesowym.
\end{proof}
\begin{wniosek}[Dualność Poincarégo jako szczególny przypadek]
Jeśli $M$ jest spójną, orientowalną rozmaitością rzeczywistą
wymiaru $n$, to dla ciała $k$ zachodzi
$H_c^q(M;k)^*\cong H^{n-q}(M;k)$ przy odpowiednim
sparowaniu dualności. Dla rozmaitości nieorientowalnej
prawa strona otrzymuje lokalne współczynniki orientacji.
\end{wniosek}
\begin{proof}
Wstawiamy $A=k_M$ do poprzedniego twierdzenia i
$\omega_M\simeq\operatorname{or}_M[n]$;
porównanie stopni kohomologicznych daje wynik.
\end{proof}
\begin{uwaga}[Uwaga o skończoności wymiarów]
Identyfikacja przez liniowy dual w powyższej postaci jest
naturalna nad ciałem, lecz nie oznacza automatycznie, że
wszystkie grupy kohomologii są skończeniewymiarowe.
W obliczeniach na skończonych kompleksach i na zwartych
odmianach algebraicznych problem ten nie występuje.
\end{uwaga}

\section{Dualność i obrazy z nośnikami właściwymi}
\begin{twierdzenie}[Wymiana obrazów przez dualność]
\label{thm:xvi-duality-exchange}
Niech $f:X\to Y$ będzie morfizmem w kategorii
konstruktybilnych kompleksów, a $A\in D^b_c(X)$.
Istnieje naturalny izomorfizm
\[
\mathbb D_Y(Rf_!A)\simeq Rf_*\mathbb D_X A.
\]
Jeśli $f$ jest właściwe, $Rf_!\simeq Rf_*$ i
$\mathbb D_Y Rf_*A\simeq Rf_*\mathbb D_XA$.
\end{twierdzenie}
\begin{proof}[Wyprowadzenie z adjunkcji]
Dla $B\in D^b_c(Y)$ rachunek formalny daje
\[
\begin{aligned}
R\operatorname{Hom}_Y(B,\mathbb D_Y Rf_!A)
&\simeq R\operatorname{Hom}_Y(B\otimes^L Rf_!A,\omega_Y)\\
&\simeq R\operatorname{Hom}_X(f^*B\otimes^L A,f^!\omega_Y)\\
&\simeq R\operatorname{Hom}_X(f^*B,\mathbb D_XA)\\
&\simeq R\operatorname{Hom}_Y(B,Rf_*\mathbb D_XA).
\end{aligned}
\]
W drugim kroku korzystamy ze wzoru projekcji i adjunkcji
$(Rf_!,f^!)$; w trzecim z kanonicznego
$f^!\omega_Y\simeq\omega_X$. Yoneda daje izomorfizm.
To formalny dowód \emph{warunkowy}: zależy od wcześniej
przyjętych twierdzeń o wzorze projekcji i istnieniu $f^!$.
\end{proof}

\section{Zmiana podstawy i wzór projekcji}
Rozważmy kartezjański kwadrat
\[
\begin{tikzpicture}[baseline=(current bounding box.center),node distance=22mm]
\node (Xp) {$X'$}; \node[right=of Xp] (X) {$X$};
\node[below=of Xp] (Yp) {$Y'$}; \node[right=of Yp] (Y) {$Y$};
\draw[->] (Xp)--node[above] {$g'$}(X);
\draw[->] (Xp)--node[left] {$f'$}(Yp);
\draw[->] (X)--node[right] {$f$}(Y);
\draw[->] (Yp)--node[below] {$g$}(Y);
\end{tikzpicture}
\]

\begin{twierdzenie}[Zmiana podstawy dla obrazu z nośnikami właściwymi]
\label{thm:xvi-basechange}
Przy standardowych założeniach formalizmu sześciu operacji
istnieje naturalny izomorfizm
$g^*Rf_!A\simeq Rf'_!(g')^*A$
dla konstruktybilnego $A$.
Jeśli $f$ jest właściwa, otrzymujemy też
$g^*Rf_*A\simeq Rf'_*(g')^*A$.
\end{twierdzenie}
\begin{proof}[Dowód w prostym przypadku i zakres twierdzenia]
Jeśli $f$ jest projekcją skończonego dyskretnego zbioru
$\{1,\dots,m\}\times Y\to Y$, to $f_!=f_*$ i
$(Rf_*A)_y\cong\bigoplus_{i=1}^m A_{(i,y)}$.
Identyfikacja pozostaje prawdziwa po dowolnym
$g:Y'\to Y$ i daje żądany izomorfizm na kiełkach.
Przypadek ogólny wymaga twierdzenia o zmianie podstawy
z właściwymi nośnikami; jest tu przyjęty, a nie
uzasadniony tym szczególnym obliczeniem.
\end{proof}

\begin{twierdzenie}[Wzór projekcji]
\label{thm:xvi-projection}
Jeżeli $f:X\to Y$ jest mapą rozważanej kategorii,
$A\in D^b_c(X)$ oraz $L$ jest lokalnym systemem
skończeniewymiarowych $k$-przestrzeni na $Y$, to
\[
Rf_!(A\otimes_k f^*L)\simeq (Rf_!A)\otimes_k L.
\]
Ogólniejszy wzór dla $B\in D^b_c(Y)$ ma postać
$Rf_!(A\otimes^L f^*B)\simeq Rf_!A\otimes^L B$,
przy zwykłych hipotezach konstruktybilności i
skończonego wymiaru kohomologicznego.
\end{twierdzenie}
\begin{proof}[Wersja dla lokalnego systemu]
Lokalnie na $Y$ mamy $L|_V\cong k_V^r$.
Dla $V$ otwartego po obu stronach dostajemy
odpowiednio $r$ kopii $Rf_!A|_V$, ponieważ
funktor $Rf_!$ jest addytywny, a tensor z $k^r$
odpowiada sumie skończonej. Izomorfizmy są naturalne
w zmianach trywializacji $L$, więc sklejają się na $Y$.
Uogólnienie na dowolne konstruktybilne $B$ nie wynika
z lokalnej trywialności $B$: jest odrębnym twierdzeniem
formalizmu sześciu operacji.
\end{proof}

\section{Włókna, obrazy pochodne i ciąg spektralny Leraya}
\begin{twierdzenie}[Kohomologia włókien przy mapie właściwej]
Niech $f:X\to Y$ będzie właściwą mapą rozważanej kategorii,
$A\in D^b_c(X;k)$, $y\in Y$, a $i_y:\{y\}\hookrightarrow Y$.
Dla każdego $q\in\Z$ istnieje naturalny izomorfizm
\[
\mathcal H^q(i_y^*Rf_*A)
\cong\mathbb H^q(f^{-1}(y);A|_{f^{-1}(y)}).
\]
W szczególności lewa strona jest kiełkiem
$\mathcal H^q(Rf_*A)$ w $y$. Dla snopu $\mathcal F$
umieszczonego w stopniu zerowym można ją zapisać
$(R^qf_*\mathcal F)_y$.
\end{twierdzenie}
\begin{proof}
Zastosujmy Twierdzenie~\ref{thm:xvi-basechange} do
$g:\{y\}\hookrightarrow Y$. Lewa strona jest
kompleksem $i_y^*Rf_*A$, a prawa
$R\Gamma(f^{-1}(y),A|_{f^{-1}(y)})$.
Po wzięciu kohomologii stopnia $q$ otrzymujemy tezę.
\end{proof}
\begin{uwaga}
Dla $f$ niewłaściwej analogiczny wzór z $Rf_*$ jest na ogół
fałszywy: wymaga dodatkowej hipotezy zmiany podstawy. Natomiast przy obrazach z nośnikami właściwymi
otrzymujemy lokalnie grupy $H_c^q(f^{-1}(y);A|_{f^{-1}(y)})$,
o ile spełniono założenia zmiany podstawy dla $Rf_!$.
\end{uwaga}
\begin{przyklad}[Złożenie $x\mapsto x^2$]
Dla właściwej mapy $f:\R\to\R$, $f(x)=x^2$, włókno
nad $y<0$ jest puste, nad $y=0$ ma jeden punkt, a nad
$y>0$ dwa punkty. Zatem snop $R^0f_*k_\R$ ma
odpowiednio rzędy $0,1,2$, a $R^qf_*k_\R=0$ dla $q>0$.
To konstruktybilny snop, który nie jest lokalnym systemem
na całej prostej.
\end{przyklad}
\begin{twierdzenie}[Ciąg Leraya]
Dla mapy $f:X\to Y$ w rozważanej kategorii i snopu
$\mathcal F$ istnieje pierwszoćwiartkowy ciąg spektralny
$E_2^{p,q}=H^p(Y,R^qf_*\mathcal F)
\Longrightarrow H^{p+q}(X,\mathcal F)$.
Dla $Rf_!$ i $R\Gamma_c$ istnieje odpowiedni ciąg
spektralny z nośnikami zwartymi, przy standardowych
założeniach skończoności.
\end{twierdzenie}
\begin{proof}[Wyjaśnienie konstrukcji]
Złożenie lewostronnie dokładnych funktorów
$\Gamma(Y,-)\circ f_*=\Gamma(X,-)$ prowadzi do
ciągu Grothendiecka. Dla snopów injektywnych
$f_*$ jest acykliczny względem $\Gamma(Y,-)$:
$\operatorname{Hom}(-,f_*I)\cong
\operatorname{Hom}(f^*(-),I)$ i dokładność $f^*$
pokazują, że $f_*I$ jest injektywny.
Reszta konstrukcji ciągu Grothendiecka to wynik
z algebry homologicznej rozwinięty wcześniej.
\end{proof}

\section{Cykle pobliskie i zanikające: konstrukcja z monodromią}
\label{sec:xvi-nearby}
Niech $f:X\to\Delta$ będzie holomorficzną funkcją,
$X_0=f^{-1}(0)$, $X^*=f^{-1}(\Delta^*)$, a
$i:X_0\hookrightarrow X$ oraz $j:X^*\hookrightarrow X$
odpowiednimi włożeniami. Oznaczmy przez
$\pi:\widetilde{\Delta^*}\to\Delta^*$ uniwersalne nakrycie
i przez $\widetilde X^*=X^*\times_{\Delta^*}
\widetilde{\Delta^*}$ odpowiednią przestrzeń. Niech
$\widetilde j:\widetilde X^*\to X$ będzie naturalną mapą.

\begin{definicja}[Cykle pobliskie i zanikające]
Dla $A\in D^b_c(X)$ określamy
$R\psi_fA=i^*R\widetilde j_*\widetilde j^*A$
na $X_0$; w tym zapisie $\widetilde j^*A$ oznacza
pullback $A$ na $\widetilde X^*$.
Morfizm specjalizacji $\mathrm{sp}:i^*A\to R\psi_fA$
wyznacza kompleks cykli zanikających
$R\phi_fA=\operatorname{Cone}(\mathrm{sp})$.
Otrzymujemy trójkąt
\[
 i^*A\xrightarrow{\mathrm{sp}}R\psi_fA
 \longrightarrow R\phi_fA\longrightarrow i^*A[1].
\]
Deck transformation nakrycia indukuje operator monodromii
na $R\psi_fA$ i $R\phi_fA$.
\end{definicja}
\begin{uwaga}[Znaczenie konwencji]
W niektórych książkach używa się przesuniętego
kompleksu $R\phi_f[-1]$ lub normalizacji perwersyjnej.
W naszym tekście pozostajemy przy stożku bez przesunięcia,
spójnym z Rozdziałem~\ref{ch:xiv-hodge-monodromia}.
\end{uwaga}
\begin{stwierdzenie}[Włókno Milnora jako kiełek cykli pobliskich]
\label{prop:xvi-milnor-stalk}
Jeśli $x\in X_0$ jest punktem, w którym $f$
ma izolowaną osobliwość krytyczną, to dla małej kuli
$B_\varepsilon(x)$ i dostatecznie małego $t\ne0$
kiełek $(R^q\psi_fk_X)_x$ można zidentyfikować
z $H^q(B_\varepsilon(x)\cap f^{-1}(t);k)$.
\end{stwierdzenie}
\begin{proof}[Status]
Wybór promieni i parametrów jest taki jak w
Twierdzeniu Milnora: najpierw ustala się
$\varepsilon$, a następnie odpowiednio małe $|t|$.
Identyfikacja wykorzystuje lokalną trywialność
fibracji Milnora i odpowiednie porównanie
pochodnego obrazu na uniwersalnym nakryciu.
Pełny dowód tego twierdzenia porównawczego byłby
powtórzeniem znacznej części teorii Milnora;
przyjmujemy go jawnie jako twierdzenie pomocnicze.
\end{proof}
\begin{przyklad}[Węzeł $zw=t$]
Dla $f(z,w)=zw$ włókno Milnora w małej kuli jest
homotopijnie równoważne $S^1$, zatem
$(R^0\psi_fk_X)_0\cong k$ i
$(R^1\psi_fk_X)_0\cong k$, a pozostałe grupy znikają.
Morfizm $i^*k_X\to R\psi_fk_X$ jest izomorfizmem
w stopniu zerowym, więc
$(R^1\phi_fk_X)_0\cong k$, pozostałe grupy znikają.
Lokalna monodromia na $H^1(S^1;k)$ jest tożsamością,
co nie przeczy nietrywialnej transwekcji na
$H^1$ globalnego gładkiego włókna zwartej rodziny.
\end{przyklad}
\begin{przyklad}[Rozgałęzienie $z^2=t$]
Dla $f:\C\to\C$, $f(z)=z^2$, włókno Milnora w
punkcie $0$ składa się z dwóch punktów.
Otrzymujemy $(R^0\psi_fk_\C)_0\cong k^2$, a
morfizm specjalizacji w stopniu zerowym jest przekątną
$k\to k^2$. Zatem $(R^0\phi_fk_\C)_0\cong k$.
Po okrążeniu $t=0$ dwa punkty włókna zamieniają się,
więc na zredukowanej $H^0$ monodromia działa przez $-1$
(gdy $\operatorname{char}k=2$, znak ten jest równy $+1$).
\end{przyklad}

\section{Lokalny schemat zależności i ograniczenia formalizmu}
\begin{center}
\begin{tikzpicture}[node distance=12mm and 13mm,>=Stealth,
 box/.style={draw,rounded corners,align=center,inner sep=5pt,font=\small}]
\node[box] (a) {$A\in D^b_c(X)$};
\node[box,below left=of a] (b) {$Rf_*,\ Rf_!$};
\node[box,below right=of a] (c) {$f^*,\ f^!$};
\node[box,below=25mm of a] (d) {$\mathbb D_X$ i $\omega_X$};
\node[box,below=of d] (e) {lokalizacja, zmiana podstawy\\cykle $R\psi_f,\ R\phi_f$};
\draw[->] (a)--(b); \draw[->] (a)--(c);
\draw[->] (b)--(d); \draw[->] (c)--(d); \draw[->] (d)--(e);
\end{tikzpicture}
\par\smallskip\small Powiązania sześciu operacji. Tensor i wewnętrzny Hom
są obecne w formule projekcji i definicji dualności.
\end{center}

\begin{uwaga}[Topologia klasyczna a topologia \emph{étale}]
W tym rozdziale pracujemy ze snopami na klasycznej przestrzeni
analitycznej nad $\C$. W geometrii algebraicznej podobne formuły
występują dla snopów \emph{étale}, lecz ich dowody i ograniczenia
(w szczególności warunki na torsję i charakterystykę) są inne.
Samo podobieństwo symboli $f^!,Rf_!$ nie upoważnia do
przeniesienia twierdzenia bez sprawdzenia kategorii współczynników.
\end{uwaga}
\begin{uwaga}[Precyzja stosowania funktorów]
Dla dowolnego kompleksu $A$ wyrażenie $R^qf_*A$ jest
skrótowe i bywa niejednoznaczne; używamy
$\mathcal H^q(Rf_*A)$. Izomorfizmy w rodzaju
$g^*Rf_!\simeq Rf'_!(g')^*$ lub
$\mathbb D_YRf_!\simeq Rf_*\mathbb D_X$ odnosimy do
kategorii i hipotez wskazanych przed twierdzeniami,
nie do wszystkich odwzorowań między dowolnymi
przestrzeniami topologicznymi. Przy pracy ze snopami
koherentnymi na schematach te same symbole mogą oznaczać
inne operacje, zwłaszcza pochodny pullback tensorowy.
\end{uwaga}
\begin{uwaga}[Status dowodów]
W rozdziale samodzielnie uzasadniono długi ciąg stożka,
krótki ciąg lokalizacyjny, podstawową adjunkcję $f^*\dashv Rf_*$,
przypadek punktu na prostej, wersję projekcji dla
lokalnych systemów oraz obliczenia dla $x^2$, $z^2$ i $zw$.
Istnienie ogólnego $f^!$, twierdzenie o zmianie podstawy,
pełna dualność Verdiera, ogólny wzór projekcji i formalne
własności cykli pobliskich wymagają niezależnych wyników
formalizmu sześciu operacji. Zostały przyjęte jawnie,
nie jako konsekwencje elementarnych rachunków.
\end{uwaga}

\chapter{Topologia gładkich zespolonych odmian projektowych: twierdzenia Lefschetza, przecięcia i rozdmuchania}
\label{ch:xvii-topologia-odmian}

\section{Motywacja, obiekty i zależności logiczne}
Wcześniejsze rozdziały dostarczyły homologii singularnej, dualności
Poincarégo, klas charakterystycznych, struktur Hodge'a, teorii snopów
oraz monodromii. Teraz pytamy, jakie \emph{ograniczenia topologiczne}
narzuca pochodzenie przestrzeni z geometrii zespolonej algebraicznej.
W odróżnieniu od samej rozmaitości gładkiej, gładka odmiana projektowa
nosi wyróżnioną klasę kohomologii stopnia dwa, co prowadzi do
silnych zależności między jej grupami kohomologii i formami przecięcia.

Przez $X\subset\mathbb P^N_\C$ rozumiemy odtąd niepustą,
gładką, \emph{spójną} zespoloną odmianę projektową wymiaru zespolonego
$n$. Rozważamy jej analityfikację jako zwartą rozmaitość
rzeczywistą wymiaru $2n$, z orientacją zespoloną. Współczynniki
całkowite oznaczamy jawnie przez $\Z$; dla rozkładu Hodge'a używamy
$\C$, a w twierdzeniu o silnym Lefschetzu początkowo $\Q$ lub $\R$.
Nie utożsamiamy odmian osobliwych z rozmaitościami gładkimi: dla
pierwszych właściwym narzędziem bywa homologia przecięć.

\begin{definicja}[Klasa hiperpłaszczyzny i operator Lefschetza]
Niech $\mathcal O_X(1)$ oznacza ograniczenie wiązki hiperpłaszczyznowej
z $\mathbb P^N_\C$ do $X$, a $h=c_1(\mathcal O_X(1))\in
H^2(X;\Z)$. Niech $\omega$ będzie formą Kählera, której klasa
de Rhama jest obrazem $h$ w $H^2(X;\R)$ (po przyjęciu stosownej
normalizacji krzywizny). Definiujemy
$L:H^k(X;\Q)\to H^{k+2}(X;\Q)$ przez $L(\alpha)=h\smile\alpha$;
na kohomologii rzeczywistej jest to mnożenie przez $[\omega]$.
\end{definicja}
\begin{uwaga}
Gdy $h$ ma znaczenie jako całkowita klasa Cherna, pozostawiamy
zapis $h$. Gdy odwołujemy się do analizy Kählerowskiej,
używamy gładkiej formy $\omega$. Przejście między nimi wymaga
porównania kohomologii singularnej i de Rhama, przedstawionego
wcześniej w skrypcie.
\end{uwaga}

\section{Słabe twierdzenie Lefschetza o hiperpłaszczyźnie}
Niech $H\subset\mathbb P^N_\C$ będzie hiperpłaszczyzną taką,
że $Y=X\cap H$ jest gładkie i niepuste. Z twierdzenia Bertiniego
wiemy, że taki $H$ można wybrać z niepustego zbioru Zariski-otwartego
w przestrzeni dualnej (dla $n\ge1$ i odpowiedniego osadzenia).

\begin{twierdzenie}[Słabe twierdzenie Lefschetza]
\label{thm:xvii-weak-lefschetz}
Jeżeli $X\subset\mathbb P^N_\C$ jest gładką spójną odmianą
projektową wymiaru zespolonego $n\ge2$, a $Y=X\cap H$ jest gładkim
przekrojem hiperpłaszczyznowym, to dla każdego $q\in\Z$
restrykcja
$i^*:H^q(X;\Z)\to H^q(Y;\Z)$ jest izomorfizmem dla
$q<n-1$ oraz monomorfizmem dla $q=n-1$.
Ponadto $\pi_j(Y)\to\pi_j(X)$ jest izomorfizmem dla
$1\le j<n-1$ oraz epimorfizmem dla $j=n-1$ (przy ustalonych
punktach bazowych). Wersje dla współczynników w ciele
wynikają z tego samego twierdzenia.
\end{twierdzenie}
\begin{proof}[Redukcja do twierdzenia o dołączaniu komórek]
Podstawą dowodu jest następujące \emph{twierdzenie geometryczne
Lefschetza}: $X$ ma typ homotopijny przestrzeni otrzymanej z $Y$
przez dołączenie komórek rzeczywistego wymiaru co najmniej $n$.
Jego dowód wymaga teorii funkcji plurisubharmonicznych i
zespolonej teorii Morse'a (lub snopowego twierdzenia o zanikaniu).
Nie utożsamiamy tego głębokiego kroku z elementarną konsekwencją
samej gładkości $X$.

Z twierdzenia o dołączaniu komórek wynika
$H_j(X,Y;\Z)=0$ dla $j<n$: w komórkowym kompleksie
względnym nie ma bowiem generatorów w tych stopniach.
Względna wersja twierdzenia o aproksymacji komórkowej daje
również $\pi_j(X,Y)=0$ dla $j<n$.
Z długiego ciągu dokładnego homotopii pary wynika odpowiednie
twierdzenie o $\pi_j(Y)\to\pi_j(X)$.

Dla kohomologii korzystamy z długiego ciągu dokładnego pary
$\cdots\to H^q(X,Y;\Z)\to H^q(X;\Z)
\xrightarrow{i^*}H^q(Y;\Z)\to H^{q+1}(X,Y;\Z)\to\cdots$.
Ponieważ względny kompleks komórkowy nie ma generatorów
w stopniach mniejszych niż $n$, mamy
$H^q(X,Y;\Z)=0$ dla $q<n$. Gdy $q<n-1$, znikają
obie grupy względne sąsiadujące z $i^*$, co daje izomorfizm;
gdy $q=n-1$, znika grupa po lewej i otrzymujemy injektywność.
Tym samym cała implikacja topologiczna jest dowiedziona;
niezależny dowód twierdzenia geometrycznego o dołączaniu
komórek pozostaje wskazanym wynikiem zewnętrznym.
\end{proof}

\begin{przyklad}[Przekrój powierzchni projektowej]
Dla $X=\mathbb P^2_\C$ i ogólnej prostej projektowej
$Y\cong\mathbb P^1_\C$ mamy $n=2$.
Słabe twierdzenie zapewnia izomorfizm w stopniu $0$ oraz
injektywność $H^1(X;\Z)\hookrightarrow H^1(Y;\Z)$.
\emph{Nie} twierdzi natomiast, że $H^2(X;\Z)\to H^2(Y;\Z)$
jest izomorfizmem, choć w tym konkretnym przykładzie oba
pierścienie mają naturalne generatory w stopniu dwa i restrykcja
jest izomorfizmem. To rozróżnienie zapobiega błędnemu
poszerzaniu zakresu twierdzenia.
\end{przyklad}

\begin{przyklad}[Krzywa jako przekrój powierzchni]
Jeśli $X\subset\mathbb P^3_\C$ jest gładką powierzchnią,
a $Y=X\cap H$ gładką krzywą, to $\pi_1(Y)\twoheadrightarrow\pi_1(X)$.
Nie oczekujemy na ogół izomorfizmu: krzywa może mieć
znacznie bogatszą grupę podstawową niż powierzchnia.
\end{przyklad}

\section{Silne twierdzenie Lefschetza i reprezentacje $\mathfrak{sl}_2$}
\begin{twierdzenie}[Silne twierdzenie Lefschetza]
\label{thm:xvii-hard-lefschetz}
Niech $X$ będzie gładką spójną zespoloną odmianą projektową
wymiaru $n$, a $h=c_1(\mathcal O_X(1))$.
Dla każdego $0\le k\le n$ mnożenie przez $h^{n-k}$ daje izomorfizm
\[
L^{n-k}:H^k(X;\Q)\xrightarrow{\ \simeq\ }H^{2n-k}(X;\Q).
\]
To samo zachodzi na kohomologii rzeczywistej dla każdej
klasy Kählerowskiej $[\omega]$.
\end{twierdzenie}

\begin{lemat}[Rachunek łańcuchów Lefschetza]
\label{lem:xvii-sl2}
Niech $V=\bigoplus_{k=0}^{2n}V^k$ będzie skończeniewymiarową
przestrzenią z dodatnio określonym iloczynem skalarnym,
$L:V^k\to V^{k+2}$, a $\Lambda=L^*$.
Załóżmy, że $[\Lambda,L]|_{V^k}=(n-k)\id$.
Wtedy dla $0\le k\le n$ operator
$L^{n-k}:V^k\to V^{2n-k}$ jest izomorfizmem.
\end{lemat}
\begin{proof}
Wprowadźmy $P^k=\ker(\Lambda:V^k\to V^{k-2})$.
Jeżeli $v\in P^k$ i $r\ge1$, to indukcja z komutatora daje
\[
\Lambda L^rv=r(n-k-r+1)L^{r-1}v.
\]
Ponieważ $\Lambda=L^*$, dla $v\ne0$ mamy
$\|L^rv\|^2=r(n-k-r+1)\|L^{r-1}v\|^2$.
Dla $r\le n-k$ współczynnik jest dodatni; przy
$r=n-k+1$ jest zerowy, zatem $L^{n-k+1}v=0$.
Ponadto każdy stopień dopuszcza prostopadły rozkład
$V^k=P^k\oplus LV^{k-2}$, gdyż $(\operatorname{im}L)^\perp=\ker\Lambda$.
Powtarzając rozkład (oraz korzystając z tożsamości komutatorów
w każdym stopniu), dostajemy sumę prostą
$V^k=\bigoplus_{r\ge0}L^rP^{k-2r}$, przy czym
$P^m=0$ dla $m>n$: gdy $v\in P^m$ i $m>n$, to
$\|Lv\|^2=(n-m)\|v\|^2<0$ dla $v\ne0$, sprzeczność.
Każdy łańcuch $v,Lv,\ldots,L^{n-m}v$ ma dokładnie
jeden niezerowy wektor w stopniach $m+2r$,
$0\le r\le n-m$. Obracając go działaniem $L^{n-k}$
na stopniu $k$ otrzymujemy bijekcję pomiędzy
odpowiednimi składnikami w stopniach $k$ i $2n-k$.
To dowodzi twierdzenia.
\end{proof}
\begin{proof}[Dowód silnego Lefschetza przy użyciu tożsamości Kählerowskich]
Rozważmy rzeczywistą przestrzeń harmonicznych form
$V^k=\mathcal H^k(X)$ dla ustalonej metryki Kählerowskiej.
Operator $L\alpha=\omega\wedge\alpha$ zachowuje formy
harmoniczne, podobnie jak jego formalny adjunkt $\Lambda$.
Istotnym analitycznym faktem są \emph{tożsamości Kählerowskie}
$[\Delta,L]=[\Delta,\Lambda]=0$ oraz
$[\Lambda,L]|_{\mathcal H^k}=(n-k)\id$.
Pełne wyprowadzenie ich w lokalnej ramce unitarnej należy do
analitycznej teorii Kählera; w tym rozdziale traktujemy je jako
jawnie przyjęty lemat geometryczny.
Lemat~\ref{lem:xvii-sl2} daje teraz izomorfizm
$L^{n-k}:\mathcal H^k(X)\to\mathcal H^{2n-k}(X)$.
Twierdzenie Hodge'a o reprezentantach harmonicznych
przenosi go na $H^k(X;\R)\to H^{2n-k}(X;\R)$.
Jeżeli $[\omega]=h$ jest klasą całkowitą,
macierz $L^{n-k}$ w bazach wymiernych ma współczynniki
w $\Q$. Jej wyznacznik nie znika nad $\R$, więc nie znika
nad $\Q$, co daje izomorfizm na kohomologii wymiernej.
\end{proof}

\begin{wniosek}[Symetria liczb Bettiego]
Dla gładkiej spójnej odmiany projektowej $X$ wymiaru $n$
$\dim_\Q H^k(X;\Q)=\dim_\Q H^{2n-k}(X;\Q)$.
Ponadto $b_{2j+1}(X)$ jest liczbą parzystą.
\end{wniosek}
\begin{proof}
Pierwsza równość wynika z silnego twierdzenia Lefschetza
(dla $k\le n$) i symetrii.
Z rozkładu Hodge'a $H^{2j+1}(X;\C)=
\bigoplus_{p+q=2j+1}H^{p,q}(X)$ i sprzężenia
$H^{p,q}=\overline{H^{q,p}}$ wynika, że suma wymiarów
jest parzysta, ponieważ $p\ne q$ przy $p+q$ nieparzystym.
\end{proof}

\section{Prymitywna kohomologia i relacje Hodge'a--Riemanna}
\begin{definicja}[Kohomologia prymitywna]
Dla $0\le k\le n$ definiujemy
$P^k(X;\C)=\ker(L^{n-k+1}:H^k(X;\C)\to H^{2n-k+2}(X;\C))$.
Korzystając z tożsamości Kählerowskich, można równoważnie
charakteryzować te klasy przez harmoniczne reprezentanty
$\alpha$ spełniające $\Lambda\alpha=0$.
Dla $p+q=k$ piszemy $P^{p,q}=P^k(X;\C)\cap H^{p,q}(X)$.
\end{definicja}
\begin{twierdzenie}[Rozkład Lefschetza]
\label{thm:xvii-decomposition}
Dla $0\le k\le n$ mamy rozkład
$H^k(X;\C)=\bigoplus_{r\ge0}L^rP^{k-2r}(X;\C)$,
przy czym składniki z ujemnym indeksem są zerowe.
\end{twierdzenie}
\begin{proof}
Lemat~\ref{lem:xvii-sl2} zastosowany do harmonicznych form
dostarcza rozkładu ortogonalnego na łańcuchy pierwotne.
Ponieważ $\omega$ jest typu $(1,1)$, mnożenie przez
$\omega$ zachowuje dwustopniowość Hodge'a, więc
rozłożenie jest zgodne z bigradowaniem.
\end{proof}
\begin{twierdzenie}[Relacje Hodge'a--Riemanna]
\label{thm:xvii-HR}
Dla $X$ i $\omega$ jak wyżej, $k=p+q\le n$ oraz
niezerowej klasy $\alpha\in P^{p,q}$ zachodzi
\[
(-1)^{k(k-1)/2}\,\mathrm i^{p-q}
\int_X\alpha\wedge\overline\alpha\wedge\omega^{n-k}>0.
\]
Całkę rozumiemy przez harmoniczne reprezentanty,
a wynik po lewej jest rzeczywistą liczbą dodatnią.
\end{twierdzenie}
\begin{proof}[Dodatniość z tożsamości gwiazdki Hodge'a]
Klasa prymitywna posiada harmonicznego reprezentanta
prymitywnego $a$ typu $(p,q)$. W ortonormalnej ramce
unitarnej dla metryki Kählerowskiej zachodzi punktowa
tożsamość
\[
*\overline a=
\frac{(-1)^{k(k-1)/2}\mathrm i^{p-q}}{(n-k)!}
\overline a\wedge\omega^{n-k},
\]
przy konwencji $\omega=\mathrm i\sum_j\theta^j\wedge
\overline{\theta^j}$ (z odpowiednio znormalizowaną ramką).
Tożsamość uzasadnia się sprawdzeniem na monomialnych
formach prymitywnych i liniowością; jest to potrzebny
lemat liniowej algebry hermitowskiej. Stąd
$(-1)^{k(k-1)/2}\mathrm i^{p-q}\int_X a\wedge
\overline a\wedge\omega^{n-k}=(n-k)!\int_X a\wedge *\overline a>0$.
Ostatnia całka jest normą $L^2$ niezerowej formy.
\end{proof}
\begin{uwaga}[Ścisły zakres dowodu]
Tożsamości Kählerowskie oraz punktową tożsamość dla
prymitywnej formy przyjmujemy jako lematy analitycznej
teorii Kählera; powyższy argument wyprowadza z nich
relacje i znaki. Nie jest to samodzielny dowód całej
teorii Hodge'a i twierdzenia Hodge'a--Riemanna od podstaw.
\end{uwaga}

\section{Formy przecięcia i twierdzenie o sygnaturze powierzchni}
\begin{definicja}[Forma przecięcia]
Dla gładkiej spójnej zespolonej powierzchni projektowej $S$
(przez \emph{powierzchnię} rozumiemy tutaj wymiar zespolony $2$)
definiujemy na $H^2(S;\R)$ formę
$Q_S(\alpha,\beta)=\int_S\alpha\wedge\beta$.
Z dualności Poincarégo jest ona niedegenerowana;
na $H^2(S;\Z)/\operatorname{Tor}$ wyznacza całkowitą
formę unimodularną.
\end{definicja}
\begin{twierdzenie}[Twierdzenie Hodge'a o indeksie]
\label{thm:xvii-hodge-index}
Jeśli $S$ jest gładką, spójną zespoloną powierzchnią projektową,
to sygnatura rzeczywistej formy $Q_S$ wynosi
\[
\bigl(b_2^+(S),b_2^-(S)\bigr)
=\bigl(1+2h^{2,0}(S),\ h^{1,1}(S)-1\bigr),
\]
a zatem $\sigma(S)=2+2h^{2,0}(S)-h^{1,1}(S)$.
W szczególności przecięcie na rzeczywistej przestrzeni
klas $(1,1)$ ma jeden kierunek dodatni, a na części
ortogonalnej do klasy Kählerowskiej jest ujemnie określone.
\end{twierdzenie}
\begin{proof}
Rozkład Hodge'a daje
$H^2(S;\C)=H^{2,0}\oplus H^{1,1}\oplus H^{0,2}$.
Przez $L$ oznaczamy mnożenie przez $[\omega]$.
Dekompozycja Lefschetza w stopniu $2$ przyjmuje postać
$H^2(S;\R)=\R[\omega]\oplus P^2(S;\R)$.
Ponieważ $\int_S\omega^2>0$, pierwsza prosta jest dodatnia.
Dla $a\in P^{1,1}\cap H^2(S;\R)$ relacja
Hodge'a--Riemanna w stopniu $k=2$ mówi
$-\int_S a\wedge a>0$ dla $a\ne0$.
Dla niezerowej klasy $b\in H^{2,0}$ mamy
$\int_S b\wedge\overline b>0$ (przypadek $k=2$, $p=2,q=0$).
Dodatniość na rzeczywistej płaszczyźnie
rozpiętej przez $\Re b,\Im b$ wynika stąd, że
$\int_S b\wedge b=0$ z powodów typu, a
$Q(\Re b,\Re b)=Q(\Im b,\Im b)
=\frac12\int_S b\wedge\overline b>0$ i składniki
mieszane są zerowe. Z powodów typu część $H^{1,1}$ jest ortogonalna do
$H^{2,0}\oplus H^{0,2}$ względem $Q$, natomiast
$H^{2,0}$ i $H^{0,2}$ \emph{nie} są wzajemnie ortogonalne:
parują ze sobą w całce $\int_S b\wedge\overline b$.
Na rzeczywistej przestrzeni $(H^{2,0}\oplus H^{0,2})\cap
H^2(S;\R)$ forma $Q$ jest dodatnio określona,
ponieważ każdy jej element można zapisać jako $b+\overline b$
i $Q(b+\overline b,b+\overline b)=2\int_S b\wedge\overline b>0$
dla $b\ne0$.
Liczba dodatnich kierunków to $1+2h^{2,0}$,
a ujemnych $h^{1,1}-1$, co kończy dowód.
\end{proof}
\begin{przyklad}[Płaszczyzna projektowa]
Dla $S=\mathbb P^2_\C$ mamy
$H^*(S;\Z)=\Z[h]/(h^3)$, $\deg h=2$,
$\int_S h^2=1$, stąd forma przecięcia na
$H^2(S;\Z)$ jest macierzą $(1)$ i $\sigma(S)=1$.
W tym przykładzie $h^{2,0}=0$ i $h^{1,1}=1$.
\end{przyklad}

\begin{przyklad}[Iloczyn dwóch prostych projektowych]
Niech $S=\mathbb P^1_\C\times\mathbb P^1_\C$.
Dla $a=\mathrm{pr}_1^*h$ i $b=\mathrm{pr}_2^*h$ mamy
$H^2(S;\Z)=\Z a\oplus\Z b$, $a^2=b^2=0$ oraz
$\int_S a\wedge b=1$ (formuła Künnetha).
Macierz przecięcia w bazie $(a,b)$ jest więc
$\begin{psmallmatrix}0&1\\1&0\end{psmallmatrix}$,
a jej wartości własne wynoszą $1$ i $-1$.
Ponieważ $h^{2,0}(S)=0$, $h^{1,1}(S)=2$,
twierdzenie o indeksie Hodge'a daje ten sam podpis $(1,1)$.
\end{przyklad}

\section{Rozdmuchanie w punkcie: geometria i topologia}
\begin{definicja}[Rozdmuchanie w punkcie]
Niech $S$ będzie gładką zespoloną powierzchnią i $p\in S$.
Rozdmuchaniem $\pi:\widetilde S=\operatorname{Bl}_pS\to S$
nazywamy odmianę, która jest izomorficzna z $S\setminus\{p\}$
poza dzielnikiem wyjątkowym $E=\pi^{-1}(p)\cong\mathbb P^1_\C$.
Lokalnie nad kartą ze współrzędnymi $(x,y)$ wokół $p$ jest to
pododmiana $\{((x,y),[s:t]):xt=ys\}$ w
$\C^2\times\mathbb P^1$ z naturalnym rzutowaniem.
\end{definicja}
\begin{lemat}[Wiązka normalna dzielnika wyjątkowego]
\label{lem:xvii-exceptional}
Dla rozdmuchania gładkiej powierzchni w punkcie zachodzi
$N_{E/\widetilde S}\simeq\mathcal O_{\mathbb P^1}(-1)$.
W konsekwencji $E\cdot E=-1$.
\end{lemat}
\begin{proof}
W pierwszej karcie rozdmuchania piszemy $y=ux$,
gdzie współrzędnymi są $(x,u)$, a $E$ jest dany przez $x=0$.
W drugiej karcie piszemy $x=vy$, gdzie współrzędnymi są
$(v,y)$ i $E=\{y=0\}$. Na przecięciu mamy $v=1/u$
oraz $y=ux$; przejście dla współrzędnej normalnej
ma więc postać $y=ux$, co jest funkcją przejścia
wiązki tautologicznej $\mathcal O(-1)$ w przyjętej konwencji
kart $\mathbb P^1$. Pierwsza klasa Cherna ma stopień $-1$.
Ponieważ $E\cdot E=\langle c_1(N_{E/\widetilde S}),[E]\rangle$,
dostajemy $E^2=-1$.
\end{proof}
\begin{twierdzenie}[Kohomologia rozdmuchania powierzchni]
\label{thm:xvii-blowup}
Jeżeli $S$ jest gładką zwartą spójną zespoloną powierzchnią,
a $\widetilde S=\operatorname{Bl}_pS$, to
\[
H^2(\widetilde S;\Z)\cong \pi^*H^2(S;\Z)\oplus\Z e,
\quad e=\operatorname{PD}_{\widetilde S}[E],
\]
przy czym $\int_{\widetilde S}e^2=-1$ i
$\int_{\widetilde S}\pi^*\alpha\wedge e=0$
dla każdego $\alpha\in H^2(S;\Z)$.
W innych stopniach grupy kohomologii nie zmieniają się.
\end{twierdzenie}
\begin{proof}
Usuwamy z $S$ małą zamkniętą kulę $B^4$ wokół $p$.
Jej przeciwobraz w $\widetilde S$ jest otoczeniem dyskowym
dzielnika wyjątkowego, którego retraktem jest $E\cong S^2$;
na brzegach obu wyciętych części otrzymujemy tę samą
sferę $S^3$ (z odpowiednią orientacją).
Niech $U=S\setminus\operatorname{int}(B^4)$.
Przestrzenie $S=U\cup B^4$ i
$\widetilde S=U\cup\widetilde B$
mają wspólne dane na przecięciu, homotopijnie równoważnym
$S^3$. Ciąg Mayera--Vietorisa zawiera w stopniu drugim
$H^1(S^3;\Z)\to H^2(S)\to H^2(U)\oplus H^2(B^4)
\to H^2(S^3;\Z)$; skrajne grupy i $H^2(B^4;\Z)$
są zerowe. Stąd $H^2(S;\Z)\cong H^2(U;\Z)$.
Analogicznie $H^2(\widetilde S;\Z)\cong
H^2(U;\Z)\oplus H^2(\widetilde B;\Z)$, przy czym
$H^2(\widetilde B;\Z)\cong H^2(E;\Z)\cong\Z$.
W stopniu trzecim trzeba sprawdzić mapę
$H^3(U;\Z)\to H^3(\partial U;\Z)$.
Jest ona zerowa: w długim ciągu pary $(U,\partial U)$
homomorfizm łączący
$H^3(\partial U;\Z)\to H^4(U,\partial U;\Z)$
jest izomorfizmem, ponieważ przesyła generator
orientacji brzegu na względną klasę orientacji $U$.
Ciąg dokładny daje więc zerową mapę ograniczenia.
Stąd $H^3(S;\Z)\cong H^3(U;\Z)\cong H^3(\widetilde S;\Z)$.
W stopniu czwartym obie przestrzenie mają grupę $\Z$
ze względu na orientację, a w stopniach zerowym i pierwszym
porównanie jest natychmiastowe.
Z identyfikacji przez $\pi^*$ i z rozszczepienia
na części pochodzące z $S$ i z $E$ otrzymujemy
wskazany rozkład całkowity.
Ponieważ $\pi(E)=p$, dla $\alpha\in H^2(S;\Z)$
z wzoru projekcji i definicji dualności Poincarégo mamy
$\int_{\widetilde S}\pi^*\alpha\wedge e
=\int_E(\pi|_E)^*\alpha=0$.
Samoprzecięcie $e^2=-1$ wynika z Lematu~\ref{lem:xvii-exceptional}.
\end{proof}
\begin{wniosek}[Zmiana sygnatury i charakterystyki Eulera]
Dla $S$ jak wyżej:
$\sigma(\widetilde S)=\sigma(S)-1$,
$\chi(\widetilde S)=\chi(S)+1$ oraz
$b_2(\widetilde S)=b_2(S)+1$.
\end{wniosek}
\begin{proof}
Forma przecięcia rozszczepia się ortogonalnie jako
$Q_S\oplus(-1)$, a z formuły grup kohomologii
nową klasę otrzymujemy tylko w stopniu drugim.
\end{proof}
\begin{przyklad}[Rozdmuchanie $\mathbb P^2$]
Dla $S=\mathbb P^2_\C$ i
$\widetilde S=\operatorname{Bl}_p\mathbb P^2$ otrzymujemy bazę
$(H,E)$ dla $H^2(\widetilde S;\Z)$ modulo torsję,
gdzie $H=\pi^*h$, $E=e$, i macierz przecięcia
\[
\begin{pmatrix}1&0\\0&-1\end{pmatrix}.
\]
Zatem $b_2=2$, $\sigma=0$, $\chi=4$.
Klasa $H$ jest całkowita, ale nie jest klasą Kählerowską
na rozdmuchaniu: jej całka po krzywej wyjątkowej wynosi zero.
Dla odpowiednio małego $\varepsilon>0$ klasa
$H-\varepsilon E$ jest Kählerowska.
\end{przyklad}

\section{Schemat geometryczny rozdmuchania}
\begin{center}
\begin{tikzpicture}[>=Stealth,scale=1,
  every node/.style={font=\small}]
  \draw[thick,rounded corners=8pt] (0,0) rectangle (4,2.3);
  \fill (2,1.13) circle (2pt);
  \node[below] at (2,1.0) {$p$};
  \node at (2,2.58) {$S$ (lokalnie)};
  \draw[thick,rounded corners=8pt] (7,0) rectangle (11,2.3);
  \draw[thick,manifoldblue] (9,1.28) ellipse (0.86 and 0.38);
  \draw[thick,manifoldblue,dashed] (8.14,1.28) arc (180:360:0.86 and 0.38);
  \node[below] at (9,0.63) {$E\cong\mathbb P^1$};
  \node at (9,2.58) {$\widetilde S$};
  \draw[->,very thick] (7,1.2) -- node[above] {$\pi$} (4,1.2);
  \node[font=\footnotesize,align=center] at (5.5,-0.55)
    {schemat: cały dzielnik $E$\ jest ściągany do $p$};
\end{tikzpicture}
\par\smallskip\small Rysunek poglądowy: $E$ jest zespoloną krzywą projektową,
a nie okręgiem w przestrzeni rzeczywistej.
\end{center}

\section{Dzielnik kanoniczny i charakterystyczne liczby Chernowskie}
\begin{twierdzenie}[Kanoniczna klasa rozdmuchania powierzchni]
\label{thm:xvii-canonical}
Jeżeli $S$ jest gładką zespoloną powierzchnią,
a $\pi:\widetilde S\to S$ jest rozdmuchaniem w punkcie,
to dla wiązek kanonicznych zachodzi
$K_{\widetilde S}\cong\pi^*K_S\otimes\mathcal O_{\widetilde S}(E)$.
Jeśli $S$ jest ponadto projektowa, to
$K_{\widetilde S}^2=K_S^2-1$.
\end{twierdzenie}
\begin{proof}
W karcie $y=ux$ różniczkowanie daje
$\dd x\wedge\dd y=x\,\dd x\wedge\dd u$.
Jacobian mapy $\pi$ ma zatem proste zero na $E$.
W drugiej karcie otrzymujemy zgodny rachunek, toteż
dzielnik ramifikacji ma krotność jeden wzdłuż $E$.
Formuła dla wiązki kanonicznej wynika z identyfikacji
lokalnych generatorów najwyższej potęgi wiązki kostycznej.
Po wzięciu pierwszej klasy Cherna mamy
$c_1(K_{\widetilde S})=\pi^*c_1(K_S)+e$.
Wymnażając kwadrat i korzystając z zerowania
$\pi^*c_1(K_S)\smile e$ oraz $e^2=-1$,
otrzymujemy $K_{\widetilde S}^2=K_S^2-1$.
\end{proof}

\section{Rozdmuchanie w gładkim centrum i formuła kohomologiczna}
\begin{twierdzenie}[Wzór na kohomologię rozdmuchania]
\label{thm:xvii-blowup-general}
Niech $X$ będzie gładką zespoloną odmianą projektową,
a $Z\subset X$ gładką domkniętą pododmianą o zespolonej
kodymensji $r\ge2$. Niech $\widetilde X=\operatorname{Bl}_Z X$,
$\pi:\widetilde X\to X$, $j:E\hookrightarrow\widetilde X$,
$p:E=\mathbb P(N_{Z/X})\to Z$ i
$\xi=c_1(\mathcal O_E(1))$.
Istnieje addytywny izomorfizm grup całkowitej kohomologii
\[
H^k(\widetilde X;\Z)\cong H^k(X;\Z)\oplus
\bigoplus_{a=1}^{r-1}H^{k-2a}(Z;\Z),
\]
realizowany przez $\pi^*$ oraz $j_*(\xi^{a-1}p^*(-))$,
gdzie $j_*$ jest homomorfizmem Gysina podnoszącym stopień o $2$.
Po przejściu na $\C$ rozkład jest zgodny ze strukturą Hodge'a,
przy czym nowy składnik $H^{p-a,q-a}(Z)$ przechodzi do
$H^{p,q}(\widetilde X)$.
\end{twierdzenie}
\begin{proof}[Wyprowadzenie z teorii par i formuły wiązki projektowej]
Mapy $\pi$ i $p$ dają diagram par $(\widetilde X,E)\to(X,Z)$.
Poza $E$ i $Z$ odwzorowanie jest izomorfizmem;
ponieważ $X$ i $\widetilde X$ są zwarte i Hausdorffa, odwzorowanie
$\pi$ indukuje homeomorfizm ilorazów
$\widetilde X/E\cong X/Z$. Istotnie, po ściągnięciu $E$ i $Z$
do jednego punktu dostajemy ciągłą bijekcję między przestrzeniami
zwartą i Hausdorffa; jest ona homeomorfizmem. Daje to
$H^k(X,Z;\Z)\cong H^k(\widetilde X,E;\Z)$.

Rozpatrzmy morfizm długich ciągów kohomologii par
$(X,Z)$ i $(\widetilde X,E)$; poziome mapy między grupami
względnymi są izomorfizmami, a na zamkniętych podprzestrzeniach
mamy $p^*:H^k(Z)\to H^k(E)$.
Twierdzenie Leraya--Hirscha dla $p:E\to Z$ zapewnia, że
$p^*$ jest monomorfizmem we wszystkich stopniach.
Lemat o wężu zastosowany do odpowiadających fragmentów
ciągów długich daje wtedy dokładny ciąg
\[
0\longrightarrow H^k(X;\Z)\xrightarrow{\pi^*}
H^k(\widetilde X;\Z)\xrightarrow{[j^*]}
H^k(E;\Z)/p^*H^k(Z;\Z)\longrightarrow0.
\]
W szczególności $\pi^*$ jest monomorfizmem, a $[j^*]$
oznacza restrykcję do $E$ po przejściu do ilorazu przez obraz
$p^*$. Dokładność tego ciągu, a nie jedynie formalna wzmianka
o pięciu homomorfizmach, jest potrzebnym krokiem porównania.

Z twierdzenia Leraya--Hirscha dla wiązki projektowej
$E=\mathbb P(N_{Z/X})\to Z$ mamy
$H^*(E;\Z)=\bigoplus_{a=0}^{r-1}\xi^a p^*H^{*-2a}(Z;\Z)$
addytywnie. Obraz $p^*$ jest jego składnikiem $a=0$.
Dla klas nowych na $\widetilde X$ wykorzystujemy
homomorfizm Gysina $j_*$ oraz tożsamość
$j^*j_*\beta=c_1(N_{E/\widetilde X})\smile\beta
=-\xi\smile\beta$.
W ilorazie $H^k(E)/p^*H^k(Z)$ klasy
$\xi^ap^*\beta$, $1\le a\le r-1$, tworzą dokładnie
pozostałe składniki rozkładu Leraya--Hirscha.
Stąd klasy $j_*(\xi^{a-1}p^*\beta)$ przechodzą przez
$[j^*]$ na $-\xi^ap^*\beta$. Zatem suma homomorfizmów
Gysina daje przekrój powyższego krótkiego ciągu dokładnego
(po zmianie znaku); wraz z $\pi^*$ otrzymujemy żądany
izomorfizm addytywny. Jest on określony przez jawne mapy.
Użytymi zewnętrznymi rezultatami pozostają formuła
Leraya--Hirscha dla wiązki projektowej, konstrukcja mapy
Gysina i wzór samoprzecięcia; nie ukrywamy tych zależności.
Zgodność typów Hodge'a wynika stąd, że $\pi^*,p^*$
są morfizmami struktur Hodge'a, $\xi$ ma typ $(1,1)$,
a $j_*$ podnosi oba indeksy o $1$.
\end{proof}
\begin{przyklad}[Rozdmuchanie w krzywej w trójwymiarowej odmianie]
Niech $X$ będzie gładką projektową trójwymiarową odmianą,
a $Z\subset X$ gładką spójną krzywą (zatem $r=2$).
Wtedy $H^3(\operatorname{Bl}_Z X;\Q)
\cong H^3(X;\Q)\oplus H^1(Z;\Q)$.
Jeśli $Z$ ma rodzaj $g$, to trzeci wymiar Bettiego rośnie o $2g$.
W szczególności operacja biracjonalna może zmienić środkową
kohomologię nawet wtedy, gdy nie zmienia dopełnienia centrum.
\end{przyklad}

\section{Powiązanie z teorią snopów perwersyjnych i monodromią}
Twierdzenia z rozdziałów~\ref{ch:ih-perverse-xiii},
\ref{ch:xiv-hodge-monodromia} i~\ref{ch:xvi-szesc-operacji}
dają alternatywne podejście do hiperpłaszczyzn i rozdmuchań.
Dla właściwego morfizmu algebraicznego $f$ kompleksy $Rf_*IC_X$
można badać przez twierdzenie o rozkładzie.
Gdy $f$ jest rozdmuchaniem gładkiego centrum,
dodatkowe składniki kohomologiczne są geometrycznie
pochodnymi obrazami dzielnika wyjątkowego; ich przesunięcia
odpowiadają stopniom nowo powstających klas.
Nie należy jednak utożsamiać wzoru kohomologicznego
na rozdmuchanie z pełnym twierdzeniem BBDG:
pierwszy daje tu konkretną dekompozycję,
a drugie jest znacznie ogólniejszym twierdzeniem
w kategorii pochodnej snopów konstruktybilnych.

\section{Dalsze przykłady i kontrola intuicji}
\begin{przyklad}[Czemu dodatniość zmienia się po rozdmuchaniu?]
Na $\operatorname{Bl}_p\mathbb P^2$ forma przecięcia w bazie
$(H,E)$ ma sygnaturę $(1,1)$. Dla klasy rzeczywistej
$\kappa=H-\varepsilon E$ z $0<\varepsilon<1$ mamy
$\kappa^2=1-\varepsilon^2>0$ oraz
$\kappa\cdot E=\varepsilon>0$.
To jest spójne z możliwością reprezentowania $\kappa$
przez formę Kählera; jednak sama dodatniość kwadratu
i dodatnie przecięcie z $E$ nie wystarczają w ogólności
do wykazania Kählerowskości. W tym przykładzie
Kählerowskość dla małego $\varepsilon$ wynika z osobnego
twierdzenia o klasach Kählerowskich rozdmuchań.
\end{przyklad}
\begin{przyklad}[Rozdzielenie trzech twierdzeń Lefschetza]
\emph{Słabe twierdzenie Lefschetza} porównuje kohomologię $X$ i gładkiego
przekroju $Y=X\cap H$. \emph{Silne twierdzenie Lefschetza} mówi o
izomorfizmach $L^{n-k}$ wewnątrz kohomologii $X$.
\emph{Względne silne twierdzenie Lefschetza} dla właściwych map algebraicznych
jest twierdzeniem o perwersyjnych snopach kohomologii
$Rf_*IC_X$; nie należy go używać bez warunków
właściwości, projektowości i normalizacji perwersyjnej.
\end{przyklad}

\section{Bilans dowodów i przygotowanie do wspólnej książki}
W tym rozdziale samodzielnie przeprowadziliśmy topologiczne
konsekwencje dołączania komórek, liniowoalgebraiczny lemat
Lefschetza, dowód indeksu Hodge'a z relacji
Hodge'a--Riemanna oraz obliczenia rozdmuchania w punkcie.
\emph{Jawnymi twierdzeniami pomocniczymi} pozostają:
Bertini, geometryczne twierdzenie Lefschetza o dołączaniu
komórek, tożsamości Kählerowskie, twierdzenie Hodge'a
o reprezentantach harmonicznych, punktowa tożsamość
prymitywnych form Hodge'a oraz formuła wiązki projektowej.
Ich pełne dowody wymagają teorii, której nie zastępujemy
kilkoma zdaniami.

W późniejszym połączeniu z geometrią algebraiczną
warto zachować w tym skrypcie \emph{skutki topologiczne},
a konstrukcję snopów koherentnych, formułę wiązki
projektowej w przecięciach algebraicznych i pełne
teorie dzielników wyjaśnić w tomie algebraicznym.

\chapter{Dalsze narzędzia homotopii: granice odwrotne, przeszkody, operacje i sklejanie Morse'a}
\label{ch:xviii-domkniecie}

\section{Cel, motywacja i poziom dowodów}
W poprzednich rozdziałach spotkaliśmy cztery zjawiska wymagające
szczególnej ostrożności: przechodzenie od skończonych przybliżeń do
przestrzeni granicznej; rozszerzanie przekrojów krok po kroku; zależności
między operacjami kohomologicznymi; oraz sklejanie trajektorii przy
złamanych liniach gradientowych. Łączy je pytanie, czy dane zgodne na
każdym skończonym etapie dają obiekt globalny. Odpowiedź może być
negatywna nawet wtedy, gdy wszystkie lokalne grupy przeszkód wydają się
proste.

Celem rozdziału jest ścisłe uzupełnienie tych czterech zagadnień.
Dowodzimy elementarnej teorii $\varprojlim^1$ dla przeliczalnych
układów grup abelowych i odpowiedniego ciągu Milnora dla układów
kompleksów łańcuchowych. Przeprowadzamy konstrukcję kolejnych
przeszkód komórkowych, wraz z kontrolą wpływu wyborów. Dla operacji
Steenroda wyprowadzamy kilka relacji z Bocksteina i podajemy dokładny
zakres twierdzenia Adema. Dla ekscyzji homotopijnej wyjaśniamy
mechanizm twierdzenia Blakersa--Masseya i jego zastosowanie do
zawieszenia. W końcu identyfikujemy precyzyjnie dodatkowy lemat
analityczny potrzebny do pełnego sklejania trajektorii Morse'a.

\section{Układy odwrotne i dwa rodzaje granic}
\begin{definicja}[Przeliczalny układ odwrotny]
Przez układ odwrotny grup abelowych rozumiemy ciąg
$A_0\xleftarrow{u_0}A_1\xleftarrow{u_1}A_2\leftarrow\cdots$,
gdzie każdy $A_j$ jest grupą abelową, a $u_j:A_{j+1}\to A_j$
homomorfizmem. Jego granicę odwrotną definiujemy jako
\[
\varprojlim_jA_j=
\{(a_j)_{j\ge0}\in\prod_{j\ge0}A_j:
 a_j=u_j(a_{j+1})\text{ dla każdego }j\}.
\]
Sytuację przedstawia schemat map przejścia:
\begin{center}
\begin{tikzpicture}[>=Stealth, node distance=2.0cm,
  every node/.style={font=\small}]
\node (a0) {$A_0$};
\node[right=of a0] (a1) {$A_1$};
\node[right=of a1] (a2) {$A_2$};
\node[right=of a2] (a3) {$\cdots$};
\draw[->] (a1) -- node[above] {$u_0$} (a0);
\draw[->] (a2) -- node[above] {$u_1$} (a1);
\draw[->] (a3) -- node[above] {$u_2$} (a2);
\end{tikzpicture}
\par\smallskip\small Rysunek: strzałki układu odwrotnego wskazują od dokładniejszych przybliżeń do wcześniejszych.
\end{center}
Na produkcie definiujemy operator
$T((a_j)_{j\ge0})=(u_j(a_{j+1}))_{j\ge0}$.
\end{definicja}
\begin{definicja}[Pierwszy pochodny funktor granicy odwrotnej]
Dla powyższego przeliczalnego układu grup abelowych kładziemy
\[
\varprojlim{}^1_jA_j:=
\operatorname{coker}\bigl(1-T:\prod_{j\ge0}A_j\to
\prod_{j\ge0}A_j\bigr).
\]
Nie jest to nowy typ granicy przestrzeni topologicznych:
$\varprojlim{}^1$ jest grupą mierzącą niezupełność dokładności
funktora $\varprojlim$ na takich układach.
\end{definicja}
\begin{lemat}[Jądro operatora różnicy]
Dla każdego przeliczalnego układu grup abelowych zachodzi
$\ker(1-T)=\varprojlim A_j$.
\end{lemat}
\begin{proof}
Równość $(1-T)(a)=0$ oznacza dokładnie układ równań
$a_j=u_j(a_{j+1})$ dla wszystkich $j$. Są to wszystkie i tylko
warunki definiujące granicę odwrotną.
\end{proof}

\begin{twierdzenie}[Warunek Mittaga--Lefflera dla układów przeliczalnych]
\label{thm:xviii-ml}
Niech $(A_j,u_j)$ będzie przeliczalnym układem grup abelowych.
Załóżmy, że dla każdego $j$ obrazy
$\operatorname{im}(A_k\to A_j)$ stabilizują się, gdy $k\to\infty$.
Wówczas $\varprojlim{}^1 A_j=0$.
W szczególności wniosek zachodzi, gdy wszystkie $u_j$ są surjektywne.
\end{twierdzenie}
\begin{proof}[Przypadek surjektywnych odwzorowań]
Pokażemy, że $1-T$ jest surjektywny. Ustalmy
$(b_j)\in\prod_jA_j$. Chcemy rozwiązać
$a_j-u_j(a_{j+1})=b_j$ dla wszystkich $j$.
Wybierzmy $a_0=0$. Jeśli $a_j$ jest już wybrany, to
surjektywność $u_j$ pozwala wybrać $a_{j+1}$ z równania
$u_j(a_{j+1})=a_j-b_j$. Indukcja daje szukany ciąg,
zatem kokernel $1-T$ jest zerowy.
\end{proof}
\begin{proof}[Redukcja ogólnego przypadku]
Dla każdego $j$ niech $B_j\subset A_j$ będzie ustabilizowanym
obrazem $A_k\to A_j$ dla dostatecznie dużych $k$.
Odwzorowania $B_{j+1}\to B_j$ są surjektywne: jeśli
$b\in B_j$, wybieramy $k$ większe od progów stabilizacji dla
$j$ i $j+1$ i element $a_k$ przechodzący na $b$. Jego obraz
w $B_{j+1}$ przechodzi na $b$. Zatem $\varprojlim{}^1 B_j=0$.

Dla układu ilorazów $Q_j=A_j/B_j$ każdy element $Q_k$
przechodzi na zero w $Q_j$ dla dostatecznie dużego $k$.
Taki układ nazywamy pro-zerowym. Dla ustalonego ciągu
$(q_j)\in\prod_j Q_j$ rozważmy równania
$x_j-v_j(x_{j+1})=q_j$. Możemy je rozwiązać przez sumę
skończoną w każdej współrzędnej:
$x_j=\sum_{k\ge j}v_{j,k}(q_k)$, gdzie
$v_{j,k}:Q_k\to Q_j$ jest odwzorowaniem przejścia
(i $v_{j,j}=\id$). Suma jest rzeczywiście skończona,
bo dla ustalonego $j$ mapy $v_{j,k}$ zerują się dla dużych $k$.
Wtedy $x_j-v_j(x_{j+1})=q_j$. Stąd $\varprojlim{}^1Q_j=0$.

Pozostaje użyć dokładności sześciowyrazowego ciągu dla
krótkiego ciągu układów $0\to B_j\to A_j\to Q_j\to0$.
Wynika ona wprost z lematu o wężu zastosowanego do
dwóch wierszy z produktami i pionowymi mapami $1-T$.
Ponieważ pochodne granice $B_j,Q_j$ są zerowe,
dokładność daje $\varprojlim{}^1A_j=0$.
\end{proof}
\begin{uwaga}
W powyższym dowodzie użyto produktów grup abelowych,
które zachowują dokładność krótkich ciągów. W ogólniejszych
kategoriach homologicznym warunkiem zerowania mogą być inne
hipotezy; symbol $\varprojlim{}^1$ nie jest wtedy automatycznie
definiowany przez ten sam elementarny iloraz.
\end{uwaga}

\section{Przykład podstawowy: $\mathbb Z_p/\mathbb Z$}
\begin{definicja}
Niech $p\ge2$ będzie liczbą pierwszą. Pierścieniem liczb
$p$-adycznych nazywamy granicę odwrotną
$\mathbb Z_p=\varprojlim_m\mathbb Z/p^m\mathbb Z$,
gdzie mapami przejścia są redukcje modulo $p^m$.
\end{definicja}
\begin{twierdzenie}[Niezerowe $\varprojlim{}^1$]
\label{thm:xviii-padic}
Dla układu $A_j=\Z$ z mapami $u_j(a)=pa$ zachodzą
izomorfizmy
\[
\varprojlim_j A_j=0,\qquad
\varprojlim{}^1_j A_j\cong\Z_p/\Z,
\]
gdzie $\Z$ zanurzamy w $\Z_p$ przez zwykłe redukcje
modulo potęg $p$.
\end{twierdzenie}
\begin{proof}
Jeśli $a_j=pa_{j+1}$ dla wszystkich $j$, to $a_0$
jest podzielne przez $p^m$ dla każdego $m$, zatem $a_0=0$.
Indukcyjnie otrzymujemy $a_j=0$.

Niech $P=\prod_{j\ge0}\Z$. Definiujemy homomorfizm
$\Theta:P\to\Z_p$ przez zbieżny szereg $p$-adyczny
$\Theta((a_j))=\sum_{j\ge0}p^ja_j$.
Surjektywność wynika z rozwinięcia każdej liczby $p$-adycznej
w cyfry $0,\ldots,p-1$.
Jeżeli $b=(1-T)a$, czyli $b_j=a_j-pa_{j+1}$,
to sumy częściowe teleskopują:
$\sum_{j=0}^{m}p^jb_j=a_0-p^{m+1}a_{m+1}$.
Ostatni wyraz dąży do zera $p$-adycznie, więc
$\Theta(b)=a_0\in\Z$. Dostajemy indukowane odwzorowanie
$\overline\Theta:\operatorname{coker}(1-T)\to\Z_p/\Z$.

Aby wykazać różnowartościowość, załóżmy, że
$\Theta(b)=q\in\Z$. Ustalmy $s_0=q$ i dla $j\ge0$ kładźmy
\[
s_j=\frac{q-\sum_{i=0}^{j-1}p^ib_i}{p^j}.
\]
Licznik jest podzielny przez $p^j$ w $\Z_p$, a ponieważ
jest liczbą całkowitą, jest podzielny przez $p^j$ w $\Z$.
Zatem $s_j\in\Z$ i $s_j-ps_{j+1}=b_j$.
Stąd $b=(1-T)s$, co dowodzi injektywności.
Surjektywność $\overline\Theta$ jest dziedziczona po $\Theta$.
\end{proof}
\begin{przyklad}[Dlaczego granice nie można przestawiać z homologią]
Weźmy kompleks $C_j$, skoncentrowany w stopniach $1,0$,
przy czym $C_{j,1}=\Z$, $C_{j,0}=\Z$, a różniczka jest mnożeniem
przez $p^j$. Bez określenia map łańcuchowych przejścia nie
otrzymujemy jeszcze układu odwrotnego kompleksów. Ponadto
surjektywność odwzorowań na poziomie \emph{homologii} nie jest tym
samym co surjektywność w każdym stopniu łańcuchowym.
Odpowiednie hipotezy pojawią się w następnym twierdzeniu.
\end{przyklad}

\section{Ciąg Milnora dla układów kompleksów}
\begin{twierdzenie}[Elementarny ciąg Milnora]
\label{thm:xviii-milnor}
Niech $\{C^{(j)}_*,u_j\}_{j\ge0}$ będzie przeliczalnym
układem odwrotnym kompleksów łańcuchowych grup abelowych,
a każde $u_j:C^{(j+1)}_q\to C^{(j)}_q$ niech będzie
surjektywne dla wszystkich $j,q$. Wówczas dla każdego $n\in\Z$
istnieje naturalny krótki ciąg dokładny
\[
0\longrightarrow\varprojlim{}^1_jH_{n+1}(C^{(j)}_*)
\longrightarrow H_n(\varprojlim_j C^{(j)}_*)
\longrightarrow\varprojlim_jH_n(C^{(j)}_*)
\longrightarrow0.
\]
Na ogół ten ciąg nie ma naturalnego rozszczepienia.
\end{twierdzenie}
\begin{proof}
Oznaczmy przez $P_*$ produkt kompleksów $\prod_j C^{(j)}_*$,
z różniczką działającą współrzędnie. Mapy łańcuchowe $u_j$
definiują endomorfizm $T:P_*\to P_*$. Surjektywność map
$u_j$ i argument indukcyjny użyty w twierdzeniu
\ref{thm:xviii-ml} pokazują, że $1-T:P_q\to P_q$
jest surjektywny w każdym stopniu. Jego jądro to z definicji
$\varprojlim C^{(j)}_q$. Otrzymujemy krótki ciąg dokładny
kompleksów $0\to\varprojlim C^{(j)}_*\to P_*
\xrightarrow{1-T}P_*\to0$.

Długi ciąg homologii zawiera fragment
\[
H_{n+1}(P_*)\xrightarrow{1-T}H_{n+1}(P_*)
\longrightarrow H_n(\varprojlim C^{(j)}_*)
\longrightarrow H_n(P_*)\xrightarrow{1-T}H_n(P_*).
\]
Ponieważ w kategorii grup abelowych produkty zachowują
jądra i obrazy różniczek, mamy
$H_q(P_*)\cong\prod_jH_q(C^{(j)}_*)$.
Jądro prawej mapy to $\varprojlim H_n(C^{(j)}_*)$,
a kokernel lewej to $\varprojlim{}^1 H_{n+1}(C^{(j)}_*)$.
Wyodrębnienie tego fragmentu daje żądany krótki ciąg.
\end{proof}
\begin{uwaga}[Czego twierdzenie nie mówi]
Nie wolno zastosować ciągu~\ref{thm:xviii-milnor}
automatycznie do $C_*(\varprojlim X_j)$, nawet gdy
$X_j$ są zwartymi wielościanami. Zwykle brak naturalnego
izomorfizmu $C_*(\varprojlim X_j)\cong\varprojlim C_*(X_j)$,
a odwzorowania na grupach łańcuchów nie muszą być surjektywne.
Jest to jedna z przyczyn, dla których homologia Čecha i homologia
singularna zachowują się odmiennie na przestrzeniach dzikich.
\end{uwaga}

\section{Wyższe przeszkody dla przekrojów i odwzorowań}
\begin{definicja}[Problem rozszerzania przekroju]
Niech $p:E\to B$ będzie fibracją Serre'a nad spójnym kompleksem CW,
niech $B^{(q)}$ oznacza jego $q$-szkielet, a $s_q:B^{(q)}\to E$
ciągły przekrój nad $B^{(q)}$. Problem polega na rozszerzeniu
$s_q$ na $B^{(q+1)}$, a ostatecznie na całe $B$.
Zakładamy, że włókno $F$ jest łukowo spójne; przy formułowaniu
przeszkód z abelowymi współczynnikami zakładamy też $q\ge2$.
\end{definicja}
\begin{stwierdzenie}[Kołańcuch przeszkody]
\label{prop:xviii-obstruction}
Niech $p:E\to B$ będzie fibracją Serre'a nad spójnym
kompleksem CW, której włókno $F$ jest $(q-1)$-spójne,
gdzie $q\ge2$. Niech $s_q$ będzie przekrojem nad $B^{(q)}$.
Dla każdej zorientowanej komórki $e^{q+1}$ wybierzmy
jej charakterystykę $\chi_e:D^{q+1}\to B$ i trywializację
fibracji nad dyskiem w sensie równoważności włóknowej
up-to-homotopy, ustaloną na brzegu. Klasa ograniczenia
$s_q\circ\chi_e|_{S^q}$ określa element $\pi_q(F)$.
Po uwzględnieniu transportu wzdłuż $1$-szkieletu otrzymujemy
kołańcuch komórkowy
$o_{q+1}(s_q)\in C^{q+1}(B;\underline{\pi_q(F)})$.
Jest on kocyklem, a jego klasa w
$H^{q+1}(B;\underline{\pi_q(F)})$ jest pierwszą przeszkodą
rozszerzenia przekroju. Jeśli ta klasa znika, można zmienić
przekrój na $q$-szkielecie i rozszerzyć go na $B^{(q+1)}$.
\end{stwierdzenie}
\begin{proof}[Konstrukcja i kontrola zależności]
Na każdej komórce dyskowej pullback fibracji jest
włóknowo homotopijnie trywialny: stosujemy własność
podnoszenia homotopii do ściągnięcia dysku do punktu.
Dzięki $(q-1)$-spójności $F$ wszystkie rozszerzenia do
$q$-szkieletu można porównywać, a element otrzymany na
brzegu komórki jest dobrze określony po wybraniu transportów
między włóknami. Monodromia $\pi_1(B)$ może działać
nietrywialnie na $\pi_q(F)$; właśnie to działanie definiuje
system lokalnych współczynników podkreślony w tezie.

Warunek kocyklu można sprawdzić na jednej komórce
$e^{q+2}$. Jej przyczepienie daje sumę granicznych map na
$q$-sferach ścian, przemnożonych przez odpowiednie
homomorfizmy transportu. Suma tych klas jest zerowa,
ponieważ jest obrazem brzegu mapy dysku po przejściu do
łańcuchów komórkowych i wykorzystaniu zgodności
przekształceń homotopii na przecięciach. Ten argument
jest standardowym \emph{lematem komórkowym o kocyklu
przeszkody}; jego pełna konstrukcja wymaga ułożenia
homotopii charakterystyk na $(q+2)$-komórkach.

Jeżeli zmienimy $s_q$ na $q$-komórkach przez klasy
$\beta\in C^q(B;\underline{\pi_q(F)})$, kołańcuch
przeszkody zmienia się o $\delta\beta$ (z odpowiednią
konwencją znaków). Zatem klasa kohomologii jest
niezależna od tej zmiany. Jej znikanie umożliwia wybór
$\beta$ zabijający kołańcuch, po czym każda mapa
$S^q\to F$ z trywialną klasą homotopii rozszerza się
na $D^{q+1}$. Z własności podnoszenia otrzymujemy
przekrój nad każdym dołączonym dyskiem i ich zgodność
na brzegu. To kończy argument pod warunkiem przyjęcia
wskazanego lematu komórkowego, którego nie zastępujemy
samym sloganem $\partial^2=0$.
\end{proof}
\begin{uwaga}[Kolejne przeszkody]
Po zniwelowaniu pierwszej klasy można przejść do następnego
szkieletu. Następna przeszkoda leży w grupie z wyższymi
homotopiami włókna, ale \emph{zależy od dokonanego wyboru
wcześniejszego rozszerzenia}. Nie ma na ogół jednej
kanonicznej listy klas niezależnych od siebie.
Ponadto przy $q=1$ współczynniki mogą być nieabelowe;
wtedy klasyczny zapis grupy kohomologii wymaga odrębnych
konstrukcji. W przypadku $q\ge2$ skręcenie współczynników
jest wyznaczone przez akcję monodromii na $\pi_q(F)$.
\end{uwaga}
\begin{przyklad}[Wiązka sferowa i klasa Eulera]
Dla zorientowanej rzeczywistej wiązki wektorowej $\xi$ rzędu
$r\ge3$ nad spójnym CW-kompleksem $B$ włóknem wiązki
jednostkowych sfer jest $S^{r-1}$, które jest $(r-2)$-spójne.
Pierwsza przeszkoda istnienia przekroju nigdzie niezerowego
leży w $H^r(B;\Z)$ i jest klasą Eulera $e(\xi)$.
Jeżeli $\dim B=r$, jej znikanie jest także wystarczające,
gdyż nie ma komórek wyższych stopni. Jeżeli
$\dim B>r$, mogą występować wyższe przeszkody, czego
nie wykrywa sama klasa Eulera.
\end{przyklad}
\begin{przyklad}[Zerowa klasa Eulera nie wystarcza]
Dla zorientowanej wiązki rzeczywistej rzędu trzy nad
$S^4$ klasa Eulera należy do $H^3(S^4;\Z)=0$.
Mimo to istnieją takie wiązki bez przekroju nigdzie
niezerowego: nietrywialne elementy
$\pi_3(SO(3))\cong\Z$ wyznaczają przez sklejanie
wiązki rzędu trzy. Przekrój wymuszałby redukcję
struktury do $SO(2)$, a wtedy klasa sklejenia pochodziłaby
z $\pi_3(SO(2))=0$.
Przykład pokazuje rzeczywistą wyższą przeszkodę.
\end{przyklad}

\section{Iloczyny Steenroda i relacje Adema: zakres elementarny}
Na kohomologii z $\F_2$ działają naturalne operacje
$\operatorname{Sq}^i:H^n(X;\F_2)\to H^{n+i}(X;\F_2)$.
Przypominamy normalizację:
$\operatorname{Sq}^0=\id$,
$\operatorname{Sq}^n(x)=x\smile x$ dla $x\in H^n$,
$\operatorname{Sq}^i(x)=0$ dla $i>n$.
Do tego dochodzą naturalność i wzór Cartana, już omówione
we wcześniejszych rozdziałach.
\begin{twierdzenie}[Relacje Adema, $p=2$]
\label{thm:xviii-adem}
Dla liczb całkowitych $0<a<2b$ obowiązuje równość
operacji naturalnych
\[
\operatorname{Sq}^a\operatorname{Sq}^b=
\sum_{j=0}^{\lfloor a/2\rfloor}
\binom{b-j-1}{a-2j}
\operatorname{Sq}^{a+b-j}\operatorname{Sq}^{j},
\]
przy czym współczynniki dwumianowe liczymy modulo dwa.
Wzór jest twierdzeniem o operacjach na \emph{wszystkich}
przestrzeniach, nie tylko o ich wartościach w jednym pierścieniu
kohomologii.
\end{twierdzenie}
\begin{uwaga}[Status dowodu]
Pełny dowód wymaga kołańcuchowej homotopii relacji Adema
lub uniwersalnej konstrukcji operacji i obliczeń w
$H^*(K(\F_2,n);\F_2)$. Nie wynika on wyłącznie z
naturalności i wzoru Cartana. W tym rozdziale przedstawiamy
jedynie samodzielne wyprowadzenia kilku konsekwencji i
kontrolę obliczeniową; twierdzenie~\ref{thm:xviii-adem}
przyjmujemy jako zewnętrzny wynik.
\end{uwaga}
\begin{stwierdzenie}[Pierwsze relacje]
Z relacji Adema wynikają
$\operatorname{Sq}^1\operatorname{Sq}^1=0$,
$\operatorname{Sq}^1\operatorname{Sq}^2=\operatorname{Sq}^3$,
$\operatorname{Sq}^2\operatorname{Sq}^2=
\operatorname{Sq}^3\operatorname{Sq}^1$.
\end{stwierdzenie}
\begin{proof}
Podstawiamy w\,~\ref{thm:xviii-adem} odpowiednio
$(a,b)=(1,1),(1,2),(2,2)$. W ostatnim przypadku
składnik $j=0$ ma współczynnik
$\binom{1}{2}=0$, a składnik $j=1$ ma
$\binom{0}{0}=1$. Dla $(1,1)$ jedyny współczynnik to
$\binom{0}{1}=0$; dla $(1,2)$ to $\binom{1}{1}=1$.
\end{proof}
\begin{przyklad}[Obliczenie na $\mathbb RP^{\infty}$]
Dla generatora $t\in H^1(\mathbb RP^\infty;\F_2)$ mamy
$\operatorname{Sq}^i(t^m)=\binom mi t^{m+i}$.
Wzór wynika z $\operatorname{Sq}(t)=t+t^2$ i
wzoru Cartana: $\operatorname{Sq}(t^m)=(t+t^2)^m$.
Na przykład
$\operatorname{Sq}^2\operatorname{Sq}^2(t^3)=0$,
a $\operatorname{Sq}^3\operatorname{Sq}^1(t^3)=0$.
Jest to zgodność z relacją, ale samo obliczenie na tej
przestrzeni \emph{nie dowodzi} relacji dla wszystkich $X$.
\end{przyklad}
\begin{przyklad}[Przypadek $\mathbb CP^\infty$]
Jeśli $u\in H^2(\mathbb CP^\infty;\F_2)$ jest redukcją
pierwszej klasy Cherna uniwersalnej wiązki liniowej, to
$\operatorname{Sq}^2(u)=u^2$ i
$\operatorname{Sq}^{2i}(u^m)=\binom mi u^{m+i}$,
podczas gdy nieparzyste kwadraty znikają na całym
$\F_2[u]$ z przyczyn parzystości stopni.
\end{przyklad}

\section{Ekscyzja homotopijna Blakersa--Masseya}
\begin{definicja}[Spójność pary]
Para przestrzeni z punktem bazowym $(X,A)$ jest $m$-spójna,
$m\ge1$, jeżeli $A$ i $X$ są łukowo spójne i
$\pi_j(X,A)=0$ dla wszystkich $1\le j\le m$.
Przyjmiemy typowe hipotezy CW, aby grupy względne i
podnoszenie homotopii miały oczekiwane własności.
\end{definicja}
\begin{twierdzenie}[Blakers--Massey, postać ekscyzyjna]
\label{thm:xviii-bm}
Niech $X=A\cup B$ będzie triadą CW, a $C=A\cap B$ niech
będzie niepustym łukowo spójnym podkompleksem.
Załóżmy, że $(A,C)$ jest $m$-spójna, $(B,C)$ jest
$n$-spójna, gdzie $m,n\ge1$.
Wówczas homomorfizm ekscyzji
\[
\pi_q(A,C)\longrightarrow\pi_q(X,B)
\]
jest izomorfizmem dla $q<m+n$ i epimorfizmem dla $q=m+n$.
W stopniu $q=1$ rozumiemy odpowiednie zbiory punktowane
lub grupy względne zgodnie ze zwykłą konwencją.
\end{twierdzenie}
\begin{uwaga}[Zakres dowodu i źródło trudności]
Dla homologii ekscyzja jest izomorfizmem bez ograniczenia
stopnia. W homotopii zakres zależy od spójności obu par.
Dowód Blakersa--Masseya polega na metodzie \emph{usuwania
przecięć}: reprezentant odwzorowania dysku do $X$ rozcina
się na obszary przechodzące do $A$ i $B$, a założenia
spójności pozwalają sukcesywnie zmieniać reprezentanta,
usuwając obszary niewłaściwe w zakresie wymiarów
$q\le m+n$. Ten geometryczny argument wymaga
technicznego lematu o uporządkowaniu przecinających się
komórek; nie uzasadnia go sama ekscyzja homologiczna.
Pełny lemat geometryczny przyjmujemy jako twierdzenie
zewnętrzne. Poniżej dowodzimy wynikający z niego
zakres stabilizacji Freudenthala.
\end{uwaga}
\begin{wniosek}[Stabilizacja Freudenthala]
\label{cor:xviii-freudenthal}
Niech $Y$ będzie $(r-1)$-spójnym, łukowo spójnym
kompleksem CW, gdzie $r\ge2$. Wówczas homomorfizm
zawieszenia
$\pi_k(Y)\to\pi_{k+1}(\Sigma Y)$ jest izomorfizmem
dla $k\le2r-2$ i epimorfizmem dla $k=2r-1$.
\end{wniosek}
\begin{proof}
Rozłóżmy $\Sigma Y$ na dwa stożki $A=CY$ i $B=CY$,
których przecięcie jest homotopijnie równoważne $Y$.
Po odpowiednim zastąpieniu przez triadę CW
obie pary $(CY,Y)$ są $r$-spójne, ponieważ
$\pi_j(CY,Y)\cong\pi_{j-1}(Y)$ dla $j\ge2$.
Twierdzenie~\ref{thm:xviii-bm} daje izomorfizm
$\pi_q(CY,Y)\to\pi_q(\Sigma Y,CY)$ dla $q<2r$
i epimorfizm dla $q=2r$.
Długie ciągi dokładne par oraz ściągalność stożków
dają $\pi_q(CY,Y)\cong\pi_{q-1}(Y)$ i
$\pi_q(\Sigma Y,CY)\cong\pi_q(\Sigma Y)$ dla $q\ge2$.
Przy tych identyfikacjach mapa ekscyzji odpowiada
homomorfizmowi zawieszenia, co daje wynik po podstawieniu
$q=k+1$.
\end{proof}
\begin{przyklad}[Zawieszenie okręgu]
Dla $Y=S^1$ powyższy wniosek przy $r=1$ nie ma zastosowania
w tej postaci; można jednak bezpośrednio obliczyć
$\Sigma S^1\cong S^2$ oraz
$\pi_1(S^1)\cong\Z$, $\pi_2(S^2)\cong\Z$.
Ważne jest niewnioskowanie ogólnego zakresu stabilizacji
z jednego przykładu.
\end{przyklad}

\section{Sklejanie trajektorii Morse'a: ścisły bilans analityczny}
Niech $M$ będzie zwartą gładką rozmaitością bez brzegu,
$f:M\to\R$ funkcją Morse'a, a $g$ metryką Riemannowską
taką, że para $(f,g)$ spełnia warunek Morse'a--Smale'a.
Dla punktów krytycznych $p,q$ oznaczmy przez
$\mathcal M(p,q)$ przestrzeń trajektorii równania
$\dot\gamma=-\nabla_g f(\gamma)$, modulo przesunięcia czasu.
Przyjmujemy $\operatorname{ind}(p)-\operatorname{ind}(q)=2$.
Chcemy opisać końce jednowymiarowej przestrzeni
$\mathcal M(p,q)$ jako złamane trajektorie przechodzące
przez punkty o indeksie pośrednim.

\begin{definicja}[Operator linearyzacji i norma ważona]
Niech $\gamma:\R\to M$ będzie trajektorią gradientową.
Po wybraniu trywializacji $\gamma^*TM$ definiujemy
$D_\gamma:W^{1,2}(\R,\R^d)\to L^2(\R,\R^d)$ przez
\[
D_\gamma\xi=\frac{\dd\xi}{\dd s}
+\nabla_\xi\nabla_g f(\gamma(s)).
\]
Dla $\delta>0$ i wybranej gładkiej wagi
$\rho_\delta(s)$ zachowującej się jak
$e^{\delta|s|}$ na końcach definiujemy
$W^{1,2}_\delta$ i $L^2_\delta$ przez mnożenie przez
$\rho_\delta$. Wartość $\delta$ wybieramy mniejszą od
modułów niezerowych wartości własnych hesjanów na końcach.
\end{definicja}
\begin{uwaga}[Symetria przesunięcia]
Wektor $\dot\gamma$ jest w jądrze $D_\gamma$.
Dlatego nawet przy transwersalności nie oczekujemy
odwrotności $D_\gamma$ na całej przestrzeni.
Wprowadza się warunek cechujący, usuwający kierunek
przesunięcia czasu, oraz odpowiedni wybór dopełnień.
\end{uwaga}
\begin{lemat}[Oszacowanie błędu trajektorii sklejonej]
\label{lem:xviii-preglue}
Niech $\gamma_-\in\mathcal M(p,r)$ i
$\gamma_+\in\mathcal M(r,q)$, a ich zbieżności do $r$
niech będą wykładnicze. Dla dostatecznie dużego parametru
$R$ można skonstruować krzywą przybliżoną $\gamma_R$
powstałą przez obcięcie i sklejenie obu trajektorii,
dla której residuum
$\mathcal F(\gamma_R):=\dot\gamma_R+\nabla f(\gamma_R)$
spełnia $\|\mathcal F(\gamma_R)\|_{L^2_\delta}
\le C e^{-cR}$ z pewnymi stałymi $C,c>0$.
\end{lemat}
\begin{proof}
W małym otoczeniu $r$ zapisujemy fragmenty trajektorii
w ustalonej karcie wykładniczej i przesuwamy ich parametry
o $R$. Poza przedziałami, gdzie zmienia się funkcja
odcinająca, krzywa $\gamma_R$ jest dokładną trajektorią.
Na obszarze przejścia pochodna funkcji odcinającej jest
jednostajnie ograniczona, a oba sklejane fragmenty
leżą w odległości $O(e^{-\mu R})$ od $r$, dla pewnego
$\mu>0$. Gładkość $\nabla f$ daje tę samą kontrolę
na błąd nieliniowego członu. Po pomnożeniu przez
wagę $e^{\delta|s|}$ z $\delta<\mu$ otrzymujemy
oszacowanie $O(e^{-(\mu-\delta)R})$; stałe nie zależą
od $R$ przy ustalonych końcowych trajektoriach.
\end{proof}
\begin{lemat}[Kluczowy krok Fredholmowski --- lemat przyjmowany]
\label{lem:xviii-uniform-right-inverse}
W sytuacji lematu~\ref{lem:xviii-preglue}, przy odpowiednim
wyborze warunków cechujących i przy założeniu
transwersalności Morse'a--Smale'a, operator linearyzacji
$D_R=D\mathcal F(\gamma_R)$ posiada prawy odwrotny
$Q_R$ na wybranym dopełnieniu jądra taki, że
$\|Q_R\|\le C$ dla wszystkich dostatecznie dużych $R$.
\end{lemat}
\begin{uwaga}
Ten lemat nie wynika z samej małości residuum.
Aby go udowodnić, skleja się lokalne prawe odwrotne,
kontroluje błąd komutatora z funkcjami odcinającymi i
stosuje argument wykluczający koncentrację przy długiej szyjce.
Wymaga to pełnej analizy operatorów Fredholma w normach
ważonych. W tej iteracji formułujemy go precyzyjnie,
lecz \emph{nie przedstawiamy jako udowodnionego}.
\end{uwaga}
\begin{twierdzenie}[Wniosek o sklejaniu, zależny od lematu]
\label{thm:xviii-gluing-conditional}
Przy założeniach powyżej i po przyjęciu
lematu~\ref{lem:xviii-uniform-right-inverse} dla dostatecznie
dużych $R$ istnieje w pobliżu $\gamma_R$ dokładna
trajektoria $\widetilde\gamma_R$ równania gradientowego.
Jej poprawka jest rzędu $O(e^{-cR})$ w odpowiedniej normie
ważonej. Przy dodatkowym twierdzeniu o lokalnej jednoznaczności
i zgodności orientacji konstrukcja daje współrzędną końca
skompaktyfikowanej przestrzeni trajektorii.
\end{twierdzenie}
\begin{proof}
W ustalonej karcie Banacha poszukujemy poprawki $\xi$
spełniającej $\mathcal F(\gamma_R+\xi)=0$. Rozpisujemy
\[
\mathcal F(\gamma_R+\xi)=e_R+D_R\xi+N_R(\xi),
\qquad e_R=\mathcal F(\gamma_R),
\]
gdzie $N_R(0)=0$, $DN_R(0)=0$, a gładkość gradientu
na zwartej rozmaitości daje, dla małych $\xi,\eta$,
$\|N_R(\xi)-N_R(\eta)\|
\le K(\|\xi\|+\|\eta\|)\|\xi-\eta\|$.
Po wyborze dopełnienia jądra i podstawieniu
$\xi=Q_Rv$ równanie sprowadza się do punktu stałego
$\xi=-Q_R(e_R+N_R(\xi))$.
Lemat zapewnia $\|Q_R\|\le C$, a poprzedni lemat
$\|e_R\|\le C_0e^{-cR}$. Dla promienia
$2CC_0e^{-cR}$ prawa strona zachowuje kulę i jest
kontrakcją, gdy $R$ jest duże. Zatem istnieje poprawka
$\xi_R$ tej wielkości. Aby wniosek dał pełny opis
brzegu przestrzeni moduli, trzeba osobno wykazać, że
wszystkie bliskie trajektorie powstają taką metodą i że
orientacje są zgodne. Tych ostatnich twierdzeń w tym
dowodzie nie przyjmujemy za automatyczne.
\end{proof}

\section{Podsumowanie zakresu i zależności dowodowych}
W tym rozdziale bez zewnętrznych głębokich twierdzeń
wyprowadziliśmy definicje $\varprojlim^1$, kryterium
Mittaga--Lefflera w sytuacji przeliczalnych układów,
przykład $\Z_p/\Z$ i ciąg Milnora dla stopniowo surjektywnych
układów kompleksów. Omówiliśmy konstrukcję
wyższych przeszkód, podkreślając zależność od lokalnych
współczynników i wyboru częściowych rozszerzeń.
Dla operacji Steenroda rachunki opierają się na
przyjętym twierdzeniu Adema; dla Freudenthala na
przyjętym Blakersie--Masseyu.
W części Morse'a udowodniliśmy oszacowanie residuum i
krok kontrakcyjny warunkowo. \emph{Jednostajny prawy odwrotny,
lokalna pełność parametrów sklejania i znaki orientacji}
pozostają precyzyjnie wskazanymi tematami do dalszego dowodu.

\appendix
\renewcommand{\chaptermark}[1]{\markboth{Dodatek \thechapter}{}}
\chapter{Minimalne przypomnienie geometrii różniczkowej}
\label{ch:dodatek-geometria}

W części poświęconej kohomologii de Rhama, dualności Poincarégo i klasom
charakterystycznym potrzebujemy pojęć geometrycznych. Ten dodatek
nie zastępuje dowodów z książki o geometrii różniczkowej, ale wyjaśnia
terminy oraz precyzuje przyjmowane twierdzenia.

\section{Rozmaitości gładkie, przestrzeń styczna i formy}
\begin{definicja}
Gładka rozmaitość wymiaru $n$ jest rozmaitością topologiczną Hausdorffa,
spełniającą drugi aksjomat przeliczalności, z atlasem map do otwartych
podzbiorów $\R^n$, w którym zmiany współrzędnych są gładkie.
Przestrzeń styczna $T_pM$ w punkcie $p$ jest przestrzenią derywacji
$C^\infty(M)\to\R$ w punkcie $p$; przestrzeń kostyczna to dualna
$T_p^*M$. Formy różniczkowe stopnia $k$ są gładkimi przekrojami
wiązki $\bigwedge^kT^*M$.
\end{definicja}
\begin{stwierdzenie}
Jeżeli $V$ jest $n$-wymiarową przestrzenią wektorową z bazą
$e_1,\dots,e_n$, a $e^1,\dots,e^n$ jest bazą dualną,
to dla każdego $0\leq k\leq n$ formy
$e^{i_1}\wedge\cdots\wedge e^{i_k}$,
$1\leq i_1<\cdots<i_k\leq n$, tworzą bazę $\bigwedge^k V^*$.
\end{stwierdzenie}
\begin{proof}
Z naprzemienności każdy prosty iloczyn jest równy zeru lub
odpowiedniemu iloczynowi o ściśle rosnących indeksach z właściwym znakiem,
więc wskazane elementy rozpinają $\bigwedge^kV^*$.
Ich niezależność wynika przez wartościowanie na zestawach
$(e_{j_1},\dots,e_{j_k})$: dany iloczyn ma wartość $1$ dla
odpowiadającego zestawu i $0$ dla pozostałych.
\end{proof}

\section{Regularne podrozmaitości, transwersalność i otoczenia tubularne}
\begin{definicja}
Gładkie odwzorowanie $f:M\to N$ jest submersją w $p$, gdy
$\mathrm df_p:T_pM\to T_{f(p)}N$ jest surjekcją.
Podrozmaitości $A,B\subseteq M$ są transwersalne w $p\in A\cap B$,
jeśli $T_pA+T_pB=T_pM$.
Jeśli $S\subset M$ jest podrozmaitością osadzoną,
\emph{wiązką normalną} $\nu(S,M)$ nazywamy $TM|_S/TS$.
\end{definicja}
\begin{twierdzenie}[Zestaw twierdzeń wykorzystywanych dalej]
Jeżeli $f:M\to N$ jest submersją, to w odpowiednich lokalnych
współrzędnych ma postać projekcji $(x,y)\mapsto x$.
Jeżeli dwie gładkie podrozmaitości $A,B$ przecinają się transwersalnie,
to $A\cap B$ jest podrozmaitością o wymiarze $\dim A+\dim B-\dim M$.
Dla osadzonej podrozmaitości $S\subset M$ istnieje dyfeomorfizm
otoczenia przekroju zerowego w $\nu(S,M)$ na otoczenie $S$ w $M$.
Dla gładkich rozmaitości spełnione są standardowe twierdzenia
Sarda i transwersalności.
\end{twierdzenie}
\begin{uwaga}
Twierdzenia te są tu przyjmowane z wcześniejszej książki.
W szczególności użycie transwersalności przy liczeniu przecięć
wymaga zwartości lub odpowiednich warunków właściwości dla
zapewnienia skończoności zbioru punktów przecięcia.
\end{uwaga}

\begin{definicja}
\emph{Uramowaniem} rzeczywistej wiązki wektorowej rzędu $r$
nazywamy wybór $r$ ciągłych przekrojów, które w każdym włóknie
tworzą bazę. Istnienie uramowania jest równoważne trywialności wiązki.
\end{definicja}

\section{Gdzie kontynuować dalej?}
Po opanowaniu teorii zawartej w skrypcie naturalnymi dalszymi etapami są
silne homologie, dalsza teoria ciągów spektralnych
(w szczególności ciąg Adamsa), operacje wtórne,
obliczenia pierścieni kobordyzmu oraz bardziej zaawansowana teoria widm.
Jeszcze dalej leżą torsja Whiteheada, teoria chirurgii
i topologia czterowymiarowa; wymagają znacznie większego przygotowania.

\end{document}
